\section{Linear Convergence of Training Dynamics}
\label{sec:convergence}\label{sec:balanced-convergence}
\label{sec:sampling}\label{sec:finite-width}

\paragraph{Convergence of gradient flow}
A positive conserved support gap $g(\Sigma_0) > 0$ guarantees global linear convergence for the Bures flow \eqref{eq:bures-G}, including some cases where the initial covariance $\Sigma_0$ is singular. Write $\Delta_t:=\ell(\Svu_t)-\ell_\star$ for the loss gap and $\cS_\star$ for the set of global predictor minimizers.

\begin{theorem}[Global convergence from a positive support gap]\label{thm:gf}\label{thm:balanced-convergence}
Under Assumption~\ref{ass:loss}, let $\Sigma_0\succeq0$ be a possibly
singular covariance with $\gap{\Sigma_0}>0$. 
Then the solution of \eqref{eq:bures-G} is global and satisfies
\begin{equation}\label{eq:gf-rate}
 \Delta_t\le\Delta_0e^{-2\kappa \gap{\Sigma_0}t},\qquad t\ge0.
\end{equation}
As $t \to \infty$, $\Sigma_t \to \Sigma_\infty\succeq0$ with
$\rankop\Sigma_\infty=\rankop\Sigma_0$ and $\Svu_\infty\in\cS_\star$.
The predictor $\Svu$ and any factor realization (as per \eqref{eq:mean-field-gradient-flow}) converge linearly with rate $\kappa\gap{\Sigma_0}$ in Frobenius norm.
\end{theorem}

The only nondegeneracy requirement is $\gap{\Sigma_0}>0$.
Positive definiteness is sufficient but not necessary: $\Sigma_0\succ0$ implies $\gap{\Sigma_0}\ge2\lambda_{\min}(\Sigma_0)>0$ and then
$\Sigma_\infty\succ0$. The exponent separates loss geometry (captured by the PL constant $\kappa$)
from initial covariance conditioning (the support gap $\gap{\Sigma_0}$), with no extra
dimension factor. The trajectory selects a global minimizer according to its initialization and dynamics. Proposition~\ref{prop:gf-lift-tail} supplies all explicit prefactors and bounds.

\begin{comment}
\emph{Proof sketch.}
\emph{Conserved conditioning.}
Proposition~\ref{lem:margins} gives
$\Suu_t\succeq \aplus{\Sigma_0}\Id$ and $\Svv_t\succeq \aminus{\Sigma_0}\Id$ throughout the
existence interval. The predictor equation
$\dot\Svu_t=-E_t\Suu_t-\Svv_tE_t$ therefore has a uniformly positive
preconditioner, with lower bound $\gap{\Sigma_0}$.
Only the sum of the two margins must be positive; either one may vanish.

\emph{Loss decay.}
Pairing this equation with $E_t=\nabla\ell(\Svu_t)$ gives
$-\dot\Delta_t\ge \gap{\Sigma_0}\norm{E_t}_{\F}^2$.
The PL inequality then yields
$-\dot\Delta_t\ge2\kappa \gap{\Sigma_0}\Delta_t$, proving \eqref{eq:gf-rate}
on the existence interval.

\emph{Global existence and convergence of the state.}
The same dissipation estimate applied to $\sqrt{\Delta_t}$ bounds
$\int_0^T\norm{G_t}_{\op}\,\dd t$ by
$\gap{\Sigma_0}^{-1}\sqrt{2\Delta_0/\kappa}$, uniformly over finite existence
intervals. Hence the fundamental matrix and its inverse remain bounded.
Covariance congruence prevents blow-up, so the flow is global. The
fundamental matrix converges to an invertible limit, preserving covariance
rank. Applying the length bound from time $t$ onward gives
exponent $\kappa \gap{\Sigma_0}$ for the remaining displacement and hence convergence
of the covariance and factors.
Appendix~\ref{app:gf} gives the full proof.\qed
\end{comment}

\begin{remark}[Singular covariance and predictor rank]\label{rem:singular-support}
The initialization $\Sigma_0=\diag(\Id_{d_u},0_{d_v})$ has
$\aplus{\Sigma_0}=1$, $\aminus{\Sigma_0}=0$, and $\gap{\Sigma_0}=1$.
Its rank is only $d_u$, yet the theorem applies. There is no rank
contradiction: a full-rank predictor needs rank $\min(d_u,d_v)$, not $m$.
More generally, $\aplus{\Sigma_0}>0$ implies $\rankop\Sigma_0\ge d_u$,
while $\aminus{\Sigma_0}>0$ implies $\rankop\Sigma_0\ge d_v$.
The condition cannot be dropped altogether: for $d_u=d_v=2$,
$\Sigma_0=\diag(1,0,1,0)$ and
$\ell(S):=\norm{S-e_2e_2^\top}_{\F}^2/2$, the flow is stationary and
nonoptimal. Here $g(\Sigma_0) = 0$, despite the nonzero signed eigenvalues $+1,-1$.
\end{remark}

\begin{remark}[Quadratic specialization]
For $\ell(S):=\frac12\ip{\quadH[S-\Svu_\star]}{S-\Svu_\star}$ with
$\quadH\succeq0$ a nonzero self-adjoint linear operator, the exponent is
$2\gap{\Sigma_0}\lambda_{\min}^{+}(\quadH)$. For a full-rank isotropic initialization
$\Sigma_0=\sigma^2\Id_m$, one has $\aplus{\Sigma_0}=\aminus{\Sigma_0}=\sigma^2$,
giving a rate of $4\sigma^2\lambda_{\min}^{+}(\quadH)$.
The operator $\quadH$ is allowed to have a nontrivial kernel: the predictor $\Svu$ converges to a point in
$\Svu_\star+\kerop\quadH$. If $\quadH$ is positive definite, the predictor converges to the teacher $\Svu_\star$.
\end{remark}

For a quadratic predictor loss, Figure~\ref{fig:bures-convergence-triptych} compares the initialization scale, loss of conditioning, and the finite-width sampling we discuss next. The dotted curves use a coarse determinant-based rate, derived for comparison in Appendix~\ref{app:gf}.
If the support gap $g(\Sigma_0)$ vanishes, trace and initial loss alone cannot guarantee a uniform positive rate (Example~\ref{ex:slow}).
Figure~\ref{fig:smooth-lasso-triptych} in Appendix~\ref{app:smooth-lasso} provides the counterpart for a non-quadratic smoothed-Lasso predictor loss.

\begin{figure}[t]
    \centering
    \resizebox{\linewidth}{!}{%% Creator: Matplotlib, PGF backend
%%
%% To include the figure in your LaTeX document, write
%%   \input{<filename>.pgf}
%%
%% Make sure the required packages are loaded in your preamble
%%   \usepackage{pgf}
%%
%% Also ensure that all the required font packages are loaded; for instance,
%% the lmodern package is sometimes necessary when using math font.
%%   \usepackage{lmodern}
%%
%% Figures using additional raster images can only be included by \input if
%% they are in the same directory as the main LaTeX file. For loading figures
%% from other directories you can use the `import` package
%%   \usepackage{import}
%%
%% and then include the figures with
%%   \import{<path to file>}{<filename>.pgf}
%%
%% Matplotlib used the following preamble
%%   \def\mathdefault#1{#1}
%%   \everymath=\expandafter{\the\everymath\displaystyle}
%%   \IfFileExists{scrextend.sty}{
%%     \usepackage[fontsize=6.000000pt]{scrextend}
%%   }{
%%     \renewcommand{\normalsize}{\fontsize{6.000000}{7.200000}\selectfont}
%%     \normalsize
%%   }
%%   \usepackage{amsfonts}
%%   \usepackage{amssymb}
%%   \ifdefined\pdftexversion\else  % non-pdftex case.
%%     \usepackage{fontspec}
%%   \fi
%%   \makeatletter\@ifpackageloaded{underscore}{}{\usepackage[strings]{underscore}}\makeatother
%%
\begingroup%
\makeatletter%
\begin{pgfpicture}%
\pgfpathrectangle{\pgfpointorigin}{\pgfqpoint{6.500000in}{1.687500in}}%
\pgfusepath{use as bounding box, clip}%
\begin{pgfscope}%
\pgfsetbuttcap%
\pgfsetmiterjoin%
\definecolor{currentfill}{rgb}{1.000000,1.000000,1.000000}%
\pgfsetfillcolor{currentfill}%
\pgfsetlinewidth{0.000000pt}%
\definecolor{currentstroke}{rgb}{1.000000,1.000000,1.000000}%
\pgfsetstrokecolor{currentstroke}%
\pgfsetdash{}{0pt}%
\pgfpathmoveto{\pgfqpoint{0.000000in}{0.000000in}}%
\pgfpathlineto{\pgfqpoint{6.500000in}{0.000000in}}%
\pgfpathlineto{\pgfqpoint{6.500000in}{1.687500in}}%
\pgfpathlineto{\pgfqpoint{0.000000in}{1.687500in}}%
\pgfpathlineto{\pgfqpoint{0.000000in}{0.000000in}}%
\pgfpathclose%
\pgfusepath{fill}%
\end{pgfscope}%
\begin{pgfscope}%
\pgfsetbuttcap%
\pgfsetmiterjoin%
\definecolor{currentfill}{rgb}{1.000000,1.000000,1.000000}%
\pgfsetfillcolor{currentfill}%
\pgfsetlinewidth{0.000000pt}%
\definecolor{currentstroke}{rgb}{0.000000,0.000000,0.000000}%
\pgfsetstrokecolor{currentstroke}%
\pgfsetstrokeopacity{0.000000}%
\pgfsetdash{}{0pt}%
\pgfpathmoveto{\pgfqpoint{0.487500in}{0.320625in}}%
\pgfpathlineto{\pgfqpoint{2.333179in}{0.320625in}}%
\pgfpathlineto{\pgfqpoint{2.333179in}{1.485000in}}%
\pgfpathlineto{\pgfqpoint{0.487500in}{1.485000in}}%
\pgfpathlineto{\pgfqpoint{0.487500in}{0.320625in}}%
\pgfpathclose%
\pgfusepath{fill}%
\end{pgfscope}%
\begin{pgfscope}%
\pgfpathrectangle{\pgfqpoint{0.487500in}{0.320625in}}{\pgfqpoint{1.845679in}{1.164375in}}%
\pgfusepath{clip}%
\pgfsetrectcap%
\pgfsetroundjoin%
\pgfsetlinewidth{0.803000pt}%
\definecolor{currentstroke}{rgb}{0.909804,0.909804,0.909804}%
\pgfsetstrokecolor{currentstroke}%
\pgfsetdash{}{0pt}%
\pgfpathmoveto{\pgfqpoint{0.571395in}{0.320625in}}%
\pgfpathlineto{\pgfqpoint{0.571395in}{1.485000in}}%
\pgfusepath{stroke}%
\end{pgfscope}%
\begin{pgfscope}%
\pgfsetbuttcap%
\pgfsetroundjoin%
\definecolor{currentfill}{rgb}{0.000000,0.000000,0.000000}%
\pgfsetfillcolor{currentfill}%
\pgfsetlinewidth{0.803000pt}%
\definecolor{currentstroke}{rgb}{0.000000,0.000000,0.000000}%
\pgfsetstrokecolor{currentstroke}%
\pgfsetdash{}{0pt}%
\pgfsys@defobject{currentmarker}{\pgfqpoint{0.000000in}{-0.048611in}}{\pgfqpoint{0.000000in}{0.000000in}}{%
\pgfpathmoveto{\pgfqpoint{0.000000in}{0.000000in}}%
\pgfpathlineto{\pgfqpoint{0.000000in}{-0.048611in}}%
\pgfusepath{stroke,fill}%
}%
\begin{pgfscope}%
\pgfsys@transformshift{0.571395in}{0.320625in}%
\pgfsys@useobject{currentmarker}{}%
\end{pgfscope}%
\end{pgfscope}%
\begin{pgfscope}%
\definecolor{textcolor}{rgb}{0.000000,0.000000,0.000000}%
\pgfsetstrokecolor{textcolor}%
\pgfsetfillcolor{textcolor}%
\pgftext[x=0.571395in,y=0.223403in,,top]{\color{textcolor}{\rmfamily\fontsize{6.000000}{7.200000}\selectfont\catcode`\^=\active\def^{\ifmmode\sp\else\^{}\fi}\catcode`\%=\active\def%{\%}0}}%
\end{pgfscope}%
\begin{pgfscope}%
\pgfpathrectangle{\pgfqpoint{0.487500in}{0.320625in}}{\pgfqpoint{1.845679in}{1.164375in}}%
\pgfusepath{clip}%
\pgfsetrectcap%
\pgfsetroundjoin%
\pgfsetlinewidth{0.803000pt}%
\definecolor{currentstroke}{rgb}{0.909804,0.909804,0.909804}%
\pgfsetstrokecolor{currentstroke}%
\pgfsetdash{}{0pt}%
\pgfpathmoveto{\pgfqpoint{0.990867in}{0.320625in}}%
\pgfpathlineto{\pgfqpoint{0.990867in}{1.485000in}}%
\pgfusepath{stroke}%
\end{pgfscope}%
\begin{pgfscope}%
\pgfsetbuttcap%
\pgfsetroundjoin%
\definecolor{currentfill}{rgb}{0.000000,0.000000,0.000000}%
\pgfsetfillcolor{currentfill}%
\pgfsetlinewidth{0.803000pt}%
\definecolor{currentstroke}{rgb}{0.000000,0.000000,0.000000}%
\pgfsetstrokecolor{currentstroke}%
\pgfsetdash{}{0pt}%
\pgfsys@defobject{currentmarker}{\pgfqpoint{0.000000in}{-0.048611in}}{\pgfqpoint{0.000000in}{0.000000in}}{%
\pgfpathmoveto{\pgfqpoint{0.000000in}{0.000000in}}%
\pgfpathlineto{\pgfqpoint{0.000000in}{-0.048611in}}%
\pgfusepath{stroke,fill}%
}%
\begin{pgfscope}%
\pgfsys@transformshift{0.990867in}{0.320625in}%
\pgfsys@useobject{currentmarker}{}%
\end{pgfscope}%
\end{pgfscope}%
\begin{pgfscope}%
\definecolor{textcolor}{rgb}{0.000000,0.000000,0.000000}%
\pgfsetstrokecolor{textcolor}%
\pgfsetfillcolor{textcolor}%
\pgftext[x=0.990867in,y=0.223403in,,top]{\color{textcolor}{\rmfamily\fontsize{6.000000}{7.200000}\selectfont\catcode`\^=\active\def^{\ifmmode\sp\else\^{}\fi}\catcode`\%=\active\def%{\%}10}}%
\end{pgfscope}%
\begin{pgfscope}%
\pgfpathrectangle{\pgfqpoint{0.487500in}{0.320625in}}{\pgfqpoint{1.845679in}{1.164375in}}%
\pgfusepath{clip}%
\pgfsetrectcap%
\pgfsetroundjoin%
\pgfsetlinewidth{0.803000pt}%
\definecolor{currentstroke}{rgb}{0.909804,0.909804,0.909804}%
\pgfsetstrokecolor{currentstroke}%
\pgfsetdash{}{0pt}%
\pgfpathmoveto{\pgfqpoint{1.410340in}{0.320625in}}%
\pgfpathlineto{\pgfqpoint{1.410340in}{1.485000in}}%
\pgfusepath{stroke}%
\end{pgfscope}%
\begin{pgfscope}%
\pgfsetbuttcap%
\pgfsetroundjoin%
\definecolor{currentfill}{rgb}{0.000000,0.000000,0.000000}%
\pgfsetfillcolor{currentfill}%
\pgfsetlinewidth{0.803000pt}%
\definecolor{currentstroke}{rgb}{0.000000,0.000000,0.000000}%
\pgfsetstrokecolor{currentstroke}%
\pgfsetdash{}{0pt}%
\pgfsys@defobject{currentmarker}{\pgfqpoint{0.000000in}{-0.048611in}}{\pgfqpoint{0.000000in}{0.000000in}}{%
\pgfpathmoveto{\pgfqpoint{0.000000in}{0.000000in}}%
\pgfpathlineto{\pgfqpoint{0.000000in}{-0.048611in}}%
\pgfusepath{stroke,fill}%
}%
\begin{pgfscope}%
\pgfsys@transformshift{1.410340in}{0.320625in}%
\pgfsys@useobject{currentmarker}{}%
\end{pgfscope}%
\end{pgfscope}%
\begin{pgfscope}%
\definecolor{textcolor}{rgb}{0.000000,0.000000,0.000000}%
\pgfsetstrokecolor{textcolor}%
\pgfsetfillcolor{textcolor}%
\pgftext[x=1.410340in,y=0.223403in,,top]{\color{textcolor}{\rmfamily\fontsize{6.000000}{7.200000}\selectfont\catcode`\^=\active\def^{\ifmmode\sp\else\^{}\fi}\catcode`\%=\active\def%{\%}20}}%
\end{pgfscope}%
\begin{pgfscope}%
\pgfpathrectangle{\pgfqpoint{0.487500in}{0.320625in}}{\pgfqpoint{1.845679in}{1.164375in}}%
\pgfusepath{clip}%
\pgfsetrectcap%
\pgfsetroundjoin%
\pgfsetlinewidth{0.803000pt}%
\definecolor{currentstroke}{rgb}{0.909804,0.909804,0.909804}%
\pgfsetstrokecolor{currentstroke}%
\pgfsetdash{}{0pt}%
\pgfpathmoveto{\pgfqpoint{1.829812in}{0.320625in}}%
\pgfpathlineto{\pgfqpoint{1.829812in}{1.485000in}}%
\pgfusepath{stroke}%
\end{pgfscope}%
\begin{pgfscope}%
\pgfsetbuttcap%
\pgfsetroundjoin%
\definecolor{currentfill}{rgb}{0.000000,0.000000,0.000000}%
\pgfsetfillcolor{currentfill}%
\pgfsetlinewidth{0.803000pt}%
\definecolor{currentstroke}{rgb}{0.000000,0.000000,0.000000}%
\pgfsetstrokecolor{currentstroke}%
\pgfsetdash{}{0pt}%
\pgfsys@defobject{currentmarker}{\pgfqpoint{0.000000in}{-0.048611in}}{\pgfqpoint{0.000000in}{0.000000in}}{%
\pgfpathmoveto{\pgfqpoint{0.000000in}{0.000000in}}%
\pgfpathlineto{\pgfqpoint{0.000000in}{-0.048611in}}%
\pgfusepath{stroke,fill}%
}%
\begin{pgfscope}%
\pgfsys@transformshift{1.829812in}{0.320625in}%
\pgfsys@useobject{currentmarker}{}%
\end{pgfscope}%
\end{pgfscope}%
\begin{pgfscope}%
\definecolor{textcolor}{rgb}{0.000000,0.000000,0.000000}%
\pgfsetstrokecolor{textcolor}%
\pgfsetfillcolor{textcolor}%
\pgftext[x=1.829812in,y=0.223403in,,top]{\color{textcolor}{\rmfamily\fontsize{6.000000}{7.200000}\selectfont\catcode`\^=\active\def^{\ifmmode\sp\else\^{}\fi}\catcode`\%=\active\def%{\%}30}}%
\end{pgfscope}%
\begin{pgfscope}%
\pgfpathrectangle{\pgfqpoint{0.487500in}{0.320625in}}{\pgfqpoint{1.845679in}{1.164375in}}%
\pgfusepath{clip}%
\pgfsetrectcap%
\pgfsetroundjoin%
\pgfsetlinewidth{0.803000pt}%
\definecolor{currentstroke}{rgb}{0.909804,0.909804,0.909804}%
\pgfsetstrokecolor{currentstroke}%
\pgfsetdash{}{0pt}%
\pgfpathmoveto{\pgfqpoint{2.249285in}{0.320625in}}%
\pgfpathlineto{\pgfqpoint{2.249285in}{1.485000in}}%
\pgfusepath{stroke}%
\end{pgfscope}%
\begin{pgfscope}%
\pgfsetbuttcap%
\pgfsetroundjoin%
\definecolor{currentfill}{rgb}{0.000000,0.000000,0.000000}%
\pgfsetfillcolor{currentfill}%
\pgfsetlinewidth{0.803000pt}%
\definecolor{currentstroke}{rgb}{0.000000,0.000000,0.000000}%
\pgfsetstrokecolor{currentstroke}%
\pgfsetdash{}{0pt}%
\pgfsys@defobject{currentmarker}{\pgfqpoint{0.000000in}{-0.048611in}}{\pgfqpoint{0.000000in}{0.000000in}}{%
\pgfpathmoveto{\pgfqpoint{0.000000in}{0.000000in}}%
\pgfpathlineto{\pgfqpoint{0.000000in}{-0.048611in}}%
\pgfusepath{stroke,fill}%
}%
\begin{pgfscope}%
\pgfsys@transformshift{2.249285in}{0.320625in}%
\pgfsys@useobject{currentmarker}{}%
\end{pgfscope}%
\end{pgfscope}%
\begin{pgfscope}%
\definecolor{textcolor}{rgb}{0.000000,0.000000,0.000000}%
\pgfsetstrokecolor{textcolor}%
\pgfsetfillcolor{textcolor}%
\pgftext[x=2.249285in,y=0.223403in,,top]{\color{textcolor}{\rmfamily\fontsize{6.000000}{7.200000}\selectfont\catcode`\^=\active\def^{\ifmmode\sp\else\^{}\fi}\catcode`\%=\active\def%{\%}40}}%
\end{pgfscope}%
\begin{pgfscope}%
\definecolor{textcolor}{rgb}{0.000000,0.000000,0.000000}%
\pgfsetstrokecolor{textcolor}%
\pgfsetfillcolor{textcolor}%
\pgftext[x=1.410340in,y=0.087200in,,top]{\color{textcolor}{\rmfamily\fontsize{6.000000}{7.200000}\selectfont\catcode`\^=\active\def^{\ifmmode\sp\else\^{}\fi}\catcode`\%=\active\def%{\%}time}}%
\end{pgfscope}%
\begin{pgfscope}%
\pgfpathrectangle{\pgfqpoint{0.487500in}{0.320625in}}{\pgfqpoint{1.845679in}{1.164375in}}%
\pgfusepath{clip}%
\pgfsetrectcap%
\pgfsetroundjoin%
\pgfsetlinewidth{0.803000pt}%
\definecolor{currentstroke}{rgb}{0.909804,0.909804,0.909804}%
\pgfsetstrokecolor{currentstroke}%
\pgfsetdash{}{0pt}%
\pgfpathmoveto{\pgfqpoint{0.487500in}{0.320625in}}%
\pgfpathlineto{\pgfqpoint{2.333179in}{0.320625in}}%
\pgfusepath{stroke}%
\end{pgfscope}%
\begin{pgfscope}%
\pgfsetbuttcap%
\pgfsetroundjoin%
\definecolor{currentfill}{rgb}{0.000000,0.000000,0.000000}%
\pgfsetfillcolor{currentfill}%
\pgfsetlinewidth{0.803000pt}%
\definecolor{currentstroke}{rgb}{0.000000,0.000000,0.000000}%
\pgfsetstrokecolor{currentstroke}%
\pgfsetdash{}{0pt}%
\pgfsys@defobject{currentmarker}{\pgfqpoint{-0.048611in}{0.000000in}}{\pgfqpoint{-0.000000in}{0.000000in}}{%
\pgfpathmoveto{\pgfqpoint{-0.000000in}{0.000000in}}%
\pgfpathlineto{\pgfqpoint{-0.048611in}{0.000000in}}%
\pgfusepath{stroke,fill}%
}%
\begin{pgfscope}%
\pgfsys@transformshift{0.487500in}{0.320625in}%
\pgfsys@useobject{currentmarker}{}%
\end{pgfscope}%
\end{pgfscope}%
\begin{pgfscope}%
\definecolor{textcolor}{rgb}{0.000000,0.000000,0.000000}%
\pgfsetstrokecolor{textcolor}%
\pgfsetfillcolor{textcolor}%
\pgftext[x=0.111728in, y=0.285187in, left, base]{\color{textcolor}{\rmfamily\fontsize{6.000000}{7.200000}\selectfont\catcode`\^=\active\def^{\ifmmode\sp\else\^{}\fi}\catcode`\%=\active\def%{\%}$\mathdefault{10^{-26}}$}}%
\end{pgfscope}%
\begin{pgfscope}%
\pgfpathrectangle{\pgfqpoint{0.487500in}{0.320625in}}{\pgfqpoint{1.845679in}{1.164375in}}%
\pgfusepath{clip}%
\pgfsetrectcap%
\pgfsetroundjoin%
\pgfsetlinewidth{0.803000pt}%
\definecolor{currentstroke}{rgb}{0.909804,0.909804,0.909804}%
\pgfsetstrokecolor{currentstroke}%
\pgfsetdash{}{0pt}%
\pgfpathmoveto{\pgfqpoint{0.487500in}{0.536250in}}%
\pgfpathlineto{\pgfqpoint{2.333179in}{0.536250in}}%
\pgfusepath{stroke}%
\end{pgfscope}%
\begin{pgfscope}%
\pgfsetbuttcap%
\pgfsetroundjoin%
\definecolor{currentfill}{rgb}{0.000000,0.000000,0.000000}%
\pgfsetfillcolor{currentfill}%
\pgfsetlinewidth{0.803000pt}%
\definecolor{currentstroke}{rgb}{0.000000,0.000000,0.000000}%
\pgfsetstrokecolor{currentstroke}%
\pgfsetdash{}{0pt}%
\pgfsys@defobject{currentmarker}{\pgfqpoint{-0.048611in}{0.000000in}}{\pgfqpoint{-0.000000in}{0.000000in}}{%
\pgfpathmoveto{\pgfqpoint{-0.000000in}{0.000000in}}%
\pgfpathlineto{\pgfqpoint{-0.048611in}{0.000000in}}%
\pgfusepath{stroke,fill}%
}%
\begin{pgfscope}%
\pgfsys@transformshift{0.487500in}{0.536250in}%
\pgfsys@useobject{currentmarker}{}%
\end{pgfscope}%
\end{pgfscope}%
\begin{pgfscope}%
\definecolor{textcolor}{rgb}{0.000000,0.000000,0.000000}%
\pgfsetstrokecolor{textcolor}%
\pgfsetfillcolor{textcolor}%
\pgftext[x=0.111728in, y=0.500812in, left, base]{\color{textcolor}{\rmfamily\fontsize{6.000000}{7.200000}\selectfont\catcode`\^=\active\def^{\ifmmode\sp\else\^{}\fi}\catcode`\%=\active\def%{\%}$\mathdefault{10^{-21}}$}}%
\end{pgfscope}%
\begin{pgfscope}%
\pgfpathrectangle{\pgfqpoint{0.487500in}{0.320625in}}{\pgfqpoint{1.845679in}{1.164375in}}%
\pgfusepath{clip}%
\pgfsetrectcap%
\pgfsetroundjoin%
\pgfsetlinewidth{0.803000pt}%
\definecolor{currentstroke}{rgb}{0.909804,0.909804,0.909804}%
\pgfsetstrokecolor{currentstroke}%
\pgfsetdash{}{0pt}%
\pgfpathmoveto{\pgfqpoint{0.487500in}{0.751875in}}%
\pgfpathlineto{\pgfqpoint{2.333179in}{0.751875in}}%
\pgfusepath{stroke}%
\end{pgfscope}%
\begin{pgfscope}%
\pgfsetbuttcap%
\pgfsetroundjoin%
\definecolor{currentfill}{rgb}{0.000000,0.000000,0.000000}%
\pgfsetfillcolor{currentfill}%
\pgfsetlinewidth{0.803000pt}%
\definecolor{currentstroke}{rgb}{0.000000,0.000000,0.000000}%
\pgfsetstrokecolor{currentstroke}%
\pgfsetdash{}{0pt}%
\pgfsys@defobject{currentmarker}{\pgfqpoint{-0.048611in}{0.000000in}}{\pgfqpoint{-0.000000in}{0.000000in}}{%
\pgfpathmoveto{\pgfqpoint{-0.000000in}{0.000000in}}%
\pgfpathlineto{\pgfqpoint{-0.048611in}{0.000000in}}%
\pgfusepath{stroke,fill}%
}%
\begin{pgfscope}%
\pgfsys@transformshift{0.487500in}{0.751875in}%
\pgfsys@useobject{currentmarker}{}%
\end{pgfscope}%
\end{pgfscope}%
\begin{pgfscope}%
\definecolor{textcolor}{rgb}{0.000000,0.000000,0.000000}%
\pgfsetstrokecolor{textcolor}%
\pgfsetfillcolor{textcolor}%
\pgftext[x=0.111728in, y=0.716437in, left, base]{\color{textcolor}{\rmfamily\fontsize{6.000000}{7.200000}\selectfont\catcode`\^=\active\def^{\ifmmode\sp\else\^{}\fi}\catcode`\%=\active\def%{\%}$\mathdefault{10^{-16}}$}}%
\end{pgfscope}%
\begin{pgfscope}%
\pgfpathrectangle{\pgfqpoint{0.487500in}{0.320625in}}{\pgfqpoint{1.845679in}{1.164375in}}%
\pgfusepath{clip}%
\pgfsetrectcap%
\pgfsetroundjoin%
\pgfsetlinewidth{0.803000pt}%
\definecolor{currentstroke}{rgb}{0.909804,0.909804,0.909804}%
\pgfsetstrokecolor{currentstroke}%
\pgfsetdash{}{0pt}%
\pgfpathmoveto{\pgfqpoint{0.487500in}{0.967500in}}%
\pgfpathlineto{\pgfqpoint{2.333179in}{0.967500in}}%
\pgfusepath{stroke}%
\end{pgfscope}%
\begin{pgfscope}%
\pgfsetbuttcap%
\pgfsetroundjoin%
\definecolor{currentfill}{rgb}{0.000000,0.000000,0.000000}%
\pgfsetfillcolor{currentfill}%
\pgfsetlinewidth{0.803000pt}%
\definecolor{currentstroke}{rgb}{0.000000,0.000000,0.000000}%
\pgfsetstrokecolor{currentstroke}%
\pgfsetdash{}{0pt}%
\pgfsys@defobject{currentmarker}{\pgfqpoint{-0.048611in}{0.000000in}}{\pgfqpoint{-0.000000in}{0.000000in}}{%
\pgfpathmoveto{\pgfqpoint{-0.000000in}{0.000000in}}%
\pgfpathlineto{\pgfqpoint{-0.048611in}{0.000000in}}%
\pgfusepath{stroke,fill}%
}%
\begin{pgfscope}%
\pgfsys@transformshift{0.487500in}{0.967500in}%
\pgfsys@useobject{currentmarker}{}%
\end{pgfscope}%
\end{pgfscope}%
\begin{pgfscope}%
\definecolor{textcolor}{rgb}{0.000000,0.000000,0.000000}%
\pgfsetstrokecolor{textcolor}%
\pgfsetfillcolor{textcolor}%
\pgftext[x=0.111728in, y=0.932062in, left, base]{\color{textcolor}{\rmfamily\fontsize{6.000000}{7.200000}\selectfont\catcode`\^=\active\def^{\ifmmode\sp\else\^{}\fi}\catcode`\%=\active\def%{\%}$\mathdefault{10^{-11}}$}}%
\end{pgfscope}%
\begin{pgfscope}%
\pgfpathrectangle{\pgfqpoint{0.487500in}{0.320625in}}{\pgfqpoint{1.845679in}{1.164375in}}%
\pgfusepath{clip}%
\pgfsetrectcap%
\pgfsetroundjoin%
\pgfsetlinewidth{0.803000pt}%
\definecolor{currentstroke}{rgb}{0.909804,0.909804,0.909804}%
\pgfsetstrokecolor{currentstroke}%
\pgfsetdash{}{0pt}%
\pgfpathmoveto{\pgfqpoint{0.487500in}{1.183125in}}%
\pgfpathlineto{\pgfqpoint{2.333179in}{1.183125in}}%
\pgfusepath{stroke}%
\end{pgfscope}%
\begin{pgfscope}%
\pgfsetbuttcap%
\pgfsetroundjoin%
\definecolor{currentfill}{rgb}{0.000000,0.000000,0.000000}%
\pgfsetfillcolor{currentfill}%
\pgfsetlinewidth{0.803000pt}%
\definecolor{currentstroke}{rgb}{0.000000,0.000000,0.000000}%
\pgfsetstrokecolor{currentstroke}%
\pgfsetdash{}{0pt}%
\pgfsys@defobject{currentmarker}{\pgfqpoint{-0.048611in}{0.000000in}}{\pgfqpoint{-0.000000in}{0.000000in}}{%
\pgfpathmoveto{\pgfqpoint{-0.000000in}{0.000000in}}%
\pgfpathlineto{\pgfqpoint{-0.048611in}{0.000000in}}%
\pgfusepath{stroke,fill}%
}%
\begin{pgfscope}%
\pgfsys@transformshift{0.487500in}{1.183125in}%
\pgfsys@useobject{currentmarker}{}%
\end{pgfscope}%
\end{pgfscope}%
\begin{pgfscope}%
\definecolor{textcolor}{rgb}{0.000000,0.000000,0.000000}%
\pgfsetstrokecolor{textcolor}%
\pgfsetfillcolor{textcolor}%
\pgftext[x=0.158989in, y=1.147687in, left, base]{\color{textcolor}{\rmfamily\fontsize{6.000000}{7.200000}\selectfont\catcode`\^=\active\def^{\ifmmode\sp\else\^{}\fi}\catcode`\%=\active\def%{\%}$\mathdefault{10^{-6}}$}}%
\end{pgfscope}%
\begin{pgfscope}%
\pgfpathrectangle{\pgfqpoint{0.487500in}{0.320625in}}{\pgfqpoint{1.845679in}{1.164375in}}%
\pgfusepath{clip}%
\pgfsetrectcap%
\pgfsetroundjoin%
\pgfsetlinewidth{0.803000pt}%
\definecolor{currentstroke}{rgb}{0.909804,0.909804,0.909804}%
\pgfsetstrokecolor{currentstroke}%
\pgfsetdash{}{0pt}%
\pgfpathmoveto{\pgfqpoint{0.487500in}{1.398750in}}%
\pgfpathlineto{\pgfqpoint{2.333179in}{1.398750in}}%
\pgfusepath{stroke}%
\end{pgfscope}%
\begin{pgfscope}%
\pgfsetbuttcap%
\pgfsetroundjoin%
\definecolor{currentfill}{rgb}{0.000000,0.000000,0.000000}%
\pgfsetfillcolor{currentfill}%
\pgfsetlinewidth{0.803000pt}%
\definecolor{currentstroke}{rgb}{0.000000,0.000000,0.000000}%
\pgfsetstrokecolor{currentstroke}%
\pgfsetdash{}{0pt}%
\pgfsys@defobject{currentmarker}{\pgfqpoint{-0.048611in}{0.000000in}}{\pgfqpoint{-0.000000in}{0.000000in}}{%
\pgfpathmoveto{\pgfqpoint{-0.000000in}{0.000000in}}%
\pgfpathlineto{\pgfqpoint{-0.048611in}{0.000000in}}%
\pgfusepath{stroke,fill}%
}%
\begin{pgfscope}%
\pgfsys@transformshift{0.487500in}{1.398750in}%
\pgfsys@useobject{currentmarker}{}%
\end{pgfscope}%
\end{pgfscope}%
\begin{pgfscope}%
\definecolor{textcolor}{rgb}{0.000000,0.000000,0.000000}%
\pgfsetstrokecolor{textcolor}%
\pgfsetfillcolor{textcolor}%
\pgftext[x=0.158989in, y=1.363312in, left, base]{\color{textcolor}{\rmfamily\fontsize{6.000000}{7.200000}\selectfont\catcode`\^=\active\def^{\ifmmode\sp\else\^{}\fi}\catcode`\%=\active\def%{\%}$\mathdefault{10^{-1}}$}}%
\end{pgfscope}%
\begin{pgfscope}%
\definecolor{textcolor}{rgb}{0.000000,0.000000,0.000000}%
\pgfsetstrokecolor{textcolor}%
\pgfsetfillcolor{textcolor}%
\pgftext[x=0.056172in,y=0.902813in,,bottom,rotate=90.000000]{\color{textcolor}{\rmfamily\fontsize{6.000000}{7.200000}\selectfont\catcode`\^=\active\def^{\ifmmode\sp\else\^{}\fi}\catcode`\%=\active\def%{\%}$\ell_t-\ell_\star$}}%
\end{pgfscope}%
\begin{pgfscope}%
\pgfpathrectangle{\pgfqpoint{0.487500in}{0.320625in}}{\pgfqpoint{1.845679in}{1.164375in}}%
\pgfusepath{clip}%
\pgfsetrectcap%
\pgfsetroundjoin%
\pgfsetlinewidth{1.254687pt}%
\definecolor{currentstroke}{rgb}{0.141176,0.345098,0.650980}%
\pgfsetstrokecolor{currentstroke}%
\pgfsetdash{}{0pt}%
\pgfpathmoveto{\pgfqpoint{0.571395in}{1.468525in}}%
\pgfpathlineto{\pgfqpoint{0.588173in}{1.465082in}}%
\pgfpathlineto{\pgfqpoint{0.598660in}{1.461710in}}%
\pgfpathlineto{\pgfqpoint{0.607050in}{1.457758in}}%
\pgfpathlineto{\pgfqpoint{0.615439in}{1.452402in}}%
\pgfpathlineto{\pgfqpoint{0.623829in}{1.445657in}}%
\pgfpathlineto{\pgfqpoint{0.636413in}{1.433755in}}%
\pgfpathlineto{\pgfqpoint{0.663678in}{1.407362in}}%
\pgfpathlineto{\pgfqpoint{0.680457in}{1.392974in}}%
\pgfpathlineto{\pgfqpoint{0.697236in}{1.380002in}}%
\pgfpathlineto{\pgfqpoint{0.718210in}{1.365204in}}%
\pgfpathlineto{\pgfqpoint{0.749670in}{1.344516in}}%
\pgfpathlineto{\pgfqpoint{1.087346in}{1.127817in}}%
\pgfpathlineto{\pgfqpoint{1.309666in}{0.987051in}}%
\pgfpathlineto{\pgfqpoint{1.731236in}{0.721516in}}%
\pgfpathlineto{\pgfqpoint{1.758502in}{0.704552in}}%
\pgfpathlineto{\pgfqpoint{1.810936in}{0.670820in}}%
\pgfpathlineto{\pgfqpoint{1.829812in}{0.660711in}}%
\pgfpathlineto{\pgfqpoint{1.846591in}{0.651559in}}%
\pgfpathlineto{\pgfqpoint{1.859175in}{0.644066in}}%
\pgfpathlineto{\pgfqpoint{1.867565in}{0.640854in}}%
\pgfpathlineto{\pgfqpoint{1.875954in}{0.637730in}}%
\pgfpathlineto{\pgfqpoint{1.880149in}{0.635243in}}%
\pgfpathlineto{\pgfqpoint{1.884343in}{0.631791in}}%
\pgfpathlineto{\pgfqpoint{1.896928in}{0.619924in}}%
\pgfpathlineto{\pgfqpoint{1.899025in}{0.619630in}}%
\pgfpathlineto{\pgfqpoint{1.901122in}{0.620253in}}%
\pgfpathlineto{\pgfqpoint{1.905317in}{0.623417in}}%
\pgfpathlineto{\pgfqpoint{1.909512in}{0.626985in}}%
\pgfpathlineto{\pgfqpoint{1.911609in}{0.628204in}}%
\pgfpathlineto{\pgfqpoint{1.913707in}{0.628753in}}%
\pgfpathlineto{\pgfqpoint{1.915804in}{0.628421in}}%
\pgfpathlineto{\pgfqpoint{1.917901in}{0.626958in}}%
\pgfpathlineto{\pgfqpoint{1.919999in}{0.624030in}}%
\pgfpathlineto{\pgfqpoint{1.922096in}{0.619308in}}%
\pgfpathlineto{\pgfqpoint{1.928388in}{0.600753in}}%
\pgfpathlineto{\pgfqpoint{1.930485in}{0.599600in}}%
\pgfpathlineto{\pgfqpoint{1.932583in}{0.602544in}}%
\pgfpathlineto{\pgfqpoint{1.936777in}{0.610531in}}%
\pgfpathlineto{\pgfqpoint{1.938875in}{0.612766in}}%
\pgfpathlineto{\pgfqpoint{1.940972in}{0.613312in}}%
\pgfpathlineto{\pgfqpoint{1.943070in}{0.612035in}}%
\pgfpathlineto{\pgfqpoint{1.945167in}{0.608767in}}%
\pgfpathlineto{\pgfqpoint{1.947264in}{0.603286in}}%
\pgfpathlineto{\pgfqpoint{1.951459in}{0.587383in}}%
\pgfpathlineto{\pgfqpoint{1.953556in}{0.585385in}}%
\pgfpathlineto{\pgfqpoint{1.955654in}{0.593400in}}%
\pgfpathlineto{\pgfqpoint{1.959848in}{0.613096in}}%
\pgfpathlineto{\pgfqpoint{1.961946in}{0.619497in}}%
\pgfpathlineto{\pgfqpoint{1.964043in}{0.623311in}}%
\pgfpathlineto{\pgfqpoint{1.966141in}{0.624363in}}%
\pgfpathlineto{\pgfqpoint{1.968238in}{0.621992in}}%
\pgfpathlineto{\pgfqpoint{1.970335in}{0.615667in}}%
\pgfpathlineto{\pgfqpoint{1.972433in}{0.605813in}}%
\pgfpathlineto{\pgfqpoint{1.976627in}{0.573871in}}%
\pgfpathlineto{\pgfqpoint{1.978725in}{0.571980in}}%
\pgfpathlineto{\pgfqpoint{1.982919in}{0.598860in}}%
\pgfpathlineto{\pgfqpoint{1.985017in}{0.606422in}}%
\pgfpathlineto{\pgfqpoint{1.987114in}{0.610625in}}%
\pgfpathlineto{\pgfqpoint{1.989212in}{0.612082in}}%
\pgfpathlineto{\pgfqpoint{1.991309in}{0.610972in}}%
\pgfpathlineto{\pgfqpoint{1.993406in}{0.607084in}}%
\pgfpathlineto{\pgfqpoint{1.995504in}{0.599678in}}%
\pgfpathlineto{\pgfqpoint{1.997601in}{0.586985in}}%
\pgfpathlineto{\pgfqpoint{1.999698in}{0.565811in}}%
\pgfpathlineto{\pgfqpoint{2.001796in}{0.558531in}}%
\pgfpathlineto{\pgfqpoint{2.003893in}{0.584413in}}%
\pgfpathlineto{\pgfqpoint{2.005990in}{0.602594in}}%
\pgfpathlineto{\pgfqpoint{2.008088in}{0.614372in}}%
\pgfpathlineto{\pgfqpoint{2.010185in}{0.622039in}}%
\pgfpathlineto{\pgfqpoint{2.012283in}{0.626508in}}%
\pgfpathlineto{\pgfqpoint{2.014380in}{0.627922in}}%
\pgfpathlineto{\pgfqpoint{2.016477in}{0.625720in}}%
\pgfpathlineto{\pgfqpoint{2.018575in}{0.618862in}}%
\pgfpathlineto{\pgfqpoint{2.020672in}{0.606524in}}%
\pgfpathlineto{\pgfqpoint{2.022769in}{0.584094in}}%
\pgfpathlineto{\pgfqpoint{2.024867in}{0.540075in}}%
\pgfpathlineto{\pgfqpoint{2.026964in}{0.579381in}}%
\pgfpathlineto{\pgfqpoint{2.029061in}{0.603488in}}%
\pgfpathlineto{\pgfqpoint{2.031159in}{0.616636in}}%
\pgfpathlineto{\pgfqpoint{2.033256in}{0.624596in}}%
\pgfpathlineto{\pgfqpoint{2.035354in}{0.629273in}}%
\pgfpathlineto{\pgfqpoint{2.037451in}{0.631441in}}%
\pgfpathlineto{\pgfqpoint{2.039548in}{0.631377in}}%
\pgfpathlineto{\pgfqpoint{2.041646in}{0.629038in}}%
\pgfpathlineto{\pgfqpoint{2.043743in}{0.624020in}}%
\pgfpathlineto{\pgfqpoint{2.045840in}{0.615284in}}%
\pgfpathlineto{\pgfqpoint{2.047938in}{0.600071in}}%
\pgfpathlineto{\pgfqpoint{2.050035in}{0.567555in}}%
\pgfpathlineto{\pgfqpoint{2.052132in}{0.552887in}}%
\pgfpathlineto{\pgfqpoint{2.054230in}{0.598028in}}%
\pgfpathlineto{\pgfqpoint{2.056327in}{0.618178in}}%
\pgfpathlineto{\pgfqpoint{2.058425in}{0.630437in}}%
\pgfpathlineto{\pgfqpoint{2.060522in}{0.638382in}}%
\pgfpathlineto{\pgfqpoint{2.062619in}{0.643190in}}%
\pgfpathlineto{\pgfqpoint{2.064717in}{0.645163in}}%
\pgfpathlineto{\pgfqpoint{2.066814in}{0.643977in}}%
\pgfpathlineto{\pgfqpoint{2.068911in}{0.638657in}}%
\pgfpathlineto{\pgfqpoint{2.071009in}{0.628969in}}%
\pgfpathlineto{\pgfqpoint{2.073106in}{0.612772in}}%
\pgfpathlineto{\pgfqpoint{2.077301in}{0.546799in}}%
\pgfpathlineto{\pgfqpoint{2.079398in}{0.599625in}}%
\pgfpathlineto{\pgfqpoint{2.081496in}{0.618811in}}%
\pgfpathlineto{\pgfqpoint{2.083593in}{0.629279in}}%
\pgfpathlineto{\pgfqpoint{2.085690in}{0.635230in}}%
\pgfpathlineto{\pgfqpoint{2.087788in}{0.638065in}}%
\pgfpathlineto{\pgfqpoint{2.089885in}{0.638284in}}%
\pgfpathlineto{\pgfqpoint{2.091982in}{0.635902in}}%
\pgfpathlineto{\pgfqpoint{2.094080in}{0.630459in}}%
\pgfpathlineto{\pgfqpoint{2.096177in}{0.620636in}}%
\pgfpathlineto{\pgfqpoint{2.098274in}{0.602508in}}%
\pgfpathlineto{\pgfqpoint{2.100372in}{0.553800in}}%
\pgfpathlineto{\pgfqpoint{2.102469in}{0.579340in}}%
\pgfpathlineto{\pgfqpoint{2.104566in}{0.614316in}}%
\pgfpathlineto{\pgfqpoint{2.106664in}{0.631743in}}%
\pgfpathlineto{\pgfqpoint{2.108761in}{0.642588in}}%
\pgfpathlineto{\pgfqpoint{2.110859in}{0.649471in}}%
\pgfpathlineto{\pgfqpoint{2.112956in}{0.653232in}}%
\pgfpathlineto{\pgfqpoint{2.115053in}{0.653916in}}%
\pgfpathlineto{\pgfqpoint{2.117151in}{0.650792in}}%
\pgfpathlineto{\pgfqpoint{2.119248in}{0.643921in}}%
\pgfpathlineto{\pgfqpoint{2.121345in}{0.633882in}}%
\pgfpathlineto{\pgfqpoint{2.123443in}{0.619116in}}%
\pgfpathlineto{\pgfqpoint{2.125540in}{0.594822in}}%
\pgfpathlineto{\pgfqpoint{2.127637in}{0.525487in}}%
\pgfpathlineto{\pgfqpoint{2.129735in}{0.574881in}}%
\pgfpathlineto{\pgfqpoint{2.131832in}{0.600946in}}%
\pgfpathlineto{\pgfqpoint{2.133930in}{0.612340in}}%
\pgfpathlineto{\pgfqpoint{2.136027in}{0.617373in}}%
\pgfpathlineto{\pgfqpoint{2.138124in}{0.618041in}}%
\pgfpathlineto{\pgfqpoint{2.140222in}{0.614624in}}%
\pgfpathlineto{\pgfqpoint{2.142319in}{0.606085in}}%
\pgfpathlineto{\pgfqpoint{2.144416in}{0.588368in}}%
\pgfpathlineto{\pgfqpoint{2.146514in}{0.534649in}}%
\pgfpathlineto{\pgfqpoint{2.150708in}{0.605276in}}%
\pgfpathlineto{\pgfqpoint{2.152806in}{0.622484in}}%
\pgfpathlineto{\pgfqpoint{2.154903in}{0.633180in}}%
\pgfpathlineto{\pgfqpoint{2.157001in}{0.639773in}}%
\pgfpathlineto{\pgfqpoint{2.159098in}{0.642954in}}%
\pgfpathlineto{\pgfqpoint{2.161195in}{0.642509in}}%
\pgfpathlineto{\pgfqpoint{2.163293in}{0.637584in}}%
\pgfpathlineto{\pgfqpoint{2.165390in}{0.629780in}}%
\pgfpathlineto{\pgfqpoint{2.167487in}{0.618958in}}%
\pgfpathlineto{\pgfqpoint{2.169585in}{0.603560in}}%
\pgfpathlineto{\pgfqpoint{2.171682in}{0.579000in}}%
\pgfpathlineto{\pgfqpoint{2.173779in}{0.514239in}}%
\pgfpathlineto{\pgfqpoint{2.175877in}{0.553707in}}%
\pgfpathlineto{\pgfqpoint{2.177974in}{0.579778in}}%
\pgfpathlineto{\pgfqpoint{2.180072in}{0.590042in}}%
\pgfpathlineto{\pgfqpoint{2.182169in}{0.593293in}}%
\pgfpathlineto{\pgfqpoint{2.184266in}{0.591235in}}%
\pgfpathlineto{\pgfqpoint{2.186364in}{0.583174in}}%
\pgfpathlineto{\pgfqpoint{2.188461in}{0.564510in}}%
\pgfpathlineto{\pgfqpoint{2.190558in}{0.493356in}}%
\pgfpathlineto{\pgfqpoint{2.192656in}{0.558714in}}%
\pgfpathlineto{\pgfqpoint{2.194753in}{0.589238in}}%
\pgfpathlineto{\pgfqpoint{2.196850in}{0.605700in}}%
\pgfpathlineto{\pgfqpoint{2.198948in}{0.616049in}}%
\pgfpathlineto{\pgfqpoint{2.201045in}{0.622367in}}%
\pgfpathlineto{\pgfqpoint{2.203143in}{0.625201in}}%
\pgfpathlineto{\pgfqpoint{2.205240in}{0.624180in}}%
\pgfpathlineto{\pgfqpoint{2.207337in}{0.618649in}}%
\pgfpathlineto{\pgfqpoint{2.209435in}{0.610514in}}%
\pgfpathlineto{\pgfqpoint{2.211532in}{0.599364in}}%
\pgfpathlineto{\pgfqpoint{2.213629in}{0.583602in}}%
\pgfpathlineto{\pgfqpoint{2.215727in}{0.558526in}}%
\pgfpathlineto{\pgfqpoint{2.217824in}{0.491163in}}%
\pgfpathlineto{\pgfqpoint{2.219921in}{0.533815in}}%
\pgfpathlineto{\pgfqpoint{2.222019in}{0.559016in}}%
\pgfpathlineto{\pgfqpoint{2.224116in}{0.568583in}}%
\pgfpathlineto{\pgfqpoint{2.226214in}{0.570966in}}%
\pgfpathlineto{\pgfqpoint{2.228311in}{0.567587in}}%
\pgfpathlineto{\pgfqpoint{2.230408in}{0.557019in}}%
\pgfpathlineto{\pgfqpoint{2.232506in}{0.530886in}}%
\pgfpathlineto{\pgfqpoint{2.234603in}{0.485308in}}%
\pgfpathlineto{\pgfqpoint{2.236700in}{0.552458in}}%
\pgfpathlineto{\pgfqpoint{2.238798in}{0.576067in}}%
\pgfpathlineto{\pgfqpoint{2.240895in}{0.589993in}}%
\pgfpathlineto{\pgfqpoint{2.242992in}{0.598832in}}%
\pgfpathlineto{\pgfqpoint{2.245090in}{0.603940in}}%
\pgfpathlineto{\pgfqpoint{2.247187in}{0.605520in}}%
\pgfpathlineto{\pgfqpoint{2.249285in}{0.602808in}}%
\pgfpathlineto{\pgfqpoint{2.249285in}{0.602808in}}%
\pgfusepath{stroke}%
\end{pgfscope}%
\begin{pgfscope}%
\pgfpathrectangle{\pgfqpoint{0.487500in}{0.320625in}}{\pgfqpoint{1.845679in}{1.164375in}}%
\pgfusepath{clip}%
\pgfsetbuttcap%
\pgfsetroundjoin%
\pgfsetlinewidth{0.752812pt}%
\definecolor{currentstroke}{rgb}{0.141176,0.345098,0.650980}%
\pgfsetstrokecolor{currentstroke}%
\pgfsetstrokeopacity{0.650000}%
\pgfsetdash{{2.775000pt}{1.200000pt}}{0.000000pt}%
\pgfpathmoveto{\pgfqpoint{0.571395in}{1.468525in}}%
\pgfpathlineto{\pgfqpoint{2.249285in}{1.168861in}}%
\pgfpathlineto{\pgfqpoint{2.249285in}{1.168861in}}%
\pgfusepath{stroke}%
\end{pgfscope}%
\begin{pgfscope}%
\pgfpathrectangle{\pgfqpoint{0.487500in}{0.320625in}}{\pgfqpoint{1.845679in}{1.164375in}}%
\pgfusepath{clip}%
\pgfsetbuttcap%
\pgfsetroundjoin%
\pgfsetlinewidth{0.803000pt}%
\definecolor{currentstroke}{rgb}{0.141176,0.345098,0.650980}%
\pgfsetstrokecolor{currentstroke}%
\pgfsetstrokeopacity{0.800000}%
\pgfsetdash{{0.800000pt}{1.320000pt}}{0.000000pt}%
\pgfpathmoveto{\pgfqpoint{0.571395in}{1.468525in}}%
\pgfpathlineto{\pgfqpoint{2.249285in}{1.468525in}}%
\pgfpathlineto{\pgfqpoint{2.249285in}{1.468525in}}%
\pgfusepath{stroke}%
\end{pgfscope}%
\begin{pgfscope}%
\pgfpathrectangle{\pgfqpoint{0.487500in}{0.320625in}}{\pgfqpoint{1.845679in}{1.164375in}}%
\pgfusepath{clip}%
\pgfsetrectcap%
\pgfsetroundjoin%
\pgfsetlinewidth{1.254687pt}%
\definecolor{currentstroke}{rgb}{0.819608,0.286275,0.356863}%
\pgfsetstrokecolor{currentstroke}%
\pgfsetdash{}{0pt}%
\pgfpathmoveto{\pgfqpoint{0.571395in}{1.468525in}}%
\pgfpathlineto{\pgfqpoint{0.581881in}{1.458695in}}%
\pgfpathlineto{\pgfqpoint{0.590271in}{1.449303in}}%
\pgfpathlineto{\pgfqpoint{0.598660in}{1.438055in}}%
\pgfpathlineto{\pgfqpoint{0.609147in}{1.421752in}}%
\pgfpathlineto{\pgfqpoint{0.659484in}{1.339277in}}%
\pgfpathlineto{\pgfqpoint{0.678360in}{1.312158in}}%
\pgfpathlineto{\pgfqpoint{0.699334in}{1.284091in}}%
\pgfpathlineto{\pgfqpoint{0.724502in}{1.252341in}}%
\pgfpathlineto{\pgfqpoint{0.755962in}{1.214487in}}%
\pgfpathlineto{\pgfqpoint{0.797910in}{1.165805in}}%
\pgfpathlineto{\pgfqpoint{0.858733in}{1.097049in}}%
\pgfpathlineto{\pgfqpoint{0.948920in}{0.996935in}}%
\pgfpathlineto{\pgfqpoint{1.083151in}{0.849760in}}%
\pgfpathlineto{\pgfqpoint{1.215285in}{0.706615in}}%
\pgfpathlineto{\pgfqpoint{1.227869in}{0.693571in}}%
\pgfpathlineto{\pgfqpoint{1.242551in}{0.676276in}}%
\pgfpathlineto{\pgfqpoint{1.257232in}{0.659096in}}%
\pgfpathlineto{\pgfqpoint{1.265621in}{0.650825in}}%
\pgfpathlineto{\pgfqpoint{1.274011in}{0.644232in}}%
\pgfpathlineto{\pgfqpoint{1.286595in}{0.635261in}}%
\pgfpathlineto{\pgfqpoint{1.290790in}{0.631044in}}%
\pgfpathlineto{\pgfqpoint{1.297082in}{0.623267in}}%
\pgfpathlineto{\pgfqpoint{1.305471in}{0.610723in}}%
\pgfpathlineto{\pgfqpoint{1.313861in}{0.595867in}}%
\pgfpathlineto{\pgfqpoint{1.328542in}{0.568981in}}%
\pgfpathlineto{\pgfqpoint{1.332737in}{0.563684in}}%
\pgfpathlineto{\pgfqpoint{1.336932in}{0.560797in}}%
\pgfpathlineto{\pgfqpoint{1.339029in}{0.560380in}}%
\pgfpathlineto{\pgfqpoint{1.341127in}{0.560634in}}%
\pgfpathlineto{\pgfqpoint{1.345321in}{0.562857in}}%
\pgfpathlineto{\pgfqpoint{1.351613in}{0.568758in}}%
\pgfpathlineto{\pgfqpoint{1.360003in}{0.576712in}}%
\pgfpathlineto{\pgfqpoint{1.364198in}{0.579265in}}%
\pgfpathlineto{\pgfqpoint{1.368392in}{0.580168in}}%
\pgfpathlineto{\pgfqpoint{1.370490in}{0.579797in}}%
\pgfpathlineto{\pgfqpoint{1.372587in}{0.578716in}}%
\pgfpathlineto{\pgfqpoint{1.376782in}{0.574038in}}%
\pgfpathlineto{\pgfqpoint{1.380976in}{0.566593in}}%
\pgfpathlineto{\pgfqpoint{1.385171in}{0.555850in}}%
\pgfpathlineto{\pgfqpoint{1.389366in}{0.540946in}}%
\pgfpathlineto{\pgfqpoint{1.393561in}{0.525670in}}%
\pgfpathlineto{\pgfqpoint{1.395658in}{0.524737in}}%
\pgfpathlineto{\pgfqpoint{1.397755in}{0.530092in}}%
\pgfpathlineto{\pgfqpoint{1.404047in}{0.551417in}}%
\pgfpathlineto{\pgfqpoint{1.408242in}{0.560987in}}%
\pgfpathlineto{\pgfqpoint{1.412437in}{0.567347in}}%
\pgfpathlineto{\pgfqpoint{1.416632in}{0.571280in}}%
\pgfpathlineto{\pgfqpoint{1.420826in}{0.573215in}}%
\pgfpathlineto{\pgfqpoint{1.425021in}{0.573331in}}%
\pgfpathlineto{\pgfqpoint{1.429216in}{0.571605in}}%
\pgfpathlineto{\pgfqpoint{1.433410in}{0.567795in}}%
\pgfpathlineto{\pgfqpoint{1.437605in}{0.561307in}}%
\pgfpathlineto{\pgfqpoint{1.441800in}{0.550787in}}%
\pgfpathlineto{\pgfqpoint{1.443897in}{0.543078in}}%
\pgfpathlineto{\pgfqpoint{1.445995in}{0.532592in}}%
\pgfpathlineto{\pgfqpoint{1.448092in}{0.517298in}}%
\pgfpathlineto{\pgfqpoint{1.450189in}{0.493295in}}%
\pgfpathlineto{\pgfqpoint{1.452287in}{0.485703in}}%
\pgfpathlineto{\pgfqpoint{1.454384in}{0.512914in}}%
\pgfpathlineto{\pgfqpoint{1.456481in}{0.531545in}}%
\pgfpathlineto{\pgfqpoint{1.460676in}{0.553973in}}%
\pgfpathlineto{\pgfqpoint{1.464871in}{0.567670in}}%
\pgfpathlineto{\pgfqpoint{1.469066in}{0.576945in}}%
\pgfpathlineto{\pgfqpoint{1.473260in}{0.583300in}}%
\pgfpathlineto{\pgfqpoint{1.477455in}{0.587333in}}%
\pgfpathlineto{\pgfqpoint{1.481650in}{0.589203in}}%
\pgfpathlineto{\pgfqpoint{1.483747in}{0.589282in}}%
\pgfpathlineto{\pgfqpoint{1.485845in}{0.588719in}}%
\pgfpathlineto{\pgfqpoint{1.487942in}{0.587415in}}%
\pgfpathlineto{\pgfqpoint{1.490039in}{0.585218in}}%
\pgfpathlineto{\pgfqpoint{1.496331in}{0.576813in}}%
\pgfpathlineto{\pgfqpoint{1.498429in}{0.577210in}}%
\pgfpathlineto{\pgfqpoint{1.500526in}{0.580317in}}%
\pgfpathlineto{\pgfqpoint{1.506818in}{0.593940in}}%
\pgfpathlineto{\pgfqpoint{1.511013in}{0.599555in}}%
\pgfpathlineto{\pgfqpoint{1.513110in}{0.600813in}}%
\pgfpathlineto{\pgfqpoint{1.515208in}{0.600991in}}%
\pgfpathlineto{\pgfqpoint{1.517305in}{0.600022in}}%
\pgfpathlineto{\pgfqpoint{1.519402in}{0.597776in}}%
\pgfpathlineto{\pgfqpoint{1.521500in}{0.594025in}}%
\pgfpathlineto{\pgfqpoint{1.523597in}{0.588368in}}%
\pgfpathlineto{\pgfqpoint{1.525694in}{0.580069in}}%
\pgfpathlineto{\pgfqpoint{1.527792in}{0.567732in}}%
\pgfpathlineto{\pgfqpoint{1.529889in}{0.549224in}}%
\pgfpathlineto{\pgfqpoint{1.531987in}{0.536199in}}%
\pgfpathlineto{\pgfqpoint{1.536181in}{0.574527in}}%
\pgfpathlineto{\pgfqpoint{1.538279in}{0.587318in}}%
\pgfpathlineto{\pgfqpoint{1.540376in}{0.596349in}}%
\pgfpathlineto{\pgfqpoint{1.542473in}{0.602757in}}%
\pgfpathlineto{\pgfqpoint{1.544571in}{0.607076in}}%
\pgfpathlineto{\pgfqpoint{1.546668in}{0.609503in}}%
\pgfpathlineto{\pgfqpoint{1.548765in}{0.609977in}}%
\pgfpathlineto{\pgfqpoint{1.550863in}{0.608132in}}%
\pgfpathlineto{\pgfqpoint{1.552960in}{0.603009in}}%
\pgfpathlineto{\pgfqpoint{1.555058in}{0.592321in}}%
\pgfpathlineto{\pgfqpoint{1.557155in}{0.572169in}}%
\pgfpathlineto{\pgfqpoint{1.559252in}{0.539881in}}%
\pgfpathlineto{\pgfqpoint{1.563447in}{0.592255in}}%
\pgfpathlineto{\pgfqpoint{1.565544in}{0.605736in}}%
\pgfpathlineto{\pgfqpoint{1.567642in}{0.614193in}}%
\pgfpathlineto{\pgfqpoint{1.569739in}{0.619435in}}%
\pgfpathlineto{\pgfqpoint{1.571836in}{0.622242in}}%
\pgfpathlineto{\pgfqpoint{1.573934in}{0.622932in}}%
\pgfpathlineto{\pgfqpoint{1.576031in}{0.621543in}}%
\pgfpathlineto{\pgfqpoint{1.578129in}{0.617831in}}%
\pgfpathlineto{\pgfqpoint{1.580226in}{0.611135in}}%
\pgfpathlineto{\pgfqpoint{1.582323in}{0.599841in}}%
\pgfpathlineto{\pgfqpoint{1.584421in}{0.579062in}}%
\pgfpathlineto{\pgfqpoint{1.586518in}{0.513327in}}%
\pgfpathlineto{\pgfqpoint{1.588615in}{0.568403in}}%
\pgfpathlineto{\pgfqpoint{1.590713in}{0.598283in}}%
\pgfpathlineto{\pgfqpoint{1.592810in}{0.614196in}}%
\pgfpathlineto{\pgfqpoint{1.594907in}{0.624282in}}%
\pgfpathlineto{\pgfqpoint{1.597005in}{0.630723in}}%
\pgfpathlineto{\pgfqpoint{1.599102in}{0.634240in}}%
\pgfpathlineto{\pgfqpoint{1.601199in}{0.634862in}}%
\pgfpathlineto{\pgfqpoint{1.603297in}{0.631928in}}%
\pgfpathlineto{\pgfqpoint{1.605394in}{0.624788in}}%
\pgfpathlineto{\pgfqpoint{1.607492in}{0.613399in}}%
\pgfpathlineto{\pgfqpoint{1.609589in}{0.595077in}}%
\pgfpathlineto{\pgfqpoint{1.611686in}{0.560621in}}%
\pgfpathlineto{\pgfqpoint{1.613784in}{0.556986in}}%
\pgfpathlineto{\pgfqpoint{1.615881in}{0.588707in}}%
\pgfpathlineto{\pgfqpoint{1.617978in}{0.603688in}}%
\pgfpathlineto{\pgfqpoint{1.620076in}{0.611530in}}%
\pgfpathlineto{\pgfqpoint{1.622173in}{0.615036in}}%
\pgfpathlineto{\pgfqpoint{1.624270in}{0.615099in}}%
\pgfpathlineto{\pgfqpoint{1.626368in}{0.611729in}}%
\pgfpathlineto{\pgfqpoint{1.628465in}{0.604050in}}%
\pgfpathlineto{\pgfqpoint{1.630563in}{0.589181in}}%
\pgfpathlineto{\pgfqpoint{1.632660in}{0.554073in}}%
\pgfpathlineto{\pgfqpoint{1.634757in}{0.545982in}}%
\pgfpathlineto{\pgfqpoint{1.636855in}{0.591182in}}%
\pgfpathlineto{\pgfqpoint{1.638952in}{0.610958in}}%
\pgfpathlineto{\pgfqpoint{1.641049in}{0.622788in}}%
\pgfpathlineto{\pgfqpoint{1.643147in}{0.630048in}}%
\pgfpathlineto{\pgfqpoint{1.645244in}{0.633713in}}%
\pgfpathlineto{\pgfqpoint{1.647341in}{0.633713in}}%
\pgfpathlineto{\pgfqpoint{1.649439in}{0.629026in}}%
\pgfpathlineto{\pgfqpoint{1.651536in}{0.620035in}}%
\pgfpathlineto{\pgfqpoint{1.653634in}{0.605666in}}%
\pgfpathlineto{\pgfqpoint{1.655731in}{0.580415in}}%
\pgfpathlineto{\pgfqpoint{1.657828in}{0.524118in}}%
\pgfpathlineto{\pgfqpoint{1.659926in}{0.575467in}}%
\pgfpathlineto{\pgfqpoint{1.662023in}{0.597955in}}%
\pgfpathlineto{\pgfqpoint{1.664120in}{0.609111in}}%
\pgfpathlineto{\pgfqpoint{1.666218in}{0.614782in}}%
\pgfpathlineto{\pgfqpoint{1.668315in}{0.616682in}}%
\pgfpathlineto{\pgfqpoint{1.670412in}{0.615263in}}%
\pgfpathlineto{\pgfqpoint{1.672510in}{0.610188in}}%
\pgfpathlineto{\pgfqpoint{1.674607in}{0.599978in}}%
\pgfpathlineto{\pgfqpoint{1.676705in}{0.579598in}}%
\pgfpathlineto{\pgfqpoint{1.678802in}{0.505737in}}%
\pgfpathlineto{\pgfqpoint{1.680899in}{0.572547in}}%
\pgfpathlineto{\pgfqpoint{1.682997in}{0.601951in}}%
\pgfpathlineto{\pgfqpoint{1.685094in}{0.617613in}}%
\pgfpathlineto{\pgfqpoint{1.687191in}{0.627270in}}%
\pgfpathlineto{\pgfqpoint{1.689289in}{0.632928in}}%
\pgfpathlineto{\pgfqpoint{1.691386in}{0.635063in}}%
\pgfpathlineto{\pgfqpoint{1.693483in}{0.633177in}}%
\pgfpathlineto{\pgfqpoint{1.695581in}{0.626381in}}%
\pgfpathlineto{\pgfqpoint{1.697678in}{0.614855in}}%
\pgfpathlineto{\pgfqpoint{1.699776in}{0.595316in}}%
\pgfpathlineto{\pgfqpoint{1.701873in}{0.550904in}}%
\pgfpathlineto{\pgfqpoint{1.703970in}{0.560907in}}%
\pgfpathlineto{\pgfqpoint{1.706068in}{0.594987in}}%
\pgfpathlineto{\pgfqpoint{1.708165in}{0.609996in}}%
\pgfpathlineto{\pgfqpoint{1.710262in}{0.618011in}}%
\pgfpathlineto{\pgfqpoint{1.712360in}{0.621862in}}%
\pgfpathlineto{\pgfqpoint{1.714457in}{0.622469in}}%
\pgfpathlineto{\pgfqpoint{1.716554in}{0.619949in}}%
\pgfpathlineto{\pgfqpoint{1.718652in}{0.613717in}}%
\pgfpathlineto{\pgfqpoint{1.720749in}{0.601866in}}%
\pgfpathlineto{\pgfqpoint{1.722847in}{0.577582in}}%
\pgfpathlineto{\pgfqpoint{1.724944in}{0.478217in}}%
\pgfpathlineto{\pgfqpoint{1.727041in}{0.584997in}}%
\pgfpathlineto{\pgfqpoint{1.729139in}{0.610172in}}%
\pgfpathlineto{\pgfqpoint{1.731236in}{0.624306in}}%
\pgfpathlineto{\pgfqpoint{1.733333in}{0.633077in}}%
\pgfpathlineto{\pgfqpoint{1.735431in}{0.638076in}}%
\pgfpathlineto{\pgfqpoint{1.737528in}{0.639608in}}%
\pgfpathlineto{\pgfqpoint{1.739625in}{0.637023in}}%
\pgfpathlineto{\pgfqpoint{1.741723in}{0.629694in}}%
\pgfpathlineto{\pgfqpoint{1.743820in}{0.617696in}}%
\pgfpathlineto{\pgfqpoint{1.745918in}{0.597483in}}%
\pgfpathlineto{\pgfqpoint{1.748015in}{0.549940in}}%
\pgfpathlineto{\pgfqpoint{1.750112in}{0.564571in}}%
\pgfpathlineto{\pgfqpoint{1.752210in}{0.597839in}}%
\pgfpathlineto{\pgfqpoint{1.754307in}{0.612349in}}%
\pgfpathlineto{\pgfqpoint{1.756404in}{0.619924in}}%
\pgfpathlineto{\pgfqpoint{1.758502in}{0.623297in}}%
\pgfpathlineto{\pgfqpoint{1.760599in}{0.623318in}}%
\pgfpathlineto{\pgfqpoint{1.762696in}{0.619989in}}%
\pgfpathlineto{\pgfqpoint{1.764794in}{0.612467in}}%
\pgfpathlineto{\pgfqpoint{1.766891in}{0.598026in}}%
\pgfpathlineto{\pgfqpoint{1.768988in}{0.564749in}}%
\pgfpathlineto{\pgfqpoint{1.771086in}{0.548000in}}%
\pgfpathlineto{\pgfqpoint{1.773183in}{0.597269in}}%
\pgfpathlineto{\pgfqpoint{1.775281in}{0.617691in}}%
\pgfpathlineto{\pgfqpoint{1.777378in}{0.629802in}}%
\pgfpathlineto{\pgfqpoint{1.779475in}{0.637242in}}%
\pgfpathlineto{\pgfqpoint{1.781573in}{0.641064in}}%
\pgfpathlineto{\pgfqpoint{1.783670in}{0.641245in}}%
\pgfpathlineto{\pgfqpoint{1.785767in}{0.636789in}}%
\pgfpathlineto{\pgfqpoint{1.787865in}{0.628198in}}%
\pgfpathlineto{\pgfqpoint{1.789962in}{0.614822in}}%
\pgfpathlineto{\pgfqpoint{1.792059in}{0.592393in}}%
\pgfpathlineto{\pgfqpoint{1.794157in}{0.533430in}}%
\pgfpathlineto{\pgfqpoint{1.798352in}{0.596993in}}%
\pgfpathlineto{\pgfqpoint{1.800449in}{0.609701in}}%
\pgfpathlineto{\pgfqpoint{1.802546in}{0.615871in}}%
\pgfpathlineto{\pgfqpoint{1.804644in}{0.617787in}}%
\pgfpathlineto{\pgfqpoint{1.806741in}{0.616005in}}%
\pgfpathlineto{\pgfqpoint{1.808838in}{0.610069in}}%
\pgfpathlineto{\pgfqpoint{1.810936in}{0.597937in}}%
\pgfpathlineto{\pgfqpoint{1.813033in}{0.571543in}}%
\pgfpathlineto{\pgfqpoint{1.815130in}{0.513578in}}%
\pgfpathlineto{\pgfqpoint{1.817228in}{0.587578in}}%
\pgfpathlineto{\pgfqpoint{1.819325in}{0.611068in}}%
\pgfpathlineto{\pgfqpoint{1.821423in}{0.624534in}}%
\pgfpathlineto{\pgfqpoint{1.823520in}{0.632818in}}%
\pgfpathlineto{\pgfqpoint{1.825617in}{0.637280in}}%
\pgfpathlineto{\pgfqpoint{1.827715in}{0.638037in}}%
\pgfpathlineto{\pgfqpoint{1.829812in}{0.634142in}}%
\pgfpathlineto{\pgfqpoint{1.831909in}{0.626031in}}%
\pgfpathlineto{\pgfqpoint{1.834007in}{0.613553in}}%
\pgfpathlineto{\pgfqpoint{1.836104in}{0.593354in}}%
\pgfpathlineto{\pgfqpoint{1.838201in}{0.548661in}}%
\pgfpathlineto{\pgfqpoint{1.840299in}{0.552355in}}%
\pgfpathlineto{\pgfqpoint{1.842396in}{0.588092in}}%
\pgfpathlineto{\pgfqpoint{1.844494in}{0.602460in}}%
\pgfpathlineto{\pgfqpoint{1.846591in}{0.609346in}}%
\pgfpathlineto{\pgfqpoint{1.848688in}{0.611634in}}%
\pgfpathlineto{\pgfqpoint{1.850786in}{0.610043in}}%
\pgfpathlineto{\pgfqpoint{1.852883in}{0.604160in}}%
\pgfpathlineto{\pgfqpoint{1.854980in}{0.591889in}}%
\pgfpathlineto{\pgfqpoint{1.857078in}{0.564675in}}%
\pgfpathlineto{\pgfqpoint{1.859175in}{0.515572in}}%
\pgfpathlineto{\pgfqpoint{1.861272in}{0.583568in}}%
\pgfpathlineto{\pgfqpoint{1.863370in}{0.606554in}}%
\pgfpathlineto{\pgfqpoint{1.865467in}{0.619814in}}%
\pgfpathlineto{\pgfqpoint{1.867565in}{0.627954in}}%
\pgfpathlineto{\pgfqpoint{1.869662in}{0.632263in}}%
\pgfpathlineto{\pgfqpoint{1.871759in}{0.632805in}}%
\pgfpathlineto{\pgfqpoint{1.873857in}{0.628580in}}%
\pgfpathlineto{\pgfqpoint{1.875954in}{0.619992in}}%
\pgfpathlineto{\pgfqpoint{1.878051in}{0.606380in}}%
\pgfpathlineto{\pgfqpoint{1.880149in}{0.582978in}}%
\pgfpathlineto{\pgfqpoint{1.882246in}{0.512183in}}%
\pgfpathlineto{\pgfqpoint{1.884343in}{0.565574in}}%
\pgfpathlineto{\pgfqpoint{1.886441in}{0.592372in}}%
\pgfpathlineto{\pgfqpoint{1.888538in}{0.604699in}}%
\pgfpathlineto{\pgfqpoint{1.890636in}{0.610854in}}%
\pgfpathlineto{\pgfqpoint{1.892733in}{0.612954in}}%
\pgfpathlineto{\pgfqpoint{1.894830in}{0.611557in}}%
\pgfpathlineto{\pgfqpoint{1.896928in}{0.606322in}}%
\pgfpathlineto{\pgfqpoint{1.899025in}{0.595638in}}%
\pgfpathlineto{\pgfqpoint{1.901122in}{0.573736in}}%
\pgfpathlineto{\pgfqpoint{1.903220in}{0.460888in}}%
\pgfpathlineto{\pgfqpoint{1.905317in}{0.573769in}}%
\pgfpathlineto{\pgfqpoint{1.907414in}{0.601121in}}%
\pgfpathlineto{\pgfqpoint{1.909512in}{0.616070in}}%
\pgfpathlineto{\pgfqpoint{1.911609in}{0.625285in}}%
\pgfpathlineto{\pgfqpoint{1.913707in}{0.630550in}}%
\pgfpathlineto{\pgfqpoint{1.915804in}{0.632221in}}%
\pgfpathlineto{\pgfqpoint{1.917901in}{0.629629in}}%
\pgfpathlineto{\pgfqpoint{1.919999in}{0.622130in}}%
\pgfpathlineto{\pgfqpoint{1.922096in}{0.609633in}}%
\pgfpathlineto{\pgfqpoint{1.924193in}{0.587907in}}%
\pgfpathlineto{\pgfqpoint{1.926291in}{0.527960in}}%
\pgfpathlineto{\pgfqpoint{1.930485in}{0.595230in}}%
\pgfpathlineto{\pgfqpoint{1.932583in}{0.608838in}}%
\pgfpathlineto{\pgfqpoint{1.934680in}{0.616095in}}%
\pgfpathlineto{\pgfqpoint{1.936777in}{0.619393in}}%
\pgfpathlineto{\pgfqpoint{1.938875in}{0.619500in}}%
\pgfpathlineto{\pgfqpoint{1.940972in}{0.616427in}}%
\pgfpathlineto{\pgfqpoint{1.943070in}{0.609438in}}%
\pgfpathlineto{\pgfqpoint{1.945167in}{0.596223in}}%
\pgfpathlineto{\pgfqpoint{1.947264in}{0.567738in}}%
\pgfpathlineto{\pgfqpoint{1.949362in}{0.524049in}}%
\pgfpathlineto{\pgfqpoint{1.951459in}{0.587987in}}%
\pgfpathlineto{\pgfqpoint{1.953556in}{0.610390in}}%
\pgfpathlineto{\pgfqpoint{1.955654in}{0.623391in}}%
\pgfpathlineto{\pgfqpoint{1.957751in}{0.631445in}}%
\pgfpathlineto{\pgfqpoint{1.959848in}{0.635840in}}%
\pgfpathlineto{\pgfqpoint{1.961946in}{0.636720in}}%
\pgfpathlineto{\pgfqpoint{1.964043in}{0.633229in}}%
\pgfpathlineto{\pgfqpoint{1.966141in}{0.625033in}}%
\pgfpathlineto{\pgfqpoint{1.968238in}{0.611454in}}%
\pgfpathlineto{\pgfqpoint{1.970335in}{0.587121in}}%
\pgfpathlineto{\pgfqpoint{1.972433in}{0.493684in}}%
\pgfpathlineto{\pgfqpoint{1.974530in}{0.577644in}}%
\pgfpathlineto{\pgfqpoint{1.976627in}{0.602823in}}%
\pgfpathlineto{\pgfqpoint{1.978725in}{0.615071in}}%
\pgfpathlineto{\pgfqpoint{1.980822in}{0.621583in}}%
\pgfpathlineto{\pgfqpoint{1.982919in}{0.624341in}}%
\pgfpathlineto{\pgfqpoint{1.985017in}{0.623962in}}%
\pgfpathlineto{\pgfqpoint{1.987114in}{0.620353in}}%
\pgfpathlineto{\pgfqpoint{1.989212in}{0.612625in}}%
\pgfpathlineto{\pgfqpoint{1.991309in}{0.598023in}}%
\pgfpathlineto{\pgfqpoint{1.993406in}{0.564578in}}%
\pgfpathlineto{\pgfqpoint{1.995504in}{0.547902in}}%
\pgfpathlineto{\pgfqpoint{1.997601in}{0.597070in}}%
\pgfpathlineto{\pgfqpoint{1.999698in}{0.617445in}}%
\pgfpathlineto{\pgfqpoint{2.001796in}{0.629540in}}%
\pgfpathlineto{\pgfqpoint{2.003893in}{0.636997in}}%
\pgfpathlineto{\pgfqpoint{2.005990in}{0.640877in}}%
\pgfpathlineto{\pgfqpoint{2.008088in}{0.641190in}}%
\pgfpathlineto{\pgfqpoint{2.010185in}{0.636965in}}%
\pgfpathlineto{\pgfqpoint{2.012283in}{0.628439in}}%
\pgfpathlineto{\pgfqpoint{2.014380in}{0.614923in}}%
\pgfpathlineto{\pgfqpoint{2.016477in}{0.591747in}}%
\pgfpathlineto{\pgfqpoint{2.018575in}{0.523035in}}%
\pgfpathlineto{\pgfqpoint{2.020672in}{0.573176in}}%
\pgfpathlineto{\pgfqpoint{2.022769in}{0.600395in}}%
\pgfpathlineto{\pgfqpoint{2.024867in}{0.612853in}}%
\pgfpathlineto{\pgfqpoint{2.026964in}{0.619081in}}%
\pgfpathlineto{\pgfqpoint{2.029061in}{0.621237in}}%
\pgfpathlineto{\pgfqpoint{2.031159in}{0.619894in}}%
\pgfpathlineto{\pgfqpoint{2.033256in}{0.614726in}}%
\pgfpathlineto{\pgfqpoint{2.035354in}{0.604150in}}%
\pgfpathlineto{\pgfqpoint{2.037451in}{0.582528in}}%
\pgfpathlineto{\pgfqpoint{2.039548in}{0.479244in}}%
\pgfpathlineto{\pgfqpoint{2.041646in}{0.581433in}}%
\pgfpathlineto{\pgfqpoint{2.043743in}{0.609114in}}%
\pgfpathlineto{\pgfqpoint{2.045840in}{0.624179in}}%
\pgfpathlineto{\pgfqpoint{2.047938in}{0.633464in}}%
\pgfpathlineto{\pgfqpoint{2.050035in}{0.638784in}}%
\pgfpathlineto{\pgfqpoint{2.052132in}{0.640511in}}%
\pgfpathlineto{\pgfqpoint{2.054230in}{0.637995in}}%
\pgfpathlineto{\pgfqpoint{2.056327in}{0.630867in}}%
\pgfpathlineto{\pgfqpoint{2.058425in}{0.619825in}}%
\pgfpathlineto{\pgfqpoint{2.060522in}{0.602518in}}%
\pgfpathlineto{\pgfqpoint{2.062619in}{0.569669in}}%
\pgfpathlineto{\pgfqpoint{2.064717in}{0.524704in}}%
\pgfpathlineto{\pgfqpoint{2.066814in}{0.582824in}}%
\pgfpathlineto{\pgfqpoint{2.068911in}{0.600842in}}%
\pgfpathlineto{\pgfqpoint{2.071009in}{0.609413in}}%
\pgfpathlineto{\pgfqpoint{2.073106in}{0.612836in}}%
\pgfpathlineto{\pgfqpoint{2.075203in}{0.612253in}}%
\pgfpathlineto{\pgfqpoint{2.077301in}{0.607540in}}%
\pgfpathlineto{\pgfqpoint{2.079398in}{0.597155in}}%
\pgfpathlineto{\pgfqpoint{2.081496in}{0.575223in}}%
\pgfpathlineto{\pgfqpoint{2.083593in}{0.444407in}}%
\pgfpathlineto{\pgfqpoint{2.085690in}{0.577056in}}%
\pgfpathlineto{\pgfqpoint{2.087788in}{0.604142in}}%
\pgfpathlineto{\pgfqpoint{2.089885in}{0.619024in}}%
\pgfpathlineto{\pgfqpoint{2.091982in}{0.628176in}}%
\pgfpathlineto{\pgfqpoint{2.094080in}{0.633326in}}%
\pgfpathlineto{\pgfqpoint{2.096177in}{0.634770in}}%
\pgfpathlineto{\pgfqpoint{2.098274in}{0.631710in}}%
\pgfpathlineto{\pgfqpoint{2.100372in}{0.624103in}}%
\pgfpathlineto{\pgfqpoint{2.102469in}{0.612242in}}%
\pgfpathlineto{\pgfqpoint{2.104566in}{0.593114in}}%
\pgfpathlineto{\pgfqpoint{2.106664in}{0.552744in}}%
\pgfpathlineto{\pgfqpoint{2.108761in}{0.543832in}}%
\pgfpathlineto{\pgfqpoint{2.110859in}{0.584252in}}%
\pgfpathlineto{\pgfqpoint{2.112956in}{0.599763in}}%
\pgfpathlineto{\pgfqpoint{2.115053in}{0.607349in}}%
\pgfpathlineto{\pgfqpoint{2.117151in}{0.610266in}}%
\pgfpathlineto{\pgfqpoint{2.119248in}{0.609400in}}%
\pgfpathlineto{\pgfqpoint{2.121345in}{0.604547in}}%
\pgfpathlineto{\pgfqpoint{2.123443in}{0.594169in}}%
\pgfpathlineto{\pgfqpoint{2.125540in}{0.572596in}}%
\pgfpathlineto{\pgfqpoint{2.127637in}{0.464956in}}%
\pgfpathlineto{\pgfqpoint{2.129735in}{0.572297in}}%
\pgfpathlineto{\pgfqpoint{2.131832in}{0.599888in}}%
\pgfpathlineto{\pgfqpoint{2.133930in}{0.614937in}}%
\pgfpathlineto{\pgfqpoint{2.136027in}{0.624197in}}%
\pgfpathlineto{\pgfqpoint{2.138124in}{0.629459in}}%
\pgfpathlineto{\pgfqpoint{2.140222in}{0.631061in}}%
\pgfpathlineto{\pgfqpoint{2.142319in}{0.628281in}}%
\pgfpathlineto{\pgfqpoint{2.144416in}{0.620655in}}%
\pgfpathlineto{\pgfqpoint{2.146514in}{0.608157in}}%
\pgfpathlineto{\pgfqpoint{2.148611in}{0.586722in}}%
\pgfpathlineto{\pgfqpoint{2.150708in}{0.530205in}}%
\pgfpathlineto{\pgfqpoint{2.154903in}{0.591856in}}%
\pgfpathlineto{\pgfqpoint{2.157001in}{0.605577in}}%
\pgfpathlineto{\pgfqpoint{2.159098in}{0.612771in}}%
\pgfpathlineto{\pgfqpoint{2.161195in}{0.615909in}}%
\pgfpathlineto{\pgfqpoint{2.163293in}{0.615754in}}%
\pgfpathlineto{\pgfqpoint{2.165390in}{0.612269in}}%
\pgfpathlineto{\pgfqpoint{2.167487in}{0.604572in}}%
\pgfpathlineto{\pgfqpoint{2.169585in}{0.589848in}}%
\pgfpathlineto{\pgfqpoint{2.171682in}{0.555555in}}%
\pgfpathlineto{\pgfqpoint{2.173779in}{0.542658in}}%
\pgfpathlineto{\pgfqpoint{2.175877in}{0.590214in}}%
\pgfpathlineto{\pgfqpoint{2.177974in}{0.610324in}}%
\pgfpathlineto{\pgfqpoint{2.180072in}{0.622288in}}%
\pgfpathlineto{\pgfqpoint{2.182169in}{0.629633in}}%
\pgfpathlineto{\pgfqpoint{2.184266in}{0.633378in}}%
\pgfpathlineto{\pgfqpoint{2.186364in}{0.633484in}}%
\pgfpathlineto{\pgfqpoint{2.188461in}{0.628919in}}%
\pgfpathlineto{\pgfqpoint{2.190558in}{0.619642in}}%
\pgfpathlineto{\pgfqpoint{2.192656in}{0.604084in}}%
\pgfpathlineto{\pgfqpoint{2.194753in}{0.573860in}}%
\pgfpathlineto{\pgfqpoint{2.196850in}{0.518785in}}%
\pgfpathlineto{\pgfqpoint{2.198948in}{0.585508in}}%
\pgfpathlineto{\pgfqpoint{2.201045in}{0.605632in}}%
\pgfpathlineto{\pgfqpoint{2.203143in}{0.616082in}}%
\pgfpathlineto{\pgfqpoint{2.205240in}{0.621621in}}%
\pgfpathlineto{\pgfqpoint{2.207337in}{0.623720in}}%
\pgfpathlineto{\pgfqpoint{2.209435in}{0.622809in}}%
\pgfpathlineto{\pgfqpoint{2.211532in}{0.618681in}}%
\pgfpathlineto{\pgfqpoint{2.213629in}{0.610308in}}%
\pgfpathlineto{\pgfqpoint{2.215727in}{0.594559in}}%
\pgfpathlineto{\pgfqpoint{2.217824in}{0.556556in}}%
\pgfpathlineto{\pgfqpoint{2.219921in}{0.555744in}}%
\pgfpathlineto{\pgfqpoint{2.222019in}{0.598432in}}%
\pgfpathlineto{\pgfqpoint{2.224116in}{0.617570in}}%
\pgfpathlineto{\pgfqpoint{2.226214in}{0.629082in}}%
\pgfpathlineto{\pgfqpoint{2.228311in}{0.636158in}}%
\pgfpathlineto{\pgfqpoint{2.230408in}{0.639730in}}%
\pgfpathlineto{\pgfqpoint{2.232506in}{0.639732in}}%
\pgfpathlineto{\pgfqpoint{2.234603in}{0.635324in}}%
\pgfpathlineto{\pgfqpoint{2.249285in}{0.594943in}}%
\pgfpathlineto{\pgfqpoint{2.249285in}{0.594943in}}%
\pgfusepath{stroke}%
\end{pgfscope}%
\begin{pgfscope}%
\pgfpathrectangle{\pgfqpoint{0.487500in}{0.320625in}}{\pgfqpoint{1.845679in}{1.164375in}}%
\pgfusepath{clip}%
\pgfsetbuttcap%
\pgfsetroundjoin%
\pgfsetlinewidth{0.752812pt}%
\definecolor{currentstroke}{rgb}{0.819608,0.286275,0.356863}%
\pgfsetstrokecolor{currentstroke}%
\pgfsetstrokeopacity{0.650000}%
\pgfsetdash{{2.775000pt}{1.200000pt}}{0.000000pt}%
\pgfpathmoveto{\pgfqpoint{0.571395in}{1.468525in}}%
\pgfpathlineto{\pgfqpoint{1.862470in}{0.315625in}}%
\pgfpathlineto{\pgfqpoint{1.862470in}{0.315625in}}%
\pgfusepath{stroke}%
\end{pgfscope}%
\begin{pgfscope}%
\pgfpathrectangle{\pgfqpoint{0.487500in}{0.320625in}}{\pgfqpoint{1.845679in}{1.164375in}}%
\pgfusepath{clip}%
\pgfsetbuttcap%
\pgfsetroundjoin%
\pgfsetlinewidth{0.803000pt}%
\definecolor{currentstroke}{rgb}{0.819608,0.286275,0.356863}%
\pgfsetstrokecolor{currentstroke}%
\pgfsetstrokeopacity{0.800000}%
\pgfsetdash{{0.800000pt}{1.320000pt}}{0.000000pt}%
\pgfpathmoveto{\pgfqpoint{0.571395in}{1.468525in}}%
\pgfpathlineto{\pgfqpoint{2.249285in}{1.468525in}}%
\pgfpathlineto{\pgfqpoint{2.249285in}{1.468525in}}%
\pgfusepath{stroke}%
\end{pgfscope}%
\begin{pgfscope}%
\pgfpathrectangle{\pgfqpoint{0.487500in}{0.320625in}}{\pgfqpoint{1.845679in}{1.164375in}}%
\pgfusepath{clip}%
\pgfsetrectcap%
\pgfsetroundjoin%
\pgfsetlinewidth{1.254687pt}%
\definecolor{currentstroke}{rgb}{0.000000,0.529412,0.423529}%
\pgfsetstrokecolor{currentstroke}%
\pgfsetdash{}{0pt}%
\pgfpathmoveto{\pgfqpoint{0.571395in}{1.468525in}}%
\pgfpathlineto{\pgfqpoint{0.581881in}{1.449020in}}%
\pgfpathlineto{\pgfqpoint{0.592368in}{1.426689in}}%
\pgfpathlineto{\pgfqpoint{0.604952in}{1.396826in}}%
\pgfpathlineto{\pgfqpoint{0.646900in}{1.295102in}}%
\pgfpathlineto{\pgfqpoint{0.669971in}{1.242905in}}%
\pgfpathlineto{\pgfqpoint{0.697236in}{1.184202in}}%
\pgfpathlineto{\pgfqpoint{0.730794in}{1.114812in}}%
\pgfpathlineto{\pgfqpoint{0.774839in}{1.026516in}}%
\pgfpathlineto{\pgfqpoint{0.835662in}{0.907336in}}%
\pgfpathlineto{\pgfqpoint{0.911167in}{0.761770in}}%
\pgfpathlineto{\pgfqpoint{0.923751in}{0.737306in}}%
\pgfpathlineto{\pgfqpoint{0.936336in}{0.713371in}}%
\pgfpathlineto{\pgfqpoint{0.948920in}{0.690581in}}%
\pgfpathlineto{\pgfqpoint{0.967796in}{0.651376in}}%
\pgfpathlineto{\pgfqpoint{0.971991in}{0.645504in}}%
\pgfpathlineto{\pgfqpoint{0.978283in}{0.639267in}}%
\pgfpathlineto{\pgfqpoint{0.982478in}{0.634500in}}%
\pgfpathlineto{\pgfqpoint{0.986672in}{0.626257in}}%
\pgfpathlineto{\pgfqpoint{0.990867in}{0.615244in}}%
\pgfpathlineto{\pgfqpoint{0.995062in}{0.601390in}}%
\pgfpathlineto{\pgfqpoint{1.001354in}{0.575320in}}%
\pgfpathlineto{\pgfqpoint{1.005549in}{0.557112in}}%
\pgfpathlineto{\pgfqpoint{1.007646in}{0.549758in}}%
\pgfpathlineto{\pgfqpoint{1.009743in}{0.544705in}}%
\pgfpathlineto{\pgfqpoint{1.011841in}{0.542704in}}%
\pgfpathlineto{\pgfqpoint{1.013938in}{0.544411in}}%
\pgfpathlineto{\pgfqpoint{1.016035in}{0.549625in}}%
\pgfpathlineto{\pgfqpoint{1.024425in}{0.577878in}}%
\pgfpathlineto{\pgfqpoint{1.026522in}{0.582692in}}%
\pgfpathlineto{\pgfqpoint{1.028620in}{0.586081in}}%
\pgfpathlineto{\pgfqpoint{1.030717in}{0.587857in}}%
\pgfpathlineto{\pgfqpoint{1.032814in}{0.587683in}}%
\pgfpathlineto{\pgfqpoint{1.034912in}{0.584901in}}%
\pgfpathlineto{\pgfqpoint{1.037009in}{0.579145in}}%
\pgfpathlineto{\pgfqpoint{1.039106in}{0.570189in}}%
\pgfpathlineto{\pgfqpoint{1.043301in}{0.544015in}}%
\pgfpathlineto{\pgfqpoint{1.045398in}{0.547323in}}%
\pgfpathlineto{\pgfqpoint{1.049593in}{0.575543in}}%
\pgfpathlineto{\pgfqpoint{1.053788in}{0.592092in}}%
\pgfpathlineto{\pgfqpoint{1.055885in}{0.597191in}}%
\pgfpathlineto{\pgfqpoint{1.057983in}{0.600781in}}%
\pgfpathlineto{\pgfqpoint{1.060080in}{0.603095in}}%
\pgfpathlineto{\pgfqpoint{1.062177in}{0.604258in}}%
\pgfpathlineto{\pgfqpoint{1.064275in}{0.604321in}}%
\pgfpathlineto{\pgfqpoint{1.066372in}{0.603260in}}%
\pgfpathlineto{\pgfqpoint{1.068469in}{0.600976in}}%
\pgfpathlineto{\pgfqpoint{1.070567in}{0.597259in}}%
\pgfpathlineto{\pgfqpoint{1.072664in}{0.591720in}}%
\pgfpathlineto{\pgfqpoint{1.074762in}{0.583604in}}%
\pgfpathlineto{\pgfqpoint{1.076859in}{0.571251in}}%
\pgfpathlineto{\pgfqpoint{1.078956in}{0.549873in}}%
\pgfpathlineto{\pgfqpoint{1.081054in}{0.493255in}}%
\pgfpathlineto{\pgfqpoint{1.083151in}{0.539372in}}%
\pgfpathlineto{\pgfqpoint{1.085248in}{0.568614in}}%
\pgfpathlineto{\pgfqpoint{1.087346in}{0.584842in}}%
\pgfpathlineto{\pgfqpoint{1.089443in}{0.595751in}}%
\pgfpathlineto{\pgfqpoint{1.093638in}{0.609282in}}%
\pgfpathlineto{\pgfqpoint{1.095735in}{0.613259in}}%
\pgfpathlineto{\pgfqpoint{1.097832in}{0.615689in}}%
\pgfpathlineto{\pgfqpoint{1.099930in}{0.616564in}}%
\pgfpathlineto{\pgfqpoint{1.102027in}{0.615690in}}%
\pgfpathlineto{\pgfqpoint{1.104125in}{0.612592in}}%
\pgfpathlineto{\pgfqpoint{1.106222in}{0.606569in}}%
\pgfpathlineto{\pgfqpoint{1.108319in}{0.598461in}}%
\pgfpathlineto{\pgfqpoint{1.110417in}{0.593906in}}%
\pgfpathlineto{\pgfqpoint{1.112514in}{0.600064in}}%
\pgfpathlineto{\pgfqpoint{1.116709in}{0.621158in}}%
\pgfpathlineto{\pgfqpoint{1.118806in}{0.628774in}}%
\pgfpathlineto{\pgfqpoint{1.120903in}{0.634258in}}%
\pgfpathlineto{\pgfqpoint{1.123001in}{0.637912in}}%
\pgfpathlineto{\pgfqpoint{1.125098in}{0.639936in}}%
\pgfpathlineto{\pgfqpoint{1.127196in}{0.640415in}}%
\pgfpathlineto{\pgfqpoint{1.129293in}{0.639315in}}%
\pgfpathlineto{\pgfqpoint{1.131390in}{0.636460in}}%
\pgfpathlineto{\pgfqpoint{1.133488in}{0.631451in}}%
\pgfpathlineto{\pgfqpoint{1.135585in}{0.623448in}}%
\pgfpathlineto{\pgfqpoint{1.137682in}{0.610481in}}%
\pgfpathlineto{\pgfqpoint{1.139780in}{0.586220in}}%
\pgfpathlineto{\pgfqpoint{1.141877in}{0.512108in}}%
\pgfpathlineto{\pgfqpoint{1.143974in}{0.589469in}}%
\pgfpathlineto{\pgfqpoint{1.146072in}{0.615099in}}%
\pgfpathlineto{\pgfqpoint{1.148169in}{0.629752in}}%
\pgfpathlineto{\pgfqpoint{1.150267in}{0.639403in}}%
\pgfpathlineto{\pgfqpoint{1.152364in}{0.645871in}}%
\pgfpathlineto{\pgfqpoint{1.154461in}{0.649839in}}%
\pgfpathlineto{\pgfqpoint{1.156559in}{0.651462in}}%
\pgfpathlineto{\pgfqpoint{1.158656in}{0.650472in}}%
\pgfpathlineto{\pgfqpoint{1.160753in}{0.645992in}}%
\pgfpathlineto{\pgfqpoint{1.162851in}{0.638158in}}%
\pgfpathlineto{\pgfqpoint{1.164948in}{0.627381in}}%
\pgfpathlineto{\pgfqpoint{1.171240in}{0.586277in}}%
\pgfpathlineto{\pgfqpoint{1.175435in}{0.596800in}}%
\pgfpathlineto{\pgfqpoint{1.177532in}{0.598099in}}%
\pgfpathlineto{\pgfqpoint{1.179630in}{0.593628in}}%
\pgfpathlineto{\pgfqpoint{1.181727in}{0.581265in}}%
\pgfpathlineto{\pgfqpoint{1.183824in}{0.557891in}}%
\pgfpathlineto{\pgfqpoint{1.185922in}{0.569312in}}%
\pgfpathlineto{\pgfqpoint{1.188019in}{0.596174in}}%
\pgfpathlineto{\pgfqpoint{1.190116in}{0.612442in}}%
\pgfpathlineto{\pgfqpoint{1.192214in}{0.622199in}}%
\pgfpathlineto{\pgfqpoint{1.194311in}{0.627024in}}%
\pgfpathlineto{\pgfqpoint{1.196409in}{0.626754in}}%
\pgfpathlineto{\pgfqpoint{1.198506in}{0.620323in}}%
\pgfpathlineto{\pgfqpoint{1.200603in}{0.609975in}}%
\pgfpathlineto{\pgfqpoint{1.202701in}{0.594794in}}%
\pgfpathlineto{\pgfqpoint{1.204798in}{0.571500in}}%
\pgfpathlineto{\pgfqpoint{1.206895in}{0.541785in}}%
\pgfpathlineto{\pgfqpoint{1.211090in}{0.574294in}}%
\pgfpathlineto{\pgfqpoint{1.213187in}{0.580408in}}%
\pgfpathlineto{\pgfqpoint{1.215285in}{0.579511in}}%
\pgfpathlineto{\pgfqpoint{1.217382in}{0.571151in}}%
\pgfpathlineto{\pgfqpoint{1.219480in}{0.549383in}}%
\pgfpathlineto{\pgfqpoint{1.221577in}{0.505861in}}%
\pgfpathlineto{\pgfqpoint{1.223674in}{0.562596in}}%
\pgfpathlineto{\pgfqpoint{1.225772in}{0.586830in}}%
\pgfpathlineto{\pgfqpoint{1.227869in}{0.600376in}}%
\pgfpathlineto{\pgfqpoint{1.229966in}{0.607948in}}%
\pgfpathlineto{\pgfqpoint{1.232064in}{0.610526in}}%
\pgfpathlineto{\pgfqpoint{1.234161in}{0.607128in}}%
\pgfpathlineto{\pgfqpoint{1.236258in}{0.596490in}}%
\pgfpathlineto{\pgfqpoint{1.238356in}{0.575190in}}%
\pgfpathlineto{\pgfqpoint{1.240453in}{0.518887in}}%
\pgfpathlineto{\pgfqpoint{1.242551in}{0.564813in}}%
\pgfpathlineto{\pgfqpoint{1.244648in}{0.590984in}}%
\pgfpathlineto{\pgfqpoint{1.246745in}{0.604019in}}%
\pgfpathlineto{\pgfqpoint{1.248843in}{0.611096in}}%
\pgfpathlineto{\pgfqpoint{1.250940in}{0.614257in}}%
\pgfpathlineto{\pgfqpoint{1.253037in}{0.614145in}}%
\pgfpathlineto{\pgfqpoint{1.255135in}{0.610671in}}%
\pgfpathlineto{\pgfqpoint{1.257232in}{0.602913in}}%
\pgfpathlineto{\pgfqpoint{1.259329in}{0.587975in}}%
\pgfpathlineto{\pgfqpoint{1.261427in}{0.552779in}}%
\pgfpathlineto{\pgfqpoint{1.263524in}{0.542958in}}%
\pgfpathlineto{\pgfqpoint{1.265621in}{0.588865in}}%
\pgfpathlineto{\pgfqpoint{1.267719in}{0.608337in}}%
\pgfpathlineto{\pgfqpoint{1.269816in}{0.619596in}}%
\pgfpathlineto{\pgfqpoint{1.271914in}{0.625960in}}%
\pgfpathlineto{\pgfqpoint{1.274011in}{0.628180in}}%
\pgfpathlineto{\pgfqpoint{1.276108in}{0.625618in}}%
\pgfpathlineto{\pgfqpoint{1.278206in}{0.617139in}}%
\pgfpathlineto{\pgfqpoint{1.280303in}{0.602126in}}%
\pgfpathlineto{\pgfqpoint{1.282400in}{0.573232in}}%
\pgfpathlineto{\pgfqpoint{1.284498in}{0.521101in}}%
\pgfpathlineto{\pgfqpoint{1.286595in}{0.580605in}}%
\pgfpathlineto{\pgfqpoint{1.288692in}{0.600055in}}%
\pgfpathlineto{\pgfqpoint{1.290790in}{0.609306in}}%
\pgfpathlineto{\pgfqpoint{1.292887in}{0.613055in}}%
\pgfpathlineto{\pgfqpoint{1.294985in}{0.612523in}}%
\pgfpathlineto{\pgfqpoint{1.297082in}{0.607538in}}%
\pgfpathlineto{\pgfqpoint{1.299179in}{0.596292in}}%
\pgfpathlineto{\pgfqpoint{1.301277in}{0.571453in}}%
\pgfpathlineto{\pgfqpoint{1.303374in}{0.497944in}}%
\pgfpathlineto{\pgfqpoint{1.305471in}{0.583484in}}%
\pgfpathlineto{\pgfqpoint{1.307569in}{0.607486in}}%
\pgfpathlineto{\pgfqpoint{1.309666in}{0.620660in}}%
\pgfpathlineto{\pgfqpoint{1.311763in}{0.628077in}}%
\pgfpathlineto{\pgfqpoint{1.313861in}{0.630873in}}%
\pgfpathlineto{\pgfqpoint{1.315958in}{0.628433in}}%
\pgfpathlineto{\pgfqpoint{1.318056in}{0.619749in}}%
\pgfpathlineto{\pgfqpoint{1.320153in}{0.604338in}}%
\pgfpathlineto{\pgfqpoint{1.322250in}{0.574027in}}%
\pgfpathlineto{\pgfqpoint{1.324348in}{0.525127in}}%
\pgfpathlineto{\pgfqpoint{1.326445in}{0.585025in}}%
\pgfpathlineto{\pgfqpoint{1.328542in}{0.603920in}}%
\pgfpathlineto{\pgfqpoint{1.330640in}{0.612917in}}%
\pgfpathlineto{\pgfqpoint{1.332737in}{0.616500in}}%
\pgfpathlineto{\pgfqpoint{1.334834in}{0.615827in}}%
\pgfpathlineto{\pgfqpoint{1.336932in}{0.610682in}}%
\pgfpathlineto{\pgfqpoint{1.339029in}{0.599186in}}%
\pgfpathlineto{\pgfqpoint{1.341127in}{0.573637in}}%
\pgfpathlineto{\pgfqpoint{1.343224in}{0.508956in}}%
\pgfpathlineto{\pgfqpoint{1.345321in}{0.587782in}}%
\pgfpathlineto{\pgfqpoint{1.347419in}{0.611299in}}%
\pgfpathlineto{\pgfqpoint{1.349516in}{0.624263in}}%
\pgfpathlineto{\pgfqpoint{1.351613in}{0.631534in}}%
\pgfpathlineto{\pgfqpoint{1.353711in}{0.634186in}}%
\pgfpathlineto{\pgfqpoint{1.355808in}{0.631549in}}%
\pgfpathlineto{\pgfqpoint{1.357905in}{0.622684in}}%
\pgfpathlineto{\pgfqpoint{1.360003in}{0.606991in}}%
\pgfpathlineto{\pgfqpoint{1.362100in}{0.575767in}}%
\pgfpathlineto{\pgfqpoint{1.364198in}{0.530180in}}%
\pgfpathlineto{\pgfqpoint{1.366295in}{0.588908in}}%
\pgfpathlineto{\pgfqpoint{1.368392in}{0.607424in}}%
\pgfpathlineto{\pgfqpoint{1.370490in}{0.616224in}}%
\pgfpathlineto{\pgfqpoint{1.372587in}{0.619654in}}%
\pgfpathlineto{\pgfqpoint{1.374684in}{0.618819in}}%
\pgfpathlineto{\pgfqpoint{1.376782in}{0.613460in}}%
\pgfpathlineto{\pgfqpoint{1.378879in}{0.601576in}}%
\pgfpathlineto{\pgfqpoint{1.380976in}{0.574822in}}%
\pgfpathlineto{\pgfqpoint{1.383074in}{0.522727in}}%
\pgfpathlineto{\pgfqpoint{1.385171in}{0.592517in}}%
\pgfpathlineto{\pgfqpoint{1.387269in}{0.615277in}}%
\pgfpathlineto{\pgfqpoint{1.389366in}{0.627910in}}%
\pgfpathlineto{\pgfqpoint{1.391463in}{0.634945in}}%
\pgfpathlineto{\pgfqpoint{1.393561in}{0.637359in}}%
\pgfpathlineto{\pgfqpoint{1.395658in}{0.634381in}}%
\pgfpathlineto{\pgfqpoint{1.397755in}{0.625426in}}%
\pgfpathlineto{\pgfqpoint{1.399853in}{0.610210in}}%
\pgfpathlineto{\pgfqpoint{1.401950in}{0.581595in}}%
\pgfpathlineto{\pgfqpoint{1.404047in}{0.503027in}}%
\pgfpathlineto{\pgfqpoint{1.406145in}{0.584416in}}%
\pgfpathlineto{\pgfqpoint{1.408242in}{0.604393in}}%
\pgfpathlineto{\pgfqpoint{1.410340in}{0.613378in}}%
\pgfpathlineto{\pgfqpoint{1.412437in}{0.616514in}}%
\pgfpathlineto{\pgfqpoint{1.414534in}{0.614970in}}%
\pgfpathlineto{\pgfqpoint{1.416632in}{0.608238in}}%
\pgfpathlineto{\pgfqpoint{1.418729in}{0.593271in}}%
\pgfpathlineto{\pgfqpoint{1.420826in}{0.553621in}}%
\pgfpathlineto{\pgfqpoint{1.422924in}{0.560557in}}%
\pgfpathlineto{\pgfqpoint{1.425021in}{0.600949in}}%
\pgfpathlineto{\pgfqpoint{1.427118in}{0.619286in}}%
\pgfpathlineto{\pgfqpoint{1.429216in}{0.629746in}}%
\pgfpathlineto{\pgfqpoint{1.431313in}{0.635075in}}%
\pgfpathlineto{\pgfqpoint{1.433410in}{0.635556in}}%
\pgfpathlineto{\pgfqpoint{1.435508in}{0.629916in}}%
\pgfpathlineto{\pgfqpoint{1.437605in}{0.618970in}}%
\pgfpathlineto{\pgfqpoint{1.439703in}{0.600756in}}%
\pgfpathlineto{\pgfqpoint{1.441800in}{0.563343in}}%
\pgfpathlineto{\pgfqpoint{1.443897in}{0.543116in}}%
\pgfpathlineto{\pgfqpoint{1.445995in}{0.587263in}}%
\pgfpathlineto{\pgfqpoint{1.448092in}{0.602571in}}%
\pgfpathlineto{\pgfqpoint{1.450189in}{0.609031in}}%
\pgfpathlineto{\pgfqpoint{1.452287in}{0.609956in}}%
\pgfpathlineto{\pgfqpoint{1.454384in}{0.605753in}}%
\pgfpathlineto{\pgfqpoint{1.456481in}{0.594716in}}%
\pgfpathlineto{\pgfqpoint{1.458579in}{0.568766in}}%
\pgfpathlineto{\pgfqpoint{1.460676in}{0.514514in}}%
\pgfpathlineto{\pgfqpoint{1.462774in}{0.586349in}}%
\pgfpathlineto{\pgfqpoint{1.464871in}{0.609429in}}%
\pgfpathlineto{\pgfqpoint{1.466968in}{0.622181in}}%
\pgfpathlineto{\pgfqpoint{1.469066in}{0.629176in}}%
\pgfpathlineto{\pgfqpoint{1.471163in}{0.631324in}}%
\pgfpathlineto{\pgfqpoint{1.473260in}{0.627643in}}%
\pgfpathlineto{\pgfqpoint{1.475358in}{0.618126in}}%
\pgfpathlineto{\pgfqpoint{1.477455in}{0.602157in}}%
\pgfpathlineto{\pgfqpoint{1.479552in}{0.571692in}}%
\pgfpathlineto{\pgfqpoint{1.481650in}{0.510415in}}%
\pgfpathlineto{\pgfqpoint{1.483747in}{0.579181in}}%
\pgfpathlineto{\pgfqpoint{1.485845in}{0.597863in}}%
\pgfpathlineto{\pgfqpoint{1.487942in}{0.606137in}}%
\pgfpathlineto{\pgfqpoint{1.490039in}{0.608631in}}%
\pgfpathlineto{\pgfqpoint{1.492137in}{0.606311in}}%
\pgfpathlineto{\pgfqpoint{1.494234in}{0.598363in}}%
\pgfpathlineto{\pgfqpoint{1.496331in}{0.580673in}}%
\pgfpathlineto{\pgfqpoint{1.498429in}{0.523463in}}%
\pgfpathlineto{\pgfqpoint{1.502623in}{0.598690in}}%
\pgfpathlineto{\pgfqpoint{1.504721in}{0.614862in}}%
\pgfpathlineto{\pgfqpoint{1.506818in}{0.624091in}}%
\pgfpathlineto{\pgfqpoint{1.508916in}{0.628369in}}%
\pgfpathlineto{\pgfqpoint{1.511013in}{0.627569in}}%
\pgfpathlineto{\pgfqpoint{1.513110in}{0.620375in}}%
\pgfpathlineto{\pgfqpoint{1.515208in}{0.607057in}}%
\pgfpathlineto{\pgfqpoint{1.517305in}{0.582492in}}%
\pgfpathlineto{\pgfqpoint{1.519402in}{0.483718in}}%
\pgfpathlineto{\pgfqpoint{1.521500in}{0.573187in}}%
\pgfpathlineto{\pgfqpoint{1.523597in}{0.597209in}}%
\pgfpathlineto{\pgfqpoint{1.525694in}{0.608141in}}%
\pgfpathlineto{\pgfqpoint{1.527792in}{0.612893in}}%
\pgfpathlineto{\pgfqpoint{1.529889in}{0.613168in}}%
\pgfpathlineto{\pgfqpoint{1.531987in}{0.609044in}}%
\pgfpathlineto{\pgfqpoint{1.534084in}{0.599085in}}%
\pgfpathlineto{\pgfqpoint{1.536181in}{0.577573in}}%
\pgfpathlineto{\pgfqpoint{1.538279in}{0.460729in}}%
\pgfpathlineto{\pgfqpoint{1.540376in}{0.578239in}}%
\pgfpathlineto{\pgfqpoint{1.542473in}{0.605243in}}%
\pgfpathlineto{\pgfqpoint{1.544571in}{0.619589in}}%
\pgfpathlineto{\pgfqpoint{1.546668in}{0.627756in}}%
\pgfpathlineto{\pgfqpoint{1.548765in}{0.631216in}}%
\pgfpathlineto{\pgfqpoint{1.550863in}{0.629597in}}%
\pgfpathlineto{\pgfqpoint{1.552960in}{0.621621in}}%
\pgfpathlineto{\pgfqpoint{1.555058in}{0.606992in}}%
\pgfpathlineto{\pgfqpoint{1.557155in}{0.578403in}}%
\pgfpathlineto{\pgfqpoint{1.559252in}{0.505664in}}%
\pgfpathlineto{\pgfqpoint{1.561350in}{0.584614in}}%
\pgfpathlineto{\pgfqpoint{1.563447in}{0.604932in}}%
\pgfpathlineto{\pgfqpoint{1.565544in}{0.614639in}}%
\pgfpathlineto{\pgfqpoint{1.567642in}{0.618805in}}%
\pgfpathlineto{\pgfqpoint{1.569739in}{0.618765in}}%
\pgfpathlineto{\pgfqpoint{1.571836in}{0.614490in}}%
\pgfpathlineto{\pgfqpoint{1.573934in}{0.604537in}}%
\pgfpathlineto{\pgfqpoint{1.576031in}{0.583390in}}%
\pgfpathlineto{\pgfqpoint{1.578129in}{0.483886in}}%
\pgfpathlineto{\pgfqpoint{1.580226in}{0.581892in}}%
\pgfpathlineto{\pgfqpoint{1.582323in}{0.609450in}}%
\pgfpathlineto{\pgfqpoint{1.584421in}{0.624006in}}%
\pgfpathlineto{\pgfqpoint{1.586518in}{0.632339in}}%
\pgfpathlineto{\pgfqpoint{1.588615in}{0.636010in}}%
\pgfpathlineto{\pgfqpoint{1.590713in}{0.634750in}}%
\pgfpathlineto{\pgfqpoint{1.592810in}{0.627298in}}%
\pgfpathlineto{\pgfqpoint{1.594907in}{0.613835in}}%
\pgfpathlineto{\pgfqpoint{1.597005in}{0.589186in}}%
\pgfpathlineto{\pgfqpoint{1.599102in}{0.491069in}}%
\pgfpathlineto{\pgfqpoint{1.601199in}{0.579340in}}%
\pgfpathlineto{\pgfqpoint{1.603297in}{0.603296in}}%
\pgfpathlineto{\pgfqpoint{1.605394in}{0.614084in}}%
\pgfpathlineto{\pgfqpoint{1.607492in}{0.618635in}}%
\pgfpathlineto{\pgfqpoint{1.609589in}{0.618627in}}%
\pgfpathlineto{\pgfqpoint{1.611686in}{0.614058in}}%
\pgfpathlineto{\pgfqpoint{1.613784in}{0.603255in}}%
\pgfpathlineto{\pgfqpoint{1.615881in}{0.579293in}}%
\pgfpathlineto{\pgfqpoint{1.617978in}{0.490331in}}%
\pgfpathlineto{\pgfqpoint{1.620076in}{0.589013in}}%
\pgfpathlineto{\pgfqpoint{1.622173in}{0.613638in}}%
\pgfpathlineto{\pgfqpoint{1.624270in}{0.627052in}}%
\pgfpathlineto{\pgfqpoint{1.626368in}{0.634588in}}%
\pgfpathlineto{\pgfqpoint{1.628465in}{0.637427in}}%
\pgfpathlineto{\pgfqpoint{1.630563in}{0.634942in}}%
\pgfpathlineto{\pgfqpoint{1.632660in}{0.626453in}}%
\pgfpathlineto{\pgfqpoint{1.634757in}{0.612285in}}%
\pgfpathlineto{\pgfqpoint{1.636855in}{0.587049in}}%
\pgfpathlineto{\pgfqpoint{1.638952in}{0.489553in}}%
\pgfpathlineto{\pgfqpoint{1.641049in}{0.575719in}}%
\pgfpathlineto{\pgfqpoint{1.643147in}{0.599045in}}%
\pgfpathlineto{\pgfqpoint{1.645244in}{0.609022in}}%
\pgfpathlineto{\pgfqpoint{1.647341in}{0.612514in}}%
\pgfpathlineto{\pgfqpoint{1.649439in}{0.610990in}}%
\pgfpathlineto{\pgfqpoint{1.651536in}{0.603929in}}%
\pgfpathlineto{\pgfqpoint{1.653634in}{0.587851in}}%
\pgfpathlineto{\pgfqpoint{1.655731in}{0.541136in}}%
\pgfpathlineto{\pgfqpoint{1.657828in}{0.565070in}}%
\pgfpathlineto{\pgfqpoint{1.659926in}{0.601003in}}%
\pgfpathlineto{\pgfqpoint{1.662023in}{0.618235in}}%
\pgfpathlineto{\pgfqpoint{1.664120in}{0.628067in}}%
\pgfpathlineto{\pgfqpoint{1.666218in}{0.632828in}}%
\pgfpathlineto{\pgfqpoint{1.668315in}{0.632559in}}%
\pgfpathlineto{\pgfqpoint{1.670412in}{0.626010in}}%
\pgfpathlineto{\pgfqpoint{1.672510in}{0.614202in}}%
\pgfpathlineto{\pgfqpoint{1.674607in}{0.594336in}}%
\pgfpathlineto{\pgfqpoint{1.676705in}{0.549811in}}%
\pgfpathlineto{\pgfqpoint{1.678802in}{0.552263in}}%
\pgfpathlineto{\pgfqpoint{1.680899in}{0.587838in}}%
\pgfpathlineto{\pgfqpoint{1.682997in}{0.601344in}}%
\pgfpathlineto{\pgfqpoint{1.685094in}{0.606847in}}%
\pgfpathlineto{\pgfqpoint{1.687191in}{0.606995in}}%
\pgfpathlineto{\pgfqpoint{1.689289in}{0.601903in}}%
\pgfpathlineto{\pgfqpoint{1.691386in}{0.589372in}}%
\pgfpathlineto{\pgfqpoint{1.695581in}{0.535919in}}%
\pgfpathlineto{\pgfqpoint{1.697678in}{0.588575in}}%
\pgfpathlineto{\pgfqpoint{1.699776in}{0.609337in}}%
\pgfpathlineto{\pgfqpoint{1.701873in}{0.621017in}}%
\pgfpathlineto{\pgfqpoint{1.703970in}{0.627228in}}%
\pgfpathlineto{\pgfqpoint{1.706068in}{0.628561in}}%
\pgfpathlineto{\pgfqpoint{1.708165in}{0.623812in}}%
\pgfpathlineto{\pgfqpoint{1.710262in}{0.613222in}}%
\pgfpathlineto{\pgfqpoint{1.712360in}{0.594919in}}%
\pgfpathlineto{\pgfqpoint{1.714457in}{0.555193in}}%
\pgfpathlineto{\pgfqpoint{1.716554in}{0.545737in}}%
\pgfpathlineto{\pgfqpoint{1.718652in}{0.586335in}}%
\pgfpathlineto{\pgfqpoint{1.720749in}{0.601539in}}%
\pgfpathlineto{\pgfqpoint{1.722847in}{0.608471in}}%
\pgfpathlineto{\pgfqpoint{1.724944in}{0.610283in}}%
\pgfpathlineto{\pgfqpoint{1.727041in}{0.607590in}}%
\pgfpathlineto{\pgfqpoint{1.729139in}{0.599482in}}%
\pgfpathlineto{\pgfqpoint{1.731236in}{0.581931in}}%
\pgfpathlineto{\pgfqpoint{1.733333in}{0.527494in}}%
\pgfpathlineto{\pgfqpoint{1.737528in}{0.598354in}}%
\pgfpathlineto{\pgfqpoint{1.739625in}{0.614728in}}%
\pgfpathlineto{\pgfqpoint{1.741723in}{0.624103in}}%
\pgfpathlineto{\pgfqpoint{1.743820in}{0.628571in}}%
\pgfpathlineto{\pgfqpoint{1.745918in}{0.628105in}}%
\pgfpathlineto{\pgfqpoint{1.748015in}{0.621314in}}%
\pgfpathlineto{\pgfqpoint{1.750112in}{0.608168in}}%
\pgfpathlineto{\pgfqpoint{1.752210in}{0.583449in}}%
\pgfpathlineto{\pgfqpoint{1.754307in}{0.470314in}}%
\pgfpathlineto{\pgfqpoint{1.756404in}{0.576517in}}%
\pgfpathlineto{\pgfqpoint{1.758502in}{0.600301in}}%
\pgfpathlineto{\pgfqpoint{1.760599in}{0.611382in}}%
\pgfpathlineto{\pgfqpoint{1.762696in}{0.616440in}}%
\pgfpathlineto{\pgfqpoint{1.764794in}{0.617188in}}%
\pgfpathlineto{\pgfqpoint{1.766891in}{0.613818in}}%
\pgfpathlineto{\pgfqpoint{1.768988in}{0.605252in}}%
\pgfpathlineto{\pgfqpoint{1.771086in}{0.587391in}}%
\pgfpathlineto{\pgfqpoint{1.773183in}{0.532883in}}%
\pgfpathlineto{\pgfqpoint{1.777378in}{0.603236in}}%
\pgfpathlineto{\pgfqpoint{1.779475in}{0.619584in}}%
\pgfpathlineto{\pgfqpoint{1.781573in}{0.629002in}}%
\pgfpathlineto{\pgfqpoint{1.783670in}{0.633618in}}%
\pgfpathlineto{\pgfqpoint{1.785767in}{0.633507in}}%
\pgfpathlineto{\pgfqpoint{1.787865in}{0.627303in}}%
\pgfpathlineto{\pgfqpoint{1.789962in}{0.614991in}}%
\pgfpathlineto{\pgfqpoint{1.792059in}{0.592476in}}%
\pgfpathlineto{\pgfqpoint{1.794157in}{0.522780in}}%
\pgfpathlineto{\pgfqpoint{1.796254in}{0.575188in}}%
\pgfpathlineto{\pgfqpoint{1.798352in}{0.602118in}}%
\pgfpathlineto{\pgfqpoint{1.800449in}{0.614181in}}%
\pgfpathlineto{\pgfqpoint{1.802546in}{0.619735in}}%
\pgfpathlineto{\pgfqpoint{1.804644in}{0.620801in}}%
\pgfpathlineto{\pgfqpoint{1.806741in}{0.617679in}}%
\pgfpathlineto{\pgfqpoint{1.808838in}{0.609361in}}%
\pgfpathlineto{\pgfqpoint{1.810936in}{0.591896in}}%
\pgfpathlineto{\pgfqpoint{1.813033in}{0.539807in}}%
\pgfpathlineto{\pgfqpoint{1.817228in}{0.606738in}}%
\pgfpathlineto{\pgfqpoint{1.819325in}{0.623313in}}%
\pgfpathlineto{\pgfqpoint{1.821423in}{0.632847in}}%
\pgfpathlineto{\pgfqpoint{1.823520in}{0.637534in}}%
\pgfpathlineto{\pgfqpoint{1.825617in}{0.637463in}}%
\pgfpathlineto{\pgfqpoint{1.827715in}{0.631334in}}%
\pgfpathlineto{\pgfqpoint{1.829812in}{0.619718in}}%
\pgfpathlineto{\pgfqpoint{1.831909in}{0.599799in}}%
\pgfpathlineto{\pgfqpoint{1.834007in}{0.553664in}}%
\pgfpathlineto{\pgfqpoint{1.836104in}{0.561362in}}%
\pgfpathlineto{\pgfqpoint{1.838201in}{0.595812in}}%
\pgfpathlineto{\pgfqpoint{1.840299in}{0.609373in}}%
\pgfpathlineto{\pgfqpoint{1.842396in}{0.615205in}}%
\pgfpathlineto{\pgfqpoint{1.844494in}{0.615926in}}%
\pgfpathlineto{\pgfqpoint{1.846591in}{0.611819in}}%
\pgfpathlineto{\pgfqpoint{1.848688in}{0.601317in}}%
\pgfpathlineto{\pgfqpoint{1.850786in}{0.577510in}}%
\pgfpathlineto{\pgfqpoint{1.852883in}{0.490395in}}%
\pgfpathlineto{\pgfqpoint{1.854980in}{0.587761in}}%
\pgfpathlineto{\pgfqpoint{1.857078in}{0.612395in}}%
\pgfpathlineto{\pgfqpoint{1.859175in}{0.625809in}}%
\pgfpathlineto{\pgfqpoint{1.861272in}{0.633296in}}%
\pgfpathlineto{\pgfqpoint{1.863370in}{0.635996in}}%
\pgfpathlineto{\pgfqpoint{1.865467in}{0.633177in}}%
\pgfpathlineto{\pgfqpoint{1.867565in}{0.624465in}}%
\pgfpathlineto{\pgfqpoint{1.869662in}{0.610195in}}%
\pgfpathlineto{\pgfqpoint{1.871759in}{0.585104in}}%
\pgfpathlineto{\pgfqpoint{1.873857in}{0.495051in}}%
\pgfpathlineto{\pgfqpoint{1.875954in}{0.571735in}}%
\pgfpathlineto{\pgfqpoint{1.878051in}{0.595285in}}%
\pgfpathlineto{\pgfqpoint{1.880149in}{0.605111in}}%
\pgfpathlineto{\pgfqpoint{1.882246in}{0.608279in}}%
\pgfpathlineto{\pgfqpoint{1.884343in}{0.606210in}}%
\pgfpathlineto{\pgfqpoint{1.886441in}{0.598145in}}%
\pgfpathlineto{\pgfqpoint{1.888538in}{0.579560in}}%
\pgfpathlineto{\pgfqpoint{1.890636in}{0.511016in}}%
\pgfpathlineto{\pgfqpoint{1.892733in}{0.571473in}}%
\pgfpathlineto{\pgfqpoint{1.894830in}{0.601756in}}%
\pgfpathlineto{\pgfqpoint{1.896928in}{0.617290in}}%
\pgfpathlineto{\pgfqpoint{1.899025in}{0.626107in}}%
\pgfpathlineto{\pgfqpoint{1.901122in}{0.629952in}}%
\pgfpathlineto{\pgfqpoint{1.903220in}{0.628495in}}%
\pgfpathlineto{\pgfqpoint{1.905317in}{0.620707in}}%
\pgfpathlineto{\pgfqpoint{1.907414in}{0.607309in}}%
\pgfpathlineto{\pgfqpoint{1.909512in}{0.583660in}}%
\pgfpathlineto{\pgfqpoint{1.911609in}{0.509584in}}%
\pgfpathlineto{\pgfqpoint{1.913707in}{0.566360in}}%
\pgfpathlineto{\pgfqpoint{1.915804in}{0.591920in}}%
\pgfpathlineto{\pgfqpoint{1.917901in}{0.602831in}}%
\pgfpathlineto{\pgfqpoint{1.919999in}{0.607053in}}%
\pgfpathlineto{\pgfqpoint{1.922096in}{0.606340in}}%
\pgfpathlineto{\pgfqpoint{1.924193in}{0.600476in}}%
\pgfpathlineto{\pgfqpoint{1.926291in}{0.586903in}}%
\pgfpathlineto{\pgfqpoint{1.928388in}{0.552808in}}%
\pgfpathlineto{\pgfqpoint{1.930485in}{0.542936in}}%
\pgfpathlineto{\pgfqpoint{1.932583in}{0.589568in}}%
\pgfpathlineto{\pgfqpoint{1.934680in}{0.609244in}}%
\pgfpathlineto{\pgfqpoint{1.936777in}{0.620391in}}%
\pgfpathlineto{\pgfqpoint{1.938875in}{0.626230in}}%
\pgfpathlineto{\pgfqpoint{1.940972in}{0.627223in}}%
\pgfpathlineto{\pgfqpoint{1.943070in}{0.622101in}}%
\pgfpathlineto{\pgfqpoint{1.945167in}{0.610874in}}%
\pgfpathlineto{\pgfqpoint{1.947264in}{0.590717in}}%
\pgfpathlineto{\pgfqpoint{1.949362in}{0.539871in}}%
\pgfpathlineto{\pgfqpoint{1.951459in}{0.560303in}}%
\pgfpathlineto{\pgfqpoint{1.953556in}{0.592392in}}%
\pgfpathlineto{\pgfqpoint{1.955654in}{0.605973in}}%
\pgfpathlineto{\pgfqpoint{1.957751in}{0.612375in}}%
\pgfpathlineto{\pgfqpoint{1.959848in}{0.614103in}}%
\pgfpathlineto{\pgfqpoint{1.961946in}{0.611662in}}%
\pgfpathlineto{\pgfqpoint{1.964043in}{0.604282in}}%
\pgfpathlineto{\pgfqpoint{1.966141in}{0.588745in}}%
\pgfpathlineto{\pgfqpoint{1.968238in}{0.547950in}}%
\pgfpathlineto{\pgfqpoint{1.970335in}{0.556936in}}%
\pgfpathlineto{\pgfqpoint{1.972433in}{0.596512in}}%
\pgfpathlineto{\pgfqpoint{1.974530in}{0.614621in}}%
\pgfpathlineto{\pgfqpoint{1.976627in}{0.625002in}}%
\pgfpathlineto{\pgfqpoint{1.978725in}{0.630373in}}%
\pgfpathlineto{\pgfqpoint{1.980822in}{0.631070in}}%
\pgfpathlineto{\pgfqpoint{1.982919in}{0.625807in}}%
\pgfpathlineto{\pgfqpoint{1.985017in}{0.614284in}}%
\pgfpathlineto{\pgfqpoint{1.987114in}{0.593090in}}%
\pgfpathlineto{\pgfqpoint{1.989212in}{0.533454in}}%
\pgfpathlineto{\pgfqpoint{1.993406in}{0.599989in}}%
\pgfpathlineto{\pgfqpoint{1.995504in}{0.612946in}}%
\pgfpathlineto{\pgfqpoint{1.997601in}{0.619182in}}%
\pgfpathlineto{\pgfqpoint{1.999698in}{0.620956in}}%
\pgfpathlineto{\pgfqpoint{2.001796in}{0.618744in}}%
\pgfpathlineto{\pgfqpoint{2.003893in}{0.611879in}}%
\pgfpathlineto{\pgfqpoint{2.005990in}{0.597600in}}%
\pgfpathlineto{\pgfqpoint{2.008088in}{0.562977in}}%
\pgfpathlineto{\pgfqpoint{2.010185in}{0.552862in}}%
\pgfpathlineto{\pgfqpoint{2.012283in}{0.599215in}}%
\pgfpathlineto{\pgfqpoint{2.014380in}{0.618761in}}%
\pgfpathlineto{\pgfqpoint{2.016477in}{0.629907in}}%
\pgfpathlineto{\pgfqpoint{2.018575in}{0.635915in}}%
\pgfpathlineto{\pgfqpoint{2.020672in}{0.637371in}}%
\pgfpathlineto{\pgfqpoint{2.022769in}{0.633202in}}%
\pgfpathlineto{\pgfqpoint{2.024867in}{0.623171in}}%
\pgfpathlineto{\pgfqpoint{2.026964in}{0.605782in}}%
\pgfpathlineto{\pgfqpoint{2.029061in}{0.569380in}}%
\pgfpathlineto{\pgfqpoint{2.031159in}{0.547953in}}%
\pgfpathlineto{\pgfqpoint{2.033256in}{0.593805in}}%
\pgfpathlineto{\pgfqpoint{2.035354in}{0.610075in}}%
\pgfpathlineto{\pgfqpoint{2.037451in}{0.617609in}}%
\pgfpathlineto{\pgfqpoint{2.039548in}{0.619939in}}%
\pgfpathlineto{\pgfqpoint{2.041646in}{0.617843in}}%
\pgfpathlineto{\pgfqpoint{2.043743in}{0.610639in}}%
\pgfpathlineto{\pgfqpoint{2.045840in}{0.595073in}}%
\pgfpathlineto{\pgfqpoint{2.047938in}{0.553240in}}%
\pgfpathlineto{\pgfqpoint{2.050035in}{0.565315in}}%
\pgfpathlineto{\pgfqpoint{2.052132in}{0.604081in}}%
\pgfpathlineto{\pgfqpoint{2.054230in}{0.622007in}}%
\pgfpathlineto{\pgfqpoint{2.056327in}{0.632263in}}%
\pgfpathlineto{\pgfqpoint{2.058425in}{0.637473in}}%
\pgfpathlineto{\pgfqpoint{2.060522in}{0.637894in}}%
\pgfpathlineto{\pgfqpoint{2.062619in}{0.632262in}}%
\pgfpathlineto{\pgfqpoint{2.064717in}{0.621404in}}%
\pgfpathlineto{\pgfqpoint{2.066814in}{0.603476in}}%
\pgfpathlineto{\pgfqpoint{2.071009in}{0.540382in}}%
\pgfpathlineto{\pgfqpoint{2.073106in}{0.587747in}}%
\pgfpathlineto{\pgfqpoint{2.075203in}{0.603460in}}%
\pgfpathlineto{\pgfqpoint{2.077301in}{0.610014in}}%
\pgfpathlineto{\pgfqpoint{2.079398in}{0.610902in}}%
\pgfpathlineto{\pgfqpoint{2.081496in}{0.606540in}}%
\pgfpathlineto{\pgfqpoint{2.083593in}{0.595098in}}%
\pgfpathlineto{\pgfqpoint{2.085690in}{0.567674in}}%
\pgfpathlineto{\pgfqpoint{2.087788in}{0.526111in}}%
\pgfpathlineto{\pgfqpoint{2.089885in}{0.589570in}}%
\pgfpathlineto{\pgfqpoint{2.091982in}{0.611814in}}%
\pgfpathlineto{\pgfqpoint{2.094080in}{0.624192in}}%
\pgfpathlineto{\pgfqpoint{2.096177in}{0.630906in}}%
\pgfpathlineto{\pgfqpoint{2.098274in}{0.632743in}}%
\pgfpathlineto{\pgfqpoint{2.100372in}{0.628621in}}%
\pgfpathlineto{\pgfqpoint{2.102469in}{0.618940in}}%
\pgfpathlineto{\pgfqpoint{2.104566in}{0.603080in}}%
\pgfpathlineto{\pgfqpoint{2.106664in}{0.573715in}}%
\pgfpathlineto{\pgfqpoint{2.108761in}{0.492587in}}%
\pgfpathlineto{\pgfqpoint{2.110859in}{0.576036in}}%
\pgfpathlineto{\pgfqpoint{2.112956in}{0.595296in}}%
\pgfpathlineto{\pgfqpoint{2.115053in}{0.603469in}}%
\pgfpathlineto{\pgfqpoint{2.117151in}{0.605547in}}%
\pgfpathlineto{\pgfqpoint{2.119248in}{0.602437in}}%
\pgfpathlineto{\pgfqpoint{2.121345in}{0.592929in}}%
\pgfpathlineto{\pgfqpoint{2.123443in}{0.570951in}}%
\pgfpathlineto{\pgfqpoint{2.125540in}{0.409298in}}%
\pgfpathlineto{\pgfqpoint{2.127637in}{0.576375in}}%
\pgfpathlineto{\pgfqpoint{2.129735in}{0.602448in}}%
\pgfpathlineto{\pgfqpoint{2.131832in}{0.616449in}}%
\pgfpathlineto{\pgfqpoint{2.133930in}{0.624288in}}%
\pgfpathlineto{\pgfqpoint{2.136027in}{0.627236in}}%
\pgfpathlineto{\pgfqpoint{2.138124in}{0.624608in}}%
\pgfpathlineto{\pgfqpoint{2.140222in}{0.615690in}}%
\pgfpathlineto{\pgfqpoint{2.142319in}{0.600202in}}%
\pgfpathlineto{\pgfqpoint{2.144416in}{0.570136in}}%
\pgfpathlineto{\pgfqpoint{2.146514in}{0.509046in}}%
\pgfpathlineto{\pgfqpoint{2.148611in}{0.578584in}}%
\pgfpathlineto{\pgfqpoint{2.150708in}{0.597761in}}%
\pgfpathlineto{\pgfqpoint{2.152806in}{0.606637in}}%
\pgfpathlineto{\pgfqpoint{2.154903in}{0.609921in}}%
\pgfpathlineto{\pgfqpoint{2.157001in}{0.608751in}}%
\pgfpathlineto{\pgfqpoint{2.159098in}{0.602756in}}%
\pgfpathlineto{\pgfqpoint{2.161195in}{0.589493in}}%
\pgfpathlineto{\pgfqpoint{2.163293in}{0.557631in}}%
\pgfpathlineto{\pgfqpoint{2.165390in}{0.537711in}}%
\pgfpathlineto{\pgfqpoint{2.167487in}{0.588689in}}%
\pgfpathlineto{\pgfqpoint{2.169585in}{0.609095in}}%
\pgfpathlineto{\pgfqpoint{2.171682in}{0.620624in}}%
\pgfpathlineto{\pgfqpoint{2.173779in}{0.626821in}}%
\pgfpathlineto{\pgfqpoint{2.175877in}{0.628316in}}%
\pgfpathlineto{\pgfqpoint{2.177974in}{0.623982in}}%
\pgfpathlineto{\pgfqpoint{2.180072in}{0.613312in}}%
\pgfpathlineto{\pgfqpoint{2.182169in}{0.593749in}}%
\pgfpathlineto{\pgfqpoint{2.184266in}{0.544610in}}%
\pgfpathlineto{\pgfqpoint{2.186364in}{0.562441in}}%
\pgfpathlineto{\pgfqpoint{2.188461in}{0.595527in}}%
\pgfpathlineto{\pgfqpoint{2.190558in}{0.609647in}}%
\pgfpathlineto{\pgfqpoint{2.192656in}{0.616578in}}%
\pgfpathlineto{\pgfqpoint{2.194753in}{0.618936in}}%
\pgfpathlineto{\pgfqpoint{2.196850in}{0.617364in}}%
\pgfpathlineto{\pgfqpoint{2.198948in}{0.611390in}}%
\pgfpathlineto{\pgfqpoint{2.201045in}{0.598814in}}%
\pgfpathlineto{\pgfqpoint{2.203143in}{0.570490in}}%
\pgfpathlineto{\pgfqpoint{2.205240in}{0.529213in}}%
\pgfpathlineto{\pgfqpoint{2.207337in}{0.591766in}}%
\pgfpathlineto{\pgfqpoint{2.209435in}{0.613731in}}%
\pgfpathlineto{\pgfqpoint{2.211532in}{0.626030in}}%
\pgfpathlineto{\pgfqpoint{2.213629in}{0.632892in}}%
\pgfpathlineto{\pgfqpoint{2.215727in}{0.635236in}}%
\pgfpathlineto{\pgfqpoint{2.217824in}{0.632313in}}%
\pgfpathlineto{\pgfqpoint{2.222019in}{0.615408in}}%
\pgfpathlineto{\pgfqpoint{2.226214in}{0.589430in}}%
\pgfpathlineto{\pgfqpoint{2.230408in}{0.551389in}}%
\pgfpathlineto{\pgfqpoint{2.232506in}{0.531114in}}%
\pgfpathlineto{\pgfqpoint{2.234603in}{0.523148in}}%
\pgfpathlineto{\pgfqpoint{2.236700in}{0.534327in}}%
\pgfpathlineto{\pgfqpoint{2.240895in}{0.567592in}}%
\pgfpathlineto{\pgfqpoint{2.242992in}{0.579213in}}%
\pgfpathlineto{\pgfqpoint{2.245090in}{0.586774in}}%
\pgfpathlineto{\pgfqpoint{2.247187in}{0.589959in}}%
\pgfpathlineto{\pgfqpoint{2.249285in}{0.587339in}}%
\pgfpathlineto{\pgfqpoint{2.249285in}{0.587339in}}%
\pgfusepath{stroke}%
\end{pgfscope}%
\begin{pgfscope}%
\pgfpathrectangle{\pgfqpoint{0.487500in}{0.320625in}}{\pgfqpoint{1.845679in}{1.164375in}}%
\pgfusepath{clip}%
\pgfsetbuttcap%
\pgfsetroundjoin%
\pgfsetlinewidth{0.752812pt}%
\definecolor{currentstroke}{rgb}{0.000000,0.529412,0.423529}%
\pgfsetstrokecolor{currentstroke}%
\pgfsetstrokeopacity{0.650000}%
\pgfsetdash{{2.775000pt}{1.200000pt}}{0.000000pt}%
\pgfpathmoveto{\pgfqpoint{0.571395in}{1.468525in}}%
\pgfpathlineto{\pgfqpoint{1.216932in}{0.315625in}}%
\pgfpathlineto{\pgfqpoint{1.216932in}{0.315625in}}%
\pgfusepath{stroke}%
\end{pgfscope}%
\begin{pgfscope}%
\pgfpathrectangle{\pgfqpoint{0.487500in}{0.320625in}}{\pgfqpoint{1.845679in}{1.164375in}}%
\pgfusepath{clip}%
\pgfsetbuttcap%
\pgfsetroundjoin%
\pgfsetlinewidth{0.803000pt}%
\definecolor{currentstroke}{rgb}{0.000000,0.529412,0.423529}%
\pgfsetstrokecolor{currentstroke}%
\pgfsetstrokeopacity{0.800000}%
\pgfsetdash{{0.800000pt}{1.320000pt}}{0.000000pt}%
\pgfpathmoveto{\pgfqpoint{0.571395in}{1.468525in}}%
\pgfpathlineto{\pgfqpoint{2.249285in}{1.468525in}}%
\pgfpathlineto{\pgfqpoint{2.249285in}{1.468525in}}%
\pgfusepath{stroke}%
\end{pgfscope}%
\begin{pgfscope}%
\pgfpathrectangle{\pgfqpoint{0.487500in}{0.320625in}}{\pgfqpoint{1.845679in}{1.164375in}}%
\pgfusepath{clip}%
\pgfsetrectcap%
\pgfsetroundjoin%
\pgfsetlinewidth{1.254687pt}%
\definecolor{currentstroke}{rgb}{0.478431,0.243137,0.615686}%
\pgfsetstrokecolor{currentstroke}%
\pgfsetdash{}{0pt}%
\pgfpathmoveto{\pgfqpoint{0.571395in}{1.468525in}}%
\pgfpathlineto{\pgfqpoint{0.598660in}{1.219935in}}%
\pgfpathlineto{\pgfqpoint{0.655289in}{0.703632in}}%
\pgfpathlineto{\pgfqpoint{0.657386in}{0.687880in}}%
\pgfpathlineto{\pgfqpoint{0.661581in}{0.635236in}}%
\pgfpathlineto{\pgfqpoint{0.663678in}{0.628303in}}%
\pgfpathlineto{\pgfqpoint{0.665776in}{0.635128in}}%
\pgfpathlineto{\pgfqpoint{0.667873in}{0.613144in}}%
\pgfpathlineto{\pgfqpoint{0.669971in}{0.545875in}}%
\pgfpathlineto{\pgfqpoint{0.672068in}{0.601897in}}%
\pgfpathlineto{\pgfqpoint{0.674165in}{0.583457in}}%
\pgfpathlineto{\pgfqpoint{0.676263in}{0.586511in}}%
\pgfpathlineto{\pgfqpoint{0.678360in}{0.624315in}}%
\pgfpathlineto{\pgfqpoint{0.680457in}{0.605347in}}%
\pgfpathlineto{\pgfqpoint{0.684652in}{0.666735in}}%
\pgfpathlineto{\pgfqpoint{0.686749in}{0.674127in}}%
\pgfpathlineto{\pgfqpoint{0.688847in}{0.663123in}}%
\pgfpathlineto{\pgfqpoint{0.690944in}{0.589738in}}%
\pgfpathlineto{\pgfqpoint{0.693042in}{0.660727in}}%
\pgfpathlineto{\pgfqpoint{0.695139in}{0.683119in}}%
\pgfpathlineto{\pgfqpoint{0.697236in}{0.676148in}}%
\pgfpathlineto{\pgfqpoint{0.699334in}{0.625452in}}%
\pgfpathlineto{\pgfqpoint{0.703528in}{0.633982in}}%
\pgfpathlineto{\pgfqpoint{0.705626in}{0.597412in}}%
\pgfpathlineto{\pgfqpoint{0.707723in}{0.662049in}}%
\pgfpathlineto{\pgfqpoint{0.709820in}{0.662056in}}%
\pgfpathlineto{\pgfqpoint{0.711918in}{0.612722in}}%
\pgfpathlineto{\pgfqpoint{0.714015in}{0.615012in}}%
\pgfpathlineto{\pgfqpoint{0.716113in}{0.620572in}}%
\pgfpathlineto{\pgfqpoint{0.718210in}{0.577875in}}%
\pgfpathlineto{\pgfqpoint{0.720307in}{0.647107in}}%
\pgfpathlineto{\pgfqpoint{0.722405in}{0.648095in}}%
\pgfpathlineto{\pgfqpoint{0.724502in}{0.532743in}}%
\pgfpathlineto{\pgfqpoint{0.726599in}{0.644753in}}%
\pgfpathlineto{\pgfqpoint{0.728697in}{0.653226in}}%
\pgfpathlineto{\pgfqpoint{0.730794in}{0.623917in}}%
\pgfpathlineto{\pgfqpoint{0.732891in}{0.632237in}}%
\pgfpathlineto{\pgfqpoint{0.734989in}{0.666421in}}%
\pgfpathlineto{\pgfqpoint{0.737086in}{0.655621in}}%
\pgfpathlineto{\pgfqpoint{0.739184in}{0.624539in}}%
\pgfpathlineto{\pgfqpoint{0.741281in}{0.674391in}}%
\pgfpathlineto{\pgfqpoint{0.743378in}{0.678589in}}%
\pgfpathlineto{\pgfqpoint{0.745476in}{0.648761in}}%
\pgfpathlineto{\pgfqpoint{0.747573in}{0.654151in}}%
\pgfpathlineto{\pgfqpoint{0.749670in}{0.689790in}}%
\pgfpathlineto{\pgfqpoint{0.751768in}{0.684393in}}%
\pgfpathlineto{\pgfqpoint{0.755962in}{0.606141in}}%
\pgfpathlineto{\pgfqpoint{0.758060in}{0.603575in}}%
\pgfpathlineto{\pgfqpoint{0.762254in}{0.664687in}}%
\pgfpathlineto{\pgfqpoint{0.764352in}{0.648469in}}%
\pgfpathlineto{\pgfqpoint{0.766449in}{0.582492in}}%
\pgfpathlineto{\pgfqpoint{0.768547in}{0.592695in}}%
\pgfpathlineto{\pgfqpoint{0.770644in}{0.519305in}}%
\pgfpathlineto{\pgfqpoint{0.772741in}{0.623417in}}%
\pgfpathlineto{\pgfqpoint{0.774839in}{0.638351in}}%
\pgfpathlineto{\pgfqpoint{0.776936in}{0.593974in}}%
\pgfpathlineto{\pgfqpoint{0.781131in}{0.629974in}}%
\pgfpathlineto{\pgfqpoint{0.783228in}{0.607548in}}%
\pgfpathlineto{\pgfqpoint{0.785325in}{0.604279in}}%
\pgfpathlineto{\pgfqpoint{0.787423in}{0.643810in}}%
\pgfpathlineto{\pgfqpoint{0.789520in}{0.634086in}}%
\pgfpathlineto{\pgfqpoint{0.791618in}{0.613281in}}%
\pgfpathlineto{\pgfqpoint{0.793715in}{0.659854in}}%
\pgfpathlineto{\pgfqpoint{0.795812in}{0.665337in}}%
\pgfpathlineto{\pgfqpoint{0.797910in}{0.642078in}}%
\pgfpathlineto{\pgfqpoint{0.800007in}{0.626189in}}%
\pgfpathlineto{\pgfqpoint{0.802104in}{0.672261in}}%
\pgfpathlineto{\pgfqpoint{0.804202in}{0.674007in}}%
\pgfpathlineto{\pgfqpoint{0.808396in}{0.595425in}}%
\pgfpathlineto{\pgfqpoint{0.810494in}{0.615360in}}%
\pgfpathlineto{\pgfqpoint{0.812591in}{0.590259in}}%
\pgfpathlineto{\pgfqpoint{0.814689in}{0.650247in}}%
\pgfpathlineto{\pgfqpoint{0.816786in}{0.648153in}}%
\pgfpathlineto{\pgfqpoint{0.818883in}{0.600571in}}%
\pgfpathlineto{\pgfqpoint{0.820981in}{0.588506in}}%
\pgfpathlineto{\pgfqpoint{0.823078in}{0.586934in}}%
\pgfpathlineto{\pgfqpoint{0.825175in}{0.592744in}}%
\pgfpathlineto{\pgfqpoint{0.827273in}{0.632724in}}%
\pgfpathlineto{\pgfqpoint{0.829370in}{0.619027in}}%
\pgfpathlineto{\pgfqpoint{0.831467in}{0.572616in}}%
\pgfpathlineto{\pgfqpoint{0.833565in}{0.624975in}}%
\pgfpathlineto{\pgfqpoint{0.835662in}{0.622212in}}%
\pgfpathlineto{\pgfqpoint{0.837760in}{0.510430in}}%
\pgfpathlineto{\pgfqpoint{0.839857in}{0.629774in}}%
\pgfpathlineto{\pgfqpoint{0.841954in}{0.642979in}}%
\pgfpathlineto{\pgfqpoint{0.844052in}{0.585247in}}%
\pgfpathlineto{\pgfqpoint{0.846149in}{0.640768in}}%
\pgfpathlineto{\pgfqpoint{0.848246in}{0.659921in}}%
\pgfpathlineto{\pgfqpoint{0.850344in}{0.651719in}}%
\pgfpathlineto{\pgfqpoint{0.852441in}{0.550710in}}%
\pgfpathlineto{\pgfqpoint{0.854538in}{0.655813in}}%
\pgfpathlineto{\pgfqpoint{0.856636in}{0.672234in}}%
\pgfpathlineto{\pgfqpoint{0.858733in}{0.650993in}}%
\pgfpathlineto{\pgfqpoint{0.860831in}{0.530303in}}%
\pgfpathlineto{\pgfqpoint{0.862928in}{0.616590in}}%
\pgfpathlineto{\pgfqpoint{0.865025in}{0.567172in}}%
\pgfpathlineto{\pgfqpoint{0.867123in}{0.632846in}}%
\pgfpathlineto{\pgfqpoint{0.869220in}{0.651187in}}%
\pgfpathlineto{\pgfqpoint{0.871317in}{0.622951in}}%
\pgfpathlineto{\pgfqpoint{0.873415in}{0.546568in}}%
\pgfpathlineto{\pgfqpoint{0.875512in}{0.596643in}}%
\pgfpathlineto{\pgfqpoint{0.877609in}{0.466427in}}%
\pgfpathlineto{\pgfqpoint{0.879707in}{0.619976in}}%
\pgfpathlineto{\pgfqpoint{0.881804in}{0.629257in}}%
\pgfpathlineto{\pgfqpoint{0.883902in}{0.569167in}}%
\pgfpathlineto{\pgfqpoint{0.885999in}{0.611208in}}%
\pgfpathlineto{\pgfqpoint{0.888096in}{0.624356in}}%
\pgfpathlineto{\pgfqpoint{0.890194in}{0.596172in}}%
\pgfpathlineto{\pgfqpoint{0.892291in}{0.605477in}}%
\pgfpathlineto{\pgfqpoint{0.894388in}{0.639049in}}%
\pgfpathlineto{\pgfqpoint{0.898583in}{0.610882in}}%
\pgfpathlineto{\pgfqpoint{0.900680in}{0.651769in}}%
\pgfpathlineto{\pgfqpoint{0.902778in}{0.654318in}}%
\pgfpathlineto{\pgfqpoint{0.904875in}{0.622249in}}%
\pgfpathlineto{\pgfqpoint{0.906973in}{0.632006in}}%
\pgfpathlineto{\pgfqpoint{0.909070in}{0.665874in}}%
\pgfpathlineto{\pgfqpoint{0.911167in}{0.659583in}}%
\pgfpathlineto{\pgfqpoint{0.913265in}{0.607846in}}%
\pgfpathlineto{\pgfqpoint{0.915362in}{0.608240in}}%
\pgfpathlineto{\pgfqpoint{0.917459in}{0.606161in}}%
\pgfpathlineto{\pgfqpoint{0.919557in}{0.602772in}}%
\pgfpathlineto{\pgfqpoint{0.921654in}{0.646633in}}%
\pgfpathlineto{\pgfqpoint{0.923751in}{0.637058in}}%
\pgfpathlineto{\pgfqpoint{0.925849in}{0.570476in}}%
\pgfpathlineto{\pgfqpoint{0.927946in}{0.594481in}}%
\pgfpathlineto{\pgfqpoint{0.930043in}{0.579993in}}%
\pgfpathlineto{\pgfqpoint{0.932141in}{0.596605in}}%
\pgfpathlineto{\pgfqpoint{0.934238in}{0.629996in}}%
\pgfpathlineto{\pgfqpoint{0.936336in}{0.610607in}}%
\pgfpathlineto{\pgfqpoint{0.938433in}{0.581135in}}%
\pgfpathlineto{\pgfqpoint{0.940530in}{0.622060in}}%
\pgfpathlineto{\pgfqpoint{0.942628in}{0.615453in}}%
\pgfpathlineto{\pgfqpoint{0.944725in}{0.529187in}}%
\pgfpathlineto{\pgfqpoint{0.946822in}{0.629317in}}%
\pgfpathlineto{\pgfqpoint{0.948920in}{0.637569in}}%
\pgfpathlineto{\pgfqpoint{0.951017in}{0.549414in}}%
\pgfpathlineto{\pgfqpoint{0.953114in}{0.638723in}}%
\pgfpathlineto{\pgfqpoint{0.955212in}{0.653755in}}%
\pgfpathlineto{\pgfqpoint{0.957309in}{0.640807in}}%
\pgfpathlineto{\pgfqpoint{0.959407in}{0.571518in}}%
\pgfpathlineto{\pgfqpoint{0.961504in}{0.655297in}}%
\pgfpathlineto{\pgfqpoint{0.963601in}{0.665002in}}%
\pgfpathlineto{\pgfqpoint{0.965699in}{0.634997in}}%
\pgfpathlineto{\pgfqpoint{0.967796in}{0.577192in}}%
\pgfpathlineto{\pgfqpoint{0.969893in}{0.616077in}}%
\pgfpathlineto{\pgfqpoint{0.971991in}{0.509204in}}%
\pgfpathlineto{\pgfqpoint{0.974088in}{0.634827in}}%
\pgfpathlineto{\pgfqpoint{0.976185in}{0.646238in}}%
\pgfpathlineto{\pgfqpoint{0.980380in}{0.574256in}}%
\pgfpathlineto{\pgfqpoint{0.982478in}{0.597489in}}%
\pgfpathlineto{\pgfqpoint{0.984575in}{0.524021in}}%
\pgfpathlineto{\pgfqpoint{0.986672in}{0.621426in}}%
\pgfpathlineto{\pgfqpoint{0.988770in}{0.626748in}}%
\pgfpathlineto{\pgfqpoint{0.990867in}{0.552636in}}%
\pgfpathlineto{\pgfqpoint{0.992964in}{0.612361in}}%
\pgfpathlineto{\pgfqpoint{0.995062in}{0.622440in}}%
\pgfpathlineto{\pgfqpoint{0.997159in}{0.588003in}}%
\pgfpathlineto{\pgfqpoint{0.999256in}{0.609851in}}%
\pgfpathlineto{\pgfqpoint{1.001354in}{0.638501in}}%
\pgfpathlineto{\pgfqpoint{1.003451in}{0.618254in}}%
\pgfpathlineto{\pgfqpoint{1.005549in}{0.613806in}}%
\pgfpathlineto{\pgfqpoint{1.007646in}{0.649054in}}%
\pgfpathlineto{\pgfqpoint{1.009743in}{0.648625in}}%
\pgfpathlineto{\pgfqpoint{1.011841in}{0.605026in}}%
\pgfpathlineto{\pgfqpoint{1.016035in}{0.663627in}}%
\pgfpathlineto{\pgfqpoint{1.018133in}{0.650575in}}%
\pgfpathlineto{\pgfqpoint{1.020230in}{0.573463in}}%
\pgfpathlineto{\pgfqpoint{1.022327in}{0.614131in}}%
\pgfpathlineto{\pgfqpoint{1.024425in}{0.599072in}}%
\pgfpathlineto{\pgfqpoint{1.026522in}{0.611916in}}%
\pgfpathlineto{\pgfqpoint{1.028620in}{0.646407in}}%
\pgfpathlineto{\pgfqpoint{1.030717in}{0.630927in}}%
\pgfpathlineto{\pgfqpoint{1.032814in}{0.526031in}}%
\pgfpathlineto{\pgfqpoint{1.034912in}{0.599419in}}%
\pgfpathlineto{\pgfqpoint{1.037009in}{0.575714in}}%
\pgfpathlineto{\pgfqpoint{1.041204in}{0.630770in}}%
\pgfpathlineto{\pgfqpoint{1.045398in}{0.587794in}}%
\pgfpathlineto{\pgfqpoint{1.047496in}{0.622311in}}%
\pgfpathlineto{\pgfqpoint{1.049593in}{0.612480in}}%
\pgfpathlineto{\pgfqpoint{1.051691in}{0.559576in}}%
\pgfpathlineto{\pgfqpoint{1.053788in}{0.631481in}}%
\pgfpathlineto{\pgfqpoint{1.055885in}{0.635558in}}%
\pgfpathlineto{\pgfqpoint{1.057983in}{0.463573in}}%
\pgfpathlineto{\pgfqpoint{1.060080in}{0.638825in}}%
\pgfpathlineto{\pgfqpoint{1.062177in}{0.650708in}}%
\pgfpathlineto{\pgfqpoint{1.064275in}{0.632934in}}%
\pgfpathlineto{\pgfqpoint{1.066372in}{0.600295in}}%
\pgfpathlineto{\pgfqpoint{1.068469in}{0.656499in}}%
\pgfpathlineto{\pgfqpoint{1.070567in}{0.660034in}}%
\pgfpathlineto{\pgfqpoint{1.072664in}{0.621006in}}%
\pgfpathlineto{\pgfqpoint{1.074762in}{0.598280in}}%
\pgfpathlineto{\pgfqpoint{1.076859in}{0.615170in}}%
\pgfpathlineto{\pgfqpoint{1.078956in}{0.543813in}}%
\pgfpathlineto{\pgfqpoint{1.081054in}{0.638137in}}%
\pgfpathlineto{\pgfqpoint{1.083151in}{0.643013in}}%
\pgfpathlineto{\pgfqpoint{1.085248in}{0.599662in}}%
\pgfpathlineto{\pgfqpoint{1.087346in}{0.587510in}}%
\pgfpathlineto{\pgfqpoint{1.089443in}{0.598169in}}%
\pgfpathlineto{\pgfqpoint{1.091540in}{0.551573in}}%
\pgfpathlineto{\pgfqpoint{1.093638in}{0.624656in}}%
\pgfpathlineto{\pgfqpoint{1.095735in}{0.625669in}}%
\pgfpathlineto{\pgfqpoint{1.097832in}{0.524856in}}%
\pgfpathlineto{\pgfqpoint{1.099930in}{0.614905in}}%
\pgfpathlineto{\pgfqpoint{1.102027in}{0.621754in}}%
\pgfpathlineto{\pgfqpoint{1.104125in}{0.578129in}}%
\pgfpathlineto{\pgfqpoint{1.106222in}{0.615563in}}%
\pgfpathlineto{\pgfqpoint{1.108319in}{0.639304in}}%
\pgfpathlineto{\pgfqpoint{1.110417in}{0.612093in}}%
\pgfpathlineto{\pgfqpoint{1.112514in}{0.618327in}}%
\pgfpathlineto{\pgfqpoint{1.114611in}{0.648048in}}%
\pgfpathlineto{\pgfqpoint{1.116709in}{0.644201in}}%
\pgfpathlineto{\pgfqpoint{1.118806in}{0.579370in}}%
\pgfpathlineto{\pgfqpoint{1.120903in}{0.642777in}}%
\pgfpathlineto{\pgfqpoint{1.123001in}{0.662380in}}%
\pgfpathlineto{\pgfqpoint{1.125098in}{0.642237in}}%
\pgfpathlineto{\pgfqpoint{1.127196in}{0.506817in}}%
\pgfpathlineto{\pgfqpoint{1.129293in}{0.617758in}}%
\pgfpathlineto{\pgfqpoint{1.131390in}{0.588983in}}%
\pgfpathlineto{\pgfqpoint{1.135585in}{0.646888in}}%
\pgfpathlineto{\pgfqpoint{1.137682in}{0.625101in}}%
\pgfpathlineto{\pgfqpoint{1.139780in}{0.529315in}}%
\pgfpathlineto{\pgfqpoint{1.141877in}{0.603360in}}%
\pgfpathlineto{\pgfqpoint{1.143974in}{0.568184in}}%
\pgfpathlineto{\pgfqpoint{1.146072in}{0.608767in}}%
\pgfpathlineto{\pgfqpoint{1.148169in}{0.632228in}}%
\pgfpathlineto{\pgfqpoint{1.150267in}{0.602176in}}%
\pgfpathlineto{\pgfqpoint{1.152364in}{0.595178in}}%
\pgfpathlineto{\pgfqpoint{1.154461in}{0.623308in}}%
\pgfpathlineto{\pgfqpoint{1.156559in}{0.609270in}}%
\pgfpathlineto{\pgfqpoint{1.158656in}{0.579341in}}%
\pgfpathlineto{\pgfqpoint{1.160753in}{0.634508in}}%
\pgfpathlineto{\pgfqpoint{1.162851in}{0.633446in}}%
\pgfpathlineto{\pgfqpoint{1.164948in}{0.551270in}}%
\pgfpathlineto{\pgfqpoint{1.167045in}{0.640038in}}%
\pgfpathlineto{\pgfqpoint{1.169143in}{0.648378in}}%
\pgfpathlineto{\pgfqpoint{1.171240in}{0.624052in}}%
\pgfpathlineto{\pgfqpoint{1.173338in}{0.616675in}}%
\pgfpathlineto{\pgfqpoint{1.175435in}{0.658081in}}%
\pgfpathlineto{\pgfqpoint{1.177532in}{0.654700in}}%
\pgfpathlineto{\pgfqpoint{1.179630in}{0.602739in}}%
\pgfpathlineto{\pgfqpoint{1.181727in}{0.609507in}}%
\pgfpathlineto{\pgfqpoint{1.183824in}{0.612765in}}%
\pgfpathlineto{\pgfqpoint{1.185922in}{0.580084in}}%
\pgfpathlineto{\pgfqpoint{1.188019in}{0.641641in}}%
\pgfpathlineto{\pgfqpoint{1.190116in}{0.639649in}}%
\pgfpathlineto{\pgfqpoint{1.192214in}{0.584173in}}%
\pgfpathlineto{\pgfqpoint{1.194311in}{0.596877in}}%
\pgfpathlineto{\pgfqpoint{1.196409in}{0.597517in}}%
\pgfpathlineto{\pgfqpoint{1.198506in}{0.572004in}}%
\pgfpathlineto{\pgfqpoint{1.200603in}{0.628374in}}%
\pgfpathlineto{\pgfqpoint{1.202701in}{0.624091in}}%
\pgfpathlineto{\pgfqpoint{1.204798in}{0.493040in}}%
\pgfpathlineto{\pgfqpoint{1.206895in}{0.617962in}}%
\pgfpathlineto{\pgfqpoint{1.208993in}{0.620917in}}%
\pgfpathlineto{\pgfqpoint{1.211090in}{0.560142in}}%
\pgfpathlineto{\pgfqpoint{1.213187in}{0.621702in}}%
\pgfpathlineto{\pgfqpoint{1.215285in}{0.640093in}}%
\pgfpathlineto{\pgfqpoint{1.217382in}{0.603107in}}%
\pgfpathlineto{\pgfqpoint{1.221577in}{0.647505in}}%
\pgfpathlineto{\pgfqpoint{1.223674in}{0.639404in}}%
\pgfpathlineto{\pgfqpoint{1.225772in}{0.418414in}}%
\pgfpathlineto{\pgfqpoint{1.227869in}{0.647710in}}%
\pgfpathlineto{\pgfqpoint{1.229966in}{0.660733in}}%
\pgfpathlineto{\pgfqpoint{1.232064in}{0.624098in}}%
\pgfpathlineto{\pgfqpoint{1.234161in}{0.631763in}}%
\pgfpathlineto{\pgfqpoint{1.236258in}{0.655077in}}%
\pgfpathlineto{\pgfqpoint{1.238356in}{0.639548in}}%
\pgfpathlineto{\pgfqpoint{1.240453in}{0.612249in}}%
\pgfpathlineto{\pgfqpoint{1.242551in}{0.665494in}}%
\pgfpathlineto{\pgfqpoint{1.244648in}{0.664539in}}%
\pgfpathlineto{\pgfqpoint{1.246745in}{0.614374in}}%
\pgfpathlineto{\pgfqpoint{1.250940in}{0.624120in}}%
\pgfpathlineto{\pgfqpoint{1.253037in}{0.581396in}}%
\pgfpathlineto{\pgfqpoint{1.255135in}{0.650370in}}%
\pgfpathlineto{\pgfqpoint{1.257232in}{0.650616in}}%
\pgfpathlineto{\pgfqpoint{1.259329in}{0.605894in}}%
\pgfpathlineto{\pgfqpoint{1.261427in}{0.587623in}}%
\pgfpathlineto{\pgfqpoint{1.263524in}{0.590819in}}%
\pgfpathlineto{\pgfqpoint{1.265621in}{0.589614in}}%
\pgfpathlineto{\pgfqpoint{1.267719in}{0.633428in}}%
\pgfpathlineto{\pgfqpoint{1.269816in}{0.622871in}}%
\pgfpathlineto{\pgfqpoint{1.271914in}{0.474768in}}%
\pgfpathlineto{\pgfqpoint{1.274011in}{0.608269in}}%
\pgfpathlineto{\pgfqpoint{1.276108in}{0.602454in}}%
\pgfpathlineto{\pgfqpoint{1.278206in}{0.561483in}}%
\pgfpathlineto{\pgfqpoint{1.280303in}{0.626091in}}%
\pgfpathlineto{\pgfqpoint{1.282400in}{0.625908in}}%
\pgfpathlineto{\pgfqpoint{1.284498in}{0.529754in}}%
\pgfpathlineto{\pgfqpoint{1.286595in}{0.629414in}}%
\pgfpathlineto{\pgfqpoint{1.288692in}{0.637304in}}%
\pgfpathlineto{\pgfqpoint{1.290790in}{0.609970in}}%
\pgfpathlineto{\pgfqpoint{1.292887in}{0.612038in}}%
\pgfpathlineto{\pgfqpoint{1.294985in}{0.648857in}}%
\pgfpathlineto{\pgfqpoint{1.297082in}{0.641036in}}%
\pgfpathlineto{\pgfqpoint{1.299179in}{0.589583in}}%
\pgfpathlineto{\pgfqpoint{1.301277in}{0.651765in}}%
\pgfpathlineto{\pgfqpoint{1.303374in}{0.656337in}}%
\pgfpathlineto{\pgfqpoint{1.305471in}{0.622666in}}%
\pgfpathlineto{\pgfqpoint{1.307569in}{0.638414in}}%
\pgfpathlineto{\pgfqpoint{1.309666in}{0.669632in}}%
\pgfpathlineto{\pgfqpoint{1.311763in}{0.658657in}}%
\pgfpathlineto{\pgfqpoint{1.313861in}{0.494093in}}%
\pgfpathlineto{\pgfqpoint{1.315958in}{0.642631in}}%
\pgfpathlineto{\pgfqpoint{1.318056in}{0.638696in}}%
\pgfpathlineto{\pgfqpoint{1.320153in}{0.587193in}}%
\pgfpathlineto{\pgfqpoint{1.322250in}{0.659911in}}%
\pgfpathlineto{\pgfqpoint{1.324348in}{0.662083in}}%
\pgfpathlineto{\pgfqpoint{1.328542in}{0.587578in}}%
\pgfpathlineto{\pgfqpoint{1.330640in}{0.597410in}}%
\pgfpathlineto{\pgfqpoint{1.332737in}{0.596846in}}%
\pgfpathlineto{\pgfqpoint{1.334834in}{0.641447in}}%
\pgfpathlineto{\pgfqpoint{1.336932in}{0.631963in}}%
\pgfpathlineto{\pgfqpoint{1.339029in}{0.557637in}}%
\pgfpathlineto{\pgfqpoint{1.341127in}{0.595607in}}%
\pgfpathlineto{\pgfqpoint{1.343224in}{0.583468in}}%
\pgfpathlineto{\pgfqpoint{1.345321in}{0.589660in}}%
\pgfpathlineto{\pgfqpoint{1.347419in}{0.627191in}}%
\pgfpathlineto{\pgfqpoint{1.349516in}{0.611349in}}%
\pgfpathlineto{\pgfqpoint{1.351613in}{0.576042in}}%
\pgfpathlineto{\pgfqpoint{1.353711in}{0.621878in}}%
\pgfpathlineto{\pgfqpoint{1.355808in}{0.618264in}}%
\pgfpathlineto{\pgfqpoint{1.357905in}{0.500081in}}%
\pgfpathlineto{\pgfqpoint{1.360003in}{0.625868in}}%
\pgfpathlineto{\pgfqpoint{1.362100in}{0.638887in}}%
\pgfpathlineto{\pgfqpoint{1.364198in}{0.584236in}}%
\pgfpathlineto{\pgfqpoint{1.366295in}{0.633585in}}%
\pgfpathlineto{\pgfqpoint{1.368392in}{0.652691in}}%
\pgfpathlineto{\pgfqpoint{1.370490in}{0.643089in}}%
\pgfpathlineto{\pgfqpoint{1.372587in}{0.494074in}}%
\pgfpathlineto{\pgfqpoint{1.374684in}{0.651714in}}%
\pgfpathlineto{\pgfqpoint{1.376782in}{0.664951in}}%
\pgfpathlineto{\pgfqpoint{1.378879in}{0.639016in}}%
\pgfpathlineto{\pgfqpoint{1.380976in}{0.551780in}}%
\pgfpathlineto{\pgfqpoint{1.383074in}{0.613882in}}%
\pgfpathlineto{\pgfqpoint{1.385171in}{0.543774in}}%
\pgfpathlineto{\pgfqpoint{1.387269in}{0.630900in}}%
\pgfpathlineto{\pgfqpoint{1.389366in}{0.645437in}}%
\pgfpathlineto{\pgfqpoint{1.391463in}{0.613104in}}%
\pgfpathlineto{\pgfqpoint{1.393561in}{0.563158in}}%
\pgfpathlineto{\pgfqpoint{1.395658in}{0.594678in}}%
\pgfpathlineto{\pgfqpoint{1.397755in}{0.503127in}}%
\pgfpathlineto{\pgfqpoint{1.399853in}{0.618004in}}%
\pgfpathlineto{\pgfqpoint{1.401950in}{0.624945in}}%
\pgfpathlineto{\pgfqpoint{1.404047in}{0.556737in}}%
\pgfpathlineto{\pgfqpoint{1.406145in}{0.609185in}}%
\pgfpathlineto{\pgfqpoint{1.408242in}{0.620348in}}%
\pgfpathlineto{\pgfqpoint{1.410340in}{0.588093in}}%
\pgfpathlineto{\pgfqpoint{1.412437in}{0.605866in}}%
\pgfpathlineto{\pgfqpoint{1.414534in}{0.635985in}}%
\pgfpathlineto{\pgfqpoint{1.416632in}{0.617261in}}%
\pgfpathlineto{\pgfqpoint{1.418729in}{0.611407in}}%
\pgfpathlineto{\pgfqpoint{1.420826in}{0.647857in}}%
\pgfpathlineto{\pgfqpoint{1.422924in}{0.648338in}}%
\pgfpathlineto{\pgfqpoint{1.425021in}{0.609098in}}%
\pgfpathlineto{\pgfqpoint{1.429216in}{0.662204in}}%
\pgfpathlineto{\pgfqpoint{1.431313in}{0.651430in}}%
\pgfpathlineto{\pgfqpoint{1.433410in}{0.584205in}}%
\pgfpathlineto{\pgfqpoint{1.435508in}{0.611438in}}%
\pgfpathlineto{\pgfqpoint{1.437605in}{0.600791in}}%
\pgfpathlineto{\pgfqpoint{1.439703in}{0.606754in}}%
\pgfpathlineto{\pgfqpoint{1.441800in}{0.644702in}}%
\pgfpathlineto{\pgfqpoint{1.443897in}{0.631477in}}%
\pgfpathlineto{\pgfqpoint{1.445995in}{0.543989in}}%
\pgfpathlineto{\pgfqpoint{1.448092in}{0.597099in}}%
\pgfpathlineto{\pgfqpoint{1.450189in}{0.577484in}}%
\pgfpathlineto{\pgfqpoint{1.452287in}{0.598104in}}%
\pgfpathlineto{\pgfqpoint{1.454384in}{0.629147in}}%
\pgfpathlineto{\pgfqpoint{1.458579in}{0.583670in}}%
\pgfpathlineto{\pgfqpoint{1.460676in}{0.621131in}}%
\pgfpathlineto{\pgfqpoint{1.462774in}{0.612916in}}%
\pgfpathlineto{\pgfqpoint{1.464871in}{0.546450in}}%
\pgfpathlineto{\pgfqpoint{1.466968in}{0.629351in}}%
\pgfpathlineto{\pgfqpoint{1.469066in}{0.635476in}}%
\pgfpathlineto{\pgfqpoint{1.471163in}{0.523972in}}%
\pgfpathlineto{\pgfqpoint{1.473260in}{0.637304in}}%
\pgfpathlineto{\pgfqpoint{1.475358in}{0.650621in}}%
\pgfpathlineto{\pgfqpoint{1.477455in}{0.634923in}}%
\pgfpathlineto{\pgfqpoint{1.479552in}{0.590549in}}%
\pgfpathlineto{\pgfqpoint{1.481650in}{0.654878in}}%
\pgfpathlineto{\pgfqpoint{1.483747in}{0.660820in}}%
\pgfpathlineto{\pgfqpoint{1.487942in}{0.592054in}}%
\pgfpathlineto{\pgfqpoint{1.490039in}{0.615033in}}%
\pgfpathlineto{\pgfqpoint{1.492137in}{0.511701in}}%
\pgfpathlineto{\pgfqpoint{1.494234in}{0.635979in}}%
\pgfpathlineto{\pgfqpoint{1.496331in}{0.643387in}}%
\pgfpathlineto{\pgfqpoint{1.498429in}{0.603116in}}%
\pgfpathlineto{\pgfqpoint{1.500526in}{0.582942in}}%
\pgfpathlineto{\pgfqpoint{1.502623in}{0.597662in}}%
\pgfpathlineto{\pgfqpoint{1.504721in}{0.540742in}}%
\pgfpathlineto{\pgfqpoint{1.506818in}{0.622669in}}%
\pgfpathlineto{\pgfqpoint{1.508916in}{0.625529in}}%
\pgfpathlineto{\pgfqpoint{1.511013in}{0.538370in}}%
\pgfpathlineto{\pgfqpoint{1.513110in}{0.613254in}}%
\pgfpathlineto{\pgfqpoint{1.515208in}{0.621428in}}%
\pgfpathlineto{\pgfqpoint{1.517305in}{0.581855in}}%
\pgfpathlineto{\pgfqpoint{1.521500in}{0.638387in}}%
\pgfpathlineto{\pgfqpoint{1.523597in}{0.613973in}}%
\pgfpathlineto{\pgfqpoint{1.525694in}{0.615938in}}%
\pgfpathlineto{\pgfqpoint{1.527792in}{0.647691in}}%
\pgfpathlineto{\pgfqpoint{1.529889in}{0.645111in}}%
\pgfpathlineto{\pgfqpoint{1.531987in}{0.589573in}}%
\pgfpathlineto{\pgfqpoint{1.534084in}{0.640414in}}%
\pgfpathlineto{\pgfqpoint{1.536181in}{0.662210in}}%
\pgfpathlineto{\pgfqpoint{1.538279in}{0.644489in}}%
\pgfpathlineto{\pgfqpoint{1.540376in}{0.510233in}}%
\pgfpathlineto{\pgfqpoint{1.542473in}{0.616470in}}%
\pgfpathlineto{\pgfqpoint{1.544571in}{0.592691in}}%
\pgfpathlineto{\pgfqpoint{1.548765in}{0.646298in}}%
\pgfpathlineto{\pgfqpoint{1.550863in}{0.626692in}}%
\pgfpathlineto{\pgfqpoint{1.552960in}{0.494367in}}%
\pgfpathlineto{\pgfqpoint{1.555058in}{0.601997in}}%
\pgfpathlineto{\pgfqpoint{1.557155in}{0.571251in}}%
\pgfpathlineto{\pgfqpoint{1.559252in}{0.606099in}}%
\pgfpathlineto{\pgfqpoint{1.561350in}{0.631437in}}%
\pgfpathlineto{\pgfqpoint{1.563447in}{0.603650in}}%
\pgfpathlineto{\pgfqpoint{1.565544in}{0.592394in}}%
\pgfpathlineto{\pgfqpoint{1.567642in}{0.622689in}}%
\pgfpathlineto{\pgfqpoint{1.569739in}{0.610195in}}%
\pgfpathlineto{\pgfqpoint{1.571836in}{0.573138in}}%
\pgfpathlineto{\pgfqpoint{1.573934in}{0.633204in}}%
\pgfpathlineto{\pgfqpoint{1.576031in}{0.633942in}}%
\pgfpathlineto{\pgfqpoint{1.578129in}{0.534303in}}%
\pgfpathlineto{\pgfqpoint{1.580226in}{0.639290in}}%
\pgfpathlineto{\pgfqpoint{1.582323in}{0.648806in}}%
\pgfpathlineto{\pgfqpoint{1.584421in}{0.626777in}}%
\pgfpathlineto{\pgfqpoint{1.586518in}{0.611972in}}%
\pgfpathlineto{\pgfqpoint{1.588615in}{0.657327in}}%
\pgfpathlineto{\pgfqpoint{1.590713in}{0.656106in}}%
\pgfpathlineto{\pgfqpoint{1.592810in}{0.608749in}}%
\pgfpathlineto{\pgfqpoint{1.594907in}{0.606540in}}%
\pgfpathlineto{\pgfqpoint{1.597005in}{0.613539in}}%
\pgfpathlineto{\pgfqpoint{1.599102in}{0.571508in}}%
\pgfpathlineto{\pgfqpoint{1.601199in}{0.640391in}}%
\pgfpathlineto{\pgfqpoint{1.603297in}{0.640506in}}%
\pgfpathlineto{\pgfqpoint{1.605394in}{0.589297in}}%
\pgfpathlineto{\pgfqpoint{1.607492in}{0.594268in}}%
\pgfpathlineto{\pgfqpoint{1.609589in}{0.597762in}}%
\pgfpathlineto{\pgfqpoint{1.611686in}{0.566252in}}%
\pgfpathlineto{\pgfqpoint{1.613784in}{0.627094in}}%
\pgfpathlineto{\pgfqpoint{1.615881in}{0.624475in}}%
\pgfpathlineto{\pgfqpoint{1.617978in}{0.395652in}}%
\pgfpathlineto{\pgfqpoint{1.620076in}{0.616923in}}%
\pgfpathlineto{\pgfqpoint{1.622173in}{0.621068in}}%
\pgfpathlineto{\pgfqpoint{1.624270in}{0.566610in}}%
\pgfpathlineto{\pgfqpoint{1.626368in}{0.619792in}}%
\pgfpathlineto{\pgfqpoint{1.628465in}{0.639758in}}%
\pgfpathlineto{\pgfqpoint{1.630563in}{0.605886in}}%
\pgfpathlineto{\pgfqpoint{1.632660in}{0.621932in}}%
\pgfpathlineto{\pgfqpoint{1.634757in}{0.647542in}}%
\pgfpathlineto{\pgfqpoint{1.636855in}{0.640732in}}%
\pgfpathlineto{\pgfqpoint{1.638952in}{0.531120in}}%
\pgfpathlineto{\pgfqpoint{1.641049in}{0.646281in}}%
\pgfpathlineto{\pgfqpoint{1.643147in}{0.661172in}}%
\pgfpathlineto{\pgfqpoint{1.645244in}{0.628382in}}%
\pgfpathlineto{\pgfqpoint{1.647341in}{0.629345in}}%
\pgfpathlineto{\pgfqpoint{1.649439in}{0.655438in}}%
\pgfpathlineto{\pgfqpoint{1.651536in}{0.642224in}}%
\pgfpathlineto{\pgfqpoint{1.653634in}{0.604362in}}%
\pgfpathlineto{\pgfqpoint{1.655731in}{0.664777in}}%
\pgfpathlineto{\pgfqpoint{1.657828in}{0.666323in}}%
\pgfpathlineto{\pgfqpoint{1.659926in}{0.620764in}}%
\pgfpathlineto{\pgfqpoint{1.662023in}{0.615880in}}%
\pgfpathlineto{\pgfqpoint{1.664120in}{0.625221in}}%
\pgfpathlineto{\pgfqpoint{1.666218in}{0.569752in}}%
\pgfpathlineto{\pgfqpoint{1.668315in}{0.649290in}}%
\pgfpathlineto{\pgfqpoint{1.670412in}{0.651846in}}%
\pgfpathlineto{\pgfqpoint{1.672510in}{0.610129in}}%
\pgfpathlineto{\pgfqpoint{1.674607in}{0.583266in}}%
\pgfpathlineto{\pgfqpoint{1.676705in}{0.591222in}}%
\pgfpathlineto{\pgfqpoint{1.678802in}{0.585878in}}%
\pgfpathlineto{\pgfqpoint{1.680899in}{0.632517in}}%
\pgfpathlineto{\pgfqpoint{1.682997in}{0.623604in}}%
\pgfpathlineto{\pgfqpoint{1.685094in}{0.461540in}}%
\pgfpathlineto{\pgfqpoint{1.687191in}{0.607140in}}%
\pgfpathlineto{\pgfqpoint{1.689289in}{0.602820in}}%
\pgfpathlineto{\pgfqpoint{1.691386in}{0.555124in}}%
\pgfpathlineto{\pgfqpoint{1.693483in}{0.624926in}}%
\pgfpathlineto{\pgfqpoint{1.695581in}{0.626163in}}%
\pgfpathlineto{\pgfqpoint{1.697678in}{0.512507in}}%
\pgfpathlineto{\pgfqpoint{1.699776in}{0.628777in}}%
\pgfpathlineto{\pgfqpoint{1.701873in}{0.637575in}}%
\pgfpathlineto{\pgfqpoint{1.703970in}{0.612236in}}%
\pgfpathlineto{\pgfqpoint{1.706068in}{0.608962in}}%
\pgfpathlineto{\pgfqpoint{1.708165in}{0.648296in}}%
\pgfpathlineto{\pgfqpoint{1.710262in}{0.642414in}}%
\pgfpathlineto{\pgfqpoint{1.712360in}{0.582061in}}%
\pgfpathlineto{\pgfqpoint{1.714457in}{0.651508in}}%
\pgfpathlineto{\pgfqpoint{1.716554in}{0.657349in}}%
\pgfpathlineto{\pgfqpoint{1.718652in}{0.627070in}}%
\pgfpathlineto{\pgfqpoint{1.720749in}{0.635432in}}%
\pgfpathlineto{\pgfqpoint{1.722847in}{0.669557in}}%
\pgfpathlineto{\pgfqpoint{1.724944in}{0.660968in}}%
\pgfpathlineto{\pgfqpoint{1.727041in}{0.556648in}}%
\pgfpathlineto{\pgfqpoint{1.729139in}{0.641895in}}%
\pgfpathlineto{\pgfqpoint{1.731236in}{0.640550in}}%
\pgfpathlineto{\pgfqpoint{1.733333in}{0.572315in}}%
\pgfpathlineto{\pgfqpoint{1.735431in}{0.658981in}}%
\pgfpathlineto{\pgfqpoint{1.737528in}{0.663640in}}%
\pgfpathlineto{\pgfqpoint{1.741723in}{0.581286in}}%
\pgfpathlineto{\pgfqpoint{1.743820in}{0.598024in}}%
\pgfpathlineto{\pgfqpoint{1.745918in}{0.592871in}}%
\pgfpathlineto{\pgfqpoint{1.748015in}{0.640719in}}%
\pgfpathlineto{\pgfqpoint{1.750112in}{0.633014in}}%
\pgfpathlineto{\pgfqpoint{1.752210in}{0.565593in}}%
\pgfpathlineto{\pgfqpoint{1.754307in}{0.593797in}}%
\pgfpathlineto{\pgfqpoint{1.756404in}{0.584166in}}%
\pgfpathlineto{\pgfqpoint{1.758502in}{0.586897in}}%
\pgfpathlineto{\pgfqpoint{1.760599in}{0.626370in}}%
\pgfpathlineto{\pgfqpoint{1.762696in}{0.611996in}}%
\pgfpathlineto{\pgfqpoint{1.764794in}{0.572680in}}%
\pgfpathlineto{\pgfqpoint{1.766891in}{0.621312in}}%
\pgfpathlineto{\pgfqpoint{1.768988in}{0.618633in}}%
\pgfpathlineto{\pgfqpoint{1.771086in}{0.523241in}}%
\pgfpathlineto{\pgfqpoint{1.773183in}{0.624536in}}%
\pgfpathlineto{\pgfqpoint{1.775281in}{0.638824in}}%
\pgfpathlineto{\pgfqpoint{1.777378in}{0.588556in}}%
\pgfpathlineto{\pgfqpoint{1.779475in}{0.632740in}}%
\pgfpathlineto{\pgfqpoint{1.781573in}{0.652977in}}%
\pgfpathlineto{\pgfqpoint{1.783670in}{0.644539in}}%
\pgfpathlineto{\pgfqpoint{1.785767in}{0.518390in}}%
\pgfpathlineto{\pgfqpoint{1.787865in}{0.650557in}}%
\pgfpathlineto{\pgfqpoint{1.789962in}{0.665495in}}%
\pgfpathlineto{\pgfqpoint{1.792059in}{0.641666in}}%
\pgfpathlineto{\pgfqpoint{1.794157in}{0.525138in}}%
\pgfpathlineto{\pgfqpoint{1.796254in}{0.613348in}}%
\pgfpathlineto{\pgfqpoint{1.798352in}{0.555281in}}%
\pgfpathlineto{\pgfqpoint{1.800449in}{0.629264in}}%
\pgfpathlineto{\pgfqpoint{1.802546in}{0.645723in}}%
\pgfpathlineto{\pgfqpoint{1.804644in}{0.615270in}}%
\pgfpathlineto{\pgfqpoint{1.806741in}{0.556055in}}%
\pgfpathlineto{\pgfqpoint{1.808838in}{0.594101in}}%
\pgfpathlineto{\pgfqpoint{1.810936in}{0.475645in}}%
\pgfpathlineto{\pgfqpoint{1.813033in}{0.616678in}}%
\pgfpathlineto{\pgfqpoint{1.815130in}{0.625016in}}%
\pgfpathlineto{\pgfqpoint{1.817228in}{0.561411in}}%
\pgfpathlineto{\pgfqpoint{1.819325in}{0.608070in}}%
\pgfpathlineto{\pgfqpoint{1.821423in}{0.620289in}}%
\pgfpathlineto{\pgfqpoint{1.823520in}{0.590186in}}%
\pgfpathlineto{\pgfqpoint{1.825617in}{0.603640in}}%
\pgfpathlineto{\pgfqpoint{1.827715in}{0.635443in}}%
\pgfpathlineto{\pgfqpoint{1.829812in}{0.618661in}}%
\pgfpathlineto{\pgfqpoint{1.831909in}{0.609404in}}%
\pgfpathlineto{\pgfqpoint{1.834007in}{0.647787in}}%
\pgfpathlineto{\pgfqpoint{1.836104in}{0.649187in}}%
\pgfpathlineto{\pgfqpoint{1.838201in}{0.613282in}}%
\pgfpathlineto{\pgfqpoint{1.840299in}{0.631346in}}%
\pgfpathlineto{\pgfqpoint{1.842396in}{0.662074in}}%
\pgfpathlineto{\pgfqpoint{1.844494in}{0.653168in}}%
\pgfpathlineto{\pgfqpoint{1.846591in}{0.592613in}}%
\pgfpathlineto{\pgfqpoint{1.848688in}{0.609658in}}%
\pgfpathlineto{\pgfqpoint{1.850786in}{0.602523in}}%
\pgfpathlineto{\pgfqpoint{1.852883in}{0.603348in}}%
\pgfpathlineto{\pgfqpoint{1.854980in}{0.644177in}}%
\pgfpathlineto{\pgfqpoint{1.857078in}{0.632725in}}%
\pgfpathlineto{\pgfqpoint{1.859175in}{0.555305in}}%
\pgfpathlineto{\pgfqpoint{1.861272in}{0.595604in}}%
\pgfpathlineto{\pgfqpoint{1.863370in}{0.578908in}}%
\pgfpathlineto{\pgfqpoint{1.865467in}{0.595722in}}%
\pgfpathlineto{\pgfqpoint{1.867565in}{0.628528in}}%
\pgfpathlineto{\pgfqpoint{1.869662in}{0.608542in}}%
\pgfpathlineto{\pgfqpoint{1.871759in}{0.580784in}}%
\pgfpathlineto{\pgfqpoint{1.873857in}{0.620671in}}%
\pgfpathlineto{\pgfqpoint{1.875954in}{0.613574in}}%
\pgfpathlineto{\pgfqpoint{1.878051in}{0.534656in}}%
\pgfpathlineto{\pgfqpoint{1.880149in}{0.628253in}}%
\pgfpathlineto{\pgfqpoint{1.882246in}{0.635777in}}%
\pgfpathlineto{\pgfqpoint{1.884343in}{0.542514in}}%
\pgfpathlineto{\pgfqpoint{1.886441in}{0.636645in}}%
\pgfpathlineto{\pgfqpoint{1.888538in}{0.651003in}}%
\pgfpathlineto{\pgfqpoint{1.890636in}{0.636732in}}%
\pgfpathlineto{\pgfqpoint{1.892733in}{0.581889in}}%
\pgfpathlineto{\pgfqpoint{1.894830in}{0.654066in}}%
\pgfpathlineto{\pgfqpoint{1.896928in}{0.661755in}}%
\pgfpathlineto{\pgfqpoint{1.901122in}{0.586640in}}%
\pgfpathlineto{\pgfqpoint{1.903220in}{0.615191in}}%
\pgfpathlineto{\pgfqpoint{1.905317in}{0.366262in}}%
\pgfpathlineto{\pgfqpoint{1.907414in}{0.634674in}}%
\pgfpathlineto{\pgfqpoint{1.909512in}{0.644074in}}%
\pgfpathlineto{\pgfqpoint{1.911609in}{0.606034in}}%
\pgfpathlineto{\pgfqpoint{1.913707in}{0.579281in}}%
\pgfpathlineto{\pgfqpoint{1.915804in}{0.597586in}}%
\pgfpathlineto{\pgfqpoint{1.917901in}{0.530341in}}%
\pgfpathlineto{\pgfqpoint{1.919999in}{0.621474in}}%
\pgfpathlineto{\pgfqpoint{1.922096in}{0.625843in}}%
\pgfpathlineto{\pgfqpoint{1.924193in}{0.547237in}}%
\pgfpathlineto{\pgfqpoint{1.926291in}{0.612227in}}%
\pgfpathlineto{\pgfqpoint{1.928388in}{0.621524in}}%
\pgfpathlineto{\pgfqpoint{1.930485in}{0.584925in}}%
\pgfpathlineto{\pgfqpoint{1.934680in}{0.637999in}}%
\pgfpathlineto{\pgfqpoint{1.936777in}{0.615867in}}%
\pgfpathlineto{\pgfqpoint{1.938875in}{0.614029in}}%
\pgfpathlineto{\pgfqpoint{1.940972in}{0.647646in}}%
\pgfpathlineto{\pgfqpoint{1.943070in}{0.646100in}}%
\pgfpathlineto{\pgfqpoint{1.945167in}{0.596476in}}%
\pgfpathlineto{\pgfqpoint{1.947264in}{0.638633in}}%
\pgfpathlineto{\pgfqpoint{1.949362in}{0.662297in}}%
\pgfpathlineto{\pgfqpoint{1.951459in}{0.646533in}}%
\pgfpathlineto{\pgfqpoint{1.953556in}{0.546001in}}%
\pgfpathlineto{\pgfqpoint{1.955654in}{0.615506in}}%
\pgfpathlineto{\pgfqpoint{1.957751in}{0.595585in}}%
\pgfpathlineto{\pgfqpoint{1.959848in}{0.614560in}}%
\pgfpathlineto{\pgfqpoint{1.961946in}{0.646021in}}%
\pgfpathlineto{\pgfqpoint{1.964043in}{0.628254in}}%
\pgfpathlineto{\pgfqpoint{1.966141in}{0.451045in}}%
\pgfpathlineto{\pgfqpoint{1.968238in}{0.601018in}}%
\pgfpathlineto{\pgfqpoint{1.970335in}{0.573826in}}%
\pgfpathlineto{\pgfqpoint{1.974530in}{0.631008in}}%
\pgfpathlineto{\pgfqpoint{1.976627in}{0.605207in}}%
\pgfpathlineto{\pgfqpoint{1.978725in}{0.589923in}}%
\pgfpathlineto{\pgfqpoint{1.980822in}{0.622343in}}%
\pgfpathlineto{\pgfqpoint{1.982919in}{0.611181in}}%
\pgfpathlineto{\pgfqpoint{1.985017in}{0.567051in}}%
\pgfpathlineto{\pgfqpoint{1.987114in}{0.632253in}}%
\pgfpathlineto{\pgfqpoint{1.989212in}{0.634575in}}%
\pgfpathlineto{\pgfqpoint{1.991309in}{0.506055in}}%
\pgfpathlineto{\pgfqpoint{1.993406in}{0.638698in}}%
\pgfpathlineto{\pgfqpoint{1.995504in}{0.649280in}}%
\pgfpathlineto{\pgfqpoint{1.999698in}{0.607449in}}%
\pgfpathlineto{\pgfqpoint{2.001796in}{0.656750in}}%
\pgfpathlineto{\pgfqpoint{2.003893in}{0.657409in}}%
\pgfpathlineto{\pgfqpoint{2.005990in}{0.613508in}}%
\pgfpathlineto{\pgfqpoint{2.008088in}{0.603758in}}%
\pgfpathlineto{\pgfqpoint{2.010185in}{0.614271in}}%
\pgfpathlineto{\pgfqpoint{2.012283in}{0.562122in}}%
\pgfpathlineto{\pgfqpoint{2.014380in}{0.639363in}}%
\pgfpathlineto{\pgfqpoint{2.016477in}{0.641414in}}%
\pgfpathlineto{\pgfqpoint{2.018575in}{0.593569in}}%
\pgfpathlineto{\pgfqpoint{2.020672in}{0.591903in}}%
\pgfpathlineto{\pgfqpoint{2.022769in}{0.598162in}}%
\pgfpathlineto{\pgfqpoint{2.024867in}{0.560242in}}%
\pgfpathlineto{\pgfqpoint{2.026964in}{0.626055in}}%
\pgfpathlineto{\pgfqpoint{2.029061in}{0.625060in}}%
\pgfpathlineto{\pgfqpoint{2.031159in}{0.496008in}}%
\pgfpathlineto{\pgfqpoint{2.033256in}{0.616021in}}%
\pgfpathlineto{\pgfqpoint{2.035354in}{0.621349in}}%
\pgfpathlineto{\pgfqpoint{2.037451in}{0.571990in}}%
\pgfpathlineto{\pgfqpoint{2.039548in}{0.618036in}}%
\pgfpathlineto{\pgfqpoint{2.041646in}{0.639568in}}%
\pgfpathlineto{\pgfqpoint{2.043743in}{0.608617in}}%
\pgfpathlineto{\pgfqpoint{2.045840in}{0.620291in}}%
\pgfpathlineto{\pgfqpoint{2.047938in}{0.647570in}}%
\pgfpathlineto{\pgfqpoint{2.050035in}{0.641948in}}%
\pgfpathlineto{\pgfqpoint{2.052132in}{0.556729in}}%
\pgfpathlineto{\pgfqpoint{2.054230in}{0.644905in}}%
\pgfpathlineto{\pgfqpoint{2.056327in}{0.661544in}}%
\pgfpathlineto{\pgfqpoint{2.058425in}{0.637782in}}%
\pgfpathlineto{\pgfqpoint{2.060522in}{0.557395in}}%
\pgfpathlineto{\pgfqpoint{2.062619in}{0.618778in}}%
\pgfpathlineto{\pgfqpoint{2.064717in}{0.581960in}}%
\pgfpathlineto{\pgfqpoint{2.066814in}{0.623972in}}%
\pgfpathlineto{\pgfqpoint{2.068911in}{0.646887in}}%
\pgfpathlineto{\pgfqpoint{2.071009in}{0.621869in}}%
\pgfpathlineto{\pgfqpoint{2.073106in}{0.553093in}}%
\pgfpathlineto{\pgfqpoint{2.075203in}{0.604802in}}%
\pgfpathlineto{\pgfqpoint{2.077301in}{0.562521in}}%
\pgfpathlineto{\pgfqpoint{2.079398in}{0.611893in}}%
\pgfpathlineto{\pgfqpoint{2.081496in}{0.632785in}}%
\pgfpathlineto{\pgfqpoint{2.083593in}{0.599256in}}%
\pgfpathlineto{\pgfqpoint{2.085690in}{0.598586in}}%
\pgfpathlineto{\pgfqpoint{2.087788in}{0.623740in}}%
\pgfpathlineto{\pgfqpoint{2.091982in}{0.586908in}}%
\pgfpathlineto{\pgfqpoint{2.094080in}{0.635932in}}%
\pgfpathlineto{\pgfqpoint{2.096177in}{0.632037in}}%
\pgfpathlineto{\pgfqpoint{2.098274in}{0.567758in}}%
\pgfpathlineto{\pgfqpoint{2.100372in}{0.640686in}}%
\pgfpathlineto{\pgfqpoint{2.102469in}{0.647163in}}%
\pgfpathlineto{\pgfqpoint{2.104566in}{0.618644in}}%
\pgfpathlineto{\pgfqpoint{2.106664in}{0.622725in}}%
\pgfpathlineto{\pgfqpoint{2.108761in}{0.658686in}}%
\pgfpathlineto{\pgfqpoint{2.110859in}{0.651177in}}%
\pgfpathlineto{\pgfqpoint{2.112956in}{0.558179in}}%
\pgfpathlineto{\pgfqpoint{2.115053in}{0.649410in}}%
\pgfpathlineto{\pgfqpoint{2.117151in}{0.652414in}}%
\pgfpathlineto{\pgfqpoint{2.119248in}{0.599186in}}%
\pgfpathlineto{\pgfqpoint{2.121345in}{0.649893in}}%
\pgfpathlineto{\pgfqpoint{2.123443in}{0.670554in}}%
\pgfpathlineto{\pgfqpoint{2.125540in}{0.648013in}}%
\pgfpathlineto{\pgfqpoint{2.127637in}{0.563414in}}%
\pgfpathlineto{\pgfqpoint{2.129735in}{0.629103in}}%
\pgfpathlineto{\pgfqpoint{2.131832in}{0.596628in}}%
\pgfpathlineto{\pgfqpoint{2.133930in}{0.631744in}}%
\pgfpathlineto{\pgfqpoint{2.136027in}{0.656712in}}%
\pgfpathlineto{\pgfqpoint{2.138124in}{0.634756in}}%
\pgfpathlineto{\pgfqpoint{2.140222in}{0.472020in}}%
\pgfpathlineto{\pgfqpoint{2.142319in}{0.599789in}}%
\pgfpathlineto{\pgfqpoint{2.144416in}{0.529219in}}%
\pgfpathlineto{\pgfqpoint{2.146514in}{0.620345in}}%
\pgfpathlineto{\pgfqpoint{2.148611in}{0.634912in}}%
\pgfpathlineto{\pgfqpoint{2.150708in}{0.596051in}}%
\pgfpathlineto{\pgfqpoint{2.152806in}{0.591195in}}%
\pgfpathlineto{\pgfqpoint{2.154903in}{0.611051in}}%
\pgfpathlineto{\pgfqpoint{2.157001in}{0.571442in}}%
\pgfpathlineto{\pgfqpoint{2.159098in}{0.608483in}}%
\pgfpathlineto{\pgfqpoint{2.161195in}{0.632119in}}%
\pgfpathlineto{\pgfqpoint{2.163293in}{0.601261in}}%
\pgfpathlineto{\pgfqpoint{2.165390in}{0.609911in}}%
\pgfpathlineto{\pgfqpoint{2.167487in}{0.637087in}}%
\pgfpathlineto{\pgfqpoint{2.169585in}{0.630023in}}%
\pgfpathlineto{\pgfqpoint{2.171682in}{0.478323in}}%
\pgfpathlineto{\pgfqpoint{2.173779in}{0.637604in}}%
\pgfpathlineto{\pgfqpoint{2.175877in}{0.650963in}}%
\pgfpathlineto{\pgfqpoint{2.177974in}{0.608272in}}%
\pgfpathlineto{\pgfqpoint{2.182169in}{0.657605in}}%
\pgfpathlineto{\pgfqpoint{2.184266in}{0.647169in}}%
\pgfpathlineto{\pgfqpoint{2.186364in}{0.564818in}}%
\pgfpathlineto{\pgfqpoint{2.188461in}{0.660222in}}%
\pgfpathlineto{\pgfqpoint{2.190558in}{0.669952in}}%
\pgfpathlineto{\pgfqpoint{2.192656in}{0.633120in}}%
\pgfpathlineto{\pgfqpoint{2.194753in}{0.624278in}}%
\pgfpathlineto{\pgfqpoint{2.196850in}{0.646258in}}%
\pgfpathlineto{\pgfqpoint{2.198948in}{0.610867in}}%
\pgfpathlineto{\pgfqpoint{2.203143in}{0.666897in}}%
\pgfpathlineto{\pgfqpoint{2.205240in}{0.647353in}}%
\pgfpathlineto{\pgfqpoint{2.207337in}{0.553574in}}%
\pgfpathlineto{\pgfqpoint{2.209435in}{0.604830in}}%
\pgfpathlineto{\pgfqpoint{2.211532in}{0.536160in}}%
\pgfpathlineto{\pgfqpoint{2.213629in}{0.628266in}}%
\pgfpathlineto{\pgfqpoint{2.215727in}{0.643206in}}%
\pgfpathlineto{\pgfqpoint{2.219921in}{0.570383in}}%
\pgfpathlineto{\pgfqpoint{2.222019in}{0.598150in}}%
\pgfpathlineto{\pgfqpoint{2.224116in}{0.460477in}}%
\pgfpathlineto{\pgfqpoint{2.226214in}{0.616367in}}%
\pgfpathlineto{\pgfqpoint{2.228311in}{0.626591in}}%
\pgfpathlineto{\pgfqpoint{2.230408in}{0.566721in}}%
\pgfpathlineto{\pgfqpoint{2.232506in}{0.610298in}}%
\pgfpathlineto{\pgfqpoint{2.234603in}{0.624217in}}%
\pgfpathlineto{\pgfqpoint{2.236700in}{0.599259in}}%
\pgfpathlineto{\pgfqpoint{2.238798in}{0.600039in}}%
\pgfpathlineto{\pgfqpoint{2.240895in}{0.637404in}}%
\pgfpathlineto{\pgfqpoint{2.242992in}{0.630163in}}%
\pgfpathlineto{\pgfqpoint{2.247187in}{0.590819in}}%
\pgfpathlineto{\pgfqpoint{2.249285in}{0.574712in}}%
\pgfpathlineto{\pgfqpoint{2.249285in}{0.574712in}}%
\pgfusepath{stroke}%
\end{pgfscope}%
\begin{pgfscope}%
\pgfpathrectangle{\pgfqpoint{0.487500in}{0.320625in}}{\pgfqpoint{1.845679in}{1.164375in}}%
\pgfusepath{clip}%
\pgfsetbuttcap%
\pgfsetroundjoin%
\pgfsetlinewidth{0.752812pt}%
\definecolor{currentstroke}{rgb}{0.478431,0.243137,0.615686}%
\pgfsetstrokecolor{currentstroke}%
\pgfsetstrokeopacity{0.650000}%
\pgfsetdash{{2.775000pt}{1.200000pt}}{0.000000pt}%
\pgfpathmoveto{\pgfqpoint{0.571395in}{1.468525in}}%
\pgfpathlineto{\pgfqpoint{0.700502in}{0.315625in}}%
\pgfpathlineto{\pgfqpoint{0.700502in}{0.315625in}}%
\pgfusepath{stroke}%
\end{pgfscope}%
\begin{pgfscope}%
\pgfpathrectangle{\pgfqpoint{0.487500in}{0.320625in}}{\pgfqpoint{1.845679in}{1.164375in}}%
\pgfusepath{clip}%
\pgfsetbuttcap%
\pgfsetroundjoin%
\pgfsetlinewidth{0.803000pt}%
\definecolor{currentstroke}{rgb}{0.478431,0.243137,0.615686}%
\pgfsetstrokecolor{currentstroke}%
\pgfsetstrokeopacity{0.800000}%
\pgfsetdash{{0.800000pt}{1.320000pt}}{0.000000pt}%
\pgfpathmoveto{\pgfqpoint{0.571395in}{1.468525in}}%
\pgfpathlineto{\pgfqpoint{2.249285in}{1.468525in}}%
\pgfpathlineto{\pgfqpoint{2.249285in}{1.468525in}}%
\pgfusepath{stroke}%
\end{pgfscope}%
\begin{pgfscope}%
\pgfsetrectcap%
\pgfsetmiterjoin%
\pgfsetlinewidth{0.803000pt}%
\definecolor{currentstroke}{rgb}{0.000000,0.000000,0.000000}%
\pgfsetstrokecolor{currentstroke}%
\pgfsetdash{}{0pt}%
\pgfpathmoveto{\pgfqpoint{0.487500in}{0.320625in}}%
\pgfpathlineto{\pgfqpoint{0.487500in}{1.485000in}}%
\pgfusepath{stroke}%
\end{pgfscope}%
\begin{pgfscope}%
\pgfsetrectcap%
\pgfsetmiterjoin%
\pgfsetlinewidth{0.803000pt}%
\definecolor{currentstroke}{rgb}{0.000000,0.000000,0.000000}%
\pgfsetstrokecolor{currentstroke}%
\pgfsetdash{}{0pt}%
\pgfpathmoveto{\pgfqpoint{2.333179in}{0.320625in}}%
\pgfpathlineto{\pgfqpoint{2.333179in}{1.485000in}}%
\pgfusepath{stroke}%
\end{pgfscope}%
\begin{pgfscope}%
\pgfsetrectcap%
\pgfsetmiterjoin%
\pgfsetlinewidth{0.803000pt}%
\definecolor{currentstroke}{rgb}{0.000000,0.000000,0.000000}%
\pgfsetstrokecolor{currentstroke}%
\pgfsetdash{}{0pt}%
\pgfpathmoveto{\pgfqpoint{0.487500in}{0.320625in}}%
\pgfpathlineto{\pgfqpoint{2.333179in}{0.320625in}}%
\pgfusepath{stroke}%
\end{pgfscope}%
\begin{pgfscope}%
\pgfsetrectcap%
\pgfsetmiterjoin%
\pgfsetlinewidth{0.803000pt}%
\definecolor{currentstroke}{rgb}{0.000000,0.000000,0.000000}%
\pgfsetstrokecolor{currentstroke}%
\pgfsetdash{}{0pt}%
\pgfpathmoveto{\pgfqpoint{0.487500in}{1.485000in}}%
\pgfpathlineto{\pgfqpoint{2.333179in}{1.485000in}}%
\pgfusepath{stroke}%
\end{pgfscope}%
\begin{pgfscope}%
\definecolor{textcolor}{rgb}{0.000000,0.000000,0.000000}%
\pgfsetstrokecolor{textcolor}%
\pgfsetfillcolor{textcolor}%
\pgftext[x=1.410340in,y=1.568333in,,base]{\color{textcolor}{\rmfamily\fontsize{8.000000}{9.600000}\selectfont\catcode`\^=\active\def^{\ifmmode\sp\else\^{}\fi}\catcode`\%=\active\def%{\%}(a) Initialization scale}}%
\end{pgfscope}%
\begin{pgfscope}%
\pgfsetbuttcap%
\pgfsetmiterjoin%
\definecolor{currentfill}{rgb}{1.000000,1.000000,1.000000}%
\pgfsetfillcolor{currentfill}%
\pgfsetfillopacity{0.800000}%
\pgfsetlinewidth{1.003750pt}%
\definecolor{currentstroke}{rgb}{0.800000,0.800000,0.800000}%
\pgfsetstrokecolor{currentstroke}%
\pgfsetstrokeopacity{0.800000}%
\pgfsetdash{}{0pt}%
\pgfpathmoveto{\pgfqpoint{1.409008in}{1.284674in}}%
\pgfpathlineto{\pgfqpoint{2.288457in}{1.284674in}}%
\pgfpathquadraticcurveto{\pgfqpoint{2.301235in}{1.284674in}}{\pgfqpoint{2.301235in}{1.297451in}}%
\pgfpathlineto{\pgfqpoint{2.301235in}{1.440278in}}%
\pgfpathquadraticcurveto{\pgfqpoint{2.301235in}{1.453056in}}{\pgfqpoint{2.288457in}{1.453056in}}%
\pgfpathlineto{\pgfqpoint{1.409008in}{1.453056in}}%
\pgfpathquadraticcurveto{\pgfqpoint{1.396230in}{1.453056in}}{\pgfqpoint{1.396230in}{1.440278in}}%
\pgfpathlineto{\pgfqpoint{1.396230in}{1.297451in}}%
\pgfpathquadraticcurveto{\pgfqpoint{1.396230in}{1.284674in}}{\pgfqpoint{1.409008in}{1.284674in}}%
\pgfpathlineto{\pgfqpoint{1.409008in}{1.284674in}}%
\pgfpathclose%
\pgfusepath{stroke,fill}%
\end{pgfscope}%
\begin{pgfscope}%
\pgfsetrectcap%
\pgfsetroundjoin%
\pgfsetlinewidth{1.254687pt}%
\definecolor{currentstroke}{rgb}{0.141176,0.345098,0.650980}%
\pgfsetstrokecolor{currentstroke}%
\pgfsetdash{}{0pt}%
\pgfpathmoveto{\pgfqpoint{1.412202in}{1.410904in}}%
\pgfpathlineto{\pgfqpoint{1.456924in}{1.410904in}}%
\pgfpathlineto{\pgfqpoint{1.501647in}{1.410904in}}%
\pgfusepath{stroke}%
\end{pgfscope}%
\begin{pgfscope}%
\definecolor{textcolor}{rgb}{0.000000,0.000000,0.000000}%
\pgfsetstrokecolor{textcolor}%
\pgfsetfillcolor{textcolor}%
\pgftext[x=1.552758in,y=1.388543in,left,base]{\color{textcolor}{\rmfamily\fontsize{4.600000}{5.520000}\selectfont\catcode`\^=\active\def^{\ifmmode\sp\else\^{}\fi}\catcode`\%=\active\def%{\%}$\alpha=0.1$}}%
\end{pgfscope}%
\begin{pgfscope}%
\pgfsetrectcap%
\pgfsetroundjoin%
\pgfsetlinewidth{1.254687pt}%
\definecolor{currentstroke}{rgb}{0.819608,0.286275,0.356863}%
\pgfsetstrokecolor{currentstroke}%
\pgfsetdash{}{0pt}%
\pgfpathmoveto{\pgfqpoint{1.412202in}{1.336296in}}%
\pgfpathlineto{\pgfqpoint{1.456924in}{1.336296in}}%
\pgfpathlineto{\pgfqpoint{1.501647in}{1.336296in}}%
\pgfusepath{stroke}%
\end{pgfscope}%
\begin{pgfscope}%
\definecolor{textcolor}{rgb}{0.000000,0.000000,0.000000}%
\pgfsetstrokecolor{textcolor}%
\pgfsetfillcolor{textcolor}%
\pgftext[x=1.552758in,y=1.313935in,left,base]{\color{textcolor}{\rmfamily\fontsize{4.600000}{5.520000}\selectfont\catcode`\^=\active\def^{\ifmmode\sp\else\^{}\fi}\catcode`\%=\active\def%{\%}$\alpha=0.5$}}%
\end{pgfscope}%
\begin{pgfscope}%
\pgfsetrectcap%
\pgfsetroundjoin%
\pgfsetlinewidth{1.254687pt}%
\definecolor{currentstroke}{rgb}{0.000000,0.529412,0.423529}%
\pgfsetstrokecolor{currentstroke}%
\pgfsetdash{}{0pt}%
\pgfpathmoveto{\pgfqpoint{1.905847in}{1.410904in}}%
\pgfpathlineto{\pgfqpoint{1.950569in}{1.410904in}}%
\pgfpathlineto{\pgfqpoint{1.995292in}{1.410904in}}%
\pgfusepath{stroke}%
\end{pgfscope}%
\begin{pgfscope}%
\definecolor{textcolor}{rgb}{0.000000,0.000000,0.000000}%
\pgfsetstrokecolor{textcolor}%
\pgfsetfillcolor{textcolor}%
\pgftext[x=2.046403in,y=1.388543in,left,base]{\color{textcolor}{\rmfamily\fontsize{4.600000}{5.520000}\selectfont\catcode`\^=\active\def^{\ifmmode\sp\else\^{}\fi}\catcode`\%=\active\def%{\%}$\alpha=1$}}%
\end{pgfscope}%
\begin{pgfscope}%
\pgfsetrectcap%
\pgfsetroundjoin%
\pgfsetlinewidth{1.254687pt}%
\definecolor{currentstroke}{rgb}{0.478431,0.243137,0.615686}%
\pgfsetstrokecolor{currentstroke}%
\pgfsetdash{}{0pt}%
\pgfpathmoveto{\pgfqpoint{1.905847in}{1.336296in}}%
\pgfpathlineto{\pgfqpoint{1.950569in}{1.336296in}}%
\pgfpathlineto{\pgfqpoint{1.995292in}{1.336296in}}%
\pgfusepath{stroke}%
\end{pgfscope}%
\begin{pgfscope}%
\definecolor{textcolor}{rgb}{0.000000,0.000000,0.000000}%
\pgfsetstrokecolor{textcolor}%
\pgfsetfillcolor{textcolor}%
\pgftext[x=2.046403in,y=1.313935in,left,base]{\color{textcolor}{\rmfamily\fontsize{4.600000}{5.520000}\selectfont\catcode`\^=\active\def^{\ifmmode\sp\else\^{}\fi}\catcode`\%=\active\def%{\%}$\alpha=5$}}%
\end{pgfscope}%
\begin{pgfscope}%
\pgfsetbuttcap%
\pgfsetmiterjoin%
\definecolor{currentfill}{rgb}{1.000000,1.000000,1.000000}%
\pgfsetfillcolor{currentfill}%
\pgfsetlinewidth{0.000000pt}%
\definecolor{currentstroke}{rgb}{0.000000,0.000000,0.000000}%
\pgfsetstrokecolor{currentstroke}%
\pgfsetstrokeopacity{0.000000}%
\pgfsetdash{}{0pt}%
\pgfpathmoveto{\pgfqpoint{2.554660in}{0.320625in}}%
\pgfpathlineto{\pgfqpoint{4.400340in}{0.320625in}}%
\pgfpathlineto{\pgfqpoint{4.400340in}{1.485000in}}%
\pgfpathlineto{\pgfqpoint{2.554660in}{1.485000in}}%
\pgfpathlineto{\pgfqpoint{2.554660in}{0.320625in}}%
\pgfpathclose%
\pgfusepath{fill}%
\end{pgfscope}%
\begin{pgfscope}%
\pgfpathrectangle{\pgfqpoint{2.554660in}{0.320625in}}{\pgfqpoint{1.845679in}{1.164375in}}%
\pgfusepath{clip}%
\pgfsetrectcap%
\pgfsetroundjoin%
\pgfsetlinewidth{0.803000pt}%
\definecolor{currentstroke}{rgb}{0.909804,0.909804,0.909804}%
\pgfsetstrokecolor{currentstroke}%
\pgfsetdash{}{0pt}%
\pgfpathmoveto{\pgfqpoint{2.638555in}{0.320625in}}%
\pgfpathlineto{\pgfqpoint{2.638555in}{1.485000in}}%
\pgfusepath{stroke}%
\end{pgfscope}%
\begin{pgfscope}%
\pgfsetbuttcap%
\pgfsetroundjoin%
\definecolor{currentfill}{rgb}{0.000000,0.000000,0.000000}%
\pgfsetfillcolor{currentfill}%
\pgfsetlinewidth{0.803000pt}%
\definecolor{currentstroke}{rgb}{0.000000,0.000000,0.000000}%
\pgfsetstrokecolor{currentstroke}%
\pgfsetdash{}{0pt}%
\pgfsys@defobject{currentmarker}{\pgfqpoint{0.000000in}{-0.048611in}}{\pgfqpoint{0.000000in}{0.000000in}}{%
\pgfpathmoveto{\pgfqpoint{0.000000in}{0.000000in}}%
\pgfpathlineto{\pgfqpoint{0.000000in}{-0.048611in}}%
\pgfusepath{stroke,fill}%
}%
\begin{pgfscope}%
\pgfsys@transformshift{2.638555in}{0.320625in}%
\pgfsys@useobject{currentmarker}{}%
\end{pgfscope}%
\end{pgfscope}%
\begin{pgfscope}%
\definecolor{textcolor}{rgb}{0.000000,0.000000,0.000000}%
\pgfsetstrokecolor{textcolor}%
\pgfsetfillcolor{textcolor}%
\pgftext[x=2.638555in,y=0.223403in,,top]{\color{textcolor}{\rmfamily\fontsize{6.000000}{7.200000}\selectfont\catcode`\^=\active\def^{\ifmmode\sp\else\^{}\fi}\catcode`\%=\active\def%{\%}0}}%
\end{pgfscope}%
\begin{pgfscope}%
\pgfpathrectangle{\pgfqpoint{2.554660in}{0.320625in}}{\pgfqpoint{1.845679in}{1.164375in}}%
\pgfusepath{clip}%
\pgfsetrectcap%
\pgfsetroundjoin%
\pgfsetlinewidth{0.803000pt}%
\definecolor{currentstroke}{rgb}{0.909804,0.909804,0.909804}%
\pgfsetstrokecolor{currentstroke}%
\pgfsetdash{}{0pt}%
\pgfpathmoveto{\pgfqpoint{3.058027in}{0.320625in}}%
\pgfpathlineto{\pgfqpoint{3.058027in}{1.485000in}}%
\pgfusepath{stroke}%
\end{pgfscope}%
\begin{pgfscope}%
\pgfsetbuttcap%
\pgfsetroundjoin%
\definecolor{currentfill}{rgb}{0.000000,0.000000,0.000000}%
\pgfsetfillcolor{currentfill}%
\pgfsetlinewidth{0.803000pt}%
\definecolor{currentstroke}{rgb}{0.000000,0.000000,0.000000}%
\pgfsetstrokecolor{currentstroke}%
\pgfsetdash{}{0pt}%
\pgfsys@defobject{currentmarker}{\pgfqpoint{0.000000in}{-0.048611in}}{\pgfqpoint{0.000000in}{0.000000in}}{%
\pgfpathmoveto{\pgfqpoint{0.000000in}{0.000000in}}%
\pgfpathlineto{\pgfqpoint{0.000000in}{-0.048611in}}%
\pgfusepath{stroke,fill}%
}%
\begin{pgfscope}%
\pgfsys@transformshift{3.058027in}{0.320625in}%
\pgfsys@useobject{currentmarker}{}%
\end{pgfscope}%
\end{pgfscope}%
\begin{pgfscope}%
\definecolor{textcolor}{rgb}{0.000000,0.000000,0.000000}%
\pgfsetstrokecolor{textcolor}%
\pgfsetfillcolor{textcolor}%
\pgftext[x=3.058027in,y=0.223403in,,top]{\color{textcolor}{\rmfamily\fontsize{6.000000}{7.200000}\selectfont\catcode`\^=\active\def^{\ifmmode\sp\else\^{}\fi}\catcode`\%=\active\def%{\%}10}}%
\end{pgfscope}%
\begin{pgfscope}%
\pgfpathrectangle{\pgfqpoint{2.554660in}{0.320625in}}{\pgfqpoint{1.845679in}{1.164375in}}%
\pgfusepath{clip}%
\pgfsetrectcap%
\pgfsetroundjoin%
\pgfsetlinewidth{0.803000pt}%
\definecolor{currentstroke}{rgb}{0.909804,0.909804,0.909804}%
\pgfsetstrokecolor{currentstroke}%
\pgfsetdash{}{0pt}%
\pgfpathmoveto{\pgfqpoint{3.477500in}{0.320625in}}%
\pgfpathlineto{\pgfqpoint{3.477500in}{1.485000in}}%
\pgfusepath{stroke}%
\end{pgfscope}%
\begin{pgfscope}%
\pgfsetbuttcap%
\pgfsetroundjoin%
\definecolor{currentfill}{rgb}{0.000000,0.000000,0.000000}%
\pgfsetfillcolor{currentfill}%
\pgfsetlinewidth{0.803000pt}%
\definecolor{currentstroke}{rgb}{0.000000,0.000000,0.000000}%
\pgfsetstrokecolor{currentstroke}%
\pgfsetdash{}{0pt}%
\pgfsys@defobject{currentmarker}{\pgfqpoint{0.000000in}{-0.048611in}}{\pgfqpoint{0.000000in}{0.000000in}}{%
\pgfpathmoveto{\pgfqpoint{0.000000in}{0.000000in}}%
\pgfpathlineto{\pgfqpoint{0.000000in}{-0.048611in}}%
\pgfusepath{stroke,fill}%
}%
\begin{pgfscope}%
\pgfsys@transformshift{3.477500in}{0.320625in}%
\pgfsys@useobject{currentmarker}{}%
\end{pgfscope}%
\end{pgfscope}%
\begin{pgfscope}%
\definecolor{textcolor}{rgb}{0.000000,0.000000,0.000000}%
\pgfsetstrokecolor{textcolor}%
\pgfsetfillcolor{textcolor}%
\pgftext[x=3.477500in,y=0.223403in,,top]{\color{textcolor}{\rmfamily\fontsize{6.000000}{7.200000}\selectfont\catcode`\^=\active\def^{\ifmmode\sp\else\^{}\fi}\catcode`\%=\active\def%{\%}20}}%
\end{pgfscope}%
\begin{pgfscope}%
\pgfpathrectangle{\pgfqpoint{2.554660in}{0.320625in}}{\pgfqpoint{1.845679in}{1.164375in}}%
\pgfusepath{clip}%
\pgfsetrectcap%
\pgfsetroundjoin%
\pgfsetlinewidth{0.803000pt}%
\definecolor{currentstroke}{rgb}{0.909804,0.909804,0.909804}%
\pgfsetstrokecolor{currentstroke}%
\pgfsetdash{}{0pt}%
\pgfpathmoveto{\pgfqpoint{3.896973in}{0.320625in}}%
\pgfpathlineto{\pgfqpoint{3.896973in}{1.485000in}}%
\pgfusepath{stroke}%
\end{pgfscope}%
\begin{pgfscope}%
\pgfsetbuttcap%
\pgfsetroundjoin%
\definecolor{currentfill}{rgb}{0.000000,0.000000,0.000000}%
\pgfsetfillcolor{currentfill}%
\pgfsetlinewidth{0.803000pt}%
\definecolor{currentstroke}{rgb}{0.000000,0.000000,0.000000}%
\pgfsetstrokecolor{currentstroke}%
\pgfsetdash{}{0pt}%
\pgfsys@defobject{currentmarker}{\pgfqpoint{0.000000in}{-0.048611in}}{\pgfqpoint{0.000000in}{0.000000in}}{%
\pgfpathmoveto{\pgfqpoint{0.000000in}{0.000000in}}%
\pgfpathlineto{\pgfqpoint{0.000000in}{-0.048611in}}%
\pgfusepath{stroke,fill}%
}%
\begin{pgfscope}%
\pgfsys@transformshift{3.896973in}{0.320625in}%
\pgfsys@useobject{currentmarker}{}%
\end{pgfscope}%
\end{pgfscope}%
\begin{pgfscope}%
\definecolor{textcolor}{rgb}{0.000000,0.000000,0.000000}%
\pgfsetstrokecolor{textcolor}%
\pgfsetfillcolor{textcolor}%
\pgftext[x=3.896973in,y=0.223403in,,top]{\color{textcolor}{\rmfamily\fontsize{6.000000}{7.200000}\selectfont\catcode`\^=\active\def^{\ifmmode\sp\else\^{}\fi}\catcode`\%=\active\def%{\%}30}}%
\end{pgfscope}%
\begin{pgfscope}%
\pgfpathrectangle{\pgfqpoint{2.554660in}{0.320625in}}{\pgfqpoint{1.845679in}{1.164375in}}%
\pgfusepath{clip}%
\pgfsetrectcap%
\pgfsetroundjoin%
\pgfsetlinewidth{0.803000pt}%
\definecolor{currentstroke}{rgb}{0.909804,0.909804,0.909804}%
\pgfsetstrokecolor{currentstroke}%
\pgfsetdash{}{0pt}%
\pgfpathmoveto{\pgfqpoint{4.316445in}{0.320625in}}%
\pgfpathlineto{\pgfqpoint{4.316445in}{1.485000in}}%
\pgfusepath{stroke}%
\end{pgfscope}%
\begin{pgfscope}%
\pgfsetbuttcap%
\pgfsetroundjoin%
\definecolor{currentfill}{rgb}{0.000000,0.000000,0.000000}%
\pgfsetfillcolor{currentfill}%
\pgfsetlinewidth{0.803000pt}%
\definecolor{currentstroke}{rgb}{0.000000,0.000000,0.000000}%
\pgfsetstrokecolor{currentstroke}%
\pgfsetdash{}{0pt}%
\pgfsys@defobject{currentmarker}{\pgfqpoint{0.000000in}{-0.048611in}}{\pgfqpoint{0.000000in}{0.000000in}}{%
\pgfpathmoveto{\pgfqpoint{0.000000in}{0.000000in}}%
\pgfpathlineto{\pgfqpoint{0.000000in}{-0.048611in}}%
\pgfusepath{stroke,fill}%
}%
\begin{pgfscope}%
\pgfsys@transformshift{4.316445in}{0.320625in}%
\pgfsys@useobject{currentmarker}{}%
\end{pgfscope}%
\end{pgfscope}%
\begin{pgfscope}%
\definecolor{textcolor}{rgb}{0.000000,0.000000,0.000000}%
\pgfsetstrokecolor{textcolor}%
\pgfsetfillcolor{textcolor}%
\pgftext[x=4.316445in,y=0.223403in,,top]{\color{textcolor}{\rmfamily\fontsize{6.000000}{7.200000}\selectfont\catcode`\^=\active\def^{\ifmmode\sp\else\^{}\fi}\catcode`\%=\active\def%{\%}40}}%
\end{pgfscope}%
\begin{pgfscope}%
\definecolor{textcolor}{rgb}{0.000000,0.000000,0.000000}%
\pgfsetstrokecolor{textcolor}%
\pgfsetfillcolor{textcolor}%
\pgftext[x=3.477500in,y=0.087200in,,top]{\color{textcolor}{\rmfamily\fontsize{6.000000}{7.200000}\selectfont\catcode`\^=\active\def^{\ifmmode\sp\else\^{}\fi}\catcode`\%=\active\def%{\%}time}}%
\end{pgfscope}%
\begin{pgfscope}%
\pgfpathrectangle{\pgfqpoint{2.554660in}{0.320625in}}{\pgfqpoint{1.845679in}{1.164375in}}%
\pgfusepath{clip}%
\pgfsetrectcap%
\pgfsetroundjoin%
\pgfsetlinewidth{0.803000pt}%
\definecolor{currentstroke}{rgb}{0.909804,0.909804,0.909804}%
\pgfsetstrokecolor{currentstroke}%
\pgfsetdash{}{0pt}%
\pgfpathmoveto{\pgfqpoint{2.554660in}{0.320625in}}%
\pgfpathlineto{\pgfqpoint{4.400340in}{0.320625in}}%
\pgfusepath{stroke}%
\end{pgfscope}%
\begin{pgfscope}%
\pgfsetbuttcap%
\pgfsetroundjoin%
\definecolor{currentfill}{rgb}{0.000000,0.000000,0.000000}%
\pgfsetfillcolor{currentfill}%
\pgfsetlinewidth{0.803000pt}%
\definecolor{currentstroke}{rgb}{0.000000,0.000000,0.000000}%
\pgfsetstrokecolor{currentstroke}%
\pgfsetdash{}{0pt}%
\pgfsys@defobject{currentmarker}{\pgfqpoint{-0.048611in}{0.000000in}}{\pgfqpoint{-0.000000in}{0.000000in}}{%
\pgfpathmoveto{\pgfqpoint{-0.000000in}{0.000000in}}%
\pgfpathlineto{\pgfqpoint{-0.048611in}{0.000000in}}%
\pgfusepath{stroke,fill}%
}%
\begin{pgfscope}%
\pgfsys@transformshift{2.554660in}{0.320625in}%
\pgfsys@useobject{currentmarker}{}%
\end{pgfscope}%
\end{pgfscope}%
\begin{pgfscope}%
\pgfpathrectangle{\pgfqpoint{2.554660in}{0.320625in}}{\pgfqpoint{1.845679in}{1.164375in}}%
\pgfusepath{clip}%
\pgfsetrectcap%
\pgfsetroundjoin%
\pgfsetlinewidth{0.803000pt}%
\definecolor{currentstroke}{rgb}{0.909804,0.909804,0.909804}%
\pgfsetstrokecolor{currentstroke}%
\pgfsetdash{}{0pt}%
\pgfpathmoveto{\pgfqpoint{2.554660in}{0.536250in}}%
\pgfpathlineto{\pgfqpoint{4.400340in}{0.536250in}}%
\pgfusepath{stroke}%
\end{pgfscope}%
\begin{pgfscope}%
\pgfsetbuttcap%
\pgfsetroundjoin%
\definecolor{currentfill}{rgb}{0.000000,0.000000,0.000000}%
\pgfsetfillcolor{currentfill}%
\pgfsetlinewidth{0.803000pt}%
\definecolor{currentstroke}{rgb}{0.000000,0.000000,0.000000}%
\pgfsetstrokecolor{currentstroke}%
\pgfsetdash{}{0pt}%
\pgfsys@defobject{currentmarker}{\pgfqpoint{-0.048611in}{0.000000in}}{\pgfqpoint{-0.000000in}{0.000000in}}{%
\pgfpathmoveto{\pgfqpoint{-0.000000in}{0.000000in}}%
\pgfpathlineto{\pgfqpoint{-0.048611in}{0.000000in}}%
\pgfusepath{stroke,fill}%
}%
\begin{pgfscope}%
\pgfsys@transformshift{2.554660in}{0.536250in}%
\pgfsys@useobject{currentmarker}{}%
\end{pgfscope}%
\end{pgfscope}%
\begin{pgfscope}%
\pgfpathrectangle{\pgfqpoint{2.554660in}{0.320625in}}{\pgfqpoint{1.845679in}{1.164375in}}%
\pgfusepath{clip}%
\pgfsetrectcap%
\pgfsetroundjoin%
\pgfsetlinewidth{0.803000pt}%
\definecolor{currentstroke}{rgb}{0.909804,0.909804,0.909804}%
\pgfsetstrokecolor{currentstroke}%
\pgfsetdash{}{0pt}%
\pgfpathmoveto{\pgfqpoint{2.554660in}{0.751875in}}%
\pgfpathlineto{\pgfqpoint{4.400340in}{0.751875in}}%
\pgfusepath{stroke}%
\end{pgfscope}%
\begin{pgfscope}%
\pgfsetbuttcap%
\pgfsetroundjoin%
\definecolor{currentfill}{rgb}{0.000000,0.000000,0.000000}%
\pgfsetfillcolor{currentfill}%
\pgfsetlinewidth{0.803000pt}%
\definecolor{currentstroke}{rgb}{0.000000,0.000000,0.000000}%
\pgfsetstrokecolor{currentstroke}%
\pgfsetdash{}{0pt}%
\pgfsys@defobject{currentmarker}{\pgfqpoint{-0.048611in}{0.000000in}}{\pgfqpoint{-0.000000in}{0.000000in}}{%
\pgfpathmoveto{\pgfqpoint{-0.000000in}{0.000000in}}%
\pgfpathlineto{\pgfqpoint{-0.048611in}{0.000000in}}%
\pgfusepath{stroke,fill}%
}%
\begin{pgfscope}%
\pgfsys@transformshift{2.554660in}{0.751875in}%
\pgfsys@useobject{currentmarker}{}%
\end{pgfscope}%
\end{pgfscope}%
\begin{pgfscope}%
\pgfpathrectangle{\pgfqpoint{2.554660in}{0.320625in}}{\pgfqpoint{1.845679in}{1.164375in}}%
\pgfusepath{clip}%
\pgfsetrectcap%
\pgfsetroundjoin%
\pgfsetlinewidth{0.803000pt}%
\definecolor{currentstroke}{rgb}{0.909804,0.909804,0.909804}%
\pgfsetstrokecolor{currentstroke}%
\pgfsetdash{}{0pt}%
\pgfpathmoveto{\pgfqpoint{2.554660in}{0.967500in}}%
\pgfpathlineto{\pgfqpoint{4.400340in}{0.967500in}}%
\pgfusepath{stroke}%
\end{pgfscope}%
\begin{pgfscope}%
\pgfsetbuttcap%
\pgfsetroundjoin%
\definecolor{currentfill}{rgb}{0.000000,0.000000,0.000000}%
\pgfsetfillcolor{currentfill}%
\pgfsetlinewidth{0.803000pt}%
\definecolor{currentstroke}{rgb}{0.000000,0.000000,0.000000}%
\pgfsetstrokecolor{currentstroke}%
\pgfsetdash{}{0pt}%
\pgfsys@defobject{currentmarker}{\pgfqpoint{-0.048611in}{0.000000in}}{\pgfqpoint{-0.000000in}{0.000000in}}{%
\pgfpathmoveto{\pgfqpoint{-0.000000in}{0.000000in}}%
\pgfpathlineto{\pgfqpoint{-0.048611in}{0.000000in}}%
\pgfusepath{stroke,fill}%
}%
\begin{pgfscope}%
\pgfsys@transformshift{2.554660in}{0.967500in}%
\pgfsys@useobject{currentmarker}{}%
\end{pgfscope}%
\end{pgfscope}%
\begin{pgfscope}%
\pgfpathrectangle{\pgfqpoint{2.554660in}{0.320625in}}{\pgfqpoint{1.845679in}{1.164375in}}%
\pgfusepath{clip}%
\pgfsetrectcap%
\pgfsetroundjoin%
\pgfsetlinewidth{0.803000pt}%
\definecolor{currentstroke}{rgb}{0.909804,0.909804,0.909804}%
\pgfsetstrokecolor{currentstroke}%
\pgfsetdash{}{0pt}%
\pgfpathmoveto{\pgfqpoint{2.554660in}{1.183125in}}%
\pgfpathlineto{\pgfqpoint{4.400340in}{1.183125in}}%
\pgfusepath{stroke}%
\end{pgfscope}%
\begin{pgfscope}%
\pgfsetbuttcap%
\pgfsetroundjoin%
\definecolor{currentfill}{rgb}{0.000000,0.000000,0.000000}%
\pgfsetfillcolor{currentfill}%
\pgfsetlinewidth{0.803000pt}%
\definecolor{currentstroke}{rgb}{0.000000,0.000000,0.000000}%
\pgfsetstrokecolor{currentstroke}%
\pgfsetdash{}{0pt}%
\pgfsys@defobject{currentmarker}{\pgfqpoint{-0.048611in}{0.000000in}}{\pgfqpoint{-0.000000in}{0.000000in}}{%
\pgfpathmoveto{\pgfqpoint{-0.000000in}{0.000000in}}%
\pgfpathlineto{\pgfqpoint{-0.048611in}{0.000000in}}%
\pgfusepath{stroke,fill}%
}%
\begin{pgfscope}%
\pgfsys@transformshift{2.554660in}{1.183125in}%
\pgfsys@useobject{currentmarker}{}%
\end{pgfscope}%
\end{pgfscope}%
\begin{pgfscope}%
\pgfpathrectangle{\pgfqpoint{2.554660in}{0.320625in}}{\pgfqpoint{1.845679in}{1.164375in}}%
\pgfusepath{clip}%
\pgfsetrectcap%
\pgfsetroundjoin%
\pgfsetlinewidth{0.803000pt}%
\definecolor{currentstroke}{rgb}{0.909804,0.909804,0.909804}%
\pgfsetstrokecolor{currentstroke}%
\pgfsetdash{}{0pt}%
\pgfpathmoveto{\pgfqpoint{2.554660in}{1.398750in}}%
\pgfpathlineto{\pgfqpoint{4.400340in}{1.398750in}}%
\pgfusepath{stroke}%
\end{pgfscope}%
\begin{pgfscope}%
\pgfsetbuttcap%
\pgfsetroundjoin%
\definecolor{currentfill}{rgb}{0.000000,0.000000,0.000000}%
\pgfsetfillcolor{currentfill}%
\pgfsetlinewidth{0.803000pt}%
\definecolor{currentstroke}{rgb}{0.000000,0.000000,0.000000}%
\pgfsetstrokecolor{currentstroke}%
\pgfsetdash{}{0pt}%
\pgfsys@defobject{currentmarker}{\pgfqpoint{-0.048611in}{0.000000in}}{\pgfqpoint{-0.000000in}{0.000000in}}{%
\pgfpathmoveto{\pgfqpoint{-0.000000in}{0.000000in}}%
\pgfpathlineto{\pgfqpoint{-0.048611in}{0.000000in}}%
\pgfusepath{stroke,fill}%
}%
\begin{pgfscope}%
\pgfsys@transformshift{2.554660in}{1.398750in}%
\pgfsys@useobject{currentmarker}{}%
\end{pgfscope}%
\end{pgfscope}%
\begin{pgfscope}%
\pgfpathrectangle{\pgfqpoint{2.554660in}{0.320625in}}{\pgfqpoint{1.845679in}{1.164375in}}%
\pgfusepath{clip}%
\pgfsetrectcap%
\pgfsetroundjoin%
\pgfsetlinewidth{1.154312pt}%
\definecolor{currentstroke}{rgb}{0.146785,0.068738,0.334011}%
\pgfsetstrokecolor{currentstroke}%
\pgfsetdash{}{0pt}%
\pgfpathmoveto{\pgfqpoint{2.638555in}{1.468525in}}%
\pgfpathlineto{\pgfqpoint{2.653237in}{1.443410in}}%
\pgfpathlineto{\pgfqpoint{2.665821in}{1.419142in}}%
\pgfpathlineto{\pgfqpoint{2.701476in}{1.348595in}}%
\pgfpathlineto{\pgfqpoint{2.716157in}{1.322658in}}%
\pgfpathlineto{\pgfqpoint{2.728742in}{1.302884in}}%
\pgfpathlineto{\pgfqpoint{2.739228in}{1.288502in}}%
\pgfpathlineto{\pgfqpoint{2.749715in}{1.276171in}}%
\pgfpathlineto{\pgfqpoint{2.760202in}{1.265775in}}%
\pgfpathlineto{\pgfqpoint{2.770689in}{1.256975in}}%
\pgfpathlineto{\pgfqpoint{2.783273in}{1.247925in}}%
\pgfpathlineto{\pgfqpoint{2.800052in}{1.237402in}}%
\pgfpathlineto{\pgfqpoint{2.825220in}{1.223142in}}%
\pgfpathlineto{\pgfqpoint{2.877654in}{1.195041in}}%
\pgfpathlineto{\pgfqpoint{3.280348in}{0.982203in}}%
\pgfpathlineto{\pgfqpoint{3.836149in}{0.689041in}}%
\pgfpathlineto{\pgfqpoint{3.842441in}{0.685643in}}%
\pgfpathlineto{\pgfqpoint{3.850831in}{0.680858in}}%
\pgfpathlineto{\pgfqpoint{3.863415in}{0.674351in}}%
\pgfpathlineto{\pgfqpoint{3.869707in}{0.670979in}}%
\pgfpathlineto{\pgfqpoint{3.880194in}{0.666874in}}%
\pgfpathlineto{\pgfqpoint{3.886486in}{0.662037in}}%
\pgfpathlineto{\pgfqpoint{3.890680in}{0.660034in}}%
\pgfpathlineto{\pgfqpoint{3.896973in}{0.657723in}}%
\pgfpathlineto{\pgfqpoint{3.901167in}{0.655092in}}%
\pgfpathlineto{\pgfqpoint{3.907459in}{0.651002in}}%
\pgfpathlineto{\pgfqpoint{3.909557in}{0.650634in}}%
\pgfpathlineto{\pgfqpoint{3.915849in}{0.651511in}}%
\pgfpathlineto{\pgfqpoint{3.917946in}{0.650640in}}%
\pgfpathlineto{\pgfqpoint{3.922141in}{0.645393in}}%
\pgfpathlineto{\pgfqpoint{3.924238in}{0.642673in}}%
\pgfpathlineto{\pgfqpoint{3.926336in}{0.640913in}}%
\pgfpathlineto{\pgfqpoint{3.928433in}{0.640399in}}%
\pgfpathlineto{\pgfqpoint{3.934725in}{0.641505in}}%
\pgfpathlineto{\pgfqpoint{3.936822in}{0.640828in}}%
\pgfpathlineto{\pgfqpoint{3.938920in}{0.639151in}}%
\pgfpathlineto{\pgfqpoint{3.945212in}{0.631011in}}%
\pgfpathlineto{\pgfqpoint{3.947309in}{0.630552in}}%
\pgfpathlineto{\pgfqpoint{3.949407in}{0.632665in}}%
\pgfpathlineto{\pgfqpoint{3.953601in}{0.639659in}}%
\pgfpathlineto{\pgfqpoint{3.955699in}{0.641515in}}%
\pgfpathlineto{\pgfqpoint{3.957796in}{0.640766in}}%
\pgfpathlineto{\pgfqpoint{3.959893in}{0.636562in}}%
\pgfpathlineto{\pgfqpoint{3.964088in}{0.623731in}}%
\pgfpathlineto{\pgfqpoint{3.966185in}{0.619870in}}%
\pgfpathlineto{\pgfqpoint{3.968283in}{0.620677in}}%
\pgfpathlineto{\pgfqpoint{3.972478in}{0.627544in}}%
\pgfpathlineto{\pgfqpoint{3.974575in}{0.629296in}}%
\pgfpathlineto{\pgfqpoint{3.976672in}{0.628941in}}%
\pgfpathlineto{\pgfqpoint{3.978770in}{0.626267in}}%
\pgfpathlineto{\pgfqpoint{3.980867in}{0.621231in}}%
\pgfpathlineto{\pgfqpoint{3.982964in}{0.614572in}}%
\pgfpathlineto{\pgfqpoint{3.985062in}{0.609883in}}%
\pgfpathlineto{\pgfqpoint{3.987159in}{0.612919in}}%
\pgfpathlineto{\pgfqpoint{3.991354in}{0.630104in}}%
\pgfpathlineto{\pgfqpoint{3.993451in}{0.635969in}}%
\pgfpathlineto{\pgfqpoint{3.995549in}{0.638454in}}%
\pgfpathlineto{\pgfqpoint{3.997646in}{0.636612in}}%
\pgfpathlineto{\pgfqpoint{3.999743in}{0.629818in}}%
\pgfpathlineto{\pgfqpoint{4.006035in}{0.598872in}}%
\pgfpathlineto{\pgfqpoint{4.008133in}{0.602209in}}%
\pgfpathlineto{\pgfqpoint{4.010230in}{0.609849in}}%
\pgfpathlineto{\pgfqpoint{4.012327in}{0.615239in}}%
\pgfpathlineto{\pgfqpoint{4.014425in}{0.617247in}}%
\pgfpathlineto{\pgfqpoint{4.016522in}{0.615772in}}%
\pgfpathlineto{\pgfqpoint{4.018620in}{0.610549in}}%
\pgfpathlineto{\pgfqpoint{4.022814in}{0.590831in}}%
\pgfpathlineto{\pgfqpoint{4.024912in}{0.593831in}}%
\pgfpathlineto{\pgfqpoint{4.029106in}{0.620684in}}%
\pgfpathlineto{\pgfqpoint{4.031204in}{0.628749in}}%
\pgfpathlineto{\pgfqpoint{4.033301in}{0.632580in}}%
\pgfpathlineto{\pgfqpoint{4.035398in}{0.631604in}}%
\pgfpathlineto{\pgfqpoint{4.037496in}{0.624811in}}%
\pgfpathlineto{\pgfqpoint{4.039593in}{0.613804in}}%
\pgfpathlineto{\pgfqpoint{4.043788in}{0.580929in}}%
\pgfpathlineto{\pgfqpoint{4.045885in}{0.582395in}}%
\pgfpathlineto{\pgfqpoint{4.047983in}{0.594632in}}%
\pgfpathlineto{\pgfqpoint{4.050080in}{0.602707in}}%
\pgfpathlineto{\pgfqpoint{4.052177in}{0.605898in}}%
\pgfpathlineto{\pgfqpoint{4.054275in}{0.604660in}}%
\pgfpathlineto{\pgfqpoint{4.056372in}{0.598677in}}%
\pgfpathlineto{\pgfqpoint{4.058469in}{0.586724in}}%
\pgfpathlineto{\pgfqpoint{4.060567in}{0.571082in}}%
\pgfpathlineto{\pgfqpoint{4.062664in}{0.579274in}}%
\pgfpathlineto{\pgfqpoint{4.064762in}{0.599248in}}%
\pgfpathlineto{\pgfqpoint{4.066859in}{0.613108in}}%
\pgfpathlineto{\pgfqpoint{4.068956in}{0.621499in}}%
\pgfpathlineto{\pgfqpoint{4.071054in}{0.625199in}}%
\pgfpathlineto{\pgfqpoint{4.073151in}{0.623658in}}%
\pgfpathlineto{\pgfqpoint{4.075248in}{0.615914in}}%
\pgfpathlineto{\pgfqpoint{4.077346in}{0.602908in}}%
\pgfpathlineto{\pgfqpoint{4.081540in}{0.558995in}}%
\pgfpathlineto{\pgfqpoint{4.085735in}{0.592544in}}%
\pgfpathlineto{\pgfqpoint{4.087832in}{0.601240in}}%
\pgfpathlineto{\pgfqpoint{4.089930in}{0.604498in}}%
\pgfpathlineto{\pgfqpoint{4.092027in}{0.603243in}}%
\pgfpathlineto{\pgfqpoint{4.094125in}{0.597098in}}%
\pgfpathlineto{\pgfqpoint{4.096222in}{0.583929in}}%
\pgfpathlineto{\pgfqpoint{4.098319in}{0.558730in}}%
\pgfpathlineto{\pgfqpoint{4.100417in}{0.561807in}}%
\pgfpathlineto{\pgfqpoint{4.102514in}{0.590666in}}%
\pgfpathlineto{\pgfqpoint{4.104611in}{0.607591in}}%
\pgfpathlineto{\pgfqpoint{4.106709in}{0.617559in}}%
\pgfpathlineto{\pgfqpoint{4.108806in}{0.622531in}}%
\pgfpathlineto{\pgfqpoint{4.110903in}{0.622551in}}%
\pgfpathlineto{\pgfqpoint{4.113001in}{0.616203in}}%
\pgfpathlineto{\pgfqpoint{4.115098in}{0.603537in}}%
\pgfpathlineto{\pgfqpoint{4.117196in}{0.580539in}}%
\pgfpathlineto{\pgfqpoint{4.119293in}{0.539317in}}%
\pgfpathlineto{\pgfqpoint{4.121390in}{0.574225in}}%
\pgfpathlineto{\pgfqpoint{4.123488in}{0.595942in}}%
\pgfpathlineto{\pgfqpoint{4.125585in}{0.606619in}}%
\pgfpathlineto{\pgfqpoint{4.127682in}{0.611502in}}%
\pgfpathlineto{\pgfqpoint{4.129780in}{0.612118in}}%
\pgfpathlineto{\pgfqpoint{4.131877in}{0.608604in}}%
\pgfpathlineto{\pgfqpoint{4.133974in}{0.599852in}}%
\pgfpathlineto{\pgfqpoint{4.136072in}{0.581834in}}%
\pgfpathlineto{\pgfqpoint{4.138169in}{0.537308in}}%
\pgfpathlineto{\pgfqpoint{4.142364in}{0.599917in}}%
\pgfpathlineto{\pgfqpoint{4.144461in}{0.615402in}}%
\pgfpathlineto{\pgfqpoint{4.146559in}{0.624312in}}%
\pgfpathlineto{\pgfqpoint{4.148656in}{0.628438in}}%
\pgfpathlineto{\pgfqpoint{4.150753in}{0.627646in}}%
\pgfpathlineto{\pgfqpoint{4.152851in}{0.620488in}}%
\pgfpathlineto{\pgfqpoint{4.154948in}{0.606603in}}%
\pgfpathlineto{\pgfqpoint{4.157045in}{0.579760in}}%
\pgfpathlineto{\pgfqpoint{4.159143in}{0.521751in}}%
\pgfpathlineto{\pgfqpoint{4.161240in}{0.583101in}}%
\pgfpathlineto{\pgfqpoint{4.163338in}{0.603877in}}%
\pgfpathlineto{\pgfqpoint{4.165435in}{0.613922in}}%
\pgfpathlineto{\pgfqpoint{4.167532in}{0.618381in}}%
\pgfpathlineto{\pgfqpoint{4.169630in}{0.618649in}}%
\pgfpathlineto{\pgfqpoint{4.171727in}{0.614770in}}%
\pgfpathlineto{\pgfqpoint{4.173824in}{0.605481in}}%
\pgfpathlineto{\pgfqpoint{4.175922in}{0.586084in}}%
\pgfpathlineto{\pgfqpoint{4.178019in}{0.524632in}}%
\pgfpathlineto{\pgfqpoint{4.180116in}{0.577773in}}%
\pgfpathlineto{\pgfqpoint{4.182214in}{0.607152in}}%
\pgfpathlineto{\pgfqpoint{4.184311in}{0.622371in}}%
\pgfpathlineto{\pgfqpoint{4.186409in}{0.631078in}}%
\pgfpathlineto{\pgfqpoint{4.188506in}{0.635026in}}%
\pgfpathlineto{\pgfqpoint{4.190603in}{0.634033in}}%
\pgfpathlineto{\pgfqpoint{4.192701in}{0.626801in}}%
\pgfpathlineto{\pgfqpoint{4.194798in}{0.613682in}}%
\pgfpathlineto{\pgfqpoint{4.196895in}{0.590146in}}%
\pgfpathlineto{\pgfqpoint{4.198993in}{0.519352in}}%
\pgfpathlineto{\pgfqpoint{4.201090in}{0.574781in}}%
\pgfpathlineto{\pgfqpoint{4.203187in}{0.599911in}}%
\pgfpathlineto{\pgfqpoint{4.205285in}{0.610880in}}%
\pgfpathlineto{\pgfqpoint{4.207382in}{0.615285in}}%
\pgfpathlineto{\pgfqpoint{4.209480in}{0.614868in}}%
\pgfpathlineto{\pgfqpoint{4.211577in}{0.609502in}}%
\pgfpathlineto{\pgfqpoint{4.213674in}{0.596990in}}%
\pgfpathlineto{\pgfqpoint{4.217869in}{0.539882in}}%
\pgfpathlineto{\pgfqpoint{4.219966in}{0.594136in}}%
\pgfpathlineto{\pgfqpoint{4.222064in}{0.615163in}}%
\pgfpathlineto{\pgfqpoint{4.224161in}{0.626927in}}%
\pgfpathlineto{\pgfqpoint{4.226258in}{0.633179in}}%
\pgfpathlineto{\pgfqpoint{4.228356in}{0.634551in}}%
\pgfpathlineto{\pgfqpoint{4.230453in}{0.629867in}}%
\pgfpathlineto{\pgfqpoint{4.232551in}{0.619761in}}%
\pgfpathlineto{\pgfqpoint{4.234648in}{0.603302in}}%
\pgfpathlineto{\pgfqpoint{4.236745in}{0.572560in}}%
\pgfpathlineto{\pgfqpoint{4.238843in}{0.510999in}}%
\pgfpathlineto{\pgfqpoint{4.240940in}{0.578158in}}%
\pgfpathlineto{\pgfqpoint{4.243037in}{0.596133in}}%
\pgfpathlineto{\pgfqpoint{4.245135in}{0.603430in}}%
\pgfpathlineto{\pgfqpoint{4.247232in}{0.604530in}}%
\pgfpathlineto{\pgfqpoint{4.249329in}{0.599998in}}%
\pgfpathlineto{\pgfqpoint{4.251427in}{0.587777in}}%
\pgfpathlineto{\pgfqpoint{4.253524in}{0.556950in}}%
\pgfpathlineto{\pgfqpoint{4.255621in}{0.536756in}}%
\pgfpathlineto{\pgfqpoint{4.257719in}{0.588112in}}%
\pgfpathlineto{\pgfqpoint{4.259816in}{0.608730in}}%
\pgfpathlineto{\pgfqpoint{4.261914in}{0.620274in}}%
\pgfpathlineto{\pgfqpoint{4.264011in}{0.626244in}}%
\pgfpathlineto{\pgfqpoint{4.266108in}{0.627101in}}%
\pgfpathlineto{\pgfqpoint{4.268206in}{0.621552in}}%
\pgfpathlineto{\pgfqpoint{4.270303in}{0.610568in}}%
\pgfpathlineto{\pgfqpoint{4.272400in}{0.592306in}}%
\pgfpathlineto{\pgfqpoint{4.274498in}{0.554923in}}%
\pgfpathlineto{\pgfqpoint{4.276595in}{0.534102in}}%
\pgfpathlineto{\pgfqpoint{4.278692in}{0.578232in}}%
\pgfpathlineto{\pgfqpoint{4.280790in}{0.593332in}}%
\pgfpathlineto{\pgfqpoint{4.282887in}{0.599472in}}%
\pgfpathlineto{\pgfqpoint{4.284985in}{0.599908in}}%
\pgfpathlineto{\pgfqpoint{4.287082in}{0.594887in}}%
\pgfpathlineto{\pgfqpoint{4.289179in}{0.582152in}}%
\pgfpathlineto{\pgfqpoint{4.291277in}{0.550040in}}%
\pgfpathlineto{\pgfqpoint{4.293374in}{0.534139in}}%
\pgfpathlineto{\pgfqpoint{4.295471in}{0.583488in}}%
\pgfpathlineto{\pgfqpoint{4.297569in}{0.603694in}}%
\pgfpathlineto{\pgfqpoint{4.299666in}{0.615032in}}%
\pgfpathlineto{\pgfqpoint{4.301763in}{0.620868in}}%
\pgfpathlineto{\pgfqpoint{4.303861in}{0.621630in}}%
\pgfpathlineto{\pgfqpoint{4.308056in}{0.609632in}}%
\pgfpathlineto{\pgfqpoint{4.316445in}{0.582819in}}%
\pgfpathlineto{\pgfqpoint{4.316445in}{0.582819in}}%
\pgfusepath{stroke}%
\end{pgfscope}%
\begin{pgfscope}%
\pgfpathrectangle{\pgfqpoint{2.554660in}{0.320625in}}{\pgfqpoint{1.845679in}{1.164375in}}%
\pgfusepath{clip}%
\pgfsetrectcap%
\pgfsetroundjoin%
\pgfsetlinewidth{1.154312pt}%
\definecolor{currentstroke}{rgb}{0.291366,0.064553,0.475462}%
\pgfsetstrokecolor{currentstroke}%
\pgfsetdash{}{0pt}%
\pgfpathmoveto{\pgfqpoint{2.638555in}{1.468525in}}%
\pgfpathlineto{\pgfqpoint{2.651139in}{1.446841in}}%
\pgfpathlineto{\pgfqpoint{2.661626in}{1.426459in}}%
\pgfpathlineto{\pgfqpoint{2.680502in}{1.386621in}}%
\pgfpathlineto{\pgfqpoint{2.701476in}{1.343293in}}%
\pgfpathlineto{\pgfqpoint{2.716157in}{1.315550in}}%
\pgfpathlineto{\pgfqpoint{2.728742in}{1.294055in}}%
\pgfpathlineto{\pgfqpoint{2.739228in}{1.277992in}}%
\pgfpathlineto{\pgfqpoint{2.749715in}{1.263706in}}%
\pgfpathlineto{\pgfqpoint{2.760202in}{1.251126in}}%
\pgfpathlineto{\pgfqpoint{2.772786in}{1.237951in}}%
\pgfpathlineto{\pgfqpoint{2.787468in}{1.224528in}}%
\pgfpathlineto{\pgfqpoint{2.804247in}{1.210775in}}%
\pgfpathlineto{\pgfqpoint{2.829415in}{1.191702in}}%
\pgfpathlineto{\pgfqpoint{2.877654in}{1.156844in}}%
\pgfpathlineto{\pgfqpoint{3.070612in}{1.019655in}}%
\pgfpathlineto{\pgfqpoint{3.515253in}{0.704447in}}%
\pgfpathlineto{\pgfqpoint{3.534129in}{0.691199in}}%
\pgfpathlineto{\pgfqpoint{3.550908in}{0.679476in}}%
\pgfpathlineto{\pgfqpoint{3.567687in}{0.668106in}}%
\pgfpathlineto{\pgfqpoint{3.578173in}{0.659320in}}%
\pgfpathlineto{\pgfqpoint{3.584465in}{0.654313in}}%
\pgfpathlineto{\pgfqpoint{3.588660in}{0.652198in}}%
\pgfpathlineto{\pgfqpoint{3.592855in}{0.651695in}}%
\pgfpathlineto{\pgfqpoint{3.597050in}{0.652691in}}%
\pgfpathlineto{\pgfqpoint{3.607536in}{0.656270in}}%
\pgfpathlineto{\pgfqpoint{3.611731in}{0.655726in}}%
\pgfpathlineto{\pgfqpoint{3.613829in}{0.654582in}}%
\pgfpathlineto{\pgfqpoint{3.615926in}{0.652666in}}%
\pgfpathlineto{\pgfqpoint{3.618023in}{0.649785in}}%
\pgfpathlineto{\pgfqpoint{3.628510in}{0.629809in}}%
\pgfpathlineto{\pgfqpoint{3.632705in}{0.625377in}}%
\pgfpathlineto{\pgfqpoint{3.638997in}{0.621046in}}%
\pgfpathlineto{\pgfqpoint{3.643192in}{0.619763in}}%
\pgfpathlineto{\pgfqpoint{3.647386in}{0.619378in}}%
\pgfpathlineto{\pgfqpoint{3.649484in}{0.618099in}}%
\pgfpathlineto{\pgfqpoint{3.659971in}{0.606128in}}%
\pgfpathlineto{\pgfqpoint{3.668360in}{0.599810in}}%
\pgfpathlineto{\pgfqpoint{3.674652in}{0.595992in}}%
\pgfpathlineto{\pgfqpoint{3.680944in}{0.592289in}}%
\pgfpathlineto{\pgfqpoint{3.685139in}{0.587889in}}%
\pgfpathlineto{\pgfqpoint{3.687236in}{0.586810in}}%
\pgfpathlineto{\pgfqpoint{3.689334in}{0.587818in}}%
\pgfpathlineto{\pgfqpoint{3.693528in}{0.594163in}}%
\pgfpathlineto{\pgfqpoint{3.697723in}{0.599739in}}%
\pgfpathlineto{\pgfqpoint{3.699820in}{0.600919in}}%
\pgfpathlineto{\pgfqpoint{3.701918in}{0.600762in}}%
\pgfpathlineto{\pgfqpoint{3.704015in}{0.599112in}}%
\pgfpathlineto{\pgfqpoint{3.706113in}{0.595764in}}%
\pgfpathlineto{\pgfqpoint{3.708210in}{0.590437in}}%
\pgfpathlineto{\pgfqpoint{3.714502in}{0.566967in}}%
\pgfpathlineto{\pgfqpoint{3.716599in}{0.571149in}}%
\pgfpathlineto{\pgfqpoint{3.720794in}{0.591932in}}%
\pgfpathlineto{\pgfqpoint{3.722891in}{0.599440in}}%
\pgfpathlineto{\pgfqpoint{3.724989in}{0.604349in}}%
\pgfpathlineto{\pgfqpoint{3.727086in}{0.606637in}}%
\pgfpathlineto{\pgfqpoint{3.729184in}{0.605939in}}%
\pgfpathlineto{\pgfqpoint{3.731281in}{0.601164in}}%
\pgfpathlineto{\pgfqpoint{3.733378in}{0.590643in}}%
\pgfpathlineto{\pgfqpoint{3.737573in}{0.550793in}}%
\pgfpathlineto{\pgfqpoint{3.741768in}{0.586900in}}%
\pgfpathlineto{\pgfqpoint{3.743865in}{0.597840in}}%
\pgfpathlineto{\pgfqpoint{3.745962in}{0.603795in}}%
\pgfpathlineto{\pgfqpoint{3.748060in}{0.606069in}}%
\pgfpathlineto{\pgfqpoint{3.750157in}{0.605046in}}%
\pgfpathlineto{\pgfqpoint{3.752254in}{0.600434in}}%
\pgfpathlineto{\pgfqpoint{3.754352in}{0.590975in}}%
\pgfpathlineto{\pgfqpoint{3.756449in}{0.572949in}}%
\pgfpathlineto{\pgfqpoint{3.758547in}{0.538428in}}%
\pgfpathlineto{\pgfqpoint{3.762741in}{0.590913in}}%
\pgfpathlineto{\pgfqpoint{3.764839in}{0.605813in}}%
\pgfpathlineto{\pgfqpoint{3.766936in}{0.614640in}}%
\pgfpathlineto{\pgfqpoint{3.769033in}{0.619034in}}%
\pgfpathlineto{\pgfqpoint{3.771131in}{0.619022in}}%
\pgfpathlineto{\pgfqpoint{3.773228in}{0.613172in}}%
\pgfpathlineto{\pgfqpoint{3.775325in}{0.600030in}}%
\pgfpathlineto{\pgfqpoint{3.777423in}{0.573167in}}%
\pgfpathlineto{\pgfqpoint{3.779520in}{0.528239in}}%
\pgfpathlineto{\pgfqpoint{3.781618in}{0.582298in}}%
\pgfpathlineto{\pgfqpoint{3.783715in}{0.602899in}}%
\pgfpathlineto{\pgfqpoint{3.785812in}{0.613571in}}%
\pgfpathlineto{\pgfqpoint{3.787910in}{0.619048in}}%
\pgfpathlineto{\pgfqpoint{3.790007in}{0.620779in}}%
\pgfpathlineto{\pgfqpoint{3.792104in}{0.619097in}}%
\pgfpathlineto{\pgfqpoint{3.794202in}{0.613542in}}%
\pgfpathlineto{\pgfqpoint{3.796299in}{0.602358in}}%
\pgfpathlineto{\pgfqpoint{3.798396in}{0.579332in}}%
\pgfpathlineto{\pgfqpoint{3.800494in}{0.506303in}}%
\pgfpathlineto{\pgfqpoint{3.802591in}{0.583656in}}%
\pgfpathlineto{\pgfqpoint{3.804689in}{0.609051in}}%
\pgfpathlineto{\pgfqpoint{3.806786in}{0.622872in}}%
\pgfpathlineto{\pgfqpoint{3.808883in}{0.630884in}}%
\pgfpathlineto{\pgfqpoint{3.810981in}{0.634501in}}%
\pgfpathlineto{\pgfqpoint{3.813078in}{0.633518in}}%
\pgfpathlineto{\pgfqpoint{3.815175in}{0.626391in}}%
\pgfpathlineto{\pgfqpoint{3.817273in}{0.611895in}}%
\pgfpathlineto{\pgfqpoint{3.819370in}{0.581682in}}%
\pgfpathlineto{\pgfqpoint{3.821467in}{0.536916in}}%
\pgfpathlineto{\pgfqpoint{3.823565in}{0.596544in}}%
\pgfpathlineto{\pgfqpoint{3.825662in}{0.615991in}}%
\pgfpathlineto{\pgfqpoint{3.827760in}{0.625837in}}%
\pgfpathlineto{\pgfqpoint{3.829857in}{0.630542in}}%
\pgfpathlineto{\pgfqpoint{3.831954in}{0.631407in}}%
\pgfpathlineto{\pgfqpoint{3.834052in}{0.628594in}}%
\pgfpathlineto{\pgfqpoint{3.836149in}{0.621297in}}%
\pgfpathlineto{\pgfqpoint{3.838246in}{0.606676in}}%
\pgfpathlineto{\pgfqpoint{3.840344in}{0.571682in}}%
\pgfpathlineto{\pgfqpoint{3.842441in}{0.562437in}}%
\pgfpathlineto{\pgfqpoint{3.844538in}{0.608121in}}%
\pgfpathlineto{\pgfqpoint{3.846636in}{0.627529in}}%
\pgfpathlineto{\pgfqpoint{3.848733in}{0.638635in}}%
\pgfpathlineto{\pgfqpoint{3.850831in}{0.644685in}}%
\pgfpathlineto{\pgfqpoint{3.852928in}{0.646303in}}%
\pgfpathlineto{\pgfqpoint{3.855025in}{0.642514in}}%
\pgfpathlineto{\pgfqpoint{3.857123in}{0.633223in}}%
\pgfpathlineto{\pgfqpoint{3.859220in}{0.618055in}}%
\pgfpathlineto{\pgfqpoint{3.861317in}{0.590770in}}%
\pgfpathlineto{\pgfqpoint{3.863415in}{0.475300in}}%
\pgfpathlineto{\pgfqpoint{3.865512in}{0.585603in}}%
\pgfpathlineto{\pgfqpoint{3.867609in}{0.606488in}}%
\pgfpathlineto{\pgfqpoint{3.869707in}{0.615168in}}%
\pgfpathlineto{\pgfqpoint{3.871804in}{0.617422in}}%
\pgfpathlineto{\pgfqpoint{3.873902in}{0.614290in}}%
\pgfpathlineto{\pgfqpoint{3.875999in}{0.604506in}}%
\pgfpathlineto{\pgfqpoint{3.878096in}{0.581446in}}%
\pgfpathlineto{\pgfqpoint{3.880194in}{0.489616in}}%
\pgfpathlineto{\pgfqpoint{3.882291in}{0.591137in}}%
\pgfpathlineto{\pgfqpoint{3.884388in}{0.616130in}}%
\pgfpathlineto{\pgfqpoint{3.886486in}{0.629670in}}%
\pgfpathlineto{\pgfqpoint{3.888583in}{0.637137in}}%
\pgfpathlineto{\pgfqpoint{3.890680in}{0.639632in}}%
\pgfpathlineto{\pgfqpoint{3.892778in}{0.636248in}}%
\pgfpathlineto{\pgfqpoint{3.894875in}{0.627555in}}%
\pgfpathlineto{\pgfqpoint{3.896973in}{0.614497in}}%
\pgfpathlineto{\pgfqpoint{3.899070in}{0.594133in}}%
\pgfpathlineto{\pgfqpoint{3.901167in}{0.553460in}}%
\pgfpathlineto{\pgfqpoint{3.903265in}{0.537039in}}%
\pgfpathlineto{\pgfqpoint{3.905362in}{0.576423in}}%
\pgfpathlineto{\pgfqpoint{3.907459in}{0.588427in}}%
\pgfpathlineto{\pgfqpoint{3.909557in}{0.590715in}}%
\pgfpathlineto{\pgfqpoint{3.911654in}{0.585207in}}%
\pgfpathlineto{\pgfqpoint{3.913751in}{0.568169in}}%
\pgfpathlineto{\pgfqpoint{3.915849in}{0.498807in}}%
\pgfpathlineto{\pgfqpoint{3.917946in}{0.563136in}}%
\pgfpathlineto{\pgfqpoint{3.920043in}{0.593695in}}%
\pgfpathlineto{\pgfqpoint{3.922141in}{0.609519in}}%
\pgfpathlineto{\pgfqpoint{3.924238in}{0.618453in}}%
\pgfpathlineto{\pgfqpoint{3.926336in}{0.622137in}}%
\pgfpathlineto{\pgfqpoint{3.928433in}{0.619970in}}%
\pgfpathlineto{\pgfqpoint{3.930530in}{0.611528in}}%
\pgfpathlineto{\pgfqpoint{3.932628in}{0.597853in}}%
\pgfpathlineto{\pgfqpoint{3.934725in}{0.574576in}}%
\pgfpathlineto{\pgfqpoint{3.936822in}{0.510744in}}%
\pgfpathlineto{\pgfqpoint{3.938920in}{0.550737in}}%
\pgfpathlineto{\pgfqpoint{3.941017in}{0.577293in}}%
\pgfpathlineto{\pgfqpoint{3.943114in}{0.587769in}}%
\pgfpathlineto{\pgfqpoint{3.945212in}{0.590969in}}%
\pgfpathlineto{\pgfqpoint{3.947309in}{0.588509in}}%
\pgfpathlineto{\pgfqpoint{3.949407in}{0.579395in}}%
\pgfpathlineto{\pgfqpoint{3.951504in}{0.557596in}}%
\pgfpathlineto{\pgfqpoint{3.953601in}{0.420463in}}%
\pgfpathlineto{\pgfqpoint{3.955699in}{0.563816in}}%
\pgfpathlineto{\pgfqpoint{3.957796in}{0.589841in}}%
\pgfpathlineto{\pgfqpoint{3.959893in}{0.603808in}}%
\pgfpathlineto{\pgfqpoint{3.961991in}{0.611546in}}%
\pgfpathlineto{\pgfqpoint{3.964088in}{0.614264in}}%
\pgfpathlineto{\pgfqpoint{3.966185in}{0.611108in}}%
\pgfpathlineto{\pgfqpoint{3.968283in}{0.601018in}}%
\pgfpathlineto{\pgfqpoint{3.970380in}{0.581655in}}%
\pgfpathlineto{\pgfqpoint{3.972478in}{0.530701in}}%
\pgfpathlineto{\pgfqpoint{3.974575in}{0.553857in}}%
\pgfpathlineto{\pgfqpoint{3.976672in}{0.586259in}}%
\pgfpathlineto{\pgfqpoint{3.978770in}{0.600552in}}%
\pgfpathlineto{\pgfqpoint{3.980867in}{0.607860in}}%
\pgfpathlineto{\pgfqpoint{3.982964in}{0.610763in}}%
\pgfpathlineto{\pgfqpoint{3.985062in}{0.609975in}}%
\pgfpathlineto{\pgfqpoint{3.987159in}{0.605251in}}%
\pgfpathlineto{\pgfqpoint{3.989256in}{0.595088in}}%
\pgfpathlineto{\pgfqpoint{3.991354in}{0.574134in}}%
\pgfpathlineto{\pgfqpoint{3.993451in}{0.485850in}}%
\pgfpathlineto{\pgfqpoint{3.995549in}{0.570268in}}%
\pgfpathlineto{\pgfqpoint{3.997646in}{0.598352in}}%
\pgfpathlineto{\pgfqpoint{3.999743in}{0.613096in}}%
\pgfpathlineto{\pgfqpoint{4.001841in}{0.621605in}}%
\pgfpathlineto{\pgfqpoint{4.003938in}{0.625542in}}%
\pgfpathlineto{\pgfqpoint{4.006035in}{0.624782in}}%
\pgfpathlineto{\pgfqpoint{4.008133in}{0.617663in}}%
\pgfpathlineto{\pgfqpoint{4.010230in}{0.602178in}}%
\pgfpathlineto{\pgfqpoint{4.012327in}{0.566306in}}%
\pgfpathlineto{\pgfqpoint{4.014425in}{0.552149in}}%
\pgfpathlineto{\pgfqpoint{4.016522in}{0.597229in}}%
\pgfpathlineto{\pgfqpoint{4.018620in}{0.614902in}}%
\pgfpathlineto{\pgfqpoint{4.020717in}{0.624334in}}%
\pgfpathlineto{\pgfqpoint{4.022814in}{0.629118in}}%
\pgfpathlineto{\pgfqpoint{4.024912in}{0.630395in}}%
\pgfpathlineto{\pgfqpoint{4.027009in}{0.628374in}}%
\pgfpathlineto{\pgfqpoint{4.029106in}{0.622524in}}%
\pgfpathlineto{\pgfqpoint{4.031204in}{0.610993in}}%
\pgfpathlineto{\pgfqpoint{4.033301in}{0.587148in}}%
\pgfpathlineto{\pgfqpoint{4.035398in}{0.470385in}}%
\pgfpathlineto{\pgfqpoint{4.037496in}{0.593084in}}%
\pgfpathlineto{\pgfqpoint{4.039593in}{0.618263in}}%
\pgfpathlineto{\pgfqpoint{4.041691in}{0.631979in}}%
\pgfpathlineto{\pgfqpoint{4.043788in}{0.639958in}}%
\pgfpathlineto{\pgfqpoint{4.045885in}{0.643611in}}%
\pgfpathlineto{\pgfqpoint{4.047983in}{0.642770in}}%
\pgfpathlineto{\pgfqpoint{4.050080in}{0.636054in}}%
\pgfpathlineto{\pgfqpoint{4.052177in}{0.623268in}}%
\pgfpathlineto{\pgfqpoint{4.054275in}{0.599847in}}%
\pgfpathlineto{\pgfqpoint{4.056372in}{0.520839in}}%
\pgfpathlineto{\pgfqpoint{4.058469in}{0.585856in}}%
\pgfpathlineto{\pgfqpoint{4.060567in}{0.611360in}}%
\pgfpathlineto{\pgfqpoint{4.062664in}{0.622783in}}%
\pgfpathlineto{\pgfqpoint{4.064762in}{0.627796in}}%
\pgfpathlineto{\pgfqpoint{4.066859in}{0.628247in}}%
\pgfpathlineto{\pgfqpoint{4.068956in}{0.624267in}}%
\pgfpathlineto{\pgfqpoint{4.071054in}{0.614464in}}%
\pgfpathlineto{\pgfqpoint{4.073151in}{0.593251in}}%
\pgfpathlineto{\pgfqpoint{4.075248in}{0.486720in}}%
\pgfpathlineto{\pgfqpoint{4.077346in}{0.592882in}}%
\pgfpathlineto{\pgfqpoint{4.079443in}{0.620181in}}%
\pgfpathlineto{\pgfqpoint{4.081540in}{0.634606in}}%
\pgfpathlineto{\pgfqpoint{4.083638in}{0.642778in}}%
\pgfpathlineto{\pgfqpoint{4.085735in}{0.646181in}}%
\pgfpathlineto{\pgfqpoint{4.087832in}{0.644402in}}%
\pgfpathlineto{\pgfqpoint{4.089930in}{0.636764in}}%
\pgfpathlineto{\pgfqpoint{4.092027in}{0.624922in}}%
\pgfpathlineto{\pgfqpoint{4.094125in}{0.606669in}}%
\pgfpathlineto{\pgfqpoint{4.096222in}{0.573087in}}%
\pgfpathlineto{\pgfqpoint{4.098319in}{0.521138in}}%
\pgfpathlineto{\pgfqpoint{4.100417in}{0.580104in}}%
\pgfpathlineto{\pgfqpoint{4.102514in}{0.595374in}}%
\pgfpathlineto{\pgfqpoint{4.104611in}{0.599762in}}%
\pgfpathlineto{\pgfqpoint{4.106709in}{0.596701in}}%
\pgfpathlineto{\pgfqpoint{4.108806in}{0.584536in}}%
\pgfpathlineto{\pgfqpoint{4.110903in}{0.549170in}}%
\pgfpathlineto{\pgfqpoint{4.113001in}{0.549289in}}%
\pgfpathlineto{\pgfqpoint{4.115098in}{0.593127in}}%
\pgfpathlineto{\pgfqpoint{4.117196in}{0.612565in}}%
\pgfpathlineto{\pgfqpoint{4.119293in}{0.623582in}}%
\pgfpathlineto{\pgfqpoint{4.121390in}{0.629078in}}%
\pgfpathlineto{\pgfqpoint{4.123488in}{0.629210in}}%
\pgfpathlineto{\pgfqpoint{4.125585in}{0.622822in}}%
\pgfpathlineto{\pgfqpoint{4.127682in}{0.612089in}}%
\pgfpathlineto{\pgfqpoint{4.129780in}{0.595773in}}%
\pgfpathlineto{\pgfqpoint{4.131877in}{0.567837in}}%
\pgfpathlineto{\pgfqpoint{4.133974in}{0.438143in}}%
\pgfpathlineto{\pgfqpoint{4.136072in}{0.556856in}}%
\pgfpathlineto{\pgfqpoint{4.138169in}{0.576602in}}%
\pgfpathlineto{\pgfqpoint{4.140267in}{0.583025in}}%
\pgfpathlineto{\pgfqpoint{4.142364in}{0.581689in}}%
\pgfpathlineto{\pgfqpoint{4.144461in}{0.572060in}}%
\pgfpathlineto{\pgfqpoint{4.146559in}{0.545444in}}%
\pgfpathlineto{\pgfqpoint{4.148656in}{0.508574in}}%
\pgfpathlineto{\pgfqpoint{4.150753in}{0.570405in}}%
\pgfpathlineto{\pgfqpoint{4.152851in}{0.592843in}}%
\pgfpathlineto{\pgfqpoint{4.154948in}{0.605364in}}%
\pgfpathlineto{\pgfqpoint{4.157045in}{0.612028in}}%
\pgfpathlineto{\pgfqpoint{4.159143in}{0.613464in}}%
\pgfpathlineto{\pgfqpoint{4.161240in}{0.608347in}}%
\pgfpathlineto{\pgfqpoint{4.163338in}{0.597007in}}%
\pgfpathlineto{\pgfqpoint{4.165435in}{0.576701in}}%
\pgfpathlineto{\pgfqpoint{4.167532in}{0.525421in}}%
\pgfpathlineto{\pgfqpoint{4.169630in}{0.546459in}}%
\pgfpathlineto{\pgfqpoint{4.171727in}{0.578234in}}%
\pgfpathlineto{\pgfqpoint{4.173824in}{0.591596in}}%
\pgfpathlineto{\pgfqpoint{4.175922in}{0.597747in}}%
\pgfpathlineto{\pgfqpoint{4.178019in}{0.599141in}}%
\pgfpathlineto{\pgfqpoint{4.180116in}{0.596205in}}%
\pgfpathlineto{\pgfqpoint{4.182214in}{0.587965in}}%
\pgfpathlineto{\pgfqpoint{4.184311in}{0.570416in}}%
\pgfpathlineto{\pgfqpoint{4.186409in}{0.516956in}}%
\pgfpathlineto{\pgfqpoint{4.190603in}{0.586085in}}%
\pgfpathlineto{\pgfqpoint{4.192701in}{0.602500in}}%
\pgfpathlineto{\pgfqpoint{4.194798in}{0.611887in}}%
\pgfpathlineto{\pgfqpoint{4.196895in}{0.616363in}}%
\pgfpathlineto{\pgfqpoint{4.198993in}{0.615920in}}%
\pgfpathlineto{\pgfqpoint{4.201090in}{0.608872in}}%
\pgfpathlineto{\pgfqpoint{4.203187in}{0.593144in}}%
\pgfpathlineto{\pgfqpoint{4.205285in}{0.555550in}}%
\pgfpathlineto{\pgfqpoint{4.207382in}{0.548265in}}%
\pgfpathlineto{\pgfqpoint{4.209480in}{0.590779in}}%
\pgfpathlineto{\pgfqpoint{4.211577in}{0.608096in}}%
\pgfpathlineto{\pgfqpoint{4.213674in}{0.617462in}}%
\pgfpathlineto{\pgfqpoint{4.215772in}{0.622293in}}%
\pgfpathlineto{\pgfqpoint{4.217869in}{0.623699in}}%
\pgfpathlineto{\pgfqpoint{4.219966in}{0.621902in}}%
\pgfpathlineto{\pgfqpoint{4.222064in}{0.616435in}}%
\pgfpathlineto{\pgfqpoint{4.224161in}{0.605658in}}%
\pgfpathlineto{\pgfqpoint{4.226258in}{0.583991in}}%
\pgfpathlineto{\pgfqpoint{4.228356in}{0.486761in}}%
\pgfpathlineto{\pgfqpoint{4.230453in}{0.581351in}}%
\pgfpathlineto{\pgfqpoint{4.232551in}{0.608892in}}%
\pgfpathlineto{\pgfqpoint{4.234648in}{0.623506in}}%
\pgfpathlineto{\pgfqpoint{4.236745in}{0.632080in}}%
\pgfpathlineto{\pgfqpoint{4.238843in}{0.636300in}}%
\pgfpathlineto{\pgfqpoint{4.240940in}{0.636198in}}%
\pgfpathlineto{\pgfqpoint{4.243037in}{0.630371in}}%
\pgfpathlineto{\pgfqpoint{4.245135in}{0.617462in}}%
\pgfpathlineto{\pgfqpoint{4.247232in}{0.591272in}}%
\pgfpathlineto{\pgfqpoint{4.249329in}{0.492719in}}%
\pgfpathlineto{\pgfqpoint{4.251427in}{0.594653in}}%
\pgfpathlineto{\pgfqpoint{4.253524in}{0.616901in}}%
\pgfpathlineto{\pgfqpoint{4.255621in}{0.627990in}}%
\pgfpathlineto{\pgfqpoint{4.257719in}{0.633565in}}%
\pgfpathlineto{\pgfqpoint{4.259816in}{0.635232in}}%
\pgfpathlineto{\pgfqpoint{4.261914in}{0.633339in}}%
\pgfpathlineto{\pgfqpoint{4.264011in}{0.627356in}}%
\pgfpathlineto{\pgfqpoint{4.266108in}{0.615247in}}%
\pgfpathlineto{\pgfqpoint{4.268206in}{0.589184in}}%
\pgfpathlineto{\pgfqpoint{4.270303in}{0.524946in}}%
\pgfpathlineto{\pgfqpoint{4.272400in}{0.603089in}}%
\pgfpathlineto{\pgfqpoint{4.274498in}{0.626441in}}%
\pgfpathlineto{\pgfqpoint{4.276595in}{0.639367in}}%
\pgfpathlineto{\pgfqpoint{4.278692in}{0.646734in}}%
\pgfpathlineto{\pgfqpoint{4.280790in}{0.649687in}}%
\pgfpathlineto{\pgfqpoint{4.282887in}{0.647765in}}%
\pgfpathlineto{\pgfqpoint{4.284985in}{0.640548in}}%
\pgfpathlineto{\pgfqpoint{4.287082in}{0.630589in}}%
\pgfpathlineto{\pgfqpoint{4.289179in}{0.617316in}}%
\pgfpathlineto{\pgfqpoint{4.291277in}{0.599158in}}%
\pgfpathlineto{\pgfqpoint{4.293374in}{0.571928in}}%
\pgfpathlineto{\pgfqpoint{4.295471in}{0.513825in}}%
\pgfpathlineto{\pgfqpoint{4.297569in}{0.521931in}}%
\pgfpathlineto{\pgfqpoint{4.299666in}{0.538867in}}%
\pgfpathlineto{\pgfqpoint{4.303861in}{0.505211in}}%
\pgfpathlineto{\pgfqpoint{4.305958in}{0.563338in}}%
\pgfpathlineto{\pgfqpoint{4.308056in}{0.587956in}}%
\pgfpathlineto{\pgfqpoint{4.310153in}{0.602642in}}%
\pgfpathlineto{\pgfqpoint{4.312250in}{0.611350in}}%
\pgfpathlineto{\pgfqpoint{4.314348in}{0.614934in}}%
\pgfpathlineto{\pgfqpoint{4.316445in}{0.612375in}}%
\pgfpathlineto{\pgfqpoint{4.316445in}{0.612375in}}%
\pgfusepath{stroke}%
\end{pgfscope}%
\begin{pgfscope}%
\pgfpathrectangle{\pgfqpoint{2.554660in}{0.320625in}}{\pgfqpoint{1.845679in}{1.164375in}}%
\pgfusepath{clip}%
\pgfsetrectcap%
\pgfsetroundjoin%
\pgfsetlinewidth{1.154312pt}%
\definecolor{currentstroke}{rgb}{0.432967,0.117855,0.506160}%
\pgfsetstrokecolor{currentstroke}%
\pgfsetdash{}{0pt}%
\pgfpathmoveto{\pgfqpoint{2.638555in}{1.468525in}}%
\pgfpathlineto{\pgfqpoint{2.651139in}{1.446491in}}%
\pgfpathlineto{\pgfqpoint{2.661626in}{1.425563in}}%
\pgfpathlineto{\pgfqpoint{2.678405in}{1.388854in}}%
\pgfpathlineto{\pgfqpoint{2.701476in}{1.338683in}}%
\pgfpathlineto{\pgfqpoint{2.716157in}{1.309203in}}%
\pgfpathlineto{\pgfqpoint{2.730839in}{1.282310in}}%
\pgfpathlineto{\pgfqpoint{2.743423in}{1.261542in}}%
\pgfpathlineto{\pgfqpoint{2.756007in}{1.242935in}}%
\pgfpathlineto{\pgfqpoint{2.768591in}{1.226335in}}%
\pgfpathlineto{\pgfqpoint{2.783273in}{1.209047in}}%
\pgfpathlineto{\pgfqpoint{2.800052in}{1.191228in}}%
\pgfpathlineto{\pgfqpoint{2.821026in}{1.170670in}}%
\pgfpathlineto{\pgfqpoint{2.854583in}{1.139567in}}%
\pgfpathlineto{\pgfqpoint{2.923796in}{1.077316in}}%
\pgfpathlineto{\pgfqpoint{3.265666in}{0.772378in}}%
\pgfpathlineto{\pgfqpoint{3.341171in}{0.705295in}}%
\pgfpathlineto{\pgfqpoint{3.372632in}{0.676453in}}%
\pgfpathlineto{\pgfqpoint{3.385216in}{0.666701in}}%
\pgfpathlineto{\pgfqpoint{3.399898in}{0.655629in}}%
\pgfpathlineto{\pgfqpoint{3.410384in}{0.645696in}}%
\pgfpathlineto{\pgfqpoint{3.422969in}{0.632113in}}%
\pgfpathlineto{\pgfqpoint{3.443942in}{0.609097in}}%
\pgfpathlineto{\pgfqpoint{3.450234in}{0.603572in}}%
\pgfpathlineto{\pgfqpoint{3.456526in}{0.599465in}}%
\pgfpathlineto{\pgfqpoint{3.462818in}{0.597105in}}%
\pgfpathlineto{\pgfqpoint{3.469111in}{0.596363in}}%
\pgfpathlineto{\pgfqpoint{3.483792in}{0.595798in}}%
\pgfpathlineto{\pgfqpoint{3.487987in}{0.594092in}}%
\pgfpathlineto{\pgfqpoint{3.492182in}{0.590682in}}%
\pgfpathlineto{\pgfqpoint{3.496376in}{0.585917in}}%
\pgfpathlineto{\pgfqpoint{3.498474in}{0.584380in}}%
\pgfpathlineto{\pgfqpoint{3.500571in}{0.584172in}}%
\pgfpathlineto{\pgfqpoint{3.502668in}{0.585210in}}%
\pgfpathlineto{\pgfqpoint{3.508960in}{0.589553in}}%
\pgfpathlineto{\pgfqpoint{3.511058in}{0.589648in}}%
\pgfpathlineto{\pgfqpoint{3.513155in}{0.588563in}}%
\pgfpathlineto{\pgfqpoint{3.515253in}{0.586096in}}%
\pgfpathlineto{\pgfqpoint{3.517350in}{0.582041in}}%
\pgfpathlineto{\pgfqpoint{3.521545in}{0.568960in}}%
\pgfpathlineto{\pgfqpoint{3.523642in}{0.562194in}}%
\pgfpathlineto{\pgfqpoint{3.525739in}{0.561092in}}%
\pgfpathlineto{\pgfqpoint{3.527837in}{0.567507in}}%
\pgfpathlineto{\pgfqpoint{3.532031in}{0.583831in}}%
\pgfpathlineto{\pgfqpoint{3.534129in}{0.589176in}}%
\pgfpathlineto{\pgfqpoint{3.536226in}{0.592055in}}%
\pgfpathlineto{\pgfqpoint{3.538324in}{0.592099in}}%
\pgfpathlineto{\pgfqpoint{3.540421in}{0.588375in}}%
\pgfpathlineto{\pgfqpoint{3.542518in}{0.576483in}}%
\pgfpathlineto{\pgfqpoint{3.544616in}{0.550056in}}%
\pgfpathlineto{\pgfqpoint{3.546713in}{0.561688in}}%
\pgfpathlineto{\pgfqpoint{3.548810in}{0.590125in}}%
\pgfpathlineto{\pgfqpoint{3.550908in}{0.606791in}}%
\pgfpathlineto{\pgfqpoint{3.553005in}{0.617388in}}%
\pgfpathlineto{\pgfqpoint{3.555102in}{0.624255in}}%
\pgfpathlineto{\pgfqpoint{3.557200in}{0.628411in}}%
\pgfpathlineto{\pgfqpoint{3.559297in}{0.630312in}}%
\pgfpathlineto{\pgfqpoint{3.561395in}{0.630104in}}%
\pgfpathlineto{\pgfqpoint{3.563492in}{0.627670in}}%
\pgfpathlineto{\pgfqpoint{3.565589in}{0.622568in}}%
\pgfpathlineto{\pgfqpoint{3.567687in}{0.613720in}}%
\pgfpathlineto{\pgfqpoint{3.569784in}{0.598286in}}%
\pgfpathlineto{\pgfqpoint{3.571881in}{0.564914in}}%
\pgfpathlineto{\pgfqpoint{3.573979in}{0.551913in}}%
\pgfpathlineto{\pgfqpoint{3.576076in}{0.596664in}}%
\pgfpathlineto{\pgfqpoint{3.578173in}{0.616453in}}%
\pgfpathlineto{\pgfqpoint{3.580271in}{0.628326in}}%
\pgfpathlineto{\pgfqpoint{3.582368in}{0.635768in}}%
\pgfpathlineto{\pgfqpoint{3.584465in}{0.639861in}}%
\pgfpathlineto{\pgfqpoint{3.586563in}{0.640730in}}%
\pgfpathlineto{\pgfqpoint{3.588660in}{0.637656in}}%
\pgfpathlineto{\pgfqpoint{3.590758in}{0.629958in}}%
\pgfpathlineto{\pgfqpoint{3.592855in}{0.618446in}}%
\pgfpathlineto{\pgfqpoint{3.594952in}{0.601290in}}%
\pgfpathlineto{\pgfqpoint{3.597050in}{0.572466in}}%
\pgfpathlineto{\pgfqpoint{3.599147in}{0.521842in}}%
\pgfpathlineto{\pgfqpoint{3.601244in}{0.566000in}}%
\pgfpathlineto{\pgfqpoint{3.603342in}{0.582455in}}%
\pgfpathlineto{\pgfqpoint{3.605439in}{0.586975in}}%
\pgfpathlineto{\pgfqpoint{3.607536in}{0.583322in}}%
\pgfpathlineto{\pgfqpoint{3.609634in}{0.569207in}}%
\pgfpathlineto{\pgfqpoint{3.611731in}{0.523929in}}%
\pgfpathlineto{\pgfqpoint{3.618023in}{0.605336in}}%
\pgfpathlineto{\pgfqpoint{3.620121in}{0.615267in}}%
\pgfpathlineto{\pgfqpoint{3.622218in}{0.619893in}}%
\pgfpathlineto{\pgfqpoint{3.624315in}{0.618983in}}%
\pgfpathlineto{\pgfqpoint{3.626413in}{0.611766in}}%
\pgfpathlineto{\pgfqpoint{3.628510in}{0.600628in}}%
\pgfpathlineto{\pgfqpoint{3.630607in}{0.584015in}}%
\pgfpathlineto{\pgfqpoint{3.632705in}{0.556214in}}%
\pgfpathlineto{\pgfqpoint{3.634802in}{0.489822in}}%
\pgfpathlineto{\pgfqpoint{3.636900in}{0.544186in}}%
\pgfpathlineto{\pgfqpoint{3.638997in}{0.562789in}}%
\pgfpathlineto{\pgfqpoint{3.641094in}{0.568239in}}%
\pgfpathlineto{\pgfqpoint{3.643192in}{0.565423in}}%
\pgfpathlineto{\pgfqpoint{3.645289in}{0.552705in}}%
\pgfpathlineto{\pgfqpoint{3.647386in}{0.513292in}}%
\pgfpathlineto{\pgfqpoint{3.649484in}{0.527652in}}%
\pgfpathlineto{\pgfqpoint{3.651581in}{0.566582in}}%
\pgfpathlineto{\pgfqpoint{3.653678in}{0.585070in}}%
\pgfpathlineto{\pgfqpoint{3.655776in}{0.595646in}}%
\pgfpathlineto{\pgfqpoint{3.657873in}{0.600811in}}%
\pgfpathlineto{\pgfqpoint{3.659971in}{0.600549in}}%
\pgfpathlineto{\pgfqpoint{3.662068in}{0.593227in}}%
\pgfpathlineto{\pgfqpoint{3.664165in}{0.577885in}}%
\pgfpathlineto{\pgfqpoint{3.666263in}{0.543927in}}%
\pgfpathlineto{\pgfqpoint{3.668360in}{0.521786in}}%
\pgfpathlineto{\pgfqpoint{3.670457in}{0.569602in}}%
\pgfpathlineto{\pgfqpoint{3.672555in}{0.587620in}}%
\pgfpathlineto{\pgfqpoint{3.674652in}{0.597059in}}%
\pgfpathlineto{\pgfqpoint{3.676749in}{0.601709in}}%
\pgfpathlineto{\pgfqpoint{3.678847in}{0.602741in}}%
\pgfpathlineto{\pgfqpoint{3.680944in}{0.600331in}}%
\pgfpathlineto{\pgfqpoint{3.683042in}{0.593825in}}%
\pgfpathlineto{\pgfqpoint{3.685139in}{0.580981in}}%
\pgfpathlineto{\pgfqpoint{3.687236in}{0.552938in}}%
\pgfpathlineto{\pgfqpoint{3.689334in}{0.507139in}}%
\pgfpathlineto{\pgfqpoint{3.691431in}{0.572222in}}%
\pgfpathlineto{\pgfqpoint{3.693528in}{0.594435in}}%
\pgfpathlineto{\pgfqpoint{3.695626in}{0.606907in}}%
\pgfpathlineto{\pgfqpoint{3.697723in}{0.614028in}}%
\pgfpathlineto{\pgfqpoint{3.699820in}{0.616840in}}%
\pgfpathlineto{\pgfqpoint{3.701918in}{0.614865in}}%
\pgfpathlineto{\pgfqpoint{3.704015in}{0.605589in}}%
\pgfpathlineto{\pgfqpoint{3.706113in}{0.582668in}}%
\pgfpathlineto{\pgfqpoint{3.708210in}{0.473059in}}%
\pgfpathlineto{\pgfqpoint{3.710307in}{0.587263in}}%
\pgfpathlineto{\pgfqpoint{3.712405in}{0.611978in}}%
\pgfpathlineto{\pgfqpoint{3.714502in}{0.625229in}}%
\pgfpathlineto{\pgfqpoint{3.716599in}{0.633103in}}%
\pgfpathlineto{\pgfqpoint{3.718697in}{0.637469in}}%
\pgfpathlineto{\pgfqpoint{3.720794in}{0.639039in}}%
\pgfpathlineto{\pgfqpoint{3.722891in}{0.637982in}}%
\pgfpathlineto{\pgfqpoint{3.724989in}{0.634033in}}%
\pgfpathlineto{\pgfqpoint{3.727086in}{0.626283in}}%
\pgfpathlineto{\pgfqpoint{3.729184in}{0.612214in}}%
\pgfpathlineto{\pgfqpoint{3.733378in}{0.547796in}}%
\pgfpathlineto{\pgfqpoint{3.735476in}{0.605449in}}%
\pgfpathlineto{\pgfqpoint{3.737573in}{0.626810in}}%
\pgfpathlineto{\pgfqpoint{3.739670in}{0.639122in}}%
\pgfpathlineto{\pgfqpoint{3.741768in}{0.646498in}}%
\pgfpathlineto{\pgfqpoint{3.743865in}{0.650064in}}%
\pgfpathlineto{\pgfqpoint{3.745962in}{0.649735in}}%
\pgfpathlineto{\pgfqpoint{3.748060in}{0.644359in}}%
\pgfpathlineto{\pgfqpoint{3.750157in}{0.635116in}}%
\pgfpathlineto{\pgfqpoint{3.752254in}{0.621925in}}%
\pgfpathlineto{\pgfqpoint{3.754352in}{0.602360in}}%
\pgfpathlineto{\pgfqpoint{3.756449in}{0.567479in}}%
\pgfpathlineto{\pgfqpoint{3.758547in}{0.514921in}}%
\pgfpathlineto{\pgfqpoint{3.760644in}{0.570706in}}%
\pgfpathlineto{\pgfqpoint{3.762741in}{0.583182in}}%
\pgfpathlineto{\pgfqpoint{3.764839in}{0.583203in}}%
\pgfpathlineto{\pgfqpoint{3.766936in}{0.571545in}}%
\pgfpathlineto{\pgfqpoint{3.769033in}{0.528639in}}%
\pgfpathlineto{\pgfqpoint{3.771131in}{0.552801in}}%
\pgfpathlineto{\pgfqpoint{3.773228in}{0.590603in}}%
\pgfpathlineto{\pgfqpoint{3.775325in}{0.609036in}}%
\pgfpathlineto{\pgfqpoint{3.777423in}{0.619687in}}%
\pgfpathlineto{\pgfqpoint{3.779520in}{0.624909in}}%
\pgfpathlineto{\pgfqpoint{3.781618in}{0.624608in}}%
\pgfpathlineto{\pgfqpoint{3.783715in}{0.617855in}}%
\pgfpathlineto{\pgfqpoint{3.785812in}{0.607479in}}%
\pgfpathlineto{\pgfqpoint{3.787910in}{0.592618in}}%
\pgfpathlineto{\pgfqpoint{3.790007in}{0.569783in}}%
\pgfpathlineto{\pgfqpoint{3.792104in}{0.522092in}}%
\pgfpathlineto{\pgfqpoint{3.794202in}{0.521789in}}%
\pgfpathlineto{\pgfqpoint{3.796299in}{0.552758in}}%
\pgfpathlineto{\pgfqpoint{3.798396in}{0.560469in}}%
\pgfpathlineto{\pgfqpoint{3.800494in}{0.557104in}}%
\pgfpathlineto{\pgfqpoint{3.802591in}{0.539870in}}%
\pgfpathlineto{\pgfqpoint{3.804689in}{0.432076in}}%
\pgfpathlineto{\pgfqpoint{3.806786in}{0.545276in}}%
\pgfpathlineto{\pgfqpoint{3.808883in}{0.573854in}}%
\pgfpathlineto{\pgfqpoint{3.810981in}{0.589272in}}%
\pgfpathlineto{\pgfqpoint{3.813078in}{0.598064in}}%
\pgfpathlineto{\pgfqpoint{3.815175in}{0.601609in}}%
\pgfpathlineto{\pgfqpoint{3.817273in}{0.599132in}}%
\pgfpathlineto{\pgfqpoint{3.819370in}{0.589560in}}%
\pgfpathlineto{\pgfqpoint{3.821467in}{0.571391in}}%
\pgfpathlineto{\pgfqpoint{3.823565in}{0.527509in}}%
\pgfpathlineto{\pgfqpoint{3.825662in}{0.534522in}}%
\pgfpathlineto{\pgfqpoint{3.827760in}{0.570949in}}%
\pgfpathlineto{\pgfqpoint{3.829857in}{0.586218in}}%
\pgfpathlineto{\pgfqpoint{3.831954in}{0.594004in}}%
\pgfpathlineto{\pgfqpoint{3.834052in}{0.597227in}}%
\pgfpathlineto{\pgfqpoint{3.836149in}{0.596713in}}%
\pgfpathlineto{\pgfqpoint{3.838246in}{0.592287in}}%
\pgfpathlineto{\pgfqpoint{3.840344in}{0.582564in}}%
\pgfpathlineto{\pgfqpoint{3.842441in}{0.562636in}}%
\pgfpathlineto{\pgfqpoint{3.844538in}{0.490386in}}%
\pgfpathlineto{\pgfqpoint{3.846636in}{0.554330in}}%
\pgfpathlineto{\pgfqpoint{3.848733in}{0.583881in}}%
\pgfpathlineto{\pgfqpoint{3.850831in}{0.599139in}}%
\pgfpathlineto{\pgfqpoint{3.852928in}{0.607968in}}%
\pgfpathlineto{\pgfqpoint{3.855025in}{0.612190in}}%
\pgfpathlineto{\pgfqpoint{3.857123in}{0.611776in}}%
\pgfpathlineto{\pgfqpoint{3.859220in}{0.604817in}}%
\pgfpathlineto{\pgfqpoint{3.861317in}{0.586439in}}%
\pgfpathlineto{\pgfqpoint{3.863415in}{0.522215in}}%
\pgfpathlineto{\pgfqpoint{3.865512in}{0.574254in}}%
\pgfpathlineto{\pgfqpoint{3.867609in}{0.604394in}}%
\pgfpathlineto{\pgfqpoint{3.869707in}{0.619567in}}%
\pgfpathlineto{\pgfqpoint{3.871804in}{0.628565in}}%
\pgfpathlineto{\pgfqpoint{3.873902in}{0.633777in}}%
\pgfpathlineto{\pgfqpoint{3.875999in}{0.636106in}}%
\pgfpathlineto{\pgfqpoint{3.878096in}{0.635843in}}%
\pgfpathlineto{\pgfqpoint{3.880194in}{0.632860in}}%
\pgfpathlineto{\pgfqpoint{3.882291in}{0.626507in}}%
\pgfpathlineto{\pgfqpoint{3.884388in}{0.615021in}}%
\pgfpathlineto{\pgfqpoint{3.886486in}{0.592631in}}%
\pgfpathlineto{\pgfqpoint{3.888583in}{0.486245in}}%
\pgfpathlineto{\pgfqpoint{3.890680in}{0.590817in}}%
\pgfpathlineto{\pgfqpoint{3.892778in}{0.617981in}}%
\pgfpathlineto{\pgfqpoint{3.894875in}{0.632617in}}%
\pgfpathlineto{\pgfqpoint{3.896973in}{0.641483in}}%
\pgfpathlineto{\pgfqpoint{3.899070in}{0.646347in}}%
\pgfpathlineto{\pgfqpoint{3.901167in}{0.647505in}}%
\pgfpathlineto{\pgfqpoint{3.903265in}{0.644128in}}%
\pgfpathlineto{\pgfqpoint{3.905362in}{0.636015in}}%
\pgfpathlineto{\pgfqpoint{3.907459in}{0.624148in}}%
\pgfpathlineto{\pgfqpoint{3.909557in}{0.606568in}}%
\pgfpathlineto{\pgfqpoint{3.911654in}{0.576623in}}%
\pgfpathlineto{\pgfqpoint{3.913751in}{0.448907in}}%
\pgfpathlineto{\pgfqpoint{3.915849in}{0.567736in}}%
\pgfpathlineto{\pgfqpoint{3.917946in}{0.584794in}}%
\pgfpathlineto{\pgfqpoint{3.920043in}{0.588503in}}%
\pgfpathlineto{\pgfqpoint{3.922141in}{0.582899in}}%
\pgfpathlineto{\pgfqpoint{3.924238in}{0.563288in}}%
\pgfpathlineto{\pgfqpoint{3.926336in}{0.431537in}}%
\pgfpathlineto{\pgfqpoint{3.928433in}{0.572715in}}%
\pgfpathlineto{\pgfqpoint{3.930530in}{0.599517in}}%
\pgfpathlineto{\pgfqpoint{3.932628in}{0.614105in}}%
\pgfpathlineto{\pgfqpoint{3.934725in}{0.622293in}}%
\pgfpathlineto{\pgfqpoint{3.936822in}{0.625263in}}%
\pgfpathlineto{\pgfqpoint{3.938920in}{0.622056in}}%
\pgfpathlineto{\pgfqpoint{3.941017in}{0.613552in}}%
\pgfpathlineto{\pgfqpoint{3.943114in}{0.601157in}}%
\pgfpathlineto{\pgfqpoint{3.945212in}{0.582651in}}%
\pgfpathlineto{\pgfqpoint{3.947309in}{0.550041in}}%
\pgfpathlineto{\pgfqpoint{3.949407in}{0.479337in}}%
\pgfpathlineto{\pgfqpoint{3.951504in}{0.549614in}}%
\pgfpathlineto{\pgfqpoint{3.953601in}{0.564459in}}%
\pgfpathlineto{\pgfqpoint{3.955699in}{0.567049in}}%
\pgfpathlineto{\pgfqpoint{3.957796in}{0.560256in}}%
\pgfpathlineto{\pgfqpoint{3.959893in}{0.537879in}}%
\pgfpathlineto{\pgfqpoint{3.961991in}{0.474636in}}%
\pgfpathlineto{\pgfqpoint{3.964088in}{0.555700in}}%
\pgfpathlineto{\pgfqpoint{3.966185in}{0.580366in}}%
\pgfpathlineto{\pgfqpoint{3.968283in}{0.594056in}}%
\pgfpathlineto{\pgfqpoint{3.970380in}{0.601618in}}%
\pgfpathlineto{\pgfqpoint{3.972478in}{0.603960in}}%
\pgfpathlineto{\pgfqpoint{3.974575in}{0.599895in}}%
\pgfpathlineto{\pgfqpoint{3.976672in}{0.588877in}}%
\pgfpathlineto{\pgfqpoint{3.978770in}{0.567817in}}%
\pgfpathlineto{\pgfqpoint{3.980867in}{0.505158in}}%
\pgfpathlineto{\pgfqpoint{3.982964in}{0.548686in}}%
\pgfpathlineto{\pgfqpoint{3.985062in}{0.577479in}}%
\pgfpathlineto{\pgfqpoint{3.987159in}{0.590675in}}%
\pgfpathlineto{\pgfqpoint{3.989256in}{0.597347in}}%
\pgfpathlineto{\pgfqpoint{3.991354in}{0.599743in}}%
\pgfpathlineto{\pgfqpoint{3.993451in}{0.598443in}}%
\pgfpathlineto{\pgfqpoint{3.995549in}{0.593065in}}%
\pgfpathlineto{\pgfqpoint{3.997646in}{0.581813in}}%
\pgfpathlineto{\pgfqpoint{3.999743in}{0.557953in}}%
\pgfpathlineto{\pgfqpoint{4.001841in}{0.452835in}}%
\pgfpathlineto{\pgfqpoint{4.003938in}{0.565201in}}%
\pgfpathlineto{\pgfqpoint{4.006035in}{0.590175in}}%
\pgfpathlineto{\pgfqpoint{4.008133in}{0.603760in}}%
\pgfpathlineto{\pgfqpoint{4.010230in}{0.611539in}}%
\pgfpathlineto{\pgfqpoint{4.012327in}{0.614830in}}%
\pgfpathlineto{\pgfqpoint{4.014425in}{0.613277in}}%
\pgfpathlineto{\pgfqpoint{4.016522in}{0.604580in}}%
\pgfpathlineto{\pgfqpoint{4.018620in}{0.583364in}}%
\pgfpathlineto{\pgfqpoint{4.020717in}{0.480708in}}%
\pgfpathlineto{\pgfqpoint{4.022814in}{0.580776in}}%
\pgfpathlineto{\pgfqpoint{4.024912in}{0.607319in}}%
\pgfpathlineto{\pgfqpoint{4.027009in}{0.621148in}}%
\pgfpathlineto{\pgfqpoint{4.029106in}{0.629286in}}%
\pgfpathlineto{\pgfqpoint{4.031204in}{0.633787in}}%
\pgfpathlineto{\pgfqpoint{4.033301in}{0.635416in}}%
\pgfpathlineto{\pgfqpoint{4.035398in}{0.634363in}}%
\pgfpathlineto{\pgfqpoint{4.037496in}{0.630360in}}%
\pgfpathlineto{\pgfqpoint{4.039593in}{0.622466in}}%
\pgfpathlineto{\pgfqpoint{4.041691in}{0.608039in}}%
\pgfpathlineto{\pgfqpoint{4.045885in}{0.550192in}}%
\pgfpathlineto{\pgfqpoint{4.047983in}{0.603510in}}%
\pgfpathlineto{\pgfqpoint{4.050080in}{0.624257in}}%
\pgfpathlineto{\pgfqpoint{4.052177in}{0.636270in}}%
\pgfpathlineto{\pgfqpoint{4.054275in}{0.643411in}}%
\pgfpathlineto{\pgfqpoint{4.056372in}{0.646719in}}%
\pgfpathlineto{\pgfqpoint{4.058469in}{0.646023in}}%
\pgfpathlineto{\pgfqpoint{4.060567in}{0.640160in}}%
\pgfpathlineto{\pgfqpoint{4.062664in}{0.630347in}}%
\pgfpathlineto{\pgfqpoint{4.064762in}{0.615973in}}%
\pgfpathlineto{\pgfqpoint{4.066859in}{0.593485in}}%
\pgfpathlineto{\pgfqpoint{4.068956in}{0.544883in}}%
\pgfpathlineto{\pgfqpoint{4.071054in}{0.549267in}}%
\pgfpathlineto{\pgfqpoint{4.073151in}{0.580381in}}%
\pgfpathlineto{\pgfqpoint{4.075248in}{0.589523in}}%
\pgfpathlineto{\pgfqpoint{4.077346in}{0.588964in}}%
\pgfpathlineto{\pgfqpoint{4.079443in}{0.578746in}}%
\pgfpathlineto{\pgfqpoint{4.081540in}{0.547112in}}%
\pgfpathlineto{\pgfqpoint{4.083638in}{0.538056in}}%
\pgfpathlineto{\pgfqpoint{4.085735in}{0.586175in}}%
\pgfpathlineto{\pgfqpoint{4.087832in}{0.606738in}}%
\pgfpathlineto{\pgfqpoint{4.089930in}{0.618445in}}%
\pgfpathlineto{\pgfqpoint{4.092027in}{0.624515in}}%
\pgfpathlineto{\pgfqpoint{4.094125in}{0.625261in}}%
\pgfpathlineto{\pgfqpoint{4.096222in}{0.619466in}}%
\pgfpathlineto{\pgfqpoint{4.098319in}{0.609494in}}%
\pgfpathlineto{\pgfqpoint{4.100417in}{0.594878in}}%
\pgfpathlineto{\pgfqpoint{4.102514in}{0.571878in}}%
\pgfpathlineto{\pgfqpoint{4.104611in}{0.520531in}}%
\pgfpathlineto{\pgfqpoint{4.106709in}{0.531343in}}%
\pgfpathlineto{\pgfqpoint{4.108806in}{0.560679in}}%
\pgfpathlineto{\pgfqpoint{4.110903in}{0.569320in}}%
\pgfpathlineto{\pgfqpoint{4.113001in}{0.568398in}}%
\pgfpathlineto{\pgfqpoint{4.115098in}{0.557671in}}%
\pgfpathlineto{\pgfqpoint{4.117196in}{0.524169in}}%
\pgfpathlineto{\pgfqpoint{4.119293in}{0.521431in}}%
\pgfpathlineto{\pgfqpoint{4.121390in}{0.566971in}}%
\pgfpathlineto{\pgfqpoint{4.123488in}{0.587014in}}%
\pgfpathlineto{\pgfqpoint{4.125585in}{0.598450in}}%
\pgfpathlineto{\pgfqpoint{4.127682in}{0.604303in}}%
\pgfpathlineto{\pgfqpoint{4.129780in}{0.604804in}}%
\pgfpathlineto{\pgfqpoint{4.131877in}{0.598430in}}%
\pgfpathlineto{\pgfqpoint{4.133974in}{0.585039in}}%
\pgfpathlineto{\pgfqpoint{4.136072in}{0.558504in}}%
\pgfpathlineto{\pgfqpoint{4.138169in}{0.457942in}}%
\pgfpathlineto{\pgfqpoint{4.140267in}{0.561053in}}%
\pgfpathlineto{\pgfqpoint{4.142364in}{0.583017in}}%
\pgfpathlineto{\pgfqpoint{4.144461in}{0.593770in}}%
\pgfpathlineto{\pgfqpoint{4.146559in}{0.598942in}}%
\pgfpathlineto{\pgfqpoint{4.148656in}{0.600093in}}%
\pgfpathlineto{\pgfqpoint{4.150753in}{0.597477in}}%
\pgfpathlineto{\pgfqpoint{4.152851in}{0.590332in}}%
\pgfpathlineto{\pgfqpoint{4.154948in}{0.575853in}}%
\pgfpathlineto{\pgfqpoint{4.157045in}{0.541083in}}%
\pgfpathlineto{\pgfqpoint{4.159143in}{0.530869in}}%
\pgfpathlineto{\pgfqpoint{4.161240in}{0.577152in}}%
\pgfpathlineto{\pgfqpoint{4.163338in}{0.596663in}}%
\pgfpathlineto{\pgfqpoint{4.165435in}{0.607810in}}%
\pgfpathlineto{\pgfqpoint{4.167532in}{0.613869in}}%
\pgfpathlineto{\pgfqpoint{4.169630in}{0.615462in}}%
\pgfpathlineto{\pgfqpoint{4.171727in}{0.611570in}}%
\pgfpathlineto{\pgfqpoint{4.173824in}{0.599290in}}%
\pgfpathlineto{\pgfqpoint{4.178019in}{0.534407in}}%
\pgfpathlineto{\pgfqpoint{4.180116in}{0.591779in}}%
\pgfpathlineto{\pgfqpoint{4.182214in}{0.612262in}}%
\pgfpathlineto{\pgfqpoint{4.184311in}{0.623634in}}%
\pgfpathlineto{\pgfqpoint{4.186409in}{0.630244in}}%
\pgfpathlineto{\pgfqpoint{4.188506in}{0.633541in}}%
\pgfpathlineto{\pgfqpoint{4.190603in}{0.634040in}}%
\pgfpathlineto{\pgfqpoint{4.192701in}{0.631742in}}%
\pgfpathlineto{\pgfqpoint{4.194798in}{0.626126in}}%
\pgfpathlineto{\pgfqpoint{4.196895in}{0.615674in}}%
\pgfpathlineto{\pgfqpoint{4.198993in}{0.595622in}}%
\pgfpathlineto{\pgfqpoint{4.201090in}{0.529479in}}%
\pgfpathlineto{\pgfqpoint{4.203187in}{0.583836in}}%
\pgfpathlineto{\pgfqpoint{4.205285in}{0.614226in}}%
\pgfpathlineto{\pgfqpoint{4.207382in}{0.629902in}}%
\pgfpathlineto{\pgfqpoint{4.209480in}{0.639298in}}%
\pgfpathlineto{\pgfqpoint{4.211577in}{0.644480in}}%
\pgfpathlineto{\pgfqpoint{4.213674in}{0.645825in}}%
\pgfpathlineto{\pgfqpoint{4.215772in}{0.642515in}}%
\pgfpathlineto{\pgfqpoint{4.217869in}{0.634327in}}%
\pgfpathlineto{\pgfqpoint{4.219966in}{0.622170in}}%
\pgfpathlineto{\pgfqpoint{4.222064in}{0.603795in}}%
\pgfpathlineto{\pgfqpoint{4.224161in}{0.570867in}}%
\pgfpathlineto{\pgfqpoint{4.226258in}{0.507966in}}%
\pgfpathlineto{\pgfqpoint{4.228356in}{0.573447in}}%
\pgfpathlineto{\pgfqpoint{4.230453in}{0.588567in}}%
\pgfpathlineto{\pgfqpoint{4.232551in}{0.592058in}}%
\pgfpathlineto{\pgfqpoint{4.234648in}{0.587164in}}%
\pgfpathlineto{\pgfqpoint{4.236745in}{0.570420in}}%
\pgfpathlineto{\pgfqpoint{4.238843in}{0.499428in}}%
\pgfpathlineto{\pgfqpoint{4.240940in}{0.566417in}}%
\pgfpathlineto{\pgfqpoint{4.243037in}{0.596977in}}%
\pgfpathlineto{\pgfqpoint{4.245135in}{0.612886in}}%
\pgfpathlineto{\pgfqpoint{4.247232in}{0.621914in}}%
\pgfpathlineto{\pgfqpoint{4.249329in}{0.625699in}}%
\pgfpathlineto{\pgfqpoint{4.251427in}{0.623660in}}%
\pgfpathlineto{\pgfqpoint{4.253524in}{0.615676in}}%
\pgfpathlineto{\pgfqpoint{4.255621in}{0.603765in}}%
\pgfpathlineto{\pgfqpoint{4.257719in}{0.585826in}}%
\pgfpathlineto{\pgfqpoint{4.259816in}{0.554214in}}%
\pgfpathlineto{\pgfqpoint{4.261914in}{0.475555in}}%
\pgfpathlineto{\pgfqpoint{4.264011in}{0.552916in}}%
\pgfpathlineto{\pgfqpoint{4.266108in}{0.568955in}}%
\pgfpathlineto{\pgfqpoint{4.268206in}{0.572913in}}%
\pgfpathlineto{\pgfqpoint{4.270303in}{0.568482in}}%
\pgfpathlineto{\pgfqpoint{4.272400in}{0.552630in}}%
\pgfpathlineto{\pgfqpoint{4.274498in}{0.491149in}}%
\pgfpathlineto{\pgfqpoint{4.276595in}{0.544720in}}%
\pgfpathlineto{\pgfqpoint{4.278692in}{0.576688in}}%
\pgfpathlineto{\pgfqpoint{4.280790in}{0.593074in}}%
\pgfpathlineto{\pgfqpoint{4.282887in}{0.602401in}}%
\pgfpathlineto{\pgfqpoint{4.284985in}{0.606459in}}%
\pgfpathlineto{\pgfqpoint{4.287082in}{0.604782in}}%
\pgfpathlineto{\pgfqpoint{4.291277in}{0.589258in}}%
\pgfpathlineto{\pgfqpoint{4.295471in}{0.568291in}}%
\pgfpathlineto{\pgfqpoint{4.301763in}{0.535032in}}%
\pgfpathlineto{\pgfqpoint{4.303861in}{0.530765in}}%
\pgfpathlineto{\pgfqpoint{4.305958in}{0.532543in}}%
\pgfpathlineto{\pgfqpoint{4.312250in}{0.549750in}}%
\pgfpathlineto{\pgfqpoint{4.314348in}{0.551823in}}%
\pgfpathlineto{\pgfqpoint{4.316445in}{0.548867in}}%
\pgfpathlineto{\pgfqpoint{4.316445in}{0.548867in}}%
\pgfusepath{stroke}%
\end{pgfscope}%
\begin{pgfscope}%
\pgfpathrectangle{\pgfqpoint{2.554660in}{0.320625in}}{\pgfqpoint{1.845679in}{1.164375in}}%
\pgfusepath{clip}%
\pgfsetrectcap%
\pgfsetroundjoin%
\pgfsetlinewidth{1.154312pt}%
\definecolor{currentstroke}{rgb}{0.569172,0.167454,0.504105}%
\pgfsetstrokecolor{currentstroke}%
\pgfsetdash{}{0pt}%
\pgfpathmoveto{\pgfqpoint{2.638555in}{1.468525in}}%
\pgfpathlineto{\pgfqpoint{2.651139in}{1.446163in}}%
\pgfpathlineto{\pgfqpoint{2.661626in}{1.424759in}}%
\pgfpathlineto{\pgfqpoint{2.676308in}{1.391774in}}%
\pgfpathlineto{\pgfqpoint{2.703573in}{1.330233in}}%
\pgfpathlineto{\pgfqpoint{2.720352in}{1.295284in}}%
\pgfpathlineto{\pgfqpoint{2.735034in}{1.267142in}}%
\pgfpathlineto{\pgfqpoint{2.749715in}{1.241403in}}%
\pgfpathlineto{\pgfqpoint{2.764397in}{1.218030in}}%
\pgfpathlineto{\pgfqpoint{2.779078in}{1.196792in}}%
\pgfpathlineto{\pgfqpoint{2.795857in}{1.174617in}}%
\pgfpathlineto{\pgfqpoint{2.814733in}{1.151543in}}%
\pgfpathlineto{\pgfqpoint{2.839902in}{1.122588in}}%
\pgfpathlineto{\pgfqpoint{2.879752in}{1.078641in}}%
\pgfpathlineto{\pgfqpoint{2.965744in}{0.985894in}}%
\pgfpathlineto{\pgfqpoint{3.227914in}{0.705023in}}%
\pgfpathlineto{\pgfqpoint{3.244693in}{0.687946in}}%
\pgfpathlineto{\pgfqpoint{3.259374in}{0.670785in}}%
\pgfpathlineto{\pgfqpoint{3.274056in}{0.653789in}}%
\pgfpathlineto{\pgfqpoint{3.282445in}{0.645631in}}%
\pgfpathlineto{\pgfqpoint{3.290835in}{0.639153in}}%
\pgfpathlineto{\pgfqpoint{3.303419in}{0.630452in}}%
\pgfpathlineto{\pgfqpoint{3.307614in}{0.626407in}}%
\pgfpathlineto{\pgfqpoint{3.313906in}{0.618816in}}%
\pgfpathlineto{\pgfqpoint{3.320198in}{0.609823in}}%
\pgfpathlineto{\pgfqpoint{3.328587in}{0.595676in}}%
\pgfpathlineto{\pgfqpoint{3.347464in}{0.560779in}}%
\pgfpathlineto{\pgfqpoint{3.351658in}{0.555886in}}%
\pgfpathlineto{\pgfqpoint{3.353756in}{0.554424in}}%
\pgfpathlineto{\pgfqpoint{3.355853in}{0.553695in}}%
\pgfpathlineto{\pgfqpoint{3.357950in}{0.553698in}}%
\pgfpathlineto{\pgfqpoint{3.362145in}{0.555651in}}%
\pgfpathlineto{\pgfqpoint{3.366340in}{0.559429in}}%
\pgfpathlineto{\pgfqpoint{3.378924in}{0.572177in}}%
\pgfpathlineto{\pgfqpoint{3.383119in}{0.574692in}}%
\pgfpathlineto{\pgfqpoint{3.387313in}{0.575513in}}%
\pgfpathlineto{\pgfqpoint{3.389411in}{0.575091in}}%
\pgfpathlineto{\pgfqpoint{3.391508in}{0.573957in}}%
\pgfpathlineto{\pgfqpoint{3.397800in}{0.568100in}}%
\pgfpathlineto{\pgfqpoint{3.399898in}{0.569539in}}%
\pgfpathlineto{\pgfqpoint{3.401995in}{0.574335in}}%
\pgfpathlineto{\pgfqpoint{3.408287in}{0.593587in}}%
\pgfpathlineto{\pgfqpoint{3.412482in}{0.602420in}}%
\pgfpathlineto{\pgfqpoint{3.414579in}{0.605299in}}%
\pgfpathlineto{\pgfqpoint{3.416676in}{0.607208in}}%
\pgfpathlineto{\pgfqpoint{3.418774in}{0.608168in}}%
\pgfpathlineto{\pgfqpoint{3.420871in}{0.608168in}}%
\pgfpathlineto{\pgfqpoint{3.422969in}{0.607159in}}%
\pgfpathlineto{\pgfqpoint{3.425066in}{0.605031in}}%
\pgfpathlineto{\pgfqpoint{3.427163in}{0.601593in}}%
\pgfpathlineto{\pgfqpoint{3.429261in}{0.596506in}}%
\pgfpathlineto{\pgfqpoint{3.431358in}{0.589153in}}%
\pgfpathlineto{\pgfqpoint{3.433455in}{0.578295in}}%
\pgfpathlineto{\pgfqpoint{3.435553in}{0.560991in}}%
\pgfpathlineto{\pgfqpoint{3.437650in}{0.530521in}}%
\pgfpathlineto{\pgfqpoint{3.439747in}{0.539227in}}%
\pgfpathlineto{\pgfqpoint{3.441845in}{0.568211in}}%
\pgfpathlineto{\pgfqpoint{3.443942in}{0.585393in}}%
\pgfpathlineto{\pgfqpoint{3.446040in}{0.596948in}}%
\pgfpathlineto{\pgfqpoint{3.450234in}{0.611315in}}%
\pgfpathlineto{\pgfqpoint{3.452332in}{0.615612in}}%
\pgfpathlineto{\pgfqpoint{3.454429in}{0.618345in}}%
\pgfpathlineto{\pgfqpoint{3.456526in}{0.619534in}}%
\pgfpathlineto{\pgfqpoint{3.458624in}{0.619017in}}%
\pgfpathlineto{\pgfqpoint{3.460721in}{0.616377in}}%
\pgfpathlineto{\pgfqpoint{3.462818in}{0.610700in}}%
\pgfpathlineto{\pgfqpoint{3.464916in}{0.600801in}}%
\pgfpathlineto{\pgfqpoint{3.467013in}{0.584853in}}%
\pgfpathlineto{\pgfqpoint{3.471208in}{0.526193in}}%
\pgfpathlineto{\pgfqpoint{3.473305in}{0.562165in}}%
\pgfpathlineto{\pgfqpoint{3.475403in}{0.578949in}}%
\pgfpathlineto{\pgfqpoint{3.477500in}{0.586350in}}%
\pgfpathlineto{\pgfqpoint{3.479597in}{0.588008in}}%
\pgfpathlineto{\pgfqpoint{3.481695in}{0.584571in}}%
\pgfpathlineto{\pgfqpoint{3.483792in}{0.574694in}}%
\pgfpathlineto{\pgfqpoint{3.485889in}{0.552081in}}%
\pgfpathlineto{\pgfqpoint{3.487987in}{0.489800in}}%
\pgfpathlineto{\pgfqpoint{3.490084in}{0.560751in}}%
\pgfpathlineto{\pgfqpoint{3.492182in}{0.585627in}}%
\pgfpathlineto{\pgfqpoint{3.494279in}{0.599219in}}%
\pgfpathlineto{\pgfqpoint{3.496376in}{0.606817in}}%
\pgfpathlineto{\pgfqpoint{3.498474in}{0.609570in}}%
\pgfpathlineto{\pgfqpoint{3.500571in}{0.606718in}}%
\pgfpathlineto{\pgfqpoint{3.502668in}{0.597937in}}%
\pgfpathlineto{\pgfqpoint{3.504766in}{0.583514in}}%
\pgfpathlineto{\pgfqpoint{3.506863in}{0.558131in}}%
\pgfpathlineto{\pgfqpoint{3.508960in}{0.495764in}}%
\pgfpathlineto{\pgfqpoint{3.511058in}{0.548912in}}%
\pgfpathlineto{\pgfqpoint{3.513155in}{0.570758in}}%
\pgfpathlineto{\pgfqpoint{3.515253in}{0.580195in}}%
\pgfpathlineto{\pgfqpoint{3.517350in}{0.583254in}}%
\pgfpathlineto{\pgfqpoint{3.519447in}{0.581221in}}%
\pgfpathlineto{\pgfqpoint{3.521545in}{0.573360in}}%
\pgfpathlineto{\pgfqpoint{3.523642in}{0.555448in}}%
\pgfpathlineto{\pgfqpoint{3.525739in}{0.495994in}}%
\pgfpathlineto{\pgfqpoint{3.527837in}{0.543857in}}%
\pgfpathlineto{\pgfqpoint{3.529934in}{0.575244in}}%
\pgfpathlineto{\pgfqpoint{3.532031in}{0.591153in}}%
\pgfpathlineto{\pgfqpoint{3.534129in}{0.600209in}}%
\pgfpathlineto{\pgfqpoint{3.536226in}{0.604287in}}%
\pgfpathlineto{\pgfqpoint{3.538324in}{0.603157in}}%
\pgfpathlineto{\pgfqpoint{3.540421in}{0.594728in}}%
\pgfpathlineto{\pgfqpoint{3.542518in}{0.574989in}}%
\pgfpathlineto{\pgfqpoint{3.544616in}{0.508210in}}%
\pgfpathlineto{\pgfqpoint{3.546713in}{0.564050in}}%
\pgfpathlineto{\pgfqpoint{3.548810in}{0.592693in}}%
\pgfpathlineto{\pgfqpoint{3.550908in}{0.607074in}}%
\pgfpathlineto{\pgfqpoint{3.553005in}{0.615392in}}%
\pgfpathlineto{\pgfqpoint{3.555102in}{0.619908in}}%
\pgfpathlineto{\pgfqpoint{3.557200in}{0.621445in}}%
\pgfpathlineto{\pgfqpoint{3.559297in}{0.620190in}}%
\pgfpathlineto{\pgfqpoint{3.561395in}{0.615833in}}%
\pgfpathlineto{\pgfqpoint{3.563492in}{0.607293in}}%
\pgfpathlineto{\pgfqpoint{3.565589in}{0.591426in}}%
\pgfpathlineto{\pgfqpoint{3.567687in}{0.553516in}}%
\pgfpathlineto{\pgfqpoint{3.569784in}{0.551646in}}%
\pgfpathlineto{\pgfqpoint{3.571881in}{0.594371in}}%
\pgfpathlineto{\pgfqpoint{3.573979in}{0.613232in}}%
\pgfpathlineto{\pgfqpoint{3.576076in}{0.624311in}}%
\pgfpathlineto{\pgfqpoint{3.578173in}{0.630742in}}%
\pgfpathlineto{\pgfqpoint{3.580271in}{0.633311in}}%
\pgfpathlineto{\pgfqpoint{3.582368in}{0.631597in}}%
\pgfpathlineto{\pgfqpoint{3.584465in}{0.624433in}}%
\pgfpathlineto{\pgfqpoint{3.586563in}{0.612273in}}%
\pgfpathlineto{\pgfqpoint{3.588660in}{0.592042in}}%
\pgfpathlineto{\pgfqpoint{3.590758in}{0.547077in}}%
\pgfpathlineto{\pgfqpoint{3.592855in}{0.550844in}}%
\pgfpathlineto{\pgfqpoint{3.594952in}{0.585256in}}%
\pgfpathlineto{\pgfqpoint{3.597050in}{0.598239in}}%
\pgfpathlineto{\pgfqpoint{3.599147in}{0.603189in}}%
\pgfpathlineto{\pgfqpoint{3.601244in}{0.602590in}}%
\pgfpathlineto{\pgfqpoint{3.603342in}{0.596282in}}%
\pgfpathlineto{\pgfqpoint{3.605439in}{0.581128in}}%
\pgfpathlineto{\pgfqpoint{3.607536in}{0.538171in}}%
\pgfpathlineto{\pgfqpoint{3.609634in}{0.554601in}}%
\pgfpathlineto{\pgfqpoint{3.611731in}{0.592494in}}%
\pgfpathlineto{\pgfqpoint{3.613829in}{0.610280in}}%
\pgfpathlineto{\pgfqpoint{3.615926in}{0.620400in}}%
\pgfpathlineto{\pgfqpoint{3.618023in}{0.625336in}}%
\pgfpathlineto{\pgfqpoint{3.620121in}{0.625162in}}%
\pgfpathlineto{\pgfqpoint{3.622218in}{0.618646in}}%
\pgfpathlineto{\pgfqpoint{3.624315in}{0.607040in}}%
\pgfpathlineto{\pgfqpoint{3.626413in}{0.587882in}}%
\pgfpathlineto{\pgfqpoint{3.628510in}{0.547538in}}%
\pgfpathlineto{\pgfqpoint{3.630607in}{0.537067in}}%
\pgfpathlineto{\pgfqpoint{3.632705in}{0.576621in}}%
\pgfpathlineto{\pgfqpoint{3.634802in}{0.590749in}}%
\pgfpathlineto{\pgfqpoint{3.636900in}{0.596321in}}%
\pgfpathlineto{\pgfqpoint{3.638997in}{0.596248in}}%
\pgfpathlineto{\pgfqpoint{3.641094in}{0.590580in}}%
\pgfpathlineto{\pgfqpoint{3.643192in}{0.576633in}}%
\pgfpathlineto{\pgfqpoint{3.645289in}{0.539475in}}%
\pgfpathlineto{\pgfqpoint{3.647386in}{0.540834in}}%
\pgfpathlineto{\pgfqpoint{3.649484in}{0.583367in}}%
\pgfpathlineto{\pgfqpoint{3.651581in}{0.602259in}}%
\pgfpathlineto{\pgfqpoint{3.653678in}{0.612978in}}%
\pgfpathlineto{\pgfqpoint{3.655776in}{0.618399in}}%
\pgfpathlineto{\pgfqpoint{3.657873in}{0.618785in}}%
\pgfpathlineto{\pgfqpoint{3.659971in}{0.612725in}}%
\pgfpathlineto{\pgfqpoint{3.662068in}{0.600077in}}%
\pgfpathlineto{\pgfqpoint{3.664165in}{0.576002in}}%
\pgfpathlineto{\pgfqpoint{3.666263in}{0.476636in}}%
\pgfpathlineto{\pgfqpoint{3.668360in}{0.568220in}}%
\pgfpathlineto{\pgfqpoint{3.670457in}{0.592659in}}%
\pgfpathlineto{\pgfqpoint{3.672555in}{0.604177in}}%
\pgfpathlineto{\pgfqpoint{3.674652in}{0.609675in}}%
\pgfpathlineto{\pgfqpoint{3.676749in}{0.610955in}}%
\pgfpathlineto{\pgfqpoint{3.678847in}{0.608331in}}%
\pgfpathlineto{\pgfqpoint{3.680944in}{0.601016in}}%
\pgfpathlineto{\pgfqpoint{3.683042in}{0.586007in}}%
\pgfpathlineto{\pgfqpoint{3.685139in}{0.548582in}}%
\pgfpathlineto{\pgfqpoint{3.687236in}{0.547554in}}%
\pgfpathlineto{\pgfqpoint{3.689334in}{0.590329in}}%
\pgfpathlineto{\pgfqpoint{3.691431in}{0.609143in}}%
\pgfpathlineto{\pgfqpoint{3.693528in}{0.619923in}}%
\pgfpathlineto{\pgfqpoint{3.695626in}{0.625672in}}%
\pgfpathlineto{\pgfqpoint{3.697723in}{0.626886in}}%
\pgfpathlineto{\pgfqpoint{3.699820in}{0.622410in}}%
\pgfpathlineto{\pgfqpoint{3.701918in}{0.611064in}}%
\pgfpathlineto{\pgfqpoint{3.704015in}{0.588694in}}%
\pgfpathlineto{\pgfqpoint{3.706113in}{0.507807in}}%
\pgfpathlineto{\pgfqpoint{3.708210in}{0.577895in}}%
\pgfpathlineto{\pgfqpoint{3.710307in}{0.604273in}}%
\pgfpathlineto{\pgfqpoint{3.712405in}{0.616921in}}%
\pgfpathlineto{\pgfqpoint{3.714502in}{0.623480in}}%
\pgfpathlineto{\pgfqpoint{3.716599in}{0.625980in}}%
\pgfpathlineto{\pgfqpoint{3.718697in}{0.624975in}}%
\pgfpathlineto{\pgfqpoint{3.720794in}{0.620171in}}%
\pgfpathlineto{\pgfqpoint{3.722891in}{0.610095in}}%
\pgfpathlineto{\pgfqpoint{3.724989in}{0.589655in}}%
\pgfpathlineto{\pgfqpoint{3.727086in}{0.512190in}}%
\pgfpathlineto{\pgfqpoint{3.729184in}{0.582891in}}%
\pgfpathlineto{\pgfqpoint{3.731281in}{0.611851in}}%
\pgfpathlineto{\pgfqpoint{3.733378in}{0.626932in}}%
\pgfpathlineto{\pgfqpoint{3.735476in}{0.635713in}}%
\pgfpathlineto{\pgfqpoint{3.737573in}{0.639985in}}%
\pgfpathlineto{\pgfqpoint{3.739670in}{0.639763in}}%
\pgfpathlineto{\pgfqpoint{3.741768in}{0.633699in}}%
\pgfpathlineto{\pgfqpoint{3.743865in}{0.621658in}}%
\pgfpathlineto{\pgfqpoint{3.745962in}{0.599893in}}%
\pgfpathlineto{\pgfqpoint{3.748060in}{0.538165in}}%
\pgfpathlineto{\pgfqpoint{3.750157in}{0.578173in}}%
\pgfpathlineto{\pgfqpoint{3.752254in}{0.606550in}}%
\pgfpathlineto{\pgfqpoint{3.754352in}{0.618980in}}%
\pgfpathlineto{\pgfqpoint{3.756449in}{0.624664in}}%
\pgfpathlineto{\pgfqpoint{3.758547in}{0.625749in}}%
\pgfpathlineto{\pgfqpoint{3.760644in}{0.622556in}}%
\pgfpathlineto{\pgfqpoint{3.762741in}{0.614032in}}%
\pgfpathlineto{\pgfqpoint{3.764839in}{0.595973in}}%
\pgfpathlineto{\pgfqpoint{3.766936in}{0.538923in}}%
\pgfpathlineto{\pgfqpoint{3.771131in}{0.613135in}}%
\pgfpathlineto{\pgfqpoint{3.773228in}{0.629271in}}%
\pgfpathlineto{\pgfqpoint{3.775325in}{0.638530in}}%
\pgfpathlineto{\pgfqpoint{3.777423in}{0.642950in}}%
\pgfpathlineto{\pgfqpoint{3.779520in}{0.642514in}}%
\pgfpathlineto{\pgfqpoint{3.781618in}{0.636046in}}%
\pgfpathlineto{\pgfqpoint{3.783715in}{0.624930in}}%
\pgfpathlineto{\pgfqpoint{3.785812in}{0.607328in}}%
\pgfpathlineto{\pgfqpoint{3.787910in}{0.574252in}}%
\pgfpathlineto{\pgfqpoint{3.790007in}{0.524672in}}%
\pgfpathlineto{\pgfqpoint{3.792104in}{0.583622in}}%
\pgfpathlineto{\pgfqpoint{3.794202in}{0.599981in}}%
\pgfpathlineto{\pgfqpoint{3.796299in}{0.605994in}}%
\pgfpathlineto{\pgfqpoint{3.798396in}{0.605555in}}%
\pgfpathlineto{\pgfqpoint{3.800494in}{0.598611in}}%
\pgfpathlineto{\pgfqpoint{3.802591in}{0.581047in}}%
\pgfpathlineto{\pgfqpoint{3.804689in}{0.517022in}}%
\pgfpathlineto{\pgfqpoint{3.806786in}{0.571987in}}%
\pgfpathlineto{\pgfqpoint{3.808883in}{0.603124in}}%
\pgfpathlineto{\pgfqpoint{3.810981in}{0.619008in}}%
\pgfpathlineto{\pgfqpoint{3.813078in}{0.628001in}}%
\pgfpathlineto{\pgfqpoint{3.815175in}{0.631882in}}%
\pgfpathlineto{\pgfqpoint{3.817273in}{0.630252in}}%
\pgfpathlineto{\pgfqpoint{3.819370in}{0.622467in}}%
\pgfpathlineto{\pgfqpoint{3.821467in}{0.610139in}}%
\pgfpathlineto{\pgfqpoint{3.823565in}{0.590467in}}%
\pgfpathlineto{\pgfqpoint{3.825662in}{0.550396in}}%
\pgfpathlineto{\pgfqpoint{3.827760in}{0.534936in}}%
\pgfpathlineto{\pgfqpoint{3.829857in}{0.575426in}}%
\pgfpathlineto{\pgfqpoint{3.831954in}{0.588758in}}%
\pgfpathlineto{\pgfqpoint{3.834052in}{0.592977in}}%
\pgfpathlineto{\pgfqpoint{3.836149in}{0.590742in}}%
\pgfpathlineto{\pgfqpoint{3.838246in}{0.581078in}}%
\pgfpathlineto{\pgfqpoint{3.840344in}{0.556439in}}%
\pgfpathlineto{\pgfqpoint{3.842441in}{0.497964in}}%
\pgfpathlineto{\pgfqpoint{3.844538in}{0.573700in}}%
\pgfpathlineto{\pgfqpoint{3.846636in}{0.597366in}}%
\pgfpathlineto{\pgfqpoint{3.848733in}{0.610416in}}%
\pgfpathlineto{\pgfqpoint{3.850831in}{0.617528in}}%
\pgfpathlineto{\pgfqpoint{3.852928in}{0.619577in}}%
\pgfpathlineto{\pgfqpoint{3.855025in}{0.615442in}}%
\pgfpathlineto{\pgfqpoint{3.857123in}{0.605215in}}%
\pgfpathlineto{\pgfqpoint{3.859220in}{0.587425in}}%
\pgfpathlineto{\pgfqpoint{3.861317in}{0.549389in}}%
\pgfpathlineto{\pgfqpoint{3.863415in}{0.534500in}}%
\pgfpathlineto{\pgfqpoint{3.865512in}{0.577448in}}%
\pgfpathlineto{\pgfqpoint{3.867609in}{0.593203in}}%
\pgfpathlineto{\pgfqpoint{3.869707in}{0.600487in}}%
\pgfpathlineto{\pgfqpoint{3.871804in}{0.602634in}}%
\pgfpathlineto{\pgfqpoint{3.873902in}{0.600355in}}%
\pgfpathlineto{\pgfqpoint{3.875999in}{0.592901in}}%
\pgfpathlineto{\pgfqpoint{3.878096in}{0.576831in}}%
\pgfpathlineto{\pgfqpoint{3.880194in}{0.532305in}}%
\pgfpathlineto{\pgfqpoint{3.882291in}{0.550672in}}%
\pgfpathlineto{\pgfqpoint{3.884388in}{0.587738in}}%
\pgfpathlineto{\pgfqpoint{3.886486in}{0.605226in}}%
\pgfpathlineto{\pgfqpoint{3.888583in}{0.615233in}}%
\pgfpathlineto{\pgfqpoint{3.890680in}{0.620231in}}%
\pgfpathlineto{\pgfqpoint{3.892778in}{0.620395in}}%
\pgfpathlineto{\pgfqpoint{3.894875in}{0.614250in}}%
\pgfpathlineto{\pgfqpoint{3.896973in}{0.600716in}}%
\pgfpathlineto{\pgfqpoint{3.899070in}{0.572608in}}%
\pgfpathlineto{\pgfqpoint{3.901167in}{0.508110in}}%
\pgfpathlineto{\pgfqpoint{3.903265in}{0.582511in}}%
\pgfpathlineto{\pgfqpoint{3.905362in}{0.603428in}}%
\pgfpathlineto{\pgfqpoint{3.907459in}{0.614087in}}%
\pgfpathlineto{\pgfqpoint{3.909557in}{0.619492in}}%
\pgfpathlineto{\pgfqpoint{3.911654in}{0.621121in}}%
\pgfpathlineto{\pgfqpoint{3.913751in}{0.619297in}}%
\pgfpathlineto{\pgfqpoint{3.915849in}{0.613522in}}%
\pgfpathlineto{\pgfqpoint{3.917946in}{0.601911in}}%
\pgfpathlineto{\pgfqpoint{3.920043in}{0.577522in}}%
\pgfpathlineto{\pgfqpoint{3.922141in}{0.483184in}}%
\pgfpathlineto{\pgfqpoint{3.924238in}{0.585835in}}%
\pgfpathlineto{\pgfqpoint{3.926336in}{0.610477in}}%
\pgfpathlineto{\pgfqpoint{3.928433in}{0.623969in}}%
\pgfpathlineto{\pgfqpoint{3.930530in}{0.631765in}}%
\pgfpathlineto{\pgfqpoint{3.932628in}{0.635194in}}%
\pgfpathlineto{\pgfqpoint{3.934725in}{0.633987in}}%
\pgfpathlineto{\pgfqpoint{3.936822in}{0.626630in}}%
\pgfpathlineto{\pgfqpoint{3.938920in}{0.612049in}}%
\pgfpathlineto{\pgfqpoint{3.941017in}{0.581952in}}%
\pgfpathlineto{\pgfqpoint{3.943114in}{0.531987in}}%
\pgfpathlineto{\pgfqpoint{3.945212in}{0.595565in}}%
\pgfpathlineto{\pgfqpoint{3.947309in}{0.615181in}}%
\pgfpathlineto{\pgfqpoint{3.949407in}{0.625031in}}%
\pgfpathlineto{\pgfqpoint{3.951504in}{0.629688in}}%
\pgfpathlineto{\pgfqpoint{3.953601in}{0.630468in}}%
\pgfpathlineto{\pgfqpoint{3.955699in}{0.627523in}}%
\pgfpathlineto{\pgfqpoint{3.957796in}{0.619998in}}%
\pgfpathlineto{\pgfqpoint{3.959893in}{0.604876in}}%
\pgfpathlineto{\pgfqpoint{3.961991in}{0.567574in}}%
\pgfpathlineto{\pgfqpoint{3.964088in}{0.565707in}}%
\pgfpathlineto{\pgfqpoint{3.966185in}{0.608708in}}%
\pgfpathlineto{\pgfqpoint{3.968283in}{0.627566in}}%
\pgfpathlineto{\pgfqpoint{3.970380in}{0.638403in}}%
\pgfpathlineto{\pgfqpoint{3.972478in}{0.644259in}}%
\pgfpathlineto{\pgfqpoint{3.974575in}{0.645685in}}%
\pgfpathlineto{\pgfqpoint{3.976672in}{0.641649in}}%
\pgfpathlineto{\pgfqpoint{3.978770in}{0.632200in}}%
\pgfpathlineto{\pgfqpoint{3.980867in}{0.616782in}}%
\pgfpathlineto{\pgfqpoint{3.982964in}{0.588782in}}%
\pgfpathlineto{\pgfqpoint{3.985062in}{0.471721in}}%
\pgfpathlineto{\pgfqpoint{3.987159in}{0.586236in}}%
\pgfpathlineto{\pgfqpoint{3.989256in}{0.606563in}}%
\pgfpathlineto{\pgfqpoint{3.991354in}{0.615093in}}%
\pgfpathlineto{\pgfqpoint{3.993451in}{0.617329in}}%
\pgfpathlineto{\pgfqpoint{3.995549in}{0.614279in}}%
\pgfpathlineto{\pgfqpoint{3.997646in}{0.604739in}}%
\pgfpathlineto{\pgfqpoint{3.999743in}{0.582498in}}%
\pgfpathlineto{\pgfqpoint{4.001841in}{0.452704in}}%
\pgfpathlineto{\pgfqpoint{4.003938in}{0.589146in}}%
\pgfpathlineto{\pgfqpoint{4.006035in}{0.614924in}}%
\pgfpathlineto{\pgfqpoint{4.008133in}{0.628804in}}%
\pgfpathlineto{\pgfqpoint{4.010230in}{0.636542in}}%
\pgfpathlineto{\pgfqpoint{4.012327in}{0.639361in}}%
\pgfpathlineto{\pgfqpoint{4.014425in}{0.636506in}}%
\pgfpathlineto{\pgfqpoint{4.016522in}{0.628096in}}%
\pgfpathlineto{\pgfqpoint{4.018620in}{0.615324in}}%
\pgfpathlineto{\pgfqpoint{4.020717in}{0.595297in}}%
\pgfpathlineto{\pgfqpoint{4.022814in}{0.555225in}}%
\pgfpathlineto{\pgfqpoint{4.024912in}{0.537424in}}%
\pgfpathlineto{\pgfqpoint{4.027009in}{0.578057in}}%
\pgfpathlineto{\pgfqpoint{4.029106in}{0.590728in}}%
\pgfpathlineto{\pgfqpoint{4.031204in}{0.593878in}}%
\pgfpathlineto{\pgfqpoint{4.033301in}{0.589831in}}%
\pgfpathlineto{\pgfqpoint{4.035398in}{0.576294in}}%
\pgfpathlineto{\pgfqpoint{4.037496in}{0.535371in}}%
\pgfpathlineto{\pgfqpoint{4.039593in}{0.549972in}}%
\pgfpathlineto{\pgfqpoint{4.041691in}{0.589096in}}%
\pgfpathlineto{\pgfqpoint{4.043788in}{0.607425in}}%
\pgfpathlineto{\pgfqpoint{4.045885in}{0.617866in}}%
\pgfpathlineto{\pgfqpoint{4.047983in}{0.622941in}}%
\pgfpathlineto{\pgfqpoint{4.050080in}{0.622657in}}%
\pgfpathlineto{\pgfqpoint{4.052177in}{0.615822in}}%
\pgfpathlineto{\pgfqpoint{4.054275in}{0.603903in}}%
\pgfpathlineto{\pgfqpoint{4.056372in}{0.584166in}}%
\pgfpathlineto{\pgfqpoint{4.058469in}{0.541326in}}%
\pgfpathlineto{\pgfqpoint{4.060567in}{0.537994in}}%
\pgfpathlineto{\pgfqpoint{4.062664in}{0.575049in}}%
\pgfpathlineto{\pgfqpoint{4.064762in}{0.588570in}}%
\pgfpathlineto{\pgfqpoint{4.066859in}{0.593776in}}%
\pgfpathlineto{\pgfqpoint{4.068956in}{0.593357in}}%
\pgfpathlineto{\pgfqpoint{4.071054in}{0.587228in}}%
\pgfpathlineto{\pgfqpoint{4.073151in}{0.572363in}}%
\pgfpathlineto{\pgfqpoint{4.075248in}{0.530745in}}%
\pgfpathlineto{\pgfqpoint{4.077346in}{0.544007in}}%
\pgfpathlineto{\pgfqpoint{4.079443in}{0.582805in}}%
\pgfpathlineto{\pgfqpoint{4.081540in}{0.600821in}}%
\pgfpathlineto{\pgfqpoint{4.083638in}{0.611065in}}%
\pgfpathlineto{\pgfqpoint{4.085735in}{0.616096in}}%
\pgfpathlineto{\pgfqpoint{4.087832in}{0.616022in}}%
\pgfpathlineto{\pgfqpoint{4.089930in}{0.609350in}}%
\pgfpathlineto{\pgfqpoint{4.092027in}{0.595558in}}%
\pgfpathlineto{\pgfqpoint{4.094125in}{0.567696in}}%
\pgfpathlineto{\pgfqpoint{4.096222in}{0.494773in}}%
\pgfpathlineto{\pgfqpoint{4.098319in}{0.574996in}}%
\pgfpathlineto{\pgfqpoint{4.100417in}{0.595991in}}%
\pgfpathlineto{\pgfqpoint{4.102514in}{0.606438in}}%
\pgfpathlineto{\pgfqpoint{4.104611in}{0.611496in}}%
\pgfpathlineto{\pgfqpoint{4.106709in}{0.612630in}}%
\pgfpathlineto{\pgfqpoint{4.108806in}{0.610079in}}%
\pgfpathlineto{\pgfqpoint{4.110903in}{0.603113in}}%
\pgfpathlineto{\pgfqpoint{4.113001in}{0.589091in}}%
\pgfpathlineto{\pgfqpoint{4.115098in}{0.556331in}}%
\pgfpathlineto{\pgfqpoint{4.117196in}{0.537995in}}%
\pgfpathlineto{\pgfqpoint{4.119293in}{0.587856in}}%
\pgfpathlineto{\pgfqpoint{4.121390in}{0.608012in}}%
\pgfpathlineto{\pgfqpoint{4.123488in}{0.619497in}}%
\pgfpathlineto{\pgfqpoint{4.125585in}{0.625841in}}%
\pgfpathlineto{\pgfqpoint{4.127682in}{0.627776in}}%
\pgfpathlineto{\pgfqpoint{4.129780in}{0.624444in}}%
\pgfpathlineto{\pgfqpoint{4.131877in}{0.614226in}}%
\pgfpathlineto{\pgfqpoint{4.133974in}{0.593649in}}%
\pgfpathlineto{\pgfqpoint{4.136072in}{0.529586in}}%
\pgfpathlineto{\pgfqpoint{4.138169in}{0.576798in}}%
\pgfpathlineto{\pgfqpoint{4.140267in}{0.605763in}}%
\pgfpathlineto{\pgfqpoint{4.142364in}{0.619460in}}%
\pgfpathlineto{\pgfqpoint{4.144461in}{0.626787in}}%
\pgfpathlineto{\pgfqpoint{4.146559in}{0.630032in}}%
\pgfpathlineto{\pgfqpoint{4.148656in}{0.629895in}}%
\pgfpathlineto{\pgfqpoint{4.150753in}{0.626285in}}%
\pgfpathlineto{\pgfqpoint{4.152851in}{0.618212in}}%
\pgfpathlineto{\pgfqpoint{4.154948in}{0.602491in}}%
\pgfpathlineto{\pgfqpoint{4.157045in}{0.563546in}}%
\pgfpathlineto{\pgfqpoint{4.159143in}{0.565870in}}%
\pgfpathlineto{\pgfqpoint{4.161240in}{0.607283in}}%
\pgfpathlineto{\pgfqpoint{4.163338in}{0.625806in}}%
\pgfpathlineto{\pgfqpoint{4.165435in}{0.636544in}}%
\pgfpathlineto{\pgfqpoint{4.167532in}{0.642456in}}%
\pgfpathlineto{\pgfqpoint{4.169630in}{0.644130in}}%
\pgfpathlineto{\pgfqpoint{4.171727in}{0.640689in}}%
\pgfpathlineto{\pgfqpoint{4.173824in}{0.631344in}}%
\pgfpathlineto{\pgfqpoint{4.175922in}{0.614980in}}%
\pgfpathlineto{\pgfqpoint{4.180116in}{0.545256in}}%
\pgfpathlineto{\pgfqpoint{4.182214in}{0.599136in}}%
\pgfpathlineto{\pgfqpoint{4.184311in}{0.616740in}}%
\pgfpathlineto{\pgfqpoint{4.186409in}{0.624996in}}%
\pgfpathlineto{\pgfqpoint{4.188506in}{0.627929in}}%
\pgfpathlineto{\pgfqpoint{4.190603in}{0.626506in}}%
\pgfpathlineto{\pgfqpoint{4.192701in}{0.620277in}}%
\pgfpathlineto{\pgfqpoint{4.194798in}{0.606694in}}%
\pgfpathlineto{\pgfqpoint{4.196895in}{0.573840in}}%
\pgfpathlineto{\pgfqpoint{4.198993in}{0.557882in}}%
\pgfpathlineto{\pgfqpoint{4.201090in}{0.606949in}}%
\pgfpathlineto{\pgfqpoint{4.203187in}{0.627013in}}%
\pgfpathlineto{\pgfqpoint{4.205285in}{0.638372in}}%
\pgfpathlineto{\pgfqpoint{4.207382in}{0.644447in}}%
\pgfpathlineto{\pgfqpoint{4.209480in}{0.645825in}}%
\pgfpathlineto{\pgfqpoint{4.211577in}{0.641396in}}%
\pgfpathlineto{\pgfqpoint{4.213674in}{0.632093in}}%
\pgfpathlineto{\pgfqpoint{4.215772in}{0.617934in}}%
\pgfpathlineto{\pgfqpoint{4.217869in}{0.594909in}}%
\pgfpathlineto{\pgfqpoint{4.219966in}{0.539049in}}%
\pgfpathlineto{\pgfqpoint{4.224161in}{0.590300in}}%
\pgfpathlineto{\pgfqpoint{4.226258in}{0.600075in}}%
\pgfpathlineto{\pgfqpoint{4.228356in}{0.601614in}}%
\pgfpathlineto{\pgfqpoint{4.230453in}{0.596077in}}%
\pgfpathlineto{\pgfqpoint{4.232551in}{0.580089in}}%
\pgfpathlineto{\pgfqpoint{4.234648in}{0.524959in}}%
\pgfpathlineto{\pgfqpoint{4.238843in}{0.600208in}}%
\pgfpathlineto{\pgfqpoint{4.240940in}{0.616830in}}%
\pgfpathlineto{\pgfqpoint{4.243037in}{0.626277in}}%
\pgfpathlineto{\pgfqpoint{4.245135in}{0.630505in}}%
\pgfpathlineto{\pgfqpoint{4.247232in}{0.629197in}}%
\pgfpathlineto{\pgfqpoint{4.249329in}{0.621647in}}%
\pgfpathlineto{\pgfqpoint{4.251427in}{0.609783in}}%
\pgfpathlineto{\pgfqpoint{4.253524in}{0.591303in}}%
\pgfpathlineto{\pgfqpoint{4.257719in}{0.514923in}}%
\pgfpathlineto{\pgfqpoint{4.259816in}{0.568014in}}%
\pgfpathlineto{\pgfqpoint{4.261914in}{0.582966in}}%
\pgfpathlineto{\pgfqpoint{4.264011in}{0.587616in}}%
\pgfpathlineto{\pgfqpoint{4.266108in}{0.585286in}}%
\pgfpathlineto{\pgfqpoint{4.268206in}{0.574882in}}%
\pgfpathlineto{\pgfqpoint{4.272400in}{0.515506in}}%
\pgfpathlineto{\pgfqpoint{4.274498in}{0.573918in}}%
\pgfpathlineto{\pgfqpoint{4.276595in}{0.595836in}}%
\pgfpathlineto{\pgfqpoint{4.278692in}{0.608113in}}%
\pgfpathlineto{\pgfqpoint{4.280790in}{0.614647in}}%
\pgfpathlineto{\pgfqpoint{4.282887in}{0.616043in}}%
\pgfpathlineto{\pgfqpoint{4.284985in}{0.611158in}}%
\pgfpathlineto{\pgfqpoint{4.287082in}{0.602797in}}%
\pgfpathlineto{\pgfqpoint{4.289179in}{0.592078in}}%
\pgfpathlineto{\pgfqpoint{4.291277in}{0.578426in}}%
\pgfpathlineto{\pgfqpoint{4.293374in}{0.560802in}}%
\pgfpathlineto{\pgfqpoint{4.295471in}{0.537282in}}%
\pgfpathlineto{\pgfqpoint{4.297569in}{0.504243in}}%
\pgfpathlineto{\pgfqpoint{4.299666in}{0.460347in}}%
\pgfpathlineto{\pgfqpoint{4.301763in}{0.461578in}}%
\pgfpathlineto{\pgfqpoint{4.303861in}{0.502350in}}%
\pgfpathlineto{\pgfqpoint{4.305958in}{0.531418in}}%
\pgfpathlineto{\pgfqpoint{4.308056in}{0.550885in}}%
\pgfpathlineto{\pgfqpoint{4.310153in}{0.563982in}}%
\pgfpathlineto{\pgfqpoint{4.312250in}{0.572130in}}%
\pgfpathlineto{\pgfqpoint{4.314348in}{0.575534in}}%
\pgfpathlineto{\pgfqpoint{4.316445in}{0.572962in}}%
\pgfpathlineto{\pgfqpoint{4.316445in}{0.572962in}}%
\pgfusepath{stroke}%
\end{pgfscope}%
\begin{pgfscope}%
\pgfpathrectangle{\pgfqpoint{2.554660in}{0.320625in}}{\pgfqpoint{1.845679in}{1.164375in}}%
\pgfusepath{clip}%
\pgfsetrectcap%
\pgfsetroundjoin%
\pgfsetlinewidth{1.154312pt}%
\definecolor{currentstroke}{rgb}{0.716387,0.214982,0.475290}%
\pgfsetstrokecolor{currentstroke}%
\pgfsetdash{}{0pt}%
\pgfpathmoveto{\pgfqpoint{2.638555in}{1.468525in}}%
\pgfpathlineto{\pgfqpoint{2.651139in}{1.445854in}}%
\pgfpathlineto{\pgfqpoint{2.661626in}{1.424042in}}%
\pgfpathlineto{\pgfqpoint{2.676308in}{1.390295in}}%
\pgfpathlineto{\pgfqpoint{2.705671in}{1.322202in}}%
\pgfpathlineto{\pgfqpoint{2.722449in}{1.286001in}}%
\pgfpathlineto{\pgfqpoint{2.739228in}{1.252283in}}%
\pgfpathlineto{\pgfqpoint{2.756007in}{1.221029in}}%
\pgfpathlineto{\pgfqpoint{2.772786in}{1.192109in}}%
\pgfpathlineto{\pgfqpoint{2.789565in}{1.165274in}}%
\pgfpathlineto{\pgfqpoint{2.808441in}{1.137126in}}%
\pgfpathlineto{\pgfqpoint{2.831512in}{1.104849in}}%
\pgfpathlineto{\pgfqpoint{2.860875in}{1.065831in}}%
\pgfpathlineto{\pgfqpoint{2.904920in}{1.009355in}}%
\pgfpathlineto{\pgfqpoint{2.993009in}{0.898608in}}%
\pgfpathlineto{\pgfqpoint{3.169188in}{0.678606in}}%
\pgfpathlineto{\pgfqpoint{3.188064in}{0.657756in}}%
\pgfpathlineto{\pgfqpoint{3.196453in}{0.645702in}}%
\pgfpathlineto{\pgfqpoint{3.217427in}{0.613787in}}%
\pgfpathlineto{\pgfqpoint{3.221622in}{0.609248in}}%
\pgfpathlineto{\pgfqpoint{3.225816in}{0.606039in}}%
\pgfpathlineto{\pgfqpoint{3.230011in}{0.604147in}}%
\pgfpathlineto{\pgfqpoint{3.236303in}{0.603011in}}%
\pgfpathlineto{\pgfqpoint{3.242595in}{0.601974in}}%
\pgfpathlineto{\pgfqpoint{3.246790in}{0.599863in}}%
\pgfpathlineto{\pgfqpoint{3.250985in}{0.595206in}}%
\pgfpathlineto{\pgfqpoint{3.255180in}{0.588502in}}%
\pgfpathlineto{\pgfqpoint{3.261472in}{0.575439in}}%
\pgfpathlineto{\pgfqpoint{3.267764in}{0.561636in}}%
\pgfpathlineto{\pgfqpoint{3.269861in}{0.558522in}}%
\pgfpathlineto{\pgfqpoint{3.271958in}{0.556855in}}%
\pgfpathlineto{\pgfqpoint{3.274056in}{0.556608in}}%
\pgfpathlineto{\pgfqpoint{3.278251in}{0.558575in}}%
\pgfpathlineto{\pgfqpoint{3.282445in}{0.560643in}}%
\pgfpathlineto{\pgfqpoint{3.284543in}{0.560984in}}%
\pgfpathlineto{\pgfqpoint{3.286640in}{0.560658in}}%
\pgfpathlineto{\pgfqpoint{3.288737in}{0.559551in}}%
\pgfpathlineto{\pgfqpoint{3.290835in}{0.557531in}}%
\pgfpathlineto{\pgfqpoint{3.292932in}{0.554419in}}%
\pgfpathlineto{\pgfqpoint{3.295029in}{0.549935in}}%
\pgfpathlineto{\pgfqpoint{3.297127in}{0.543608in}}%
\pgfpathlineto{\pgfqpoint{3.299224in}{0.534553in}}%
\pgfpathlineto{\pgfqpoint{3.301322in}{0.520871in}}%
\pgfpathlineto{\pgfqpoint{3.305516in}{0.477283in}}%
\pgfpathlineto{\pgfqpoint{3.307614in}{0.509183in}}%
\pgfpathlineto{\pgfqpoint{3.309711in}{0.530723in}}%
\pgfpathlineto{\pgfqpoint{3.311808in}{0.544885in}}%
\pgfpathlineto{\pgfqpoint{3.316003in}{0.563136in}}%
\pgfpathlineto{\pgfqpoint{3.320198in}{0.574378in}}%
\pgfpathlineto{\pgfqpoint{3.324393in}{0.581256in}}%
\pgfpathlineto{\pgfqpoint{3.326490in}{0.583332in}}%
\pgfpathlineto{\pgfqpoint{3.328587in}{0.584516in}}%
\pgfpathlineto{\pgfqpoint{3.330685in}{0.584751in}}%
\pgfpathlineto{\pgfqpoint{3.332782in}{0.583911in}}%
\pgfpathlineto{\pgfqpoint{3.334879in}{0.581769in}}%
\pgfpathlineto{\pgfqpoint{3.336977in}{0.577873in}}%
\pgfpathlineto{\pgfqpoint{3.341171in}{0.565252in}}%
\pgfpathlineto{\pgfqpoint{3.343269in}{0.564450in}}%
\pgfpathlineto{\pgfqpoint{3.345366in}{0.570921in}}%
\pgfpathlineto{\pgfqpoint{3.349561in}{0.586290in}}%
\pgfpathlineto{\pgfqpoint{3.351658in}{0.591166in}}%
\pgfpathlineto{\pgfqpoint{3.353756in}{0.593959in}}%
\pgfpathlineto{\pgfqpoint{3.355853in}{0.594732in}}%
\pgfpathlineto{\pgfqpoint{3.357950in}{0.593426in}}%
\pgfpathlineto{\pgfqpoint{3.360048in}{0.589761in}}%
\pgfpathlineto{\pgfqpoint{3.362145in}{0.583079in}}%
\pgfpathlineto{\pgfqpoint{3.364242in}{0.571903in}}%
\pgfpathlineto{\pgfqpoint{3.366340in}{0.552759in}}%
\pgfpathlineto{\pgfqpoint{3.368437in}{0.528500in}}%
\pgfpathlineto{\pgfqpoint{3.372632in}{0.574980in}}%
\pgfpathlineto{\pgfqpoint{3.374729in}{0.588673in}}%
\pgfpathlineto{\pgfqpoint{3.376827in}{0.597410in}}%
\pgfpathlineto{\pgfqpoint{3.378924in}{0.602573in}}%
\pgfpathlineto{\pgfqpoint{3.381021in}{0.604503in}}%
\pgfpathlineto{\pgfqpoint{3.383119in}{0.602740in}}%
\pgfpathlineto{\pgfqpoint{3.385216in}{0.595000in}}%
\pgfpathlineto{\pgfqpoint{3.387313in}{0.573549in}}%
\pgfpathlineto{\pgfqpoint{3.389411in}{0.530712in}}%
\pgfpathlineto{\pgfqpoint{3.391508in}{0.583476in}}%
\pgfpathlineto{\pgfqpoint{3.393605in}{0.607367in}}%
\pgfpathlineto{\pgfqpoint{3.395703in}{0.621018in}}%
\pgfpathlineto{\pgfqpoint{3.397800in}{0.629611in}}%
\pgfpathlineto{\pgfqpoint{3.399898in}{0.634917in}}%
\pgfpathlineto{\pgfqpoint{3.401995in}{0.637681in}}%
\pgfpathlineto{\pgfqpoint{3.404092in}{0.638192in}}%
\pgfpathlineto{\pgfqpoint{3.406190in}{0.636430in}}%
\pgfpathlineto{\pgfqpoint{3.408287in}{0.632051in}}%
\pgfpathlineto{\pgfqpoint{3.410384in}{0.624145in}}%
\pgfpathlineto{\pgfqpoint{3.412482in}{0.610353in}}%
\pgfpathlineto{\pgfqpoint{3.414579in}{0.581941in}}%
\pgfpathlineto{\pgfqpoint{3.416676in}{0.533428in}}%
\pgfpathlineto{\pgfqpoint{3.418774in}{0.599588in}}%
\pgfpathlineto{\pgfqpoint{3.420871in}{0.622058in}}%
\pgfpathlineto{\pgfqpoint{3.422969in}{0.635035in}}%
\pgfpathlineto{\pgfqpoint{3.425066in}{0.643113in}}%
\pgfpathlineto{\pgfqpoint{3.427163in}{0.647644in}}%
\pgfpathlineto{\pgfqpoint{3.429261in}{0.648859in}}%
\pgfpathlineto{\pgfqpoint{3.431358in}{0.646102in}}%
\pgfpathlineto{\pgfqpoint{3.433455in}{0.638666in}}%
\pgfpathlineto{\pgfqpoint{3.435553in}{0.627657in}}%
\pgfpathlineto{\pgfqpoint{3.437650in}{0.611616in}}%
\pgfpathlineto{\pgfqpoint{3.439747in}{0.586053in}}%
\pgfpathlineto{\pgfqpoint{3.441845in}{0.536431in}}%
\pgfpathlineto{\pgfqpoint{3.443942in}{0.565415in}}%
\pgfpathlineto{\pgfqpoint{3.446040in}{0.584798in}}%
\pgfpathlineto{\pgfqpoint{3.448137in}{0.590256in}}%
\pgfpathlineto{\pgfqpoint{3.450234in}{0.586773in}}%
\pgfpathlineto{\pgfqpoint{3.452332in}{0.571951in}}%
\pgfpathlineto{\pgfqpoint{3.454429in}{0.519918in}}%
\pgfpathlineto{\pgfqpoint{3.458624in}{0.595283in}}%
\pgfpathlineto{\pgfqpoint{3.460721in}{0.612035in}}%
\pgfpathlineto{\pgfqpoint{3.462818in}{0.621673in}}%
\pgfpathlineto{\pgfqpoint{3.464916in}{0.626072in}}%
\pgfpathlineto{\pgfqpoint{3.467013in}{0.624884in}}%
\pgfpathlineto{\pgfqpoint{3.469111in}{0.617523in}}%
\pgfpathlineto{\pgfqpoint{3.471208in}{0.606462in}}%
\pgfpathlineto{\pgfqpoint{3.473305in}{0.590238in}}%
\pgfpathlineto{\pgfqpoint{3.475403in}{0.563900in}}%
\pgfpathlineto{\pgfqpoint{3.477500in}{0.502650in}}%
\pgfpathlineto{\pgfqpoint{3.479597in}{0.544018in}}%
\pgfpathlineto{\pgfqpoint{3.481695in}{0.564147in}}%
\pgfpathlineto{\pgfqpoint{3.483792in}{0.569726in}}%
\pgfpathlineto{\pgfqpoint{3.485889in}{0.566341in}}%
\pgfpathlineto{\pgfqpoint{3.487987in}{0.551708in}}%
\pgfpathlineto{\pgfqpoint{3.490084in}{0.499049in}}%
\pgfpathlineto{\pgfqpoint{3.494279in}{0.574061in}}%
\pgfpathlineto{\pgfqpoint{3.496376in}{0.591020in}}%
\pgfpathlineto{\pgfqpoint{3.498474in}{0.600760in}}%
\pgfpathlineto{\pgfqpoint{3.500571in}{0.605232in}}%
\pgfpathlineto{\pgfqpoint{3.502668in}{0.604122in}}%
\pgfpathlineto{\pgfqpoint{3.504766in}{0.595901in}}%
\pgfpathlineto{\pgfqpoint{3.506863in}{0.579618in}}%
\pgfpathlineto{\pgfqpoint{3.508960in}{0.543060in}}%
\pgfpathlineto{\pgfqpoint{3.511058in}{0.530447in}}%
\pgfpathlineto{\pgfqpoint{3.513155in}{0.573706in}}%
\pgfpathlineto{\pgfqpoint{3.515253in}{0.590758in}}%
\pgfpathlineto{\pgfqpoint{3.517350in}{0.599636in}}%
\pgfpathlineto{\pgfqpoint{3.519447in}{0.603818in}}%
\pgfpathlineto{\pgfqpoint{3.521545in}{0.604360in}}%
\pgfpathlineto{\pgfqpoint{3.523642in}{0.601335in}}%
\pgfpathlineto{\pgfqpoint{3.525739in}{0.593901in}}%
\pgfpathlineto{\pgfqpoint{3.527837in}{0.579207in}}%
\pgfpathlineto{\pgfqpoint{3.529934in}{0.544244in}}%
\pgfpathlineto{\pgfqpoint{3.532031in}{0.534153in}}%
\pgfpathlineto{\pgfqpoint{3.534129in}{0.580293in}}%
\pgfpathlineto{\pgfqpoint{3.536226in}{0.599787in}}%
\pgfpathlineto{\pgfqpoint{3.538324in}{0.610984in}}%
\pgfpathlineto{\pgfqpoint{3.540421in}{0.617181in}}%
\pgfpathlineto{\pgfqpoint{3.542518in}{0.619061in}}%
\pgfpathlineto{\pgfqpoint{3.544616in}{0.615785in}}%
\pgfpathlineto{\pgfqpoint{3.546713in}{0.604373in}}%
\pgfpathlineto{\pgfqpoint{3.548810in}{0.576309in}}%
\pgfpathlineto{\pgfqpoint{3.550908in}{0.533635in}}%
\pgfpathlineto{\pgfqpoint{3.553005in}{0.595649in}}%
\pgfpathlineto{\pgfqpoint{3.555102in}{0.616958in}}%
\pgfpathlineto{\pgfqpoint{3.557200in}{0.628834in}}%
\pgfpathlineto{\pgfqpoint{3.559297in}{0.635879in}}%
\pgfpathlineto{\pgfqpoint{3.561395in}{0.639624in}}%
\pgfpathlineto{\pgfqpoint{3.563492in}{0.640644in}}%
\pgfpathlineto{\pgfqpoint{3.565589in}{0.639020in}}%
\pgfpathlineto{\pgfqpoint{3.567687in}{0.634383in}}%
\pgfpathlineto{\pgfqpoint{3.569784in}{0.625617in}}%
\pgfpathlineto{\pgfqpoint{3.571881in}{0.609529in}}%
\pgfpathlineto{\pgfqpoint{3.573979in}{0.571125in}}%
\pgfpathlineto{\pgfqpoint{3.576076in}{0.570431in}}%
\pgfpathlineto{\pgfqpoint{3.578173in}{0.612706in}}%
\pgfpathlineto{\pgfqpoint{3.580271in}{0.631501in}}%
\pgfpathlineto{\pgfqpoint{3.582368in}{0.642614in}}%
\pgfpathlineto{\pgfqpoint{3.584465in}{0.649168in}}%
\pgfpathlineto{\pgfqpoint{3.586563in}{0.651990in}}%
\pgfpathlineto{\pgfqpoint{3.588660in}{0.650776in}}%
\pgfpathlineto{\pgfqpoint{3.590758in}{0.644453in}}%
\pgfpathlineto{\pgfqpoint{3.592855in}{0.634521in}}%
\pgfpathlineto{\pgfqpoint{3.594952in}{0.620280in}}%
\pgfpathlineto{\pgfqpoint{3.597050in}{0.598613in}}%
\pgfpathlineto{\pgfqpoint{3.599147in}{0.555860in}}%
\pgfpathlineto{\pgfqpoint{3.601244in}{0.542468in}}%
\pgfpathlineto{\pgfqpoint{3.603342in}{0.578167in}}%
\pgfpathlineto{\pgfqpoint{3.605439in}{0.587502in}}%
\pgfpathlineto{\pgfqpoint{3.607536in}{0.585854in}}%
\pgfpathlineto{\pgfqpoint{3.609634in}{0.572311in}}%
\pgfpathlineto{\pgfqpoint{3.611731in}{0.520143in}}%
\pgfpathlineto{\pgfqpoint{3.615926in}{0.595393in}}%
\pgfpathlineto{\pgfqpoint{3.618023in}{0.612750in}}%
\pgfpathlineto{\pgfqpoint{3.620121in}{0.622771in}}%
\pgfpathlineto{\pgfqpoint{3.622218in}{0.627477in}}%
\pgfpathlineto{\pgfqpoint{3.624315in}{0.626594in}}%
\pgfpathlineto{\pgfqpoint{3.626413in}{0.619438in}}%
\pgfpathlineto{\pgfqpoint{3.628510in}{0.608725in}}%
\pgfpathlineto{\pgfqpoint{3.630607in}{0.593290in}}%
\pgfpathlineto{\pgfqpoint{3.632705in}{0.569106in}}%
\pgfpathlineto{\pgfqpoint{3.634802in}{0.513651in}}%
\pgfpathlineto{\pgfqpoint{3.636900in}{0.531944in}}%
\pgfpathlineto{\pgfqpoint{3.638997in}{0.558164in}}%
\pgfpathlineto{\pgfqpoint{3.641094in}{0.564814in}}%
\pgfpathlineto{\pgfqpoint{3.643192in}{0.561042in}}%
\pgfpathlineto{\pgfqpoint{3.645289in}{0.543768in}}%
\pgfpathlineto{\pgfqpoint{3.647386in}{0.444665in}}%
\pgfpathlineto{\pgfqpoint{3.649484in}{0.547914in}}%
\pgfpathlineto{\pgfqpoint{3.651581in}{0.576687in}}%
\pgfpathlineto{\pgfqpoint{3.653678in}{0.592156in}}%
\pgfpathlineto{\pgfqpoint{3.655776in}{0.600990in}}%
\pgfpathlineto{\pgfqpoint{3.657873in}{0.604612in}}%
\pgfpathlineto{\pgfqpoint{3.659971in}{0.602306in}}%
\pgfpathlineto{\pgfqpoint{3.662068in}{0.592961in}}%
\pgfpathlineto{\pgfqpoint{3.664165in}{0.575306in}}%
\pgfpathlineto{\pgfqpoint{3.666263in}{0.534077in}}%
\pgfpathlineto{\pgfqpoint{3.668360in}{0.533900in}}%
\pgfpathlineto{\pgfqpoint{3.670457in}{0.572546in}}%
\pgfpathlineto{\pgfqpoint{3.672555in}{0.588319in}}%
\pgfpathlineto{\pgfqpoint{3.674652in}{0.596359in}}%
\pgfpathlineto{\pgfqpoint{3.676749in}{0.599766in}}%
\pgfpathlineto{\pgfqpoint{3.678847in}{0.599427in}}%
\pgfpathlineto{\pgfqpoint{3.680944in}{0.595211in}}%
\pgfpathlineto{\pgfqpoint{3.683042in}{0.585818in}}%
\pgfpathlineto{\pgfqpoint{3.685139in}{0.566682in}}%
\pgfpathlineto{\pgfqpoint{3.687236in}{0.503327in}}%
\pgfpathlineto{\pgfqpoint{3.689334in}{0.554538in}}%
\pgfpathlineto{\pgfqpoint{3.691431in}{0.585440in}}%
\pgfpathlineto{\pgfqpoint{3.693528in}{0.601150in}}%
\pgfpathlineto{\pgfqpoint{3.695626in}{0.610263in}}%
\pgfpathlineto{\pgfqpoint{3.697723in}{0.614745in}}%
\pgfpathlineto{\pgfqpoint{3.699820in}{0.614658in}}%
\pgfpathlineto{\pgfqpoint{3.701918in}{0.608280in}}%
\pgfpathlineto{\pgfqpoint{3.704015in}{0.591337in}}%
\pgfpathlineto{\pgfqpoint{3.706113in}{0.539665in}}%
\pgfpathlineto{\pgfqpoint{3.710307in}{0.603966in}}%
\pgfpathlineto{\pgfqpoint{3.712405in}{0.619992in}}%
\pgfpathlineto{\pgfqpoint{3.714502in}{0.629405in}}%
\pgfpathlineto{\pgfqpoint{3.716599in}{0.634875in}}%
\pgfpathlineto{\pgfqpoint{3.718697in}{0.637394in}}%
\pgfpathlineto{\pgfqpoint{3.720794in}{0.637294in}}%
\pgfpathlineto{\pgfqpoint{3.722891in}{0.634472in}}%
\pgfpathlineto{\pgfqpoint{3.724989in}{0.628313in}}%
\pgfpathlineto{\pgfqpoint{3.727086in}{0.617132in}}%
\pgfpathlineto{\pgfqpoint{3.729184in}{0.595500in}}%
\pgfpathlineto{\pgfqpoint{3.731281in}{0.508000in}}%
\pgfpathlineto{\pgfqpoint{3.733378in}{0.590562in}}%
\pgfpathlineto{\pgfqpoint{3.735476in}{0.618671in}}%
\pgfpathlineto{\pgfqpoint{3.737573in}{0.633635in}}%
\pgfpathlineto{\pgfqpoint{3.739670in}{0.642690in}}%
\pgfpathlineto{\pgfqpoint{3.741768in}{0.647699in}}%
\pgfpathlineto{\pgfqpoint{3.743865in}{0.648997in}}%
\pgfpathlineto{\pgfqpoint{3.745962in}{0.645801in}}%
\pgfpathlineto{\pgfqpoint{3.748060in}{0.637808in}}%
\pgfpathlineto{\pgfqpoint{3.750157in}{0.626071in}}%
\pgfpathlineto{\pgfqpoint{3.752254in}{0.608673in}}%
\pgfpathlineto{\pgfqpoint{3.754352in}{0.579106in}}%
\pgfpathlineto{\pgfqpoint{3.756449in}{0.453982in}}%
\pgfpathlineto{\pgfqpoint{3.758547in}{0.569520in}}%
\pgfpathlineto{\pgfqpoint{3.760644in}{0.587045in}}%
\pgfpathlineto{\pgfqpoint{3.762741in}{0.591205in}}%
\pgfpathlineto{\pgfqpoint{3.764839in}{0.586340in}}%
\pgfpathlineto{\pgfqpoint{3.766936in}{0.568720in}}%
\pgfpathlineto{\pgfqpoint{3.769033in}{0.477363in}}%
\pgfpathlineto{\pgfqpoint{3.771131in}{0.570379in}}%
\pgfpathlineto{\pgfqpoint{3.773228in}{0.599385in}}%
\pgfpathlineto{\pgfqpoint{3.775325in}{0.614821in}}%
\pgfpathlineto{\pgfqpoint{3.777423in}{0.623591in}}%
\pgfpathlineto{\pgfqpoint{3.779520in}{0.627158in}}%
\pgfpathlineto{\pgfqpoint{3.781618in}{0.624830in}}%
\pgfpathlineto{\pgfqpoint{3.783715in}{0.616759in}}%
\pgfpathlineto{\pgfqpoint{3.785812in}{0.604921in}}%
\pgfpathlineto{\pgfqpoint{3.787910in}{0.587347in}}%
\pgfpathlineto{\pgfqpoint{3.790007in}{0.557295in}}%
\pgfpathlineto{\pgfqpoint{3.792104in}{0.442491in}}%
\pgfpathlineto{\pgfqpoint{3.794202in}{0.549382in}}%
\pgfpathlineto{\pgfqpoint{3.796299in}{0.566441in}}%
\pgfpathlineto{\pgfqpoint{3.798396in}{0.570382in}}%
\pgfpathlineto{\pgfqpoint{3.800494in}{0.565302in}}%
\pgfpathlineto{\pgfqpoint{3.802591in}{0.547249in}}%
\pgfpathlineto{\pgfqpoint{3.804689in}{0.444929in}}%
\pgfpathlineto{\pgfqpoint{3.806786in}{0.550513in}}%
\pgfpathlineto{\pgfqpoint{3.808883in}{0.579014in}}%
\pgfpathlineto{\pgfqpoint{3.810981in}{0.594261in}}%
\pgfpathlineto{\pgfqpoint{3.813078in}{0.602907in}}%
\pgfpathlineto{\pgfqpoint{3.815175in}{0.606357in}}%
\pgfpathlineto{\pgfqpoint{3.817273in}{0.603867in}}%
\pgfpathlineto{\pgfqpoint{3.819370in}{0.594381in}}%
\pgfpathlineto{\pgfqpoint{3.821467in}{0.576468in}}%
\pgfpathlineto{\pgfqpoint{3.823565in}{0.534032in}}%
\pgfpathlineto{\pgfqpoint{3.825662in}{0.537107in}}%
\pgfpathlineto{\pgfqpoint{3.827760in}{0.574722in}}%
\pgfpathlineto{\pgfqpoint{3.829857in}{0.590257in}}%
\pgfpathlineto{\pgfqpoint{3.831954in}{0.598174in}}%
\pgfpathlineto{\pgfqpoint{3.834052in}{0.601492in}}%
\pgfpathlineto{\pgfqpoint{3.836149in}{0.601069in}}%
\pgfpathlineto{\pgfqpoint{3.838246in}{0.596754in}}%
\pgfpathlineto{\pgfqpoint{3.840344in}{0.587207in}}%
\pgfpathlineto{\pgfqpoint{3.842441in}{0.567709in}}%
\pgfpathlineto{\pgfqpoint{3.844538in}{0.500561in}}%
\pgfpathlineto{\pgfqpoint{3.846636in}{0.557364in}}%
\pgfpathlineto{\pgfqpoint{3.848733in}{0.587626in}}%
\pgfpathlineto{\pgfqpoint{3.850831in}{0.603129in}}%
\pgfpathlineto{\pgfqpoint{3.852928in}{0.612118in}}%
\pgfpathlineto{\pgfqpoint{3.855025in}{0.616494in}}%
\pgfpathlineto{\pgfqpoint{3.857123in}{0.616283in}}%
\pgfpathlineto{\pgfqpoint{3.859220in}{0.609751in}}%
\pgfpathlineto{\pgfqpoint{3.861317in}{0.593026in}}%
\pgfpathlineto{\pgfqpoint{3.863415in}{0.544601in}}%
\pgfpathlineto{\pgfqpoint{3.865512in}{0.568630in}}%
\pgfpathlineto{\pgfqpoint{3.867609in}{0.603046in}}%
\pgfpathlineto{\pgfqpoint{3.869707in}{0.619254in}}%
\pgfpathlineto{\pgfqpoint{3.871804in}{0.628652in}}%
\pgfpathlineto{\pgfqpoint{3.873902in}{0.634020in}}%
\pgfpathlineto{\pgfqpoint{3.875999in}{0.636367in}}%
\pgfpathlineto{\pgfqpoint{3.878096in}{0.636014in}}%
\pgfpathlineto{\pgfqpoint{3.880194in}{0.632816in}}%
\pgfpathlineto{\pgfqpoint{3.882291in}{0.626049in}}%
\pgfpathlineto{\pgfqpoint{3.884388in}{0.613701in}}%
\pgfpathlineto{\pgfqpoint{3.886486in}{0.588713in}}%
\pgfpathlineto{\pgfqpoint{3.888583in}{0.496538in}}%
\pgfpathlineto{\pgfqpoint{3.890680in}{0.596848in}}%
\pgfpathlineto{\pgfqpoint{3.892778in}{0.621390in}}%
\pgfpathlineto{\pgfqpoint{3.894875in}{0.635014in}}%
\pgfpathlineto{\pgfqpoint{3.896973in}{0.643211in}}%
\pgfpathlineto{\pgfqpoint{3.899070in}{0.647448in}}%
\pgfpathlineto{\pgfqpoint{3.901167in}{0.647825in}}%
\pgfpathlineto{\pgfqpoint{3.903265in}{0.643258in}}%
\pgfpathlineto{\pgfqpoint{3.905362in}{0.634349in}}%
\pgfpathlineto{\pgfqpoint{3.907459in}{0.621269in}}%
\pgfpathlineto{\pgfqpoint{3.909557in}{0.601300in}}%
\pgfpathlineto{\pgfqpoint{3.913751in}{0.533312in}}%
\pgfpathlineto{\pgfqpoint{3.915849in}{0.578304in}}%
\pgfpathlineto{\pgfqpoint{3.917946in}{0.590673in}}%
\pgfpathlineto{\pgfqpoint{3.920043in}{0.592462in}}%
\pgfpathlineto{\pgfqpoint{3.922141in}{0.585522in}}%
\pgfpathlineto{\pgfqpoint{3.924238in}{0.564009in}}%
\pgfpathlineto{\pgfqpoint{3.926336in}{0.481364in}}%
\pgfpathlineto{\pgfqpoint{3.928433in}{0.577610in}}%
\pgfpathlineto{\pgfqpoint{3.930530in}{0.603042in}}%
\pgfpathlineto{\pgfqpoint{3.932628in}{0.617024in}}%
\pgfpathlineto{\pgfqpoint{3.934725in}{0.624804in}}%
\pgfpathlineto{\pgfqpoint{3.936822in}{0.627416in}}%
\pgfpathlineto{\pgfqpoint{3.938920in}{0.623821in}}%
\pgfpathlineto{\pgfqpoint{3.941017in}{0.615107in}}%
\pgfpathlineto{\pgfqpoint{3.943114in}{0.602305in}}%
\pgfpathlineto{\pgfqpoint{3.945212in}{0.582859in}}%
\pgfpathlineto{\pgfqpoint{3.949407in}{0.507096in}}%
\pgfpathlineto{\pgfqpoint{3.951504in}{0.557296in}}%
\pgfpathlineto{\pgfqpoint{3.953601in}{0.570486in}}%
\pgfpathlineto{\pgfqpoint{3.955699in}{0.572754in}}%
\pgfpathlineto{\pgfqpoint{3.957796in}{0.566356in}}%
\pgfpathlineto{\pgfqpoint{3.959893in}{0.546115in}}%
\pgfpathlineto{\pgfqpoint{3.961991in}{0.425225in}}%
\pgfpathlineto{\pgfqpoint{3.964088in}{0.555706in}}%
\pgfpathlineto{\pgfqpoint{3.966185in}{0.582184in}}%
\pgfpathlineto{\pgfqpoint{3.968283in}{0.596598in}}%
\pgfpathlineto{\pgfqpoint{3.970380in}{0.604676in}}%
\pgfpathlineto{\pgfqpoint{3.972478in}{0.607583in}}%
\pgfpathlineto{\pgfqpoint{3.974575in}{0.604371in}}%
\pgfpathlineto{\pgfqpoint{3.976672in}{0.594290in}}%
\pgfpathlineto{\pgfqpoint{3.978770in}{0.575367in}}%
\pgfpathlineto{\pgfqpoint{3.980867in}{0.528170in}}%
\pgfpathlineto{\pgfqpoint{3.982964in}{0.542708in}}%
\pgfpathlineto{\pgfqpoint{3.985062in}{0.576960in}}%
\pgfpathlineto{\pgfqpoint{3.987159in}{0.591626in}}%
\pgfpathlineto{\pgfqpoint{3.989256in}{0.599047in}}%
\pgfpathlineto{\pgfqpoint{3.991354in}{0.601963in}}%
\pgfpathlineto{\pgfqpoint{3.993451in}{0.601120in}}%
\pgfpathlineto{\pgfqpoint{3.995549in}{0.596259in}}%
\pgfpathlineto{\pgfqpoint{3.997646in}{0.585803in}}%
\pgfpathlineto{\pgfqpoint{3.999743in}{0.563992in}}%
\pgfpathlineto{\pgfqpoint{4.001841in}{0.450222in}}%
\pgfpathlineto{\pgfqpoint{4.003938in}{0.563895in}}%
\pgfpathlineto{\pgfqpoint{4.006035in}{0.590891in}}%
\pgfpathlineto{\pgfqpoint{4.008133in}{0.605258in}}%
\pgfpathlineto{\pgfqpoint{4.010230in}{0.613541in}}%
\pgfpathlineto{\pgfqpoint{4.012327in}{0.617287in}}%
\pgfpathlineto{\pgfqpoint{4.014425in}{0.616296in}}%
\pgfpathlineto{\pgfqpoint{4.016522in}{0.608577in}}%
\pgfpathlineto{\pgfqpoint{4.018620in}{0.589999in}}%
\pgfpathlineto{\pgfqpoint{4.020717in}{0.530200in}}%
\pgfpathlineto{\pgfqpoint{4.022814in}{0.574042in}}%
\pgfpathlineto{\pgfqpoint{4.024912in}{0.604760in}}%
\pgfpathlineto{\pgfqpoint{4.027009in}{0.619823in}}%
\pgfpathlineto{\pgfqpoint{4.029106in}{0.628551in}}%
\pgfpathlineto{\pgfqpoint{4.031204in}{0.633393in}}%
\pgfpathlineto{\pgfqpoint{4.033301in}{0.635244in}}%
\pgfpathlineto{\pgfqpoint{4.035398in}{0.634345in}}%
\pgfpathlineto{\pgfqpoint{4.037496in}{0.630452in}}%
\pgfpathlineto{\pgfqpoint{4.039593in}{0.622638in}}%
\pgfpathlineto{\pgfqpoint{4.041691in}{0.608271in}}%
\pgfpathlineto{\pgfqpoint{4.045885in}{0.550475in}}%
\pgfpathlineto{\pgfqpoint{4.047983in}{0.603816in}}%
\pgfpathlineto{\pgfqpoint{4.050080in}{0.624563in}}%
\pgfpathlineto{\pgfqpoint{4.052177in}{0.636556in}}%
\pgfpathlineto{\pgfqpoint{4.054275in}{0.643650in}}%
\pgfpathlineto{\pgfqpoint{4.056372in}{0.646872in}}%
\pgfpathlineto{\pgfqpoint{4.058469in}{0.646010in}}%
\pgfpathlineto{\pgfqpoint{4.060567in}{0.639916in}}%
\pgfpathlineto{\pgfqpoint{4.062664in}{0.629830in}}%
\pgfpathlineto{\pgfqpoint{4.064762in}{0.614879in}}%
\pgfpathlineto{\pgfqpoint{4.066859in}{0.590835in}}%
\pgfpathlineto{\pgfqpoint{4.068956in}{0.531403in}}%
\pgfpathlineto{\pgfqpoint{4.073151in}{0.585450in}}%
\pgfpathlineto{\pgfqpoint{4.075248in}{0.593667in}}%
\pgfpathlineto{\pgfqpoint{4.077346in}{0.593079in}}%
\pgfpathlineto{\pgfqpoint{4.079443in}{0.583664in}}%
\pgfpathlineto{\pgfqpoint{4.083638in}{0.526790in}}%
\pgfpathlineto{\pgfqpoint{4.085735in}{0.584697in}}%
\pgfpathlineto{\pgfqpoint{4.087832in}{0.606780in}}%
\pgfpathlineto{\pgfqpoint{4.089930in}{0.619220in}}%
\pgfpathlineto{\pgfqpoint{4.092027in}{0.625878in}}%
\pgfpathlineto{\pgfqpoint{4.094125in}{0.627340in}}%
\pgfpathlineto{\pgfqpoint{4.096222in}{0.622373in}}%
\pgfpathlineto{\pgfqpoint{4.098319in}{0.612850in}}%
\pgfpathlineto{\pgfqpoint{4.100417in}{0.598760in}}%
\pgfpathlineto{\pgfqpoint{4.102514in}{0.576583in}}%
\pgfpathlineto{\pgfqpoint{4.104611in}{0.528455in}}%
\pgfpathlineto{\pgfqpoint{4.106709in}{0.532810in}}%
\pgfpathlineto{\pgfqpoint{4.108806in}{0.564524in}}%
\pgfpathlineto{\pgfqpoint{4.110903in}{0.574382in}}%
\pgfpathlineto{\pgfqpoint{4.113001in}{0.574955in}}%
\pgfpathlineto{\pgfqpoint{4.115098in}{0.567045in}}%
\pgfpathlineto{\pgfqpoint{4.117196in}{0.543870in}}%
\pgfpathlineto{\pgfqpoint{4.119293in}{0.481680in}}%
\pgfpathlineto{\pgfqpoint{4.121390in}{0.561283in}}%
\pgfpathlineto{\pgfqpoint{4.123488in}{0.585626in}}%
\pgfpathlineto{\pgfqpoint{4.125585in}{0.599103in}}%
\pgfpathlineto{\pgfqpoint{4.127682in}{0.606519in}}%
\pgfpathlineto{\pgfqpoint{4.129780in}{0.608778in}}%
\pgfpathlineto{\pgfqpoint{4.131877in}{0.604720in}}%
\pgfpathlineto{\pgfqpoint{4.133974in}{0.593902in}}%
\pgfpathlineto{\pgfqpoint{4.136072in}{0.573588in}}%
\pgfpathlineto{\pgfqpoint{4.138169in}{0.518089in}}%
\pgfpathlineto{\pgfqpoint{4.142364in}{0.579750in}}%
\pgfpathlineto{\pgfqpoint{4.144461in}{0.593389in}}%
\pgfpathlineto{\pgfqpoint{4.146559in}{0.600226in}}%
\pgfpathlineto{\pgfqpoint{4.148656in}{0.602674in}}%
\pgfpathlineto{\pgfqpoint{4.150753in}{0.601351in}}%
\pgfpathlineto{\pgfqpoint{4.152851in}{0.595868in}}%
\pgfpathlineto{\pgfqpoint{4.154948in}{0.584350in}}%
\pgfpathlineto{\pgfqpoint{4.157045in}{0.559599in}}%
\pgfpathlineto{\pgfqpoint{4.159143in}{0.477951in}}%
\pgfpathlineto{\pgfqpoint{4.161240in}{0.570027in}}%
\pgfpathlineto{\pgfqpoint{4.163338in}{0.594256in}}%
\pgfpathlineto{\pgfqpoint{4.165435in}{0.607543in}}%
\pgfpathlineto{\pgfqpoint{4.167532in}{0.615119in}}%
\pgfpathlineto{\pgfqpoint{4.169630in}{0.618209in}}%
\pgfpathlineto{\pgfqpoint{4.171727in}{0.616374in}}%
\pgfpathlineto{\pgfqpoint{4.173824in}{0.607428in}}%
\pgfpathlineto{\pgfqpoint{4.175922in}{0.586810in}}%
\pgfpathlineto{\pgfqpoint{4.178019in}{0.506864in}}%
\pgfpathlineto{\pgfqpoint{4.180116in}{0.578843in}}%
\pgfpathlineto{\pgfqpoint{4.182214in}{0.606502in}}%
\pgfpathlineto{\pgfqpoint{4.184311in}{0.620501in}}%
\pgfpathlineto{\pgfqpoint{4.186409in}{0.628576in}}%
\pgfpathlineto{\pgfqpoint{4.188506in}{0.632893in}}%
\pgfpathlineto{\pgfqpoint{4.190603in}{0.634233in}}%
\pgfpathlineto{\pgfqpoint{4.192701in}{0.632754in}}%
\pgfpathlineto{\pgfqpoint{4.194798in}{0.628098in}}%
\pgfpathlineto{\pgfqpoint{4.196895in}{0.619080in}}%
\pgfpathlineto{\pgfqpoint{4.198993in}{0.602176in}}%
\pgfpathlineto{\pgfqpoint{4.201090in}{0.559106in}}%
\pgfpathlineto{\pgfqpoint{4.203187in}{0.571406in}}%
\pgfpathlineto{\pgfqpoint{4.205285in}{0.609620in}}%
\pgfpathlineto{\pgfqpoint{4.207382in}{0.627461in}}%
\pgfpathlineto{\pgfqpoint{4.209480in}{0.638014in}}%
\pgfpathlineto{\pgfqpoint{4.211577in}{0.644047in}}%
\pgfpathlineto{\pgfqpoint{4.213674in}{0.646211in}}%
\pgfpathlineto{\pgfqpoint{4.215772in}{0.643948in}}%
\pgfpathlineto{\pgfqpoint{4.217869in}{0.636480in}}%
\pgfpathlineto{\pgfqpoint{4.219966in}{0.625056in}}%
\pgfpathlineto{\pgfqpoint{4.222064in}{0.607797in}}%
\pgfpathlineto{\pgfqpoint{4.224161in}{0.577725in}}%
\pgfpathlineto{\pgfqpoint{4.226258in}{0.473722in}}%
\pgfpathlineto{\pgfqpoint{4.228356in}{0.572940in}}%
\pgfpathlineto{\pgfqpoint{4.230453in}{0.590489in}}%
\pgfpathlineto{\pgfqpoint{4.232551in}{0.595638in}}%
\pgfpathlineto{\pgfqpoint{4.234648in}{0.592874in}}%
\pgfpathlineto{\pgfqpoint{4.236745in}{0.580672in}}%
\pgfpathlineto{\pgfqpoint{4.238843in}{0.544030in}}%
\pgfpathlineto{\pgfqpoint{4.240940in}{0.548673in}}%
\pgfpathlineto{\pgfqpoint{4.243037in}{0.591122in}}%
\pgfpathlineto{\pgfqpoint{4.245135in}{0.610347in}}%
\pgfpathlineto{\pgfqpoint{4.247232in}{0.621314in}}%
\pgfpathlineto{\pgfqpoint{4.249329in}{0.626818in}}%
\pgfpathlineto{\pgfqpoint{4.251427in}{0.627005in}}%
\pgfpathlineto{\pgfqpoint{4.253524in}{0.620731in}}%
\pgfpathlineto{\pgfqpoint{4.255621in}{0.610272in}}%
\pgfpathlineto{\pgfqpoint{4.257719in}{0.594621in}}%
\pgfpathlineto{\pgfqpoint{4.259816in}{0.568742in}}%
\pgfpathlineto{\pgfqpoint{4.261914in}{0.490591in}}%
\pgfpathlineto{\pgfqpoint{4.264011in}{0.547751in}}%
\pgfpathlineto{\pgfqpoint{4.266108in}{0.570264in}}%
\pgfpathlineto{\pgfqpoint{4.268206in}{0.577479in}}%
\pgfpathlineto{\pgfqpoint{4.270303in}{0.576436in}}%
\pgfpathlineto{\pgfqpoint{4.272400in}{0.566793in}}%
\pgfpathlineto{\pgfqpoint{4.276595in}{0.508705in}}%
\pgfpathlineto{\pgfqpoint{4.278692in}{0.567128in}}%
\pgfpathlineto{\pgfqpoint{4.280790in}{0.589221in}}%
\pgfpathlineto{\pgfqpoint{4.282887in}{0.601644in}}%
\pgfpathlineto{\pgfqpoint{4.284985in}{0.608288in}}%
\pgfpathlineto{\pgfqpoint{4.287082in}{0.609751in}}%
\pgfpathlineto{\pgfqpoint{4.289179in}{0.604923in}}%
\pgfpathlineto{\pgfqpoint{4.293374in}{0.588748in}}%
\pgfpathlineto{\pgfqpoint{4.301763in}{0.549653in}}%
\pgfpathlineto{\pgfqpoint{4.303861in}{0.543607in}}%
\pgfpathlineto{\pgfqpoint{4.305958in}{0.541653in}}%
\pgfpathlineto{\pgfqpoint{4.308056in}{0.543323in}}%
\pgfpathlineto{\pgfqpoint{4.312250in}{0.549985in}}%
\pgfpathlineto{\pgfqpoint{4.314348in}{0.551059in}}%
\pgfpathlineto{\pgfqpoint{4.316445in}{0.547781in}}%
\pgfpathlineto{\pgfqpoint{4.316445in}{0.547781in}}%
\pgfusepath{stroke}%
\end{pgfscope}%
\begin{pgfscope}%
\pgfpathrectangle{\pgfqpoint{2.554660in}{0.320625in}}{\pgfqpoint{1.845679in}{1.164375in}}%
\pgfusepath{clip}%
\pgfsetrectcap%
\pgfsetroundjoin%
\pgfsetlinewidth{1.154312pt}%
\definecolor{currentstroke}{rgb}{0.852126,0.276106,0.418573}%
\pgfsetstrokecolor{currentstroke}%
\pgfsetdash{}{0pt}%
\pgfpathmoveto{\pgfqpoint{2.638555in}{1.468525in}}%
\pgfpathlineto{\pgfqpoint{2.649042in}{1.449644in}}%
\pgfpathlineto{\pgfqpoint{2.659529in}{1.428050in}}%
\pgfpathlineto{\pgfqpoint{2.672113in}{1.399087in}}%
\pgfpathlineto{\pgfqpoint{2.711963in}{1.305175in}}%
\pgfpathlineto{\pgfqpoint{2.730839in}{1.264095in}}%
\pgfpathlineto{\pgfqpoint{2.749715in}{1.225519in}}%
\pgfpathlineto{\pgfqpoint{2.770689in}{1.185267in}}%
\pgfpathlineto{\pgfqpoint{2.791662in}{1.147413in}}%
\pgfpathlineto{\pgfqpoint{2.814733in}{1.108125in}}%
\pgfpathlineto{\pgfqpoint{2.839902in}{1.067513in}}%
\pgfpathlineto{\pgfqpoint{2.869265in}{1.022269in}}%
\pgfpathlineto{\pgfqpoint{2.909115in}{0.963121in}}%
\pgfpathlineto{\pgfqpoint{2.967841in}{0.878191in}}%
\pgfpathlineto{\pgfqpoint{3.062222in}{0.743874in}}%
\pgfpathlineto{\pgfqpoint{3.076904in}{0.722611in}}%
\pgfpathlineto{\pgfqpoint{3.091585in}{0.701551in}}%
\pgfpathlineto{\pgfqpoint{3.114656in}{0.670328in}}%
\pgfpathlineto{\pgfqpoint{3.137727in}{0.634044in}}%
\pgfpathlineto{\pgfqpoint{3.141922in}{0.629569in}}%
\pgfpathlineto{\pgfqpoint{3.146117in}{0.626225in}}%
\pgfpathlineto{\pgfqpoint{3.158701in}{0.617911in}}%
\pgfpathlineto{\pgfqpoint{3.162896in}{0.612135in}}%
\pgfpathlineto{\pgfqpoint{3.167090in}{0.604532in}}%
\pgfpathlineto{\pgfqpoint{3.171285in}{0.595194in}}%
\pgfpathlineto{\pgfqpoint{3.177577in}{0.577621in}}%
\pgfpathlineto{\pgfqpoint{3.190161in}{0.538413in}}%
\pgfpathlineto{\pgfqpoint{3.196453in}{0.527277in}}%
\pgfpathlineto{\pgfqpoint{3.200648in}{0.520443in}}%
\pgfpathlineto{\pgfqpoint{3.202746in}{0.519628in}}%
\pgfpathlineto{\pgfqpoint{3.204843in}{0.523163in}}%
\pgfpathlineto{\pgfqpoint{3.209038in}{0.539036in}}%
\pgfpathlineto{\pgfqpoint{3.213232in}{0.554727in}}%
\pgfpathlineto{\pgfqpoint{3.217427in}{0.566353in}}%
\pgfpathlineto{\pgfqpoint{3.221622in}{0.573986in}}%
\pgfpathlineto{\pgfqpoint{3.223719in}{0.576346in}}%
\pgfpathlineto{\pgfqpoint{3.225816in}{0.577678in}}%
\pgfpathlineto{\pgfqpoint{3.227914in}{0.577863in}}%
\pgfpathlineto{\pgfqpoint{3.230011in}{0.576672in}}%
\pgfpathlineto{\pgfqpoint{3.232109in}{0.573717in}}%
\pgfpathlineto{\pgfqpoint{3.234206in}{0.569788in}}%
\pgfpathlineto{\pgfqpoint{3.236303in}{0.570973in}}%
\pgfpathlineto{\pgfqpoint{3.240498in}{0.594455in}}%
\pgfpathlineto{\pgfqpoint{3.244693in}{0.615555in}}%
\pgfpathlineto{\pgfqpoint{3.248887in}{0.629285in}}%
\pgfpathlineto{\pgfqpoint{3.253082in}{0.637885in}}%
\pgfpathlineto{\pgfqpoint{3.257277in}{0.642641in}}%
\pgfpathlineto{\pgfqpoint{3.259374in}{0.643767in}}%
\pgfpathlineto{\pgfqpoint{3.261472in}{0.644088in}}%
\pgfpathlineto{\pgfqpoint{3.263569in}{0.643586in}}%
\pgfpathlineto{\pgfqpoint{3.265666in}{0.642205in}}%
\pgfpathlineto{\pgfqpoint{3.267764in}{0.639841in}}%
\pgfpathlineto{\pgfqpoint{3.269861in}{0.636312in}}%
\pgfpathlineto{\pgfqpoint{3.271958in}{0.631313in}}%
\pgfpathlineto{\pgfqpoint{3.274056in}{0.624299in}}%
\pgfpathlineto{\pgfqpoint{3.276153in}{0.614186in}}%
\pgfpathlineto{\pgfqpoint{3.278251in}{0.598354in}}%
\pgfpathlineto{\pgfqpoint{3.282445in}{0.541344in}}%
\pgfpathlineto{\pgfqpoint{3.284543in}{0.591149in}}%
\pgfpathlineto{\pgfqpoint{3.286640in}{0.612533in}}%
\pgfpathlineto{\pgfqpoint{3.288737in}{0.625973in}}%
\pgfpathlineto{\pgfqpoint{3.292932in}{0.642647in}}%
\pgfpathlineto{\pgfqpoint{3.297127in}{0.651970in}}%
\pgfpathlineto{\pgfqpoint{3.299224in}{0.654664in}}%
\pgfpathlineto{\pgfqpoint{3.301322in}{0.656138in}}%
\pgfpathlineto{\pgfqpoint{3.303419in}{0.656317in}}%
\pgfpathlineto{\pgfqpoint{3.305516in}{0.654992in}}%
\pgfpathlineto{\pgfqpoint{3.307614in}{0.651730in}}%
\pgfpathlineto{\pgfqpoint{3.309711in}{0.645837in}}%
\pgfpathlineto{\pgfqpoint{3.313906in}{0.628108in}}%
\pgfpathlineto{\pgfqpoint{3.318100in}{0.602874in}}%
\pgfpathlineto{\pgfqpoint{3.322295in}{0.578604in}}%
\pgfpathlineto{\pgfqpoint{3.326490in}{0.567329in}}%
\pgfpathlineto{\pgfqpoint{3.328587in}{0.564701in}}%
\pgfpathlineto{\pgfqpoint{3.330685in}{0.571442in}}%
\pgfpathlineto{\pgfqpoint{3.334879in}{0.596928in}}%
\pgfpathlineto{\pgfqpoint{3.336977in}{0.605459in}}%
\pgfpathlineto{\pgfqpoint{3.339074in}{0.609897in}}%
\pgfpathlineto{\pgfqpoint{3.341171in}{0.609322in}}%
\pgfpathlineto{\pgfqpoint{3.343269in}{0.602827in}}%
\pgfpathlineto{\pgfqpoint{3.347464in}{0.582531in}}%
\pgfpathlineto{\pgfqpoint{3.351658in}{0.550875in}}%
\pgfpathlineto{\pgfqpoint{3.353756in}{0.532092in}}%
\pgfpathlineto{\pgfqpoint{3.355853in}{0.519726in}}%
\pgfpathlineto{\pgfqpoint{3.360048in}{0.509710in}}%
\pgfpathlineto{\pgfqpoint{3.362145in}{0.510666in}}%
\pgfpathlineto{\pgfqpoint{3.366340in}{0.547149in}}%
\pgfpathlineto{\pgfqpoint{3.368437in}{0.560154in}}%
\pgfpathlineto{\pgfqpoint{3.370535in}{0.568359in}}%
\pgfpathlineto{\pgfqpoint{3.372632in}{0.571920in}}%
\pgfpathlineto{\pgfqpoint{3.374729in}{0.569750in}}%
\pgfpathlineto{\pgfqpoint{3.376827in}{0.558364in}}%
\pgfpathlineto{\pgfqpoint{3.378924in}{0.527714in}}%
\pgfpathlineto{\pgfqpoint{3.381021in}{0.515297in}}%
\pgfpathlineto{\pgfqpoint{3.383119in}{0.559669in}}%
\pgfpathlineto{\pgfqpoint{3.385216in}{0.579374in}}%
\pgfpathlineto{\pgfqpoint{3.387313in}{0.591014in}}%
\pgfpathlineto{\pgfqpoint{3.389411in}{0.598272in}}%
\pgfpathlineto{\pgfqpoint{3.391508in}{0.602498in}}%
\pgfpathlineto{\pgfqpoint{3.393605in}{0.604253in}}%
\pgfpathlineto{\pgfqpoint{3.395703in}{0.603694in}}%
\pgfpathlineto{\pgfqpoint{3.397800in}{0.600654in}}%
\pgfpathlineto{\pgfqpoint{3.399898in}{0.594528in}}%
\pgfpathlineto{\pgfqpoint{3.401995in}{0.583780in}}%
\pgfpathlineto{\pgfqpoint{3.404092in}{0.563755in}}%
\pgfpathlineto{\pgfqpoint{3.406190in}{0.500996in}}%
\pgfpathlineto{\pgfqpoint{3.408287in}{0.549974in}}%
\pgfpathlineto{\pgfqpoint{3.410384in}{0.580971in}}%
\pgfpathlineto{\pgfqpoint{3.412482in}{0.596975in}}%
\pgfpathlineto{\pgfqpoint{3.414579in}{0.606787in}}%
\pgfpathlineto{\pgfqpoint{3.416676in}{0.612611in}}%
\pgfpathlineto{\pgfqpoint{3.418774in}{0.615032in}}%
\pgfpathlineto{\pgfqpoint{3.420871in}{0.613712in}}%
\pgfpathlineto{\pgfqpoint{3.422969in}{0.607214in}}%
\pgfpathlineto{\pgfqpoint{3.425066in}{0.593911in}}%
\pgfpathlineto{\pgfqpoint{3.427163in}{0.567391in}}%
\pgfpathlineto{\pgfqpoint{3.429261in}{0.502924in}}%
\pgfpathlineto{\pgfqpoint{3.431358in}{0.571934in}}%
\pgfpathlineto{\pgfqpoint{3.433455in}{0.593472in}}%
\pgfpathlineto{\pgfqpoint{3.435553in}{0.604273in}}%
\pgfpathlineto{\pgfqpoint{3.437650in}{0.609648in}}%
\pgfpathlineto{\pgfqpoint{3.439747in}{0.611134in}}%
\pgfpathlineto{\pgfqpoint{3.441845in}{0.609043in}}%
\pgfpathlineto{\pgfqpoint{3.443942in}{0.602791in}}%
\pgfpathlineto{\pgfqpoint{3.446040in}{0.590197in}}%
\pgfpathlineto{\pgfqpoint{3.448137in}{0.562611in}}%
\pgfpathlineto{\pgfqpoint{3.450234in}{0.513985in}}%
\pgfpathlineto{\pgfqpoint{3.452332in}{0.581026in}}%
\pgfpathlineto{\pgfqpoint{3.454429in}{0.603461in}}%
\pgfpathlineto{\pgfqpoint{3.456526in}{0.616021in}}%
\pgfpathlineto{\pgfqpoint{3.458624in}{0.623169in}}%
\pgfpathlineto{\pgfqpoint{3.460721in}{0.625954in}}%
\pgfpathlineto{\pgfqpoint{3.462818in}{0.623864in}}%
\pgfpathlineto{\pgfqpoint{3.464916in}{0.615381in}}%
\pgfpathlineto{\pgfqpoint{3.467013in}{0.598925in}}%
\pgfpathlineto{\pgfqpoint{3.469111in}{0.562162in}}%
\pgfpathlineto{\pgfqpoint{3.471208in}{0.549868in}}%
\pgfpathlineto{\pgfqpoint{3.473305in}{0.592803in}}%
\pgfpathlineto{\pgfqpoint{3.475403in}{0.609707in}}%
\pgfpathlineto{\pgfqpoint{3.477500in}{0.618443in}}%
\pgfpathlineto{\pgfqpoint{3.479597in}{0.622462in}}%
\pgfpathlineto{\pgfqpoint{3.481695in}{0.622797in}}%
\pgfpathlineto{\pgfqpoint{3.483792in}{0.619478in}}%
\pgfpathlineto{\pgfqpoint{3.485889in}{0.611556in}}%
\pgfpathlineto{\pgfqpoint{3.487987in}{0.595801in}}%
\pgfpathlineto{\pgfqpoint{3.490084in}{0.555944in}}%
\pgfpathlineto{\pgfqpoint{3.492182in}{0.561319in}}%
\pgfpathlineto{\pgfqpoint{3.494279in}{0.601803in}}%
\pgfpathlineto{\pgfqpoint{3.496376in}{0.620116in}}%
\pgfpathlineto{\pgfqpoint{3.498474in}{0.630702in}}%
\pgfpathlineto{\pgfqpoint{3.500571in}{0.636418in}}%
\pgfpathlineto{\pgfqpoint{3.502668in}{0.637768in}}%
\pgfpathlineto{\pgfqpoint{3.504766in}{0.633711in}}%
\pgfpathlineto{\pgfqpoint{3.506863in}{0.623363in}}%
\pgfpathlineto{\pgfqpoint{3.508960in}{0.604210in}}%
\pgfpathlineto{\pgfqpoint{3.511058in}{0.556888in}}%
\pgfpathlineto{\pgfqpoint{3.513155in}{0.571836in}}%
\pgfpathlineto{\pgfqpoint{3.515253in}{0.605582in}}%
\pgfpathlineto{\pgfqpoint{3.517350in}{0.620020in}}%
\pgfpathlineto{\pgfqpoint{3.519447in}{0.627225in}}%
\pgfpathlineto{\pgfqpoint{3.521545in}{0.629889in}}%
\pgfpathlineto{\pgfqpoint{3.523642in}{0.628713in}}%
\pgfpathlineto{\pgfqpoint{3.525739in}{0.623346in}}%
\pgfpathlineto{\pgfqpoint{3.527837in}{0.611944in}}%
\pgfpathlineto{\pgfqpoint{3.529934in}{0.587399in}}%
\pgfpathlineto{\pgfqpoint{3.532031in}{0.502763in}}%
\pgfpathlineto{\pgfqpoint{3.534129in}{0.597350in}}%
\pgfpathlineto{\pgfqpoint{3.536226in}{0.621716in}}%
\pgfpathlineto{\pgfqpoint{3.538324in}{0.635057in}}%
\pgfpathlineto{\pgfqpoint{3.540421in}{0.642658in}}%
\pgfpathlineto{\pgfqpoint{3.542518in}{0.645751in}}%
\pgfpathlineto{\pgfqpoint{3.544616in}{0.643890in}}%
\pgfpathlineto{\pgfqpoint{3.546713in}{0.636122in}}%
\pgfpathlineto{\pgfqpoint{3.548810in}{0.623174in}}%
\pgfpathlineto{\pgfqpoint{3.550908in}{0.601207in}}%
\pgfpathlineto{\pgfqpoint{3.553005in}{0.545924in}}%
\pgfpathlineto{\pgfqpoint{3.557200in}{0.601275in}}%
\pgfpathlineto{\pgfqpoint{3.559297in}{0.612908in}}%
\pgfpathlineto{\pgfqpoint{3.561395in}{0.617158in}}%
\pgfpathlineto{\pgfqpoint{3.563492in}{0.616040in}}%
\pgfpathlineto{\pgfqpoint{3.565589in}{0.609182in}}%
\pgfpathlineto{\pgfqpoint{3.567687in}{0.593084in}}%
\pgfpathlineto{\pgfqpoint{3.569784in}{0.545060in}}%
\pgfpathlineto{\pgfqpoint{3.571881in}{0.572103in}}%
\pgfpathlineto{\pgfqpoint{3.573979in}{0.607445in}}%
\pgfpathlineto{\pgfqpoint{3.576076in}{0.624540in}}%
\pgfpathlineto{\pgfqpoint{3.578173in}{0.634282in}}%
\pgfpathlineto{\pgfqpoint{3.580271in}{0.638925in}}%
\pgfpathlineto{\pgfqpoint{3.582368in}{0.638441in}}%
\pgfpathlineto{\pgfqpoint{3.584465in}{0.631747in}}%
\pgfpathlineto{\pgfqpoint{3.586563in}{0.620632in}}%
\pgfpathlineto{\pgfqpoint{3.588660in}{0.603365in}}%
\pgfpathlineto{\pgfqpoint{3.590758in}{0.572081in}}%
\pgfpathlineto{\pgfqpoint{3.592855in}{0.503805in}}%
\pgfpathlineto{\pgfqpoint{3.594952in}{0.574768in}}%
\pgfpathlineto{\pgfqpoint{3.597050in}{0.591901in}}%
\pgfpathlineto{\pgfqpoint{3.599147in}{0.597862in}}%
\pgfpathlineto{\pgfqpoint{3.601244in}{0.596927in}}%
\pgfpathlineto{\pgfqpoint{3.603342in}{0.588808in}}%
\pgfpathlineto{\pgfqpoint{3.605439in}{0.567802in}}%
\pgfpathlineto{\pgfqpoint{3.607536in}{0.419323in}}%
\pgfpathlineto{\pgfqpoint{3.609634in}{0.574295in}}%
\pgfpathlineto{\pgfqpoint{3.611731in}{0.600739in}}%
\pgfpathlineto{\pgfqpoint{3.613829in}{0.614975in}}%
\pgfpathlineto{\pgfqpoint{3.615926in}{0.622917in}}%
\pgfpathlineto{\pgfqpoint{3.618023in}{0.625791in}}%
\pgfpathlineto{\pgfqpoint{3.620121in}{0.622749in}}%
\pgfpathlineto{\pgfqpoint{3.622218in}{0.613729in}}%
\pgfpathlineto{\pgfqpoint{3.624315in}{0.598897in}}%
\pgfpathlineto{\pgfqpoint{3.626413in}{0.572244in}}%
\pgfpathlineto{\pgfqpoint{3.628510in}{0.440571in}}%
\pgfpathlineto{\pgfqpoint{3.630607in}{0.565623in}}%
\pgfpathlineto{\pgfqpoint{3.632705in}{0.587282in}}%
\pgfpathlineto{\pgfqpoint{3.634802in}{0.596485in}}%
\pgfpathlineto{\pgfqpoint{3.636900in}{0.599325in}}%
\pgfpathlineto{\pgfqpoint{3.638997in}{0.597030in}}%
\pgfpathlineto{\pgfqpoint{3.641094in}{0.588756in}}%
\pgfpathlineto{\pgfqpoint{3.643192in}{0.569843in}}%
\pgfpathlineto{\pgfqpoint{3.645289in}{0.498110in}}%
\pgfpathlineto{\pgfqpoint{3.647386in}{0.562881in}}%
\pgfpathlineto{\pgfqpoint{3.649484in}{0.592724in}}%
\pgfpathlineto{\pgfqpoint{3.651581in}{0.608103in}}%
\pgfpathlineto{\pgfqpoint{3.653678in}{0.616829in}}%
\pgfpathlineto{\pgfqpoint{3.655776in}{0.620607in}}%
\pgfpathlineto{\pgfqpoint{3.657873in}{0.619088in}}%
\pgfpathlineto{\pgfqpoint{3.659971in}{0.610865in}}%
\pgfpathlineto{\pgfqpoint{3.662068in}{0.595107in}}%
\pgfpathlineto{\pgfqpoint{3.666263in}{0.533151in}}%
\pgfpathlineto{\pgfqpoint{3.668360in}{0.583796in}}%
\pgfpathlineto{\pgfqpoint{3.670457in}{0.601744in}}%
\pgfpathlineto{\pgfqpoint{3.672555in}{0.610810in}}%
\pgfpathlineto{\pgfqpoint{3.674652in}{0.614913in}}%
\pgfpathlineto{\pgfqpoint{3.676749in}{0.615182in}}%
\pgfpathlineto{\pgfqpoint{3.678847in}{0.611642in}}%
\pgfpathlineto{\pgfqpoint{3.680944in}{0.603236in}}%
\pgfpathlineto{\pgfqpoint{3.683042in}{0.586264in}}%
\pgfpathlineto{\pgfqpoint{3.685139in}{0.539348in}}%
\pgfpathlineto{\pgfqpoint{3.687236in}{0.561759in}}%
\pgfpathlineto{\pgfqpoint{3.689334in}{0.597621in}}%
\pgfpathlineto{\pgfqpoint{3.691431in}{0.614783in}}%
\pgfpathlineto{\pgfqpoint{3.693528in}{0.624721in}}%
\pgfpathlineto{\pgfqpoint{3.695626in}{0.629874in}}%
\pgfpathlineto{\pgfqpoint{3.697723in}{0.630541in}}%
\pgfpathlineto{\pgfqpoint{3.699820in}{0.625437in}}%
\pgfpathlineto{\pgfqpoint{3.701918in}{0.613387in}}%
\pgfpathlineto{\pgfqpoint{3.704015in}{0.589356in}}%
\pgfpathlineto{\pgfqpoint{3.706113in}{0.466424in}}%
\pgfpathlineto{\pgfqpoint{3.708210in}{0.585353in}}%
\pgfpathlineto{\pgfqpoint{3.710307in}{0.609684in}}%
\pgfpathlineto{\pgfqpoint{3.712405in}{0.621634in}}%
\pgfpathlineto{\pgfqpoint{3.714502in}{0.627804in}}%
\pgfpathlineto{\pgfqpoint{3.716599in}{0.630022in}}%
\pgfpathlineto{\pgfqpoint{3.718697in}{0.628764in}}%
\pgfpathlineto{\pgfqpoint{3.720794in}{0.623676in}}%
\pgfpathlineto{\pgfqpoint{3.722891in}{0.613174in}}%
\pgfpathlineto{\pgfqpoint{3.724989in}{0.591713in}}%
\pgfpathlineto{\pgfqpoint{3.727086in}{0.494857in}}%
\pgfpathlineto{\pgfqpoint{3.729184in}{0.589266in}}%
\pgfpathlineto{\pgfqpoint{3.731281in}{0.616863in}}%
\pgfpathlineto{\pgfqpoint{3.733378in}{0.631476in}}%
\pgfpathlineto{\pgfqpoint{3.735476in}{0.639983in}}%
\pgfpathlineto{\pgfqpoint{3.737573in}{0.644038in}}%
\pgfpathlineto{\pgfqpoint{3.739670in}{0.643585in}}%
\pgfpathlineto{\pgfqpoint{3.741768in}{0.637294in}}%
\pgfpathlineto{\pgfqpoint{3.743865in}{0.625279in}}%
\pgfpathlineto{\pgfqpoint{3.745962in}{0.604046in}}%
\pgfpathlineto{\pgfqpoint{3.748060in}{0.548370in}}%
\pgfpathlineto{\pgfqpoint{3.752254in}{0.607107in}}%
\pgfpathlineto{\pgfqpoint{3.754352in}{0.619762in}}%
\pgfpathlineto{\pgfqpoint{3.756449in}{0.625369in}}%
\pgfpathlineto{\pgfqpoint{3.758547in}{0.626188in}}%
\pgfpathlineto{\pgfqpoint{3.760644in}{0.622491in}}%
\pgfpathlineto{\pgfqpoint{3.762741in}{0.612976in}}%
\pgfpathlineto{\pgfqpoint{3.764839in}{0.592264in}}%
\pgfpathlineto{\pgfqpoint{3.766936in}{0.497428in}}%
\pgfpathlineto{\pgfqpoint{3.769033in}{0.590376in}}%
\pgfpathlineto{\pgfqpoint{3.771131in}{0.618224in}}%
\pgfpathlineto{\pgfqpoint{3.773228in}{0.632878in}}%
\pgfpathlineto{\pgfqpoint{3.775325in}{0.641215in}}%
\pgfpathlineto{\pgfqpoint{3.777423in}{0.644782in}}%
\pgfpathlineto{\pgfqpoint{3.779520in}{0.643222in}}%
\pgfpathlineto{\pgfqpoint{3.781618in}{0.635733in}}%
\pgfpathlineto{\pgfqpoint{3.783715in}{0.623845in}}%
\pgfpathlineto{\pgfqpoint{3.785812in}{0.605156in}}%
\pgfpathlineto{\pgfqpoint{3.790007in}{0.535479in}}%
\pgfpathlineto{\pgfqpoint{3.792104in}{0.584474in}}%
\pgfpathlineto{\pgfqpoint{3.794202in}{0.599125in}}%
\pgfpathlineto{\pgfqpoint{3.796299in}{0.603897in}}%
\pgfpathlineto{\pgfqpoint{3.798396in}{0.601986in}}%
\pgfpathlineto{\pgfqpoint{3.800494in}{0.592579in}}%
\pgfpathlineto{\pgfqpoint{3.802591in}{0.568313in}}%
\pgfpathlineto{\pgfqpoint{3.804689in}{0.507049in}}%
\pgfpathlineto{\pgfqpoint{3.806786in}{0.584966in}}%
\pgfpathlineto{\pgfqpoint{3.808883in}{0.608851in}}%
\pgfpathlineto{\pgfqpoint{3.810981in}{0.622018in}}%
\pgfpathlineto{\pgfqpoint{3.813078in}{0.629221in}}%
\pgfpathlineto{\pgfqpoint{3.815175in}{0.631363in}}%
\pgfpathlineto{\pgfqpoint{3.817273in}{0.627358in}}%
\pgfpathlineto{\pgfqpoint{3.819370in}{0.618067in}}%
\pgfpathlineto{\pgfqpoint{3.821467in}{0.603604in}}%
\pgfpathlineto{\pgfqpoint{3.823565in}{0.579318in}}%
\pgfpathlineto{\pgfqpoint{3.825662in}{0.510235in}}%
\pgfpathlineto{\pgfqpoint{3.827760in}{0.556852in}}%
\pgfpathlineto{\pgfqpoint{3.829857in}{0.581902in}}%
\pgfpathlineto{\pgfqpoint{3.831954in}{0.591312in}}%
\pgfpathlineto{\pgfqpoint{3.834052in}{0.593258in}}%
\pgfpathlineto{\pgfqpoint{3.836149in}{0.588918in}}%
\pgfpathlineto{\pgfqpoint{3.838246in}{0.576113in}}%
\pgfpathlineto{\pgfqpoint{3.840344in}{0.541214in}}%
\pgfpathlineto{\pgfqpoint{3.842441in}{0.537519in}}%
\pgfpathlineto{\pgfqpoint{3.844538in}{0.582329in}}%
\pgfpathlineto{\pgfqpoint{3.846636in}{0.601841in}}%
\pgfpathlineto{\pgfqpoint{3.848733in}{0.612902in}}%
\pgfpathlineto{\pgfqpoint{3.850831in}{0.618535in}}%
\pgfpathlineto{\pgfqpoint{3.852928in}{0.619023in}}%
\pgfpathlineto{\pgfqpoint{3.855025in}{0.612984in}}%
\pgfpathlineto{\pgfqpoint{3.857123in}{0.600766in}}%
\pgfpathlineto{\pgfqpoint{3.859220in}{0.578494in}}%
\pgfpathlineto{\pgfqpoint{3.861317in}{0.511463in}}%
\pgfpathlineto{\pgfqpoint{3.863415in}{0.559917in}}%
\pgfpathlineto{\pgfqpoint{3.865512in}{0.587278in}}%
\pgfpathlineto{\pgfqpoint{3.867609in}{0.599465in}}%
\pgfpathlineto{\pgfqpoint{3.869707in}{0.605079in}}%
\pgfpathlineto{\pgfqpoint{3.871804in}{0.606180in}}%
\pgfpathlineto{\pgfqpoint{3.873902in}{0.603081in}}%
\pgfpathlineto{\pgfqpoint{3.875999in}{0.594780in}}%
\pgfpathlineto{\pgfqpoint{3.878096in}{0.577330in}}%
\pgfpathlineto{\pgfqpoint{3.880194in}{0.525279in}}%
\pgfpathlineto{\pgfqpoint{3.884388in}{0.592172in}}%
\pgfpathlineto{\pgfqpoint{3.886486in}{0.608750in}}%
\pgfpathlineto{\pgfqpoint{3.888583in}{0.618283in}}%
\pgfpathlineto{\pgfqpoint{3.890680in}{0.622965in}}%
\pgfpathlineto{\pgfqpoint{3.892778in}{0.622880in}}%
\pgfpathlineto{\pgfqpoint{3.894875in}{0.616511in}}%
\pgfpathlineto{\pgfqpoint{3.896973in}{0.602499in}}%
\pgfpathlineto{\pgfqpoint{3.899070in}{0.572511in}}%
\pgfpathlineto{\pgfqpoint{3.901167in}{0.526827in}}%
\pgfpathlineto{\pgfqpoint{3.903265in}{0.588499in}}%
\pgfpathlineto{\pgfqpoint{3.905362in}{0.608403in}}%
\pgfpathlineto{\pgfqpoint{3.907459in}{0.618751in}}%
\pgfpathlineto{\pgfqpoint{3.909557in}{0.624062in}}%
\pgfpathlineto{\pgfqpoint{3.911654in}{0.625716in}}%
\pgfpathlineto{\pgfqpoint{3.913751in}{0.624023in}}%
\pgfpathlineto{\pgfqpoint{3.915849in}{0.618528in}}%
\pgfpathlineto{\pgfqpoint{3.917946in}{0.607522in}}%
\pgfpathlineto{\pgfqpoint{3.920043in}{0.584988in}}%
\pgfpathlineto{\pgfqpoint{3.922141in}{0.445185in}}%
\pgfpathlineto{\pgfqpoint{3.924238in}{0.586346in}}%
\pgfpathlineto{\pgfqpoint{3.926336in}{0.612791in}}%
\pgfpathlineto{\pgfqpoint{3.928433in}{0.626999in}}%
\pgfpathlineto{\pgfqpoint{3.930530in}{0.635294in}}%
\pgfpathlineto{\pgfqpoint{3.932628in}{0.639224in}}%
\pgfpathlineto{\pgfqpoint{3.934725in}{0.638710in}}%
\pgfpathlineto{\pgfqpoint{3.936822in}{0.632323in}}%
\pgfpathlineto{\pgfqpoint{3.938920in}{0.619277in}}%
\pgfpathlineto{\pgfqpoint{3.941017in}{0.593894in}}%
\pgfpathlineto{\pgfqpoint{3.943114in}{0.412802in}}%
\pgfpathlineto{\pgfqpoint{3.945212in}{0.591952in}}%
\pgfpathlineto{\pgfqpoint{3.947309in}{0.614976in}}%
\pgfpathlineto{\pgfqpoint{3.949407in}{0.626096in}}%
\pgfpathlineto{\pgfqpoint{3.951504in}{0.631464in}}%
\pgfpathlineto{\pgfqpoint{3.953601in}{0.632749in}}%
\pgfpathlineto{\pgfqpoint{3.955699in}{0.630247in}}%
\pgfpathlineto{\pgfqpoint{3.957796in}{0.623229in}}%
\pgfpathlineto{\pgfqpoint{3.959893in}{0.608958in}}%
\pgfpathlineto{\pgfqpoint{3.961991in}{0.574952in}}%
\pgfpathlineto{\pgfqpoint{3.964088in}{0.561970in}}%
\pgfpathlineto{\pgfqpoint{3.966185in}{0.609455in}}%
\pgfpathlineto{\pgfqpoint{3.968283in}{0.629202in}}%
\pgfpathlineto{\pgfqpoint{3.970380in}{0.640476in}}%
\pgfpathlineto{\pgfqpoint{3.972478in}{0.646639in}}%
\pgfpathlineto{\pgfqpoint{3.974575in}{0.648348in}}%
\pgfpathlineto{\pgfqpoint{3.976672in}{0.644648in}}%
\pgfpathlineto{\pgfqpoint{3.978770in}{0.635606in}}%
\pgfpathlineto{\pgfqpoint{3.980867in}{0.621208in}}%
\pgfpathlineto{\pgfqpoint{3.982964in}{0.596495in}}%
\pgfpathlineto{\pgfqpoint{3.985062in}{0.518536in}}%
\pgfpathlineto{\pgfqpoint{3.987159in}{0.578704in}}%
\pgfpathlineto{\pgfqpoint{3.989256in}{0.602898in}}%
\pgfpathlineto{\pgfqpoint{3.991354in}{0.612484in}}%
\pgfpathlineto{\pgfqpoint{3.993451in}{0.615029in}}%
\pgfpathlineto{\pgfqpoint{3.995549in}{0.611846in}}%
\pgfpathlineto{\pgfqpoint{3.997646in}{0.601573in}}%
\pgfpathlineto{\pgfqpoint{3.999743in}{0.576502in}}%
\pgfpathlineto{\pgfqpoint{4.001841in}{0.517924in}}%
\pgfpathlineto{\pgfqpoint{4.003938in}{0.593312in}}%
\pgfpathlineto{\pgfqpoint{4.006035in}{0.616821in}}%
\pgfpathlineto{\pgfqpoint{4.008133in}{0.629778in}}%
\pgfpathlineto{\pgfqpoint{4.010230in}{0.636863in}}%
\pgfpathlineto{\pgfqpoint{4.012327in}{0.638981in}}%
\pgfpathlineto{\pgfqpoint{4.014425in}{0.635097in}}%
\pgfpathlineto{\pgfqpoint{4.016522in}{0.626149in}}%
\pgfpathlineto{\pgfqpoint{4.018620in}{0.612699in}}%
\pgfpathlineto{\pgfqpoint{4.020717in}{0.591479in}}%
\pgfpathlineto{\pgfqpoint{4.022814in}{0.546718in}}%
\pgfpathlineto{\pgfqpoint{4.024912in}{0.543493in}}%
\pgfpathlineto{\pgfqpoint{4.027009in}{0.578267in}}%
\pgfpathlineto{\pgfqpoint{4.029106in}{0.589469in}}%
\pgfpathlineto{\pgfqpoint{4.031204in}{0.591539in}}%
\pgfpathlineto{\pgfqpoint{4.033301in}{0.586129in}}%
\pgfpathlineto{\pgfqpoint{4.035398in}{0.569754in}}%
\pgfpathlineto{\pgfqpoint{4.037496in}{0.509004in}}%
\pgfpathlineto{\pgfqpoint{4.039593in}{0.560509in}}%
\pgfpathlineto{\pgfqpoint{4.041691in}{0.592474in}}%
\pgfpathlineto{\pgfqpoint{4.043788in}{0.608776in}}%
\pgfpathlineto{\pgfqpoint{4.045885in}{0.618051in}}%
\pgfpathlineto{\pgfqpoint{4.047983in}{0.622122in}}%
\pgfpathlineto{\pgfqpoint{4.050080in}{0.620586in}}%
\pgfpathlineto{\pgfqpoint{4.052177in}{0.612596in}}%
\pgfpathlineto{\pgfqpoint{4.054275in}{0.599172in}}%
\pgfpathlineto{\pgfqpoint{4.056372in}{0.575854in}}%
\pgfpathlineto{\pgfqpoint{4.058469in}{0.507817in}}%
\pgfpathlineto{\pgfqpoint{4.060567in}{0.555291in}}%
\pgfpathlineto{\pgfqpoint{4.062664in}{0.581560in}}%
\pgfpathlineto{\pgfqpoint{4.064762in}{0.592465in}}%
\pgfpathlineto{\pgfqpoint{4.066859in}{0.596455in}}%
\pgfpathlineto{\pgfqpoint{4.068956in}{0.595296in}}%
\pgfpathlineto{\pgfqpoint{4.071054in}{0.588604in}}%
\pgfpathlineto{\pgfqpoint{4.073151in}{0.573123in}}%
\pgfpathlineto{\pgfqpoint{4.075248in}{0.529443in}}%
\pgfpathlineto{\pgfqpoint{4.077346in}{0.546972in}}%
\pgfpathlineto{\pgfqpoint{4.079443in}{0.584494in}}%
\pgfpathlineto{\pgfqpoint{4.081540in}{0.602158in}}%
\pgfpathlineto{\pgfqpoint{4.083638in}{0.612227in}}%
\pgfpathlineto{\pgfqpoint{4.085735in}{0.617169in}}%
\pgfpathlineto{\pgfqpoint{4.087832in}{0.617086in}}%
\pgfpathlineto{\pgfqpoint{4.089930in}{0.610450in}}%
\pgfpathlineto{\pgfqpoint{4.092027in}{0.596344in}}%
\pgfpathlineto{\pgfqpoint{4.094125in}{0.566767in}}%
\pgfpathlineto{\pgfqpoint{4.096222in}{0.515694in}}%
\pgfpathlineto{\pgfqpoint{4.098319in}{0.580527in}}%
\pgfpathlineto{\pgfqpoint{4.100417in}{0.600564in}}%
\pgfpathlineto{\pgfqpoint{4.102514in}{0.610822in}}%
\pgfpathlineto{\pgfqpoint{4.104611in}{0.615948in}}%
\pgfpathlineto{\pgfqpoint{4.106709in}{0.617324in}}%
\pgfpathlineto{\pgfqpoint{4.108806in}{0.615217in}}%
\pgfpathlineto{\pgfqpoint{4.110903in}{0.609055in}}%
\pgfpathlineto{\pgfqpoint{4.113001in}{0.596747in}}%
\pgfpathlineto{\pgfqpoint{4.115098in}{0.570268in}}%
\pgfpathlineto{\pgfqpoint{4.117196in}{0.510195in}}%
\pgfpathlineto{\pgfqpoint{4.119293in}{0.585220in}}%
\pgfpathlineto{\pgfqpoint{4.121390in}{0.608348in}}%
\pgfpathlineto{\pgfqpoint{4.123488in}{0.621215in}}%
\pgfpathlineto{\pgfqpoint{4.125585in}{0.628598in}}%
\pgfpathlineto{\pgfqpoint{4.127682in}{0.631643in}}%
\pgfpathlineto{\pgfqpoint{4.129780in}{0.629948in}}%
\pgfpathlineto{\pgfqpoint{4.131877in}{0.621821in}}%
\pgfpathlineto{\pgfqpoint{4.133974in}{0.605157in}}%
\pgfpathlineto{\pgfqpoint{4.136072in}{0.565437in}}%
\pgfpathlineto{\pgfqpoint{4.138169in}{0.563209in}}%
\pgfpathlineto{\pgfqpoint{4.140267in}{0.603425in}}%
\pgfpathlineto{\pgfqpoint{4.142364in}{0.620031in}}%
\pgfpathlineto{\pgfqpoint{4.144461in}{0.628883in}}%
\pgfpathlineto{\pgfqpoint{4.146559in}{0.633223in}}%
\pgfpathlineto{\pgfqpoint{4.148656in}{0.634076in}}%
\pgfpathlineto{\pgfqpoint{4.150753in}{0.631569in}}%
\pgfpathlineto{\pgfqpoint{4.152851in}{0.625040in}}%
\pgfpathlineto{\pgfqpoint{4.154948in}{0.612276in}}%
\pgfpathlineto{\pgfqpoint{4.157045in}{0.584697in}}%
\pgfpathlineto{\pgfqpoint{4.159143in}{0.534340in}}%
\pgfpathlineto{\pgfqpoint{4.161240in}{0.602454in}}%
\pgfpathlineto{\pgfqpoint{4.163338in}{0.624967in}}%
\pgfpathlineto{\pgfqpoint{4.165435in}{0.637596in}}%
\pgfpathlineto{\pgfqpoint{4.167532in}{0.644869in}}%
\pgfpathlineto{\pgfqpoint{4.169630in}{0.647888in}}%
\pgfpathlineto{\pgfqpoint{4.171727in}{0.646269in}}%
\pgfpathlineto{\pgfqpoint{4.173824in}{0.638789in}}%
\pgfpathlineto{\pgfqpoint{4.175922in}{0.625600in}}%
\pgfpathlineto{\pgfqpoint{4.178019in}{0.602086in}}%
\pgfpathlineto{\pgfqpoint{4.180116in}{0.527289in}}%
\pgfpathlineto{\pgfqpoint{4.182214in}{0.585555in}}%
\pgfpathlineto{\pgfqpoint{4.184311in}{0.611219in}}%
\pgfpathlineto{\pgfqpoint{4.186409in}{0.622359in}}%
\pgfpathlineto{\pgfqpoint{4.188506in}{0.626902in}}%
\pgfpathlineto{\pgfqpoint{4.190603in}{0.626657in}}%
\pgfpathlineto{\pgfqpoint{4.192701in}{0.621562in}}%
\pgfpathlineto{\pgfqpoint{4.194798in}{0.609590in}}%
\pgfpathlineto{\pgfqpoint{4.196895in}{0.581660in}}%
\pgfpathlineto{\pgfqpoint{4.198993in}{0.541037in}}%
\pgfpathlineto{\pgfqpoint{4.201090in}{0.603586in}}%
\pgfpathlineto{\pgfqpoint{4.203187in}{0.625637in}}%
\pgfpathlineto{\pgfqpoint{4.205285in}{0.637931in}}%
\pgfpathlineto{\pgfqpoint{4.207382in}{0.644650in}}%
\pgfpathlineto{\pgfqpoint{4.209480in}{0.646618in}}%
\pgfpathlineto{\pgfqpoint{4.211577in}{0.642848in}}%
\pgfpathlineto{\pgfqpoint{4.213674in}{0.634060in}}%
\pgfpathlineto{\pgfqpoint{4.215772in}{0.620913in}}%
\pgfpathlineto{\pgfqpoint{4.217869in}{0.600414in}}%
\pgfpathlineto{\pgfqpoint{4.219966in}{0.559179in}}%
\pgfpathlineto{\pgfqpoint{4.222064in}{0.544369in}}%
\pgfpathlineto{\pgfqpoint{4.224161in}{0.583283in}}%
\pgfpathlineto{\pgfqpoint{4.226258in}{0.595229in}}%
\pgfpathlineto{\pgfqpoint{4.228356in}{0.597561in}}%
\pgfpathlineto{\pgfqpoint{4.230453in}{0.592204in}}%
\pgfpathlineto{\pgfqpoint{4.232551in}{0.575610in}}%
\pgfpathlineto{\pgfqpoint{4.234648in}{0.511390in}}%
\pgfpathlineto{\pgfqpoint{4.236745in}{0.568264in}}%
\pgfpathlineto{\pgfqpoint{4.238843in}{0.599653in}}%
\pgfpathlineto{\pgfqpoint{4.240940in}{0.615789in}}%
\pgfpathlineto{\pgfqpoint{4.243037in}{0.624966in}}%
\pgfpathlineto{\pgfqpoint{4.245135in}{0.628942in}}%
\pgfpathlineto{\pgfqpoint{4.247232in}{0.627262in}}%
\pgfpathlineto{\pgfqpoint{4.249329in}{0.619407in}}%
\pgfpathlineto{\pgfqpoint{4.251427in}{0.607122in}}%
\pgfpathlineto{\pgfqpoint{4.253524in}{0.587714in}}%
\pgfpathlineto{\pgfqpoint{4.255621in}{0.549150in}}%
\pgfpathlineto{\pgfqpoint{4.257719in}{0.527114in}}%
\pgfpathlineto{\pgfqpoint{4.259816in}{0.570246in}}%
\pgfpathlineto{\pgfqpoint{4.261914in}{0.583879in}}%
\pgfpathlineto{\pgfqpoint{4.264011in}{0.588067in}}%
\pgfpathlineto{\pgfqpoint{4.266108in}{0.585587in}}%
\pgfpathlineto{\pgfqpoint{4.268206in}{0.575296in}}%
\pgfpathlineto{\pgfqpoint{4.270303in}{0.548377in}}%
\pgfpathlineto{\pgfqpoint{4.272400in}{0.510019in}}%
\pgfpathlineto{\pgfqpoint{4.274498in}{0.572337in}}%
\pgfpathlineto{\pgfqpoint{4.276595in}{0.594714in}}%
\pgfpathlineto{\pgfqpoint{4.278692in}{0.607203in}}%
\pgfpathlineto{\pgfqpoint{4.280790in}{0.613914in}}%
\pgfpathlineto{\pgfqpoint{4.282887in}{0.615541in}}%
\pgfpathlineto{\pgfqpoint{4.284985in}{0.610956in}}%
\pgfpathlineto{\pgfqpoint{4.287082in}{0.602759in}}%
\pgfpathlineto{\pgfqpoint{4.291277in}{0.579072in}}%
\pgfpathlineto{\pgfqpoint{4.293374in}{0.562207in}}%
\pgfpathlineto{\pgfqpoint{4.295471in}{0.540236in}}%
\pgfpathlineto{\pgfqpoint{4.299666in}{0.477282in}}%
\pgfpathlineto{\pgfqpoint{4.301763in}{0.474537in}}%
\pgfpathlineto{\pgfqpoint{4.305958in}{0.531294in}}%
\pgfpathlineto{\pgfqpoint{4.308056in}{0.549927in}}%
\pgfpathlineto{\pgfqpoint{4.310153in}{0.562714in}}%
\pgfpathlineto{\pgfqpoint{4.312250in}{0.570735in}}%
\pgfpathlineto{\pgfqpoint{4.314348in}{0.574089in}}%
\pgfpathlineto{\pgfqpoint{4.316445in}{0.571509in}}%
\pgfpathlineto{\pgfqpoint{4.316445in}{0.571509in}}%
\pgfusepath{stroke}%
\end{pgfscope}%
\begin{pgfscope}%
\pgfpathrectangle{\pgfqpoint{2.554660in}{0.320625in}}{\pgfqpoint{1.845679in}{1.164375in}}%
\pgfusepath{clip}%
\pgfsetrectcap%
\pgfsetroundjoin%
\pgfsetlinewidth{1.154312pt}%
\definecolor{currentstroke}{rgb}{0.950210,0.390820,0.362468}%
\pgfsetstrokecolor{currentstroke}%
\pgfsetdash{}{0pt}%
\pgfpathmoveto{\pgfqpoint{2.638555in}{1.468525in}}%
\pgfpathlineto{\pgfqpoint{2.649042in}{1.449425in}}%
\pgfpathlineto{\pgfqpoint{2.659529in}{1.427545in}}%
\pgfpathlineto{\pgfqpoint{2.672113in}{1.398195in}}%
\pgfpathlineto{\pgfqpoint{2.714060in}{1.298024in}}%
\pgfpathlineto{\pgfqpoint{2.735034in}{1.251605in}}%
\pgfpathlineto{\pgfqpoint{2.758105in}{1.203355in}}%
\pgfpathlineto{\pgfqpoint{2.783273in}{1.153309in}}%
\pgfpathlineto{\pgfqpoint{2.812636in}{1.097504in}}%
\pgfpathlineto{\pgfqpoint{2.844097in}{1.040115in}}%
\pgfpathlineto{\pgfqpoint{2.879752in}{0.977452in}}%
\pgfpathlineto{\pgfqpoint{2.919602in}{0.909725in}}%
\pgfpathlineto{\pgfqpoint{2.965744in}{0.833485in}}%
\pgfpathlineto{\pgfqpoint{3.047541in}{0.701210in}}%
\pgfpathlineto{\pgfqpoint{3.066417in}{0.672950in}}%
\pgfpathlineto{\pgfqpoint{3.089488in}{0.632424in}}%
\pgfpathlineto{\pgfqpoint{3.093683in}{0.627903in}}%
\pgfpathlineto{\pgfqpoint{3.097877in}{0.624801in}}%
\pgfpathlineto{\pgfqpoint{3.106267in}{0.619444in}}%
\pgfpathlineto{\pgfqpoint{3.108364in}{0.617049in}}%
\pgfpathlineto{\pgfqpoint{3.112559in}{0.609696in}}%
\pgfpathlineto{\pgfqpoint{3.116754in}{0.599987in}}%
\pgfpathlineto{\pgfqpoint{3.120948in}{0.587840in}}%
\pgfpathlineto{\pgfqpoint{3.131435in}{0.554821in}}%
\pgfpathlineto{\pgfqpoint{3.133533in}{0.551543in}}%
\pgfpathlineto{\pgfqpoint{3.139825in}{0.545517in}}%
\pgfpathlineto{\pgfqpoint{3.141922in}{0.541986in}}%
\pgfpathlineto{\pgfqpoint{3.148214in}{0.525543in}}%
\pgfpathlineto{\pgfqpoint{3.150311in}{0.527760in}}%
\pgfpathlineto{\pgfqpoint{3.158701in}{0.565907in}}%
\pgfpathlineto{\pgfqpoint{3.162896in}{0.577179in}}%
\pgfpathlineto{\pgfqpoint{3.164993in}{0.580713in}}%
\pgfpathlineto{\pgfqpoint{3.167090in}{0.582880in}}%
\pgfpathlineto{\pgfqpoint{3.169188in}{0.583573in}}%
\pgfpathlineto{\pgfqpoint{3.171285in}{0.582524in}}%
\pgfpathlineto{\pgfqpoint{3.173382in}{0.579056in}}%
\pgfpathlineto{\pgfqpoint{3.175480in}{0.572621in}}%
\pgfpathlineto{\pgfqpoint{3.177577in}{0.574650in}}%
\pgfpathlineto{\pgfqpoint{3.183869in}{0.621376in}}%
\pgfpathlineto{\pgfqpoint{3.188064in}{0.640333in}}%
\pgfpathlineto{\pgfqpoint{3.192259in}{0.652584in}}%
\pgfpathlineto{\pgfqpoint{3.196453in}{0.660503in}}%
\pgfpathlineto{\pgfqpoint{3.200648in}{0.665184in}}%
\pgfpathlineto{\pgfqpoint{3.202746in}{0.666476in}}%
\pgfpathlineto{\pgfqpoint{3.204843in}{0.667106in}}%
\pgfpathlineto{\pgfqpoint{3.206940in}{0.667071in}}%
\pgfpathlineto{\pgfqpoint{3.209038in}{0.666350in}}%
\pgfpathlineto{\pgfqpoint{3.211135in}{0.664888in}}%
\pgfpathlineto{\pgfqpoint{3.213232in}{0.662598in}}%
\pgfpathlineto{\pgfqpoint{3.215330in}{0.659337in}}%
\pgfpathlineto{\pgfqpoint{3.217427in}{0.654874in}}%
\pgfpathlineto{\pgfqpoint{3.219524in}{0.648822in}}%
\pgfpathlineto{\pgfqpoint{3.221622in}{0.640478in}}%
\pgfpathlineto{\pgfqpoint{3.223719in}{0.628379in}}%
\pgfpathlineto{\pgfqpoint{3.225816in}{0.608707in}}%
\pgfpathlineto{\pgfqpoint{3.227914in}{0.568904in}}%
\pgfpathlineto{\pgfqpoint{3.230011in}{0.591948in}}%
\pgfpathlineto{\pgfqpoint{3.232109in}{0.621590in}}%
\pgfpathlineto{\pgfqpoint{3.234206in}{0.638452in}}%
\pgfpathlineto{\pgfqpoint{3.238401in}{0.658506in}}%
\pgfpathlineto{\pgfqpoint{3.242595in}{0.670157in}}%
\pgfpathlineto{\pgfqpoint{3.246790in}{0.676898in}}%
\pgfpathlineto{\pgfqpoint{3.248887in}{0.678749in}}%
\pgfpathlineto{\pgfqpoint{3.250985in}{0.679590in}}%
\pgfpathlineto{\pgfqpoint{3.253082in}{0.679330in}}%
\pgfpathlineto{\pgfqpoint{3.255180in}{0.677773in}}%
\pgfpathlineto{\pgfqpoint{3.257277in}{0.674550in}}%
\pgfpathlineto{\pgfqpoint{3.261472in}{0.662579in}}%
\pgfpathlineto{\pgfqpoint{3.274056in}{0.620213in}}%
\pgfpathlineto{\pgfqpoint{3.276153in}{0.616348in}}%
\pgfpathlineto{\pgfqpoint{3.278251in}{0.614367in}}%
\pgfpathlineto{\pgfqpoint{3.280348in}{0.614475in}}%
\pgfpathlineto{\pgfqpoint{3.284543in}{0.617909in}}%
\pgfpathlineto{\pgfqpoint{3.286640in}{0.618134in}}%
\pgfpathlineto{\pgfqpoint{3.288737in}{0.614689in}}%
\pgfpathlineto{\pgfqpoint{3.292932in}{0.600888in}}%
\pgfpathlineto{\pgfqpoint{3.301322in}{0.569147in}}%
\pgfpathlineto{\pgfqpoint{3.305516in}{0.558939in}}%
\pgfpathlineto{\pgfqpoint{3.307614in}{0.556681in}}%
\pgfpathlineto{\pgfqpoint{3.309711in}{0.556803in}}%
\pgfpathlineto{\pgfqpoint{3.316003in}{0.563188in}}%
\pgfpathlineto{\pgfqpoint{3.318100in}{0.561556in}}%
\pgfpathlineto{\pgfqpoint{3.320198in}{0.553606in}}%
\pgfpathlineto{\pgfqpoint{3.324393in}{0.521172in}}%
\pgfpathlineto{\pgfqpoint{3.326490in}{0.548929in}}%
\pgfpathlineto{\pgfqpoint{3.328587in}{0.568735in}}%
\pgfpathlineto{\pgfqpoint{3.330685in}{0.581172in}}%
\pgfpathlineto{\pgfqpoint{3.332782in}{0.589256in}}%
\pgfpathlineto{\pgfqpoint{3.334879in}{0.594333in}}%
\pgfpathlineto{\pgfqpoint{3.336977in}{0.597023in}}%
\pgfpathlineto{\pgfqpoint{3.339074in}{0.597574in}}%
\pgfpathlineto{\pgfqpoint{3.341171in}{0.595966in}}%
\pgfpathlineto{\pgfqpoint{3.343269in}{0.591890in}}%
\pgfpathlineto{\pgfqpoint{3.345366in}{0.584551in}}%
\pgfpathlineto{\pgfqpoint{3.347464in}{0.571961in}}%
\pgfpathlineto{\pgfqpoint{3.349561in}{0.547508in}}%
\pgfpathlineto{\pgfqpoint{3.351658in}{0.461959in}}%
\pgfpathlineto{\pgfqpoint{3.353756in}{0.552563in}}%
\pgfpathlineto{\pgfqpoint{3.355853in}{0.577759in}}%
\pgfpathlineto{\pgfqpoint{3.357950in}{0.591888in}}%
\pgfpathlineto{\pgfqpoint{3.360048in}{0.600755in}}%
\pgfpathlineto{\pgfqpoint{3.362145in}{0.606027in}}%
\pgfpathlineto{\pgfqpoint{3.364242in}{0.608129in}}%
\pgfpathlineto{\pgfqpoint{3.366340in}{0.606698in}}%
\pgfpathlineto{\pgfqpoint{3.368437in}{0.600266in}}%
\pgfpathlineto{\pgfqpoint{3.370535in}{0.586423in}}%
\pgfpathlineto{\pgfqpoint{3.374729in}{0.526809in}}%
\pgfpathlineto{\pgfqpoint{3.376827in}{0.574594in}}%
\pgfpathlineto{\pgfqpoint{3.378924in}{0.593880in}}%
\pgfpathlineto{\pgfqpoint{3.381021in}{0.604192in}}%
\pgfpathlineto{\pgfqpoint{3.383119in}{0.609642in}}%
\pgfpathlineto{\pgfqpoint{3.385216in}{0.611556in}}%
\pgfpathlineto{\pgfqpoint{3.387313in}{0.610272in}}%
\pgfpathlineto{\pgfqpoint{3.389411in}{0.605443in}}%
\pgfpathlineto{\pgfqpoint{3.391508in}{0.595691in}}%
\pgfpathlineto{\pgfqpoint{3.393605in}{0.576528in}}%
\pgfpathlineto{\pgfqpoint{3.395703in}{0.516264in}}%
\pgfpathlineto{\pgfqpoint{3.397800in}{0.562325in}}%
\pgfpathlineto{\pgfqpoint{3.399898in}{0.593755in}}%
\pgfpathlineto{\pgfqpoint{3.401995in}{0.609692in}}%
\pgfpathlineto{\pgfqpoint{3.404092in}{0.619084in}}%
\pgfpathlineto{\pgfqpoint{3.406190in}{0.624025in}}%
\pgfpathlineto{\pgfqpoint{3.408287in}{0.624784in}}%
\pgfpathlineto{\pgfqpoint{3.410384in}{0.620199in}}%
\pgfpathlineto{\pgfqpoint{3.412482in}{0.608454in}}%
\pgfpathlineto{\pgfqpoint{3.414579in}{0.584091in}}%
\pgfpathlineto{\pgfqpoint{3.416676in}{0.508786in}}%
\pgfpathlineto{\pgfqpoint{3.418774in}{0.585198in}}%
\pgfpathlineto{\pgfqpoint{3.420871in}{0.608639in}}%
\pgfpathlineto{\pgfqpoint{3.422969in}{0.620643in}}%
\pgfpathlineto{\pgfqpoint{3.425066in}{0.627148in}}%
\pgfpathlineto{\pgfqpoint{3.427163in}{0.629909in}}%
\pgfpathlineto{\pgfqpoint{3.429261in}{0.629461in}}%
\pgfpathlineto{\pgfqpoint{3.431358in}{0.625640in}}%
\pgfpathlineto{\pgfqpoint{3.433455in}{0.617427in}}%
\pgfpathlineto{\pgfqpoint{3.435553in}{0.601633in}}%
\pgfpathlineto{\pgfqpoint{3.437650in}{0.562804in}}%
\pgfpathlineto{\pgfqpoint{3.439747in}{0.564540in}}%
\pgfpathlineto{\pgfqpoint{3.441845in}{0.606108in}}%
\pgfpathlineto{\pgfqpoint{3.443942in}{0.624684in}}%
\pgfpathlineto{\pgfqpoint{3.446040in}{0.635498in}}%
\pgfpathlineto{\pgfqpoint{3.448137in}{0.641540in}}%
\pgfpathlineto{\pgfqpoint{3.450234in}{0.643452in}}%
\pgfpathlineto{\pgfqpoint{3.452332in}{0.640489in}}%
\pgfpathlineto{\pgfqpoint{3.454429in}{0.631323in}}%
\pgfpathlineto{\pgfqpoint{3.456526in}{0.614213in}}%
\pgfpathlineto{\pgfqpoint{3.458624in}{0.576037in}}%
\pgfpathlineto{\pgfqpoint{3.460721in}{0.567339in}}%
\pgfpathlineto{\pgfqpoint{3.462818in}{0.608433in}}%
\pgfpathlineto{\pgfqpoint{3.464916in}{0.624717in}}%
\pgfpathlineto{\pgfqpoint{3.467013in}{0.632952in}}%
\pgfpathlineto{\pgfqpoint{3.469111in}{0.636439in}}%
\pgfpathlineto{\pgfqpoint{3.471208in}{0.636113in}}%
\pgfpathlineto{\pgfqpoint{3.473305in}{0.631851in}}%
\pgfpathlineto{\pgfqpoint{3.475403in}{0.622315in}}%
\pgfpathlineto{\pgfqpoint{3.477500in}{0.602743in}}%
\pgfpathlineto{\pgfqpoint{3.479597in}{0.534386in}}%
\pgfpathlineto{\pgfqpoint{3.481695in}{0.593026in}}%
\pgfpathlineto{\pgfqpoint{3.483792in}{0.623100in}}%
\pgfpathlineto{\pgfqpoint{3.485889in}{0.638543in}}%
\pgfpathlineto{\pgfqpoint{3.487987in}{0.647490in}}%
\pgfpathlineto{\pgfqpoint{3.490084in}{0.651814in}}%
\pgfpathlineto{\pgfqpoint{3.492182in}{0.651520in}}%
\pgfpathlineto{\pgfqpoint{3.494279in}{0.645411in}}%
\pgfpathlineto{\pgfqpoint{3.496376in}{0.634766in}}%
\pgfpathlineto{\pgfqpoint{3.498474in}{0.618177in}}%
\pgfpathlineto{\pgfqpoint{3.500571in}{0.588597in}}%
\pgfpathlineto{\pgfqpoint{3.502668in}{0.507765in}}%
\pgfpathlineto{\pgfqpoint{3.504766in}{0.587066in}}%
\pgfpathlineto{\pgfqpoint{3.506863in}{0.605481in}}%
\pgfpathlineto{\pgfqpoint{3.508960in}{0.612257in}}%
\pgfpathlineto{\pgfqpoint{3.511058in}{0.612219in}}%
\pgfpathlineto{\pgfqpoint{3.513155in}{0.605527in}}%
\pgfpathlineto{\pgfqpoint{3.515253in}{0.588154in}}%
\pgfpathlineto{\pgfqpoint{3.517350in}{0.524772in}}%
\pgfpathlineto{\pgfqpoint{3.519447in}{0.579192in}}%
\pgfpathlineto{\pgfqpoint{3.521545in}{0.610428in}}%
\pgfpathlineto{\pgfqpoint{3.523642in}{0.626381in}}%
\pgfpathlineto{\pgfqpoint{3.525739in}{0.635432in}}%
\pgfpathlineto{\pgfqpoint{3.527837in}{0.639365in}}%
\pgfpathlineto{\pgfqpoint{3.529934in}{0.637791in}}%
\pgfpathlineto{\pgfqpoint{3.532031in}{0.630231in}}%
\pgfpathlineto{\pgfqpoint{3.534129in}{0.618864in}}%
\pgfpathlineto{\pgfqpoint{3.536226in}{0.601963in}}%
\pgfpathlineto{\pgfqpoint{3.538324in}{0.573440in}}%
\pgfpathlineto{\pgfqpoint{3.540421in}{0.467941in}}%
\pgfpathlineto{\pgfqpoint{3.542518in}{0.561447in}}%
\pgfpathlineto{\pgfqpoint{3.544616in}{0.580165in}}%
\pgfpathlineto{\pgfqpoint{3.546713in}{0.585279in}}%
\pgfpathlineto{\pgfqpoint{3.548810in}{0.581809in}}%
\pgfpathlineto{\pgfqpoint{3.550908in}{0.567514in}}%
\pgfpathlineto{\pgfqpoint{3.553005in}{0.517555in}}%
\pgfpathlineto{\pgfqpoint{3.557200in}{0.587456in}}%
\pgfpathlineto{\pgfqpoint{3.559297in}{0.604783in}}%
\pgfpathlineto{\pgfqpoint{3.561395in}{0.614715in}}%
\pgfpathlineto{\pgfqpoint{3.563492in}{0.619350in}}%
\pgfpathlineto{\pgfqpoint{3.565589in}{0.618456in}}%
\pgfpathlineto{\pgfqpoint{3.567687in}{0.610957in}}%
\pgfpathlineto{\pgfqpoint{3.569784in}{0.598062in}}%
\pgfpathlineto{\pgfqpoint{3.571881in}{0.575844in}}%
\pgfpathlineto{\pgfqpoint{3.573979in}{0.517059in}}%
\pgfpathlineto{\pgfqpoint{3.578173in}{0.578381in}}%
\pgfpathlineto{\pgfqpoint{3.580271in}{0.590020in}}%
\pgfpathlineto{\pgfqpoint{3.582368in}{0.594518in}}%
\pgfpathlineto{\pgfqpoint{3.584465in}{0.593876in}}%
\pgfpathlineto{\pgfqpoint{3.586563in}{0.587906in}}%
\pgfpathlineto{\pgfqpoint{3.588660in}{0.573895in}}%
\pgfpathlineto{\pgfqpoint{3.590758in}{0.537507in}}%
\pgfpathlineto{\pgfqpoint{3.592855in}{0.535755in}}%
\pgfpathlineto{\pgfqpoint{3.594952in}{0.579265in}}%
\pgfpathlineto{\pgfqpoint{3.597050in}{0.598333in}}%
\pgfpathlineto{\pgfqpoint{3.599147in}{0.609163in}}%
\pgfpathlineto{\pgfqpoint{3.601244in}{0.614730in}}%
\pgfpathlineto{\pgfqpoint{3.603342in}{0.615381in}}%
\pgfpathlineto{\pgfqpoint{3.605439in}{0.609668in}}%
\pgfpathlineto{\pgfqpoint{3.607536in}{0.596271in}}%
\pgfpathlineto{\pgfqpoint{3.609634in}{0.567667in}}%
\pgfpathlineto{\pgfqpoint{3.611731in}{0.512334in}}%
\pgfpathlineto{\pgfqpoint{3.613829in}{0.580618in}}%
\pgfpathlineto{\pgfqpoint{3.615926in}{0.601322in}}%
\pgfpathlineto{\pgfqpoint{3.618023in}{0.612105in}}%
\pgfpathlineto{\pgfqpoint{3.620121in}{0.617776in}}%
\pgfpathlineto{\pgfqpoint{3.622218in}{0.619802in}}%
\pgfpathlineto{\pgfqpoint{3.624315in}{0.618563in}}%
\pgfpathlineto{\pgfqpoint{3.626413in}{0.613719in}}%
\pgfpathlineto{\pgfqpoint{3.628510in}{0.603864in}}%
\pgfpathlineto{\pgfqpoint{3.630607in}{0.584350in}}%
\pgfpathlineto{\pgfqpoint{3.632705in}{0.520104in}}%
\pgfpathlineto{\pgfqpoint{3.634802in}{0.572189in}}%
\pgfpathlineto{\pgfqpoint{3.636900in}{0.602906in}}%
\pgfpathlineto{\pgfqpoint{3.638997in}{0.618598in}}%
\pgfpathlineto{\pgfqpoint{3.641094in}{0.627815in}}%
\pgfpathlineto{\pgfqpoint{3.643192in}{0.632567in}}%
\pgfpathlineto{\pgfqpoint{3.645289in}{0.633056in}}%
\pgfpathlineto{\pgfqpoint{3.647386in}{0.628005in}}%
\pgfpathlineto{\pgfqpoint{3.649484in}{0.615546in}}%
\pgfpathlineto{\pgfqpoint{3.651581in}{0.589132in}}%
\pgfpathlineto{\pgfqpoint{3.653678in}{0.509028in}}%
\pgfpathlineto{\pgfqpoint{3.655776in}{0.596344in}}%
\pgfpathlineto{\pgfqpoint{3.657873in}{0.618505in}}%
\pgfpathlineto{\pgfqpoint{3.659971in}{0.629971in}}%
\pgfpathlineto{\pgfqpoint{3.662068in}{0.636139in}}%
\pgfpathlineto{\pgfqpoint{3.664165in}{0.638632in}}%
\pgfpathlineto{\pgfqpoint{3.666263in}{0.637921in}}%
\pgfpathlineto{\pgfqpoint{3.668360in}{0.633786in}}%
\pgfpathlineto{\pgfqpoint{3.670457in}{0.625102in}}%
\pgfpathlineto{\pgfqpoint{3.672555in}{0.608321in}}%
\pgfpathlineto{\pgfqpoint{3.674652in}{0.564503in}}%
\pgfpathlineto{\pgfqpoint{3.676749in}{0.579225in}}%
\pgfpathlineto{\pgfqpoint{3.678847in}{0.616862in}}%
\pgfpathlineto{\pgfqpoint{3.680944in}{0.634513in}}%
\pgfpathlineto{\pgfqpoint{3.683042in}{0.644840in}}%
\pgfpathlineto{\pgfqpoint{3.685139in}{0.650506in}}%
\pgfpathlineto{\pgfqpoint{3.687236in}{0.652023in}}%
\pgfpathlineto{\pgfqpoint{3.689334in}{0.648495in}}%
\pgfpathlineto{\pgfqpoint{3.691431in}{0.639394in}}%
\pgfpathlineto{\pgfqpoint{3.693528in}{0.624151in}}%
\pgfpathlineto{\pgfqpoint{3.695626in}{0.595855in}}%
\pgfpathlineto{\pgfqpoint{3.697723in}{0.500952in}}%
\pgfpathlineto{\pgfqpoint{3.699820in}{0.596458in}}%
\pgfpathlineto{\pgfqpoint{3.701918in}{0.616686in}}%
\pgfpathlineto{\pgfqpoint{3.704015in}{0.625584in}}%
\pgfpathlineto{\pgfqpoint{3.706113in}{0.628481in}}%
\pgfpathlineto{\pgfqpoint{3.708210in}{0.626521in}}%
\pgfpathlineto{\pgfqpoint{3.710307in}{0.619027in}}%
\pgfpathlineto{\pgfqpoint{3.712405in}{0.602250in}}%
\pgfpathlineto{\pgfqpoint{3.714502in}{0.551666in}}%
\pgfpathlineto{\pgfqpoint{3.718697in}{0.617340in}}%
\pgfpathlineto{\pgfqpoint{3.720794in}{0.634111in}}%
\pgfpathlineto{\pgfqpoint{3.722891in}{0.643688in}}%
\pgfpathlineto{\pgfqpoint{3.724989in}{0.648258in}}%
\pgfpathlineto{\pgfqpoint{3.727086in}{0.647801in}}%
\pgfpathlineto{\pgfqpoint{3.729184in}{0.641301in}}%
\pgfpathlineto{\pgfqpoint{3.731281in}{0.630986in}}%
\pgfpathlineto{\pgfqpoint{3.733378in}{0.615953in}}%
\pgfpathlineto{\pgfqpoint{3.735476in}{0.592234in}}%
\pgfpathlineto{\pgfqpoint{3.737573in}{0.537366in}}%
\pgfpathlineto{\pgfqpoint{3.739670in}{0.555155in}}%
\pgfpathlineto{\pgfqpoint{3.741768in}{0.582537in}}%
\pgfpathlineto{\pgfqpoint{3.743865in}{0.590348in}}%
\pgfpathlineto{\pgfqpoint{3.745962in}{0.588520in}}%
\pgfpathlineto{\pgfqpoint{3.748060in}{0.576069in}}%
\pgfpathlineto{\pgfqpoint{3.750157in}{0.534186in}}%
\pgfpathlineto{\pgfqpoint{3.752254in}{0.553960in}}%
\pgfpathlineto{\pgfqpoint{3.754352in}{0.592413in}}%
\pgfpathlineto{\pgfqpoint{3.756449in}{0.610845in}}%
\pgfpathlineto{\pgfqpoint{3.758547in}{0.621445in}}%
\pgfpathlineto{\pgfqpoint{3.760644in}{0.626649in}}%
\pgfpathlineto{\pgfqpoint{3.762741in}{0.626421in}}%
\pgfpathlineto{\pgfqpoint{3.764839in}{0.619685in}}%
\pgfpathlineto{\pgfqpoint{3.766936in}{0.608632in}}%
\pgfpathlineto{\pgfqpoint{3.769033in}{0.591659in}}%
\pgfpathlineto{\pgfqpoint{3.771131in}{0.561631in}}%
\pgfpathlineto{\pgfqpoint{3.773228in}{0.467744in}}%
\pgfpathlineto{\pgfqpoint{3.775325in}{0.559166in}}%
\pgfpathlineto{\pgfqpoint{3.777423in}{0.577110in}}%
\pgfpathlineto{\pgfqpoint{3.779520in}{0.583133in}}%
\pgfpathlineto{\pgfqpoint{3.781618in}{0.581910in}}%
\pgfpathlineto{\pgfqpoint{3.783715in}{0.572993in}}%
\pgfpathlineto{\pgfqpoint{3.785812in}{0.549306in}}%
\pgfpathlineto{\pgfqpoint{3.787910in}{0.484270in}}%
\pgfpathlineto{\pgfqpoint{3.790007in}{0.565302in}}%
\pgfpathlineto{\pgfqpoint{3.792104in}{0.589503in}}%
\pgfpathlineto{\pgfqpoint{3.794202in}{0.602837in}}%
\pgfpathlineto{\pgfqpoint{3.796299in}{0.610154in}}%
\pgfpathlineto{\pgfqpoint{3.798396in}{0.612383in}}%
\pgfpathlineto{\pgfqpoint{3.800494in}{0.608414in}}%
\pgfpathlineto{\pgfqpoint{3.802591in}{0.597696in}}%
\pgfpathlineto{\pgfqpoint{3.804689in}{0.577520in}}%
\pgfpathlineto{\pgfqpoint{3.806786in}{0.522698in}}%
\pgfpathlineto{\pgfqpoint{3.810981in}{0.583504in}}%
\pgfpathlineto{\pgfqpoint{3.813078in}{0.597256in}}%
\pgfpathlineto{\pgfqpoint{3.815175in}{0.604190in}}%
\pgfpathlineto{\pgfqpoint{3.817273in}{0.606744in}}%
\pgfpathlineto{\pgfqpoint{3.819370in}{0.605560in}}%
\pgfpathlineto{\pgfqpoint{3.821467in}{0.600287in}}%
\pgfpathlineto{\pgfqpoint{3.823565in}{0.589164in}}%
\pgfpathlineto{\pgfqpoint{3.825662in}{0.565590in}}%
\pgfpathlineto{\pgfqpoint{3.827760in}{0.450445in}}%
\pgfpathlineto{\pgfqpoint{3.829857in}{0.571987in}}%
\pgfpathlineto{\pgfqpoint{3.831954in}{0.597205in}}%
\pgfpathlineto{\pgfqpoint{3.834052in}{0.610896in}}%
\pgfpathlineto{\pgfqpoint{3.836149in}{0.618752in}}%
\pgfpathlineto{\pgfqpoint{3.838246in}{0.622118in}}%
\pgfpathlineto{\pgfqpoint{3.840344in}{0.620664in}}%
\pgfpathlineto{\pgfqpoint{3.842441in}{0.612454in}}%
\pgfpathlineto{\pgfqpoint{3.844538in}{0.594140in}}%
\pgfpathlineto{\pgfqpoint{3.846636in}{0.541051in}}%
\pgfpathlineto{\pgfqpoint{3.850831in}{0.604284in}}%
\pgfpathlineto{\pgfqpoint{3.852928in}{0.619491in}}%
\pgfpathlineto{\pgfqpoint{3.855025in}{0.628051in}}%
\pgfpathlineto{\pgfqpoint{3.857123in}{0.632563in}}%
\pgfpathlineto{\pgfqpoint{3.859220in}{0.633929in}}%
\pgfpathlineto{\pgfqpoint{3.861317in}{0.632332in}}%
\pgfpathlineto{\pgfqpoint{3.863415in}{0.627377in}}%
\pgfpathlineto{\pgfqpoint{3.865512in}{0.617722in}}%
\pgfpathlineto{\pgfqpoint{3.867609in}{0.599182in}}%
\pgfpathlineto{\pgfqpoint{3.869707in}{0.545217in}}%
\pgfpathlineto{\pgfqpoint{3.873902in}{0.613369in}}%
\pgfpathlineto{\pgfqpoint{3.875999in}{0.629845in}}%
\pgfpathlineto{\pgfqpoint{3.878096in}{0.639613in}}%
\pgfpathlineto{\pgfqpoint{3.880194in}{0.644968in}}%
\pgfpathlineto{\pgfqpoint{3.882291in}{0.646332in}}%
\pgfpathlineto{\pgfqpoint{3.884388in}{0.642822in}}%
\pgfpathlineto{\pgfqpoint{3.886486in}{0.634397in}}%
\pgfpathlineto{\pgfqpoint{3.888583in}{0.621712in}}%
\pgfpathlineto{\pgfqpoint{3.890680in}{0.602020in}}%
\pgfpathlineto{\pgfqpoint{3.892778in}{0.563699in}}%
\pgfpathlineto{\pgfqpoint{3.894875in}{0.538049in}}%
\pgfpathlineto{\pgfqpoint{3.896973in}{0.581992in}}%
\pgfpathlineto{\pgfqpoint{3.899070in}{0.595058in}}%
\pgfpathlineto{\pgfqpoint{3.901167in}{0.598205in}}%
\pgfpathlineto{\pgfqpoint{3.903265in}{0.593888in}}%
\pgfpathlineto{\pgfqpoint{3.905362in}{0.579521in}}%
\pgfpathlineto{\pgfqpoint{3.907459in}{0.533210in}}%
\pgfpathlineto{\pgfqpoint{3.909557in}{0.559919in}}%
\pgfpathlineto{\pgfqpoint{3.911654in}{0.596126in}}%
\pgfpathlineto{\pgfqpoint{3.913751in}{0.613720in}}%
\pgfpathlineto{\pgfqpoint{3.915849in}{0.623768in}}%
\pgfpathlineto{\pgfqpoint{3.917946in}{0.628519in}}%
\pgfpathlineto{\pgfqpoint{3.920043in}{0.627841in}}%
\pgfpathlineto{\pgfqpoint{3.922141in}{0.620821in}}%
\pgfpathlineto{\pgfqpoint{3.924238in}{0.609682in}}%
\pgfpathlineto{\pgfqpoint{3.926336in}{0.592784in}}%
\pgfpathlineto{\pgfqpoint{3.928433in}{0.563435in}}%
\pgfpathlineto{\pgfqpoint{3.930530in}{0.439224in}}%
\pgfpathlineto{\pgfqpoint{3.932628in}{0.557186in}}%
\pgfpathlineto{\pgfqpoint{3.934725in}{0.575528in}}%
\pgfpathlineto{\pgfqpoint{3.936822in}{0.581320in}}%
\pgfpathlineto{\pgfqpoint{3.938920in}{0.579444in}}%
\pgfpathlineto{\pgfqpoint{3.941017in}{0.569027in}}%
\pgfpathlineto{\pgfqpoint{3.945212in}{0.516919in}}%
\pgfpathlineto{\pgfqpoint{3.947309in}{0.570992in}}%
\pgfpathlineto{\pgfqpoint{3.949407in}{0.592384in}}%
\pgfpathlineto{\pgfqpoint{3.951504in}{0.604452in}}%
\pgfpathlineto{\pgfqpoint{3.953601in}{0.610837in}}%
\pgfpathlineto{\pgfqpoint{3.955699in}{0.612049in}}%
\pgfpathlineto{\pgfqpoint{3.957796in}{0.606716in}}%
\pgfpathlineto{\pgfqpoint{3.959893in}{0.594795in}}%
\pgfpathlineto{\pgfqpoint{3.961991in}{0.572395in}}%
\pgfpathlineto{\pgfqpoint{3.964088in}{0.499382in}}%
\pgfpathlineto{\pgfqpoint{3.966185in}{0.557468in}}%
\pgfpathlineto{\pgfqpoint{3.968283in}{0.584267in}}%
\pgfpathlineto{\pgfqpoint{3.970380in}{0.596618in}}%
\pgfpathlineto{\pgfqpoint{3.972478in}{0.602633in}}%
\pgfpathlineto{\pgfqpoint{3.974575in}{0.604355in}}%
\pgfpathlineto{\pgfqpoint{3.976672in}{0.602216in}}%
\pgfpathlineto{\pgfqpoint{3.978770in}{0.595587in}}%
\pgfpathlineto{\pgfqpoint{3.980867in}{0.581925in}}%
\pgfpathlineto{\pgfqpoint{3.982964in}{0.550014in}}%
\pgfpathlineto{\pgfqpoint{3.985062in}{0.528546in}}%
\pgfpathlineto{\pgfqpoint{3.987159in}{0.580009in}}%
\pgfpathlineto{\pgfqpoint{3.989256in}{0.600446in}}%
\pgfpathlineto{\pgfqpoint{3.991354in}{0.612044in}}%
\pgfpathlineto{\pgfqpoint{3.993451in}{0.618422in}}%
\pgfpathlineto{\pgfqpoint{3.995549in}{0.620313in}}%
\pgfpathlineto{\pgfqpoint{3.997646in}{0.616804in}}%
\pgfpathlineto{\pgfqpoint{3.999743in}{0.605686in}}%
\pgfpathlineto{\pgfqpoint{4.001841in}{0.581101in}}%
\pgfpathlineto{\pgfqpoint{4.003938in}{0.473892in}}%
\pgfpathlineto{\pgfqpoint{4.006035in}{0.585645in}}%
\pgfpathlineto{\pgfqpoint{4.008133in}{0.609220in}}%
\pgfpathlineto{\pgfqpoint{4.010230in}{0.621572in}}%
\pgfpathlineto{\pgfqpoint{4.012327in}{0.628551in}}%
\pgfpathlineto{\pgfqpoint{4.014425in}{0.631926in}}%
\pgfpathlineto{\pgfqpoint{4.016522in}{0.632298in}}%
\pgfpathlineto{\pgfqpoint{4.018620in}{0.629661in}}%
\pgfpathlineto{\pgfqpoint{4.020717in}{0.623382in}}%
\pgfpathlineto{\pgfqpoint{4.022814in}{0.611566in}}%
\pgfpathlineto{\pgfqpoint{4.024912in}{0.587604in}}%
\pgfpathlineto{\pgfqpoint{4.027009in}{0.464981in}}%
\pgfpathlineto{\pgfqpoint{4.029106in}{0.592912in}}%
\pgfpathlineto{\pgfqpoint{4.031204in}{0.618194in}}%
\pgfpathlineto{\pgfqpoint{4.033301in}{0.632040in}}%
\pgfpathlineto{\pgfqpoint{4.035398in}{0.640262in}}%
\pgfpathlineto{\pgfqpoint{4.037496in}{0.644357in}}%
\pgfpathlineto{\pgfqpoint{4.039593in}{0.644353in}}%
\pgfpathlineto{\pgfqpoint{4.041691in}{0.639072in}}%
\pgfpathlineto{\pgfqpoint{4.043788in}{0.629376in}}%
\pgfpathlineto{\pgfqpoint{4.045885in}{0.614672in}}%
\pgfpathlineto{\pgfqpoint{4.047983in}{0.590526in}}%
\pgfpathlineto{\pgfqpoint{4.050080in}{0.526735in}}%
\pgfpathlineto{\pgfqpoint{4.052177in}{0.563653in}}%
\pgfpathlineto{\pgfqpoint{4.054275in}{0.589246in}}%
\pgfpathlineto{\pgfqpoint{4.056372in}{0.598140in}}%
\pgfpathlineto{\pgfqpoint{4.058469in}{0.598998in}}%
\pgfpathlineto{\pgfqpoint{4.060567in}{0.592590in}}%
\pgfpathlineto{\pgfqpoint{4.062664in}{0.574702in}}%
\pgfpathlineto{\pgfqpoint{4.064762in}{0.499505in}}%
\pgfpathlineto{\pgfqpoint{4.066859in}{0.570849in}}%
\pgfpathlineto{\pgfqpoint{4.068956in}{0.600703in}}%
\pgfpathlineto{\pgfqpoint{4.071054in}{0.616255in}}%
\pgfpathlineto{\pgfqpoint{4.073151in}{0.625057in}}%
\pgfpathlineto{\pgfqpoint{4.075248in}{0.628716in}}%
\pgfpathlineto{\pgfqpoint{4.077346in}{0.626681in}}%
\pgfpathlineto{\pgfqpoint{4.079443in}{0.618669in}}%
\pgfpathlineto{\pgfqpoint{4.081540in}{0.606394in}}%
\pgfpathlineto{\pgfqpoint{4.083638in}{0.587281in}}%
\pgfpathlineto{\pgfqpoint{4.087832in}{0.518683in}}%
\pgfpathlineto{\pgfqpoint{4.089930in}{0.566160in}}%
\pgfpathlineto{\pgfqpoint{4.092027in}{0.580132in}}%
\pgfpathlineto{\pgfqpoint{4.094125in}{0.584105in}}%
\pgfpathlineto{\pgfqpoint{4.096222in}{0.580970in}}%
\pgfpathlineto{\pgfqpoint{4.098319in}{0.569100in}}%
\pgfpathlineto{\pgfqpoint{4.100417in}{0.535729in}}%
\pgfpathlineto{\pgfqpoint{4.102514in}{0.528777in}}%
\pgfpathlineto{\pgfqpoint{4.104611in}{0.575254in}}%
\pgfpathlineto{\pgfqpoint{4.106709in}{0.595233in}}%
\pgfpathlineto{\pgfqpoint{4.108806in}{0.606570in}}%
\pgfpathlineto{\pgfqpoint{4.110903in}{0.612393in}}%
\pgfpathlineto{\pgfqpoint{4.113001in}{0.613010in}}%
\pgfpathlineto{\pgfqpoint{4.115098in}{0.606987in}}%
\pgfpathlineto{\pgfqpoint{4.117196in}{0.594342in}}%
\pgfpathlineto{\pgfqpoint{4.119293in}{0.570269in}}%
\pgfpathlineto{\pgfqpoint{4.121390in}{0.468236in}}%
\pgfpathlineto{\pgfqpoint{4.123488in}{0.562500in}}%
\pgfpathlineto{\pgfqpoint{4.125585in}{0.586977in}}%
\pgfpathlineto{\pgfqpoint{4.127682in}{0.598527in}}%
\pgfpathlineto{\pgfqpoint{4.129780in}{0.604064in}}%
\pgfpathlineto{\pgfqpoint{4.131877in}{0.605398in}}%
\pgfpathlineto{\pgfqpoint{4.133974in}{0.602854in}}%
\pgfpathlineto{\pgfqpoint{4.136072in}{0.595677in}}%
\pgfpathlineto{\pgfqpoint{4.138169in}{0.580972in}}%
\pgfpathlineto{\pgfqpoint{4.140267in}{0.544929in}}%
\pgfpathlineto{\pgfqpoint{4.142364in}{0.539440in}}%
\pgfpathlineto{\pgfqpoint{4.144461in}{0.583880in}}%
\pgfpathlineto{\pgfqpoint{4.146559in}{0.603045in}}%
\pgfpathlineto{\pgfqpoint{4.148656in}{0.614016in}}%
\pgfpathlineto{\pgfqpoint{4.150753in}{0.619926in}}%
\pgfpathlineto{\pgfqpoint{4.152851in}{0.621332in}}%
\pgfpathlineto{\pgfqpoint{4.154948in}{0.617137in}}%
\pgfpathlineto{\pgfqpoint{4.157045in}{0.605229in}}%
\pgfpathlineto{\pgfqpoint{4.159143in}{0.579064in}}%
\pgfpathlineto{\pgfqpoint{4.161240in}{0.502877in}}%
\pgfpathlineto{\pgfqpoint{4.163338in}{0.587785in}}%
\pgfpathlineto{\pgfqpoint{4.165435in}{0.610152in}}%
\pgfpathlineto{\pgfqpoint{4.167532in}{0.621946in}}%
\pgfpathlineto{\pgfqpoint{4.169630in}{0.628521in}}%
\pgfpathlineto{\pgfqpoint{4.171727in}{0.631522in}}%
\pgfpathlineto{\pgfqpoint{4.173824in}{0.631486in}}%
\pgfpathlineto{\pgfqpoint{4.175922in}{0.628330in}}%
\pgfpathlineto{\pgfqpoint{4.178019in}{0.621274in}}%
\pgfpathlineto{\pgfqpoint{4.180116in}{0.607987in}}%
\pgfpathlineto{\pgfqpoint{4.182214in}{0.579436in}}%
\pgfpathlineto{\pgfqpoint{4.184311in}{0.535167in}}%
\pgfpathlineto{\pgfqpoint{4.186409in}{0.599126in}}%
\pgfpathlineto{\pgfqpoint{4.188506in}{0.621213in}}%
\pgfpathlineto{\pgfqpoint{4.190603in}{0.633731in}}%
\pgfpathlineto{\pgfqpoint{4.192701in}{0.641060in}}%
\pgfpathlineto{\pgfqpoint{4.194798in}{0.644322in}}%
\pgfpathlineto{\pgfqpoint{4.196895in}{0.643275in}}%
\pgfpathlineto{\pgfqpoint{4.198993in}{0.636783in}}%
\pgfpathlineto{\pgfqpoint{4.201090in}{0.626045in}}%
\pgfpathlineto{\pgfqpoint{4.203187in}{0.609585in}}%
\pgfpathlineto{\pgfqpoint{4.205285in}{0.580928in}}%
\pgfpathlineto{\pgfqpoint{4.207382in}{0.430027in}}%
\pgfpathlineto{\pgfqpoint{4.209480in}{0.573980in}}%
\pgfpathlineto{\pgfqpoint{4.211577in}{0.593173in}}%
\pgfpathlineto{\pgfqpoint{4.213674in}{0.599786in}}%
\pgfpathlineto{\pgfqpoint{4.215772in}{0.599087in}}%
\pgfpathlineto{\pgfqpoint{4.217869in}{0.590925in}}%
\pgfpathlineto{\pgfqpoint{4.219966in}{0.569276in}}%
\pgfpathlineto{\pgfqpoint{4.222064in}{0.462782in}}%
\pgfpathlineto{\pgfqpoint{4.224161in}{0.578843in}}%
\pgfpathlineto{\pgfqpoint{4.226258in}{0.604605in}}%
\pgfpathlineto{\pgfqpoint{4.228356in}{0.618592in}}%
\pgfpathlineto{\pgfqpoint{4.230453in}{0.626364in}}%
\pgfpathlineto{\pgfqpoint{4.232551in}{0.629050in}}%
\pgfpathlineto{\pgfqpoint{4.234648in}{0.625714in}}%
\pgfpathlineto{\pgfqpoint{4.236745in}{0.616954in}}%
\pgfpathlineto{\pgfqpoint{4.238843in}{0.603599in}}%
\pgfpathlineto{\pgfqpoint{4.240940in}{0.582271in}}%
\pgfpathlineto{\pgfqpoint{4.243037in}{0.535977in}}%
\pgfpathlineto{\pgfqpoint{4.245135in}{0.538012in}}%
\pgfpathlineto{\pgfqpoint{4.247232in}{0.571572in}}%
\pgfpathlineto{\pgfqpoint{4.249329in}{0.582897in}}%
\pgfpathlineto{\pgfqpoint{4.251427in}{0.585484in}}%
\pgfpathlineto{\pgfqpoint{4.253524in}{0.581141in}}%
\pgfpathlineto{\pgfqpoint{4.255621in}{0.567456in}}%
\pgfpathlineto{\pgfqpoint{4.257719in}{0.526682in}}%
\pgfpathlineto{\pgfqpoint{4.259816in}{0.540494in}}%
\pgfpathlineto{\pgfqpoint{4.261914in}{0.579746in}}%
\pgfpathlineto{\pgfqpoint{4.264011in}{0.598080in}}%
\pgfpathlineto{\pgfqpoint{4.266108in}{0.608529in}}%
\pgfpathlineto{\pgfqpoint{4.268206in}{0.613636in}}%
\pgfpathlineto{\pgfqpoint{4.270303in}{0.613439in}}%
\pgfpathlineto{\pgfqpoint{4.272400in}{0.606493in}}%
\pgfpathlineto{\pgfqpoint{4.274498in}{0.592831in}}%
\pgfpathlineto{\pgfqpoint{4.276595in}{0.566278in}}%
\pgfpathlineto{\pgfqpoint{4.278692in}{0.456191in}}%
\pgfpathlineto{\pgfqpoint{4.280790in}{0.567355in}}%
\pgfpathlineto{\pgfqpoint{4.282887in}{0.589307in}}%
\pgfpathlineto{\pgfqpoint{4.284985in}{0.599877in}}%
\pgfpathlineto{\pgfqpoint{4.287082in}{0.604782in}}%
\pgfpathlineto{\pgfqpoint{4.289179in}{0.605563in}}%
\pgfpathlineto{\pgfqpoint{4.291277in}{0.602401in}}%
\pgfpathlineto{\pgfqpoint{4.293374in}{0.594329in}}%
\pgfpathlineto{\pgfqpoint{4.295471in}{0.577782in}}%
\pgfpathlineto{\pgfqpoint{4.297569in}{0.532625in}}%
\pgfpathlineto{\pgfqpoint{4.299666in}{0.551531in}}%
\pgfpathlineto{\pgfqpoint{4.301763in}{0.588288in}}%
\pgfpathlineto{\pgfqpoint{4.303861in}{0.605687in}}%
\pgfpathlineto{\pgfqpoint{4.305958in}{0.615722in}}%
\pgfpathlineto{\pgfqpoint{4.308056in}{0.620895in}}%
\pgfpathlineto{\pgfqpoint{4.310153in}{0.621494in}}%
\pgfpathlineto{\pgfqpoint{4.312250in}{0.616519in}}%
\pgfpathlineto{\pgfqpoint{4.316445in}{0.603256in}}%
\pgfpathlineto{\pgfqpoint{4.316445in}{0.603256in}}%
\pgfusepath{stroke}%
\end{pgfscope}%
\begin{pgfscope}%
\pgfpathrectangle{\pgfqpoint{2.554660in}{0.320625in}}{\pgfqpoint{1.845679in}{1.164375in}}%
\pgfusepath{clip}%
\pgfsetrectcap%
\pgfsetroundjoin%
\pgfsetlinewidth{1.154312pt}%
\definecolor{currentstroke}{rgb}{0.989363,0.557873,0.391671}%
\pgfsetstrokecolor{currentstroke}%
\pgfsetdash{}{0pt}%
\pgfpathmoveto{\pgfqpoint{2.638555in}{1.468525in}}%
\pgfpathlineto{\pgfqpoint{2.649042in}{1.449217in}}%
\pgfpathlineto{\pgfqpoint{2.659529in}{1.427093in}}%
\pgfpathlineto{\pgfqpoint{2.672113in}{1.397448in}}%
\pgfpathlineto{\pgfqpoint{2.716157in}{1.291455in}}%
\pgfpathlineto{\pgfqpoint{2.739228in}{1.239960in}}%
\pgfpathlineto{\pgfqpoint{2.764397in}{1.186560in}}%
\pgfpathlineto{\pgfqpoint{2.795857in}{1.122540in}}%
\pgfpathlineto{\pgfqpoint{2.835707in}{1.044156in}}%
\pgfpathlineto{\pgfqpoint{2.886044in}{0.947761in}}%
\pgfpathlineto{\pgfqpoint{2.948965in}{0.829875in}}%
\pgfpathlineto{\pgfqpoint{3.030762in}{0.679922in}}%
\pgfpathlineto{\pgfqpoint{3.043346in}{0.660579in}}%
\pgfpathlineto{\pgfqpoint{3.049638in}{0.647198in}}%
\pgfpathlineto{\pgfqpoint{3.064320in}{0.613034in}}%
\pgfpathlineto{\pgfqpoint{3.068514in}{0.606861in}}%
\pgfpathlineto{\pgfqpoint{3.070612in}{0.605072in}}%
\pgfpathlineto{\pgfqpoint{3.072709in}{0.604158in}}%
\pgfpathlineto{\pgfqpoint{3.076904in}{0.604339in}}%
\pgfpathlineto{\pgfqpoint{3.081098in}{0.605317in}}%
\pgfpathlineto{\pgfqpoint{3.083196in}{0.605215in}}%
\pgfpathlineto{\pgfqpoint{3.085293in}{0.604151in}}%
\pgfpathlineto{\pgfqpoint{3.087391in}{0.601536in}}%
\pgfpathlineto{\pgfqpoint{3.091585in}{0.591627in}}%
\pgfpathlineto{\pgfqpoint{3.095780in}{0.576075in}}%
\pgfpathlineto{\pgfqpoint{3.102072in}{0.544057in}}%
\pgfpathlineto{\pgfqpoint{3.104169in}{0.543118in}}%
\pgfpathlineto{\pgfqpoint{3.110462in}{0.561802in}}%
\pgfpathlineto{\pgfqpoint{3.112559in}{0.564976in}}%
\pgfpathlineto{\pgfqpoint{3.114656in}{0.566206in}}%
\pgfpathlineto{\pgfqpoint{3.116754in}{0.565530in}}%
\pgfpathlineto{\pgfqpoint{3.118851in}{0.562809in}}%
\pgfpathlineto{\pgfqpoint{3.120948in}{0.557620in}}%
\pgfpathlineto{\pgfqpoint{3.123046in}{0.548984in}}%
\pgfpathlineto{\pgfqpoint{3.125143in}{0.534502in}}%
\pgfpathlineto{\pgfqpoint{3.127240in}{0.507285in}}%
\pgfpathlineto{\pgfqpoint{3.129338in}{0.498422in}}%
\pgfpathlineto{\pgfqpoint{3.131435in}{0.533483in}}%
\pgfpathlineto{\pgfqpoint{3.133533in}{0.553377in}}%
\pgfpathlineto{\pgfqpoint{3.135630in}{0.566352in}}%
\pgfpathlineto{\pgfqpoint{3.137727in}{0.575531in}}%
\pgfpathlineto{\pgfqpoint{3.139825in}{0.582152in}}%
\pgfpathlineto{\pgfqpoint{3.141922in}{0.586768in}}%
\pgfpathlineto{\pgfqpoint{3.144019in}{0.589597in}}%
\pgfpathlineto{\pgfqpoint{3.146117in}{0.590625in}}%
\pgfpathlineto{\pgfqpoint{3.148214in}{0.589590in}}%
\pgfpathlineto{\pgfqpoint{3.150311in}{0.585576in}}%
\pgfpathlineto{\pgfqpoint{3.152409in}{0.576323in}}%
\pgfpathlineto{\pgfqpoint{3.154506in}{0.577521in}}%
\pgfpathlineto{\pgfqpoint{3.158701in}{0.616553in}}%
\pgfpathlineto{\pgfqpoint{3.162896in}{0.641299in}}%
\pgfpathlineto{\pgfqpoint{3.167090in}{0.656445in}}%
\pgfpathlineto{\pgfqpoint{3.171285in}{0.666082in}}%
\pgfpathlineto{\pgfqpoint{3.175480in}{0.671901in}}%
\pgfpathlineto{\pgfqpoint{3.177577in}{0.673631in}}%
\pgfpathlineto{\pgfqpoint{3.179675in}{0.674636in}}%
\pgfpathlineto{\pgfqpoint{3.181772in}{0.674930in}}%
\pgfpathlineto{\pgfqpoint{3.183869in}{0.674498in}}%
\pgfpathlineto{\pgfqpoint{3.185967in}{0.673297in}}%
\pgfpathlineto{\pgfqpoint{3.188064in}{0.671249in}}%
\pgfpathlineto{\pgfqpoint{3.190161in}{0.668218in}}%
\pgfpathlineto{\pgfqpoint{3.192259in}{0.663987in}}%
\pgfpathlineto{\pgfqpoint{3.194356in}{0.658188in}}%
\pgfpathlineto{\pgfqpoint{3.196453in}{0.650155in}}%
\pgfpathlineto{\pgfqpoint{3.198551in}{0.638525in}}%
\pgfpathlineto{\pgfqpoint{3.200648in}{0.619815in}}%
\pgfpathlineto{\pgfqpoint{3.202746in}{0.582224in}}%
\pgfpathlineto{\pgfqpoint{3.204843in}{0.597297in}}%
\pgfpathlineto{\pgfqpoint{3.206940in}{0.628780in}}%
\pgfpathlineto{\pgfqpoint{3.209038in}{0.646266in}}%
\pgfpathlineto{\pgfqpoint{3.213232in}{0.666717in}}%
\pgfpathlineto{\pgfqpoint{3.217427in}{0.678368in}}%
\pgfpathlineto{\pgfqpoint{3.221622in}{0.684846in}}%
\pgfpathlineto{\pgfqpoint{3.223719in}{0.686448in}}%
\pgfpathlineto{\pgfqpoint{3.225816in}{0.686931in}}%
\pgfpathlineto{\pgfqpoint{3.227914in}{0.686155in}}%
\pgfpathlineto{\pgfqpoint{3.230011in}{0.683833in}}%
\pgfpathlineto{\pgfqpoint{3.232109in}{0.679414in}}%
\pgfpathlineto{\pgfqpoint{3.236303in}{0.666606in}}%
\pgfpathlineto{\pgfqpoint{3.246790in}{0.632473in}}%
\pgfpathlineto{\pgfqpoint{3.248887in}{0.627951in}}%
\pgfpathlineto{\pgfqpoint{3.250985in}{0.624870in}}%
\pgfpathlineto{\pgfqpoint{3.253082in}{0.623549in}}%
\pgfpathlineto{\pgfqpoint{3.255180in}{0.623924in}}%
\pgfpathlineto{\pgfqpoint{3.259374in}{0.625714in}}%
\pgfpathlineto{\pgfqpoint{3.261472in}{0.623839in}}%
\pgfpathlineto{\pgfqpoint{3.265666in}{0.611386in}}%
\pgfpathlineto{\pgfqpoint{3.276153in}{0.574575in}}%
\pgfpathlineto{\pgfqpoint{3.278251in}{0.569588in}}%
\pgfpathlineto{\pgfqpoint{3.280348in}{0.566186in}}%
\pgfpathlineto{\pgfqpoint{3.282445in}{0.564780in}}%
\pgfpathlineto{\pgfqpoint{3.284543in}{0.565509in}}%
\pgfpathlineto{\pgfqpoint{3.288737in}{0.569401in}}%
\pgfpathlineto{\pgfqpoint{3.290835in}{0.569155in}}%
\pgfpathlineto{\pgfqpoint{3.292932in}{0.564477in}}%
\pgfpathlineto{\pgfqpoint{3.295029in}{0.552018in}}%
\pgfpathlineto{\pgfqpoint{3.297127in}{0.530933in}}%
\pgfpathlineto{\pgfqpoint{3.299224in}{0.538195in}}%
\pgfpathlineto{\pgfqpoint{3.301322in}{0.561815in}}%
\pgfpathlineto{\pgfqpoint{3.303419in}{0.577104in}}%
\pgfpathlineto{\pgfqpoint{3.305516in}{0.586912in}}%
\pgfpathlineto{\pgfqpoint{3.307614in}{0.593142in}}%
\pgfpathlineto{\pgfqpoint{3.309711in}{0.596685in}}%
\pgfpathlineto{\pgfqpoint{3.311808in}{0.597927in}}%
\pgfpathlineto{\pgfqpoint{3.313906in}{0.596939in}}%
\pgfpathlineto{\pgfqpoint{3.316003in}{0.593495in}}%
\pgfpathlineto{\pgfqpoint{3.318100in}{0.586921in}}%
\pgfpathlineto{\pgfqpoint{3.320198in}{0.575535in}}%
\pgfpathlineto{\pgfqpoint{3.322295in}{0.554118in}}%
\pgfpathlineto{\pgfqpoint{3.324393in}{0.481660in}}%
\pgfpathlineto{\pgfqpoint{3.326490in}{0.547337in}}%
\pgfpathlineto{\pgfqpoint{3.328587in}{0.575846in}}%
\pgfpathlineto{\pgfqpoint{3.330685in}{0.591123in}}%
\pgfpathlineto{\pgfqpoint{3.332782in}{0.600596in}}%
\pgfpathlineto{\pgfqpoint{3.334879in}{0.606251in}}%
\pgfpathlineto{\pgfqpoint{3.336977in}{0.608618in}}%
\pgfpathlineto{\pgfqpoint{3.339074in}{0.607372in}}%
\pgfpathlineto{\pgfqpoint{3.341171in}{0.601057in}}%
\pgfpathlineto{\pgfqpoint{3.343269in}{0.587332in}}%
\pgfpathlineto{\pgfqpoint{3.347464in}{0.529210in}}%
\pgfpathlineto{\pgfqpoint{3.349561in}{0.575625in}}%
\pgfpathlineto{\pgfqpoint{3.351658in}{0.594902in}}%
\pgfpathlineto{\pgfqpoint{3.353756in}{0.605264in}}%
\pgfpathlineto{\pgfqpoint{3.355853in}{0.610783in}}%
\pgfpathlineto{\pgfqpoint{3.357950in}{0.612788in}}%
\pgfpathlineto{\pgfqpoint{3.360048in}{0.611628in}}%
\pgfpathlineto{\pgfqpoint{3.362145in}{0.606986in}}%
\pgfpathlineto{\pgfqpoint{3.364242in}{0.597568in}}%
\pgfpathlineto{\pgfqpoint{3.366340in}{0.579226in}}%
\pgfpathlineto{\pgfqpoint{3.368437in}{0.525799in}}%
\pgfpathlineto{\pgfqpoint{3.372632in}{0.593527in}}%
\pgfpathlineto{\pgfqpoint{3.374729in}{0.610010in}}%
\pgfpathlineto{\pgfqpoint{3.376827in}{0.619724in}}%
\pgfpathlineto{\pgfqpoint{3.378924in}{0.624943in}}%
\pgfpathlineto{\pgfqpoint{3.381021in}{0.626031in}}%
\pgfpathlineto{\pgfqpoint{3.383119in}{0.621946in}}%
\pgfpathlineto{\pgfqpoint{3.385216in}{0.610816in}}%
\pgfpathlineto{\pgfqpoint{3.387313in}{0.587777in}}%
\pgfpathlineto{\pgfqpoint{3.389411in}{0.514528in}}%
\pgfpathlineto{\pgfqpoint{3.391508in}{0.584591in}}%
\pgfpathlineto{\pgfqpoint{3.393605in}{0.609261in}}%
\pgfpathlineto{\pgfqpoint{3.395703in}{0.621779in}}%
\pgfpathlineto{\pgfqpoint{3.397800in}{0.628627in}}%
\pgfpathlineto{\pgfqpoint{3.399898in}{0.631688in}}%
\pgfpathlineto{\pgfqpoint{3.401995in}{0.631557in}}%
\pgfpathlineto{\pgfqpoint{3.404092in}{0.628138in}}%
\pgfpathlineto{\pgfqpoint{3.406190in}{0.620543in}}%
\pgfpathlineto{\pgfqpoint{3.408287in}{0.606029in}}%
\pgfpathlineto{\pgfqpoint{3.410384in}{0.572727in}}%
\pgfpathlineto{\pgfqpoint{3.412482in}{0.555252in}}%
\pgfpathlineto{\pgfqpoint{3.414579in}{0.604502in}}%
\pgfpathlineto{\pgfqpoint{3.416676in}{0.624571in}}%
\pgfpathlineto{\pgfqpoint{3.418774in}{0.636131in}}%
\pgfpathlineto{\pgfqpoint{3.420871in}{0.642736in}}%
\pgfpathlineto{\pgfqpoint{3.422969in}{0.645228in}}%
\pgfpathlineto{\pgfqpoint{3.425066in}{0.643078in}}%
\pgfpathlineto{\pgfqpoint{3.427163in}{0.634885in}}%
\pgfpathlineto{\pgfqpoint{3.429261in}{0.619580in}}%
\pgfpathlineto{\pgfqpoint{3.433455in}{0.548829in}}%
\pgfpathlineto{\pgfqpoint{3.435553in}{0.604162in}}%
\pgfpathlineto{\pgfqpoint{3.437650in}{0.622821in}}%
\pgfpathlineto{\pgfqpoint{3.439747in}{0.632074in}}%
\pgfpathlineto{\pgfqpoint{3.441845in}{0.636168in}}%
\pgfpathlineto{\pgfqpoint{3.443942in}{0.636296in}}%
\pgfpathlineto{\pgfqpoint{3.446040in}{0.632454in}}%
\pgfpathlineto{\pgfqpoint{3.448137in}{0.623443in}}%
\pgfpathlineto{\pgfqpoint{3.450234in}{0.604927in}}%
\pgfpathlineto{\pgfqpoint{3.452332in}{0.546342in}}%
\pgfpathlineto{\pgfqpoint{3.454429in}{0.590327in}}%
\pgfpathlineto{\pgfqpoint{3.456526in}{0.622218in}}%
\pgfpathlineto{\pgfqpoint{3.458624in}{0.638238in}}%
\pgfpathlineto{\pgfqpoint{3.460721in}{0.647507in}}%
\pgfpathlineto{\pgfqpoint{3.462818in}{0.652071in}}%
\pgfpathlineto{\pgfqpoint{3.464916in}{0.652007in}}%
\pgfpathlineto{\pgfqpoint{3.467013in}{0.646121in}}%
\pgfpathlineto{\pgfqpoint{3.469111in}{0.635802in}}%
\pgfpathlineto{\pgfqpoint{3.471208in}{0.619988in}}%
\pgfpathlineto{\pgfqpoint{3.473305in}{0.592967in}}%
\pgfpathlineto{\pgfqpoint{3.475403in}{0.504467in}}%
\pgfpathlineto{\pgfqpoint{3.477500in}{0.580835in}}%
\pgfpathlineto{\pgfqpoint{3.479597in}{0.601577in}}%
\pgfpathlineto{\pgfqpoint{3.481695in}{0.608980in}}%
\pgfpathlineto{\pgfqpoint{3.483792in}{0.609032in}}%
\pgfpathlineto{\pgfqpoint{3.485889in}{0.601987in}}%
\pgfpathlineto{\pgfqpoint{3.487987in}{0.583197in}}%
\pgfpathlineto{\pgfqpoint{3.490084in}{0.497963in}}%
\pgfpathlineto{\pgfqpoint{3.492182in}{0.582033in}}%
\pgfpathlineto{\pgfqpoint{3.494279in}{0.610833in}}%
\pgfpathlineto{\pgfqpoint{3.496376in}{0.625984in}}%
\pgfpathlineto{\pgfqpoint{3.498474in}{0.634533in}}%
\pgfpathlineto{\pgfqpoint{3.500571in}{0.637977in}}%
\pgfpathlineto{\pgfqpoint{3.502668in}{0.635699in}}%
\pgfpathlineto{\pgfqpoint{3.504766in}{0.627696in}}%
\pgfpathlineto{\pgfqpoint{3.506863in}{0.615870in}}%
\pgfpathlineto{\pgfqpoint{3.508960in}{0.598214in}}%
\pgfpathlineto{\pgfqpoint{3.511058in}{0.567715in}}%
\pgfpathlineto{\pgfqpoint{3.513155in}{0.480158in}}%
\pgfpathlineto{\pgfqpoint{3.515253in}{0.562446in}}%
\pgfpathlineto{\pgfqpoint{3.517350in}{0.579153in}}%
\pgfpathlineto{\pgfqpoint{3.519447in}{0.583262in}}%
\pgfpathlineto{\pgfqpoint{3.521545in}{0.578756in}}%
\pgfpathlineto{\pgfqpoint{3.523642in}{0.562440in}}%
\pgfpathlineto{\pgfqpoint{3.525739in}{0.495149in}}%
\pgfpathlineto{\pgfqpoint{3.527837in}{0.557535in}}%
\pgfpathlineto{\pgfqpoint{3.529934in}{0.588548in}}%
\pgfpathlineto{\pgfqpoint{3.532031in}{0.604668in}}%
\pgfpathlineto{\pgfqpoint{3.534129in}{0.613869in}}%
\pgfpathlineto{\pgfqpoint{3.536226in}{0.617851in}}%
\pgfpathlineto{\pgfqpoint{3.538324in}{0.616116in}}%
\pgfpathlineto{\pgfqpoint{3.540421in}{0.607855in}}%
\pgfpathlineto{\pgfqpoint{3.542518in}{0.593813in}}%
\pgfpathlineto{\pgfqpoint{3.544616in}{0.568580in}}%
\pgfpathlineto{\pgfqpoint{3.546713in}{0.470473in}}%
\pgfpathlineto{\pgfqpoint{3.548810in}{0.558411in}}%
\pgfpathlineto{\pgfqpoint{3.550908in}{0.581681in}}%
\pgfpathlineto{\pgfqpoint{3.553005in}{0.591818in}}%
\pgfpathlineto{\pgfqpoint{3.555102in}{0.595583in}}%
\pgfpathlineto{\pgfqpoint{3.557200in}{0.594488in}}%
\pgfpathlineto{\pgfqpoint{3.559297in}{0.588179in}}%
\pgfpathlineto{\pgfqpoint{3.561395in}{0.573826in}}%
\pgfpathlineto{\pgfqpoint{3.563492in}{0.536561in}}%
\pgfpathlineto{\pgfqpoint{3.565589in}{0.537032in}}%
\pgfpathlineto{\pgfqpoint{3.567687in}{0.579648in}}%
\pgfpathlineto{\pgfqpoint{3.569784in}{0.598501in}}%
\pgfpathlineto{\pgfqpoint{3.571881in}{0.609227in}}%
\pgfpathlineto{\pgfqpoint{3.573979in}{0.614744in}}%
\pgfpathlineto{\pgfqpoint{3.576076in}{0.615394in}}%
\pgfpathlineto{\pgfqpoint{3.578173in}{0.609732in}}%
\pgfpathlineto{\pgfqpoint{3.580271in}{0.596188in}}%
\pgfpathlineto{\pgfqpoint{3.582368in}{0.566611in}}%
\pgfpathlineto{\pgfqpoint{3.584465in}{0.521022in}}%
\pgfpathlineto{\pgfqpoint{3.586563in}{0.583087in}}%
\pgfpathlineto{\pgfqpoint{3.588660in}{0.603312in}}%
\pgfpathlineto{\pgfqpoint{3.590758in}{0.613997in}}%
\pgfpathlineto{\pgfqpoint{3.592855in}{0.619697in}}%
\pgfpathlineto{\pgfqpoint{3.594952in}{0.621833in}}%
\pgfpathlineto{\pgfqpoint{3.597050in}{0.620789in}}%
\pgfpathlineto{\pgfqpoint{3.599147in}{0.616276in}}%
\pgfpathlineto{\pgfqpoint{3.601244in}{0.607039in}}%
\pgfpathlineto{\pgfqpoint{3.603342in}{0.589083in}}%
\pgfpathlineto{\pgfqpoint{3.605439in}{0.538196in}}%
\pgfpathlineto{\pgfqpoint{3.609634in}{0.601994in}}%
\pgfpathlineto{\pgfqpoint{3.611731in}{0.618735in}}%
\pgfpathlineto{\pgfqpoint{3.613829in}{0.628575in}}%
\pgfpathlineto{\pgfqpoint{3.615926in}{0.633870in}}%
\pgfpathlineto{\pgfqpoint{3.618023in}{0.635010in}}%
\pgfpathlineto{\pgfqpoint{3.620121in}{0.630963in}}%
\pgfpathlineto{\pgfqpoint{3.622218in}{0.619899in}}%
\pgfpathlineto{\pgfqpoint{3.624315in}{0.597050in}}%
\pgfpathlineto{\pgfqpoint{3.626413in}{0.495668in}}%
\pgfpathlineto{\pgfqpoint{3.628510in}{0.591576in}}%
\pgfpathlineto{\pgfqpoint{3.630607in}{0.617064in}}%
\pgfpathlineto{\pgfqpoint{3.632705in}{0.629752in}}%
\pgfpathlineto{\pgfqpoint{3.634802in}{0.636628in}}%
\pgfpathlineto{\pgfqpoint{3.636900in}{0.639647in}}%
\pgfpathlineto{\pgfqpoint{3.638997in}{0.639416in}}%
\pgfpathlineto{\pgfqpoint{3.641094in}{0.635818in}}%
\pgfpathlineto{\pgfqpoint{3.643192in}{0.627898in}}%
\pgfpathlineto{\pgfqpoint{3.645289in}{0.612657in}}%
\pgfpathlineto{\pgfqpoint{3.647386in}{0.576214in}}%
\pgfpathlineto{\pgfqpoint{3.649484in}{0.570549in}}%
\pgfpathlineto{\pgfqpoint{3.651581in}{0.614819in}}%
\pgfpathlineto{\pgfqpoint{3.653678in}{0.633954in}}%
\pgfpathlineto{\pgfqpoint{3.655776in}{0.645039in}}%
\pgfpathlineto{\pgfqpoint{3.657873in}{0.651265in}}%
\pgfpathlineto{\pgfqpoint{3.659971in}{0.653339in}}%
\pgfpathlineto{\pgfqpoint{3.662068in}{0.650565in}}%
\pgfpathlineto{\pgfqpoint{3.664165in}{0.642165in}}%
\pgfpathlineto{\pgfqpoint{3.666263in}{0.628257in}}%
\pgfpathlineto{\pgfqpoint{3.668360in}{0.603881in}}%
\pgfpathlineto{\pgfqpoint{3.670457in}{0.523962in}}%
\pgfpathlineto{\pgfqpoint{3.672555in}{0.588033in}}%
\pgfpathlineto{\pgfqpoint{3.674652in}{0.612578in}}%
\pgfpathlineto{\pgfqpoint{3.676749in}{0.622861in}}%
\pgfpathlineto{\pgfqpoint{3.678847in}{0.626423in}}%
\pgfpathlineto{\pgfqpoint{3.680944in}{0.624834in}}%
\pgfpathlineto{\pgfqpoint{3.683042in}{0.617538in}}%
\pgfpathlineto{\pgfqpoint{3.685139in}{0.600778in}}%
\pgfpathlineto{\pgfqpoint{3.687236in}{0.549090in}}%
\pgfpathlineto{\pgfqpoint{3.691431in}{0.616797in}}%
\pgfpathlineto{\pgfqpoint{3.693528in}{0.633481in}}%
\pgfpathlineto{\pgfqpoint{3.695626in}{0.642994in}}%
\pgfpathlineto{\pgfqpoint{3.697723in}{0.647465in}}%
\pgfpathlineto{\pgfqpoint{3.699820in}{0.646817in}}%
\pgfpathlineto{\pgfqpoint{3.701918in}{0.640121in}}%
\pgfpathlineto{\pgfqpoint{3.704015in}{0.629745in}}%
\pgfpathlineto{\pgfqpoint{3.706113in}{0.614747in}}%
\pgfpathlineto{\pgfqpoint{3.708210in}{0.591347in}}%
\pgfpathlineto{\pgfqpoint{3.710307in}{0.539416in}}%
\pgfpathlineto{\pgfqpoint{3.712405in}{0.550162in}}%
\pgfpathlineto{\pgfqpoint{3.714502in}{0.578780in}}%
\pgfpathlineto{\pgfqpoint{3.716599in}{0.586503in}}%
\pgfpathlineto{\pgfqpoint{3.718697in}{0.584041in}}%
\pgfpathlineto{\pgfqpoint{3.720794in}{0.569764in}}%
\pgfpathlineto{\pgfqpoint{3.722891in}{0.514130in}}%
\pgfpathlineto{\pgfqpoint{3.724989in}{0.560134in}}%
\pgfpathlineto{\pgfqpoint{3.727086in}{0.593547in}}%
\pgfpathlineto{\pgfqpoint{3.729184in}{0.610584in}}%
\pgfpathlineto{\pgfqpoint{3.731281in}{0.620395in}}%
\pgfpathlineto{\pgfqpoint{3.733378in}{0.624918in}}%
\pgfpathlineto{\pgfqpoint{3.735476in}{0.623825in}}%
\pgfpathlineto{\pgfqpoint{3.737573in}{0.616347in}}%
\pgfpathlineto{\pgfqpoint{3.739670in}{0.604510in}}%
\pgfpathlineto{\pgfqpoint{3.741768in}{0.586003in}}%
\pgfpathlineto{\pgfqpoint{3.745962in}{0.511838in}}%
\pgfpathlineto{\pgfqpoint{3.748060in}{0.563656in}}%
\pgfpathlineto{\pgfqpoint{3.750157in}{0.578591in}}%
\pgfpathlineto{\pgfqpoint{3.752254in}{0.583390in}}%
\pgfpathlineto{\pgfqpoint{3.754352in}{0.581371in}}%
\pgfpathlineto{\pgfqpoint{3.756449in}{0.571645in}}%
\pgfpathlineto{\pgfqpoint{3.758547in}{0.546214in}}%
\pgfpathlineto{\pgfqpoint{3.760644in}{0.497202in}}%
\pgfpathlineto{\pgfqpoint{3.762741in}{0.566471in}}%
\pgfpathlineto{\pgfqpoint{3.764839in}{0.589631in}}%
\pgfpathlineto{\pgfqpoint{3.766936in}{0.602487in}}%
\pgfpathlineto{\pgfqpoint{3.769033in}{0.609466in}}%
\pgfpathlineto{\pgfqpoint{3.771131in}{0.611369in}}%
\pgfpathlineto{\pgfqpoint{3.773228in}{0.606998in}}%
\pgfpathlineto{\pgfqpoint{3.775325in}{0.595695in}}%
\pgfpathlineto{\pgfqpoint{3.777423in}{0.573818in}}%
\pgfpathlineto{\pgfqpoint{3.779520in}{0.501728in}}%
\pgfpathlineto{\pgfqpoint{3.781618in}{0.559526in}}%
\pgfpathlineto{\pgfqpoint{3.783715in}{0.586809in}}%
\pgfpathlineto{\pgfqpoint{3.785812in}{0.599642in}}%
\pgfpathlineto{\pgfqpoint{3.787910in}{0.606213in}}%
\pgfpathlineto{\pgfqpoint{3.790007in}{0.608634in}}%
\pgfpathlineto{\pgfqpoint{3.792104in}{0.607465in}}%
\pgfpathlineto{\pgfqpoint{3.794202in}{0.602370in}}%
\pgfpathlineto{\pgfqpoint{3.796299in}{0.591720in}}%
\pgfpathlineto{\pgfqpoint{3.798396in}{0.569681in}}%
\pgfpathlineto{\pgfqpoint{3.800494in}{0.450437in}}%
\pgfpathlineto{\pgfqpoint{3.802591in}{0.569948in}}%
\pgfpathlineto{\pgfqpoint{3.804689in}{0.596785in}}%
\pgfpathlineto{\pgfqpoint{3.806786in}{0.611110in}}%
\pgfpathlineto{\pgfqpoint{3.808883in}{0.619409in}}%
\pgfpathlineto{\pgfqpoint{3.810981in}{0.623236in}}%
\pgfpathlineto{\pgfqpoint{3.813078in}{0.622441in}}%
\pgfpathlineto{\pgfqpoint{3.815175in}{0.615239in}}%
\pgfpathlineto{\pgfqpoint{3.817273in}{0.598561in}}%
\pgfpathlineto{\pgfqpoint{3.819370in}{0.554412in}}%
\pgfpathlineto{\pgfqpoint{3.821467in}{0.567088in}}%
\pgfpathlineto{\pgfqpoint{3.823565in}{0.603650in}}%
\pgfpathlineto{\pgfqpoint{3.825662in}{0.620025in}}%
\pgfpathlineto{\pgfqpoint{3.827760in}{0.629224in}}%
\pgfpathlineto{\pgfqpoint{3.829857in}{0.634213in}}%
\pgfpathlineto{\pgfqpoint{3.831954in}{0.636015in}}%
\pgfpathlineto{\pgfqpoint{3.834052in}{0.634894in}}%
\pgfpathlineto{\pgfqpoint{3.836149in}{0.630555in}}%
\pgfpathlineto{\pgfqpoint{3.838246in}{0.621885in}}%
\pgfpathlineto{\pgfqpoint{3.840344in}{0.605546in}}%
\pgfpathlineto{\pgfqpoint{3.842441in}{0.564976in}}%
\pgfpathlineto{\pgfqpoint{3.844538in}{0.571054in}}%
\pgfpathlineto{\pgfqpoint{3.846636in}{0.611136in}}%
\pgfpathlineto{\pgfqpoint{3.848733in}{0.629402in}}%
\pgfpathlineto{\pgfqpoint{3.850831in}{0.640135in}}%
\pgfpathlineto{\pgfqpoint{3.852928in}{0.646242in}}%
\pgfpathlineto{\pgfqpoint{3.855025in}{0.648397in}}%
\pgfpathlineto{\pgfqpoint{3.857123in}{0.646008in}}%
\pgfpathlineto{\pgfqpoint{3.859220in}{0.638422in}}%
\pgfpathlineto{\pgfqpoint{3.861317in}{0.626967in}}%
\pgfpathlineto{\pgfqpoint{3.863415in}{0.609791in}}%
\pgfpathlineto{\pgfqpoint{3.865512in}{0.580243in}}%
\pgfpathlineto{\pgfqpoint{3.867609in}{0.448333in}}%
\pgfpathlineto{\pgfqpoint{3.869707in}{0.572712in}}%
\pgfpathlineto{\pgfqpoint{3.871804in}{0.590609in}}%
\pgfpathlineto{\pgfqpoint{3.873902in}{0.595635in}}%
\pgfpathlineto{\pgfqpoint{3.875999in}{0.592432in}}%
\pgfpathlineto{\pgfqpoint{3.878096in}{0.579072in}}%
\pgfpathlineto{\pgfqpoint{3.880194in}{0.535881in}}%
\pgfpathlineto{\pgfqpoint{3.882291in}{0.557070in}}%
\pgfpathlineto{\pgfqpoint{3.884388in}{0.594814in}}%
\pgfpathlineto{\pgfqpoint{3.886486in}{0.612922in}}%
\pgfpathlineto{\pgfqpoint{3.888583in}{0.623291in}}%
\pgfpathlineto{\pgfqpoint{3.890680in}{0.628304in}}%
\pgfpathlineto{\pgfqpoint{3.892778in}{0.627897in}}%
\pgfpathlineto{\pgfqpoint{3.894875in}{0.621101in}}%
\pgfpathlineto{\pgfqpoint{3.896973in}{0.610327in}}%
\pgfpathlineto{\pgfqpoint{3.899070in}{0.594289in}}%
\pgfpathlineto{\pgfqpoint{3.901167in}{0.567705in}}%
\pgfpathlineto{\pgfqpoint{3.903265in}{0.481752in}}%
\pgfpathlineto{\pgfqpoint{3.905362in}{0.548725in}}%
\pgfpathlineto{\pgfqpoint{3.907459in}{0.570140in}}%
\pgfpathlineto{\pgfqpoint{3.909557in}{0.576655in}}%
\pgfpathlineto{\pgfqpoint{3.911654in}{0.574735in}}%
\pgfpathlineto{\pgfqpoint{3.913751in}{0.563413in}}%
\pgfpathlineto{\pgfqpoint{3.915849in}{0.529427in}}%
\pgfpathlineto{\pgfqpoint{3.917946in}{0.526818in}}%
\pgfpathlineto{\pgfqpoint{3.920043in}{0.572038in}}%
\pgfpathlineto{\pgfqpoint{3.922141in}{0.591928in}}%
\pgfpathlineto{\pgfqpoint{3.924238in}{0.603273in}}%
\pgfpathlineto{\pgfqpoint{3.926336in}{0.609089in}}%
\pgfpathlineto{\pgfqpoint{3.928433in}{0.609627in}}%
\pgfpathlineto{\pgfqpoint{3.930530in}{0.603416in}}%
\pgfpathlineto{\pgfqpoint{3.932628in}{0.590495in}}%
\pgfpathlineto{\pgfqpoint{3.934725in}{0.565644in}}%
\pgfpathlineto{\pgfqpoint{3.936822in}{0.429244in}}%
\pgfpathlineto{\pgfqpoint{3.938920in}{0.561180in}}%
\pgfpathlineto{\pgfqpoint{3.941017in}{0.584762in}}%
\pgfpathlineto{\pgfqpoint{3.943114in}{0.596023in}}%
\pgfpathlineto{\pgfqpoint{3.945212in}{0.601416in}}%
\pgfpathlineto{\pgfqpoint{3.947309in}{0.602664in}}%
\pgfpathlineto{\pgfqpoint{3.949407in}{0.600063in}}%
\pgfpathlineto{\pgfqpoint{3.951504in}{0.592844in}}%
\pgfpathlineto{\pgfqpoint{3.953601in}{0.578104in}}%
\pgfpathlineto{\pgfqpoint{3.955699in}{0.542008in}}%
\pgfpathlineto{\pgfqpoint{3.957796in}{0.536622in}}%
\pgfpathlineto{\pgfqpoint{3.959893in}{0.581009in}}%
\pgfpathlineto{\pgfqpoint{3.961991in}{0.600163in}}%
\pgfpathlineto{\pgfqpoint{3.964088in}{0.611133in}}%
\pgfpathlineto{\pgfqpoint{3.966185in}{0.617053in}}%
\pgfpathlineto{\pgfqpoint{3.968283in}{0.618484in}}%
\pgfpathlineto{\pgfqpoint{3.970380in}{0.614338in}}%
\pgfpathlineto{\pgfqpoint{3.972478in}{0.602068in}}%
\pgfpathlineto{\pgfqpoint{3.974575in}{0.573605in}}%
\pgfpathlineto{\pgfqpoint{3.976672in}{0.527055in}}%
\pgfpathlineto{\pgfqpoint{3.978770in}{0.590979in}}%
\pgfpathlineto{\pgfqpoint{3.980867in}{0.612011in}}%
\pgfpathlineto{\pgfqpoint{3.982964in}{0.623477in}}%
\pgfpathlineto{\pgfqpoint{3.985062in}{0.630034in}}%
\pgfpathlineto{\pgfqpoint{3.987159in}{0.633190in}}%
\pgfpathlineto{\pgfqpoint{3.989256in}{0.633464in}}%
\pgfpathlineto{\pgfqpoint{3.991354in}{0.630825in}}%
\pgfpathlineto{\pgfqpoint{3.993451in}{0.624659in}}%
\pgfpathlineto{\pgfqpoint{3.995549in}{0.613172in}}%
\pgfpathlineto{\pgfqpoint{3.997646in}{0.590314in}}%
\pgfpathlineto{\pgfqpoint{3.999743in}{0.452327in}}%
\pgfpathlineto{\pgfqpoint{4.001841in}{0.591205in}}%
\pgfpathlineto{\pgfqpoint{4.003938in}{0.617648in}}%
\pgfpathlineto{\pgfqpoint{4.006035in}{0.631958in}}%
\pgfpathlineto{\pgfqpoint{4.008133in}{0.640513in}}%
\pgfpathlineto{\pgfqpoint{4.010230in}{0.644959in}}%
\pgfpathlineto{\pgfqpoint{4.012327in}{0.645442in}}%
\pgfpathlineto{\pgfqpoint{4.014425in}{0.640866in}}%
\pgfpathlineto{\pgfqpoint{4.016522in}{0.631768in}}%
\pgfpathlineto{\pgfqpoint{4.018620in}{0.618105in}}%
\pgfpathlineto{\pgfqpoint{4.020717in}{0.596453in}}%
\pgfpathlineto{\pgfqpoint{4.022814in}{0.549646in}}%
\pgfpathlineto{\pgfqpoint{4.024912in}{0.551935in}}%
\pgfpathlineto{\pgfqpoint{4.027009in}{0.584920in}}%
\pgfpathlineto{\pgfqpoint{4.029106in}{0.595674in}}%
\pgfpathlineto{\pgfqpoint{4.031204in}{0.597458in}}%
\pgfpathlineto{\pgfqpoint{4.033301in}{0.591729in}}%
\pgfpathlineto{\pgfqpoint{4.035398in}{0.574711in}}%
\pgfpathlineto{\pgfqpoint{4.037496in}{0.507456in}}%
\pgfpathlineto{\pgfqpoint{4.039593in}{0.568432in}}%
\pgfpathlineto{\pgfqpoint{4.041691in}{0.599320in}}%
\pgfpathlineto{\pgfqpoint{4.043788in}{0.615266in}}%
\pgfpathlineto{\pgfqpoint{4.045885in}{0.624323in}}%
\pgfpathlineto{\pgfqpoint{4.047983in}{0.628201in}}%
\pgfpathlineto{\pgfqpoint{4.050080in}{0.626418in}}%
\pgfpathlineto{\pgfqpoint{4.052177in}{0.618589in}}%
\pgfpathlineto{\pgfqpoint{4.054275in}{0.606646in}}%
\pgfpathlineto{\pgfqpoint{4.056372in}{0.588324in}}%
\pgfpathlineto{\pgfqpoint{4.058469in}{0.554722in}}%
\pgfpathlineto{\pgfqpoint{4.060567in}{0.502017in}}%
\pgfpathlineto{\pgfqpoint{4.062664in}{0.561460in}}%
\pgfpathlineto{\pgfqpoint{4.064762in}{0.576717in}}%
\pgfpathlineto{\pgfqpoint{4.066859in}{0.581063in}}%
\pgfpathlineto{\pgfqpoint{4.068956in}{0.577930in}}%
\pgfpathlineto{\pgfqpoint{4.071054in}{0.565611in}}%
\pgfpathlineto{\pgfqpoint{4.073151in}{0.529505in}}%
\pgfpathlineto{\pgfqpoint{4.075248in}{0.531926in}}%
\pgfpathlineto{\pgfqpoint{4.077346in}{0.574996in}}%
\pgfpathlineto{\pgfqpoint{4.079443in}{0.594311in}}%
\pgfpathlineto{\pgfqpoint{4.081540in}{0.605315in}}%
\pgfpathlineto{\pgfqpoint{4.083638in}{0.610865in}}%
\pgfpathlineto{\pgfqpoint{4.085735in}{0.611142in}}%
\pgfpathlineto{\pgfqpoint{4.087832in}{0.604674in}}%
\pgfpathlineto{\pgfqpoint{4.089930in}{0.591431in}}%
\pgfpathlineto{\pgfqpoint{4.092027in}{0.565701in}}%
\pgfpathlineto{\pgfqpoint{4.094125in}{0.416105in}}%
\pgfpathlineto{\pgfqpoint{4.096222in}{0.564617in}}%
\pgfpathlineto{\pgfqpoint{4.098319in}{0.587314in}}%
\pgfpathlineto{\pgfqpoint{4.100417in}{0.598264in}}%
\pgfpathlineto{\pgfqpoint{4.102514in}{0.603486in}}%
\pgfpathlineto{\pgfqpoint{4.104611in}{0.604613in}}%
\pgfpathlineto{\pgfqpoint{4.106709in}{0.601904in}}%
\pgfpathlineto{\pgfqpoint{4.108806in}{0.594560in}}%
\pgfpathlineto{\pgfqpoint{4.110903in}{0.579604in}}%
\pgfpathlineto{\pgfqpoint{4.113001in}{0.542636in}}%
\pgfpathlineto{\pgfqpoint{4.115098in}{0.540020in}}%
\pgfpathlineto{\pgfqpoint{4.117196in}{0.583352in}}%
\pgfpathlineto{\pgfqpoint{4.119293in}{0.602288in}}%
\pgfpathlineto{\pgfqpoint{4.121390in}{0.613150in}}%
\pgfpathlineto{\pgfqpoint{4.123488in}{0.618994in}}%
\pgfpathlineto{\pgfqpoint{4.125585in}{0.620355in}}%
\pgfpathlineto{\pgfqpoint{4.127682in}{0.616125in}}%
\pgfpathlineto{\pgfqpoint{4.129780in}{0.603972in}}%
\pgfpathlineto{\pgfqpoint{4.131877in}{0.576540in}}%
\pgfpathlineto{\pgfqpoint{4.133974in}{0.518801in}}%
\pgfpathlineto{\pgfqpoint{4.136072in}{0.590027in}}%
\pgfpathlineto{\pgfqpoint{4.138169in}{0.611592in}}%
\pgfpathlineto{\pgfqpoint{4.140267in}{0.623170in}}%
\pgfpathlineto{\pgfqpoint{4.142364in}{0.629704in}}%
\pgfpathlineto{\pgfqpoint{4.144461in}{0.632758in}}%
\pgfpathlineto{\pgfqpoint{4.146559in}{0.632853in}}%
\pgfpathlineto{\pgfqpoint{4.148656in}{0.629930in}}%
\pgfpathlineto{\pgfqpoint{4.150753in}{0.623289in}}%
\pgfpathlineto{\pgfqpoint{4.152851in}{0.610865in}}%
\pgfpathlineto{\pgfqpoint{4.154948in}{0.585193in}}%
\pgfpathlineto{\pgfqpoint{4.157045in}{0.509487in}}%
\pgfpathlineto{\pgfqpoint{4.159143in}{0.596191in}}%
\pgfpathlineto{\pgfqpoint{4.161240in}{0.620089in}}%
\pgfpathlineto{\pgfqpoint{4.163338in}{0.633394in}}%
\pgfpathlineto{\pgfqpoint{4.165435in}{0.641279in}}%
\pgfpathlineto{\pgfqpoint{4.167532in}{0.645095in}}%
\pgfpathlineto{\pgfqpoint{4.169630in}{0.644787in}}%
\pgfpathlineto{\pgfqpoint{4.171727in}{0.639166in}}%
\pgfpathlineto{\pgfqpoint{4.173824in}{0.629260in}}%
\pgfpathlineto{\pgfqpoint{4.175922in}{0.614306in}}%
\pgfpathlineto{\pgfqpoint{4.178019in}{0.589726in}}%
\pgfpathlineto{\pgfqpoint{4.180116in}{0.522939in}}%
\pgfpathlineto{\pgfqpoint{4.182214in}{0.564325in}}%
\pgfpathlineto{\pgfqpoint{4.184311in}{0.589099in}}%
\pgfpathlineto{\pgfqpoint{4.186409in}{0.597556in}}%
\pgfpathlineto{\pgfqpoint{4.188506in}{0.597929in}}%
\pgfpathlineto{\pgfqpoint{4.190603in}{0.590731in}}%
\pgfpathlineto{\pgfqpoint{4.192701in}{0.570791in}}%
\pgfpathlineto{\pgfqpoint{4.194798in}{0.428031in}}%
\pgfpathlineto{\pgfqpoint{4.196895in}{0.575560in}}%
\pgfpathlineto{\pgfqpoint{4.198993in}{0.602807in}}%
\pgfpathlineto{\pgfqpoint{4.201090in}{0.617416in}}%
\pgfpathlineto{\pgfqpoint{4.203187in}{0.625609in}}%
\pgfpathlineto{\pgfqpoint{4.205285in}{0.628686in}}%
\pgfpathlineto{\pgfqpoint{4.207382in}{0.625836in}}%
\pgfpathlineto{\pgfqpoint{4.209480in}{0.617382in}}%
\pgfpathlineto{\pgfqpoint{4.211577in}{0.604597in}}%
\pgfpathlineto{\pgfqpoint{4.213674in}{0.584618in}}%
\pgfpathlineto{\pgfqpoint{4.215772in}{0.544921in}}%
\pgfpathlineto{\pgfqpoint{4.217869in}{0.525412in}}%
\pgfpathlineto{\pgfqpoint{4.219966in}{0.566702in}}%
\pgfpathlineto{\pgfqpoint{4.222064in}{0.579430in}}%
\pgfpathlineto{\pgfqpoint{4.224161in}{0.582544in}}%
\pgfpathlineto{\pgfqpoint{4.226258in}{0.578382in}}%
\pgfpathlineto{\pgfqpoint{4.228356in}{0.564540in}}%
\pgfpathlineto{\pgfqpoint{4.230453in}{0.521800in}}%
\pgfpathlineto{\pgfqpoint{4.232551in}{0.540715in}}%
\pgfpathlineto{\pgfqpoint{4.234648in}{0.578718in}}%
\pgfpathlineto{\pgfqpoint{4.236745in}{0.596779in}}%
\pgfpathlineto{\pgfqpoint{4.238843in}{0.607087in}}%
\pgfpathlineto{\pgfqpoint{4.240940in}{0.612065in}}%
\pgfpathlineto{\pgfqpoint{4.243037in}{0.611681in}}%
\pgfpathlineto{\pgfqpoint{4.245135in}{0.604485in}}%
\pgfpathlineto{\pgfqpoint{4.247232in}{0.590430in}}%
\pgfpathlineto{\pgfqpoint{4.249329in}{0.562658in}}%
\pgfpathlineto{\pgfqpoint{4.251427in}{0.482918in}}%
\pgfpathlineto{\pgfqpoint{4.253524in}{0.568035in}}%
\pgfpathlineto{\pgfqpoint{4.255621in}{0.589033in}}%
\pgfpathlineto{\pgfqpoint{4.257719in}{0.599269in}}%
\pgfpathlineto{\pgfqpoint{4.259816in}{0.604005in}}%
\pgfpathlineto{\pgfqpoint{4.261914in}{0.604686in}}%
\pgfpathlineto{\pgfqpoint{4.264011in}{0.601459in}}%
\pgfpathlineto{\pgfqpoint{4.266108in}{0.593345in}}%
\pgfpathlineto{\pgfqpoint{4.268206in}{0.576777in}}%
\pgfpathlineto{\pgfqpoint{4.270303in}{0.531670in}}%
\pgfpathlineto{\pgfqpoint{4.272400in}{0.550395in}}%
\pgfpathlineto{\pgfqpoint{4.274498in}{0.587184in}}%
\pgfpathlineto{\pgfqpoint{4.276595in}{0.604591in}}%
\pgfpathlineto{\pgfqpoint{4.278692in}{0.614639in}}%
\pgfpathlineto{\pgfqpoint{4.280790in}{0.619836in}}%
\pgfpathlineto{\pgfqpoint{4.282887in}{0.620482in}}%
\pgfpathlineto{\pgfqpoint{4.284985in}{0.615471in}}%
\pgfpathlineto{\pgfqpoint{4.287082in}{0.607174in}}%
\pgfpathlineto{\pgfqpoint{4.291277in}{0.583123in}}%
\pgfpathlineto{\pgfqpoint{4.293374in}{0.565884in}}%
\pgfpathlineto{\pgfqpoint{4.295471in}{0.543186in}}%
\pgfpathlineto{\pgfqpoint{4.299666in}{0.474866in}}%
\pgfpathlineto{\pgfqpoint{4.301763in}{0.474063in}}%
\pgfpathlineto{\pgfqpoint{4.305958in}{0.536032in}}%
\pgfpathlineto{\pgfqpoint{4.308056in}{0.554983in}}%
\pgfpathlineto{\pgfqpoint{4.310153in}{0.567874in}}%
\pgfpathlineto{\pgfqpoint{4.312250in}{0.575930in}}%
\pgfpathlineto{\pgfqpoint{4.314348in}{0.579295in}}%
\pgfpathlineto{\pgfqpoint{4.316445in}{0.576717in}}%
\pgfpathlineto{\pgfqpoint{4.316445in}{0.576717in}}%
\pgfusepath{stroke}%
\end{pgfscope}%
\begin{pgfscope}%
\pgfpathrectangle{\pgfqpoint{2.554660in}{0.320625in}}{\pgfqpoint{1.845679in}{1.164375in}}%
\pgfusepath{clip}%
\pgfsetrectcap%
\pgfsetroundjoin%
\pgfsetlinewidth{1.154312pt}%
\definecolor{currentstroke}{rgb}{0.997351,0.719089,0.493755}%
\pgfsetstrokecolor{currentstroke}%
\pgfsetdash{}{0pt}%
\pgfpathmoveto{\pgfqpoint{2.638555in}{1.468525in}}%
\pgfpathlineto{\pgfqpoint{2.649042in}{1.449020in}}%
\pgfpathlineto{\pgfqpoint{2.659529in}{1.426689in}}%
\pgfpathlineto{\pgfqpoint{2.672113in}{1.396826in}}%
\pgfpathlineto{\pgfqpoint{2.714060in}{1.295102in}}%
\pgfpathlineto{\pgfqpoint{2.737131in}{1.242905in}}%
\pgfpathlineto{\pgfqpoint{2.764397in}{1.184202in}}%
\pgfpathlineto{\pgfqpoint{2.797955in}{1.114812in}}%
\pgfpathlineto{\pgfqpoint{2.841999in}{1.026516in}}%
\pgfpathlineto{\pgfqpoint{2.902823in}{0.907336in}}%
\pgfpathlineto{\pgfqpoint{2.978328in}{0.761770in}}%
\pgfpathlineto{\pgfqpoint{2.990912in}{0.737306in}}%
\pgfpathlineto{\pgfqpoint{3.003496in}{0.713371in}}%
\pgfpathlineto{\pgfqpoint{3.016080in}{0.690581in}}%
\pgfpathlineto{\pgfqpoint{3.034957in}{0.651376in}}%
\pgfpathlineto{\pgfqpoint{3.039151in}{0.645504in}}%
\pgfpathlineto{\pgfqpoint{3.045443in}{0.639267in}}%
\pgfpathlineto{\pgfqpoint{3.049638in}{0.634500in}}%
\pgfpathlineto{\pgfqpoint{3.053833in}{0.626257in}}%
\pgfpathlineto{\pgfqpoint{3.058027in}{0.615244in}}%
\pgfpathlineto{\pgfqpoint{3.062222in}{0.601390in}}%
\pgfpathlineto{\pgfqpoint{3.068514in}{0.575320in}}%
\pgfpathlineto{\pgfqpoint{3.072709in}{0.557112in}}%
\pgfpathlineto{\pgfqpoint{3.074806in}{0.549758in}}%
\pgfpathlineto{\pgfqpoint{3.076904in}{0.544705in}}%
\pgfpathlineto{\pgfqpoint{3.079001in}{0.542704in}}%
\pgfpathlineto{\pgfqpoint{3.081098in}{0.544411in}}%
\pgfpathlineto{\pgfqpoint{3.083196in}{0.549625in}}%
\pgfpathlineto{\pgfqpoint{3.091585in}{0.577878in}}%
\pgfpathlineto{\pgfqpoint{3.093683in}{0.582692in}}%
\pgfpathlineto{\pgfqpoint{3.095780in}{0.586081in}}%
\pgfpathlineto{\pgfqpoint{3.097877in}{0.587857in}}%
\pgfpathlineto{\pgfqpoint{3.099975in}{0.587683in}}%
\pgfpathlineto{\pgfqpoint{3.102072in}{0.584901in}}%
\pgfpathlineto{\pgfqpoint{3.104169in}{0.579145in}}%
\pgfpathlineto{\pgfqpoint{3.106267in}{0.570189in}}%
\pgfpathlineto{\pgfqpoint{3.110462in}{0.544015in}}%
\pgfpathlineto{\pgfqpoint{3.112559in}{0.547323in}}%
\pgfpathlineto{\pgfqpoint{3.116754in}{0.575543in}}%
\pgfpathlineto{\pgfqpoint{3.120948in}{0.592092in}}%
\pgfpathlineto{\pgfqpoint{3.123046in}{0.597191in}}%
\pgfpathlineto{\pgfqpoint{3.125143in}{0.600781in}}%
\pgfpathlineto{\pgfqpoint{3.127240in}{0.603095in}}%
\pgfpathlineto{\pgfqpoint{3.129338in}{0.604258in}}%
\pgfpathlineto{\pgfqpoint{3.131435in}{0.604321in}}%
\pgfpathlineto{\pgfqpoint{3.133533in}{0.603260in}}%
\pgfpathlineto{\pgfqpoint{3.135630in}{0.600976in}}%
\pgfpathlineto{\pgfqpoint{3.137727in}{0.597259in}}%
\pgfpathlineto{\pgfqpoint{3.139825in}{0.591720in}}%
\pgfpathlineto{\pgfqpoint{3.141922in}{0.583604in}}%
\pgfpathlineto{\pgfqpoint{3.144019in}{0.571251in}}%
\pgfpathlineto{\pgfqpoint{3.146117in}{0.549873in}}%
\pgfpathlineto{\pgfqpoint{3.148214in}{0.493255in}}%
\pgfpathlineto{\pgfqpoint{3.150311in}{0.539372in}}%
\pgfpathlineto{\pgfqpoint{3.152409in}{0.568614in}}%
\pgfpathlineto{\pgfqpoint{3.154506in}{0.584842in}}%
\pgfpathlineto{\pgfqpoint{3.156604in}{0.595751in}}%
\pgfpathlineto{\pgfqpoint{3.160798in}{0.609282in}}%
\pgfpathlineto{\pgfqpoint{3.162896in}{0.613259in}}%
\pgfpathlineto{\pgfqpoint{3.164993in}{0.615689in}}%
\pgfpathlineto{\pgfqpoint{3.167090in}{0.616564in}}%
\pgfpathlineto{\pgfqpoint{3.169188in}{0.615690in}}%
\pgfpathlineto{\pgfqpoint{3.171285in}{0.612592in}}%
\pgfpathlineto{\pgfqpoint{3.173382in}{0.606569in}}%
\pgfpathlineto{\pgfqpoint{3.175480in}{0.598461in}}%
\pgfpathlineto{\pgfqpoint{3.177577in}{0.593906in}}%
\pgfpathlineto{\pgfqpoint{3.179675in}{0.600064in}}%
\pgfpathlineto{\pgfqpoint{3.183869in}{0.621158in}}%
\pgfpathlineto{\pgfqpoint{3.185967in}{0.628774in}}%
\pgfpathlineto{\pgfqpoint{3.188064in}{0.634258in}}%
\pgfpathlineto{\pgfqpoint{3.190161in}{0.637912in}}%
\pgfpathlineto{\pgfqpoint{3.192259in}{0.639936in}}%
\pgfpathlineto{\pgfqpoint{3.194356in}{0.640415in}}%
\pgfpathlineto{\pgfqpoint{3.196453in}{0.639315in}}%
\pgfpathlineto{\pgfqpoint{3.198551in}{0.636460in}}%
\pgfpathlineto{\pgfqpoint{3.200648in}{0.631451in}}%
\pgfpathlineto{\pgfqpoint{3.202746in}{0.623448in}}%
\pgfpathlineto{\pgfqpoint{3.204843in}{0.610481in}}%
\pgfpathlineto{\pgfqpoint{3.206940in}{0.586220in}}%
\pgfpathlineto{\pgfqpoint{3.209038in}{0.512108in}}%
\pgfpathlineto{\pgfqpoint{3.211135in}{0.589469in}}%
\pgfpathlineto{\pgfqpoint{3.213232in}{0.615099in}}%
\pgfpathlineto{\pgfqpoint{3.215330in}{0.629752in}}%
\pgfpathlineto{\pgfqpoint{3.217427in}{0.639403in}}%
\pgfpathlineto{\pgfqpoint{3.219524in}{0.645871in}}%
\pgfpathlineto{\pgfqpoint{3.221622in}{0.649839in}}%
\pgfpathlineto{\pgfqpoint{3.223719in}{0.651462in}}%
\pgfpathlineto{\pgfqpoint{3.225816in}{0.650472in}}%
\pgfpathlineto{\pgfqpoint{3.227914in}{0.645992in}}%
\pgfpathlineto{\pgfqpoint{3.230011in}{0.638158in}}%
\pgfpathlineto{\pgfqpoint{3.232109in}{0.627381in}}%
\pgfpathlineto{\pgfqpoint{3.238401in}{0.586277in}}%
\pgfpathlineto{\pgfqpoint{3.242595in}{0.596800in}}%
\pgfpathlineto{\pgfqpoint{3.244693in}{0.598099in}}%
\pgfpathlineto{\pgfqpoint{3.246790in}{0.593628in}}%
\pgfpathlineto{\pgfqpoint{3.248887in}{0.581265in}}%
\pgfpathlineto{\pgfqpoint{3.250985in}{0.557891in}}%
\pgfpathlineto{\pgfqpoint{3.253082in}{0.569312in}}%
\pgfpathlineto{\pgfqpoint{3.255180in}{0.596174in}}%
\pgfpathlineto{\pgfqpoint{3.257277in}{0.612442in}}%
\pgfpathlineto{\pgfqpoint{3.259374in}{0.622199in}}%
\pgfpathlineto{\pgfqpoint{3.261472in}{0.627024in}}%
\pgfpathlineto{\pgfqpoint{3.263569in}{0.626754in}}%
\pgfpathlineto{\pgfqpoint{3.265666in}{0.620323in}}%
\pgfpathlineto{\pgfqpoint{3.267764in}{0.609975in}}%
\pgfpathlineto{\pgfqpoint{3.269861in}{0.594794in}}%
\pgfpathlineto{\pgfqpoint{3.271958in}{0.571500in}}%
\pgfpathlineto{\pgfqpoint{3.274056in}{0.541785in}}%
\pgfpathlineto{\pgfqpoint{3.278251in}{0.574294in}}%
\pgfpathlineto{\pgfqpoint{3.280348in}{0.580408in}}%
\pgfpathlineto{\pgfqpoint{3.282445in}{0.579511in}}%
\pgfpathlineto{\pgfqpoint{3.284543in}{0.571151in}}%
\pgfpathlineto{\pgfqpoint{3.286640in}{0.549383in}}%
\pgfpathlineto{\pgfqpoint{3.288737in}{0.505861in}}%
\pgfpathlineto{\pgfqpoint{3.290835in}{0.562596in}}%
\pgfpathlineto{\pgfqpoint{3.292932in}{0.586830in}}%
\pgfpathlineto{\pgfqpoint{3.295029in}{0.600376in}}%
\pgfpathlineto{\pgfqpoint{3.297127in}{0.607948in}}%
\pgfpathlineto{\pgfqpoint{3.299224in}{0.610526in}}%
\pgfpathlineto{\pgfqpoint{3.301322in}{0.607128in}}%
\pgfpathlineto{\pgfqpoint{3.303419in}{0.596490in}}%
\pgfpathlineto{\pgfqpoint{3.305516in}{0.575190in}}%
\pgfpathlineto{\pgfqpoint{3.307614in}{0.518887in}}%
\pgfpathlineto{\pgfqpoint{3.309711in}{0.564813in}}%
\pgfpathlineto{\pgfqpoint{3.311808in}{0.590984in}}%
\pgfpathlineto{\pgfqpoint{3.313906in}{0.604019in}}%
\pgfpathlineto{\pgfqpoint{3.316003in}{0.611096in}}%
\pgfpathlineto{\pgfqpoint{3.318100in}{0.614257in}}%
\pgfpathlineto{\pgfqpoint{3.320198in}{0.614145in}}%
\pgfpathlineto{\pgfqpoint{3.322295in}{0.610671in}}%
\pgfpathlineto{\pgfqpoint{3.324393in}{0.602913in}}%
\pgfpathlineto{\pgfqpoint{3.326490in}{0.587975in}}%
\pgfpathlineto{\pgfqpoint{3.328587in}{0.552779in}}%
\pgfpathlineto{\pgfqpoint{3.330685in}{0.542958in}}%
\pgfpathlineto{\pgfqpoint{3.332782in}{0.588865in}}%
\pgfpathlineto{\pgfqpoint{3.334879in}{0.608337in}}%
\pgfpathlineto{\pgfqpoint{3.336977in}{0.619596in}}%
\pgfpathlineto{\pgfqpoint{3.339074in}{0.625960in}}%
\pgfpathlineto{\pgfqpoint{3.341171in}{0.628180in}}%
\pgfpathlineto{\pgfqpoint{3.343269in}{0.625618in}}%
\pgfpathlineto{\pgfqpoint{3.345366in}{0.617139in}}%
\pgfpathlineto{\pgfqpoint{3.347464in}{0.602126in}}%
\pgfpathlineto{\pgfqpoint{3.349561in}{0.573232in}}%
\pgfpathlineto{\pgfqpoint{3.351658in}{0.521101in}}%
\pgfpathlineto{\pgfqpoint{3.353756in}{0.580605in}}%
\pgfpathlineto{\pgfqpoint{3.355853in}{0.600055in}}%
\pgfpathlineto{\pgfqpoint{3.357950in}{0.609306in}}%
\pgfpathlineto{\pgfqpoint{3.360048in}{0.613055in}}%
\pgfpathlineto{\pgfqpoint{3.362145in}{0.612523in}}%
\pgfpathlineto{\pgfqpoint{3.364242in}{0.607538in}}%
\pgfpathlineto{\pgfqpoint{3.366340in}{0.596292in}}%
\pgfpathlineto{\pgfqpoint{3.368437in}{0.571453in}}%
\pgfpathlineto{\pgfqpoint{3.370535in}{0.497944in}}%
\pgfpathlineto{\pgfqpoint{3.372632in}{0.583484in}}%
\pgfpathlineto{\pgfqpoint{3.374729in}{0.607486in}}%
\pgfpathlineto{\pgfqpoint{3.376827in}{0.620660in}}%
\pgfpathlineto{\pgfqpoint{3.378924in}{0.628077in}}%
\pgfpathlineto{\pgfqpoint{3.381021in}{0.630873in}}%
\pgfpathlineto{\pgfqpoint{3.383119in}{0.628433in}}%
\pgfpathlineto{\pgfqpoint{3.385216in}{0.619749in}}%
\pgfpathlineto{\pgfqpoint{3.387313in}{0.604338in}}%
\pgfpathlineto{\pgfqpoint{3.389411in}{0.574027in}}%
\pgfpathlineto{\pgfqpoint{3.391508in}{0.525127in}}%
\pgfpathlineto{\pgfqpoint{3.393605in}{0.585025in}}%
\pgfpathlineto{\pgfqpoint{3.395703in}{0.603920in}}%
\pgfpathlineto{\pgfqpoint{3.397800in}{0.612917in}}%
\pgfpathlineto{\pgfqpoint{3.399898in}{0.616500in}}%
\pgfpathlineto{\pgfqpoint{3.401995in}{0.615827in}}%
\pgfpathlineto{\pgfqpoint{3.404092in}{0.610682in}}%
\pgfpathlineto{\pgfqpoint{3.406190in}{0.599186in}}%
\pgfpathlineto{\pgfqpoint{3.408287in}{0.573637in}}%
\pgfpathlineto{\pgfqpoint{3.410384in}{0.508956in}}%
\pgfpathlineto{\pgfqpoint{3.412482in}{0.587782in}}%
\pgfpathlineto{\pgfqpoint{3.414579in}{0.611299in}}%
\pgfpathlineto{\pgfqpoint{3.416676in}{0.624263in}}%
\pgfpathlineto{\pgfqpoint{3.418774in}{0.631534in}}%
\pgfpathlineto{\pgfqpoint{3.420871in}{0.634186in}}%
\pgfpathlineto{\pgfqpoint{3.422969in}{0.631549in}}%
\pgfpathlineto{\pgfqpoint{3.425066in}{0.622684in}}%
\pgfpathlineto{\pgfqpoint{3.427163in}{0.606991in}}%
\pgfpathlineto{\pgfqpoint{3.429261in}{0.575767in}}%
\pgfpathlineto{\pgfqpoint{3.431358in}{0.530180in}}%
\pgfpathlineto{\pgfqpoint{3.433455in}{0.588908in}}%
\pgfpathlineto{\pgfqpoint{3.435553in}{0.607424in}}%
\pgfpathlineto{\pgfqpoint{3.437650in}{0.616224in}}%
\pgfpathlineto{\pgfqpoint{3.439747in}{0.619654in}}%
\pgfpathlineto{\pgfqpoint{3.441845in}{0.618819in}}%
\pgfpathlineto{\pgfqpoint{3.443942in}{0.613460in}}%
\pgfpathlineto{\pgfqpoint{3.446040in}{0.601576in}}%
\pgfpathlineto{\pgfqpoint{3.448137in}{0.574822in}}%
\pgfpathlineto{\pgfqpoint{3.450234in}{0.522727in}}%
\pgfpathlineto{\pgfqpoint{3.452332in}{0.592517in}}%
\pgfpathlineto{\pgfqpoint{3.454429in}{0.615277in}}%
\pgfpathlineto{\pgfqpoint{3.456526in}{0.627910in}}%
\pgfpathlineto{\pgfqpoint{3.458624in}{0.634945in}}%
\pgfpathlineto{\pgfqpoint{3.460721in}{0.637359in}}%
\pgfpathlineto{\pgfqpoint{3.462818in}{0.634381in}}%
\pgfpathlineto{\pgfqpoint{3.464916in}{0.625426in}}%
\pgfpathlineto{\pgfqpoint{3.467013in}{0.610210in}}%
\pgfpathlineto{\pgfqpoint{3.469111in}{0.581595in}}%
\pgfpathlineto{\pgfqpoint{3.471208in}{0.503027in}}%
\pgfpathlineto{\pgfqpoint{3.473305in}{0.584416in}}%
\pgfpathlineto{\pgfqpoint{3.475403in}{0.604393in}}%
\pgfpathlineto{\pgfqpoint{3.477500in}{0.613378in}}%
\pgfpathlineto{\pgfqpoint{3.479597in}{0.616514in}}%
\pgfpathlineto{\pgfqpoint{3.481695in}{0.614970in}}%
\pgfpathlineto{\pgfqpoint{3.483792in}{0.608238in}}%
\pgfpathlineto{\pgfqpoint{3.485889in}{0.593271in}}%
\pgfpathlineto{\pgfqpoint{3.487987in}{0.553621in}}%
\pgfpathlineto{\pgfqpoint{3.490084in}{0.560557in}}%
\pgfpathlineto{\pgfqpoint{3.492182in}{0.600949in}}%
\pgfpathlineto{\pgfqpoint{3.494279in}{0.619286in}}%
\pgfpathlineto{\pgfqpoint{3.496376in}{0.629746in}}%
\pgfpathlineto{\pgfqpoint{3.498474in}{0.635075in}}%
\pgfpathlineto{\pgfqpoint{3.500571in}{0.635556in}}%
\pgfpathlineto{\pgfqpoint{3.502668in}{0.629916in}}%
\pgfpathlineto{\pgfqpoint{3.504766in}{0.618970in}}%
\pgfpathlineto{\pgfqpoint{3.506863in}{0.600756in}}%
\pgfpathlineto{\pgfqpoint{3.508960in}{0.563343in}}%
\pgfpathlineto{\pgfqpoint{3.511058in}{0.543116in}}%
\pgfpathlineto{\pgfqpoint{3.513155in}{0.587263in}}%
\pgfpathlineto{\pgfqpoint{3.515253in}{0.602571in}}%
\pgfpathlineto{\pgfqpoint{3.517350in}{0.609031in}}%
\pgfpathlineto{\pgfqpoint{3.519447in}{0.609956in}}%
\pgfpathlineto{\pgfqpoint{3.521545in}{0.605753in}}%
\pgfpathlineto{\pgfqpoint{3.523642in}{0.594716in}}%
\pgfpathlineto{\pgfqpoint{3.525739in}{0.568766in}}%
\pgfpathlineto{\pgfqpoint{3.527837in}{0.514514in}}%
\pgfpathlineto{\pgfqpoint{3.529934in}{0.586349in}}%
\pgfpathlineto{\pgfqpoint{3.532031in}{0.609429in}}%
\pgfpathlineto{\pgfqpoint{3.534129in}{0.622181in}}%
\pgfpathlineto{\pgfqpoint{3.536226in}{0.629176in}}%
\pgfpathlineto{\pgfqpoint{3.538324in}{0.631324in}}%
\pgfpathlineto{\pgfqpoint{3.540421in}{0.627643in}}%
\pgfpathlineto{\pgfqpoint{3.542518in}{0.618126in}}%
\pgfpathlineto{\pgfqpoint{3.544616in}{0.602157in}}%
\pgfpathlineto{\pgfqpoint{3.546713in}{0.571692in}}%
\pgfpathlineto{\pgfqpoint{3.548810in}{0.510415in}}%
\pgfpathlineto{\pgfqpoint{3.550908in}{0.579181in}}%
\pgfpathlineto{\pgfqpoint{3.553005in}{0.597863in}}%
\pgfpathlineto{\pgfqpoint{3.555102in}{0.606137in}}%
\pgfpathlineto{\pgfqpoint{3.557200in}{0.608631in}}%
\pgfpathlineto{\pgfqpoint{3.559297in}{0.606311in}}%
\pgfpathlineto{\pgfqpoint{3.561395in}{0.598363in}}%
\pgfpathlineto{\pgfqpoint{3.563492in}{0.580673in}}%
\pgfpathlineto{\pgfqpoint{3.565589in}{0.523463in}}%
\pgfpathlineto{\pgfqpoint{3.569784in}{0.598690in}}%
\pgfpathlineto{\pgfqpoint{3.571881in}{0.614862in}}%
\pgfpathlineto{\pgfqpoint{3.573979in}{0.624091in}}%
\pgfpathlineto{\pgfqpoint{3.576076in}{0.628369in}}%
\pgfpathlineto{\pgfqpoint{3.578173in}{0.627569in}}%
\pgfpathlineto{\pgfqpoint{3.580271in}{0.620375in}}%
\pgfpathlineto{\pgfqpoint{3.582368in}{0.607057in}}%
\pgfpathlineto{\pgfqpoint{3.584465in}{0.582492in}}%
\pgfpathlineto{\pgfqpoint{3.586563in}{0.483718in}}%
\pgfpathlineto{\pgfqpoint{3.588660in}{0.573187in}}%
\pgfpathlineto{\pgfqpoint{3.590758in}{0.597209in}}%
\pgfpathlineto{\pgfqpoint{3.592855in}{0.608141in}}%
\pgfpathlineto{\pgfqpoint{3.594952in}{0.612893in}}%
\pgfpathlineto{\pgfqpoint{3.597050in}{0.613168in}}%
\pgfpathlineto{\pgfqpoint{3.599147in}{0.609044in}}%
\pgfpathlineto{\pgfqpoint{3.601244in}{0.599085in}}%
\pgfpathlineto{\pgfqpoint{3.603342in}{0.577573in}}%
\pgfpathlineto{\pgfqpoint{3.605439in}{0.460729in}}%
\pgfpathlineto{\pgfqpoint{3.607536in}{0.578239in}}%
\pgfpathlineto{\pgfqpoint{3.609634in}{0.605243in}}%
\pgfpathlineto{\pgfqpoint{3.611731in}{0.619589in}}%
\pgfpathlineto{\pgfqpoint{3.613829in}{0.627756in}}%
\pgfpathlineto{\pgfqpoint{3.615926in}{0.631216in}}%
\pgfpathlineto{\pgfqpoint{3.618023in}{0.629597in}}%
\pgfpathlineto{\pgfqpoint{3.620121in}{0.621621in}}%
\pgfpathlineto{\pgfqpoint{3.622218in}{0.606992in}}%
\pgfpathlineto{\pgfqpoint{3.624315in}{0.578403in}}%
\pgfpathlineto{\pgfqpoint{3.626413in}{0.505664in}}%
\pgfpathlineto{\pgfqpoint{3.628510in}{0.584614in}}%
\pgfpathlineto{\pgfqpoint{3.630607in}{0.604932in}}%
\pgfpathlineto{\pgfqpoint{3.632705in}{0.614639in}}%
\pgfpathlineto{\pgfqpoint{3.634802in}{0.618805in}}%
\pgfpathlineto{\pgfqpoint{3.636900in}{0.618765in}}%
\pgfpathlineto{\pgfqpoint{3.638997in}{0.614490in}}%
\pgfpathlineto{\pgfqpoint{3.641094in}{0.604537in}}%
\pgfpathlineto{\pgfqpoint{3.643192in}{0.583390in}}%
\pgfpathlineto{\pgfqpoint{3.645289in}{0.483886in}}%
\pgfpathlineto{\pgfqpoint{3.647386in}{0.581892in}}%
\pgfpathlineto{\pgfqpoint{3.649484in}{0.609450in}}%
\pgfpathlineto{\pgfqpoint{3.651581in}{0.624006in}}%
\pgfpathlineto{\pgfqpoint{3.653678in}{0.632339in}}%
\pgfpathlineto{\pgfqpoint{3.655776in}{0.636010in}}%
\pgfpathlineto{\pgfqpoint{3.657873in}{0.634750in}}%
\pgfpathlineto{\pgfqpoint{3.659971in}{0.627298in}}%
\pgfpathlineto{\pgfqpoint{3.662068in}{0.613835in}}%
\pgfpathlineto{\pgfqpoint{3.664165in}{0.589186in}}%
\pgfpathlineto{\pgfqpoint{3.666263in}{0.491069in}}%
\pgfpathlineto{\pgfqpoint{3.668360in}{0.579340in}}%
\pgfpathlineto{\pgfqpoint{3.670457in}{0.603296in}}%
\pgfpathlineto{\pgfqpoint{3.672555in}{0.614084in}}%
\pgfpathlineto{\pgfqpoint{3.674652in}{0.618635in}}%
\pgfpathlineto{\pgfqpoint{3.676749in}{0.618627in}}%
\pgfpathlineto{\pgfqpoint{3.678847in}{0.614058in}}%
\pgfpathlineto{\pgfqpoint{3.680944in}{0.603255in}}%
\pgfpathlineto{\pgfqpoint{3.683042in}{0.579293in}}%
\pgfpathlineto{\pgfqpoint{3.685139in}{0.490331in}}%
\pgfpathlineto{\pgfqpoint{3.687236in}{0.589013in}}%
\pgfpathlineto{\pgfqpoint{3.689334in}{0.613638in}}%
\pgfpathlineto{\pgfqpoint{3.691431in}{0.627052in}}%
\pgfpathlineto{\pgfqpoint{3.693528in}{0.634588in}}%
\pgfpathlineto{\pgfqpoint{3.695626in}{0.637427in}}%
\pgfpathlineto{\pgfqpoint{3.697723in}{0.634942in}}%
\pgfpathlineto{\pgfqpoint{3.699820in}{0.626453in}}%
\pgfpathlineto{\pgfqpoint{3.701918in}{0.612285in}}%
\pgfpathlineto{\pgfqpoint{3.704015in}{0.587049in}}%
\pgfpathlineto{\pgfqpoint{3.706113in}{0.489553in}}%
\pgfpathlineto{\pgfqpoint{3.708210in}{0.575719in}}%
\pgfpathlineto{\pgfqpoint{3.710307in}{0.599045in}}%
\pgfpathlineto{\pgfqpoint{3.712405in}{0.609022in}}%
\pgfpathlineto{\pgfqpoint{3.714502in}{0.612514in}}%
\pgfpathlineto{\pgfqpoint{3.716599in}{0.610990in}}%
\pgfpathlineto{\pgfqpoint{3.718697in}{0.603929in}}%
\pgfpathlineto{\pgfqpoint{3.720794in}{0.587851in}}%
\pgfpathlineto{\pgfqpoint{3.722891in}{0.541136in}}%
\pgfpathlineto{\pgfqpoint{3.724989in}{0.565070in}}%
\pgfpathlineto{\pgfqpoint{3.727086in}{0.601003in}}%
\pgfpathlineto{\pgfqpoint{3.729184in}{0.618235in}}%
\pgfpathlineto{\pgfqpoint{3.731281in}{0.628067in}}%
\pgfpathlineto{\pgfqpoint{3.733378in}{0.632828in}}%
\pgfpathlineto{\pgfqpoint{3.735476in}{0.632559in}}%
\pgfpathlineto{\pgfqpoint{3.737573in}{0.626010in}}%
\pgfpathlineto{\pgfqpoint{3.739670in}{0.614202in}}%
\pgfpathlineto{\pgfqpoint{3.741768in}{0.594336in}}%
\pgfpathlineto{\pgfqpoint{3.743865in}{0.549811in}}%
\pgfpathlineto{\pgfqpoint{3.745962in}{0.552263in}}%
\pgfpathlineto{\pgfqpoint{3.748060in}{0.587838in}}%
\pgfpathlineto{\pgfqpoint{3.750157in}{0.601344in}}%
\pgfpathlineto{\pgfqpoint{3.752254in}{0.606847in}}%
\pgfpathlineto{\pgfqpoint{3.754352in}{0.606995in}}%
\pgfpathlineto{\pgfqpoint{3.756449in}{0.601903in}}%
\pgfpathlineto{\pgfqpoint{3.758547in}{0.589372in}}%
\pgfpathlineto{\pgfqpoint{3.762741in}{0.535919in}}%
\pgfpathlineto{\pgfqpoint{3.764839in}{0.588575in}}%
\pgfpathlineto{\pgfqpoint{3.766936in}{0.609337in}}%
\pgfpathlineto{\pgfqpoint{3.769033in}{0.621017in}}%
\pgfpathlineto{\pgfqpoint{3.771131in}{0.627228in}}%
\pgfpathlineto{\pgfqpoint{3.773228in}{0.628561in}}%
\pgfpathlineto{\pgfqpoint{3.775325in}{0.623812in}}%
\pgfpathlineto{\pgfqpoint{3.777423in}{0.613222in}}%
\pgfpathlineto{\pgfqpoint{3.779520in}{0.594919in}}%
\pgfpathlineto{\pgfqpoint{3.781618in}{0.555193in}}%
\pgfpathlineto{\pgfqpoint{3.783715in}{0.545737in}}%
\pgfpathlineto{\pgfqpoint{3.785812in}{0.586335in}}%
\pgfpathlineto{\pgfqpoint{3.787910in}{0.601539in}}%
\pgfpathlineto{\pgfqpoint{3.790007in}{0.608471in}}%
\pgfpathlineto{\pgfqpoint{3.792104in}{0.610283in}}%
\pgfpathlineto{\pgfqpoint{3.794202in}{0.607590in}}%
\pgfpathlineto{\pgfqpoint{3.796299in}{0.599482in}}%
\pgfpathlineto{\pgfqpoint{3.798396in}{0.581931in}}%
\pgfpathlineto{\pgfqpoint{3.800494in}{0.527494in}}%
\pgfpathlineto{\pgfqpoint{3.804689in}{0.598354in}}%
\pgfpathlineto{\pgfqpoint{3.806786in}{0.614728in}}%
\pgfpathlineto{\pgfqpoint{3.808883in}{0.624103in}}%
\pgfpathlineto{\pgfqpoint{3.810981in}{0.628571in}}%
\pgfpathlineto{\pgfqpoint{3.813078in}{0.628105in}}%
\pgfpathlineto{\pgfqpoint{3.815175in}{0.621314in}}%
\pgfpathlineto{\pgfqpoint{3.817273in}{0.608168in}}%
\pgfpathlineto{\pgfqpoint{3.819370in}{0.583449in}}%
\pgfpathlineto{\pgfqpoint{3.821467in}{0.470314in}}%
\pgfpathlineto{\pgfqpoint{3.823565in}{0.576517in}}%
\pgfpathlineto{\pgfqpoint{3.825662in}{0.600301in}}%
\pgfpathlineto{\pgfqpoint{3.827760in}{0.611382in}}%
\pgfpathlineto{\pgfqpoint{3.829857in}{0.616440in}}%
\pgfpathlineto{\pgfqpoint{3.831954in}{0.617188in}}%
\pgfpathlineto{\pgfqpoint{3.834052in}{0.613818in}}%
\pgfpathlineto{\pgfqpoint{3.836149in}{0.605252in}}%
\pgfpathlineto{\pgfqpoint{3.838246in}{0.587391in}}%
\pgfpathlineto{\pgfqpoint{3.840344in}{0.532883in}}%
\pgfpathlineto{\pgfqpoint{3.844538in}{0.603236in}}%
\pgfpathlineto{\pgfqpoint{3.846636in}{0.619584in}}%
\pgfpathlineto{\pgfqpoint{3.848733in}{0.629002in}}%
\pgfpathlineto{\pgfqpoint{3.850831in}{0.633618in}}%
\pgfpathlineto{\pgfqpoint{3.852928in}{0.633507in}}%
\pgfpathlineto{\pgfqpoint{3.855025in}{0.627303in}}%
\pgfpathlineto{\pgfqpoint{3.857123in}{0.614991in}}%
\pgfpathlineto{\pgfqpoint{3.859220in}{0.592476in}}%
\pgfpathlineto{\pgfqpoint{3.861317in}{0.522780in}}%
\pgfpathlineto{\pgfqpoint{3.863415in}{0.575188in}}%
\pgfpathlineto{\pgfqpoint{3.865512in}{0.602118in}}%
\pgfpathlineto{\pgfqpoint{3.867609in}{0.614181in}}%
\pgfpathlineto{\pgfqpoint{3.869707in}{0.619735in}}%
\pgfpathlineto{\pgfqpoint{3.871804in}{0.620801in}}%
\pgfpathlineto{\pgfqpoint{3.873902in}{0.617679in}}%
\pgfpathlineto{\pgfqpoint{3.875999in}{0.609361in}}%
\pgfpathlineto{\pgfqpoint{3.878096in}{0.591896in}}%
\pgfpathlineto{\pgfqpoint{3.880194in}{0.539807in}}%
\pgfpathlineto{\pgfqpoint{3.884388in}{0.606738in}}%
\pgfpathlineto{\pgfqpoint{3.886486in}{0.623313in}}%
\pgfpathlineto{\pgfqpoint{3.888583in}{0.632847in}}%
\pgfpathlineto{\pgfqpoint{3.890680in}{0.637534in}}%
\pgfpathlineto{\pgfqpoint{3.892778in}{0.637463in}}%
\pgfpathlineto{\pgfqpoint{3.894875in}{0.631334in}}%
\pgfpathlineto{\pgfqpoint{3.896973in}{0.619718in}}%
\pgfpathlineto{\pgfqpoint{3.899070in}{0.599799in}}%
\pgfpathlineto{\pgfqpoint{3.901167in}{0.553664in}}%
\pgfpathlineto{\pgfqpoint{3.903265in}{0.561362in}}%
\pgfpathlineto{\pgfqpoint{3.905362in}{0.595812in}}%
\pgfpathlineto{\pgfqpoint{3.907459in}{0.609373in}}%
\pgfpathlineto{\pgfqpoint{3.909557in}{0.615205in}}%
\pgfpathlineto{\pgfqpoint{3.911654in}{0.615926in}}%
\pgfpathlineto{\pgfqpoint{3.913751in}{0.611819in}}%
\pgfpathlineto{\pgfqpoint{3.915849in}{0.601317in}}%
\pgfpathlineto{\pgfqpoint{3.917946in}{0.577510in}}%
\pgfpathlineto{\pgfqpoint{3.920043in}{0.490395in}}%
\pgfpathlineto{\pgfqpoint{3.922141in}{0.587761in}}%
\pgfpathlineto{\pgfqpoint{3.924238in}{0.612395in}}%
\pgfpathlineto{\pgfqpoint{3.926336in}{0.625809in}}%
\pgfpathlineto{\pgfqpoint{3.928433in}{0.633296in}}%
\pgfpathlineto{\pgfqpoint{3.930530in}{0.635996in}}%
\pgfpathlineto{\pgfqpoint{3.932628in}{0.633177in}}%
\pgfpathlineto{\pgfqpoint{3.934725in}{0.624465in}}%
\pgfpathlineto{\pgfqpoint{3.936822in}{0.610195in}}%
\pgfpathlineto{\pgfqpoint{3.938920in}{0.585104in}}%
\pgfpathlineto{\pgfqpoint{3.941017in}{0.495051in}}%
\pgfpathlineto{\pgfqpoint{3.943114in}{0.571735in}}%
\pgfpathlineto{\pgfqpoint{3.945212in}{0.595285in}}%
\pgfpathlineto{\pgfqpoint{3.947309in}{0.605111in}}%
\pgfpathlineto{\pgfqpoint{3.949407in}{0.608279in}}%
\pgfpathlineto{\pgfqpoint{3.951504in}{0.606210in}}%
\pgfpathlineto{\pgfqpoint{3.953601in}{0.598145in}}%
\pgfpathlineto{\pgfqpoint{3.955699in}{0.579560in}}%
\pgfpathlineto{\pgfqpoint{3.957796in}{0.511016in}}%
\pgfpathlineto{\pgfqpoint{3.959893in}{0.571473in}}%
\pgfpathlineto{\pgfqpoint{3.961991in}{0.601756in}}%
\pgfpathlineto{\pgfqpoint{3.964088in}{0.617290in}}%
\pgfpathlineto{\pgfqpoint{3.966185in}{0.626107in}}%
\pgfpathlineto{\pgfqpoint{3.968283in}{0.629952in}}%
\pgfpathlineto{\pgfqpoint{3.970380in}{0.628495in}}%
\pgfpathlineto{\pgfqpoint{3.972478in}{0.620707in}}%
\pgfpathlineto{\pgfqpoint{3.974575in}{0.607309in}}%
\pgfpathlineto{\pgfqpoint{3.976672in}{0.583660in}}%
\pgfpathlineto{\pgfqpoint{3.978770in}{0.509584in}}%
\pgfpathlineto{\pgfqpoint{3.980867in}{0.566360in}}%
\pgfpathlineto{\pgfqpoint{3.982964in}{0.591920in}}%
\pgfpathlineto{\pgfqpoint{3.985062in}{0.602831in}}%
\pgfpathlineto{\pgfqpoint{3.987159in}{0.607053in}}%
\pgfpathlineto{\pgfqpoint{3.989256in}{0.606340in}}%
\pgfpathlineto{\pgfqpoint{3.991354in}{0.600476in}}%
\pgfpathlineto{\pgfqpoint{3.993451in}{0.586903in}}%
\pgfpathlineto{\pgfqpoint{3.995549in}{0.552808in}}%
\pgfpathlineto{\pgfqpoint{3.997646in}{0.542936in}}%
\pgfpathlineto{\pgfqpoint{3.999743in}{0.589568in}}%
\pgfpathlineto{\pgfqpoint{4.001841in}{0.609244in}}%
\pgfpathlineto{\pgfqpoint{4.003938in}{0.620391in}}%
\pgfpathlineto{\pgfqpoint{4.006035in}{0.626230in}}%
\pgfpathlineto{\pgfqpoint{4.008133in}{0.627223in}}%
\pgfpathlineto{\pgfqpoint{4.010230in}{0.622101in}}%
\pgfpathlineto{\pgfqpoint{4.012327in}{0.610874in}}%
\pgfpathlineto{\pgfqpoint{4.014425in}{0.590717in}}%
\pgfpathlineto{\pgfqpoint{4.016522in}{0.539871in}}%
\pgfpathlineto{\pgfqpoint{4.018620in}{0.560303in}}%
\pgfpathlineto{\pgfqpoint{4.020717in}{0.592392in}}%
\pgfpathlineto{\pgfqpoint{4.022814in}{0.605973in}}%
\pgfpathlineto{\pgfqpoint{4.024912in}{0.612375in}}%
\pgfpathlineto{\pgfqpoint{4.027009in}{0.614103in}}%
\pgfpathlineto{\pgfqpoint{4.029106in}{0.611662in}}%
\pgfpathlineto{\pgfqpoint{4.031204in}{0.604282in}}%
\pgfpathlineto{\pgfqpoint{4.033301in}{0.588745in}}%
\pgfpathlineto{\pgfqpoint{4.035398in}{0.547950in}}%
\pgfpathlineto{\pgfqpoint{4.037496in}{0.556936in}}%
\pgfpathlineto{\pgfqpoint{4.039593in}{0.596512in}}%
\pgfpathlineto{\pgfqpoint{4.041691in}{0.614621in}}%
\pgfpathlineto{\pgfqpoint{4.043788in}{0.625002in}}%
\pgfpathlineto{\pgfqpoint{4.045885in}{0.630373in}}%
\pgfpathlineto{\pgfqpoint{4.047983in}{0.631070in}}%
\pgfpathlineto{\pgfqpoint{4.050080in}{0.625807in}}%
\pgfpathlineto{\pgfqpoint{4.052177in}{0.614284in}}%
\pgfpathlineto{\pgfqpoint{4.054275in}{0.593090in}}%
\pgfpathlineto{\pgfqpoint{4.056372in}{0.533454in}}%
\pgfpathlineto{\pgfqpoint{4.060567in}{0.599989in}}%
\pgfpathlineto{\pgfqpoint{4.062664in}{0.612946in}}%
\pgfpathlineto{\pgfqpoint{4.064762in}{0.619182in}}%
\pgfpathlineto{\pgfqpoint{4.066859in}{0.620956in}}%
\pgfpathlineto{\pgfqpoint{4.068956in}{0.618744in}}%
\pgfpathlineto{\pgfqpoint{4.071054in}{0.611879in}}%
\pgfpathlineto{\pgfqpoint{4.073151in}{0.597600in}}%
\pgfpathlineto{\pgfqpoint{4.075248in}{0.562977in}}%
\pgfpathlineto{\pgfqpoint{4.077346in}{0.552862in}}%
\pgfpathlineto{\pgfqpoint{4.079443in}{0.599215in}}%
\pgfpathlineto{\pgfqpoint{4.081540in}{0.618761in}}%
\pgfpathlineto{\pgfqpoint{4.083638in}{0.629907in}}%
\pgfpathlineto{\pgfqpoint{4.085735in}{0.635915in}}%
\pgfpathlineto{\pgfqpoint{4.087832in}{0.637371in}}%
\pgfpathlineto{\pgfqpoint{4.089930in}{0.633202in}}%
\pgfpathlineto{\pgfqpoint{4.092027in}{0.623171in}}%
\pgfpathlineto{\pgfqpoint{4.094125in}{0.605782in}}%
\pgfpathlineto{\pgfqpoint{4.096222in}{0.569380in}}%
\pgfpathlineto{\pgfqpoint{4.098319in}{0.547953in}}%
\pgfpathlineto{\pgfqpoint{4.100417in}{0.593805in}}%
\pgfpathlineto{\pgfqpoint{4.102514in}{0.610075in}}%
\pgfpathlineto{\pgfqpoint{4.104611in}{0.617609in}}%
\pgfpathlineto{\pgfqpoint{4.106709in}{0.619939in}}%
\pgfpathlineto{\pgfqpoint{4.108806in}{0.617843in}}%
\pgfpathlineto{\pgfqpoint{4.110903in}{0.610639in}}%
\pgfpathlineto{\pgfqpoint{4.113001in}{0.595073in}}%
\pgfpathlineto{\pgfqpoint{4.115098in}{0.553240in}}%
\pgfpathlineto{\pgfqpoint{4.117196in}{0.565315in}}%
\pgfpathlineto{\pgfqpoint{4.119293in}{0.604081in}}%
\pgfpathlineto{\pgfqpoint{4.121390in}{0.622007in}}%
\pgfpathlineto{\pgfqpoint{4.123488in}{0.632263in}}%
\pgfpathlineto{\pgfqpoint{4.125585in}{0.637473in}}%
\pgfpathlineto{\pgfqpoint{4.127682in}{0.637894in}}%
\pgfpathlineto{\pgfqpoint{4.129780in}{0.632262in}}%
\pgfpathlineto{\pgfqpoint{4.131877in}{0.621404in}}%
\pgfpathlineto{\pgfqpoint{4.133974in}{0.603476in}}%
\pgfpathlineto{\pgfqpoint{4.138169in}{0.540382in}}%
\pgfpathlineto{\pgfqpoint{4.140267in}{0.587747in}}%
\pgfpathlineto{\pgfqpoint{4.142364in}{0.603460in}}%
\pgfpathlineto{\pgfqpoint{4.144461in}{0.610014in}}%
\pgfpathlineto{\pgfqpoint{4.146559in}{0.610902in}}%
\pgfpathlineto{\pgfqpoint{4.148656in}{0.606540in}}%
\pgfpathlineto{\pgfqpoint{4.150753in}{0.595098in}}%
\pgfpathlineto{\pgfqpoint{4.152851in}{0.567674in}}%
\pgfpathlineto{\pgfqpoint{4.154948in}{0.526111in}}%
\pgfpathlineto{\pgfqpoint{4.157045in}{0.589570in}}%
\pgfpathlineto{\pgfqpoint{4.159143in}{0.611814in}}%
\pgfpathlineto{\pgfqpoint{4.161240in}{0.624192in}}%
\pgfpathlineto{\pgfqpoint{4.163338in}{0.630906in}}%
\pgfpathlineto{\pgfqpoint{4.165435in}{0.632743in}}%
\pgfpathlineto{\pgfqpoint{4.167532in}{0.628621in}}%
\pgfpathlineto{\pgfqpoint{4.169630in}{0.618940in}}%
\pgfpathlineto{\pgfqpoint{4.171727in}{0.603080in}}%
\pgfpathlineto{\pgfqpoint{4.173824in}{0.573715in}}%
\pgfpathlineto{\pgfqpoint{4.175922in}{0.492587in}}%
\pgfpathlineto{\pgfqpoint{4.178019in}{0.576036in}}%
\pgfpathlineto{\pgfqpoint{4.180116in}{0.595296in}}%
\pgfpathlineto{\pgfqpoint{4.182214in}{0.603469in}}%
\pgfpathlineto{\pgfqpoint{4.184311in}{0.605547in}}%
\pgfpathlineto{\pgfqpoint{4.186409in}{0.602437in}}%
\pgfpathlineto{\pgfqpoint{4.188506in}{0.592929in}}%
\pgfpathlineto{\pgfqpoint{4.190603in}{0.570951in}}%
\pgfpathlineto{\pgfqpoint{4.192701in}{0.409298in}}%
\pgfpathlineto{\pgfqpoint{4.194798in}{0.576375in}}%
\pgfpathlineto{\pgfqpoint{4.196895in}{0.602448in}}%
\pgfpathlineto{\pgfqpoint{4.198993in}{0.616449in}}%
\pgfpathlineto{\pgfqpoint{4.201090in}{0.624288in}}%
\pgfpathlineto{\pgfqpoint{4.203187in}{0.627236in}}%
\pgfpathlineto{\pgfqpoint{4.205285in}{0.624608in}}%
\pgfpathlineto{\pgfqpoint{4.207382in}{0.615690in}}%
\pgfpathlineto{\pgfqpoint{4.209480in}{0.600202in}}%
\pgfpathlineto{\pgfqpoint{4.211577in}{0.570136in}}%
\pgfpathlineto{\pgfqpoint{4.213674in}{0.509046in}}%
\pgfpathlineto{\pgfqpoint{4.215772in}{0.578584in}}%
\pgfpathlineto{\pgfqpoint{4.217869in}{0.597761in}}%
\pgfpathlineto{\pgfqpoint{4.219966in}{0.606637in}}%
\pgfpathlineto{\pgfqpoint{4.222064in}{0.609921in}}%
\pgfpathlineto{\pgfqpoint{4.224161in}{0.608751in}}%
\pgfpathlineto{\pgfqpoint{4.226258in}{0.602756in}}%
\pgfpathlineto{\pgfqpoint{4.228356in}{0.589493in}}%
\pgfpathlineto{\pgfqpoint{4.230453in}{0.557631in}}%
\pgfpathlineto{\pgfqpoint{4.232551in}{0.537711in}}%
\pgfpathlineto{\pgfqpoint{4.234648in}{0.588689in}}%
\pgfpathlineto{\pgfqpoint{4.236745in}{0.609095in}}%
\pgfpathlineto{\pgfqpoint{4.238843in}{0.620624in}}%
\pgfpathlineto{\pgfqpoint{4.240940in}{0.626821in}}%
\pgfpathlineto{\pgfqpoint{4.243037in}{0.628316in}}%
\pgfpathlineto{\pgfqpoint{4.245135in}{0.623982in}}%
\pgfpathlineto{\pgfqpoint{4.247232in}{0.613312in}}%
\pgfpathlineto{\pgfqpoint{4.249329in}{0.593749in}}%
\pgfpathlineto{\pgfqpoint{4.251427in}{0.544610in}}%
\pgfpathlineto{\pgfqpoint{4.253524in}{0.562441in}}%
\pgfpathlineto{\pgfqpoint{4.255621in}{0.595527in}}%
\pgfpathlineto{\pgfqpoint{4.257719in}{0.609647in}}%
\pgfpathlineto{\pgfqpoint{4.259816in}{0.616578in}}%
\pgfpathlineto{\pgfqpoint{4.261914in}{0.618936in}}%
\pgfpathlineto{\pgfqpoint{4.264011in}{0.617364in}}%
\pgfpathlineto{\pgfqpoint{4.266108in}{0.611390in}}%
\pgfpathlineto{\pgfqpoint{4.268206in}{0.598814in}}%
\pgfpathlineto{\pgfqpoint{4.270303in}{0.570490in}}%
\pgfpathlineto{\pgfqpoint{4.272400in}{0.529213in}}%
\pgfpathlineto{\pgfqpoint{4.274498in}{0.591766in}}%
\pgfpathlineto{\pgfqpoint{4.276595in}{0.613731in}}%
\pgfpathlineto{\pgfqpoint{4.278692in}{0.626030in}}%
\pgfpathlineto{\pgfqpoint{4.280790in}{0.632892in}}%
\pgfpathlineto{\pgfqpoint{4.282887in}{0.635236in}}%
\pgfpathlineto{\pgfqpoint{4.284985in}{0.632313in}}%
\pgfpathlineto{\pgfqpoint{4.289179in}{0.615408in}}%
\pgfpathlineto{\pgfqpoint{4.293374in}{0.589430in}}%
\pgfpathlineto{\pgfqpoint{4.297569in}{0.551389in}}%
\pgfpathlineto{\pgfqpoint{4.299666in}{0.531114in}}%
\pgfpathlineto{\pgfqpoint{4.301763in}{0.523148in}}%
\pgfpathlineto{\pgfqpoint{4.303861in}{0.534327in}}%
\pgfpathlineto{\pgfqpoint{4.308056in}{0.567592in}}%
\pgfpathlineto{\pgfqpoint{4.310153in}{0.579213in}}%
\pgfpathlineto{\pgfqpoint{4.312250in}{0.586774in}}%
\pgfpathlineto{\pgfqpoint{4.314348in}{0.589959in}}%
\pgfpathlineto{\pgfqpoint{4.316445in}{0.587339in}}%
\pgfpathlineto{\pgfqpoint{4.316445in}{0.587339in}}%
\pgfusepath{stroke}%
\end{pgfscope}%
\begin{pgfscope}%
\pgfpathrectangle{\pgfqpoint{2.554660in}{0.320625in}}{\pgfqpoint{1.845679in}{1.164375in}}%
\pgfusepath{clip}%
\pgfsetbuttcap%
\pgfsetroundjoin%
\pgfsetlinewidth{0.803000pt}%
\definecolor{currentstroke}{rgb}{0.000000,0.000000,0.000000}%
\pgfsetstrokecolor{currentstroke}%
\pgfsetstrokeopacity{0.700000}%
\pgfsetdash{{2.960000pt}{1.280000pt}}{0.000000pt}%
\pgfpathmoveto{\pgfqpoint{2.638555in}{1.468525in}}%
\pgfpathlineto{\pgfqpoint{3.284093in}{0.315625in}}%
\pgfpathlineto{\pgfqpoint{3.284093in}{0.315625in}}%
\pgfusepath{stroke}%
\end{pgfscope}%
\begin{pgfscope}%
\pgfsetrectcap%
\pgfsetmiterjoin%
\pgfsetlinewidth{0.803000pt}%
\definecolor{currentstroke}{rgb}{0.000000,0.000000,0.000000}%
\pgfsetstrokecolor{currentstroke}%
\pgfsetdash{}{0pt}%
\pgfpathmoveto{\pgfqpoint{2.554660in}{0.320625in}}%
\pgfpathlineto{\pgfqpoint{2.554660in}{1.485000in}}%
\pgfusepath{stroke}%
\end{pgfscope}%
\begin{pgfscope}%
\pgfsetrectcap%
\pgfsetmiterjoin%
\pgfsetlinewidth{0.803000pt}%
\definecolor{currentstroke}{rgb}{0.000000,0.000000,0.000000}%
\pgfsetstrokecolor{currentstroke}%
\pgfsetdash{}{0pt}%
\pgfpathmoveto{\pgfqpoint{4.400340in}{0.320625in}}%
\pgfpathlineto{\pgfqpoint{4.400340in}{1.485000in}}%
\pgfusepath{stroke}%
\end{pgfscope}%
\begin{pgfscope}%
\pgfsetrectcap%
\pgfsetmiterjoin%
\pgfsetlinewidth{0.803000pt}%
\definecolor{currentstroke}{rgb}{0.000000,0.000000,0.000000}%
\pgfsetstrokecolor{currentstroke}%
\pgfsetdash{}{0pt}%
\pgfpathmoveto{\pgfqpoint{2.554660in}{0.320625in}}%
\pgfpathlineto{\pgfqpoint{4.400340in}{0.320625in}}%
\pgfusepath{stroke}%
\end{pgfscope}%
\begin{pgfscope}%
\pgfsetrectcap%
\pgfsetmiterjoin%
\pgfsetlinewidth{0.803000pt}%
\definecolor{currentstroke}{rgb}{0.000000,0.000000,0.000000}%
\pgfsetstrokecolor{currentstroke}%
\pgfsetdash{}{0pt}%
\pgfpathmoveto{\pgfqpoint{2.554660in}{1.485000in}}%
\pgfpathlineto{\pgfqpoint{4.400340in}{1.485000in}}%
\pgfusepath{stroke}%
\end{pgfscope}%
\begin{pgfscope}%
\definecolor{textcolor}{rgb}{0.000000,0.000000,0.000000}%
\pgfsetstrokecolor{textcolor}%
\pgfsetfillcolor{textcolor}%
\pgftext[x=3.477500in,y=1.568333in,,base]{\color{textcolor}{\rmfamily\fontsize{8.000000}{9.600000}\selectfont\catcode`\^=\active\def^{\ifmmode\sp\else\^{}\fi}\catcode`\%=\active\def%{\%}(b) Covariance conditioning}}%
\end{pgfscope}%
\begin{pgfscope}%
\pgfsetbuttcap%
\pgfsetmiterjoin%
\definecolor{currentfill}{rgb}{1.000000,1.000000,1.000000}%
\pgfsetfillcolor{currentfill}%
\pgfsetfillopacity{0.800000}%
\pgfsetlinewidth{1.003750pt}%
\definecolor{currentstroke}{rgb}{0.800000,0.800000,0.800000}%
\pgfsetstrokecolor{currentstroke}%
\pgfsetstrokeopacity{0.800000}%
\pgfsetdash{}{0pt}%
\pgfpathmoveto{\pgfqpoint{3.905118in}{1.210066in}}%
\pgfpathlineto{\pgfqpoint{4.355617in}{1.210066in}}%
\pgfpathquadraticcurveto{\pgfqpoint{4.368395in}{1.210066in}}{\pgfqpoint{4.368395in}{1.222844in}}%
\pgfpathlineto{\pgfqpoint{4.368395in}{1.440278in}}%
\pgfpathquadraticcurveto{\pgfqpoint{4.368395in}{1.453056in}}{\pgfqpoint{4.355617in}{1.453056in}}%
\pgfpathlineto{\pgfqpoint{3.905118in}{1.453056in}}%
\pgfpathquadraticcurveto{\pgfqpoint{3.892340in}{1.453056in}}{\pgfqpoint{3.892340in}{1.440278in}}%
\pgfpathlineto{\pgfqpoint{3.892340in}{1.222844in}}%
\pgfpathquadraticcurveto{\pgfqpoint{3.892340in}{1.210066in}}{\pgfqpoint{3.905118in}{1.210066in}}%
\pgfpathlineto{\pgfqpoint{3.905118in}{1.210066in}}%
\pgfpathclose%
\pgfusepath{stroke,fill}%
\end{pgfscope}%
\begin{pgfscope}%
\pgfsetrectcap%
\pgfsetroundjoin%
\pgfsetlinewidth{1.154312pt}%
\definecolor{currentstroke}{rgb}{0.146785,0.068738,0.334011}%
\pgfsetstrokecolor{currentstroke}%
\pgfsetdash{}{0pt}%
\pgfpathmoveto{\pgfqpoint{3.908312in}{1.410904in}}%
\pgfpathlineto{\pgfqpoint{3.953034in}{1.410904in}}%
\pgfpathlineto{\pgfqpoint{3.997757in}{1.410904in}}%
\pgfusepath{stroke}%
\end{pgfscope}%
\begin{pgfscope}%
\definecolor{textcolor}{rgb}{0.000000,0.000000,0.000000}%
\pgfsetstrokecolor{textcolor}%
\pgfsetfillcolor{textcolor}%
\pgftext[x=4.048868in,y=1.388543in,left,base]{\color{textcolor}{\rmfamily\fontsize{4.600000}{5.520000}\selectfont\catcode`\^=\active\def^{\ifmmode\sp\else\^{}\fi}\catcode`\%=\active\def%{\%}$r=0$}}%
\end{pgfscope}%
\begin{pgfscope}%
\pgfsetrectcap%
\pgfsetroundjoin%
\pgfsetlinewidth{1.154312pt}%
\definecolor{currentstroke}{rgb}{0.716387,0.214982,0.475290}%
\pgfsetstrokecolor{currentstroke}%
\pgfsetdash{}{0pt}%
\pgfpathmoveto{\pgfqpoint{3.908312in}{1.336296in}}%
\pgfpathlineto{\pgfqpoint{3.953034in}{1.336296in}}%
\pgfpathlineto{\pgfqpoint{3.997757in}{1.336296in}}%
\pgfusepath{stroke}%
\end{pgfscope}%
\begin{pgfscope}%
\definecolor{textcolor}{rgb}{0.000000,0.000000,0.000000}%
\pgfsetstrokecolor{textcolor}%
\pgfsetfillcolor{textcolor}%
\pgftext[x=4.048868in,y=1.313935in,left,base]{\color{textcolor}{\rmfamily\fontsize{4.600000}{5.520000}\selectfont\catcode`\^=\active\def^{\ifmmode\sp\else\^{}\fi}\catcode`\%=\active\def%{\%}$r=0.5$}}%
\end{pgfscope}%
\begin{pgfscope}%
\pgfsetrectcap%
\pgfsetroundjoin%
\pgfsetlinewidth{1.154312pt}%
\definecolor{currentstroke}{rgb}{0.997351,0.719089,0.493755}%
\pgfsetstrokecolor{currentstroke}%
\pgfsetdash{}{0pt}%
\pgfpathmoveto{\pgfqpoint{3.908312in}{1.261688in}}%
\pgfpathlineto{\pgfqpoint{3.953034in}{1.261688in}}%
\pgfpathlineto{\pgfqpoint{3.997757in}{1.261688in}}%
\pgfusepath{stroke}%
\end{pgfscope}%
\begin{pgfscope}%
\definecolor{textcolor}{rgb}{0.000000,0.000000,0.000000}%
\pgfsetstrokecolor{textcolor}%
\pgfsetfillcolor{textcolor}%
\pgftext[x=4.048868in,y=1.239327in,left,base]{\color{textcolor}{\rmfamily\fontsize{4.600000}{5.520000}\selectfont\catcode`\^=\active\def^{\ifmmode\sp\else\^{}\fi}\catcode`\%=\active\def%{\%}$r=1$}}%
\end{pgfscope}%
\begin{pgfscope}%
\pgfsetbuttcap%
\pgfsetmiterjoin%
\definecolor{currentfill}{rgb}{1.000000,1.000000,1.000000}%
\pgfsetfillcolor{currentfill}%
\pgfsetlinewidth{0.000000pt}%
\definecolor{currentstroke}{rgb}{0.000000,0.000000,0.000000}%
\pgfsetstrokecolor{currentstroke}%
\pgfsetstrokeopacity{0.000000}%
\pgfsetdash{}{0pt}%
\pgfpathmoveto{\pgfqpoint{4.621821in}{0.320625in}}%
\pgfpathlineto{\pgfqpoint{6.467500in}{0.320625in}}%
\pgfpathlineto{\pgfqpoint{6.467500in}{1.485000in}}%
\pgfpathlineto{\pgfqpoint{4.621821in}{1.485000in}}%
\pgfpathlineto{\pgfqpoint{4.621821in}{0.320625in}}%
\pgfpathclose%
\pgfusepath{fill}%
\end{pgfscope}%
\begin{pgfscope}%
\pgfpathrectangle{\pgfqpoint{4.621821in}{0.320625in}}{\pgfqpoint{1.845679in}{1.164375in}}%
\pgfusepath{clip}%
\pgfsetrectcap%
\pgfsetroundjoin%
\pgfsetlinewidth{0.803000pt}%
\definecolor{currentstroke}{rgb}{0.909804,0.909804,0.909804}%
\pgfsetstrokecolor{currentstroke}%
\pgfsetdash{}{0pt}%
\pgfpathmoveto{\pgfqpoint{4.705715in}{0.320625in}}%
\pgfpathlineto{\pgfqpoint{4.705715in}{1.485000in}}%
\pgfusepath{stroke}%
\end{pgfscope}%
\begin{pgfscope}%
\pgfsetbuttcap%
\pgfsetroundjoin%
\definecolor{currentfill}{rgb}{0.000000,0.000000,0.000000}%
\pgfsetfillcolor{currentfill}%
\pgfsetlinewidth{0.803000pt}%
\definecolor{currentstroke}{rgb}{0.000000,0.000000,0.000000}%
\pgfsetstrokecolor{currentstroke}%
\pgfsetdash{}{0pt}%
\pgfsys@defobject{currentmarker}{\pgfqpoint{0.000000in}{-0.048611in}}{\pgfqpoint{0.000000in}{0.000000in}}{%
\pgfpathmoveto{\pgfqpoint{0.000000in}{0.000000in}}%
\pgfpathlineto{\pgfqpoint{0.000000in}{-0.048611in}}%
\pgfusepath{stroke,fill}%
}%
\begin{pgfscope}%
\pgfsys@transformshift{4.705715in}{0.320625in}%
\pgfsys@useobject{currentmarker}{}%
\end{pgfscope}%
\end{pgfscope}%
\begin{pgfscope}%
\definecolor{textcolor}{rgb}{0.000000,0.000000,0.000000}%
\pgfsetstrokecolor{textcolor}%
\pgfsetfillcolor{textcolor}%
\pgftext[x=4.705715in,y=0.223403in,,top]{\color{textcolor}{\rmfamily\fontsize{6.000000}{7.200000}\selectfont\catcode`\^=\active\def^{\ifmmode\sp\else\^{}\fi}\catcode`\%=\active\def%{\%}0}}%
\end{pgfscope}%
\begin{pgfscope}%
\pgfpathrectangle{\pgfqpoint{4.621821in}{0.320625in}}{\pgfqpoint{1.845679in}{1.164375in}}%
\pgfusepath{clip}%
\pgfsetrectcap%
\pgfsetroundjoin%
\pgfsetlinewidth{0.803000pt}%
\definecolor{currentstroke}{rgb}{0.909804,0.909804,0.909804}%
\pgfsetstrokecolor{currentstroke}%
\pgfsetdash{}{0pt}%
\pgfpathmoveto{\pgfqpoint{5.125188in}{0.320625in}}%
\pgfpathlineto{\pgfqpoint{5.125188in}{1.485000in}}%
\pgfusepath{stroke}%
\end{pgfscope}%
\begin{pgfscope}%
\pgfsetbuttcap%
\pgfsetroundjoin%
\definecolor{currentfill}{rgb}{0.000000,0.000000,0.000000}%
\pgfsetfillcolor{currentfill}%
\pgfsetlinewidth{0.803000pt}%
\definecolor{currentstroke}{rgb}{0.000000,0.000000,0.000000}%
\pgfsetstrokecolor{currentstroke}%
\pgfsetdash{}{0pt}%
\pgfsys@defobject{currentmarker}{\pgfqpoint{0.000000in}{-0.048611in}}{\pgfqpoint{0.000000in}{0.000000in}}{%
\pgfpathmoveto{\pgfqpoint{0.000000in}{0.000000in}}%
\pgfpathlineto{\pgfqpoint{0.000000in}{-0.048611in}}%
\pgfusepath{stroke,fill}%
}%
\begin{pgfscope}%
\pgfsys@transformshift{5.125188in}{0.320625in}%
\pgfsys@useobject{currentmarker}{}%
\end{pgfscope}%
\end{pgfscope}%
\begin{pgfscope}%
\definecolor{textcolor}{rgb}{0.000000,0.000000,0.000000}%
\pgfsetstrokecolor{textcolor}%
\pgfsetfillcolor{textcolor}%
\pgftext[x=5.125188in,y=0.223403in,,top]{\color{textcolor}{\rmfamily\fontsize{6.000000}{7.200000}\selectfont\catcode`\^=\active\def^{\ifmmode\sp\else\^{}\fi}\catcode`\%=\active\def%{\%}10}}%
\end{pgfscope}%
\begin{pgfscope}%
\pgfpathrectangle{\pgfqpoint{4.621821in}{0.320625in}}{\pgfqpoint{1.845679in}{1.164375in}}%
\pgfusepath{clip}%
\pgfsetrectcap%
\pgfsetroundjoin%
\pgfsetlinewidth{0.803000pt}%
\definecolor{currentstroke}{rgb}{0.909804,0.909804,0.909804}%
\pgfsetstrokecolor{currentstroke}%
\pgfsetdash{}{0pt}%
\pgfpathmoveto{\pgfqpoint{5.544660in}{0.320625in}}%
\pgfpathlineto{\pgfqpoint{5.544660in}{1.485000in}}%
\pgfusepath{stroke}%
\end{pgfscope}%
\begin{pgfscope}%
\pgfsetbuttcap%
\pgfsetroundjoin%
\definecolor{currentfill}{rgb}{0.000000,0.000000,0.000000}%
\pgfsetfillcolor{currentfill}%
\pgfsetlinewidth{0.803000pt}%
\definecolor{currentstroke}{rgb}{0.000000,0.000000,0.000000}%
\pgfsetstrokecolor{currentstroke}%
\pgfsetdash{}{0pt}%
\pgfsys@defobject{currentmarker}{\pgfqpoint{0.000000in}{-0.048611in}}{\pgfqpoint{0.000000in}{0.000000in}}{%
\pgfpathmoveto{\pgfqpoint{0.000000in}{0.000000in}}%
\pgfpathlineto{\pgfqpoint{0.000000in}{-0.048611in}}%
\pgfusepath{stroke,fill}%
}%
\begin{pgfscope}%
\pgfsys@transformshift{5.544660in}{0.320625in}%
\pgfsys@useobject{currentmarker}{}%
\end{pgfscope}%
\end{pgfscope}%
\begin{pgfscope}%
\definecolor{textcolor}{rgb}{0.000000,0.000000,0.000000}%
\pgfsetstrokecolor{textcolor}%
\pgfsetfillcolor{textcolor}%
\pgftext[x=5.544660in,y=0.223403in,,top]{\color{textcolor}{\rmfamily\fontsize{6.000000}{7.200000}\selectfont\catcode`\^=\active\def^{\ifmmode\sp\else\^{}\fi}\catcode`\%=\active\def%{\%}20}}%
\end{pgfscope}%
\begin{pgfscope}%
\pgfpathrectangle{\pgfqpoint{4.621821in}{0.320625in}}{\pgfqpoint{1.845679in}{1.164375in}}%
\pgfusepath{clip}%
\pgfsetrectcap%
\pgfsetroundjoin%
\pgfsetlinewidth{0.803000pt}%
\definecolor{currentstroke}{rgb}{0.909804,0.909804,0.909804}%
\pgfsetstrokecolor{currentstroke}%
\pgfsetdash{}{0pt}%
\pgfpathmoveto{\pgfqpoint{5.964133in}{0.320625in}}%
\pgfpathlineto{\pgfqpoint{5.964133in}{1.485000in}}%
\pgfusepath{stroke}%
\end{pgfscope}%
\begin{pgfscope}%
\pgfsetbuttcap%
\pgfsetroundjoin%
\definecolor{currentfill}{rgb}{0.000000,0.000000,0.000000}%
\pgfsetfillcolor{currentfill}%
\pgfsetlinewidth{0.803000pt}%
\definecolor{currentstroke}{rgb}{0.000000,0.000000,0.000000}%
\pgfsetstrokecolor{currentstroke}%
\pgfsetdash{}{0pt}%
\pgfsys@defobject{currentmarker}{\pgfqpoint{0.000000in}{-0.048611in}}{\pgfqpoint{0.000000in}{0.000000in}}{%
\pgfpathmoveto{\pgfqpoint{0.000000in}{0.000000in}}%
\pgfpathlineto{\pgfqpoint{0.000000in}{-0.048611in}}%
\pgfusepath{stroke,fill}%
}%
\begin{pgfscope}%
\pgfsys@transformshift{5.964133in}{0.320625in}%
\pgfsys@useobject{currentmarker}{}%
\end{pgfscope}%
\end{pgfscope}%
\begin{pgfscope}%
\definecolor{textcolor}{rgb}{0.000000,0.000000,0.000000}%
\pgfsetstrokecolor{textcolor}%
\pgfsetfillcolor{textcolor}%
\pgftext[x=5.964133in,y=0.223403in,,top]{\color{textcolor}{\rmfamily\fontsize{6.000000}{7.200000}\selectfont\catcode`\^=\active\def^{\ifmmode\sp\else\^{}\fi}\catcode`\%=\active\def%{\%}30}}%
\end{pgfscope}%
\begin{pgfscope}%
\pgfpathrectangle{\pgfqpoint{4.621821in}{0.320625in}}{\pgfqpoint{1.845679in}{1.164375in}}%
\pgfusepath{clip}%
\pgfsetrectcap%
\pgfsetroundjoin%
\pgfsetlinewidth{0.803000pt}%
\definecolor{currentstroke}{rgb}{0.909804,0.909804,0.909804}%
\pgfsetstrokecolor{currentstroke}%
\pgfsetdash{}{0pt}%
\pgfpathmoveto{\pgfqpoint{6.383605in}{0.320625in}}%
\pgfpathlineto{\pgfqpoint{6.383605in}{1.485000in}}%
\pgfusepath{stroke}%
\end{pgfscope}%
\begin{pgfscope}%
\pgfsetbuttcap%
\pgfsetroundjoin%
\definecolor{currentfill}{rgb}{0.000000,0.000000,0.000000}%
\pgfsetfillcolor{currentfill}%
\pgfsetlinewidth{0.803000pt}%
\definecolor{currentstroke}{rgb}{0.000000,0.000000,0.000000}%
\pgfsetstrokecolor{currentstroke}%
\pgfsetdash{}{0pt}%
\pgfsys@defobject{currentmarker}{\pgfqpoint{0.000000in}{-0.048611in}}{\pgfqpoint{0.000000in}{0.000000in}}{%
\pgfpathmoveto{\pgfqpoint{0.000000in}{0.000000in}}%
\pgfpathlineto{\pgfqpoint{0.000000in}{-0.048611in}}%
\pgfusepath{stroke,fill}%
}%
\begin{pgfscope}%
\pgfsys@transformshift{6.383605in}{0.320625in}%
\pgfsys@useobject{currentmarker}{}%
\end{pgfscope}%
\end{pgfscope}%
\begin{pgfscope}%
\definecolor{textcolor}{rgb}{0.000000,0.000000,0.000000}%
\pgfsetstrokecolor{textcolor}%
\pgfsetfillcolor{textcolor}%
\pgftext[x=6.383605in,y=0.223403in,,top]{\color{textcolor}{\rmfamily\fontsize{6.000000}{7.200000}\selectfont\catcode`\^=\active\def^{\ifmmode\sp\else\^{}\fi}\catcode`\%=\active\def%{\%}40}}%
\end{pgfscope}%
\begin{pgfscope}%
\definecolor{textcolor}{rgb}{0.000000,0.000000,0.000000}%
\pgfsetstrokecolor{textcolor}%
\pgfsetfillcolor{textcolor}%
\pgftext[x=5.544660in,y=0.087200in,,top]{\color{textcolor}{\rmfamily\fontsize{6.000000}{7.200000}\selectfont\catcode`\^=\active\def^{\ifmmode\sp\else\^{}\fi}\catcode`\%=\active\def%{\%}time}}%
\end{pgfscope}%
\begin{pgfscope}%
\pgfpathrectangle{\pgfqpoint{4.621821in}{0.320625in}}{\pgfqpoint{1.845679in}{1.164375in}}%
\pgfusepath{clip}%
\pgfsetrectcap%
\pgfsetroundjoin%
\pgfsetlinewidth{0.803000pt}%
\definecolor{currentstroke}{rgb}{0.909804,0.909804,0.909804}%
\pgfsetstrokecolor{currentstroke}%
\pgfsetdash{}{0pt}%
\pgfpathmoveto{\pgfqpoint{4.621821in}{0.320625in}}%
\pgfpathlineto{\pgfqpoint{6.467500in}{0.320625in}}%
\pgfusepath{stroke}%
\end{pgfscope}%
\begin{pgfscope}%
\pgfsetbuttcap%
\pgfsetroundjoin%
\definecolor{currentfill}{rgb}{0.000000,0.000000,0.000000}%
\pgfsetfillcolor{currentfill}%
\pgfsetlinewidth{0.803000pt}%
\definecolor{currentstroke}{rgb}{0.000000,0.000000,0.000000}%
\pgfsetstrokecolor{currentstroke}%
\pgfsetdash{}{0pt}%
\pgfsys@defobject{currentmarker}{\pgfqpoint{-0.048611in}{0.000000in}}{\pgfqpoint{-0.000000in}{0.000000in}}{%
\pgfpathmoveto{\pgfqpoint{-0.000000in}{0.000000in}}%
\pgfpathlineto{\pgfqpoint{-0.048611in}{0.000000in}}%
\pgfusepath{stroke,fill}%
}%
\begin{pgfscope}%
\pgfsys@transformshift{4.621821in}{0.320625in}%
\pgfsys@useobject{currentmarker}{}%
\end{pgfscope}%
\end{pgfscope}%
\begin{pgfscope}%
\pgfpathrectangle{\pgfqpoint{4.621821in}{0.320625in}}{\pgfqpoint{1.845679in}{1.164375in}}%
\pgfusepath{clip}%
\pgfsetrectcap%
\pgfsetroundjoin%
\pgfsetlinewidth{0.803000pt}%
\definecolor{currentstroke}{rgb}{0.909804,0.909804,0.909804}%
\pgfsetstrokecolor{currentstroke}%
\pgfsetdash{}{0pt}%
\pgfpathmoveto{\pgfqpoint{4.621821in}{0.536250in}}%
\pgfpathlineto{\pgfqpoint{6.467500in}{0.536250in}}%
\pgfusepath{stroke}%
\end{pgfscope}%
\begin{pgfscope}%
\pgfsetbuttcap%
\pgfsetroundjoin%
\definecolor{currentfill}{rgb}{0.000000,0.000000,0.000000}%
\pgfsetfillcolor{currentfill}%
\pgfsetlinewidth{0.803000pt}%
\definecolor{currentstroke}{rgb}{0.000000,0.000000,0.000000}%
\pgfsetstrokecolor{currentstroke}%
\pgfsetdash{}{0pt}%
\pgfsys@defobject{currentmarker}{\pgfqpoint{-0.048611in}{0.000000in}}{\pgfqpoint{-0.000000in}{0.000000in}}{%
\pgfpathmoveto{\pgfqpoint{-0.000000in}{0.000000in}}%
\pgfpathlineto{\pgfqpoint{-0.048611in}{0.000000in}}%
\pgfusepath{stroke,fill}%
}%
\begin{pgfscope}%
\pgfsys@transformshift{4.621821in}{0.536250in}%
\pgfsys@useobject{currentmarker}{}%
\end{pgfscope}%
\end{pgfscope}%
\begin{pgfscope}%
\pgfpathrectangle{\pgfqpoint{4.621821in}{0.320625in}}{\pgfqpoint{1.845679in}{1.164375in}}%
\pgfusepath{clip}%
\pgfsetrectcap%
\pgfsetroundjoin%
\pgfsetlinewidth{0.803000pt}%
\definecolor{currentstroke}{rgb}{0.909804,0.909804,0.909804}%
\pgfsetstrokecolor{currentstroke}%
\pgfsetdash{}{0pt}%
\pgfpathmoveto{\pgfqpoint{4.621821in}{0.751875in}}%
\pgfpathlineto{\pgfqpoint{6.467500in}{0.751875in}}%
\pgfusepath{stroke}%
\end{pgfscope}%
\begin{pgfscope}%
\pgfsetbuttcap%
\pgfsetroundjoin%
\definecolor{currentfill}{rgb}{0.000000,0.000000,0.000000}%
\pgfsetfillcolor{currentfill}%
\pgfsetlinewidth{0.803000pt}%
\definecolor{currentstroke}{rgb}{0.000000,0.000000,0.000000}%
\pgfsetstrokecolor{currentstroke}%
\pgfsetdash{}{0pt}%
\pgfsys@defobject{currentmarker}{\pgfqpoint{-0.048611in}{0.000000in}}{\pgfqpoint{-0.000000in}{0.000000in}}{%
\pgfpathmoveto{\pgfqpoint{-0.000000in}{0.000000in}}%
\pgfpathlineto{\pgfqpoint{-0.048611in}{0.000000in}}%
\pgfusepath{stroke,fill}%
}%
\begin{pgfscope}%
\pgfsys@transformshift{4.621821in}{0.751875in}%
\pgfsys@useobject{currentmarker}{}%
\end{pgfscope}%
\end{pgfscope}%
\begin{pgfscope}%
\pgfpathrectangle{\pgfqpoint{4.621821in}{0.320625in}}{\pgfqpoint{1.845679in}{1.164375in}}%
\pgfusepath{clip}%
\pgfsetrectcap%
\pgfsetroundjoin%
\pgfsetlinewidth{0.803000pt}%
\definecolor{currentstroke}{rgb}{0.909804,0.909804,0.909804}%
\pgfsetstrokecolor{currentstroke}%
\pgfsetdash{}{0pt}%
\pgfpathmoveto{\pgfqpoint{4.621821in}{0.967500in}}%
\pgfpathlineto{\pgfqpoint{6.467500in}{0.967500in}}%
\pgfusepath{stroke}%
\end{pgfscope}%
\begin{pgfscope}%
\pgfsetbuttcap%
\pgfsetroundjoin%
\definecolor{currentfill}{rgb}{0.000000,0.000000,0.000000}%
\pgfsetfillcolor{currentfill}%
\pgfsetlinewidth{0.803000pt}%
\definecolor{currentstroke}{rgb}{0.000000,0.000000,0.000000}%
\pgfsetstrokecolor{currentstroke}%
\pgfsetdash{}{0pt}%
\pgfsys@defobject{currentmarker}{\pgfqpoint{-0.048611in}{0.000000in}}{\pgfqpoint{-0.000000in}{0.000000in}}{%
\pgfpathmoveto{\pgfqpoint{-0.000000in}{0.000000in}}%
\pgfpathlineto{\pgfqpoint{-0.048611in}{0.000000in}}%
\pgfusepath{stroke,fill}%
}%
\begin{pgfscope}%
\pgfsys@transformshift{4.621821in}{0.967500in}%
\pgfsys@useobject{currentmarker}{}%
\end{pgfscope}%
\end{pgfscope}%
\begin{pgfscope}%
\pgfpathrectangle{\pgfqpoint{4.621821in}{0.320625in}}{\pgfqpoint{1.845679in}{1.164375in}}%
\pgfusepath{clip}%
\pgfsetrectcap%
\pgfsetroundjoin%
\pgfsetlinewidth{0.803000pt}%
\definecolor{currentstroke}{rgb}{0.909804,0.909804,0.909804}%
\pgfsetstrokecolor{currentstroke}%
\pgfsetdash{}{0pt}%
\pgfpathmoveto{\pgfqpoint{4.621821in}{1.183125in}}%
\pgfpathlineto{\pgfqpoint{6.467500in}{1.183125in}}%
\pgfusepath{stroke}%
\end{pgfscope}%
\begin{pgfscope}%
\pgfsetbuttcap%
\pgfsetroundjoin%
\definecolor{currentfill}{rgb}{0.000000,0.000000,0.000000}%
\pgfsetfillcolor{currentfill}%
\pgfsetlinewidth{0.803000pt}%
\definecolor{currentstroke}{rgb}{0.000000,0.000000,0.000000}%
\pgfsetstrokecolor{currentstroke}%
\pgfsetdash{}{0pt}%
\pgfsys@defobject{currentmarker}{\pgfqpoint{-0.048611in}{0.000000in}}{\pgfqpoint{-0.000000in}{0.000000in}}{%
\pgfpathmoveto{\pgfqpoint{-0.000000in}{0.000000in}}%
\pgfpathlineto{\pgfqpoint{-0.048611in}{0.000000in}}%
\pgfusepath{stroke,fill}%
}%
\begin{pgfscope}%
\pgfsys@transformshift{4.621821in}{1.183125in}%
\pgfsys@useobject{currentmarker}{}%
\end{pgfscope}%
\end{pgfscope}%
\begin{pgfscope}%
\pgfpathrectangle{\pgfqpoint{4.621821in}{0.320625in}}{\pgfqpoint{1.845679in}{1.164375in}}%
\pgfusepath{clip}%
\pgfsetrectcap%
\pgfsetroundjoin%
\pgfsetlinewidth{0.803000pt}%
\definecolor{currentstroke}{rgb}{0.909804,0.909804,0.909804}%
\pgfsetstrokecolor{currentstroke}%
\pgfsetdash{}{0pt}%
\pgfpathmoveto{\pgfqpoint{4.621821in}{1.398750in}}%
\pgfpathlineto{\pgfqpoint{6.467500in}{1.398750in}}%
\pgfusepath{stroke}%
\end{pgfscope}%
\begin{pgfscope}%
\pgfsetbuttcap%
\pgfsetroundjoin%
\definecolor{currentfill}{rgb}{0.000000,0.000000,0.000000}%
\pgfsetfillcolor{currentfill}%
\pgfsetlinewidth{0.803000pt}%
\definecolor{currentstroke}{rgb}{0.000000,0.000000,0.000000}%
\pgfsetstrokecolor{currentstroke}%
\pgfsetdash{}{0pt}%
\pgfsys@defobject{currentmarker}{\pgfqpoint{-0.048611in}{0.000000in}}{\pgfqpoint{-0.000000in}{0.000000in}}{%
\pgfpathmoveto{\pgfqpoint{-0.000000in}{0.000000in}}%
\pgfpathlineto{\pgfqpoint{-0.048611in}{0.000000in}}%
\pgfusepath{stroke,fill}%
}%
\begin{pgfscope}%
\pgfsys@transformshift{4.621821in}{1.398750in}%
\pgfsys@useobject{currentmarker}{}%
\end{pgfscope}%
\end{pgfscope}%
\begin{pgfscope}%
\pgfpathrectangle{\pgfqpoint{4.621821in}{0.320625in}}{\pgfqpoint{1.845679in}{1.164375in}}%
\pgfusepath{clip}%
\pgfsetrectcap%
\pgfsetroundjoin%
\pgfsetlinewidth{1.304875pt}%
\definecolor{currentstroke}{rgb}{0.000000,0.000000,0.000000}%
\pgfsetstrokecolor{currentstroke}%
\pgfsetdash{}{0pt}%
\pgfpathmoveto{\pgfqpoint{4.705715in}{1.468525in}}%
\pgfpathlineto{\pgfqpoint{4.715783in}{1.449849in}}%
\pgfpathlineto{\pgfqpoint{4.725850in}{1.428582in}}%
\pgfpathlineto{\pgfqpoint{4.739273in}{1.396826in}}%
\pgfpathlineto{\pgfqpoint{4.779543in}{1.299011in}}%
\pgfpathlineto{\pgfqpoint{4.803033in}{1.245684in}}%
\pgfpathlineto{\pgfqpoint{4.829879in}{1.187742in}}%
\pgfpathlineto{\pgfqpoint{4.863437in}{1.118228in}}%
\pgfpathlineto{\pgfqpoint{4.907062in}{1.030672in}}%
\pgfpathlineto{\pgfqpoint{4.967466in}{0.912223in}}%
\pgfpathlineto{\pgfqpoint{5.044649in}{0.763330in}}%
\pgfpathlineto{\pgfqpoint{5.058072in}{0.737306in}}%
\pgfpathlineto{\pgfqpoint{5.068140in}{0.717867in}}%
\pgfpathlineto{\pgfqpoint{5.084919in}{0.687078in}}%
\pgfpathlineto{\pgfqpoint{5.101698in}{0.652063in}}%
\pgfpathlineto{\pgfqpoint{5.105053in}{0.647076in}}%
\pgfpathlineto{\pgfqpoint{5.111765in}{0.640040in}}%
\pgfpathlineto{\pgfqpoint{5.115121in}{0.636691in}}%
\pgfpathlineto{\pgfqpoint{5.118476in}{0.631615in}}%
\pgfpathlineto{\pgfqpoint{5.121832in}{0.624268in}}%
\pgfpathlineto{\pgfqpoint{5.125188in}{0.615244in}}%
\pgfpathlineto{\pgfqpoint{5.128544in}{0.604400in}}%
\pgfpathlineto{\pgfqpoint{5.135255in}{0.577200in}}%
\pgfpathlineto{\pgfqpoint{5.141967in}{0.549758in}}%
\pgfpathlineto{\pgfqpoint{5.145323in}{0.543089in}}%
\pgfpathlineto{\pgfqpoint{5.148678in}{0.545206in}}%
\pgfpathlineto{\pgfqpoint{5.152034in}{0.555371in}}%
\pgfpathlineto{\pgfqpoint{5.158746in}{0.577878in}}%
\pgfpathlineto{\pgfqpoint{5.162102in}{0.584907in}}%
\pgfpathlineto{\pgfqpoint{5.165457in}{0.587993in}}%
\pgfpathlineto{\pgfqpoint{5.168813in}{0.585714in}}%
\pgfpathlineto{\pgfqpoint{5.172169in}{0.575987in}}%
\pgfpathlineto{\pgfqpoint{5.175525in}{0.557412in}}%
\pgfpathlineto{\pgfqpoint{5.178880in}{0.543212in}}%
\pgfpathlineto{\pgfqpoint{5.182236in}{0.565484in}}%
\pgfpathlineto{\pgfqpoint{5.185592in}{0.583410in}}%
\pgfpathlineto{\pgfqpoint{5.188948in}{0.594334in}}%
\pgfpathlineto{\pgfqpoint{5.192304in}{0.600781in}}%
\pgfpathlineto{\pgfqpoint{5.195659in}{0.603926in}}%
\pgfpathlineto{\pgfqpoint{5.199015in}{0.604201in}}%
\pgfpathlineto{\pgfqpoint{5.202371in}{0.601539in}}%
\pgfpathlineto{\pgfqpoint{5.205727in}{0.595294in}}%
\pgfpathlineto{\pgfqpoint{5.209082in}{0.583604in}}%
\pgfpathlineto{\pgfqpoint{5.212438in}{0.560150in}}%
\pgfpathlineto{\pgfqpoint{5.215794in}{0.482906in}}%
\pgfpathlineto{\pgfqpoint{5.219150in}{0.564283in}}%
\pgfpathlineto{\pgfqpoint{5.222506in}{0.589670in}}%
\pgfpathlineto{\pgfqpoint{5.225861in}{0.603589in}}%
\pgfpathlineto{\pgfqpoint{5.229217in}{0.611859in}}%
\pgfpathlineto{\pgfqpoint{5.232573in}{0.615991in}}%
\pgfpathlineto{\pgfqpoint{5.235929in}{0.616022in}}%
\pgfpathlineto{\pgfqpoint{5.239285in}{0.610546in}}%
\pgfpathlineto{\pgfqpoint{5.242640in}{0.598461in}}%
\pgfpathlineto{\pgfqpoint{5.245996in}{0.596439in}}%
\pgfpathlineto{\pgfqpoint{5.249352in}{0.613372in}}%
\pgfpathlineto{\pgfqpoint{5.252708in}{0.627432in}}%
\pgfpathlineto{\pgfqpoint{5.256063in}{0.635926in}}%
\pgfpathlineto{\pgfqpoint{5.259419in}{0.639936in}}%
\pgfpathlineto{\pgfqpoint{5.262775in}{0.639952in}}%
\pgfpathlineto{\pgfqpoint{5.266131in}{0.635649in}}%
\pgfpathlineto{\pgfqpoint{5.269487in}{0.625356in}}%
\pgfpathlineto{\pgfqpoint{5.272842in}{0.602859in}}%
\pgfpathlineto{\pgfqpoint{5.276198in}{0.512108in}}%
\pgfpathlineto{\pgfqpoint{5.279554in}{0.606847in}}%
\pgfpathlineto{\pgfqpoint{5.282910in}{0.631998in}}%
\pgfpathlineto{\pgfqpoint{5.286265in}{0.644790in}}%
\pgfpathlineto{\pgfqpoint{5.289621in}{0.650774in}}%
\pgfpathlineto{\pgfqpoint{5.292977in}{0.650472in}}%
\pgfpathlineto{\pgfqpoint{5.296333in}{0.641618in}}%
\pgfpathlineto{\pgfqpoint{5.299689in}{0.624825in}}%
\pgfpathlineto{\pgfqpoint{5.303044in}{0.599643in}}%
\pgfpathlineto{\pgfqpoint{5.306400in}{0.587013in}}%
\pgfpathlineto{\pgfqpoint{5.309756in}{0.596800in}}%
\pgfpathlineto{\pgfqpoint{5.313112in}{0.596191in}}%
\pgfpathlineto{\pgfqpoint{5.316467in}{0.577447in}}%
\pgfpathlineto{\pgfqpoint{5.319823in}{0.562834in}}%
\pgfpathlineto{\pgfqpoint{5.323179in}{0.603664in}}%
\pgfpathlineto{\pgfqpoint{5.326535in}{0.622199in}}%
\pgfpathlineto{\pgfqpoint{5.329891in}{0.627556in}}%
\pgfpathlineto{\pgfqpoint{5.333246in}{0.618560in}}%
\pgfpathlineto{\pgfqpoint{5.336602in}{0.598332in}}%
\pgfpathlineto{\pgfqpoint{5.339958in}{0.558815in}}%
\pgfpathlineto{\pgfqpoint{5.343314in}{0.558340in}}%
\pgfpathlineto{\pgfqpoint{5.346669in}{0.578882in}}%
\pgfpathlineto{\pgfqpoint{5.350025in}{0.578510in}}%
\pgfpathlineto{\pgfqpoint{5.356737in}{0.535103in}}%
\pgfpathlineto{\pgfqpoint{5.360093in}{0.586830in}}%
\pgfpathlineto{\pgfqpoint{5.363448in}{0.605528in}}%
\pgfpathlineto{\pgfqpoint{5.366804in}{0.610396in}}%
\pgfpathlineto{\pgfqpoint{5.370160in}{0.599243in}}%
\pgfpathlineto{\pgfqpoint{5.373516in}{0.560179in}}%
\pgfpathlineto{\pgfqpoint{5.376871in}{0.564813in}}%
\pgfpathlineto{\pgfqpoint{5.380227in}{0.599722in}}%
\pgfpathlineto{\pgfqpoint{5.383583in}{0.612008in}}%
\pgfpathlineto{\pgfqpoint{5.386939in}{0.614424in}}%
\pgfpathlineto{\pgfqpoint{5.390295in}{0.608176in}}%
\pgfpathlineto{\pgfqpoint{5.393650in}{0.587975in}}%
\pgfpathlineto{\pgfqpoint{5.397006in}{0.447069in}}%
\pgfpathlineto{\pgfqpoint{5.400362in}{0.593784in}}%
\pgfpathlineto{\pgfqpoint{5.403718in}{0.617788in}}%
\pgfpathlineto{\pgfqpoint{5.407074in}{0.627349in}}%
\pgfpathlineto{\pgfqpoint{5.410429in}{0.625618in}}%
\pgfpathlineto{\pgfqpoint{5.413785in}{0.609150in}}%
\pgfpathlineto{\pgfqpoint{5.417141in}{0.563415in}}%
\pgfpathlineto{\pgfqpoint{5.420497in}{0.574230in}}%
\pgfpathlineto{\pgfqpoint{5.423852in}{0.604582in}}%
\pgfpathlineto{\pgfqpoint{5.427208in}{0.613055in}}%
\pgfpathlineto{\pgfqpoint{5.430564in}{0.610130in}}%
\pgfpathlineto{\pgfqpoint{5.433920in}{0.592877in}}%
\pgfpathlineto{\pgfqpoint{5.437276in}{0.488073in}}%
\pgfpathlineto{\pgfqpoint{5.440631in}{0.595125in}}%
\pgfpathlineto{\pgfqpoint{5.443987in}{0.620660in}}%
\pgfpathlineto{\pgfqpoint{5.447343in}{0.630310in}}%
\pgfpathlineto{\pgfqpoint{5.450699in}{0.627125in}}%
\pgfpathlineto{\pgfqpoint{5.454054in}{0.608164in}}%
\pgfpathlineto{\pgfqpoint{5.457410in}{0.548672in}}%
\pgfpathlineto{\pgfqpoint{5.460766in}{0.585025in}}%
\pgfpathlineto{\pgfqpoint{5.464122in}{0.610098in}}%
\pgfpathlineto{\pgfqpoint{5.467478in}{0.616696in}}%
\pgfpathlineto{\pgfqpoint{5.470833in}{0.612124in}}%
\pgfpathlineto{\pgfqpoint{5.474189in}{0.591603in}}%
\pgfpathlineto{\pgfqpoint{5.477545in}{0.508956in}}%
\pgfpathlineto{\pgfqpoint{5.480901in}{0.603721in}}%
\pgfpathlineto{\pgfqpoint{5.484256in}{0.626110in}}%
\pgfpathlineto{\pgfqpoint{5.487612in}{0.634032in}}%
\pgfpathlineto{\pgfqpoint{5.490968in}{0.628646in}}%
\pgfpathlineto{\pgfqpoint{5.494324in}{0.606991in}}%
\pgfpathlineto{\pgfqpoint{5.497680in}{0.522114in}}%
\pgfpathlineto{\pgfqpoint{5.501035in}{0.593890in}}%
\pgfpathlineto{\pgfqpoint{5.504391in}{0.614962in}}%
\pgfpathlineto{\pgfqpoint{5.507747in}{0.619819in}}%
\pgfpathlineto{\pgfqpoint{5.511103in}{0.613460in}}%
\pgfpathlineto{\pgfqpoint{5.514458in}{0.588647in}}%
\pgfpathlineto{\pgfqpoint{5.517814in}{0.549940in}}%
\pgfpathlineto{\pgfqpoint{5.521170in}{0.611809in}}%
\pgfpathlineto{\pgfqpoint{5.524526in}{0.631299in}}%
\pgfpathlineto{\pgfqpoint{5.527882in}{0.637359in}}%
\pgfpathlineto{\pgfqpoint{5.531237in}{0.629612in}}%
\pgfpathlineto{\pgfqpoint{5.534593in}{0.606011in}}%
\pgfpathlineto{\pgfqpoint{5.537949in}{0.502674in}}%
\pgfpathlineto{\pgfqpoint{5.541305in}{0.594478in}}%
\pgfpathlineto{\pgfqpoint{5.544660in}{0.613378in}}%
\pgfpathlineto{\pgfqpoint{5.548016in}{0.616146in}}%
\pgfpathlineto{\pgfqpoint{5.551372in}{0.606084in}}%
\pgfpathlineto{\pgfqpoint{5.554728in}{0.566419in}}%
\pgfpathlineto{\pgfqpoint{5.558084in}{0.582002in}}%
\pgfpathlineto{\pgfqpoint{5.561439in}{0.619286in}}%
\pgfpathlineto{\pgfqpoint{5.564795in}{0.633504in}}%
\pgfpathlineto{\pgfqpoint{5.568151in}{0.634970in}}%
\pgfpathlineto{\pgfqpoint{5.571507in}{0.621609in}}%
\pgfpathlineto{\pgfqpoint{5.574863in}{0.589703in}}%
\pgfpathlineto{\pgfqpoint{5.578218in}{0.543116in}}%
\pgfpathlineto{\pgfqpoint{5.581574in}{0.597888in}}%
\pgfpathlineto{\pgfqpoint{5.584930in}{0.609619in}}%
\pgfpathlineto{\pgfqpoint{5.588286in}{0.607046in}}%
\pgfpathlineto{\pgfqpoint{5.591641in}{0.587146in}}%
\pgfpathlineto{\pgfqpoint{5.594997in}{0.514514in}}%
\pgfpathlineto{\pgfqpoint{5.598353in}{0.601965in}}%
\pgfpathlineto{\pgfqpoint{5.601709in}{0.623982in}}%
\pgfpathlineto{\pgfqpoint{5.605065in}{0.631299in}}%
\pgfpathlineto{\pgfqpoint{5.608420in}{0.624434in}}%
\pgfpathlineto{\pgfqpoint{5.611776in}{0.602157in}}%
\pgfpathlineto{\pgfqpoint{5.615132in}{0.522924in}}%
\pgfpathlineto{\pgfqpoint{5.618488in}{0.584306in}}%
\pgfpathlineto{\pgfqpoint{5.621843in}{0.605012in}}%
\pgfpathlineto{\pgfqpoint{5.625199in}{0.608286in}}%
\pgfpathlineto{\pgfqpoint{5.628555in}{0.598363in}}%
\pgfpathlineto{\pgfqpoint{5.631911in}{0.558152in}}%
\pgfpathlineto{\pgfqpoint{5.635267in}{0.575430in}}%
\pgfpathlineto{\pgfqpoint{5.638622in}{0.612304in}}%
\pgfpathlineto{\pgfqpoint{5.641978in}{0.626366in}}%
\pgfpathlineto{\pgfqpoint{5.645334in}{0.627569in}}%
\pgfpathlineto{\pgfqpoint{5.648690in}{0.613272in}}%
\pgfpathlineto{\pgfqpoint{5.652045in}{0.574661in}}%
\pgfpathlineto{\pgfqpoint{5.655401in}{0.564593in}}%
\pgfpathlineto{\pgfqpoint{5.658757in}{0.602550in}}%
\pgfpathlineto{\pgfqpoint{5.662113in}{0.612893in}}%
\pgfpathlineto{\pgfqpoint{5.665469in}{0.611269in}}%
\pgfpathlineto{\pgfqpoint{5.668824in}{0.596058in}}%
\pgfpathlineto{\pgfqpoint{5.672180in}{0.525101in}}%
\pgfpathlineto{\pgfqpoint{5.675536in}{0.591550in}}%
\pgfpathlineto{\pgfqpoint{5.678892in}{0.619589in}}%
\pgfpathlineto{\pgfqpoint{5.682247in}{0.630384in}}%
\pgfpathlineto{\pgfqpoint{5.685603in}{0.628486in}}%
\pgfpathlineto{\pgfqpoint{5.688959in}{0.610631in}}%
\pgfpathlineto{\pgfqpoint{5.692315in}{0.555079in}}%
\pgfpathlineto{\pgfqpoint{5.699026in}{0.611567in}}%
\pgfpathlineto{\pgfqpoint{5.702382in}{0.619118in}}%
\pgfpathlineto{\pgfqpoint{5.705738in}{0.615725in}}%
\pgfpathlineto{\pgfqpoint{5.709094in}{0.598062in}}%
\pgfpathlineto{\pgfqpoint{5.712449in}{0.483886in}}%
\pgfpathlineto{\pgfqpoint{5.715805in}{0.600811in}}%
\pgfpathlineto{\pgfqpoint{5.719161in}{0.626086in}}%
\pgfpathlineto{\pgfqpoint{5.722517in}{0.635639in}}%
\pgfpathlineto{\pgfqpoint{5.725873in}{0.632444in}}%
\pgfpathlineto{\pgfqpoint{5.729228in}{0.613835in}}%
\pgfpathlineto{\pgfqpoint{5.732584in}{0.558212in}}%
\pgfpathlineto{\pgfqpoint{5.739296in}{0.612524in}}%
\pgfpathlineto{\pgfqpoint{5.742652in}{0.619150in}}%
\pgfpathlineto{\pgfqpoint{5.746007in}{0.614058in}}%
\pgfpathlineto{\pgfqpoint{5.749363in}{0.591524in}}%
\pgfpathlineto{\pgfqpoint{5.752719in}{0.537536in}}%
\pgfpathlineto{\pgfqpoint{5.756075in}{0.609938in}}%
\pgfpathlineto{\pgfqpoint{5.759430in}{0.630659in}}%
\pgfpathlineto{\pgfqpoint{5.762786in}{0.637427in}}%
\pgfpathlineto{\pgfqpoint{5.766142in}{0.630415in}}%
\pgfpathlineto{\pgfqpoint{5.769498in}{0.608451in}}%
\pgfpathlineto{\pgfqpoint{5.772854in}{0.535465in}}%
\pgfpathlineto{\pgfqpoint{5.776209in}{0.587722in}}%
\pgfpathlineto{\pgfqpoint{5.779565in}{0.609022in}}%
\pgfpathlineto{\pgfqpoint{5.782921in}{0.612194in}}%
\pgfpathlineto{\pgfqpoint{5.786277in}{0.601645in}}%
\pgfpathlineto{\pgfqpoint{5.789632in}{0.557223in}}%
\pgfpathlineto{\pgfqpoint{5.792988in}{0.583664in}}%
\pgfpathlineto{\pgfqpoint{5.796344in}{0.618235in}}%
\pgfpathlineto{\pgfqpoint{5.799700in}{0.631491in}}%
\pgfpathlineto{\pgfqpoint{5.803056in}{0.631770in}}%
\pgfpathlineto{\pgfqpoint{5.806411in}{0.617047in}}%
\pgfpathlineto{\pgfqpoint{5.813123in}{0.552263in}}%
\pgfpathlineto{\pgfqpoint{5.816479in}{0.597205in}}%
\pgfpathlineto{\pgfqpoint{5.819834in}{0.607275in}}%
\pgfpathlineto{\pgfqpoint{5.823190in}{0.603398in}}%
\pgfpathlineto{\pgfqpoint{5.826546in}{0.580692in}}%
\pgfpathlineto{\pgfqpoint{5.829902in}{0.535919in}}%
\pgfpathlineto{\pgfqpoint{5.833258in}{0.602533in}}%
\pgfpathlineto{\pgfqpoint{5.836613in}{0.622651in}}%
\pgfpathlineto{\pgfqpoint{5.839969in}{0.628716in}}%
\pgfpathlineto{\pgfqpoint{5.843325in}{0.620251in}}%
\pgfpathlineto{\pgfqpoint{5.846681in}{0.594919in}}%
\pgfpathlineto{\pgfqpoint{5.850036in}{0.464888in}}%
\pgfpathlineto{\pgfqpoint{5.853392in}{0.590390in}}%
\pgfpathlineto{\pgfqpoint{5.856748in}{0.607542in}}%
\pgfpathlineto{\pgfqpoint{5.860104in}{0.609760in}}%
\pgfpathlineto{\pgfqpoint{5.863460in}{0.599482in}}%
\pgfpathlineto{\pgfqpoint{5.866815in}{0.560016in}}%
\pgfpathlineto{\pgfqpoint{5.870171in}{0.574650in}}%
\pgfpathlineto{\pgfqpoint{5.873527in}{0.612138in}}%
\pgfpathlineto{\pgfqpoint{5.876883in}{0.626444in}}%
\pgfpathlineto{\pgfqpoint{5.880238in}{0.628105in}}%
\pgfpathlineto{\pgfqpoint{5.883594in}{0.614337in}}%
\pgfpathlineto{\pgfqpoint{5.886950in}{0.575464in}}%
\pgfpathlineto{\pgfqpoint{5.890306in}{0.568138in}}%
\pgfpathlineto{\pgfqpoint{5.893662in}{0.605679in}}%
\pgfpathlineto{\pgfqpoint{5.897017in}{0.616440in}}%
\pgfpathlineto{\pgfqpoint{5.900373in}{0.615694in}}%
\pgfpathlineto{\pgfqpoint{5.903729in}{0.602664in}}%
\pgfpathlineto{\pgfqpoint{5.907085in}{0.552754in}}%
\pgfpathlineto{\pgfqpoint{5.913796in}{0.619584in}}%
\pgfpathlineto{\pgfqpoint{5.917152in}{0.632309in}}%
\pgfpathlineto{\pgfqpoint{5.920508in}{0.632804in}}%
\pgfpathlineto{\pgfqpoint{5.923864in}{0.618034in}}%
\pgfpathlineto{\pgfqpoint{5.927219in}{0.576907in}}%
\pgfpathlineto{\pgfqpoint{5.930575in}{0.575188in}}%
\pgfpathlineto{\pgfqpoint{5.933931in}{0.610352in}}%
\pgfpathlineto{\pgfqpoint{5.937287in}{0.620281in}}%
\pgfpathlineto{\pgfqpoint{5.940643in}{0.618663in}}%
\pgfpathlineto{\pgfqpoint{5.943998in}{0.603912in}}%
\pgfpathlineto{\pgfqpoint{5.947354in}{0.539807in}}%
\pgfpathlineto{\pgfqpoint{5.950710in}{0.596627in}}%
\pgfpathlineto{\pgfqpoint{5.954066in}{0.625668in}}%
\pgfpathlineto{\pgfqpoint{5.957421in}{0.636958in}}%
\pgfpathlineto{\pgfqpoint{5.960777in}{0.635784in}}%
\pgfpathlineto{\pgfqpoint{5.964133in}{0.619718in}}%
\pgfpathlineto{\pgfqpoint{5.967489in}{0.578838in}}%
\pgfpathlineto{\pgfqpoint{5.970845in}{0.572048in}}%
\pgfpathlineto{\pgfqpoint{5.974200in}{0.607434in}}%
\pgfpathlineto{\pgfqpoint{5.977556in}{0.616055in}}%
\pgfpathlineto{\pgfqpoint{5.980912in}{0.611819in}}%
\pgfpathlineto{\pgfqpoint{5.984268in}{0.589701in}}%
\pgfpathlineto{\pgfqpoint{5.987623in}{0.536531in}}%
\pgfpathlineto{\pgfqpoint{5.990979in}{0.608692in}}%
\pgfpathlineto{\pgfqpoint{5.994335in}{0.629404in}}%
\pgfpathlineto{\pgfqpoint{5.997691in}{0.635996in}}%
\pgfpathlineto{\pgfqpoint{6.001047in}{0.628498in}}%
\pgfpathlineto{\pgfqpoint{6.004402in}{0.606360in}}%
\pgfpathlineto{\pgfqpoint{6.007758in}{0.535341in}}%
\pgfpathlineto{\pgfqpoint{6.011114in}{0.583918in}}%
\pgfpathlineto{\pgfqpoint{6.014470in}{0.605111in}}%
\pgfpathlineto{\pgfqpoint{6.017825in}{0.607668in}}%
\pgfpathlineto{\pgfqpoint{6.021181in}{0.595549in}}%
\pgfpathlineto{\pgfqpoint{6.024537in}{0.539262in}}%
\pgfpathlineto{\pgfqpoint{6.027893in}{0.586646in}}%
\pgfpathlineto{\pgfqpoint{6.031249in}{0.617290in}}%
\pgfpathlineto{\pgfqpoint{6.034604in}{0.628991in}}%
\pgfpathlineto{\pgfqpoint{6.037960in}{0.627382in}}%
\pgfpathlineto{\pgfqpoint{6.041316in}{0.610562in}}%
\pgfpathlineto{\pgfqpoint{6.044672in}{0.567444in}}%
\pgfpathlineto{\pgfqpoint{6.048027in}{0.566360in}}%
\pgfpathlineto{\pgfqpoint{6.051383in}{0.599471in}}%
\pgfpathlineto{\pgfqpoint{6.054739in}{0.607289in}}%
\pgfpathlineto{\pgfqpoint{6.058095in}{0.602132in}}%
\pgfpathlineto{\pgfqpoint{6.061451in}{0.577499in}}%
\pgfpathlineto{\pgfqpoint{6.064806in}{0.542936in}}%
\pgfpathlineto{\pgfqpoint{6.068162in}{0.602760in}}%
\pgfpathlineto{\pgfqpoint{6.071518in}{0.621943in}}%
\pgfpathlineto{\pgfqpoint{6.074874in}{0.627448in}}%
\pgfpathlineto{\pgfqpoint{6.078230in}{0.618358in}}%
\pgfpathlineto{\pgfqpoint{6.081585in}{0.590717in}}%
\pgfpathlineto{\pgfqpoint{6.084941in}{0.524879in}}%
\pgfpathlineto{\pgfqpoint{6.088297in}{0.595950in}}%
\pgfpathlineto{\pgfqpoint{6.091653in}{0.611509in}}%
\pgfpathlineto{\pgfqpoint{6.095008in}{0.613639in}}%
\pgfpathlineto{\pgfqpoint{6.098364in}{0.604282in}}%
\pgfpathlineto{\pgfqpoint{6.101720in}{0.570660in}}%
\pgfpathlineto{\pgfqpoint{6.105076in}{0.568819in}}%
\pgfpathlineto{\pgfqpoint{6.108432in}{0.611778in}}%
\pgfpathlineto{\pgfqpoint{6.111787in}{0.627697in}}%
\pgfpathlineto{\pgfqpoint{6.115143in}{0.631070in}}%
\pgfpathlineto{\pgfqpoint{6.118499in}{0.619709in}}%
\pgfpathlineto{\pgfqpoint{6.121855in}{0.586599in}}%
\pgfpathlineto{\pgfqpoint{6.125210in}{0.559430in}}%
\pgfpathlineto{\pgfqpoint{6.128566in}{0.606270in}}%
\pgfpathlineto{\pgfqpoint{6.131922in}{0.619182in}}%
\pgfpathlineto{\pgfqpoint{6.135278in}{0.620122in}}%
\pgfpathlineto{\pgfqpoint{6.138634in}{0.609773in}}%
\pgfpathlineto{\pgfqpoint{6.141989in}{0.573585in}}%
\pgfpathlineto{\pgfqpoint{6.145345in}{0.578333in}}%
\pgfpathlineto{\pgfqpoint{6.148701in}{0.618761in}}%
\pgfpathlineto{\pgfqpoint{6.152057in}{0.634055in}}%
\pgfpathlineto{\pgfqpoint{6.155412in}{0.637055in}}%
\pgfpathlineto{\pgfqpoint{6.158768in}{0.625640in}}%
\pgfpathlineto{\pgfqpoint{6.162124in}{0.595077in}}%
\pgfpathlineto{\pgfqpoint{6.165480in}{0.547953in}}%
\pgfpathlineto{\pgfqpoint{6.168836in}{0.605007in}}%
\pgfpathlineto{\pgfqpoint{6.172191in}{0.618443in}}%
\pgfpathlineto{\pgfqpoint{6.175547in}{0.618627in}}%
\pgfpathlineto{\pgfqpoint{6.178903in}{0.605774in}}%
\pgfpathlineto{\pgfqpoint{6.182259in}{0.553240in}}%
\pgfpathlineto{\pgfqpoint{6.185614in}{0.592905in}}%
\pgfpathlineto{\pgfqpoint{6.188970in}{0.624536in}}%
\pgfpathlineto{\pgfqpoint{6.192326in}{0.636799in}}%
\pgfpathlineto{\pgfqpoint{6.195682in}{0.636439in}}%
\pgfpathlineto{\pgfqpoint{6.199038in}{0.621404in}}%
\pgfpathlineto{\pgfqpoint{6.202393in}{0.585914in}}%
\pgfpathlineto{\pgfqpoint{6.205749in}{0.556946in}}%
\pgfpathlineto{\pgfqpoint{6.209105in}{0.601264in}}%
\pgfpathlineto{\pgfqpoint{6.212461in}{0.610983in}}%
\pgfpathlineto{\pgfqpoint{6.215816in}{0.606540in}}%
\pgfpathlineto{\pgfqpoint{6.219172in}{0.582037in}}%
\pgfpathlineto{\pgfqpoint{6.222528in}{0.549526in}}%
\pgfpathlineto{\pgfqpoint{6.225884in}{0.608409in}}%
\pgfpathlineto{\pgfqpoint{6.229240in}{0.627470in}}%
\pgfpathlineto{\pgfqpoint{6.232595in}{0.632743in}}%
\pgfpathlineto{\pgfqpoint{6.235951in}{0.623399in}}%
\pgfpathlineto{\pgfqpoint{6.239307in}{0.598751in}}%
\pgfpathlineto{\pgfqpoint{6.242663in}{0.486561in}}%
\pgfpathlineto{\pgfqpoint{6.246019in}{0.585813in}}%
\pgfpathlineto{\pgfqpoint{6.249374in}{0.603469in}}%
\pgfpathlineto{\pgfqpoint{6.252730in}{0.604328in}}%
\pgfpathlineto{\pgfqpoint{6.256086in}{0.589914in}}%
\pgfpathlineto{\pgfqpoint{6.259442in}{0.511172in}}%
\pgfpathlineto{\pgfqpoint{6.262797in}{0.589144in}}%
\pgfpathlineto{\pgfqpoint{6.266153in}{0.616449in}}%
\pgfpathlineto{\pgfqpoint{6.269509in}{0.626644in}}%
\pgfpathlineto{\pgfqpoint{6.272865in}{0.623231in}}%
\pgfpathlineto{\pgfqpoint{6.276221in}{0.604029in}}%
\pgfpathlineto{\pgfqpoint{6.279576in}{0.544999in}}%
\pgfpathlineto{\pgfqpoint{6.282932in}{0.578584in}}%
\pgfpathlineto{\pgfqpoint{6.286288in}{0.603891in}}%
\pgfpathlineto{\pgfqpoint{6.289644in}{0.610036in}}%
\pgfpathlineto{\pgfqpoint{6.292999in}{0.604412in}}%
\pgfpathlineto{\pgfqpoint{6.296355in}{0.580498in}}%
\pgfpathlineto{\pgfqpoint{6.299711in}{0.537711in}}%
\pgfpathlineto{\pgfqpoint{6.303067in}{0.602397in}}%
\pgfpathlineto{\pgfqpoint{6.306423in}{0.622244in}}%
\pgfpathlineto{\pgfqpoint{6.309778in}{0.628419in}}%
\pgfpathlineto{\pgfqpoint{6.313134in}{0.620470in}}%
\pgfpathlineto{\pgfqpoint{6.316490in}{0.593749in}}%
\pgfpathlineto{\pgfqpoint{6.319846in}{0.525069in}}%
\pgfpathlineto{\pgfqpoint{6.323201in}{0.599202in}}%
\pgfpathlineto{\pgfqpoint{6.326557in}{0.615600in}}%
\pgfpathlineto{\pgfqpoint{6.329913in}{0.618781in}}%
\pgfpathlineto{\pgfqpoint{6.333269in}{0.611390in}}%
\pgfpathlineto{\pgfqpoint{6.336625in}{0.585205in}}%
\pgfpathlineto{\pgfqpoint{6.339980in}{0.552242in}}%
\pgfpathlineto{\pgfqpoint{6.343336in}{0.610365in}}%
\pgfpathlineto{\pgfqpoint{6.346692in}{0.629336in}}%
\pgfpathlineto{\pgfqpoint{6.350048in}{0.635236in}}%
\pgfpathlineto{\pgfqpoint{6.353403in}{0.628049in}}%
\pgfpathlineto{\pgfqpoint{6.356759in}{0.613266in}}%
\pgfpathlineto{\pgfqpoint{6.360115in}{0.592531in}}%
\pgfpathlineto{\pgfqpoint{6.366827in}{0.531114in}}%
\pgfpathlineto{\pgfqpoint{6.370182in}{0.528216in}}%
\pgfpathlineto{\pgfqpoint{6.373538in}{0.555432in}}%
\pgfpathlineto{\pgfqpoint{6.376894in}{0.577212in}}%
\pgfpathlineto{\pgfqpoint{6.380250in}{0.588616in}}%
\pgfpathlineto{\pgfqpoint{6.383605in}{0.587339in}}%
\pgfpathlineto{\pgfqpoint{6.383605in}{0.587339in}}%
\pgfusepath{stroke}%
\end{pgfscope}%
\begin{pgfscope}%
\pgfpathrectangle{\pgfqpoint{4.621821in}{0.320625in}}{\pgfqpoint{1.845679in}{1.164375in}}%
\pgfusepath{clip}%
\pgfsetrectcap%
\pgfsetroundjoin%
\pgfsetlinewidth{0.602250pt}%
\definecolor{currentstroke}{rgb}{0.141176,0.345098,0.650980}%
\pgfsetstrokecolor{currentstroke}%
\pgfsetstrokeopacity{0.120000}%
\pgfsetdash{}{0pt}%
\pgfpathmoveto{\pgfqpoint{4.705715in}{1.479263in}}%
\pgfpathlineto{\pgfqpoint{4.712427in}{1.467492in}}%
\pgfpathlineto{\pgfqpoint{4.722494in}{1.453345in}}%
\pgfpathlineto{\pgfqpoint{4.742629in}{1.427571in}}%
\pgfpathlineto{\pgfqpoint{4.766120in}{1.399100in}}%
\pgfpathlineto{\pgfqpoint{4.782898in}{1.380408in}}%
\pgfpathlineto{\pgfqpoint{4.799677in}{1.363349in}}%
\pgfpathlineto{\pgfqpoint{4.819812in}{1.344689in}}%
\pgfpathlineto{\pgfqpoint{4.843302in}{1.324650in}}%
\pgfpathlineto{\pgfqpoint{4.873504in}{1.300485in}}%
\pgfpathlineto{\pgfqpoint{4.917130in}{1.267139in}}%
\pgfpathlineto{\pgfqpoint{4.990957in}{1.212326in}}%
\pgfpathlineto{\pgfqpoint{5.128544in}{1.111783in}}%
\pgfpathlineto{\pgfqpoint{5.410429in}{0.907360in}}%
\pgfpathlineto{\pgfqpoint{5.578218in}{0.786396in}}%
\pgfpathlineto{\pgfqpoint{5.584930in}{0.784073in}}%
\pgfpathlineto{\pgfqpoint{5.591641in}{0.776572in}}%
\pgfpathlineto{\pgfqpoint{5.594997in}{0.774523in}}%
\pgfpathlineto{\pgfqpoint{5.598353in}{0.773502in}}%
\pgfpathlineto{\pgfqpoint{5.601709in}{0.771288in}}%
\pgfpathlineto{\pgfqpoint{5.605065in}{0.767409in}}%
\pgfpathlineto{\pgfqpoint{5.608420in}{0.764612in}}%
\pgfpathlineto{\pgfqpoint{5.615132in}{0.768018in}}%
\pgfpathlineto{\pgfqpoint{5.618488in}{0.763483in}}%
\pgfpathlineto{\pgfqpoint{5.621843in}{0.756215in}}%
\pgfpathlineto{\pgfqpoint{5.625199in}{0.752248in}}%
\pgfpathlineto{\pgfqpoint{5.628555in}{0.751104in}}%
\pgfpathlineto{\pgfqpoint{5.631911in}{0.748776in}}%
\pgfpathlineto{\pgfqpoint{5.635267in}{0.744981in}}%
\pgfpathlineto{\pgfqpoint{5.638622in}{0.744668in}}%
\pgfpathlineto{\pgfqpoint{5.641978in}{0.748569in}}%
\pgfpathlineto{\pgfqpoint{5.645334in}{0.747722in}}%
\pgfpathlineto{\pgfqpoint{5.648690in}{0.739467in}}%
\pgfpathlineto{\pgfqpoint{5.652045in}{0.732964in}}%
\pgfpathlineto{\pgfqpoint{5.658757in}{0.730017in}}%
\pgfpathlineto{\pgfqpoint{5.662113in}{0.726177in}}%
\pgfpathlineto{\pgfqpoint{5.665469in}{0.723872in}}%
\pgfpathlineto{\pgfqpoint{5.668824in}{0.728076in}}%
\pgfpathlineto{\pgfqpoint{5.672180in}{0.731017in}}%
\pgfpathlineto{\pgfqpoint{5.675536in}{0.723491in}}%
\pgfpathlineto{\pgfqpoint{5.678892in}{0.713717in}}%
\pgfpathlineto{\pgfqpoint{5.682247in}{0.714554in}}%
\pgfpathlineto{\pgfqpoint{5.685603in}{0.718450in}}%
\pgfpathlineto{\pgfqpoint{5.688959in}{0.716062in}}%
\pgfpathlineto{\pgfqpoint{5.692315in}{0.706477in}}%
\pgfpathlineto{\pgfqpoint{5.695671in}{0.704159in}}%
\pgfpathlineto{\pgfqpoint{5.699026in}{0.718108in}}%
\pgfpathlineto{\pgfqpoint{5.702382in}{0.725287in}}%
\pgfpathlineto{\pgfqpoint{5.705738in}{0.715011in}}%
\pgfpathlineto{\pgfqpoint{5.709094in}{0.691667in}}%
\pgfpathlineto{\pgfqpoint{5.712449in}{0.711415in}}%
\pgfpathlineto{\pgfqpoint{5.715805in}{0.726449in}}%
\pgfpathlineto{\pgfqpoint{5.719161in}{0.728020in}}%
\pgfpathlineto{\pgfqpoint{5.722517in}{0.715791in}}%
\pgfpathlineto{\pgfqpoint{5.725873in}{0.681349in}}%
\pgfpathlineto{\pgfqpoint{5.729228in}{0.712631in}}%
\pgfpathlineto{\pgfqpoint{5.732584in}{0.735708in}}%
\pgfpathlineto{\pgfqpoint{5.735940in}{0.740075in}}%
\pgfpathlineto{\pgfqpoint{5.739296in}{0.721127in}}%
\pgfpathlineto{\pgfqpoint{5.742652in}{0.669872in}}%
\pgfpathlineto{\pgfqpoint{5.746007in}{0.723302in}}%
\pgfpathlineto{\pgfqpoint{5.749363in}{0.737739in}}%
\pgfpathlineto{\pgfqpoint{5.752719in}{0.735506in}}%
\pgfpathlineto{\pgfqpoint{5.756075in}{0.712559in}}%
\pgfpathlineto{\pgfqpoint{5.759430in}{0.683870in}}%
\pgfpathlineto{\pgfqpoint{5.762786in}{0.736857in}}%
\pgfpathlineto{\pgfqpoint{5.766142in}{0.752117in}}%
\pgfpathlineto{\pgfqpoint{5.769498in}{0.746943in}}%
\pgfpathlineto{\pgfqpoint{5.772854in}{0.718675in}}%
\pgfpathlineto{\pgfqpoint{5.776209in}{0.656417in}}%
\pgfpathlineto{\pgfqpoint{5.779565in}{0.713388in}}%
\pgfpathlineto{\pgfqpoint{5.782921in}{0.718149in}}%
\pgfpathlineto{\pgfqpoint{5.786277in}{0.692395in}}%
\pgfpathlineto{\pgfqpoint{5.789632in}{0.689421in}}%
\pgfpathlineto{\pgfqpoint{5.792988in}{0.732355in}}%
\pgfpathlineto{\pgfqpoint{5.796344in}{0.743258in}}%
\pgfpathlineto{\pgfqpoint{5.799700in}{0.730818in}}%
\pgfpathlineto{\pgfqpoint{5.803056in}{0.697432in}}%
\pgfpathlineto{\pgfqpoint{5.806411in}{0.640716in}}%
\pgfpathlineto{\pgfqpoint{5.809767in}{0.686139in}}%
\pgfpathlineto{\pgfqpoint{5.813123in}{0.676465in}}%
\pgfpathlineto{\pgfqpoint{5.816479in}{0.643518in}}%
\pgfpathlineto{\pgfqpoint{5.819834in}{0.705367in}}%
\pgfpathlineto{\pgfqpoint{5.823190in}{0.721942in}}%
\pgfpathlineto{\pgfqpoint{5.826546in}{0.714007in}}%
\pgfpathlineto{\pgfqpoint{5.833258in}{0.648828in}}%
\pgfpathlineto{\pgfqpoint{5.836613in}{0.688292in}}%
\pgfpathlineto{\pgfqpoint{5.839969in}{0.689059in}}%
\pgfpathlineto{\pgfqpoint{5.843325in}{0.655970in}}%
\pgfpathlineto{\pgfqpoint{5.846681in}{0.669456in}}%
\pgfpathlineto{\pgfqpoint{5.850036in}{0.706238in}}%
\pgfpathlineto{\pgfqpoint{5.853392in}{0.714580in}}%
\pgfpathlineto{\pgfqpoint{5.856748in}{0.696283in}}%
\pgfpathlineto{\pgfqpoint{5.860104in}{0.608287in}}%
\pgfpathlineto{\pgfqpoint{5.863460in}{0.699105in}}%
\pgfpathlineto{\pgfqpoint{5.866815in}{0.714803in}}%
\pgfpathlineto{\pgfqpoint{5.870171in}{0.713662in}}%
\pgfpathlineto{\pgfqpoint{5.873527in}{0.693234in}}%
\pgfpathlineto{\pgfqpoint{5.876883in}{0.637873in}}%
\pgfpathlineto{\pgfqpoint{5.880238in}{0.710638in}}%
\pgfpathlineto{\pgfqpoint{5.883594in}{0.728433in}}%
\pgfpathlineto{\pgfqpoint{5.886950in}{0.726128in}}%
\pgfpathlineto{\pgfqpoint{5.890306in}{0.688138in}}%
\pgfpathlineto{\pgfqpoint{5.893662in}{0.696340in}}%
\pgfpathlineto{\pgfqpoint{5.897017in}{0.729848in}}%
\pgfpathlineto{\pgfqpoint{5.900373in}{0.738219in}}%
\pgfpathlineto{\pgfqpoint{5.903729in}{0.732551in}}%
\pgfpathlineto{\pgfqpoint{5.907085in}{0.703741in}}%
\pgfpathlineto{\pgfqpoint{5.910441in}{0.689995in}}%
\pgfpathlineto{\pgfqpoint{5.913796in}{0.736969in}}%
\pgfpathlineto{\pgfqpoint{5.917152in}{0.750595in}}%
\pgfpathlineto{\pgfqpoint{5.920508in}{0.744657in}}%
\pgfpathlineto{\pgfqpoint{5.923864in}{0.710719in}}%
\pgfpathlineto{\pgfqpoint{5.927219in}{0.688261in}}%
\pgfpathlineto{\pgfqpoint{5.930575in}{0.728253in}}%
\pgfpathlineto{\pgfqpoint{5.933931in}{0.733241in}}%
\pgfpathlineto{\pgfqpoint{5.937287in}{0.716535in}}%
\pgfpathlineto{\pgfqpoint{5.940643in}{0.646523in}}%
\pgfpathlineto{\pgfqpoint{5.943998in}{0.732131in}}%
\pgfpathlineto{\pgfqpoint{5.947354in}{0.750613in}}%
\pgfpathlineto{\pgfqpoint{5.950710in}{0.746929in}}%
\pgfpathlineto{\pgfqpoint{5.954066in}{0.722763in}}%
\pgfpathlineto{\pgfqpoint{5.957421in}{0.649250in}}%
\pgfpathlineto{\pgfqpoint{5.960777in}{0.687829in}}%
\pgfpathlineto{\pgfqpoint{5.964133in}{0.694580in}}%
\pgfpathlineto{\pgfqpoint{5.967489in}{0.640173in}}%
\pgfpathlineto{\pgfqpoint{5.970845in}{0.698543in}}%
\pgfpathlineto{\pgfqpoint{5.974200in}{0.726909in}}%
\pgfpathlineto{\pgfqpoint{5.977556in}{0.729824in}}%
\pgfpathlineto{\pgfqpoint{5.980912in}{0.709312in}}%
\pgfpathlineto{\pgfqpoint{5.984268in}{0.654149in}}%
\pgfpathlineto{\pgfqpoint{5.987623in}{0.665216in}}%
\pgfpathlineto{\pgfqpoint{5.990979in}{0.680844in}}%
\pgfpathlineto{\pgfqpoint{5.994335in}{0.654554in}}%
\pgfpathlineto{\pgfqpoint{5.997691in}{0.662928in}}%
\pgfpathlineto{\pgfqpoint{6.001047in}{0.703690in}}%
\pgfpathlineto{\pgfqpoint{6.004402in}{0.713448in}}%
\pgfpathlineto{\pgfqpoint{6.007758in}{0.697041in}}%
\pgfpathlineto{\pgfqpoint{6.011114in}{0.609815in}}%
\pgfpathlineto{\pgfqpoint{6.014470in}{0.683839in}}%
\pgfpathlineto{\pgfqpoint{6.017825in}{0.701575in}}%
\pgfpathlineto{\pgfqpoint{6.021181in}{0.698919in}}%
\pgfpathlineto{\pgfqpoint{6.024537in}{0.670091in}}%
\pgfpathlineto{\pgfqpoint{6.027893in}{0.663150in}}%
\pgfpathlineto{\pgfqpoint{6.031249in}{0.707424in}}%
\pgfpathlineto{\pgfqpoint{6.034604in}{0.719249in}}%
\pgfpathlineto{\pgfqpoint{6.037960in}{0.707594in}}%
\pgfpathlineto{\pgfqpoint{6.041316in}{0.584525in}}%
\pgfpathlineto{\pgfqpoint{6.044672in}{0.707221in}}%
\pgfpathlineto{\pgfqpoint{6.048027in}{0.727454in}}%
\pgfpathlineto{\pgfqpoint{6.051383in}{0.730882in}}%
\pgfpathlineto{\pgfqpoint{6.054739in}{0.720361in}}%
\pgfpathlineto{\pgfqpoint{6.058095in}{0.674569in}}%
\pgfpathlineto{\pgfqpoint{6.064806in}{0.734852in}}%
\pgfpathlineto{\pgfqpoint{6.068162in}{0.743678in}}%
\pgfpathlineto{\pgfqpoint{6.071518in}{0.732675in}}%
\pgfpathlineto{\pgfqpoint{6.074874in}{0.698103in}}%
\pgfpathlineto{\pgfqpoint{6.078230in}{0.656371in}}%
\pgfpathlineto{\pgfqpoint{6.081585in}{0.699740in}}%
\pgfpathlineto{\pgfqpoint{6.084941in}{0.697026in}}%
\pgfpathlineto{\pgfqpoint{6.088297in}{0.633204in}}%
\pgfpathlineto{\pgfqpoint{6.091653in}{0.699085in}}%
\pgfpathlineto{\pgfqpoint{6.095008in}{0.725532in}}%
\pgfpathlineto{\pgfqpoint{6.098364in}{0.727271in}}%
\pgfpathlineto{\pgfqpoint{6.101720in}{0.705283in}}%
\pgfpathlineto{\pgfqpoint{6.105076in}{0.632467in}}%
\pgfpathlineto{\pgfqpoint{6.108432in}{0.676157in}}%
\pgfpathlineto{\pgfqpoint{6.111787in}{0.688450in}}%
\pgfpathlineto{\pgfqpoint{6.115143in}{0.667431in}}%
\pgfpathlineto{\pgfqpoint{6.118499in}{0.653246in}}%
\pgfpathlineto{\pgfqpoint{6.121855in}{0.703245in}}%
\pgfpathlineto{\pgfqpoint{6.125210in}{0.716048in}}%
\pgfpathlineto{\pgfqpoint{6.128566in}{0.703792in}}%
\pgfpathlineto{\pgfqpoint{6.131922in}{0.641678in}}%
\pgfpathlineto{\pgfqpoint{6.135278in}{0.683534in}}%
\pgfpathlineto{\pgfqpoint{6.138634in}{0.705428in}}%
\pgfpathlineto{\pgfqpoint{6.141989in}{0.705787in}}%
\pgfpathlineto{\pgfqpoint{6.145345in}{0.684787in}}%
\pgfpathlineto{\pgfqpoint{6.148701in}{0.641963in}}%
\pgfpathlineto{\pgfqpoint{6.152057in}{0.706481in}}%
\pgfpathlineto{\pgfqpoint{6.155412in}{0.722795in}}%
\pgfpathlineto{\pgfqpoint{6.158768in}{0.717516in}}%
\pgfpathlineto{\pgfqpoint{6.162124in}{0.668765in}}%
\pgfpathlineto{\pgfqpoint{6.168836in}{0.725199in}}%
\pgfpathlineto{\pgfqpoint{6.172191in}{0.731474in}}%
\pgfpathlineto{\pgfqpoint{6.175547in}{0.723532in}}%
\pgfpathlineto{\pgfqpoint{6.178903in}{0.687097in}}%
\pgfpathlineto{\pgfqpoint{6.182259in}{0.695950in}}%
\pgfpathlineto{\pgfqpoint{6.185614in}{0.733652in}}%
\pgfpathlineto{\pgfqpoint{6.188970in}{0.744401in}}%
\pgfpathlineto{\pgfqpoint{6.192326in}{0.734606in}}%
\pgfpathlineto{\pgfqpoint{6.195682in}{0.684066in}}%
\pgfpathlineto{\pgfqpoint{6.202393in}{0.734042in}}%
\pgfpathlineto{\pgfqpoint{6.205749in}{0.736761in}}%
\pgfpathlineto{\pgfqpoint{6.209105in}{0.720579in}}%
\pgfpathlineto{\pgfqpoint{6.212461in}{0.609371in}}%
\pgfpathlineto{\pgfqpoint{6.215816in}{0.730310in}}%
\pgfpathlineto{\pgfqpoint{6.219172in}{0.750770in}}%
\pgfpathlineto{\pgfqpoint{6.222528in}{0.750137in}}%
\pgfpathlineto{\pgfqpoint{6.225884in}{0.727451in}}%
\pgfpathlineto{\pgfqpoint{6.229240in}{0.649207in}}%
\pgfpathlineto{\pgfqpoint{6.232595in}{0.700083in}}%
\pgfpathlineto{\pgfqpoint{6.235951in}{0.711166in}}%
\pgfpathlineto{\pgfqpoint{6.239307in}{0.688164in}}%
\pgfpathlineto{\pgfqpoint{6.242663in}{0.680741in}}%
\pgfpathlineto{\pgfqpoint{6.246019in}{0.727371in}}%
\pgfpathlineto{\pgfqpoint{6.249374in}{0.739180in}}%
\pgfpathlineto{\pgfqpoint{6.252730in}{0.727167in}}%
\pgfpathlineto{\pgfqpoint{6.256086in}{0.693893in}}%
\pgfpathlineto{\pgfqpoint{6.259442in}{0.626727in}}%
\pgfpathlineto{\pgfqpoint{6.262797in}{0.682014in}}%
\pgfpathlineto{\pgfqpoint{6.266153in}{0.673032in}}%
\pgfpathlineto{\pgfqpoint{6.269509in}{0.630489in}}%
\pgfpathlineto{\pgfqpoint{6.272865in}{0.700716in}}%
\pgfpathlineto{\pgfqpoint{6.276221in}{0.718007in}}%
\pgfpathlineto{\pgfqpoint{6.279576in}{0.710646in}}%
\pgfpathlineto{\pgfqpoint{6.282932in}{0.674043in}}%
\pgfpathlineto{\pgfqpoint{6.286288in}{0.655304in}}%
\pgfpathlineto{\pgfqpoint{6.289644in}{0.691596in}}%
\pgfpathlineto{\pgfqpoint{6.292999in}{0.693625in}}%
\pgfpathlineto{\pgfqpoint{6.296355in}{0.668509in}}%
\pgfpathlineto{\pgfqpoint{6.299711in}{0.655404in}}%
\pgfpathlineto{\pgfqpoint{6.303067in}{0.703267in}}%
\pgfpathlineto{\pgfqpoint{6.306423in}{0.715676in}}%
\pgfpathlineto{\pgfqpoint{6.309778in}{0.703621in}}%
\pgfpathlineto{\pgfqpoint{6.313134in}{0.608362in}}%
\pgfpathlineto{\pgfqpoint{6.316490in}{0.698311in}}%
\pgfpathlineto{\pgfqpoint{6.319846in}{0.718593in}}%
\pgfpathlineto{\pgfqpoint{6.323201in}{0.721060in}}%
\pgfpathlineto{\pgfqpoint{6.326557in}{0.707935in}}%
\pgfpathlineto{\pgfqpoint{6.329913in}{0.638827in}}%
\pgfpathlineto{\pgfqpoint{6.333269in}{0.703395in}}%
\pgfpathlineto{\pgfqpoint{6.336625in}{0.729504in}}%
\pgfpathlineto{\pgfqpoint{6.339980in}{0.734541in}}%
\pgfpathlineto{\pgfqpoint{6.343336in}{0.714613in}}%
\pgfpathlineto{\pgfqpoint{6.346692in}{0.653387in}}%
\pgfpathlineto{\pgfqpoint{6.350048in}{0.724586in}}%
\pgfpathlineto{\pgfqpoint{6.353403in}{0.739478in}}%
\pgfpathlineto{\pgfqpoint{6.356759in}{0.738707in}}%
\pgfpathlineto{\pgfqpoint{6.360115in}{0.720381in}}%
\pgfpathlineto{\pgfqpoint{6.363471in}{0.629151in}}%
\pgfpathlineto{\pgfqpoint{6.366827in}{0.731152in}}%
\pgfpathlineto{\pgfqpoint{6.370182in}{0.751264in}}%
\pgfpathlineto{\pgfqpoint{6.373538in}{0.751845in}}%
\pgfpathlineto{\pgfqpoint{6.383605in}{0.710788in}}%
\pgfpathlineto{\pgfqpoint{6.383605in}{0.710788in}}%
\pgfusepath{stroke}%
\end{pgfscope}%
\begin{pgfscope}%
\pgfpathrectangle{\pgfqpoint{4.621821in}{0.320625in}}{\pgfqpoint{1.845679in}{1.164375in}}%
\pgfusepath{clip}%
\pgfsetrectcap%
\pgfsetroundjoin%
\pgfsetlinewidth{0.602250pt}%
\definecolor{currentstroke}{rgb}{0.141176,0.345098,0.650980}%
\pgfsetstrokecolor{currentstroke}%
\pgfsetstrokeopacity{0.120000}%
\pgfsetdash{}{0pt}%
\pgfpathmoveto{\pgfqpoint{4.705715in}{1.478401in}}%
\pgfpathlineto{\pgfqpoint{4.712427in}{1.466304in}}%
\pgfpathlineto{\pgfqpoint{4.722494in}{1.451315in}}%
\pgfpathlineto{\pgfqpoint{4.742629in}{1.423879in}}%
\pgfpathlineto{\pgfqpoint{4.759408in}{1.402982in}}%
\pgfpathlineto{\pgfqpoint{4.772831in}{1.388001in}}%
\pgfpathlineto{\pgfqpoint{4.786254in}{1.374521in}}%
\pgfpathlineto{\pgfqpoint{4.803033in}{1.359412in}}%
\pgfpathlineto{\pgfqpoint{4.823168in}{1.343101in}}%
\pgfpathlineto{\pgfqpoint{4.846658in}{1.325590in}}%
\pgfpathlineto{\pgfqpoint{4.880216in}{1.302021in}}%
\pgfpathlineto{\pgfqpoint{4.930553in}{1.268096in}}%
\pgfpathlineto{\pgfqpoint{5.001024in}{1.221994in}}%
\pgfpathlineto{\pgfqpoint{5.088274in}{1.166303in}}%
\pgfpathlineto{\pgfqpoint{5.195659in}{1.099148in}}%
\pgfpathlineto{\pgfqpoint{5.343314in}{1.008243in}}%
\pgfpathlineto{\pgfqpoint{5.581574in}{0.863012in}}%
\pgfpathlineto{\pgfqpoint{5.722517in}{0.778033in}}%
\pgfpathlineto{\pgfqpoint{5.725873in}{0.776765in}}%
\pgfpathlineto{\pgfqpoint{5.732584in}{0.771506in}}%
\pgfpathlineto{\pgfqpoint{5.739296in}{0.771404in}}%
\pgfpathlineto{\pgfqpoint{5.742652in}{0.768050in}}%
\pgfpathlineto{\pgfqpoint{5.746007in}{0.763269in}}%
\pgfpathlineto{\pgfqpoint{5.749363in}{0.763684in}}%
\pgfpathlineto{\pgfqpoint{5.752719in}{0.768343in}}%
\pgfpathlineto{\pgfqpoint{5.756075in}{0.766583in}}%
\pgfpathlineto{\pgfqpoint{5.759430in}{0.757651in}}%
\pgfpathlineto{\pgfqpoint{5.762786in}{0.752999in}}%
\pgfpathlineto{\pgfqpoint{5.766142in}{0.752438in}}%
\pgfpathlineto{\pgfqpoint{5.772854in}{0.747364in}}%
\pgfpathlineto{\pgfqpoint{5.776209in}{0.751384in}}%
\pgfpathlineto{\pgfqpoint{5.779565in}{0.752081in}}%
\pgfpathlineto{\pgfqpoint{5.782921in}{0.743769in}}%
\pgfpathlineto{\pgfqpoint{5.786277in}{0.738694in}}%
\pgfpathlineto{\pgfqpoint{5.789632in}{0.737924in}}%
\pgfpathlineto{\pgfqpoint{5.796344in}{0.732809in}}%
\pgfpathlineto{\pgfqpoint{5.799700in}{0.735820in}}%
\pgfpathlineto{\pgfqpoint{5.803056in}{0.737365in}}%
\pgfpathlineto{\pgfqpoint{5.806411in}{0.728971in}}%
\pgfpathlineto{\pgfqpoint{5.809767in}{0.725355in}}%
\pgfpathlineto{\pgfqpoint{5.813123in}{0.729603in}}%
\pgfpathlineto{\pgfqpoint{5.816479in}{0.729046in}}%
\pgfpathlineto{\pgfqpoint{5.819834in}{0.720824in}}%
\pgfpathlineto{\pgfqpoint{5.823190in}{0.719018in}}%
\pgfpathlineto{\pgfqpoint{5.826546in}{0.731374in}}%
\pgfpathlineto{\pgfqpoint{5.829902in}{0.733731in}}%
\pgfpathlineto{\pgfqpoint{5.833258in}{0.714465in}}%
\pgfpathlineto{\pgfqpoint{5.836613in}{0.721342in}}%
\pgfpathlineto{\pgfqpoint{5.839969in}{0.739839in}}%
\pgfpathlineto{\pgfqpoint{5.843325in}{0.742231in}}%
\pgfpathlineto{\pgfqpoint{5.846681in}{0.727448in}}%
\pgfpathlineto{\pgfqpoint{5.850036in}{0.701748in}}%
\pgfpathlineto{\pgfqpoint{5.853392in}{0.739099in}}%
\pgfpathlineto{\pgfqpoint{5.856748in}{0.753535in}}%
\pgfpathlineto{\pgfqpoint{5.860104in}{0.743179in}}%
\pgfpathlineto{\pgfqpoint{5.863460in}{0.707979in}}%
\pgfpathlineto{\pgfqpoint{5.866815in}{0.701128in}}%
\pgfpathlineto{\pgfqpoint{5.870171in}{0.712343in}}%
\pgfpathlineto{\pgfqpoint{5.873527in}{0.692095in}}%
\pgfpathlineto{\pgfqpoint{5.876883in}{0.709730in}}%
\pgfpathlineto{\pgfqpoint{5.880238in}{0.735214in}}%
\pgfpathlineto{\pgfqpoint{5.883594in}{0.730445in}}%
\pgfpathlineto{\pgfqpoint{5.886950in}{0.698955in}}%
\pgfpathlineto{\pgfqpoint{5.890306in}{0.682480in}}%
\pgfpathlineto{\pgfqpoint{5.893662in}{0.697239in}}%
\pgfpathlineto{\pgfqpoint{5.897017in}{0.681556in}}%
\pgfpathlineto{\pgfqpoint{5.900373in}{0.688524in}}%
\pgfpathlineto{\pgfqpoint{5.903729in}{0.717802in}}%
\pgfpathlineto{\pgfqpoint{5.907085in}{0.717271in}}%
\pgfpathlineto{\pgfqpoint{5.910441in}{0.678971in}}%
\pgfpathlineto{\pgfqpoint{5.913796in}{0.692407in}}%
\pgfpathlineto{\pgfqpoint{5.917152in}{0.712893in}}%
\pgfpathlineto{\pgfqpoint{5.920508in}{0.708289in}}%
\pgfpathlineto{\pgfqpoint{5.923864in}{0.667359in}}%
\pgfpathlineto{\pgfqpoint{5.927219in}{0.702141in}}%
\pgfpathlineto{\pgfqpoint{5.930575in}{0.727055in}}%
\pgfpathlineto{\pgfqpoint{5.933931in}{0.722266in}}%
\pgfpathlineto{\pgfqpoint{5.937287in}{0.650662in}}%
\pgfpathlineto{\pgfqpoint{5.940643in}{0.720938in}}%
\pgfpathlineto{\pgfqpoint{5.943998in}{0.738721in}}%
\pgfpathlineto{\pgfqpoint{5.947354in}{0.736180in}}%
\pgfpathlineto{\pgfqpoint{5.950710in}{0.706315in}}%
\pgfpathlineto{\pgfqpoint{5.954066in}{0.704187in}}%
\pgfpathlineto{\pgfqpoint{5.957421in}{0.744538in}}%
\pgfpathlineto{\pgfqpoint{5.960777in}{0.751832in}}%
\pgfpathlineto{\pgfqpoint{5.964133in}{0.731217in}}%
\pgfpathlineto{\pgfqpoint{5.967489in}{0.653821in}}%
\pgfpathlineto{\pgfqpoint{5.970845in}{0.701724in}}%
\pgfpathlineto{\pgfqpoint{5.974200in}{0.703555in}}%
\pgfpathlineto{\pgfqpoint{5.977556in}{0.622878in}}%
\pgfpathlineto{\pgfqpoint{5.980912in}{0.718981in}}%
\pgfpathlineto{\pgfqpoint{5.984268in}{0.736616in}}%
\pgfpathlineto{\pgfqpoint{5.987623in}{0.721328in}}%
\pgfpathlineto{\pgfqpoint{5.990979in}{0.663795in}}%
\pgfpathlineto{\pgfqpoint{5.994335in}{0.681878in}}%
\pgfpathlineto{\pgfqpoint{5.997691in}{0.691114in}}%
\pgfpathlineto{\pgfqpoint{6.001047in}{0.633989in}}%
\pgfpathlineto{\pgfqpoint{6.004402in}{0.697720in}}%
\pgfpathlineto{\pgfqpoint{6.007758in}{0.720877in}}%
\pgfpathlineto{\pgfqpoint{6.011114in}{0.709691in}}%
\pgfpathlineto{\pgfqpoint{6.014470in}{0.612253in}}%
\pgfpathlineto{\pgfqpoint{6.017825in}{0.700350in}}%
\pgfpathlineto{\pgfqpoint{6.021181in}{0.713320in}}%
\pgfpathlineto{\pgfqpoint{6.024537in}{0.700459in}}%
\pgfpathlineto{\pgfqpoint{6.027893in}{0.617172in}}%
\pgfpathlineto{\pgfqpoint{6.031249in}{0.713217in}}%
\pgfpathlineto{\pgfqpoint{6.034604in}{0.729279in}}%
\pgfpathlineto{\pgfqpoint{6.037960in}{0.712491in}}%
\pgfpathlineto{\pgfqpoint{6.041316in}{0.671383in}}%
\pgfpathlineto{\pgfqpoint{6.044672in}{0.728473in}}%
\pgfpathlineto{\pgfqpoint{6.048027in}{0.739057in}}%
\pgfpathlineto{\pgfqpoint{6.051383in}{0.730071in}}%
\pgfpathlineto{\pgfqpoint{6.054739in}{0.671006in}}%
\pgfpathlineto{\pgfqpoint{6.058095in}{0.723598in}}%
\pgfpathlineto{\pgfqpoint{6.061451in}{0.749147in}}%
\pgfpathlineto{\pgfqpoint{6.064806in}{0.747198in}}%
\pgfpathlineto{\pgfqpoint{6.068162in}{0.715753in}}%
\pgfpathlineto{\pgfqpoint{6.071518in}{0.659971in}}%
\pgfpathlineto{\pgfqpoint{6.074874in}{0.707779in}}%
\pgfpathlineto{\pgfqpoint{6.078230in}{0.693979in}}%
\pgfpathlineto{\pgfqpoint{6.081585in}{0.676865in}}%
\pgfpathlineto{\pgfqpoint{6.084941in}{0.728397in}}%
\pgfpathlineto{\pgfqpoint{6.088297in}{0.735536in}}%
\pgfpathlineto{\pgfqpoint{6.091653in}{0.709948in}}%
\pgfpathlineto{\pgfqpoint{6.095008in}{0.582671in}}%
\pgfpathlineto{\pgfqpoint{6.098364in}{0.691309in}}%
\pgfpathlineto{\pgfqpoint{6.101720in}{0.686090in}}%
\pgfpathlineto{\pgfqpoint{6.105076in}{0.632219in}}%
\pgfpathlineto{\pgfqpoint{6.108432in}{0.709362in}}%
\pgfpathlineto{\pgfqpoint{6.111787in}{0.722396in}}%
\pgfpathlineto{\pgfqpoint{6.115143in}{0.699009in}}%
\pgfpathlineto{\pgfqpoint{6.118499in}{0.655119in}}%
\pgfpathlineto{\pgfqpoint{6.121855in}{0.707766in}}%
\pgfpathlineto{\pgfqpoint{6.125210in}{0.712677in}}%
\pgfpathlineto{\pgfqpoint{6.128566in}{0.688587in}}%
\pgfpathlineto{\pgfqpoint{6.131922in}{0.675420in}}%
\pgfpathlineto{\pgfqpoint{6.135278in}{0.721795in}}%
\pgfpathlineto{\pgfqpoint{6.138634in}{0.729021in}}%
\pgfpathlineto{\pgfqpoint{6.141989in}{0.695432in}}%
\pgfpathlineto{\pgfqpoint{6.145345in}{0.700911in}}%
\pgfpathlineto{\pgfqpoint{6.148701in}{0.733500in}}%
\pgfpathlineto{\pgfqpoint{6.152057in}{0.737719in}}%
\pgfpathlineto{\pgfqpoint{6.155412in}{0.720398in}}%
\pgfpathlineto{\pgfqpoint{6.158768in}{0.655430in}}%
\pgfpathlineto{\pgfqpoint{6.162124in}{0.735659in}}%
\pgfpathlineto{\pgfqpoint{6.165480in}{0.751118in}}%
\pgfpathlineto{\pgfqpoint{6.168836in}{0.738046in}}%
\pgfpathlineto{\pgfqpoint{6.172191in}{0.516782in}}%
\pgfpathlineto{\pgfqpoint{6.175547in}{0.733122in}}%
\pgfpathlineto{\pgfqpoint{6.178903in}{0.745608in}}%
\pgfpathlineto{\pgfqpoint{6.182259in}{0.733559in}}%
\pgfpathlineto{\pgfqpoint{6.185614in}{0.598002in}}%
\pgfpathlineto{\pgfqpoint{6.188970in}{0.742918in}}%
\pgfpathlineto{\pgfqpoint{6.192326in}{0.760635in}}%
\pgfpathlineto{\pgfqpoint{6.195682in}{0.748975in}}%
\pgfpathlineto{\pgfqpoint{6.199038in}{0.699971in}}%
\pgfpathlineto{\pgfqpoint{6.202393in}{0.702311in}}%
\pgfpathlineto{\pgfqpoint{6.205749in}{0.718715in}}%
\pgfpathlineto{\pgfqpoint{6.209105in}{0.682113in}}%
\pgfpathlineto{\pgfqpoint{6.215816in}{0.744267in}}%
\pgfpathlineto{\pgfqpoint{6.219172in}{0.740111in}}%
\pgfpathlineto{\pgfqpoint{6.222528in}{0.706750in}}%
\pgfpathlineto{\pgfqpoint{6.225884in}{0.641957in}}%
\pgfpathlineto{\pgfqpoint{6.229240in}{0.688975in}}%
\pgfpathlineto{\pgfqpoint{6.232595in}{0.652613in}}%
\pgfpathlineto{\pgfqpoint{6.235951in}{0.693750in}}%
\pgfpathlineto{\pgfqpoint{6.239307in}{0.723299in}}%
\pgfpathlineto{\pgfqpoint{6.242663in}{0.717710in}}%
\pgfpathlineto{\pgfqpoint{6.246019in}{0.669202in}}%
\pgfpathlineto{\pgfqpoint{6.249374in}{0.680548in}}%
\pgfpathlineto{\pgfqpoint{6.252730in}{0.700925in}}%
\pgfpathlineto{\pgfqpoint{6.256086in}{0.686171in}}%
\pgfpathlineto{\pgfqpoint{6.259442in}{0.648299in}}%
\pgfpathlineto{\pgfqpoint{6.262797in}{0.710809in}}%
\pgfpathlineto{\pgfqpoint{6.266153in}{0.721485in}}%
\pgfpathlineto{\pgfqpoint{6.269509in}{0.693066in}}%
\pgfpathlineto{\pgfqpoint{6.272865in}{0.684796in}}%
\pgfpathlineto{\pgfqpoint{6.276221in}{0.722309in}}%
\pgfpathlineto{\pgfqpoint{6.279576in}{0.727643in}}%
\pgfpathlineto{\pgfqpoint{6.282932in}{0.711172in}}%
\pgfpathlineto{\pgfqpoint{6.286288in}{0.639972in}}%
\pgfpathlineto{\pgfqpoint{6.289644in}{0.725381in}}%
\pgfpathlineto{\pgfqpoint{6.292999in}{0.741166in}}%
\pgfpathlineto{\pgfqpoint{6.296355in}{0.727326in}}%
\pgfpathlineto{\pgfqpoint{6.299711in}{0.661340in}}%
\pgfpathlineto{\pgfqpoint{6.303067in}{0.736060in}}%
\pgfpathlineto{\pgfqpoint{6.306423in}{0.748285in}}%
\pgfpathlineto{\pgfqpoint{6.309778in}{0.740208in}}%
\pgfpathlineto{\pgfqpoint{6.313134in}{0.684707in}}%
\pgfpathlineto{\pgfqpoint{6.316490in}{0.732372in}}%
\pgfpathlineto{\pgfqpoint{6.319846in}{0.758667in}}%
\pgfpathlineto{\pgfqpoint{6.323201in}{0.756889in}}%
\pgfpathlineto{\pgfqpoint{6.326557in}{0.720361in}}%
\pgfpathlineto{\pgfqpoint{6.329913in}{0.704397in}}%
\pgfpathlineto{\pgfqpoint{6.333269in}{0.735722in}}%
\pgfpathlineto{\pgfqpoint{6.336625in}{0.728113in}}%
\pgfpathlineto{\pgfqpoint{6.339980in}{0.600994in}}%
\pgfpathlineto{\pgfqpoint{6.343336in}{0.740271in}}%
\pgfpathlineto{\pgfqpoint{6.346692in}{0.757243in}}%
\pgfpathlineto{\pgfqpoint{6.350048in}{0.743160in}}%
\pgfpathlineto{\pgfqpoint{6.353403in}{0.698460in}}%
\pgfpathlineto{\pgfqpoint{6.356759in}{0.676767in}}%
\pgfpathlineto{\pgfqpoint{6.360115in}{0.691605in}}%
\pgfpathlineto{\pgfqpoint{6.363471in}{0.622934in}}%
\pgfpathlineto{\pgfqpoint{6.366827in}{0.718060in}}%
\pgfpathlineto{\pgfqpoint{6.370182in}{0.733967in}}%
\pgfpathlineto{\pgfqpoint{6.373538in}{0.720717in}}%
\pgfpathlineto{\pgfqpoint{6.380250in}{0.688485in}}%
\pgfpathlineto{\pgfqpoint{6.383605in}{0.673850in}}%
\pgfpathlineto{\pgfqpoint{6.383605in}{0.673850in}}%
\pgfusepath{stroke}%
\end{pgfscope}%
\begin{pgfscope}%
\pgfpathrectangle{\pgfqpoint{4.621821in}{0.320625in}}{\pgfqpoint{1.845679in}{1.164375in}}%
\pgfusepath{clip}%
\pgfsetrectcap%
\pgfsetroundjoin%
\pgfsetlinewidth{0.602250pt}%
\definecolor{currentstroke}{rgb}{0.141176,0.345098,0.650980}%
\pgfsetstrokecolor{currentstroke}%
\pgfsetstrokeopacity{0.120000}%
\pgfsetdash{}{0pt}%
\pgfpathmoveto{\pgfqpoint{4.705715in}{1.477290in}}%
\pgfpathlineto{\pgfqpoint{4.712427in}{1.464886in}}%
\pgfpathlineto{\pgfqpoint{4.722494in}{1.450126in}}%
\pgfpathlineto{\pgfqpoint{4.735918in}{1.433216in}}%
\pgfpathlineto{\pgfqpoint{4.752696in}{1.414004in}}%
\pgfpathlineto{\pgfqpoint{4.772831in}{1.392624in}}%
\pgfpathlineto{\pgfqpoint{4.792966in}{1.372946in}}%
\pgfpathlineto{\pgfqpoint{4.813100in}{1.354784in}}%
\pgfpathlineto{\pgfqpoint{4.836591in}{1.335036in}}%
\pgfpathlineto{\pgfqpoint{4.866793in}{1.311192in}}%
\pgfpathlineto{\pgfqpoint{4.903707in}{1.283586in}}%
\pgfpathlineto{\pgfqpoint{4.947332in}{1.252409in}}%
\pgfpathlineto{\pgfqpoint{4.997668in}{1.217844in}}%
\pgfpathlineto{\pgfqpoint{5.051361in}{1.182319in}}%
\pgfpathlineto{\pgfqpoint{5.108409in}{1.145936in}}%
\pgfpathlineto{\pgfqpoint{5.165457in}{1.110870in}}%
\pgfpathlineto{\pgfqpoint{5.225861in}{1.075073in}}%
\pgfpathlineto{\pgfqpoint{5.289621in}{1.038593in}}%
\pgfpathlineto{\pgfqpoint{5.363448in}{0.997694in}}%
\pgfpathlineto{\pgfqpoint{5.450699in}{0.950675in}}%
\pgfpathlineto{\pgfqpoint{5.568151in}{0.888701in}}%
\pgfpathlineto{\pgfqpoint{5.722517in}{0.808354in}}%
\pgfpathlineto{\pgfqpoint{5.729228in}{0.805682in}}%
\pgfpathlineto{\pgfqpoint{5.735940in}{0.801376in}}%
\pgfpathlineto{\pgfqpoint{5.742652in}{0.798499in}}%
\pgfpathlineto{\pgfqpoint{5.749363in}{0.794500in}}%
\pgfpathlineto{\pgfqpoint{5.756075in}{0.793122in}}%
\pgfpathlineto{\pgfqpoint{5.759430in}{0.789893in}}%
\pgfpathlineto{\pgfqpoint{5.762786in}{0.787451in}}%
\pgfpathlineto{\pgfqpoint{5.769498in}{0.784908in}}%
\pgfpathlineto{\pgfqpoint{5.772854in}{0.782400in}}%
\pgfpathlineto{\pgfqpoint{5.776209in}{0.780937in}}%
\pgfpathlineto{\pgfqpoint{5.779565in}{0.781690in}}%
\pgfpathlineto{\pgfqpoint{5.782921in}{0.780440in}}%
\pgfpathlineto{\pgfqpoint{5.786277in}{0.776029in}}%
\pgfpathlineto{\pgfqpoint{5.789632in}{0.773633in}}%
\pgfpathlineto{\pgfqpoint{5.792988in}{0.772880in}}%
\pgfpathlineto{\pgfqpoint{5.796344in}{0.771009in}}%
\pgfpathlineto{\pgfqpoint{5.799700in}{0.768326in}}%
\pgfpathlineto{\pgfqpoint{5.803056in}{0.768304in}}%
\pgfpathlineto{\pgfqpoint{5.806411in}{0.769896in}}%
\pgfpathlineto{\pgfqpoint{5.809767in}{0.766330in}}%
\pgfpathlineto{\pgfqpoint{5.813123in}{0.761521in}}%
\pgfpathlineto{\pgfqpoint{5.816479in}{0.760567in}}%
\pgfpathlineto{\pgfqpoint{5.819834in}{0.760204in}}%
\pgfpathlineto{\pgfqpoint{5.826546in}{0.754612in}}%
\pgfpathlineto{\pgfqpoint{5.829902in}{0.757314in}}%
\pgfpathlineto{\pgfqpoint{5.833258in}{0.759273in}}%
\pgfpathlineto{\pgfqpoint{5.836613in}{0.752515in}}%
\pgfpathlineto{\pgfqpoint{5.839969in}{0.747604in}}%
\pgfpathlineto{\pgfqpoint{5.843325in}{0.750023in}}%
\pgfpathlineto{\pgfqpoint{5.846681in}{0.750352in}}%
\pgfpathlineto{\pgfqpoint{5.850036in}{0.744758in}}%
\pgfpathlineto{\pgfqpoint{5.853392in}{0.741367in}}%
\pgfpathlineto{\pgfqpoint{5.856748in}{0.749853in}}%
\pgfpathlineto{\pgfqpoint{5.860104in}{0.754081in}}%
\pgfpathlineto{\pgfqpoint{5.863460in}{0.742628in}}%
\pgfpathlineto{\pgfqpoint{5.866815in}{0.734264in}}%
\pgfpathlineto{\pgfqpoint{5.870171in}{0.743273in}}%
\pgfpathlineto{\pgfqpoint{5.873527in}{0.745715in}}%
\pgfpathlineto{\pgfqpoint{5.876883in}{0.735924in}}%
\pgfpathlineto{\pgfqpoint{5.880238in}{0.728347in}}%
\pgfpathlineto{\pgfqpoint{5.883594in}{0.747134in}}%
\pgfpathlineto{\pgfqpoint{5.886950in}{0.755393in}}%
\pgfpathlineto{\pgfqpoint{5.890306in}{0.740791in}}%
\pgfpathlineto{\pgfqpoint{5.893662in}{0.719790in}}%
\pgfpathlineto{\pgfqpoint{5.897017in}{0.735327in}}%
\pgfpathlineto{\pgfqpoint{5.900373in}{0.740233in}}%
\pgfpathlineto{\pgfqpoint{5.903729in}{0.726570in}}%
\pgfpathlineto{\pgfqpoint{5.907085in}{0.718503in}}%
\pgfpathlineto{\pgfqpoint{5.910441in}{0.746663in}}%
\pgfpathlineto{\pgfqpoint{5.913796in}{0.754206in}}%
\pgfpathlineto{\pgfqpoint{5.917152in}{0.735078in}}%
\pgfpathlineto{\pgfqpoint{5.920508in}{0.705905in}}%
\pgfpathlineto{\pgfqpoint{5.923864in}{0.725414in}}%
\pgfpathlineto{\pgfqpoint{5.927219in}{0.727913in}}%
\pgfpathlineto{\pgfqpoint{5.930575in}{0.706560in}}%
\pgfpathlineto{\pgfqpoint{5.933931in}{0.720347in}}%
\pgfpathlineto{\pgfqpoint{5.937287in}{0.744950in}}%
\pgfpathlineto{\pgfqpoint{5.940643in}{0.742809in}}%
\pgfpathlineto{\pgfqpoint{5.943998in}{0.712931in}}%
\pgfpathlineto{\pgfqpoint{5.947354in}{0.701646in}}%
\pgfpathlineto{\pgfqpoint{5.950710in}{0.720921in}}%
\pgfpathlineto{\pgfqpoint{5.954066in}{0.714160in}}%
\pgfpathlineto{\pgfqpoint{5.957421in}{0.687138in}}%
\pgfpathlineto{\pgfqpoint{5.960777in}{0.725650in}}%
\pgfpathlineto{\pgfqpoint{5.964133in}{0.740840in}}%
\pgfpathlineto{\pgfqpoint{5.967489in}{0.726555in}}%
\pgfpathlineto{\pgfqpoint{5.970845in}{0.680506in}}%
\pgfpathlineto{\pgfqpoint{5.974200in}{0.714311in}}%
\pgfpathlineto{\pgfqpoint{5.977556in}{0.724307in}}%
\pgfpathlineto{\pgfqpoint{5.980912in}{0.708709in}}%
\pgfpathlineto{\pgfqpoint{5.984268in}{0.683692in}}%
\pgfpathlineto{\pgfqpoint{5.987623in}{0.730350in}}%
\pgfpathlineto{\pgfqpoint{5.990979in}{0.741580in}}%
\pgfpathlineto{\pgfqpoint{5.994335in}{0.721373in}}%
\pgfpathlineto{\pgfqpoint{5.997691in}{0.676469in}}%
\pgfpathlineto{\pgfqpoint{6.001047in}{0.725804in}}%
\pgfpathlineto{\pgfqpoint{6.004402in}{0.733312in}}%
\pgfpathlineto{\pgfqpoint{6.007758in}{0.716474in}}%
\pgfpathlineto{\pgfqpoint{6.011114in}{0.676814in}}%
\pgfpathlineto{\pgfqpoint{6.014470in}{0.736010in}}%
\pgfpathlineto{\pgfqpoint{6.017825in}{0.749255in}}%
\pgfpathlineto{\pgfqpoint{6.021181in}{0.732591in}}%
\pgfpathlineto{\pgfqpoint{6.024537in}{0.653443in}}%
\pgfpathlineto{\pgfqpoint{6.027893in}{0.727364in}}%
\pgfpathlineto{\pgfqpoint{6.031249in}{0.737647in}}%
\pgfpathlineto{\pgfqpoint{6.034604in}{0.722232in}}%
\pgfpathlineto{\pgfqpoint{6.037960in}{0.670480in}}%
\pgfpathlineto{\pgfqpoint{6.041316in}{0.740189in}}%
\pgfpathlineto{\pgfqpoint{6.044672in}{0.754364in}}%
\pgfpathlineto{\pgfqpoint{6.048027in}{0.738970in}}%
\pgfpathlineto{\pgfqpoint{6.051383in}{0.663283in}}%
\pgfpathlineto{\pgfqpoint{6.054739in}{0.718311in}}%
\pgfpathlineto{\pgfqpoint{6.058095in}{0.729621in}}%
\pgfpathlineto{\pgfqpoint{6.061451in}{0.707812in}}%
\pgfpathlineto{\pgfqpoint{6.064806in}{0.696231in}}%
\pgfpathlineto{\pgfqpoint{6.068162in}{0.742379in}}%
\pgfpathlineto{\pgfqpoint{6.071518in}{0.749500in}}%
\pgfpathlineto{\pgfqpoint{6.074874in}{0.725192in}}%
\pgfpathlineto{\pgfqpoint{6.078230in}{0.635927in}}%
\pgfpathlineto{\pgfqpoint{6.081585in}{0.713630in}}%
\pgfpathlineto{\pgfqpoint{6.084941in}{0.714962in}}%
\pgfpathlineto{\pgfqpoint{6.088297in}{0.659077in}}%
\pgfpathlineto{\pgfqpoint{6.091653in}{0.715815in}}%
\pgfpathlineto{\pgfqpoint{6.095008in}{0.740400in}}%
\pgfpathlineto{\pgfqpoint{6.098364in}{0.733319in}}%
\pgfpathlineto{\pgfqpoint{6.101720in}{0.684857in}}%
\pgfpathlineto{\pgfqpoint{6.105076in}{0.697188in}}%
\pgfpathlineto{\pgfqpoint{6.108432in}{0.718154in}}%
\pgfpathlineto{\pgfqpoint{6.111787in}{0.706336in}}%
\pgfpathlineto{\pgfqpoint{6.115143in}{0.643028in}}%
\pgfpathlineto{\pgfqpoint{6.118499in}{0.723964in}}%
\pgfpathlineto{\pgfqpoint{6.121855in}{0.738684in}}%
\pgfpathlineto{\pgfqpoint{6.125210in}{0.721363in}}%
\pgfpathlineto{\pgfqpoint{6.128566in}{0.625485in}}%
\pgfpathlineto{\pgfqpoint{6.131922in}{0.718638in}}%
\pgfpathlineto{\pgfqpoint{6.135278in}{0.728575in}}%
\pgfpathlineto{\pgfqpoint{6.138634in}{0.713051in}}%
\pgfpathlineto{\pgfqpoint{6.141989in}{0.654015in}}%
\pgfpathlineto{\pgfqpoint{6.145345in}{0.730518in}}%
\pgfpathlineto{\pgfqpoint{6.148701in}{0.745062in}}%
\pgfpathlineto{\pgfqpoint{6.152057in}{0.729119in}}%
\pgfpathlineto{\pgfqpoint{6.155412in}{0.622274in}}%
\pgfpathlineto{\pgfqpoint{6.158768in}{0.726820in}}%
\pgfpathlineto{\pgfqpoint{6.162124in}{0.738038in}}%
\pgfpathlineto{\pgfqpoint{6.165480in}{0.725286in}}%
\pgfpathlineto{\pgfqpoint{6.168836in}{0.611827in}}%
\pgfpathlineto{\pgfqpoint{6.172191in}{0.735517in}}%
\pgfpathlineto{\pgfqpoint{6.175547in}{0.753327in}}%
\pgfpathlineto{\pgfqpoint{6.178903in}{0.742331in}}%
\pgfpathlineto{\pgfqpoint{6.182259in}{0.676690in}}%
\pgfpathlineto{\pgfqpoint{6.185614in}{0.719451in}}%
\pgfpathlineto{\pgfqpoint{6.188970in}{0.735503in}}%
\pgfpathlineto{\pgfqpoint{6.192326in}{0.722515in}}%
\pgfpathlineto{\pgfqpoint{6.195682in}{0.653932in}}%
\pgfpathlineto{\pgfqpoint{6.199038in}{0.739019in}}%
\pgfpathlineto{\pgfqpoint{6.202393in}{0.754095in}}%
\pgfpathlineto{\pgfqpoint{6.205749in}{0.739147in}}%
\pgfpathlineto{\pgfqpoint{6.209105in}{0.673481in}}%
\pgfpathlineto{\pgfqpoint{6.212461in}{0.709454in}}%
\pgfpathlineto{\pgfqpoint{6.215816in}{0.721590in}}%
\pgfpathlineto{\pgfqpoint{6.219172in}{0.693937in}}%
\pgfpathlineto{\pgfqpoint{6.222528in}{0.702072in}}%
\pgfpathlineto{\pgfqpoint{6.225884in}{0.740125in}}%
\pgfpathlineto{\pgfqpoint{6.229240in}{0.742793in}}%
\pgfpathlineto{\pgfqpoint{6.232595in}{0.712655in}}%
\pgfpathlineto{\pgfqpoint{6.235951in}{0.664729in}}%
\pgfpathlineto{\pgfqpoint{6.239307in}{0.712695in}}%
\pgfpathlineto{\pgfqpoint{6.242663in}{0.708927in}}%
\pgfpathlineto{\pgfqpoint{6.246019in}{0.607582in}}%
\pgfpathlineto{\pgfqpoint{6.249374in}{0.717339in}}%
\pgfpathlineto{\pgfqpoint{6.252730in}{0.737447in}}%
\pgfpathlineto{\pgfqpoint{6.256086in}{0.725890in}}%
\pgfpathlineto{\pgfqpoint{6.259442in}{0.649040in}}%
\pgfpathlineto{\pgfqpoint{6.262797in}{0.708078in}}%
\pgfpathlineto{\pgfqpoint{6.266153in}{0.722761in}}%
\pgfpathlineto{\pgfqpoint{6.269509in}{0.709845in}}%
\pgfpathlineto{\pgfqpoint{6.272865in}{0.632814in}}%
\pgfpathlineto{\pgfqpoint{6.276221in}{0.724725in}}%
\pgfpathlineto{\pgfqpoint{6.279576in}{0.740492in}}%
\pgfpathlineto{\pgfqpoint{6.282932in}{0.725200in}}%
\pgfpathlineto{\pgfqpoint{6.286288in}{0.618124in}}%
\pgfpathlineto{\pgfqpoint{6.289644in}{0.723825in}}%
\pgfpathlineto{\pgfqpoint{6.292999in}{0.735509in}}%
\pgfpathlineto{\pgfqpoint{6.296355in}{0.723974in}}%
\pgfpathlineto{\pgfqpoint{6.299711in}{0.586256in}}%
\pgfpathlineto{\pgfqpoint{6.303067in}{0.730466in}}%
\pgfpathlineto{\pgfqpoint{6.306423in}{0.750018in}}%
\pgfpathlineto{\pgfqpoint{6.309778in}{0.741040in}}%
\pgfpathlineto{\pgfqpoint{6.313134in}{0.674616in}}%
\pgfpathlineto{\pgfqpoint{6.316490in}{0.721981in}}%
\pgfpathlineto{\pgfqpoint{6.319846in}{0.739547in}}%
\pgfpathlineto{\pgfqpoint{6.323201in}{0.731040in}}%
\pgfpathlineto{\pgfqpoint{6.326557in}{0.646399in}}%
\pgfpathlineto{\pgfqpoint{6.329913in}{0.733776in}}%
\pgfpathlineto{\pgfqpoint{6.333269in}{0.755327in}}%
\pgfpathlineto{\pgfqpoint{6.336625in}{0.747800in}}%
\pgfpathlineto{\pgfqpoint{6.339980in}{0.702063in}}%
\pgfpathlineto{\pgfqpoint{6.343336in}{0.706018in}}%
\pgfpathlineto{\pgfqpoint{6.346692in}{0.728505in}}%
\pgfpathlineto{\pgfqpoint{6.350048in}{0.715339in}}%
\pgfpathlineto{\pgfqpoint{6.353403in}{0.668388in}}%
\pgfpathlineto{\pgfqpoint{6.356759in}{0.737689in}}%
\pgfpathlineto{\pgfqpoint{6.360115in}{0.750422in}}%
\pgfpathlineto{\pgfqpoint{6.363471in}{0.734366in}}%
\pgfpathlineto{\pgfqpoint{6.366827in}{0.697956in}}%
\pgfpathlineto{\pgfqpoint{6.370182in}{0.525856in}}%
\pgfpathlineto{\pgfqpoint{6.373538in}{0.621262in}}%
\pgfpathlineto{\pgfqpoint{6.376894in}{0.672253in}}%
\pgfpathlineto{\pgfqpoint{6.380250in}{0.709287in}}%
\pgfpathlineto{\pgfqpoint{6.383605in}{0.712891in}}%
\pgfpathlineto{\pgfqpoint{6.383605in}{0.712891in}}%
\pgfusepath{stroke}%
\end{pgfscope}%
\begin{pgfscope}%
\pgfpathrectangle{\pgfqpoint{4.621821in}{0.320625in}}{\pgfqpoint{1.845679in}{1.164375in}}%
\pgfusepath{clip}%
\pgfsetrectcap%
\pgfsetroundjoin%
\pgfsetlinewidth{0.602250pt}%
\definecolor{currentstroke}{rgb}{0.141176,0.345098,0.650980}%
\pgfsetstrokecolor{currentstroke}%
\pgfsetstrokeopacity{0.120000}%
\pgfsetdash{}{0pt}%
\pgfpathmoveto{\pgfqpoint{4.705715in}{1.473458in}}%
\pgfpathlineto{\pgfqpoint{4.712427in}{1.460817in}}%
\pgfpathlineto{\pgfqpoint{4.722494in}{1.444775in}}%
\pgfpathlineto{\pgfqpoint{4.739273in}{1.420598in}}%
\pgfpathlineto{\pgfqpoint{4.772831in}{1.374575in}}%
\pgfpathlineto{\pgfqpoint{4.792966in}{1.348740in}}%
\pgfpathlineto{\pgfqpoint{4.809745in}{1.328851in}}%
\pgfpathlineto{\pgfqpoint{4.826524in}{1.310490in}}%
\pgfpathlineto{\pgfqpoint{4.846658in}{1.290181in}}%
\pgfpathlineto{\pgfqpoint{4.870149in}{1.268251in}}%
\pgfpathlineto{\pgfqpoint{4.896995in}{1.244725in}}%
\pgfpathlineto{\pgfqpoint{4.933909in}{1.213999in}}%
\pgfpathlineto{\pgfqpoint{4.977534in}{1.179171in}}%
\pgfpathlineto{\pgfqpoint{5.027870in}{1.140408in}}%
\pgfpathlineto{\pgfqpoint{5.078207in}{1.103029in}}%
\pgfpathlineto{\pgfqpoint{5.125188in}{1.069460in}}%
\pgfpathlineto{\pgfqpoint{5.172169in}{1.037248in}}%
\pgfpathlineto{\pgfqpoint{5.219150in}{1.006391in}}%
\pgfpathlineto{\pgfqpoint{5.269487in}{0.974683in}}%
\pgfpathlineto{\pgfqpoint{5.326535in}{0.940104in}}%
\pgfpathlineto{\pgfqpoint{5.400362in}{0.896762in}}%
\pgfpathlineto{\pgfqpoint{5.487612in}{0.846869in}}%
\pgfpathlineto{\pgfqpoint{5.527882in}{0.823930in}}%
\pgfpathlineto{\pgfqpoint{5.578218in}{0.795651in}}%
\pgfpathlineto{\pgfqpoint{5.584930in}{0.792578in}}%
\pgfpathlineto{\pgfqpoint{5.591641in}{0.787847in}}%
\pgfpathlineto{\pgfqpoint{5.601709in}{0.782859in}}%
\pgfpathlineto{\pgfqpoint{5.608420in}{0.778500in}}%
\pgfpathlineto{\pgfqpoint{5.615132in}{0.777501in}}%
\pgfpathlineto{\pgfqpoint{5.621843in}{0.770862in}}%
\pgfpathlineto{\pgfqpoint{5.625199in}{0.769575in}}%
\pgfpathlineto{\pgfqpoint{5.628555in}{0.769278in}}%
\pgfpathlineto{\pgfqpoint{5.631911in}{0.767501in}}%
\pgfpathlineto{\pgfqpoint{5.635267in}{0.763920in}}%
\pgfpathlineto{\pgfqpoint{5.638622in}{0.761634in}}%
\pgfpathlineto{\pgfqpoint{5.645334in}{0.765961in}}%
\pgfpathlineto{\pgfqpoint{5.648690in}{0.761354in}}%
\pgfpathlineto{\pgfqpoint{5.652045in}{0.753999in}}%
\pgfpathlineto{\pgfqpoint{5.655401in}{0.754923in}}%
\pgfpathlineto{\pgfqpoint{5.658757in}{0.759277in}}%
\pgfpathlineto{\pgfqpoint{5.662113in}{0.758749in}}%
\pgfpathlineto{\pgfqpoint{5.668824in}{0.744402in}}%
\pgfpathlineto{\pgfqpoint{5.675536in}{0.763800in}}%
\pgfpathlineto{\pgfqpoint{5.678892in}{0.761507in}}%
\pgfpathlineto{\pgfqpoint{5.682247in}{0.744405in}}%
\pgfpathlineto{\pgfqpoint{5.685603in}{0.736870in}}%
\pgfpathlineto{\pgfqpoint{5.688959in}{0.749122in}}%
\pgfpathlineto{\pgfqpoint{5.692315in}{0.752562in}}%
\pgfpathlineto{\pgfqpoint{5.695671in}{0.743263in}}%
\pgfpathlineto{\pgfqpoint{5.699026in}{0.727284in}}%
\pgfpathlineto{\pgfqpoint{5.702382in}{0.747053in}}%
\pgfpathlineto{\pgfqpoint{5.705738in}{0.763168in}}%
\pgfpathlineto{\pgfqpoint{5.709094in}{0.761250in}}%
\pgfpathlineto{\pgfqpoint{5.715805in}{0.717947in}}%
\pgfpathlineto{\pgfqpoint{5.719161in}{0.732233in}}%
\pgfpathlineto{\pgfqpoint{5.722517in}{0.735639in}}%
\pgfpathlineto{\pgfqpoint{5.725873in}{0.721295in}}%
\pgfpathlineto{\pgfqpoint{5.729228in}{0.717957in}}%
\pgfpathlineto{\pgfqpoint{5.732584in}{0.745418in}}%
\pgfpathlineto{\pgfqpoint{5.735940in}{0.754945in}}%
\pgfpathlineto{\pgfqpoint{5.739296in}{0.742350in}}%
\pgfpathlineto{\pgfqpoint{5.742652in}{0.712577in}}%
\pgfpathlineto{\pgfqpoint{5.746007in}{0.708787in}}%
\pgfpathlineto{\pgfqpoint{5.749363in}{0.721875in}}%
\pgfpathlineto{\pgfqpoint{5.752719in}{0.715647in}}%
\pgfpathlineto{\pgfqpoint{5.756075in}{0.695161in}}%
\pgfpathlineto{\pgfqpoint{5.759430in}{0.723234in}}%
\pgfpathlineto{\pgfqpoint{5.762786in}{0.741415in}}%
\pgfpathlineto{\pgfqpoint{5.766142in}{0.737817in}}%
\pgfpathlineto{\pgfqpoint{5.769498in}{0.709928in}}%
\pgfpathlineto{\pgfqpoint{5.772854in}{0.694685in}}%
\pgfpathlineto{\pgfqpoint{5.776209in}{0.720679in}}%
\pgfpathlineto{\pgfqpoint{5.779565in}{0.724171in}}%
\pgfpathlineto{\pgfqpoint{5.782921in}{0.707259in}}%
\pgfpathlineto{\pgfqpoint{5.786277in}{0.685308in}}%
\pgfpathlineto{\pgfqpoint{5.789632in}{0.727391in}}%
\pgfpathlineto{\pgfqpoint{5.792988in}{0.741935in}}%
\pgfpathlineto{\pgfqpoint{5.796344in}{0.733612in}}%
\pgfpathlineto{\pgfqpoint{5.799700in}{0.687747in}}%
\pgfpathlineto{\pgfqpoint{5.806411in}{0.736535in}}%
\pgfpathlineto{\pgfqpoint{5.809767in}{0.738980in}}%
\pgfpathlineto{\pgfqpoint{5.813123in}{0.723397in}}%
\pgfpathlineto{\pgfqpoint{5.816479in}{0.661682in}}%
\pgfpathlineto{\pgfqpoint{5.819834in}{0.731540in}}%
\pgfpathlineto{\pgfqpoint{5.823190in}{0.751922in}}%
\pgfpathlineto{\pgfqpoint{5.826546in}{0.751183in}}%
\pgfpathlineto{\pgfqpoint{5.829902in}{0.720951in}}%
\pgfpathlineto{\pgfqpoint{5.833258in}{0.699614in}}%
\pgfpathlineto{\pgfqpoint{5.836613in}{0.741204in}}%
\pgfpathlineto{\pgfqpoint{5.839969in}{0.749267in}}%
\pgfpathlineto{\pgfqpoint{5.843325in}{0.739773in}}%
\pgfpathlineto{\pgfqpoint{5.846681in}{0.687718in}}%
\pgfpathlineto{\pgfqpoint{5.850036in}{0.730308in}}%
\pgfpathlineto{\pgfqpoint{5.853392in}{0.758771in}}%
\pgfpathlineto{\pgfqpoint{5.856748in}{0.763739in}}%
\pgfpathlineto{\pgfqpoint{5.860104in}{0.744871in}}%
\pgfpathlineto{\pgfqpoint{5.863460in}{0.670771in}}%
\pgfpathlineto{\pgfqpoint{5.866815in}{0.723560in}}%
\pgfpathlineto{\pgfqpoint{5.870171in}{0.738870in}}%
\pgfpathlineto{\pgfqpoint{5.873527in}{0.729092in}}%
\pgfpathlineto{\pgfqpoint{5.876883in}{0.641280in}}%
\pgfpathlineto{\pgfqpoint{5.880238in}{0.734576in}}%
\pgfpathlineto{\pgfqpoint{5.883594in}{0.757309in}}%
\pgfpathlineto{\pgfqpoint{5.886950in}{0.756049in}}%
\pgfpathlineto{\pgfqpoint{5.890306in}{0.731370in}}%
\pgfpathlineto{\pgfqpoint{5.893662in}{0.633286in}}%
\pgfpathlineto{\pgfqpoint{5.897017in}{0.709838in}}%
\pgfpathlineto{\pgfqpoint{5.900373in}{0.716702in}}%
\pgfpathlineto{\pgfqpoint{5.903729in}{0.684405in}}%
\pgfpathlineto{\pgfqpoint{5.907085in}{0.701756in}}%
\pgfpathlineto{\pgfqpoint{5.910441in}{0.738104in}}%
\pgfpathlineto{\pgfqpoint{5.913796in}{0.745330in}}%
\pgfpathlineto{\pgfqpoint{5.917152in}{0.726859in}}%
\pgfpathlineto{\pgfqpoint{5.920508in}{0.660680in}}%
\pgfpathlineto{\pgfqpoint{5.923864in}{0.699716in}}%
\pgfpathlineto{\pgfqpoint{5.927219in}{0.715794in}}%
\pgfpathlineto{\pgfqpoint{5.930575in}{0.704376in}}%
\pgfpathlineto{\pgfqpoint{5.933931in}{0.607218in}}%
\pgfpathlineto{\pgfqpoint{5.937287in}{0.716978in}}%
\pgfpathlineto{\pgfqpoint{5.940643in}{0.737020in}}%
\pgfpathlineto{\pgfqpoint{5.943998in}{0.732744in}}%
\pgfpathlineto{\pgfqpoint{5.947354in}{0.693694in}}%
\pgfpathlineto{\pgfqpoint{5.950710in}{0.694241in}}%
\pgfpathlineto{\pgfqpoint{5.954066in}{0.726570in}}%
\pgfpathlineto{\pgfqpoint{5.957421in}{0.731222in}}%
\pgfpathlineto{\pgfqpoint{5.960777in}{0.716837in}}%
\pgfpathlineto{\pgfqpoint{5.964133in}{0.583134in}}%
\pgfpathlineto{\pgfqpoint{5.967489in}{0.723189in}}%
\pgfpathlineto{\pgfqpoint{5.970845in}{0.744741in}}%
\pgfpathlineto{\pgfqpoint{5.974200in}{0.744256in}}%
\pgfpathlineto{\pgfqpoint{5.977556in}{0.711520in}}%
\pgfpathlineto{\pgfqpoint{5.980912in}{0.701449in}}%
\pgfpathlineto{\pgfqpoint{5.984268in}{0.740055in}}%
\pgfpathlineto{\pgfqpoint{5.987623in}{0.748272in}}%
\pgfpathlineto{\pgfqpoint{5.990979in}{0.740338in}}%
\pgfpathlineto{\pgfqpoint{5.994335in}{0.699111in}}%
\pgfpathlineto{\pgfqpoint{5.997691in}{0.720341in}}%
\pgfpathlineto{\pgfqpoint{6.001047in}{0.754110in}}%
\pgfpathlineto{\pgfqpoint{6.004402in}{0.762298in}}%
\pgfpathlineto{\pgfqpoint{6.007758in}{0.747481in}}%
\pgfpathlineto{\pgfqpoint{6.011114in}{0.679499in}}%
\pgfpathlineto{\pgfqpoint{6.014470in}{0.727432in}}%
\pgfpathlineto{\pgfqpoint{6.017825in}{0.746766in}}%
\pgfpathlineto{\pgfqpoint{6.021181in}{0.743644in}}%
\pgfpathlineto{\pgfqpoint{6.024537in}{0.710317in}}%
\pgfpathlineto{\pgfqpoint{6.027893in}{0.717733in}}%
\pgfpathlineto{\pgfqpoint{6.031249in}{0.756161in}}%
\pgfpathlineto{\pgfqpoint{6.034604in}{0.765298in}}%
\pgfpathlineto{\pgfqpoint{6.037960in}{0.751115in}}%
\pgfpathlineto{\pgfqpoint{6.041316in}{0.707776in}}%
\pgfpathlineto{\pgfqpoint{6.044672in}{0.698448in}}%
\pgfpathlineto{\pgfqpoint{6.048027in}{0.725196in}}%
\pgfpathlineto{\pgfqpoint{6.051383in}{0.715922in}}%
\pgfpathlineto{\pgfqpoint{6.054739in}{0.619168in}}%
\pgfpathlineto{\pgfqpoint{6.058095in}{0.731269in}}%
\pgfpathlineto{\pgfqpoint{6.061451in}{0.751096in}}%
\pgfpathlineto{\pgfqpoint{6.064806in}{0.746156in}}%
\pgfpathlineto{\pgfqpoint{6.068162in}{0.717375in}}%
\pgfpathlineto{\pgfqpoint{6.071518in}{0.628175in}}%
\pgfpathlineto{\pgfqpoint{6.074874in}{0.706248in}}%
\pgfpathlineto{\pgfqpoint{6.078230in}{0.707962in}}%
\pgfpathlineto{\pgfqpoint{6.081585in}{0.661162in}}%
\pgfpathlineto{\pgfqpoint{6.084941in}{0.703668in}}%
\pgfpathlineto{\pgfqpoint{6.088297in}{0.733596in}}%
\pgfpathlineto{\pgfqpoint{6.091653in}{0.737054in}}%
\pgfpathlineto{\pgfqpoint{6.095008in}{0.712677in}}%
\pgfpathlineto{\pgfqpoint{6.098364in}{0.634762in}}%
\pgfpathlineto{\pgfqpoint{6.101720in}{0.711983in}}%
\pgfpathlineto{\pgfqpoint{6.105076in}{0.722213in}}%
\pgfpathlineto{\pgfqpoint{6.108432in}{0.711094in}}%
\pgfpathlineto{\pgfqpoint{6.111787in}{0.621063in}}%
\pgfpathlineto{\pgfqpoint{6.115143in}{0.714049in}}%
\pgfpathlineto{\pgfqpoint{6.118499in}{0.737392in}}%
\pgfpathlineto{\pgfqpoint{6.121855in}{0.737425in}}%
\pgfpathlineto{\pgfqpoint{6.125210in}{0.704228in}}%
\pgfpathlineto{\pgfqpoint{6.128566in}{0.696524in}}%
\pgfpathlineto{\pgfqpoint{6.131922in}{0.734428in}}%
\pgfpathlineto{\pgfqpoint{6.135278in}{0.742686in}}%
\pgfpathlineto{\pgfqpoint{6.138634in}{0.735085in}}%
\pgfpathlineto{\pgfqpoint{6.141989in}{0.695923in}}%
\pgfpathlineto{\pgfqpoint{6.145345in}{0.712583in}}%
\pgfpathlineto{\pgfqpoint{6.148701in}{0.747676in}}%
\pgfpathlineto{\pgfqpoint{6.152057in}{0.756545in}}%
\pgfpathlineto{\pgfqpoint{6.155412in}{0.742230in}}%
\pgfpathlineto{\pgfqpoint{6.158768in}{0.651436in}}%
\pgfpathlineto{\pgfqpoint{6.162124in}{0.732739in}}%
\pgfpathlineto{\pgfqpoint{6.165480in}{0.751411in}}%
\pgfpathlineto{\pgfqpoint{6.168836in}{0.750712in}}%
\pgfpathlineto{\pgfqpoint{6.172191in}{0.728717in}}%
\pgfpathlineto{\pgfqpoint{6.175547in}{0.689914in}}%
\pgfpathlineto{\pgfqpoint{6.178903in}{0.751472in}}%
\pgfpathlineto{\pgfqpoint{6.182259in}{0.767176in}}%
\pgfpathlineto{\pgfqpoint{6.185614in}{0.761298in}}%
\pgfpathlineto{\pgfqpoint{6.188970in}{0.729807in}}%
\pgfpathlineto{\pgfqpoint{6.192326in}{0.684348in}}%
\pgfpathlineto{\pgfqpoint{6.195682in}{0.733468in}}%
\pgfpathlineto{\pgfqpoint{6.199038in}{0.736052in}}%
\pgfpathlineto{\pgfqpoint{6.202393in}{0.706203in}}%
\pgfpathlineto{\pgfqpoint{6.205749in}{0.711296in}}%
\pgfpathlineto{\pgfqpoint{6.209105in}{0.751354in}}%
\pgfpathlineto{\pgfqpoint{6.212461in}{0.760456in}}%
\pgfpathlineto{\pgfqpoint{6.215816in}{0.745449in}}%
\pgfpathlineto{\pgfqpoint{6.219172in}{0.704345in}}%
\pgfpathlineto{\pgfqpoint{6.222528in}{0.682153in}}%
\pgfpathlineto{\pgfqpoint{6.225884in}{0.711088in}}%
\pgfpathlineto{\pgfqpoint{6.229240in}{0.696257in}}%
\pgfpathlineto{\pgfqpoint{6.232595in}{0.671042in}}%
\pgfpathlineto{\pgfqpoint{6.235951in}{0.728489in}}%
\pgfpathlineto{\pgfqpoint{6.239307in}{0.742790in}}%
\pgfpathlineto{\pgfqpoint{6.242663in}{0.731661in}}%
\pgfpathlineto{\pgfqpoint{6.246019in}{0.689749in}}%
\pgfpathlineto{\pgfqpoint{6.249374in}{0.682199in}}%
\pgfpathlineto{\pgfqpoint{6.252730in}{0.711205in}}%
\pgfpathlineto{\pgfqpoint{6.256086in}{0.707437in}}%
\pgfpathlineto{\pgfqpoint{6.259442in}{0.655207in}}%
\pgfpathlineto{\pgfqpoint{6.262797in}{0.702402in}}%
\pgfpathlineto{\pgfqpoint{6.266153in}{0.730889in}}%
\pgfpathlineto{\pgfqpoint{6.269509in}{0.733839in}}%
\pgfpathlineto{\pgfqpoint{6.276221in}{0.674649in}}%
\pgfpathlineto{\pgfqpoint{6.279576in}{0.723410in}}%
\pgfpathlineto{\pgfqpoint{6.282932in}{0.733267in}}%
\pgfpathlineto{\pgfqpoint{6.286288in}{0.726114in}}%
\pgfpathlineto{\pgfqpoint{6.289644in}{0.686089in}}%
\pgfpathlineto{\pgfqpoint{6.292999in}{0.705508in}}%
\pgfpathlineto{\pgfqpoint{6.296355in}{0.739806in}}%
\pgfpathlineto{\pgfqpoint{6.299711in}{0.747965in}}%
\pgfpathlineto{\pgfqpoint{6.303067in}{0.731229in}}%
\pgfpathlineto{\pgfqpoint{6.306423in}{0.635649in}}%
\pgfpathlineto{\pgfqpoint{6.309778in}{0.734972in}}%
\pgfpathlineto{\pgfqpoint{6.313134in}{0.751235in}}%
\pgfpathlineto{\pgfqpoint{6.316490in}{0.750705in}}%
\pgfpathlineto{\pgfqpoint{6.319846in}{0.731693in}}%
\pgfpathlineto{\pgfqpoint{6.323201in}{0.659006in}}%
\pgfpathlineto{\pgfqpoint{6.326557in}{0.745330in}}%
\pgfpathlineto{\pgfqpoint{6.329913in}{0.764181in}}%
\pgfpathlineto{\pgfqpoint{6.333269in}{0.762668in}}%
\pgfpathlineto{\pgfqpoint{6.336625in}{0.734505in}}%
\pgfpathlineto{\pgfqpoint{6.339980in}{0.688017in}}%
\pgfpathlineto{\pgfqpoint{6.343336in}{0.742091in}}%
\pgfpathlineto{\pgfqpoint{6.346692in}{0.749795in}}%
\pgfpathlineto{\pgfqpoint{6.350048in}{0.736147in}}%
\pgfpathlineto{\pgfqpoint{6.353403in}{0.582363in}}%
\pgfpathlineto{\pgfqpoint{6.356759in}{0.745089in}}%
\pgfpathlineto{\pgfqpoint{6.360115in}{0.765842in}}%
\pgfpathlineto{\pgfqpoint{6.363471in}{0.763648in}}%
\pgfpathlineto{\pgfqpoint{6.366827in}{0.747273in}}%
\pgfpathlineto{\pgfqpoint{6.373538in}{0.710293in}}%
\pgfpathlineto{\pgfqpoint{6.376894in}{0.703472in}}%
\pgfpathlineto{\pgfqpoint{6.380250in}{0.706011in}}%
\pgfpathlineto{\pgfqpoint{6.383605in}{0.702602in}}%
\pgfpathlineto{\pgfqpoint{6.383605in}{0.702602in}}%
\pgfusepath{stroke}%
\end{pgfscope}%
\begin{pgfscope}%
\pgfpathrectangle{\pgfqpoint{4.621821in}{0.320625in}}{\pgfqpoint{1.845679in}{1.164375in}}%
\pgfusepath{clip}%
\pgfsetrectcap%
\pgfsetroundjoin%
\pgfsetlinewidth{0.602250pt}%
\definecolor{currentstroke}{rgb}{0.141176,0.345098,0.650980}%
\pgfsetstrokecolor{currentstroke}%
\pgfsetstrokeopacity{0.120000}%
\pgfsetdash{}{0pt}%
\pgfpathmoveto{\pgfqpoint{4.705715in}{1.473855in}}%
\pgfpathlineto{\pgfqpoint{4.712427in}{1.462409in}}%
\pgfpathlineto{\pgfqpoint{4.725850in}{1.443133in}}%
\pgfpathlineto{\pgfqpoint{4.749341in}{1.411474in}}%
\pgfpathlineto{\pgfqpoint{4.766120in}{1.390691in}}%
\pgfpathlineto{\pgfqpoint{4.779543in}{1.375643in}}%
\pgfpathlineto{\pgfqpoint{4.796322in}{1.358544in}}%
\pgfpathlineto{\pgfqpoint{4.816456in}{1.339850in}}%
\pgfpathlineto{\pgfqpoint{4.839947in}{1.319712in}}%
\pgfpathlineto{\pgfqpoint{4.870149in}{1.295469in}}%
\pgfpathlineto{\pgfqpoint{4.910418in}{1.264771in}}%
\pgfpathlineto{\pgfqpoint{4.967466in}{1.222837in}}%
\pgfpathlineto{\pgfqpoint{5.051361in}{1.162677in}}%
\pgfpathlineto{\pgfqpoint{5.158746in}{1.087097in}}%
\pgfpathlineto{\pgfqpoint{5.276198in}{1.005797in}}%
\pgfpathlineto{\pgfqpoint{5.397006in}{0.923533in}}%
\pgfpathlineto{\pgfqpoint{5.487612in}{0.862758in}}%
\pgfpathlineto{\pgfqpoint{5.564795in}{0.812011in}}%
\pgfpathlineto{\pgfqpoint{5.568151in}{0.810483in}}%
\pgfpathlineto{\pgfqpoint{5.578218in}{0.803097in}}%
\pgfpathlineto{\pgfqpoint{5.581574in}{0.802147in}}%
\pgfpathlineto{\pgfqpoint{5.584930in}{0.800330in}}%
\pgfpathlineto{\pgfqpoint{5.591641in}{0.793788in}}%
\pgfpathlineto{\pgfqpoint{5.594997in}{0.793816in}}%
\pgfpathlineto{\pgfqpoint{5.598353in}{0.795632in}}%
\pgfpathlineto{\pgfqpoint{5.601709in}{0.793211in}}%
\pgfpathlineto{\pgfqpoint{5.605065in}{0.786555in}}%
\pgfpathlineto{\pgfqpoint{5.608420in}{0.782842in}}%
\pgfpathlineto{\pgfqpoint{5.611776in}{0.782681in}}%
\pgfpathlineto{\pgfqpoint{5.615132in}{0.781214in}}%
\pgfpathlineto{\pgfqpoint{5.618488in}{0.776900in}}%
\pgfpathlineto{\pgfqpoint{5.621843in}{0.774789in}}%
\pgfpathlineto{\pgfqpoint{5.625199in}{0.779309in}}%
\pgfpathlineto{\pgfqpoint{5.628555in}{0.781649in}}%
\pgfpathlineto{\pgfqpoint{5.635267in}{0.766041in}}%
\pgfpathlineto{\pgfqpoint{5.638622in}{0.763270in}}%
\pgfpathlineto{\pgfqpoint{5.641978in}{0.762198in}}%
\pgfpathlineto{\pgfqpoint{5.645334in}{0.758970in}}%
\pgfpathlineto{\pgfqpoint{5.648690in}{0.757234in}}%
\pgfpathlineto{\pgfqpoint{5.652045in}{0.761309in}}%
\pgfpathlineto{\pgfqpoint{5.655401in}{0.762916in}}%
\pgfpathlineto{\pgfqpoint{5.662113in}{0.747753in}}%
\pgfpathlineto{\pgfqpoint{5.665469in}{0.746785in}}%
\pgfpathlineto{\pgfqpoint{5.668824in}{0.746402in}}%
\pgfpathlineto{\pgfqpoint{5.672180in}{0.742363in}}%
\pgfpathlineto{\pgfqpoint{5.675536in}{0.739292in}}%
\pgfpathlineto{\pgfqpoint{5.682247in}{0.748914in}}%
\pgfpathlineto{\pgfqpoint{5.685603in}{0.740851in}}%
\pgfpathlineto{\pgfqpoint{5.688959in}{0.730050in}}%
\pgfpathlineto{\pgfqpoint{5.695671in}{0.742797in}}%
\pgfpathlineto{\pgfqpoint{5.699026in}{0.740128in}}%
\pgfpathlineto{\pgfqpoint{5.702382in}{0.726800in}}%
\pgfpathlineto{\pgfqpoint{5.705738in}{0.723507in}}%
\pgfpathlineto{\pgfqpoint{5.709094in}{0.743291in}}%
\pgfpathlineto{\pgfqpoint{5.712449in}{0.751829in}}%
\pgfpathlineto{\pgfqpoint{5.715805in}{0.739800in}}%
\pgfpathlineto{\pgfqpoint{5.719161in}{0.710638in}}%
\pgfpathlineto{\pgfqpoint{5.722517in}{0.742094in}}%
\pgfpathlineto{\pgfqpoint{5.725873in}{0.756959in}}%
\pgfpathlineto{\pgfqpoint{5.729228in}{0.756750in}}%
\pgfpathlineto{\pgfqpoint{5.732584in}{0.739888in}}%
\pgfpathlineto{\pgfqpoint{5.735940in}{0.701178in}}%
\pgfpathlineto{\pgfqpoint{5.739296in}{0.750382in}}%
\pgfpathlineto{\pgfqpoint{5.742652in}{0.768742in}}%
\pgfpathlineto{\pgfqpoint{5.746007in}{0.766866in}}%
\pgfpathlineto{\pgfqpoint{5.749363in}{0.736821in}}%
\pgfpathlineto{\pgfqpoint{5.752719in}{0.715229in}}%
\pgfpathlineto{\pgfqpoint{5.756075in}{0.752901in}}%
\pgfpathlineto{\pgfqpoint{5.759430in}{0.758521in}}%
\pgfpathlineto{\pgfqpoint{5.762786in}{0.743252in}}%
\pgfpathlineto{\pgfqpoint{5.766142in}{0.684858in}}%
\pgfpathlineto{\pgfqpoint{5.769498in}{0.755727in}}%
\pgfpathlineto{\pgfqpoint{5.772854in}{0.774176in}}%
\pgfpathlineto{\pgfqpoint{5.776209in}{0.769223in}}%
\pgfpathlineto{\pgfqpoint{5.779565in}{0.741553in}}%
\pgfpathlineto{\pgfqpoint{5.782921in}{0.669775in}}%
\pgfpathlineto{\pgfqpoint{5.786277in}{0.723439in}}%
\pgfpathlineto{\pgfqpoint{5.789632in}{0.722303in}}%
\pgfpathlineto{\pgfqpoint{5.792988in}{0.663614in}}%
\pgfpathlineto{\pgfqpoint{5.796344in}{0.734068in}}%
\pgfpathlineto{\pgfqpoint{5.799700in}{0.756088in}}%
\pgfpathlineto{\pgfqpoint{5.803056in}{0.751823in}}%
\pgfpathlineto{\pgfqpoint{5.806411in}{0.724697in}}%
\pgfpathlineto{\pgfqpoint{5.809767in}{0.652269in}}%
\pgfpathlineto{\pgfqpoint{5.813123in}{0.701057in}}%
\pgfpathlineto{\pgfqpoint{5.816479in}{0.698822in}}%
\pgfpathlineto{\pgfqpoint{5.819834in}{0.648207in}}%
\pgfpathlineto{\pgfqpoint{5.823190in}{0.717898in}}%
\pgfpathlineto{\pgfqpoint{5.826546in}{0.737234in}}%
\pgfpathlineto{\pgfqpoint{5.829902in}{0.729950in}}%
\pgfpathlineto{\pgfqpoint{5.833258in}{0.687626in}}%
\pgfpathlineto{\pgfqpoint{5.836613in}{0.691780in}}%
\pgfpathlineto{\pgfqpoint{5.839969in}{0.718955in}}%
\pgfpathlineto{\pgfqpoint{5.843325in}{0.718387in}}%
\pgfpathlineto{\pgfqpoint{5.846681in}{0.688488in}}%
\pgfpathlineto{\pgfqpoint{5.850036in}{0.690731in}}%
\pgfpathlineto{\pgfqpoint{5.853392in}{0.730496in}}%
\pgfpathlineto{\pgfqpoint{5.856748in}{0.739377in}}%
\pgfpathlineto{\pgfqpoint{5.860104in}{0.719498in}}%
\pgfpathlineto{\pgfqpoint{5.863460in}{0.674755in}}%
\pgfpathlineto{\pgfqpoint{5.866815in}{0.734561in}}%
\pgfpathlineto{\pgfqpoint{5.870171in}{0.748073in}}%
\pgfpathlineto{\pgfqpoint{5.873527in}{0.745781in}}%
\pgfpathlineto{\pgfqpoint{5.876883in}{0.723349in}}%
\pgfpathlineto{\pgfqpoint{5.880238in}{0.682122in}}%
\pgfpathlineto{\pgfqpoint{5.883594in}{0.744678in}}%
\pgfpathlineto{\pgfqpoint{5.886950in}{0.760702in}}%
\pgfpathlineto{\pgfqpoint{5.890306in}{0.755742in}}%
\pgfpathlineto{\pgfqpoint{5.893662in}{0.725617in}}%
\pgfpathlineto{\pgfqpoint{5.897017in}{0.669103in}}%
\pgfpathlineto{\pgfqpoint{5.900373in}{0.724087in}}%
\pgfpathlineto{\pgfqpoint{5.903729in}{0.725614in}}%
\pgfpathlineto{\pgfqpoint{5.907085in}{0.686877in}}%
\pgfpathlineto{\pgfqpoint{5.910441in}{0.713422in}}%
\pgfpathlineto{\pgfqpoint{5.913796in}{0.746303in}}%
\pgfpathlineto{\pgfqpoint{5.917152in}{0.750848in}}%
\pgfpathlineto{\pgfqpoint{5.920508in}{0.728921in}}%
\pgfpathlineto{\pgfqpoint{5.923864in}{0.633452in}}%
\pgfpathlineto{\pgfqpoint{5.927219in}{0.708736in}}%
\pgfpathlineto{\pgfqpoint{5.930575in}{0.718985in}}%
\pgfpathlineto{\pgfqpoint{5.933931in}{0.697828in}}%
\pgfpathlineto{\pgfqpoint{5.937287in}{0.680972in}}%
\pgfpathlineto{\pgfqpoint{5.940643in}{0.731429in}}%
\pgfpathlineto{\pgfqpoint{5.943998in}{0.743173in}}%
\pgfpathlineto{\pgfqpoint{5.947354in}{0.727730in}}%
\pgfpathlineto{\pgfqpoint{5.950710in}{0.633866in}}%
\pgfpathlineto{\pgfqpoint{5.954066in}{0.716097in}}%
\pgfpathlineto{\pgfqpoint{5.957421in}{0.732641in}}%
\pgfpathlineto{\pgfqpoint{5.960777in}{0.727969in}}%
\pgfpathlineto{\pgfqpoint{5.964133in}{0.689950in}}%
\pgfpathlineto{\pgfqpoint{5.967489in}{0.707690in}}%
\pgfpathlineto{\pgfqpoint{5.970845in}{0.742358in}}%
\pgfpathlineto{\pgfqpoint{5.974200in}{0.749293in}}%
\pgfpathlineto{\pgfqpoint{5.977556in}{0.728207in}}%
\pgfpathlineto{\pgfqpoint{5.980912in}{0.671629in}}%
\pgfpathlineto{\pgfqpoint{5.984268in}{0.737547in}}%
\pgfpathlineto{\pgfqpoint{5.987623in}{0.750112in}}%
\pgfpathlineto{\pgfqpoint{5.990979in}{0.745264in}}%
\pgfpathlineto{\pgfqpoint{5.994335in}{0.712984in}}%
\pgfpathlineto{\pgfqpoint{5.997691in}{0.714125in}}%
\pgfpathlineto{\pgfqpoint{6.001047in}{0.754461in}}%
\pgfpathlineto{\pgfqpoint{6.004402in}{0.764578in}}%
\pgfpathlineto{\pgfqpoint{6.007758in}{0.750845in}}%
\pgfpathlineto{\pgfqpoint{6.011114in}{0.679186in}}%
\pgfpathlineto{\pgfqpoint{6.014470in}{0.733460in}}%
\pgfpathlineto{\pgfqpoint{6.017825in}{0.752247in}}%
\pgfpathlineto{\pgfqpoint{6.021181in}{0.748776in}}%
\pgfpathlineto{\pgfqpoint{6.024537in}{0.714095in}}%
\pgfpathlineto{\pgfqpoint{6.027893in}{0.724718in}}%
\pgfpathlineto{\pgfqpoint{6.031249in}{0.761737in}}%
\pgfpathlineto{\pgfqpoint{6.034604in}{0.769678in}}%
\pgfpathlineto{\pgfqpoint{6.037960in}{0.753126in}}%
\pgfpathlineto{\pgfqpoint{6.041316in}{0.698566in}}%
\pgfpathlineto{\pgfqpoint{6.048027in}{0.733763in}}%
\pgfpathlineto{\pgfqpoint{6.051383in}{0.719276in}}%
\pgfpathlineto{\pgfqpoint{6.054739in}{0.675187in}}%
\pgfpathlineto{\pgfqpoint{6.058095in}{0.743264in}}%
\pgfpathlineto{\pgfqpoint{6.061451in}{0.758700in}}%
\pgfpathlineto{\pgfqpoint{6.064806in}{0.748368in}}%
\pgfpathlineto{\pgfqpoint{6.071518in}{0.678687in}}%
\pgfpathlineto{\pgfqpoint{6.074874in}{0.715017in}}%
\pgfpathlineto{\pgfqpoint{6.078230in}{0.706416in}}%
\pgfpathlineto{\pgfqpoint{6.081585in}{0.633505in}}%
\pgfpathlineto{\pgfqpoint{6.084941in}{0.725492in}}%
\pgfpathlineto{\pgfqpoint{6.088297in}{0.743711in}}%
\pgfpathlineto{\pgfqpoint{6.091653in}{0.735463in}}%
\pgfpathlineto{\pgfqpoint{6.095008in}{0.692190in}}%
\pgfpathlineto{\pgfqpoint{6.098364in}{0.694331in}}%
\pgfpathlineto{\pgfqpoint{6.101720in}{0.722136in}}%
\pgfpathlineto{\pgfqpoint{6.105076in}{0.720494in}}%
\pgfpathlineto{\pgfqpoint{6.108432in}{0.684523in}}%
\pgfpathlineto{\pgfqpoint{6.111787in}{0.701024in}}%
\pgfpathlineto{\pgfqpoint{6.115143in}{0.736363in}}%
\pgfpathlineto{\pgfqpoint{6.118499in}{0.742899in}}%
\pgfpathlineto{\pgfqpoint{6.121855in}{0.719911in}}%
\pgfpathlineto{\pgfqpoint{6.125210in}{0.673999in}}%
\pgfpathlineto{\pgfqpoint{6.128566in}{0.732260in}}%
\pgfpathlineto{\pgfqpoint{6.131922in}{0.743743in}}%
\pgfpathlineto{\pgfqpoint{6.135278in}{0.737871in}}%
\pgfpathlineto{\pgfqpoint{6.138634in}{0.702011in}}%
\pgfpathlineto{\pgfqpoint{6.141989in}{0.712353in}}%
\pgfpathlineto{\pgfqpoint{6.145345in}{0.749298in}}%
\pgfpathlineto{\pgfqpoint{6.148701in}{0.758074in}}%
\pgfpathlineto{\pgfqpoint{6.152057in}{0.741702in}}%
\pgfpathlineto{\pgfqpoint{6.155412in}{0.598022in}}%
\pgfpathlineto{\pgfqpoint{6.158768in}{0.740131in}}%
\pgfpathlineto{\pgfqpoint{6.162124in}{0.755820in}}%
\pgfpathlineto{\pgfqpoint{6.165480in}{0.752866in}}%
\pgfpathlineto{\pgfqpoint{6.168836in}{0.724958in}}%
\pgfpathlineto{\pgfqpoint{6.172191in}{0.713393in}}%
\pgfpathlineto{\pgfqpoint{6.175547in}{0.759342in}}%
\pgfpathlineto{\pgfqpoint{6.178903in}{0.771131in}}%
\pgfpathlineto{\pgfqpoint{6.182259in}{0.759755in}}%
\pgfpathlineto{\pgfqpoint{6.185614in}{0.713707in}}%
\pgfpathlineto{\pgfqpoint{6.188970in}{0.719233in}}%
\pgfpathlineto{\pgfqpoint{6.192326in}{0.744483in}}%
\pgfpathlineto{\pgfqpoint{6.195682in}{0.739866in}}%
\pgfpathlineto{\pgfqpoint{6.199038in}{0.686287in}}%
\pgfpathlineto{\pgfqpoint{6.202393in}{0.734726in}}%
\pgfpathlineto{\pgfqpoint{6.205749in}{0.762513in}}%
\pgfpathlineto{\pgfqpoint{6.209105in}{0.764135in}}%
\pgfpathlineto{\pgfqpoint{6.212461in}{0.741342in}}%
\pgfpathlineto{\pgfqpoint{6.215816in}{0.666875in}}%
\pgfpathlineto{\pgfqpoint{6.219172in}{0.709297in}}%
\pgfpathlineto{\pgfqpoint{6.222528in}{0.717197in}}%
\pgfpathlineto{\pgfqpoint{6.225884in}{0.673963in}}%
\pgfpathlineto{\pgfqpoint{6.229240in}{0.714298in}}%
\pgfpathlineto{\pgfqpoint{6.232595in}{0.744824in}}%
\pgfpathlineto{\pgfqpoint{6.235951in}{0.747336in}}%
\pgfpathlineto{\pgfqpoint{6.239307in}{0.723029in}}%
\pgfpathlineto{\pgfqpoint{6.242663in}{0.554900in}}%
\pgfpathlineto{\pgfqpoint{6.246019in}{0.707314in}}%
\pgfpathlineto{\pgfqpoint{6.249374in}{0.714563in}}%
\pgfpathlineto{\pgfqpoint{6.252730in}{0.687588in}}%
\pgfpathlineto{\pgfqpoint{6.256086in}{0.689292in}}%
\pgfpathlineto{\pgfqpoint{6.259442in}{0.730926in}}%
\pgfpathlineto{\pgfqpoint{6.262797in}{0.739669in}}%
\pgfpathlineto{\pgfqpoint{6.266153in}{0.719689in}}%
\pgfpathlineto{\pgfqpoint{6.269509in}{0.634234in}}%
\pgfpathlineto{\pgfqpoint{6.272865in}{0.721060in}}%
\pgfpathlineto{\pgfqpoint{6.276221in}{0.733971in}}%
\pgfpathlineto{\pgfqpoint{6.279576in}{0.727695in}}%
\pgfpathlineto{\pgfqpoint{6.282932in}{0.686461in}}%
\pgfpathlineto{\pgfqpoint{6.286288in}{0.709718in}}%
\pgfpathlineto{\pgfqpoint{6.289644in}{0.742744in}}%
\pgfpathlineto{\pgfqpoint{6.292999in}{0.749200in}}%
\pgfpathlineto{\pgfqpoint{6.296355in}{0.727214in}}%
\pgfpathlineto{\pgfqpoint{6.299711in}{0.684145in}}%
\pgfpathlineto{\pgfqpoint{6.303067in}{0.741897in}}%
\pgfpathlineto{\pgfqpoint{6.306423in}{0.754235in}}%
\pgfpathlineto{\pgfqpoint{6.309778in}{0.750157in}}%
\pgfpathlineto{\pgfqpoint{6.313134in}{0.721931in}}%
\pgfpathlineto{\pgfqpoint{6.316490in}{0.709224in}}%
\pgfpathlineto{\pgfqpoint{6.319846in}{0.755597in}}%
\pgfpathlineto{\pgfqpoint{6.323201in}{0.767921in}}%
\pgfpathlineto{\pgfqpoint{6.326557in}{0.757758in}}%
\pgfpathlineto{\pgfqpoint{6.329913in}{0.705298in}}%
\pgfpathlineto{\pgfqpoint{6.336625in}{0.754012in}}%
\pgfpathlineto{\pgfqpoint{6.339980in}{0.753067in}}%
\pgfpathlineto{\pgfqpoint{6.343336in}{0.724961in}}%
\pgfpathlineto{\pgfqpoint{6.346692in}{0.718268in}}%
\pgfpathlineto{\pgfqpoint{6.350048in}{0.762260in}}%
\pgfpathlineto{\pgfqpoint{6.353403in}{0.772804in}}%
\pgfpathlineto{\pgfqpoint{6.356759in}{0.759314in}}%
\pgfpathlineto{\pgfqpoint{6.360115in}{0.717474in}}%
\pgfpathlineto{\pgfqpoint{6.363471in}{0.701531in}}%
\pgfpathlineto{\pgfqpoint{6.366827in}{0.729069in}}%
\pgfpathlineto{\pgfqpoint{6.370182in}{0.716030in}}%
\pgfpathlineto{\pgfqpoint{6.373538in}{0.676967in}}%
\pgfpathlineto{\pgfqpoint{6.376894in}{0.742609in}}%
\pgfpathlineto{\pgfqpoint{6.380250in}{0.757724in}}%
\pgfpathlineto{\pgfqpoint{6.383605in}{0.748501in}}%
\pgfpathlineto{\pgfqpoint{6.383605in}{0.748501in}}%
\pgfusepath{stroke}%
\end{pgfscope}%
\begin{pgfscope}%
\pgfpathrectangle{\pgfqpoint{4.621821in}{0.320625in}}{\pgfqpoint{1.845679in}{1.164375in}}%
\pgfusepath{clip}%
\pgfsetrectcap%
\pgfsetroundjoin%
\pgfsetlinewidth{1.154312pt}%
\definecolor{currentstroke}{rgb}{0.141176,0.345098,0.650980}%
\pgfsetstrokecolor{currentstroke}%
\pgfsetdash{}{0pt}%
\pgfpathmoveto{\pgfqpoint{4.705715in}{1.476601in}}%
\pgfpathlineto{\pgfqpoint{4.712427in}{1.464541in}}%
\pgfpathlineto{\pgfqpoint{4.722494in}{1.449699in}}%
\pgfpathlineto{\pgfqpoint{4.739273in}{1.427338in}}%
\pgfpathlineto{\pgfqpoint{4.759408in}{1.402349in}}%
\pgfpathlineto{\pgfqpoint{4.776187in}{1.383316in}}%
\pgfpathlineto{\pgfqpoint{4.792966in}{1.366051in}}%
\pgfpathlineto{\pgfqpoint{4.809745in}{1.350323in}}%
\pgfpathlineto{\pgfqpoint{4.829879in}{1.332966in}}%
\pgfpathlineto{\pgfqpoint{4.856726in}{1.311514in}}%
\pgfpathlineto{\pgfqpoint{4.890283in}{1.286286in}}%
\pgfpathlineto{\pgfqpoint{4.933909in}{1.254938in}}%
\pgfpathlineto{\pgfqpoint{4.990957in}{1.215389in}}%
\pgfpathlineto{\pgfqpoint{5.054717in}{1.172556in}}%
\pgfpathlineto{\pgfqpoint{5.121832in}{1.128781in}}%
\pgfpathlineto{\pgfqpoint{5.192304in}{1.084148in}}%
\pgfpathlineto{\pgfqpoint{5.266131in}{1.038756in}}%
\pgfpathlineto{\pgfqpoint{5.339958in}{0.994661in}}%
\pgfpathlineto{\pgfqpoint{5.420497in}{0.947892in}}%
\pgfpathlineto{\pgfqpoint{5.507747in}{0.898588in}}%
\pgfpathlineto{\pgfqpoint{5.594997in}{0.850483in}}%
\pgfpathlineto{\pgfqpoint{5.638622in}{0.826785in}}%
\pgfpathlineto{\pgfqpoint{5.702382in}{0.793017in}}%
\pgfpathlineto{\pgfqpoint{5.712449in}{0.787962in}}%
\pgfpathlineto{\pgfqpoint{5.719161in}{0.783819in}}%
\pgfpathlineto{\pgfqpoint{5.746007in}{0.772593in}}%
\pgfpathlineto{\pgfqpoint{5.749363in}{0.769385in}}%
\pgfpathlineto{\pgfqpoint{5.756075in}{0.768985in}}%
\pgfpathlineto{\pgfqpoint{5.762786in}{0.763691in}}%
\pgfpathlineto{\pgfqpoint{5.766142in}{0.762592in}}%
\pgfpathlineto{\pgfqpoint{5.769498in}{0.762522in}}%
\pgfpathlineto{\pgfqpoint{5.772854in}{0.763689in}}%
\pgfpathlineto{\pgfqpoint{5.776209in}{0.761622in}}%
\pgfpathlineto{\pgfqpoint{5.786277in}{0.749548in}}%
\pgfpathlineto{\pgfqpoint{5.789632in}{0.748532in}}%
\pgfpathlineto{\pgfqpoint{5.792988in}{0.749639in}}%
\pgfpathlineto{\pgfqpoint{5.796344in}{0.750042in}}%
\pgfpathlineto{\pgfqpoint{5.799700in}{0.749655in}}%
\pgfpathlineto{\pgfqpoint{5.806411in}{0.745668in}}%
\pgfpathlineto{\pgfqpoint{5.809767in}{0.741952in}}%
\pgfpathlineto{\pgfqpoint{5.813123in}{0.737175in}}%
\pgfpathlineto{\pgfqpoint{5.816479in}{0.734298in}}%
\pgfpathlineto{\pgfqpoint{5.819834in}{0.736290in}}%
\pgfpathlineto{\pgfqpoint{5.823190in}{0.741503in}}%
\pgfpathlineto{\pgfqpoint{5.826546in}{0.742591in}}%
\pgfpathlineto{\pgfqpoint{5.833258in}{0.731867in}}%
\pgfpathlineto{\pgfqpoint{5.836613in}{0.733458in}}%
\pgfpathlineto{\pgfqpoint{5.839969in}{0.738124in}}%
\pgfpathlineto{\pgfqpoint{5.843325in}{0.736505in}}%
\pgfpathlineto{\pgfqpoint{5.846681in}{0.726237in}}%
\pgfpathlineto{\pgfqpoint{5.850036in}{0.725060in}}%
\pgfpathlineto{\pgfqpoint{5.853392in}{0.742164in}}%
\pgfpathlineto{\pgfqpoint{5.856748in}{0.749646in}}%
\pgfpathlineto{\pgfqpoint{5.860104in}{0.739768in}}%
\pgfpathlineto{\pgfqpoint{5.863460in}{0.717447in}}%
\pgfpathlineto{\pgfqpoint{5.866815in}{0.725340in}}%
\pgfpathlineto{\pgfqpoint{5.870171in}{0.736486in}}%
\pgfpathlineto{\pgfqpoint{5.873527in}{0.732979in}}%
\pgfpathlineto{\pgfqpoint{5.876883in}{0.716471in}}%
\pgfpathlineto{\pgfqpoint{5.880238in}{0.725569in}}%
\pgfpathlineto{\pgfqpoint{5.883594in}{0.744643in}}%
\pgfpathlineto{\pgfqpoint{5.886950in}{0.749355in}}%
\pgfpathlineto{\pgfqpoint{5.890306in}{0.736286in}}%
\pgfpathlineto{\pgfqpoint{5.893662in}{0.709976in}}%
\pgfpathlineto{\pgfqpoint{5.900373in}{0.728471in}}%
\pgfpathlineto{\pgfqpoint{5.903729in}{0.722660in}}%
\pgfpathlineto{\pgfqpoint{5.907085in}{0.708840in}}%
\pgfpathlineto{\pgfqpoint{5.910441in}{0.728324in}}%
\pgfpathlineto{\pgfqpoint{5.913796in}{0.742760in}}%
\pgfpathlineto{\pgfqpoint{5.917152in}{0.740168in}}%
\pgfpathlineto{\pgfqpoint{5.920508in}{0.724591in}}%
\pgfpathlineto{\pgfqpoint{5.923864in}{0.705883in}}%
\pgfpathlineto{\pgfqpoint{5.930575in}{0.719547in}}%
\pgfpathlineto{\pgfqpoint{5.933931in}{0.717959in}}%
\pgfpathlineto{\pgfqpoint{5.937287in}{0.722191in}}%
\pgfpathlineto{\pgfqpoint{5.940643in}{0.730530in}}%
\pgfpathlineto{\pgfqpoint{5.943998in}{0.734310in}}%
\pgfpathlineto{\pgfqpoint{5.947354in}{0.732278in}}%
\pgfpathlineto{\pgfqpoint{5.950710in}{0.723434in}}%
\pgfpathlineto{\pgfqpoint{5.954066in}{0.718246in}}%
\pgfpathlineto{\pgfqpoint{5.957421in}{0.728056in}}%
\pgfpathlineto{\pgfqpoint{5.960777in}{0.731780in}}%
\pgfpathlineto{\pgfqpoint{5.967489in}{0.711443in}}%
\pgfpathlineto{\pgfqpoint{5.970845in}{0.728479in}}%
\pgfpathlineto{\pgfqpoint{5.974200in}{0.734817in}}%
\pgfpathlineto{\pgfqpoint{5.980912in}{0.707066in}}%
\pgfpathlineto{\pgfqpoint{5.984268in}{0.728970in}}%
\pgfpathlineto{\pgfqpoint{5.987623in}{0.736982in}}%
\pgfpathlineto{\pgfqpoint{5.990979in}{0.733268in}}%
\pgfpathlineto{\pgfqpoint{5.994335in}{0.705065in}}%
\pgfpathlineto{\pgfqpoint{5.997691in}{0.703827in}}%
\pgfpathlineto{\pgfqpoint{6.001047in}{0.739639in}}%
\pgfpathlineto{\pgfqpoint{6.004402in}{0.748920in}}%
\pgfpathlineto{\pgfqpoint{6.007758in}{0.735920in}}%
\pgfpathlineto{\pgfqpoint{6.011114in}{0.688035in}}%
\pgfpathlineto{\pgfqpoint{6.014470in}{0.723533in}}%
\pgfpathlineto{\pgfqpoint{6.017825in}{0.740900in}}%
\pgfpathlineto{\pgfqpoint{6.021181in}{0.735046in}}%
\pgfpathlineto{\pgfqpoint{6.024537in}{0.700618in}}%
\pgfpathlineto{\pgfqpoint{6.027893in}{0.714393in}}%
\pgfpathlineto{\pgfqpoint{6.031249in}{0.745920in}}%
\pgfpathlineto{\pgfqpoint{6.034604in}{0.753018in}}%
\pgfpathlineto{\pgfqpoint{6.037960in}{0.736990in}}%
\pgfpathlineto{\pgfqpoint{6.041316in}{0.715120in}}%
\pgfpathlineto{\pgfqpoint{6.044672in}{0.731899in}}%
\pgfpathlineto{\pgfqpoint{6.048027in}{0.733741in}}%
\pgfpathlineto{\pgfqpoint{6.051383in}{0.721143in}}%
\pgfpathlineto{\pgfqpoint{6.054739in}{0.703773in}}%
\pgfpathlineto{\pgfqpoint{6.058095in}{0.729401in}}%
\pgfpathlineto{\pgfqpoint{6.061451in}{0.744820in}}%
\pgfpathlineto{\pgfqpoint{6.064806in}{0.741011in}}%
\pgfpathlineto{\pgfqpoint{6.068162in}{0.731311in}}%
\pgfpathlineto{\pgfqpoint{6.071518in}{0.726186in}}%
\pgfpathlineto{\pgfqpoint{6.074874in}{0.712761in}}%
\pgfpathlineto{\pgfqpoint{6.078230in}{0.694828in}}%
\pgfpathlineto{\pgfqpoint{6.081585in}{0.693329in}}%
\pgfpathlineto{\pgfqpoint{6.084941in}{0.717529in}}%
\pgfpathlineto{\pgfqpoint{6.088297in}{0.728694in}}%
\pgfpathlineto{\pgfqpoint{6.095008in}{0.720803in}}%
\pgfpathlineto{\pgfqpoint{6.098364in}{0.715780in}}%
\pgfpathlineto{\pgfqpoint{6.101720in}{0.707352in}}%
\pgfpathlineto{\pgfqpoint{6.108432in}{0.705787in}}%
\pgfpathlineto{\pgfqpoint{6.111787in}{0.704386in}}%
\pgfpathlineto{\pgfqpoint{6.115143in}{0.713460in}}%
\pgfpathlineto{\pgfqpoint{6.118499in}{0.726888in}}%
\pgfpathlineto{\pgfqpoint{6.121855in}{0.726756in}}%
\pgfpathlineto{\pgfqpoint{6.125210in}{0.710915in}}%
\pgfpathlineto{\pgfqpoint{6.128566in}{0.709301in}}%
\pgfpathlineto{\pgfqpoint{6.131922in}{0.725623in}}%
\pgfpathlineto{\pgfqpoint{6.135278in}{0.730542in}}%
\pgfpathlineto{\pgfqpoint{6.138634in}{0.721401in}}%
\pgfpathlineto{\pgfqpoint{6.141989in}{0.699886in}}%
\pgfpathlineto{\pgfqpoint{6.145345in}{0.728134in}}%
\pgfpathlineto{\pgfqpoint{6.148701in}{0.743889in}}%
\pgfpathlineto{\pgfqpoint{6.152057in}{0.740466in}}%
\pgfpathlineto{\pgfqpoint{6.155412in}{0.721670in}}%
\pgfpathlineto{\pgfqpoint{6.158768in}{0.721099in}}%
\pgfpathlineto{\pgfqpoint{6.162124in}{0.738928in}}%
\pgfpathlineto{\pgfqpoint{6.165480in}{0.744015in}}%
\pgfpathlineto{\pgfqpoint{6.172191in}{0.724713in}}%
\pgfpathlineto{\pgfqpoint{6.175547in}{0.743485in}}%
\pgfpathlineto{\pgfqpoint{6.178903in}{0.752328in}}%
\pgfpathlineto{\pgfqpoint{6.182259in}{0.748749in}}%
\pgfpathlineto{\pgfqpoint{6.185614in}{0.737650in}}%
\pgfpathlineto{\pgfqpoint{6.188970in}{0.736421in}}%
\pgfpathlineto{\pgfqpoint{6.192326in}{0.741697in}}%
\pgfpathlineto{\pgfqpoint{6.199038in}{0.723904in}}%
\pgfpathlineto{\pgfqpoint{6.202393in}{0.735360in}}%
\pgfpathlineto{\pgfqpoint{6.205749in}{0.742324in}}%
\pgfpathlineto{\pgfqpoint{6.209105in}{0.743068in}}%
\pgfpathlineto{\pgfqpoint{6.215816in}{0.733764in}}%
\pgfpathlineto{\pgfqpoint{6.219172in}{0.731710in}}%
\pgfpathlineto{\pgfqpoint{6.225884in}{0.720516in}}%
\pgfpathlineto{\pgfqpoint{6.229240in}{0.718478in}}%
\pgfpathlineto{\pgfqpoint{6.232595in}{0.719565in}}%
\pgfpathlineto{\pgfqpoint{6.235951in}{0.725759in}}%
\pgfpathlineto{\pgfqpoint{6.239307in}{0.725212in}}%
\pgfpathlineto{\pgfqpoint{6.242663in}{0.712917in}}%
\pgfpathlineto{\pgfqpoint{6.246019in}{0.705110in}}%
\pgfpathlineto{\pgfqpoint{6.249374in}{0.718660in}}%
\pgfpathlineto{\pgfqpoint{6.252730in}{0.720620in}}%
\pgfpathlineto{\pgfqpoint{6.256086in}{0.706922in}}%
\pgfpathlineto{\pgfqpoint{6.259442in}{0.701627in}}%
\pgfpathlineto{\pgfqpoint{6.262797in}{0.718117in}}%
\pgfpathlineto{\pgfqpoint{6.266153in}{0.720352in}}%
\pgfpathlineto{\pgfqpoint{6.269509in}{0.709998in}}%
\pgfpathlineto{\pgfqpoint{6.272865in}{0.703397in}}%
\pgfpathlineto{\pgfqpoint{6.276221in}{0.721835in}}%
\pgfpathlineto{\pgfqpoint{6.279576in}{0.728341in}}%
\pgfpathlineto{\pgfqpoint{6.282932in}{0.716858in}}%
\pgfpathlineto{\pgfqpoint{6.286288in}{0.702964in}}%
\pgfpathlineto{\pgfqpoint{6.289644in}{0.724360in}}%
\pgfpathlineto{\pgfqpoint{6.292999in}{0.734502in}}%
\pgfpathlineto{\pgfqpoint{6.299711in}{0.718740in}}%
\pgfpathlineto{\pgfqpoint{6.303067in}{0.732131in}}%
\pgfpathlineto{\pgfqpoint{6.306423in}{0.742384in}}%
\pgfpathlineto{\pgfqpoint{6.309778in}{0.738826in}}%
\pgfpathlineto{\pgfqpoint{6.313134in}{0.725347in}}%
\pgfpathlineto{\pgfqpoint{6.316490in}{0.731167in}}%
\pgfpathlineto{\pgfqpoint{6.319846in}{0.746156in}}%
\pgfpathlineto{\pgfqpoint{6.323201in}{0.748570in}}%
\pgfpathlineto{\pgfqpoint{6.326557in}{0.737809in}}%
\pgfpathlineto{\pgfqpoint{6.329913in}{0.738700in}}%
\pgfpathlineto{\pgfqpoint{6.333269in}{0.746722in}}%
\pgfpathlineto{\pgfqpoint{6.336625in}{0.741705in}}%
\pgfpathlineto{\pgfqpoint{6.339980in}{0.730125in}}%
\pgfpathlineto{\pgfqpoint{6.343336in}{0.730441in}}%
\pgfpathlineto{\pgfqpoint{6.346692in}{0.740233in}}%
\pgfpathlineto{\pgfqpoint{6.353403in}{0.745942in}}%
\pgfpathlineto{\pgfqpoint{6.356759in}{0.743315in}}%
\pgfpathlineto{\pgfqpoint{6.360115in}{0.744756in}}%
\pgfpathlineto{\pgfqpoint{6.366827in}{0.730875in}}%
\pgfpathlineto{\pgfqpoint{6.370182in}{0.732504in}}%
\pgfpathlineto{\pgfqpoint{6.373538in}{0.726870in}}%
\pgfpathlineto{\pgfqpoint{6.376894in}{0.726035in}}%
\pgfpathlineto{\pgfqpoint{6.380250in}{0.732956in}}%
\pgfpathlineto{\pgfqpoint{6.383605in}{0.724494in}}%
\pgfpathlineto{\pgfqpoint{6.383605in}{0.724494in}}%
\pgfusepath{stroke}%
\end{pgfscope}%
\begin{pgfscope}%
\pgfpathrectangle{\pgfqpoint{4.621821in}{0.320625in}}{\pgfqpoint{1.845679in}{1.164375in}}%
\pgfusepath{clip}%
\pgfsetrectcap%
\pgfsetroundjoin%
\pgfsetlinewidth{0.602250pt}%
\definecolor{currentstroke}{rgb}{0.819608,0.286275,0.356863}%
\pgfsetstrokecolor{currentstroke}%
\pgfsetstrokeopacity{0.120000}%
\pgfsetdash{}{0pt}%
\pgfpathmoveto{\pgfqpoint{4.705715in}{1.471660in}}%
\pgfpathlineto{\pgfqpoint{4.729206in}{1.426942in}}%
\pgfpathlineto{\pgfqpoint{4.766120in}{1.351819in}}%
\pgfpathlineto{\pgfqpoint{4.782898in}{1.321199in}}%
\pgfpathlineto{\pgfqpoint{4.799677in}{1.292882in}}%
\pgfpathlineto{\pgfqpoint{4.819812in}{1.261068in}}%
\pgfpathlineto{\pgfqpoint{4.846658in}{1.221020in}}%
\pgfpathlineto{\pgfqpoint{4.876860in}{1.177984in}}%
\pgfpathlineto{\pgfqpoint{4.913774in}{1.127288in}}%
\pgfpathlineto{\pgfqpoint{4.954043in}{1.073792in}}%
\pgfpathlineto{\pgfqpoint{4.997668in}{1.017664in}}%
\pgfpathlineto{\pgfqpoint{5.044649in}{0.959093in}}%
\pgfpathlineto{\pgfqpoint{5.094986in}{0.898148in}}%
\pgfpathlineto{\pgfqpoint{5.145323in}{0.838542in}}%
\pgfpathlineto{\pgfqpoint{5.219150in}{0.752874in}}%
\pgfpathlineto{\pgfqpoint{5.235929in}{0.732196in}}%
\pgfpathlineto{\pgfqpoint{5.245996in}{0.722004in}}%
\pgfpathlineto{\pgfqpoint{5.266131in}{0.704358in}}%
\pgfpathlineto{\pgfqpoint{5.272842in}{0.695393in}}%
\pgfpathlineto{\pgfqpoint{5.282910in}{0.679176in}}%
\pgfpathlineto{\pgfqpoint{5.296333in}{0.655961in}}%
\pgfpathlineto{\pgfqpoint{5.303044in}{0.646529in}}%
\pgfpathlineto{\pgfqpoint{5.306400in}{0.642989in}}%
\pgfpathlineto{\pgfqpoint{5.309756in}{0.640592in}}%
\pgfpathlineto{\pgfqpoint{5.313112in}{0.639697in}}%
\pgfpathlineto{\pgfqpoint{5.316467in}{0.640500in}}%
\pgfpathlineto{\pgfqpoint{5.319823in}{0.642785in}}%
\pgfpathlineto{\pgfqpoint{5.329891in}{0.652200in}}%
\pgfpathlineto{\pgfqpoint{5.333246in}{0.654199in}}%
\pgfpathlineto{\pgfqpoint{5.336602in}{0.654857in}}%
\pgfpathlineto{\pgfqpoint{5.339958in}{0.653642in}}%
\pgfpathlineto{\pgfqpoint{5.343314in}{0.649668in}}%
\pgfpathlineto{\pgfqpoint{5.346669in}{0.651547in}}%
\pgfpathlineto{\pgfqpoint{5.353381in}{0.691051in}}%
\pgfpathlineto{\pgfqpoint{5.356737in}{0.703662in}}%
\pgfpathlineto{\pgfqpoint{5.360093in}{0.711768in}}%
\pgfpathlineto{\pgfqpoint{5.363448in}{0.716379in}}%
\pgfpathlineto{\pgfqpoint{5.366804in}{0.717947in}}%
\pgfpathlineto{\pgfqpoint{5.370160in}{0.716475in}}%
\pgfpathlineto{\pgfqpoint{5.373516in}{0.711434in}}%
\pgfpathlineto{\pgfqpoint{5.376871in}{0.701250in}}%
\pgfpathlineto{\pgfqpoint{5.380227in}{0.680858in}}%
\pgfpathlineto{\pgfqpoint{5.383583in}{0.614338in}}%
\pgfpathlineto{\pgfqpoint{5.386939in}{0.674250in}}%
\pgfpathlineto{\pgfqpoint{5.390295in}{0.702734in}}%
\pgfpathlineto{\pgfqpoint{5.393650in}{0.717543in}}%
\pgfpathlineto{\pgfqpoint{5.397006in}{0.725953in}}%
\pgfpathlineto{\pgfqpoint{5.400362in}{0.729614in}}%
\pgfpathlineto{\pgfqpoint{5.403718in}{0.728250in}}%
\pgfpathlineto{\pgfqpoint{5.407074in}{0.719173in}}%
\pgfpathlineto{\pgfqpoint{5.410429in}{0.696281in}}%
\pgfpathlineto{\pgfqpoint{5.413785in}{0.642767in}}%
\pgfpathlineto{\pgfqpoint{5.417141in}{0.679860in}}%
\pgfpathlineto{\pgfqpoint{5.420497in}{0.694773in}}%
\pgfpathlineto{\pgfqpoint{5.423852in}{0.691082in}}%
\pgfpathlineto{\pgfqpoint{5.427208in}{0.661663in}}%
\pgfpathlineto{\pgfqpoint{5.430564in}{0.657645in}}%
\pgfpathlineto{\pgfqpoint{5.433920in}{0.701604in}}%
\pgfpathlineto{\pgfqpoint{5.437276in}{0.716304in}}%
\pgfpathlineto{\pgfqpoint{5.440631in}{0.714428in}}%
\pgfpathlineto{\pgfqpoint{5.443987in}{0.694457in}}%
\pgfpathlineto{\pgfqpoint{5.447343in}{0.640327in}}%
\pgfpathlineto{\pgfqpoint{5.450699in}{0.669132in}}%
\pgfpathlineto{\pgfqpoint{5.454054in}{0.689096in}}%
\pgfpathlineto{\pgfqpoint{5.457410in}{0.688814in}}%
\pgfpathlineto{\pgfqpoint{5.460766in}{0.667840in}}%
\pgfpathlineto{\pgfqpoint{5.464122in}{0.624721in}}%
\pgfpathlineto{\pgfqpoint{5.467478in}{0.690210in}}%
\pgfpathlineto{\pgfqpoint{5.470833in}{0.709008in}}%
\pgfpathlineto{\pgfqpoint{5.474189in}{0.711075in}}%
\pgfpathlineto{\pgfqpoint{5.477545in}{0.688449in}}%
\pgfpathlineto{\pgfqpoint{5.480901in}{0.657259in}}%
\pgfpathlineto{\pgfqpoint{5.484256in}{0.711172in}}%
\pgfpathlineto{\pgfqpoint{5.487612in}{0.728311in}}%
\pgfpathlineto{\pgfqpoint{5.490968in}{0.733750in}}%
\pgfpathlineto{\pgfqpoint{5.494324in}{0.730440in}}%
\pgfpathlineto{\pgfqpoint{5.497680in}{0.715469in}}%
\pgfpathlineto{\pgfqpoint{5.501035in}{0.659463in}}%
\pgfpathlineto{\pgfqpoint{5.504391in}{0.702226in}}%
\pgfpathlineto{\pgfqpoint{5.507747in}{0.732617in}}%
\pgfpathlineto{\pgfqpoint{5.511103in}{0.744103in}}%
\pgfpathlineto{\pgfqpoint{5.514458in}{0.742827in}}%
\pgfpathlineto{\pgfqpoint{5.517814in}{0.725441in}}%
\pgfpathlineto{\pgfqpoint{5.521170in}{0.680751in}}%
\pgfpathlineto{\pgfqpoint{5.524526in}{0.683820in}}%
\pgfpathlineto{\pgfqpoint{5.527882in}{0.710337in}}%
\pgfpathlineto{\pgfqpoint{5.531237in}{0.711526in}}%
\pgfpathlineto{\pgfqpoint{5.534593in}{0.690442in}}%
\pgfpathlineto{\pgfqpoint{5.537949in}{0.653502in}}%
\pgfpathlineto{\pgfqpoint{5.541305in}{0.715508in}}%
\pgfpathlineto{\pgfqpoint{5.544660in}{0.733618in}}%
\pgfpathlineto{\pgfqpoint{5.548016in}{0.734414in}}%
\pgfpathlineto{\pgfqpoint{5.551372in}{0.716398in}}%
\pgfpathlineto{\pgfqpoint{5.554728in}{0.668015in}}%
\pgfpathlineto{\pgfqpoint{5.558084in}{0.679159in}}%
\pgfpathlineto{\pgfqpoint{5.561439in}{0.703911in}}%
\pgfpathlineto{\pgfqpoint{5.564795in}{0.704477in}}%
\pgfpathlineto{\pgfqpoint{5.568151in}{0.682606in}}%
\pgfpathlineto{\pgfqpoint{5.571507in}{0.649082in}}%
\pgfpathlineto{\pgfqpoint{5.574863in}{0.708878in}}%
\pgfpathlineto{\pgfqpoint{5.578218in}{0.726552in}}%
\pgfpathlineto{\pgfqpoint{5.581574in}{0.726863in}}%
\pgfpathlineto{\pgfqpoint{5.584930in}{0.705273in}}%
\pgfpathlineto{\pgfqpoint{5.588286in}{0.618063in}}%
\pgfpathlineto{\pgfqpoint{5.591641in}{0.703581in}}%
\pgfpathlineto{\pgfqpoint{5.594997in}{0.721055in}}%
\pgfpathlineto{\pgfqpoint{5.598353in}{0.723058in}}%
\pgfpathlineto{\pgfqpoint{5.601709in}{0.711421in}}%
\pgfpathlineto{\pgfqpoint{5.605065in}{0.662163in}}%
\pgfpathlineto{\pgfqpoint{5.611776in}{0.728222in}}%
\pgfpathlineto{\pgfqpoint{5.615132in}{0.738747in}}%
\pgfpathlineto{\pgfqpoint{5.618488in}{0.732822in}}%
\pgfpathlineto{\pgfqpoint{5.621843in}{0.699161in}}%
\pgfpathlineto{\pgfqpoint{5.625199in}{0.688656in}}%
\pgfpathlineto{\pgfqpoint{5.628555in}{0.728868in}}%
\pgfpathlineto{\pgfqpoint{5.631911in}{0.740510in}}%
\pgfpathlineto{\pgfqpoint{5.635267in}{0.739930in}}%
\pgfpathlineto{\pgfqpoint{5.638622in}{0.726079in}}%
\pgfpathlineto{\pgfqpoint{5.641978in}{0.666828in}}%
\pgfpathlineto{\pgfqpoint{5.645334in}{0.716331in}}%
\pgfpathlineto{\pgfqpoint{5.648690in}{0.745444in}}%
\pgfpathlineto{\pgfqpoint{5.652045in}{0.755154in}}%
\pgfpathlineto{\pgfqpoint{5.655401in}{0.749083in}}%
\pgfpathlineto{\pgfqpoint{5.658757in}{0.721413in}}%
\pgfpathlineto{\pgfqpoint{5.662113in}{0.666609in}}%
\pgfpathlineto{\pgfqpoint{5.665469in}{0.727797in}}%
\pgfpathlineto{\pgfqpoint{5.668824in}{0.740386in}}%
\pgfpathlineto{\pgfqpoint{5.672180in}{0.737480in}}%
\pgfpathlineto{\pgfqpoint{5.675536in}{0.714791in}}%
\pgfpathlineto{\pgfqpoint{5.678892in}{0.673448in}}%
\pgfpathlineto{\pgfqpoint{5.682247in}{0.737169in}}%
\pgfpathlineto{\pgfqpoint{5.685603in}{0.755403in}}%
\pgfpathlineto{\pgfqpoint{5.688959in}{0.757070in}}%
\pgfpathlineto{\pgfqpoint{5.692315in}{0.740935in}}%
\pgfpathlineto{\pgfqpoint{5.695671in}{0.702618in}}%
\pgfpathlineto{\pgfqpoint{5.699026in}{0.677159in}}%
\pgfpathlineto{\pgfqpoint{5.702382in}{0.713703in}}%
\pgfpathlineto{\pgfqpoint{5.705738in}{0.714601in}}%
\pgfpathlineto{\pgfqpoint{5.709094in}{0.685521in}}%
\pgfpathlineto{\pgfqpoint{5.712449in}{0.687029in}}%
\pgfpathlineto{\pgfqpoint{5.715805in}{0.729973in}}%
\pgfpathlineto{\pgfqpoint{5.719161in}{0.743743in}}%
\pgfpathlineto{\pgfqpoint{5.722517in}{0.739047in}}%
\pgfpathlineto{\pgfqpoint{5.725873in}{0.717103in}}%
\pgfpathlineto{\pgfqpoint{5.729228in}{0.656066in}}%
\pgfpathlineto{\pgfqpoint{5.732584in}{0.683157in}}%
\pgfpathlineto{\pgfqpoint{5.735940in}{0.701286in}}%
\pgfpathlineto{\pgfqpoint{5.739296in}{0.693021in}}%
\pgfpathlineto{\pgfqpoint{5.742652in}{0.615693in}}%
\pgfpathlineto{\pgfqpoint{5.746007in}{0.696754in}}%
\pgfpathlineto{\pgfqpoint{5.749363in}{0.723297in}}%
\pgfpathlineto{\pgfqpoint{5.752719in}{0.729851in}}%
\pgfpathlineto{\pgfqpoint{5.756075in}{0.716342in}}%
\pgfpathlineto{\pgfqpoint{5.759430in}{0.671702in}}%
\pgfpathlineto{\pgfqpoint{5.762786in}{0.678498in}}%
\pgfpathlineto{\pgfqpoint{5.766142in}{0.708704in}}%
\pgfpathlineto{\pgfqpoint{5.769498in}{0.714286in}}%
\pgfpathlineto{\pgfqpoint{5.772854in}{0.704444in}}%
\pgfpathlineto{\pgfqpoint{5.776209in}{0.657421in}}%
\pgfpathlineto{\pgfqpoint{5.782921in}{0.721918in}}%
\pgfpathlineto{\pgfqpoint{5.786277in}{0.732167in}}%
\pgfpathlineto{\pgfqpoint{5.789632in}{0.724599in}}%
\pgfpathlineto{\pgfqpoint{5.792988in}{0.684789in}}%
\pgfpathlineto{\pgfqpoint{5.796344in}{0.691099in}}%
\pgfpathlineto{\pgfqpoint{5.799700in}{0.725516in}}%
\pgfpathlineto{\pgfqpoint{5.803056in}{0.735705in}}%
\pgfpathlineto{\pgfqpoint{5.806411in}{0.734127in}}%
\pgfpathlineto{\pgfqpoint{5.809767in}{0.718786in}}%
\pgfpathlineto{\pgfqpoint{5.813123in}{0.645488in}}%
\pgfpathlineto{\pgfqpoint{5.816479in}{0.714418in}}%
\pgfpathlineto{\pgfqpoint{5.819834in}{0.741268in}}%
\pgfpathlineto{\pgfqpoint{5.823190in}{0.749886in}}%
\pgfpathlineto{\pgfqpoint{5.826546in}{0.742302in}}%
\pgfpathlineto{\pgfqpoint{5.829902in}{0.707135in}}%
\pgfpathlineto{\pgfqpoint{5.833258in}{0.696832in}}%
\pgfpathlineto{\pgfqpoint{5.836613in}{0.736189in}}%
\pgfpathlineto{\pgfqpoint{5.839969in}{0.746755in}}%
\pgfpathlineto{\pgfqpoint{5.843325in}{0.744471in}}%
\pgfpathlineto{\pgfqpoint{5.846681in}{0.726437in}}%
\pgfpathlineto{\pgfqpoint{5.850036in}{0.571158in}}%
\pgfpathlineto{\pgfqpoint{5.853392in}{0.733199in}}%
\pgfpathlineto{\pgfqpoint{5.856748in}{0.756044in}}%
\pgfpathlineto{\pgfqpoint{5.860104in}{0.761959in}}%
\pgfpathlineto{\pgfqpoint{5.863460in}{0.750731in}}%
\pgfpathlineto{\pgfqpoint{5.866815in}{0.719042in}}%
\pgfpathlineto{\pgfqpoint{5.870171in}{0.668352in}}%
\pgfpathlineto{\pgfqpoint{5.873527in}{0.722793in}}%
\pgfpathlineto{\pgfqpoint{5.876883in}{0.730450in}}%
\pgfpathlineto{\pgfqpoint{5.880238in}{0.717933in}}%
\pgfpathlineto{\pgfqpoint{5.883594in}{0.613424in}}%
\pgfpathlineto{\pgfqpoint{5.886950in}{0.723290in}}%
\pgfpathlineto{\pgfqpoint{5.890306in}{0.747901in}}%
\pgfpathlineto{\pgfqpoint{5.893662in}{0.753470in}}%
\pgfpathlineto{\pgfqpoint{5.897017in}{0.740587in}}%
\pgfpathlineto{\pgfqpoint{5.900373in}{0.711143in}}%
\pgfpathlineto{\pgfqpoint{5.903729in}{0.585186in}}%
\pgfpathlineto{\pgfqpoint{5.907085in}{0.695098in}}%
\pgfpathlineto{\pgfqpoint{5.910441in}{0.700281in}}%
\pgfpathlineto{\pgfqpoint{5.913796in}{0.670287in}}%
\pgfpathlineto{\pgfqpoint{5.917152in}{0.680424in}}%
\pgfpathlineto{\pgfqpoint{5.920508in}{0.721130in}}%
\pgfpathlineto{\pgfqpoint{5.923864in}{0.734278in}}%
\pgfpathlineto{\pgfqpoint{5.927219in}{0.728430in}}%
\pgfpathlineto{\pgfqpoint{5.930575in}{0.703275in}}%
\pgfpathlineto{\pgfqpoint{5.933931in}{0.555311in}}%
\pgfpathlineto{\pgfqpoint{5.937287in}{0.690698in}}%
\pgfpathlineto{\pgfqpoint{5.940643in}{0.704117in}}%
\pgfpathlineto{\pgfqpoint{5.943998in}{0.697210in}}%
\pgfpathlineto{\pgfqpoint{5.947354in}{0.651619in}}%
\pgfpathlineto{\pgfqpoint{5.954066in}{0.717279in}}%
\pgfpathlineto{\pgfqpoint{5.957421in}{0.727230in}}%
\pgfpathlineto{\pgfqpoint{5.960777in}{0.718127in}}%
\pgfpathlineto{\pgfqpoint{5.964133in}{0.676369in}}%
\pgfpathlineto{\pgfqpoint{5.967489in}{0.684319in}}%
\pgfpathlineto{\pgfqpoint{5.970845in}{0.717289in}}%
\pgfpathlineto{\pgfqpoint{5.974200in}{0.726326in}}%
\pgfpathlineto{\pgfqpoint{5.977556in}{0.723017in}}%
\pgfpathlineto{\pgfqpoint{5.980912in}{0.703335in}}%
\pgfpathlineto{\pgfqpoint{5.984268in}{0.614558in}}%
\pgfpathlineto{\pgfqpoint{5.987623in}{0.714313in}}%
\pgfpathlineto{\pgfqpoint{5.990979in}{0.735765in}}%
\pgfpathlineto{\pgfqpoint{5.994335in}{0.740848in}}%
\pgfpathlineto{\pgfqpoint{5.997691in}{0.727110in}}%
\pgfpathlineto{\pgfqpoint{6.001047in}{0.655641in}}%
\pgfpathlineto{\pgfqpoint{6.004402in}{0.716199in}}%
\pgfpathlineto{\pgfqpoint{6.007758in}{0.740330in}}%
\pgfpathlineto{\pgfqpoint{6.011114in}{0.747201in}}%
\pgfpathlineto{\pgfqpoint{6.014470in}{0.743063in}}%
\pgfpathlineto{\pgfqpoint{6.017825in}{0.723359in}}%
\pgfpathlineto{\pgfqpoint{6.021181in}{0.620039in}}%
\pgfpathlineto{\pgfqpoint{6.024537in}{0.732337in}}%
\pgfpathlineto{\pgfqpoint{6.027893in}{0.754420in}}%
\pgfpathlineto{\pgfqpoint{6.031249in}{0.760439in}}%
\pgfpathlineto{\pgfqpoint{6.034604in}{0.749801in}}%
\pgfpathlineto{\pgfqpoint{6.037960in}{0.713997in}}%
\pgfpathlineto{\pgfqpoint{6.041316in}{0.695705in}}%
\pgfpathlineto{\pgfqpoint{6.044672in}{0.735564in}}%
\pgfpathlineto{\pgfqpoint{6.048027in}{0.743796in}}%
\pgfpathlineto{\pgfqpoint{6.051383in}{0.736597in}}%
\pgfpathlineto{\pgfqpoint{6.054739in}{0.700726in}}%
\pgfpathlineto{\pgfqpoint{6.058095in}{0.709899in}}%
\pgfpathlineto{\pgfqpoint{6.061451in}{0.748557in}}%
\pgfpathlineto{\pgfqpoint{6.064806in}{0.761188in}}%
\pgfpathlineto{\pgfqpoint{6.068162in}{0.756751in}}%
\pgfpathlineto{\pgfqpoint{6.071518in}{0.736221in}}%
\pgfpathlineto{\pgfqpoint{6.074874in}{0.686624in}}%
\pgfpathlineto{\pgfqpoint{6.078230in}{0.687715in}}%
\pgfpathlineto{\pgfqpoint{6.081585in}{0.710612in}}%
\pgfpathlineto{\pgfqpoint{6.084941in}{0.701383in}}%
\pgfpathlineto{\pgfqpoint{6.088297in}{0.543954in}}%
\pgfpathlineto{\pgfqpoint{6.091653in}{0.714652in}}%
\pgfpathlineto{\pgfqpoint{6.095008in}{0.738326in}}%
\pgfpathlineto{\pgfqpoint{6.098364in}{0.742816in}}%
\pgfpathlineto{\pgfqpoint{6.101720in}{0.728233in}}%
\pgfpathlineto{\pgfqpoint{6.105076in}{0.695037in}}%
\pgfpathlineto{\pgfqpoint{6.108432in}{0.635240in}}%
\pgfpathlineto{\pgfqpoint{6.111787in}{0.690985in}}%
\pgfpathlineto{\pgfqpoint{6.115143in}{0.693359in}}%
\pgfpathlineto{\pgfqpoint{6.118499in}{0.661277in}}%
\pgfpathlineto{\pgfqpoint{6.121855in}{0.673971in}}%
\pgfpathlineto{\pgfqpoint{6.125210in}{0.713397in}}%
\pgfpathlineto{\pgfqpoint{6.128566in}{0.726030in}}%
\pgfpathlineto{\pgfqpoint{6.131922in}{0.719571in}}%
\pgfpathlineto{\pgfqpoint{6.135278in}{0.688886in}}%
\pgfpathlineto{\pgfqpoint{6.138634in}{0.652548in}}%
\pgfpathlineto{\pgfqpoint{6.141989in}{0.702838in}}%
\pgfpathlineto{\pgfqpoint{6.145345in}{0.713917in}}%
\pgfpathlineto{\pgfqpoint{6.148701in}{0.710110in}}%
\pgfpathlineto{\pgfqpoint{6.152057in}{0.685827in}}%
\pgfpathlineto{\pgfqpoint{6.155412in}{0.653101in}}%
\pgfpathlineto{\pgfqpoint{6.158768in}{0.711292in}}%
\pgfpathlineto{\pgfqpoint{6.162124in}{0.728667in}}%
\pgfpathlineto{\pgfqpoint{6.165480in}{0.729602in}}%
\pgfpathlineto{\pgfqpoint{6.168836in}{0.706979in}}%
\pgfpathlineto{\pgfqpoint{6.172191in}{0.652982in}}%
\pgfpathlineto{\pgfqpoint{6.175547in}{0.720170in}}%
\pgfpathlineto{\pgfqpoint{6.178903in}{0.736998in}}%
\pgfpathlineto{\pgfqpoint{6.182259in}{0.740638in}}%
\pgfpathlineto{\pgfqpoint{6.185614in}{0.733545in}}%
\pgfpathlineto{\pgfqpoint{6.188970in}{0.707484in}}%
\pgfpathlineto{\pgfqpoint{6.192326in}{0.674161in}}%
\pgfpathlineto{\pgfqpoint{6.195682in}{0.731906in}}%
\pgfpathlineto{\pgfqpoint{6.199038in}{0.749761in}}%
\pgfpathlineto{\pgfqpoint{6.202393in}{0.752954in}}%
\pgfpathlineto{\pgfqpoint{6.205749in}{0.739638in}}%
\pgfpathlineto{\pgfqpoint{6.209105in}{0.706363in}}%
\pgfpathlineto{\pgfqpoint{6.212461in}{0.658276in}}%
\pgfpathlineto{\pgfqpoint{6.215816in}{0.709136in}}%
\pgfpathlineto{\pgfqpoint{6.219172in}{0.714242in}}%
\pgfpathlineto{\pgfqpoint{6.222528in}{0.695132in}}%
\pgfpathlineto{\pgfqpoint{6.225884in}{0.654873in}}%
\pgfpathlineto{\pgfqpoint{6.229240in}{0.720027in}}%
\pgfpathlineto{\pgfqpoint{6.232595in}{0.738463in}}%
\pgfpathlineto{\pgfqpoint{6.235951in}{0.738807in}}%
\pgfpathlineto{\pgfqpoint{6.239307in}{0.720434in}}%
\pgfpathlineto{\pgfqpoint{6.242663in}{0.673456in}}%
\pgfpathlineto{\pgfqpoint{6.246019in}{0.674540in}}%
\pgfpathlineto{\pgfqpoint{6.249374in}{0.701281in}}%
\pgfpathlineto{\pgfqpoint{6.252730in}{0.699661in}}%
\pgfpathlineto{\pgfqpoint{6.256086in}{0.666669in}}%
\pgfpathlineto{\pgfqpoint{6.259442in}{0.676532in}}%
\pgfpathlineto{\pgfqpoint{6.262797in}{0.716097in}}%
\pgfpathlineto{\pgfqpoint{6.266153in}{0.728719in}}%
\pgfpathlineto{\pgfqpoint{6.269509in}{0.722656in}}%
\pgfpathlineto{\pgfqpoint{6.276221in}{0.659110in}}%
\pgfpathlineto{\pgfqpoint{6.279576in}{0.708015in}}%
\pgfpathlineto{\pgfqpoint{6.282932in}{0.719381in}}%
\pgfpathlineto{\pgfqpoint{6.286288in}{0.716405in}}%
\pgfpathlineto{\pgfqpoint{6.289644in}{0.694976in}}%
\pgfpathlineto{\pgfqpoint{6.292999in}{0.640427in}}%
\pgfpathlineto{\pgfqpoint{6.296355in}{0.713086in}}%
\pgfpathlineto{\pgfqpoint{6.299711in}{0.732443in}}%
\pgfpathlineto{\pgfqpoint{6.303067in}{0.735328in}}%
\pgfpathlineto{\pgfqpoint{6.306423in}{0.717162in}}%
\pgfpathlineto{\pgfqpoint{6.309778in}{0.475464in}}%
\pgfpathlineto{\pgfqpoint{6.313134in}{0.716463in}}%
\pgfpathlineto{\pgfqpoint{6.316490in}{0.736393in}}%
\pgfpathlineto{\pgfqpoint{6.319846in}{0.741130in}}%
\pgfpathlineto{\pgfqpoint{6.323201in}{0.734597in}}%
\pgfpathlineto{\pgfqpoint{6.326557in}{0.708925in}}%
\pgfpathlineto{\pgfqpoint{6.329913in}{0.675163in}}%
\pgfpathlineto{\pgfqpoint{6.333269in}{0.733283in}}%
\pgfpathlineto{\pgfqpoint{6.336625in}{0.751063in}}%
\pgfpathlineto{\pgfqpoint{6.339980in}{0.753819in}}%
\pgfpathlineto{\pgfqpoint{6.343336in}{0.738079in}}%
\pgfpathlineto{\pgfqpoint{6.346692in}{0.677286in}}%
\pgfpathlineto{\pgfqpoint{6.350048in}{0.717130in}}%
\pgfpathlineto{\pgfqpoint{6.353403in}{0.741414in}}%
\pgfpathlineto{\pgfqpoint{6.356759in}{0.746179in}}%
\pgfpathlineto{\pgfqpoint{6.360115in}{0.737266in}}%
\pgfpathlineto{\pgfqpoint{6.363471in}{0.699239in}}%
\pgfpathlineto{\pgfqpoint{6.366827in}{0.711562in}}%
\pgfpathlineto{\pgfqpoint{6.370182in}{0.748967in}}%
\pgfpathlineto{\pgfqpoint{6.373538in}{0.761441in}}%
\pgfpathlineto{\pgfqpoint{6.376894in}{0.757639in}}%
\pgfpathlineto{\pgfqpoint{6.383605in}{0.733396in}}%
\pgfpathlineto{\pgfqpoint{6.383605in}{0.733396in}}%
\pgfusepath{stroke}%
\end{pgfscope}%
\begin{pgfscope}%
\pgfpathrectangle{\pgfqpoint{4.621821in}{0.320625in}}{\pgfqpoint{1.845679in}{1.164375in}}%
\pgfusepath{clip}%
\pgfsetrectcap%
\pgfsetroundjoin%
\pgfsetlinewidth{0.602250pt}%
\definecolor{currentstroke}{rgb}{0.819608,0.286275,0.356863}%
\pgfsetstrokecolor{currentstroke}%
\pgfsetstrokeopacity{0.120000}%
\pgfsetdash{}{0pt}%
\pgfpathmoveto{\pgfqpoint{4.705715in}{1.471521in}}%
\pgfpathlineto{\pgfqpoint{4.729206in}{1.429543in}}%
\pgfpathlineto{\pgfqpoint{4.766120in}{1.360190in}}%
\pgfpathlineto{\pgfqpoint{4.786254in}{1.325552in}}%
\pgfpathlineto{\pgfqpoint{4.809745in}{1.287455in}}%
\pgfpathlineto{\pgfqpoint{4.839947in}{1.240860in}}%
\pgfpathlineto{\pgfqpoint{4.873504in}{1.191214in}}%
\pgfpathlineto{\pgfqpoint{4.913774in}{1.133723in}}%
\pgfpathlineto{\pgfqpoint{4.960755in}{1.068699in}}%
\pgfpathlineto{\pgfqpoint{5.017803in}{0.991748in}}%
\pgfpathlineto{\pgfqpoint{5.088274in}{0.898544in}}%
\pgfpathlineto{\pgfqpoint{5.145323in}{0.823984in}}%
\pgfpathlineto{\pgfqpoint{5.162102in}{0.802211in}}%
\pgfpathlineto{\pgfqpoint{5.182236in}{0.775408in}}%
\pgfpathlineto{\pgfqpoint{5.195659in}{0.760052in}}%
\pgfpathlineto{\pgfqpoint{5.202371in}{0.751874in}}%
\pgfpathlineto{\pgfqpoint{5.212438in}{0.737278in}}%
\pgfpathlineto{\pgfqpoint{5.229217in}{0.711447in}}%
\pgfpathlineto{\pgfqpoint{5.235929in}{0.703464in}}%
\pgfpathlineto{\pgfqpoint{5.242640in}{0.698332in}}%
\pgfpathlineto{\pgfqpoint{5.249352in}{0.695543in}}%
\pgfpathlineto{\pgfqpoint{5.256063in}{0.692717in}}%
\pgfpathlineto{\pgfqpoint{5.259419in}{0.690004in}}%
\pgfpathlineto{\pgfqpoint{5.262775in}{0.685608in}}%
\pgfpathlineto{\pgfqpoint{5.266131in}{0.680088in}}%
\pgfpathlineto{\pgfqpoint{5.272842in}{0.665602in}}%
\pgfpathlineto{\pgfqpoint{5.279554in}{0.645838in}}%
\pgfpathlineto{\pgfqpoint{5.286265in}{0.623894in}}%
\pgfpathlineto{\pgfqpoint{5.289621in}{0.617519in}}%
\pgfpathlineto{\pgfqpoint{5.292977in}{0.615884in}}%
\pgfpathlineto{\pgfqpoint{5.296333in}{0.615612in}}%
\pgfpathlineto{\pgfqpoint{5.299689in}{0.613628in}}%
\pgfpathlineto{\pgfqpoint{5.303044in}{0.608194in}}%
\pgfpathlineto{\pgfqpoint{5.309756in}{0.586513in}}%
\pgfpathlineto{\pgfqpoint{5.313112in}{0.594445in}}%
\pgfpathlineto{\pgfqpoint{5.316467in}{0.613691in}}%
\pgfpathlineto{\pgfqpoint{5.319823in}{0.629075in}}%
\pgfpathlineto{\pgfqpoint{5.323179in}{0.640355in}}%
\pgfpathlineto{\pgfqpoint{5.326535in}{0.648543in}}%
\pgfpathlineto{\pgfqpoint{5.329891in}{0.654225in}}%
\pgfpathlineto{\pgfqpoint{5.333246in}{0.657614in}}%
\pgfpathlineto{\pgfqpoint{5.336602in}{0.658603in}}%
\pgfpathlineto{\pgfqpoint{5.339958in}{0.656652in}}%
\pgfpathlineto{\pgfqpoint{5.343314in}{0.651641in}}%
\pgfpathlineto{\pgfqpoint{5.346669in}{0.656125in}}%
\pgfpathlineto{\pgfqpoint{5.353381in}{0.689628in}}%
\pgfpathlineto{\pgfqpoint{5.356737in}{0.700738in}}%
\pgfpathlineto{\pgfqpoint{5.360093in}{0.708125in}}%
\pgfpathlineto{\pgfqpoint{5.363448in}{0.712518in}}%
\pgfpathlineto{\pgfqpoint{5.366804in}{0.714276in}}%
\pgfpathlineto{\pgfqpoint{5.370160in}{0.713435in}}%
\pgfpathlineto{\pgfqpoint{5.373516in}{0.709667in}}%
\pgfpathlineto{\pgfqpoint{5.376871in}{0.702009in}}%
\pgfpathlineto{\pgfqpoint{5.380227in}{0.687835in}}%
\pgfpathlineto{\pgfqpoint{5.383583in}{0.656720in}}%
\pgfpathlineto{\pgfqpoint{5.386939in}{0.632437in}}%
\pgfpathlineto{\pgfqpoint{5.390295in}{0.683709in}}%
\pgfpathlineto{\pgfqpoint{5.393650in}{0.704383in}}%
\pgfpathlineto{\pgfqpoint{5.397006in}{0.716316in}}%
\pgfpathlineto{\pgfqpoint{5.400362in}{0.723276in}}%
\pgfpathlineto{\pgfqpoint{5.403718in}{0.726216in}}%
\pgfpathlineto{\pgfqpoint{5.407074in}{0.724775in}}%
\pgfpathlineto{\pgfqpoint{5.410429in}{0.716743in}}%
\pgfpathlineto{\pgfqpoint{5.413785in}{0.697855in}}%
\pgfpathlineto{\pgfqpoint{5.417141in}{0.659537in}}%
\pgfpathlineto{\pgfqpoint{5.420497in}{0.665686in}}%
\pgfpathlineto{\pgfqpoint{5.423852in}{0.685470in}}%
\pgfpathlineto{\pgfqpoint{5.427208in}{0.686624in}}%
\pgfpathlineto{\pgfqpoint{5.430564in}{0.669568in}}%
\pgfpathlineto{\pgfqpoint{5.433920in}{0.609086in}}%
\pgfpathlineto{\pgfqpoint{5.437276in}{0.682940in}}%
\pgfpathlineto{\pgfqpoint{5.440631in}{0.705134in}}%
\pgfpathlineto{\pgfqpoint{5.443987in}{0.711421in}}%
\pgfpathlineto{\pgfqpoint{5.447343in}{0.701588in}}%
\pgfpathlineto{\pgfqpoint{5.450699in}{0.675371in}}%
\pgfpathlineto{\pgfqpoint{5.454054in}{0.617403in}}%
\pgfpathlineto{\pgfqpoint{5.457410in}{0.667863in}}%
\pgfpathlineto{\pgfqpoint{5.460766in}{0.680974in}}%
\pgfpathlineto{\pgfqpoint{5.464122in}{0.677146in}}%
\pgfpathlineto{\pgfqpoint{5.467478in}{0.650072in}}%
\pgfpathlineto{\pgfqpoint{5.470833in}{0.635436in}}%
\pgfpathlineto{\pgfqpoint{5.474189in}{0.684449in}}%
\pgfpathlineto{\pgfqpoint{5.477545in}{0.701081in}}%
\pgfpathlineto{\pgfqpoint{5.480901in}{0.702770in}}%
\pgfpathlineto{\pgfqpoint{5.484256in}{0.681905in}}%
\pgfpathlineto{\pgfqpoint{5.487612in}{0.637574in}}%
\pgfpathlineto{\pgfqpoint{5.490968in}{0.700709in}}%
\pgfpathlineto{\pgfqpoint{5.494324in}{0.719857in}}%
\pgfpathlineto{\pgfqpoint{5.497680in}{0.727295in}}%
\pgfpathlineto{\pgfqpoint{5.501035in}{0.726977in}}%
\pgfpathlineto{\pgfqpoint{5.504391in}{0.718238in}}%
\pgfpathlineto{\pgfqpoint{5.507747in}{0.693227in}}%
\pgfpathlineto{\pgfqpoint{5.511103in}{0.642251in}}%
\pgfpathlineto{\pgfqpoint{5.514458in}{0.711104in}}%
\pgfpathlineto{\pgfqpoint{5.517814in}{0.731539in}}%
\pgfpathlineto{\pgfqpoint{5.521170in}{0.739098in}}%
\pgfpathlineto{\pgfqpoint{5.524526in}{0.734874in}}%
\pgfpathlineto{\pgfqpoint{5.527882in}{0.715927in}}%
\pgfpathlineto{\pgfqpoint{5.531237in}{0.665219in}}%
\pgfpathlineto{\pgfqpoint{5.534593in}{0.680194in}}%
\pgfpathlineto{\pgfqpoint{5.537949in}{0.705822in}}%
\pgfpathlineto{\pgfqpoint{5.541305in}{0.708597in}}%
\pgfpathlineto{\pgfqpoint{5.544660in}{0.693931in}}%
\pgfpathlineto{\pgfqpoint{5.548016in}{0.584144in}}%
\pgfpathlineto{\pgfqpoint{5.551372in}{0.699110in}}%
\pgfpathlineto{\pgfqpoint{5.554728in}{0.724055in}}%
\pgfpathlineto{\pgfqpoint{5.558084in}{0.731741in}}%
\pgfpathlineto{\pgfqpoint{5.561439in}{0.723192in}}%
\pgfpathlineto{\pgfqpoint{5.564795in}{0.697646in}}%
\pgfpathlineto{\pgfqpoint{5.568151in}{0.592670in}}%
\pgfpathlineto{\pgfqpoint{5.571507in}{0.687205in}}%
\pgfpathlineto{\pgfqpoint{5.574863in}{0.702173in}}%
\pgfpathlineto{\pgfqpoint{5.578218in}{0.699430in}}%
\pgfpathlineto{\pgfqpoint{5.584930in}{0.648821in}}%
\pgfpathlineto{\pgfqpoint{5.588286in}{0.704074in}}%
\pgfpathlineto{\pgfqpoint{5.591641in}{0.721809in}}%
\pgfpathlineto{\pgfqpoint{5.594997in}{0.724336in}}%
\pgfpathlineto{\pgfqpoint{5.598353in}{0.708574in}}%
\pgfpathlineto{\pgfqpoint{5.601709in}{0.653200in}}%
\pgfpathlineto{\pgfqpoint{5.608420in}{0.711797in}}%
\pgfpathlineto{\pgfqpoint{5.611776in}{0.719107in}}%
\pgfpathlineto{\pgfqpoint{5.615132in}{0.715043in}}%
\pgfpathlineto{\pgfqpoint{5.618488in}{0.695199in}}%
\pgfpathlineto{\pgfqpoint{5.621843in}{0.599677in}}%
\pgfpathlineto{\pgfqpoint{5.625199in}{0.705384in}}%
\pgfpathlineto{\pgfqpoint{5.628555in}{0.727729in}}%
\pgfpathlineto{\pgfqpoint{5.631911in}{0.734710in}}%
\pgfpathlineto{\pgfqpoint{5.635267in}{0.726409in}}%
\pgfpathlineto{\pgfqpoint{5.638622in}{0.692188in}}%
\pgfpathlineto{\pgfqpoint{5.641978in}{0.679389in}}%
\pgfpathlineto{\pgfqpoint{5.645334in}{0.720940in}}%
\pgfpathlineto{\pgfqpoint{5.648690in}{0.733495in}}%
\pgfpathlineto{\pgfqpoint{5.652045in}{0.734572in}}%
\pgfpathlineto{\pgfqpoint{5.655401in}{0.724675in}}%
\pgfpathlineto{\pgfqpoint{5.658757in}{0.690651in}}%
\pgfpathlineto{\pgfqpoint{5.662113in}{0.688642in}}%
\pgfpathlineto{\pgfqpoint{5.665469in}{0.731171in}}%
\pgfpathlineto{\pgfqpoint{5.668824in}{0.746505in}}%
\pgfpathlineto{\pgfqpoint{5.672180in}{0.748692in}}%
\pgfpathlineto{\pgfqpoint{5.675536in}{0.734885in}}%
\pgfpathlineto{\pgfqpoint{5.678892in}{0.694580in}}%
\pgfpathlineto{\pgfqpoint{5.682247in}{0.690166in}}%
\pgfpathlineto{\pgfqpoint{5.685603in}{0.725096in}}%
\pgfpathlineto{\pgfqpoint{5.688959in}{0.733567in}}%
\pgfpathlineto{\pgfqpoint{5.692315in}{0.728828in}}%
\pgfpathlineto{\pgfqpoint{5.695671in}{0.704741in}}%
\pgfpathlineto{\pgfqpoint{5.699026in}{0.667225in}}%
\pgfpathlineto{\pgfqpoint{5.702382in}{0.728495in}}%
\pgfpathlineto{\pgfqpoint{5.705738in}{0.747214in}}%
\pgfpathlineto{\pgfqpoint{5.709094in}{0.751173in}}%
\pgfpathlineto{\pgfqpoint{5.712449in}{0.739308in}}%
\pgfpathlineto{\pgfqpoint{5.715805in}{0.711658in}}%
\pgfpathlineto{\pgfqpoint{5.719161in}{0.593333in}}%
\pgfpathlineto{\pgfqpoint{5.722517in}{0.699210in}}%
\pgfpathlineto{\pgfqpoint{5.725873in}{0.710915in}}%
\pgfpathlineto{\pgfqpoint{5.729228in}{0.701934in}}%
\pgfpathlineto{\pgfqpoint{5.732584in}{0.641919in}}%
\pgfpathlineto{\pgfqpoint{5.735940in}{0.698590in}}%
\pgfpathlineto{\pgfqpoint{5.739296in}{0.728050in}}%
\pgfpathlineto{\pgfqpoint{5.742652in}{0.737934in}}%
\pgfpathlineto{\pgfqpoint{5.746007in}{0.731425in}}%
\pgfpathlineto{\pgfqpoint{5.749363in}{0.709913in}}%
\pgfpathlineto{\pgfqpoint{5.752719in}{0.652568in}}%
\pgfpathlineto{\pgfqpoint{5.759430in}{0.696099in}}%
\pgfpathlineto{\pgfqpoint{5.762786in}{0.693565in}}%
\pgfpathlineto{\pgfqpoint{5.766142in}{0.662101in}}%
\pgfpathlineto{\pgfqpoint{5.769498in}{0.666430in}}%
\pgfpathlineto{\pgfqpoint{5.772854in}{0.708399in}}%
\pgfpathlineto{\pgfqpoint{5.776209in}{0.723076in}}%
\pgfpathlineto{\pgfqpoint{5.779565in}{0.722017in}}%
\pgfpathlineto{\pgfqpoint{5.782921in}{0.702479in}}%
\pgfpathlineto{\pgfqpoint{5.786277in}{0.634054in}}%
\pgfpathlineto{\pgfqpoint{5.789632in}{0.682797in}}%
\pgfpathlineto{\pgfqpoint{5.792988in}{0.705250in}}%
\pgfpathlineto{\pgfqpoint{5.796344in}{0.709313in}}%
\pgfpathlineto{\pgfqpoint{5.799700in}{0.700042in}}%
\pgfpathlineto{\pgfqpoint{5.803056in}{0.661798in}}%
\pgfpathlineto{\pgfqpoint{5.806411in}{0.674508in}}%
\pgfpathlineto{\pgfqpoint{5.809767in}{0.712265in}}%
\pgfpathlineto{\pgfqpoint{5.813123in}{0.725793in}}%
\pgfpathlineto{\pgfqpoint{5.816479in}{0.724933in}}%
\pgfpathlineto{\pgfqpoint{5.819834in}{0.703857in}}%
\pgfpathlineto{\pgfqpoint{5.823190in}{0.590383in}}%
\pgfpathlineto{\pgfqpoint{5.826546in}{0.705102in}}%
\pgfpathlineto{\pgfqpoint{5.829902in}{0.724579in}}%
\pgfpathlineto{\pgfqpoint{5.833258in}{0.729896in}}%
\pgfpathlineto{\pgfqpoint{5.836613in}{0.725177in}}%
\pgfpathlineto{\pgfqpoint{5.839969in}{0.705827in}}%
\pgfpathlineto{\pgfqpoint{5.843325in}{0.553966in}}%
\pgfpathlineto{\pgfqpoint{5.846681in}{0.712041in}}%
\pgfpathlineto{\pgfqpoint{5.850036in}{0.735525in}}%
\pgfpathlineto{\pgfqpoint{5.853392in}{0.743587in}}%
\pgfpathlineto{\pgfqpoint{5.856748in}{0.737608in}}%
\pgfpathlineto{\pgfqpoint{5.860104in}{0.710485in}}%
\pgfpathlineto{\pgfqpoint{5.863460in}{0.659425in}}%
\pgfpathlineto{\pgfqpoint{5.866815in}{0.721467in}}%
\pgfpathlineto{\pgfqpoint{5.870171in}{0.736653in}}%
\pgfpathlineto{\pgfqpoint{5.873527in}{0.738710in}}%
\pgfpathlineto{\pgfqpoint{5.876883in}{0.729213in}}%
\pgfpathlineto{\pgfqpoint{5.880238in}{0.694885in}}%
\pgfpathlineto{\pgfqpoint{5.883594in}{0.694887in}}%
\pgfpathlineto{\pgfqpoint{5.886950in}{0.736735in}}%
\pgfpathlineto{\pgfqpoint{5.890306in}{0.751734in}}%
\pgfpathlineto{\pgfqpoint{5.893662in}{0.753202in}}%
\pgfpathlineto{\pgfqpoint{5.897017in}{0.738948in}}%
\pgfpathlineto{\pgfqpoint{5.903729in}{0.668916in}}%
\pgfpathlineto{\pgfqpoint{5.907085in}{0.714856in}}%
\pgfpathlineto{\pgfqpoint{5.910441in}{0.722696in}}%
\pgfpathlineto{\pgfqpoint{5.913796in}{0.713254in}}%
\pgfpathlineto{\pgfqpoint{5.917152in}{0.662348in}}%
\pgfpathlineto{\pgfqpoint{5.920508in}{0.703072in}}%
\pgfpathlineto{\pgfqpoint{5.923864in}{0.734598in}}%
\pgfpathlineto{\pgfqpoint{5.927219in}{0.745498in}}%
\pgfpathlineto{\pgfqpoint{5.930575in}{0.740555in}}%
\pgfpathlineto{\pgfqpoint{5.933931in}{0.721348in}}%
\pgfpathlineto{\pgfqpoint{5.937287in}{0.678437in}}%
\pgfpathlineto{\pgfqpoint{5.940643in}{0.663494in}}%
\pgfpathlineto{\pgfqpoint{5.943998in}{0.694525in}}%
\pgfpathlineto{\pgfqpoint{5.947354in}{0.693057in}}%
\pgfpathlineto{\pgfqpoint{5.950710in}{0.656921in}}%
\pgfpathlineto{\pgfqpoint{5.954066in}{0.676184in}}%
\pgfpathlineto{\pgfqpoint{5.957421in}{0.713847in}}%
\pgfpathlineto{\pgfqpoint{5.960777in}{0.727202in}}%
\pgfpathlineto{\pgfqpoint{5.964133in}{0.724320in}}%
\pgfpathlineto{\pgfqpoint{5.967489in}{0.704533in}}%
\pgfpathlineto{\pgfqpoint{5.970845in}{0.648395in}}%
\pgfpathlineto{\pgfqpoint{5.977556in}{0.696469in}}%
\pgfpathlineto{\pgfqpoint{5.980912in}{0.698353in}}%
\pgfpathlineto{\pgfqpoint{5.984268in}{0.682191in}}%
\pgfpathlineto{\pgfqpoint{5.987623in}{0.549664in}}%
\pgfpathlineto{\pgfqpoint{5.990979in}{0.691781in}}%
\pgfpathlineto{\pgfqpoint{5.994335in}{0.715152in}}%
\pgfpathlineto{\pgfqpoint{5.997691in}{0.721792in}}%
\pgfpathlineto{\pgfqpoint{6.001047in}{0.711220in}}%
\pgfpathlineto{\pgfqpoint{6.004402in}{0.671293in}}%
\pgfpathlineto{\pgfqpoint{6.007758in}{0.673422in}}%
\pgfpathlineto{\pgfqpoint{6.011114in}{0.708773in}}%
\pgfpathlineto{\pgfqpoint{6.014470in}{0.719449in}}%
\pgfpathlineto{\pgfqpoint{6.017825in}{0.718832in}}%
\pgfpathlineto{\pgfqpoint{6.021181in}{0.706056in}}%
\pgfpathlineto{\pgfqpoint{6.024537in}{0.657803in}}%
\pgfpathlineto{\pgfqpoint{6.031249in}{0.721543in}}%
\pgfpathlineto{\pgfqpoint{6.034604in}{0.733711in}}%
\pgfpathlineto{\pgfqpoint{6.037960in}{0.732467in}}%
\pgfpathlineto{\pgfqpoint{6.041316in}{0.711288in}}%
\pgfpathlineto{\pgfqpoint{6.044672in}{0.609380in}}%
\pgfpathlineto{\pgfqpoint{6.048027in}{0.714413in}}%
\pgfpathlineto{\pgfqpoint{6.051383in}{0.733742in}}%
\pgfpathlineto{\pgfqpoint{6.054739in}{0.739202in}}%
\pgfpathlineto{\pgfqpoint{6.058095in}{0.734864in}}%
\pgfpathlineto{\pgfqpoint{6.061451in}{0.716638in}}%
\pgfpathlineto{\pgfqpoint{6.064806in}{0.601113in}}%
\pgfpathlineto{\pgfqpoint{6.068162in}{0.718726in}}%
\pgfpathlineto{\pgfqpoint{6.071518in}{0.743554in}}%
\pgfpathlineto{\pgfqpoint{6.074874in}{0.752428in}}%
\pgfpathlineto{\pgfqpoint{6.078230in}{0.747720in}}%
\pgfpathlineto{\pgfqpoint{6.081585in}{0.725656in}}%
\pgfpathlineto{\pgfqpoint{6.084941in}{0.630125in}}%
\pgfpathlineto{\pgfqpoint{6.088297in}{0.712877in}}%
\pgfpathlineto{\pgfqpoint{6.091653in}{0.732004in}}%
\pgfpathlineto{\pgfqpoint{6.095008in}{0.734056in}}%
\pgfpathlineto{\pgfqpoint{6.098364in}{0.721671in}}%
\pgfpathlineto{\pgfqpoint{6.101720in}{0.664987in}}%
\pgfpathlineto{\pgfqpoint{6.105076in}{0.712276in}}%
\pgfpathlineto{\pgfqpoint{6.108432in}{0.742294in}}%
\pgfpathlineto{\pgfqpoint{6.111787in}{0.752802in}}%
\pgfpathlineto{\pgfqpoint{6.115143in}{0.748241in}}%
\pgfpathlineto{\pgfqpoint{6.118499in}{0.729495in}}%
\pgfpathlineto{\pgfqpoint{6.121855in}{0.687999in}}%
\pgfpathlineto{\pgfqpoint{6.125210in}{0.668634in}}%
\pgfpathlineto{\pgfqpoint{6.128566in}{0.701786in}}%
\pgfpathlineto{\pgfqpoint{6.131922in}{0.701212in}}%
\pgfpathlineto{\pgfqpoint{6.135278in}{0.668112in}}%
\pgfpathlineto{\pgfqpoint{6.138634in}{0.680444in}}%
\pgfpathlineto{\pgfqpoint{6.141989in}{0.720247in}}%
\pgfpathlineto{\pgfqpoint{6.145345in}{0.734324in}}%
\pgfpathlineto{\pgfqpoint{6.148701in}{0.732267in}}%
\pgfpathlineto{\pgfqpoint{6.152057in}{0.714598in}}%
\pgfpathlineto{\pgfqpoint{6.155412in}{0.675377in}}%
\pgfpathlineto{\pgfqpoint{6.158768in}{0.650191in}}%
\pgfpathlineto{\pgfqpoint{6.162124in}{0.687212in}}%
\pgfpathlineto{\pgfqpoint{6.165480in}{0.689206in}}%
\pgfpathlineto{\pgfqpoint{6.168836in}{0.665203in}}%
\pgfpathlineto{\pgfqpoint{6.172191in}{0.648493in}}%
\pgfpathlineto{\pgfqpoint{6.175547in}{0.700245in}}%
\pgfpathlineto{\pgfqpoint{6.178903in}{0.717441in}}%
\pgfpathlineto{\pgfqpoint{6.182259in}{0.718810in}}%
\pgfpathlineto{\pgfqpoint{6.185614in}{0.701411in}}%
\pgfpathlineto{\pgfqpoint{6.188970in}{0.642972in}}%
\pgfpathlineto{\pgfqpoint{6.192326in}{0.677494in}}%
\pgfpathlineto{\pgfqpoint{6.195682in}{0.702771in}}%
\pgfpathlineto{\pgfqpoint{6.199038in}{0.708583in}}%
\pgfpathlineto{\pgfqpoint{6.202393in}{0.701945in}}%
\pgfpathlineto{\pgfqpoint{6.205749in}{0.674027in}}%
\pgfpathlineto{\pgfqpoint{6.209105in}{0.654150in}}%
\pgfpathlineto{\pgfqpoint{6.212461in}{0.705130in}}%
\pgfpathlineto{\pgfqpoint{6.215816in}{0.722095in}}%
\pgfpathlineto{\pgfqpoint{6.219172in}{0.724771in}}%
\pgfpathlineto{\pgfqpoint{6.222528in}{0.708943in}}%
\pgfpathlineto{\pgfqpoint{6.225884in}{0.628095in}}%
\pgfpathlineto{\pgfqpoint{6.229240in}{0.700916in}}%
\pgfpathlineto{\pgfqpoint{6.232595in}{0.724837in}}%
\pgfpathlineto{\pgfqpoint{6.235951in}{0.732788in}}%
\pgfpathlineto{\pgfqpoint{6.239307in}{0.731177in}}%
\pgfpathlineto{\pgfqpoint{6.242663in}{0.718393in}}%
\pgfpathlineto{\pgfqpoint{6.246019in}{0.674138in}}%
\pgfpathlineto{\pgfqpoint{6.249374in}{0.696228in}}%
\pgfpathlineto{\pgfqpoint{6.252730in}{0.731008in}}%
\pgfpathlineto{\pgfqpoint{6.256086in}{0.744233in}}%
\pgfpathlineto{\pgfqpoint{6.259442in}{0.744923in}}%
\pgfpathlineto{\pgfqpoint{6.262797in}{0.728597in}}%
\pgfpathlineto{\pgfqpoint{6.266153in}{0.665949in}}%
\pgfpathlineto{\pgfqpoint{6.269509in}{0.710307in}}%
\pgfpathlineto{\pgfqpoint{6.272865in}{0.735543in}}%
\pgfpathlineto{\pgfqpoint{6.276221in}{0.742656in}}%
\pgfpathlineto{\pgfqpoint{6.279576in}{0.738818in}}%
\pgfpathlineto{\pgfqpoint{6.282932in}{0.720033in}}%
\pgfpathlineto{\pgfqpoint{6.286288in}{0.553858in}}%
\pgfpathlineto{\pgfqpoint{6.289644in}{0.726262in}}%
\pgfpathlineto{\pgfqpoint{6.292999in}{0.749794in}}%
\pgfpathlineto{\pgfqpoint{6.296355in}{0.757574in}}%
\pgfpathlineto{\pgfqpoint{6.299711in}{0.750847in}}%
\pgfpathlineto{\pgfqpoint{6.303067in}{0.729269in}}%
\pgfpathlineto{\pgfqpoint{6.306423in}{0.669045in}}%
\pgfpathlineto{\pgfqpoint{6.309778in}{0.696964in}}%
\pgfpathlineto{\pgfqpoint{6.313134in}{0.718041in}}%
\pgfpathlineto{\pgfqpoint{6.316490in}{0.716180in}}%
\pgfpathlineto{\pgfqpoint{6.319846in}{0.688420in}}%
\pgfpathlineto{\pgfqpoint{6.323201in}{0.680350in}}%
\pgfpathlineto{\pgfqpoint{6.326557in}{0.727229in}}%
\pgfpathlineto{\pgfqpoint{6.329913in}{0.743188in}}%
\pgfpathlineto{\pgfqpoint{6.333269in}{0.743554in}}%
\pgfpathlineto{\pgfqpoint{6.336625in}{0.728301in}}%
\pgfpathlineto{\pgfqpoint{6.339980in}{0.698397in}}%
\pgfpathlineto{\pgfqpoint{6.343336in}{0.551701in}}%
\pgfpathlineto{\pgfqpoint{6.346692in}{0.680284in}}%
\pgfpathlineto{\pgfqpoint{6.350048in}{0.686098in}}%
\pgfpathlineto{\pgfqpoint{6.353403in}{0.658121in}}%
\pgfpathlineto{\pgfqpoint{6.356759in}{0.663514in}}%
\pgfpathlineto{\pgfqpoint{6.360115in}{0.706775in}}%
\pgfpathlineto{\pgfqpoint{6.363471in}{0.722147in}}%
\pgfpathlineto{\pgfqpoint{6.366827in}{0.721352in}}%
\pgfpathlineto{\pgfqpoint{6.380250in}{0.676277in}}%
\pgfpathlineto{\pgfqpoint{6.383605in}{0.665727in}}%
\pgfpathlineto{\pgfqpoint{6.383605in}{0.665727in}}%
\pgfusepath{stroke}%
\end{pgfscope}%
\begin{pgfscope}%
\pgfpathrectangle{\pgfqpoint{4.621821in}{0.320625in}}{\pgfqpoint{1.845679in}{1.164375in}}%
\pgfusepath{clip}%
\pgfsetrectcap%
\pgfsetroundjoin%
\pgfsetlinewidth{0.602250pt}%
\definecolor{currentstroke}{rgb}{0.819608,0.286275,0.356863}%
\pgfsetstrokecolor{currentstroke}%
\pgfsetstrokeopacity{0.120000}%
\pgfsetdash{}{0pt}%
\pgfpathmoveto{\pgfqpoint{4.705715in}{1.473414in}}%
\pgfpathlineto{\pgfqpoint{4.729206in}{1.430786in}}%
\pgfpathlineto{\pgfqpoint{4.776187in}{1.339404in}}%
\pgfpathlineto{\pgfqpoint{4.792966in}{1.310245in}}%
\pgfpathlineto{\pgfqpoint{4.813100in}{1.277468in}}%
\pgfpathlineto{\pgfqpoint{4.839947in}{1.236225in}}%
\pgfpathlineto{\pgfqpoint{4.873504in}{1.186906in}}%
\pgfpathlineto{\pgfqpoint{4.920485in}{1.119923in}}%
\pgfpathlineto{\pgfqpoint{4.987601in}{1.026222in}}%
\pgfpathlineto{\pgfqpoint{5.068140in}{0.915610in}}%
\pgfpathlineto{\pgfqpoint{5.168813in}{0.779574in}}%
\pgfpathlineto{\pgfqpoint{5.185592in}{0.759155in}}%
\pgfpathlineto{\pgfqpoint{5.195659in}{0.744208in}}%
\pgfpathlineto{\pgfqpoint{5.209082in}{0.723408in}}%
\pgfpathlineto{\pgfqpoint{5.215794in}{0.714760in}}%
\pgfpathlineto{\pgfqpoint{5.222506in}{0.708578in}}%
\pgfpathlineto{\pgfqpoint{5.229217in}{0.704636in}}%
\pgfpathlineto{\pgfqpoint{5.235929in}{0.700840in}}%
\pgfpathlineto{\pgfqpoint{5.239285in}{0.697715in}}%
\pgfpathlineto{\pgfqpoint{5.242640in}{0.693053in}}%
\pgfpathlineto{\pgfqpoint{5.249352in}{0.680765in}}%
\pgfpathlineto{\pgfqpoint{5.256063in}{0.664175in}}%
\pgfpathlineto{\pgfqpoint{5.269487in}{0.623483in}}%
\pgfpathlineto{\pgfqpoint{5.272842in}{0.617351in}}%
\pgfpathlineto{\pgfqpoint{5.282910in}{0.603903in}}%
\pgfpathlineto{\pgfqpoint{5.286265in}{0.603748in}}%
\pgfpathlineto{\pgfqpoint{5.289621in}{0.612587in}}%
\pgfpathlineto{\pgfqpoint{5.296333in}{0.636615in}}%
\pgfpathlineto{\pgfqpoint{5.299689in}{0.645751in}}%
\pgfpathlineto{\pgfqpoint{5.303044in}{0.652591in}}%
\pgfpathlineto{\pgfqpoint{5.306400in}{0.657227in}}%
\pgfpathlineto{\pgfqpoint{5.309756in}{0.659590in}}%
\pgfpathlineto{\pgfqpoint{5.313112in}{0.659307in}}%
\pgfpathlineto{\pgfqpoint{5.319823in}{0.652443in}}%
\pgfpathlineto{\pgfqpoint{5.323179in}{0.664588in}}%
\pgfpathlineto{\pgfqpoint{5.326535in}{0.682446in}}%
\pgfpathlineto{\pgfqpoint{5.329891in}{0.696580in}}%
\pgfpathlineto{\pgfqpoint{5.333246in}{0.706813in}}%
\pgfpathlineto{\pgfqpoint{5.336602in}{0.714011in}}%
\pgfpathlineto{\pgfqpoint{5.339958in}{0.718782in}}%
\pgfpathlineto{\pgfqpoint{5.343314in}{0.721468in}}%
\pgfpathlineto{\pgfqpoint{5.346669in}{0.722207in}}%
\pgfpathlineto{\pgfqpoint{5.350025in}{0.720947in}}%
\pgfpathlineto{\pgfqpoint{5.353381in}{0.717407in}}%
\pgfpathlineto{\pgfqpoint{5.356737in}{0.710907in}}%
\pgfpathlineto{\pgfqpoint{5.360093in}{0.699839in}}%
\pgfpathlineto{\pgfqpoint{5.363448in}{0.679393in}}%
\pgfpathlineto{\pgfqpoint{5.366804in}{0.616024in}}%
\pgfpathlineto{\pgfqpoint{5.370160in}{0.668114in}}%
\pgfpathlineto{\pgfqpoint{5.373516in}{0.698349in}}%
\pgfpathlineto{\pgfqpoint{5.376871in}{0.714413in}}%
\pgfpathlineto{\pgfqpoint{5.380227in}{0.724575in}}%
\pgfpathlineto{\pgfqpoint{5.383583in}{0.731030in}}%
\pgfpathlineto{\pgfqpoint{5.386939in}{0.734484in}}%
\pgfpathlineto{\pgfqpoint{5.390295in}{0.734920in}}%
\pgfpathlineto{\pgfqpoint{5.393650in}{0.731564in}}%
\pgfpathlineto{\pgfqpoint{5.397006in}{0.722228in}}%
\pgfpathlineto{\pgfqpoint{5.400362in}{0.705525in}}%
\pgfpathlineto{\pgfqpoint{5.407074in}{0.653085in}}%
\pgfpathlineto{\pgfqpoint{5.410429in}{0.659247in}}%
\pgfpathlineto{\pgfqpoint{5.413785in}{0.654609in}}%
\pgfpathlineto{\pgfqpoint{5.417141in}{0.628700in}}%
\pgfpathlineto{\pgfqpoint{5.420497in}{0.665344in}}%
\pgfpathlineto{\pgfqpoint{5.423852in}{0.690502in}}%
\pgfpathlineto{\pgfqpoint{5.427208in}{0.700105in}}%
\pgfpathlineto{\pgfqpoint{5.430564in}{0.694432in}}%
\pgfpathlineto{\pgfqpoint{5.433920in}{0.677237in}}%
\pgfpathlineto{\pgfqpoint{5.437276in}{0.647262in}}%
\pgfpathlineto{\pgfqpoint{5.440631in}{0.610061in}}%
\pgfpathlineto{\pgfqpoint{5.443987in}{0.630918in}}%
\pgfpathlineto{\pgfqpoint{5.447343in}{0.629116in}}%
\pgfpathlineto{\pgfqpoint{5.450699in}{0.588173in}}%
\pgfpathlineto{\pgfqpoint{5.454054in}{0.633373in}}%
\pgfpathlineto{\pgfqpoint{5.457410in}{0.662774in}}%
\pgfpathlineto{\pgfqpoint{5.460766in}{0.673930in}}%
\pgfpathlineto{\pgfqpoint{5.464122in}{0.669144in}}%
\pgfpathlineto{\pgfqpoint{5.467478in}{0.622146in}}%
\pgfpathlineto{\pgfqpoint{5.470833in}{0.660845in}}%
\pgfpathlineto{\pgfqpoint{5.474189in}{0.692364in}}%
\pgfpathlineto{\pgfqpoint{5.477545in}{0.706280in}}%
\pgfpathlineto{\pgfqpoint{5.480901in}{0.712013in}}%
\pgfpathlineto{\pgfqpoint{5.484256in}{0.711514in}}%
\pgfpathlineto{\pgfqpoint{5.487612in}{0.704123in}}%
\pgfpathlineto{\pgfqpoint{5.490968in}{0.684827in}}%
\pgfpathlineto{\pgfqpoint{5.494324in}{0.595944in}}%
\pgfpathlineto{\pgfqpoint{5.497680in}{0.681997in}}%
\pgfpathlineto{\pgfqpoint{5.501035in}{0.708699in}}%
\pgfpathlineto{\pgfqpoint{5.504391in}{0.720422in}}%
\pgfpathlineto{\pgfqpoint{5.507747in}{0.722751in}}%
\pgfpathlineto{\pgfqpoint{5.511103in}{0.713111in}}%
\pgfpathlineto{\pgfqpoint{5.514458in}{0.685328in}}%
\pgfpathlineto{\pgfqpoint{5.517814in}{0.634593in}}%
\pgfpathlineto{\pgfqpoint{5.521170in}{0.690324in}}%
\pgfpathlineto{\pgfqpoint{5.524526in}{0.704189in}}%
\pgfpathlineto{\pgfqpoint{5.527882in}{0.704030in}}%
\pgfpathlineto{\pgfqpoint{5.531237in}{0.689577in}}%
\pgfpathlineto{\pgfqpoint{5.534593in}{0.619549in}}%
\pgfpathlineto{\pgfqpoint{5.537949in}{0.685507in}}%
\pgfpathlineto{\pgfqpoint{5.541305in}{0.713275in}}%
\pgfpathlineto{\pgfqpoint{5.544660in}{0.723206in}}%
\pgfpathlineto{\pgfqpoint{5.548016in}{0.718780in}}%
\pgfpathlineto{\pgfqpoint{5.551372in}{0.696856in}}%
\pgfpathlineto{\pgfqpoint{5.554728in}{0.622411in}}%
\pgfpathlineto{\pgfqpoint{5.558084in}{0.686013in}}%
\pgfpathlineto{\pgfqpoint{5.561439in}{0.705181in}}%
\pgfpathlineto{\pgfqpoint{5.564795in}{0.708224in}}%
\pgfpathlineto{\pgfqpoint{5.568151in}{0.698423in}}%
\pgfpathlineto{\pgfqpoint{5.571507in}{0.659611in}}%
\pgfpathlineto{\pgfqpoint{5.574863in}{0.673132in}}%
\pgfpathlineto{\pgfqpoint{5.578218in}{0.710660in}}%
\pgfpathlineto{\pgfqpoint{5.581574in}{0.724399in}}%
\pgfpathlineto{\pgfqpoint{5.584930in}{0.724383in}}%
\pgfpathlineto{\pgfqpoint{5.588286in}{0.706189in}}%
\pgfpathlineto{\pgfqpoint{5.591641in}{0.636732in}}%
\pgfpathlineto{\pgfqpoint{5.594997in}{0.693974in}}%
\pgfpathlineto{\pgfqpoint{5.598353in}{0.716815in}}%
\pgfpathlineto{\pgfqpoint{5.601709in}{0.723297in}}%
\pgfpathlineto{\pgfqpoint{5.605065in}{0.719289in}}%
\pgfpathlineto{\pgfqpoint{5.608420in}{0.700806in}}%
\pgfpathlineto{\pgfqpoint{5.611776in}{0.551570in}}%
\pgfpathlineto{\pgfqpoint{5.615132in}{0.705267in}}%
\pgfpathlineto{\pgfqpoint{5.618488in}{0.729474in}}%
\pgfpathlineto{\pgfqpoint{5.621843in}{0.737955in}}%
\pgfpathlineto{\pgfqpoint{5.625199in}{0.732550in}}%
\pgfpathlineto{\pgfqpoint{5.628555in}{0.708344in}}%
\pgfpathlineto{\pgfqpoint{5.631911in}{0.630998in}}%
\pgfpathlineto{\pgfqpoint{5.635267in}{0.708196in}}%
\pgfpathlineto{\pgfqpoint{5.638622in}{0.725195in}}%
\pgfpathlineto{\pgfqpoint{5.641978in}{0.727758in}}%
\pgfpathlineto{\pgfqpoint{5.645334in}{0.718237in}}%
\pgfpathlineto{\pgfqpoint{5.648690in}{0.682463in}}%
\pgfpathlineto{\pgfqpoint{5.652045in}{0.687386in}}%
\pgfpathlineto{\pgfqpoint{5.655401in}{0.727586in}}%
\pgfpathlineto{\pgfqpoint{5.658757in}{0.742162in}}%
\pgfpathlineto{\pgfqpoint{5.662113in}{0.743176in}}%
\pgfpathlineto{\pgfqpoint{5.665469in}{0.728488in}}%
\pgfpathlineto{\pgfqpoint{5.668824in}{0.693203in}}%
\pgfpathlineto{\pgfqpoint{5.672180in}{0.664879in}}%
\pgfpathlineto{\pgfqpoint{5.675536in}{0.706392in}}%
\pgfpathlineto{\pgfqpoint{5.678892in}{0.714079in}}%
\pgfpathlineto{\pgfqpoint{5.682247in}{0.705287in}}%
\pgfpathlineto{\pgfqpoint{5.685603in}{0.660349in}}%
\pgfpathlineto{\pgfqpoint{5.692315in}{0.723991in}}%
\pgfpathlineto{\pgfqpoint{5.695671in}{0.735997in}}%
\pgfpathlineto{\pgfqpoint{5.699026in}{0.732745in}}%
\pgfpathlineto{\pgfqpoint{5.702382in}{0.714440in}}%
\pgfpathlineto{\pgfqpoint{5.705738in}{0.672226in}}%
\pgfpathlineto{\pgfqpoint{5.709094in}{0.661954in}}%
\pgfpathlineto{\pgfqpoint{5.712449in}{0.692893in}}%
\pgfpathlineto{\pgfqpoint{5.715805in}{0.694917in}}%
\pgfpathlineto{\pgfqpoint{5.719161in}{0.674416in}}%
\pgfpathlineto{\pgfqpoint{5.722517in}{0.636260in}}%
\pgfpathlineto{\pgfqpoint{5.725873in}{0.699929in}}%
\pgfpathlineto{\pgfqpoint{5.729228in}{0.719290in}}%
\pgfpathlineto{\pgfqpoint{5.732584in}{0.722978in}}%
\pgfpathlineto{\pgfqpoint{5.735940in}{0.709499in}}%
\pgfpathlineto{\pgfqpoint{5.739296in}{0.674492in}}%
\pgfpathlineto{\pgfqpoint{5.742652in}{0.647641in}}%
\pgfpathlineto{\pgfqpoint{5.746007in}{0.690094in}}%
\pgfpathlineto{\pgfqpoint{5.749363in}{0.698951in}}%
\pgfpathlineto{\pgfqpoint{5.752719in}{0.692684in}}%
\pgfpathlineto{\pgfqpoint{5.756075in}{0.660938in}}%
\pgfpathlineto{\pgfqpoint{5.759430in}{0.659013in}}%
\pgfpathlineto{\pgfqpoint{5.762786in}{0.702355in}}%
\pgfpathlineto{\pgfqpoint{5.766142in}{0.717511in}}%
\pgfpathlineto{\pgfqpoint{5.769498in}{0.717980in}}%
\pgfpathlineto{\pgfqpoint{5.772854in}{0.699736in}}%
\pgfpathlineto{\pgfqpoint{5.776209in}{0.627389in}}%
\pgfpathlineto{\pgfqpoint{5.779565in}{0.685938in}}%
\pgfpathlineto{\pgfqpoint{5.782921in}{0.708919in}}%
\pgfpathlineto{\pgfqpoint{5.786277in}{0.715010in}}%
\pgfpathlineto{\pgfqpoint{5.789632in}{0.710112in}}%
\pgfpathlineto{\pgfqpoint{5.792988in}{0.688972in}}%
\pgfpathlineto{\pgfqpoint{5.796344in}{0.614837in}}%
\pgfpathlineto{\pgfqpoint{5.799700in}{0.702189in}}%
\pgfpathlineto{\pgfqpoint{5.803056in}{0.723649in}}%
\pgfpathlineto{\pgfqpoint{5.806411in}{0.730236in}}%
\pgfpathlineto{\pgfqpoint{5.809767in}{0.721719in}}%
\pgfpathlineto{\pgfqpoint{5.813123in}{0.688670in}}%
\pgfpathlineto{\pgfqpoint{5.816479in}{0.669855in}}%
\pgfpathlineto{\pgfqpoint{5.819834in}{0.713483in}}%
\pgfpathlineto{\pgfqpoint{5.823190in}{0.726172in}}%
\pgfpathlineto{\pgfqpoint{5.826546in}{0.727026in}}%
\pgfpathlineto{\pgfqpoint{5.829902in}{0.716443in}}%
\pgfpathlineto{\pgfqpoint{5.833258in}{0.678992in}}%
\pgfpathlineto{\pgfqpoint{5.836613in}{0.686907in}}%
\pgfpathlineto{\pgfqpoint{5.839969in}{0.725913in}}%
\pgfpathlineto{\pgfqpoint{5.843325in}{0.740236in}}%
\pgfpathlineto{\pgfqpoint{5.846681in}{0.741367in}}%
\pgfpathlineto{\pgfqpoint{5.850036in}{0.726437in}}%
\pgfpathlineto{\pgfqpoint{5.853392in}{0.685575in}}%
\pgfpathlineto{\pgfqpoint{5.856748in}{0.680451in}}%
\pgfpathlineto{\pgfqpoint{5.860104in}{0.714607in}}%
\pgfpathlineto{\pgfqpoint{5.863460in}{0.722110in}}%
\pgfpathlineto{\pgfqpoint{5.866815in}{0.715561in}}%
\pgfpathlineto{\pgfqpoint{5.870171in}{0.684919in}}%
\pgfpathlineto{\pgfqpoint{5.873527in}{0.677997in}}%
\pgfpathlineto{\pgfqpoint{5.876883in}{0.723214in}}%
\pgfpathlineto{\pgfqpoint{5.880238in}{0.738856in}}%
\pgfpathlineto{\pgfqpoint{5.883594in}{0.739982in}}%
\pgfpathlineto{\pgfqpoint{5.886950in}{0.725268in}}%
\pgfpathlineto{\pgfqpoint{5.890306in}{0.692944in}}%
\pgfpathlineto{\pgfqpoint{5.893662in}{0.635633in}}%
\pgfpathlineto{\pgfqpoint{5.897017in}{0.692390in}}%
\pgfpathlineto{\pgfqpoint{5.900373in}{0.699737in}}%
\pgfpathlineto{\pgfqpoint{5.903729in}{0.686116in}}%
\pgfpathlineto{\pgfqpoint{5.907085in}{0.545809in}}%
\pgfpathlineto{\pgfqpoint{5.910441in}{0.696408in}}%
\pgfpathlineto{\pgfqpoint{5.913796in}{0.720546in}}%
\pgfpathlineto{\pgfqpoint{5.917152in}{0.727531in}}%
\pgfpathlineto{\pgfqpoint{5.920508in}{0.717789in}}%
\pgfpathlineto{\pgfqpoint{5.923864in}{0.692026in}}%
\pgfpathlineto{\pgfqpoint{5.927219in}{0.585188in}}%
\pgfpathlineto{\pgfqpoint{5.930575in}{0.677094in}}%
\pgfpathlineto{\pgfqpoint{5.933931in}{0.691477in}}%
\pgfpathlineto{\pgfqpoint{5.937287in}{0.686370in}}%
\pgfpathlineto{\pgfqpoint{5.940643in}{0.651130in}}%
\pgfpathlineto{\pgfqpoint{5.943998in}{0.662685in}}%
\pgfpathlineto{\pgfqpoint{5.947354in}{0.701754in}}%
\pgfpathlineto{\pgfqpoint{5.950710in}{0.715553in}}%
\pgfpathlineto{\pgfqpoint{5.954066in}{0.713839in}}%
\pgfpathlineto{\pgfqpoint{5.957421in}{0.693277in}}%
\pgfpathlineto{\pgfqpoint{5.960777in}{0.607537in}}%
\pgfpathlineto{\pgfqpoint{5.964133in}{0.680665in}}%
\pgfpathlineto{\pgfqpoint{5.967489in}{0.701419in}}%
\pgfpathlineto{\pgfqpoint{5.970845in}{0.705419in}}%
\pgfpathlineto{\pgfqpoint{5.974200in}{0.696998in}}%
\pgfpathlineto{\pgfqpoint{5.977556in}{0.663906in}}%
\pgfpathlineto{\pgfqpoint{5.980912in}{0.661961in}}%
\pgfpathlineto{\pgfqpoint{5.984268in}{0.704771in}}%
\pgfpathlineto{\pgfqpoint{5.987623in}{0.719907in}}%
\pgfpathlineto{\pgfqpoint{5.990979in}{0.721073in}}%
\pgfpathlineto{\pgfqpoint{5.994335in}{0.703672in}}%
\pgfpathlineto{\pgfqpoint{5.997691in}{0.622906in}}%
\pgfpathlineto{\pgfqpoint{6.001047in}{0.693848in}}%
\pgfpathlineto{\pgfqpoint{6.004402in}{0.716977in}}%
\pgfpathlineto{\pgfqpoint{6.007758in}{0.723935in}}%
\pgfpathlineto{\pgfqpoint{6.011114in}{0.720754in}}%
\pgfpathlineto{\pgfqpoint{6.014470in}{0.704379in}}%
\pgfpathlineto{\pgfqpoint{6.017825in}{0.627409in}}%
\pgfpathlineto{\pgfqpoint{6.021181in}{0.700614in}}%
\pgfpathlineto{\pgfqpoint{6.024537in}{0.727599in}}%
\pgfpathlineto{\pgfqpoint{6.027893in}{0.737496in}}%
\pgfpathlineto{\pgfqpoint{6.031249in}{0.733948in}}%
\pgfpathlineto{\pgfqpoint{6.034604in}{0.711959in}}%
\pgfpathlineto{\pgfqpoint{6.037960in}{0.584584in}}%
\pgfpathlineto{\pgfqpoint{6.041316in}{0.706114in}}%
\pgfpathlineto{\pgfqpoint{6.044672in}{0.725158in}}%
\pgfpathlineto{\pgfqpoint{6.048027in}{0.728731in}}%
\pgfpathlineto{\pgfqpoint{6.051383in}{0.720330in}}%
\pgfpathlineto{\pgfqpoint{6.054739in}{0.688320in}}%
\pgfpathlineto{\pgfqpoint{6.058095in}{0.682281in}}%
\pgfpathlineto{\pgfqpoint{6.061451in}{0.726679in}}%
\pgfpathlineto{\pgfqpoint{6.064806in}{0.742272in}}%
\pgfpathlineto{\pgfqpoint{6.068162in}{0.744043in}}%
\pgfpathlineto{\pgfqpoint{6.071518in}{0.730117in}}%
\pgfpathlineto{\pgfqpoint{6.074874in}{0.697554in}}%
\pgfpathlineto{\pgfqpoint{6.078230in}{0.647294in}}%
\pgfpathlineto{\pgfqpoint{6.081585in}{0.701373in}}%
\pgfpathlineto{\pgfqpoint{6.084941in}{0.709804in}}%
\pgfpathlineto{\pgfqpoint{6.088297in}{0.699830in}}%
\pgfpathlineto{\pgfqpoint{6.091653in}{0.641270in}}%
\pgfpathlineto{\pgfqpoint{6.095008in}{0.694553in}}%
\pgfpathlineto{\pgfqpoint{6.098364in}{0.724245in}}%
\pgfpathlineto{\pgfqpoint{6.101720in}{0.734323in}}%
\pgfpathlineto{\pgfqpoint{6.105076in}{0.728352in}}%
\pgfpathlineto{\pgfqpoint{6.108432in}{0.707996in}}%
\pgfpathlineto{\pgfqpoint{6.111787in}{0.659116in}}%
\pgfpathlineto{\pgfqpoint{6.115143in}{0.661859in}}%
\pgfpathlineto{\pgfqpoint{6.118499in}{0.687406in}}%
\pgfpathlineto{\pgfqpoint{6.121855in}{0.685231in}}%
\pgfpathlineto{\pgfqpoint{6.125210in}{0.650978in}}%
\pgfpathlineto{\pgfqpoint{6.128566in}{0.664299in}}%
\pgfpathlineto{\pgfqpoint{6.131922in}{0.703506in}}%
\pgfpathlineto{\pgfqpoint{6.135278in}{0.717398in}}%
\pgfpathlineto{\pgfqpoint{6.138634in}{0.715355in}}%
\pgfpathlineto{\pgfqpoint{6.141989in}{0.695781in}}%
\pgfpathlineto{\pgfqpoint{6.145345in}{0.635926in}}%
\pgfpathlineto{\pgfqpoint{6.148701in}{0.668748in}}%
\pgfpathlineto{\pgfqpoint{6.152057in}{0.692235in}}%
\pgfpathlineto{\pgfqpoint{6.155412in}{0.695350in}}%
\pgfpathlineto{\pgfqpoint{6.158768in}{0.683029in}}%
\pgfpathlineto{\pgfqpoint{6.162124in}{0.621592in}}%
\pgfpathlineto{\pgfqpoint{6.165480in}{0.676757in}}%
\pgfpathlineto{\pgfqpoint{6.168836in}{0.705804in}}%
\pgfpathlineto{\pgfqpoint{6.172191in}{0.715784in}}%
\pgfpathlineto{\pgfqpoint{6.175547in}{0.710145in}}%
\pgfpathlineto{\pgfqpoint{6.178903in}{0.681264in}}%
\pgfpathlineto{\pgfqpoint{6.182259in}{0.646124in}}%
\pgfpathlineto{\pgfqpoint{6.185614in}{0.699132in}}%
\pgfpathlineto{\pgfqpoint{6.188970in}{0.713759in}}%
\pgfpathlineto{\pgfqpoint{6.192326in}{0.716139in}}%
\pgfpathlineto{\pgfqpoint{6.195682in}{0.707868in}}%
\pgfpathlineto{\pgfqpoint{6.199038in}{0.679151in}}%
\pgfpathlineto{\pgfqpoint{6.202393in}{0.658236in}}%
\pgfpathlineto{\pgfqpoint{6.205749in}{0.709334in}}%
\pgfpathlineto{\pgfqpoint{6.209105in}{0.726597in}}%
\pgfpathlineto{\pgfqpoint{6.212461in}{0.730409in}}%
\pgfpathlineto{\pgfqpoint{6.215816in}{0.718075in}}%
\pgfpathlineto{\pgfqpoint{6.219172in}{0.671479in}}%
\pgfpathlineto{\pgfqpoint{6.222528in}{0.689841in}}%
\pgfpathlineto{\pgfqpoint{6.225884in}{0.720850in}}%
\pgfpathlineto{\pgfqpoint{6.229240in}{0.730609in}}%
\pgfpathlineto{\pgfqpoint{6.232595in}{0.729682in}}%
\pgfpathlineto{\pgfqpoint{6.235951in}{0.716951in}}%
\pgfpathlineto{\pgfqpoint{6.239307in}{0.670371in}}%
\pgfpathlineto{\pgfqpoint{6.242663in}{0.697846in}}%
\pgfpathlineto{\pgfqpoint{6.246019in}{0.731383in}}%
\pgfpathlineto{\pgfqpoint{6.249374in}{0.743968in}}%
\pgfpathlineto{\pgfqpoint{6.252730in}{0.743494in}}%
\pgfpathlineto{\pgfqpoint{6.256086in}{0.726962in}}%
\pgfpathlineto{\pgfqpoint{6.259442in}{0.682048in}}%
\pgfpathlineto{\pgfqpoint{6.262797in}{0.685551in}}%
\pgfpathlineto{\pgfqpoint{6.266153in}{0.716270in}}%
\pgfpathlineto{\pgfqpoint{6.269509in}{0.722096in}}%
\pgfpathlineto{\pgfqpoint{6.272865in}{0.713037in}}%
\pgfpathlineto{\pgfqpoint{6.276221in}{0.671386in}}%
\pgfpathlineto{\pgfqpoint{6.279576in}{0.693286in}}%
\pgfpathlineto{\pgfqpoint{6.282932in}{0.728671in}}%
\pgfpathlineto{\pgfqpoint{6.286288in}{0.741274in}}%
\pgfpathlineto{\pgfqpoint{6.289644in}{0.738879in}}%
\pgfpathlineto{\pgfqpoint{6.292999in}{0.721718in}}%
\pgfpathlineto{\pgfqpoint{6.299711in}{0.644930in}}%
\pgfpathlineto{\pgfqpoint{6.303067in}{0.689119in}}%
\pgfpathlineto{\pgfqpoint{6.306423in}{0.691618in}}%
\pgfpathlineto{\pgfqpoint{6.309778in}{0.665907in}}%
\pgfpathlineto{\pgfqpoint{6.313134in}{0.658040in}}%
\pgfpathlineto{\pgfqpoint{6.316490in}{0.705889in}}%
\pgfpathlineto{\pgfqpoint{6.319846in}{0.722321in}}%
\pgfpathlineto{\pgfqpoint{6.323201in}{0.722948in}}%
\pgfpathlineto{\pgfqpoint{6.326557in}{0.706848in}}%
\pgfpathlineto{\pgfqpoint{6.329913in}{0.669986in}}%
\pgfpathlineto{\pgfqpoint{6.333269in}{0.639930in}}%
\pgfpathlineto{\pgfqpoint{6.336625in}{0.681406in}}%
\pgfpathlineto{\pgfqpoint{6.339980in}{0.686742in}}%
\pgfpathlineto{\pgfqpoint{6.343336in}{0.672093in}}%
\pgfpathlineto{\pgfqpoint{6.346692in}{0.535297in}}%
\pgfpathlineto{\pgfqpoint{6.350048in}{0.682737in}}%
\pgfpathlineto{\pgfqpoint{6.353403in}{0.706418in}}%
\pgfpathlineto{\pgfqpoint{6.356759in}{0.713118in}}%
\pgfpathlineto{\pgfqpoint{6.360115in}{0.704247in}}%
\pgfpathlineto{\pgfqpoint{6.366827in}{0.674755in}}%
\pgfpathlineto{\pgfqpoint{6.370182in}{0.659920in}}%
\pgfpathlineto{\pgfqpoint{6.373538in}{0.650719in}}%
\pgfpathlineto{\pgfqpoint{6.376894in}{0.649632in}}%
\pgfpathlineto{\pgfqpoint{6.380250in}{0.651686in}}%
\pgfpathlineto{\pgfqpoint{6.383605in}{0.647765in}}%
\pgfpathlineto{\pgfqpoint{6.383605in}{0.647765in}}%
\pgfusepath{stroke}%
\end{pgfscope}%
\begin{pgfscope}%
\pgfpathrectangle{\pgfqpoint{4.621821in}{0.320625in}}{\pgfqpoint{1.845679in}{1.164375in}}%
\pgfusepath{clip}%
\pgfsetrectcap%
\pgfsetroundjoin%
\pgfsetlinewidth{0.602250pt}%
\definecolor{currentstroke}{rgb}{0.819608,0.286275,0.356863}%
\pgfsetstrokecolor{currentstroke}%
\pgfsetstrokeopacity{0.120000}%
\pgfsetdash{}{0pt}%
\pgfpathmoveto{\pgfqpoint{4.705715in}{1.472285in}}%
\pgfpathlineto{\pgfqpoint{4.732562in}{1.424217in}}%
\pgfpathlineto{\pgfqpoint{4.769475in}{1.354185in}}%
\pgfpathlineto{\pgfqpoint{4.786254in}{1.325127in}}%
\pgfpathlineto{\pgfqpoint{4.803033in}{1.298083in}}%
\pgfpathlineto{\pgfqpoint{4.823168in}{1.267905in}}%
\pgfpathlineto{\pgfqpoint{4.843302in}{1.239770in}}%
\pgfpathlineto{\pgfqpoint{4.866793in}{1.209015in}}%
\pgfpathlineto{\pgfqpoint{4.890283in}{1.180023in}}%
\pgfpathlineto{\pgfqpoint{4.920485in}{1.144678in}}%
\pgfpathlineto{\pgfqpoint{4.954043in}{1.107162in}}%
\pgfpathlineto{\pgfqpoint{5.001024in}{1.056484in}}%
\pgfpathlineto{\pgfqpoint{5.068140in}{0.985916in}}%
\pgfpathlineto{\pgfqpoint{5.199015in}{0.850392in}}%
\pgfpathlineto{\pgfqpoint{5.269487in}{0.779158in}}%
\pgfpathlineto{\pgfqpoint{5.279554in}{0.769200in}}%
\pgfpathlineto{\pgfqpoint{5.292977in}{0.754016in}}%
\pgfpathlineto{\pgfqpoint{5.309756in}{0.734630in}}%
\pgfpathlineto{\pgfqpoint{5.316467in}{0.727959in}}%
\pgfpathlineto{\pgfqpoint{5.323179in}{0.722494in}}%
\pgfpathlineto{\pgfqpoint{5.329891in}{0.718274in}}%
\pgfpathlineto{\pgfqpoint{5.343314in}{0.710919in}}%
\pgfpathlineto{\pgfqpoint{5.346669in}{0.708226in}}%
\pgfpathlineto{\pgfqpoint{5.353381in}{0.701229in}}%
\pgfpathlineto{\pgfqpoint{5.356737in}{0.699461in}}%
\pgfpathlineto{\pgfqpoint{5.360093in}{0.701267in}}%
\pgfpathlineto{\pgfqpoint{5.363448in}{0.706223in}}%
\pgfpathlineto{\pgfqpoint{5.370160in}{0.717803in}}%
\pgfpathlineto{\pgfqpoint{5.373516in}{0.722257in}}%
\pgfpathlineto{\pgfqpoint{5.376871in}{0.725430in}}%
\pgfpathlineto{\pgfqpoint{5.380227in}{0.727282in}}%
\pgfpathlineto{\pgfqpoint{5.383583in}{0.727784in}}%
\pgfpathlineto{\pgfqpoint{5.386939in}{0.726854in}}%
\pgfpathlineto{\pgfqpoint{5.390295in}{0.724310in}}%
\pgfpathlineto{\pgfqpoint{5.393650in}{0.719789in}}%
\pgfpathlineto{\pgfqpoint{5.397006in}{0.712575in}}%
\pgfpathlineto{\pgfqpoint{5.400362in}{0.701087in}}%
\pgfpathlineto{\pgfqpoint{5.403718in}{0.680907in}}%
\pgfpathlineto{\pgfqpoint{5.407074in}{0.634232in}}%
\pgfpathlineto{\pgfqpoint{5.410429in}{0.670074in}}%
\pgfpathlineto{\pgfqpoint{5.413785in}{0.699052in}}%
\pgfpathlineto{\pgfqpoint{5.417141in}{0.715258in}}%
\pgfpathlineto{\pgfqpoint{5.420497in}{0.725975in}}%
\pgfpathlineto{\pgfqpoint{5.423852in}{0.733390in}}%
\pgfpathlineto{\pgfqpoint{5.427208in}{0.738359in}}%
\pgfpathlineto{\pgfqpoint{5.430564in}{0.741194in}}%
\pgfpathlineto{\pgfqpoint{5.433920in}{0.741859in}}%
\pgfpathlineto{\pgfqpoint{5.437276in}{0.739944in}}%
\pgfpathlineto{\pgfqpoint{5.440631in}{0.734360in}}%
\pgfpathlineto{\pgfqpoint{5.443987in}{0.723071in}}%
\pgfpathlineto{\pgfqpoint{5.447343in}{0.706114in}}%
\pgfpathlineto{\pgfqpoint{5.454054in}{0.659079in}}%
\pgfpathlineto{\pgfqpoint{5.457410in}{0.648569in}}%
\pgfpathlineto{\pgfqpoint{5.460766in}{0.649398in}}%
\pgfpathlineto{\pgfqpoint{5.467478in}{0.685206in}}%
\pgfpathlineto{\pgfqpoint{5.470833in}{0.692053in}}%
\pgfpathlineto{\pgfqpoint{5.474189in}{0.684024in}}%
\pgfpathlineto{\pgfqpoint{5.477545in}{0.668159in}}%
\pgfpathlineto{\pgfqpoint{5.484256in}{0.624212in}}%
\pgfpathlineto{\pgfqpoint{5.487612in}{0.614106in}}%
\pgfpathlineto{\pgfqpoint{5.490968in}{0.612250in}}%
\pgfpathlineto{\pgfqpoint{5.494324in}{0.625087in}}%
\pgfpathlineto{\pgfqpoint{5.497680in}{0.642481in}}%
\pgfpathlineto{\pgfqpoint{5.501035in}{0.651107in}}%
\pgfpathlineto{\pgfqpoint{5.504391in}{0.644094in}}%
\pgfpathlineto{\pgfqpoint{5.507747in}{0.593606in}}%
\pgfpathlineto{\pgfqpoint{5.511103in}{0.659692in}}%
\pgfpathlineto{\pgfqpoint{5.514458in}{0.685311in}}%
\pgfpathlineto{\pgfqpoint{5.517814in}{0.698060in}}%
\pgfpathlineto{\pgfqpoint{5.521170in}{0.703877in}}%
\pgfpathlineto{\pgfqpoint{5.524526in}{0.704295in}}%
\pgfpathlineto{\pgfqpoint{5.527882in}{0.699020in}}%
\pgfpathlineto{\pgfqpoint{5.531237in}{0.685245in}}%
\pgfpathlineto{\pgfqpoint{5.534593in}{0.648066in}}%
\pgfpathlineto{\pgfqpoint{5.537949in}{0.648984in}}%
\pgfpathlineto{\pgfqpoint{5.541305in}{0.690222in}}%
\pgfpathlineto{\pgfqpoint{5.544660in}{0.707126in}}%
\pgfpathlineto{\pgfqpoint{5.548016in}{0.714348in}}%
\pgfpathlineto{\pgfqpoint{5.551372in}{0.713008in}}%
\pgfpathlineto{\pgfqpoint{5.554728in}{0.698437in}}%
\pgfpathlineto{\pgfqpoint{5.558084in}{0.653357in}}%
\pgfpathlineto{\pgfqpoint{5.561439in}{0.665806in}}%
\pgfpathlineto{\pgfqpoint{5.564795in}{0.696333in}}%
\pgfpathlineto{\pgfqpoint{5.568151in}{0.705024in}}%
\pgfpathlineto{\pgfqpoint{5.571507in}{0.702247in}}%
\pgfpathlineto{\pgfqpoint{5.574863in}{0.685074in}}%
\pgfpathlineto{\pgfqpoint{5.578218in}{0.580664in}}%
\pgfpathlineto{\pgfqpoint{5.581574in}{0.687008in}}%
\pgfpathlineto{\pgfqpoint{5.584930in}{0.712207in}}%
\pgfpathlineto{\pgfqpoint{5.588286in}{0.721047in}}%
\pgfpathlineto{\pgfqpoint{5.591641in}{0.715728in}}%
\pgfpathlineto{\pgfqpoint{5.594997in}{0.691029in}}%
\pgfpathlineto{\pgfqpoint{5.598353in}{0.621503in}}%
\pgfpathlineto{\pgfqpoint{5.601709in}{0.693927in}}%
\pgfpathlineto{\pgfqpoint{5.605065in}{0.710662in}}%
\pgfpathlineto{\pgfqpoint{5.608420in}{0.713499in}}%
\pgfpathlineto{\pgfqpoint{5.611776in}{0.704854in}}%
\pgfpathlineto{\pgfqpoint{5.615132in}{0.673269in}}%
\pgfpathlineto{\pgfqpoint{5.618488in}{0.664972in}}%
\pgfpathlineto{\pgfqpoint{5.621843in}{0.710304in}}%
\pgfpathlineto{\pgfqpoint{5.625199in}{0.726318in}}%
\pgfpathlineto{\pgfqpoint{5.628555in}{0.728909in}}%
\pgfpathlineto{\pgfqpoint{5.631911in}{0.714594in}}%
\pgfpathlineto{\pgfqpoint{5.635267in}{0.661830in}}%
\pgfpathlineto{\pgfqpoint{5.641978in}{0.720768in}}%
\pgfpathlineto{\pgfqpoint{5.645334in}{0.729539in}}%
\pgfpathlineto{\pgfqpoint{5.648690in}{0.728000in}}%
\pgfpathlineto{\pgfqpoint{5.652045in}{0.714556in}}%
\pgfpathlineto{\pgfqpoint{5.655401in}{0.664517in}}%
\pgfpathlineto{\pgfqpoint{5.662113in}{0.730500in}}%
\pgfpathlineto{\pgfqpoint{5.665469in}{0.742751in}}%
\pgfpathlineto{\pgfqpoint{5.668824in}{0.742204in}}%
\pgfpathlineto{\pgfqpoint{5.672180in}{0.725266in}}%
\pgfpathlineto{\pgfqpoint{5.675536in}{0.674962in}}%
\pgfpathlineto{\pgfqpoint{5.678892in}{0.695834in}}%
\pgfpathlineto{\pgfqpoint{5.682247in}{0.722650in}}%
\pgfpathlineto{\pgfqpoint{5.685603in}{0.728847in}}%
\pgfpathlineto{\pgfqpoint{5.688959in}{0.722308in}}%
\pgfpathlineto{\pgfqpoint{5.692315in}{0.694147in}}%
\pgfpathlineto{\pgfqpoint{5.695671in}{0.676025in}}%
\pgfpathlineto{\pgfqpoint{5.699026in}{0.726232in}}%
\pgfpathlineto{\pgfqpoint{5.702382in}{0.743138in}}%
\pgfpathlineto{\pgfqpoint{5.705738in}{0.745871in}}%
\pgfpathlineto{\pgfqpoint{5.709094in}{0.733039in}}%
\pgfpathlineto{\pgfqpoint{5.712449in}{0.704479in}}%
\pgfpathlineto{\pgfqpoint{5.715805in}{0.632247in}}%
\pgfpathlineto{\pgfqpoint{5.719161in}{0.694995in}}%
\pgfpathlineto{\pgfqpoint{5.722517in}{0.705295in}}%
\pgfpathlineto{\pgfqpoint{5.725873in}{0.695225in}}%
\pgfpathlineto{\pgfqpoint{5.729228in}{0.627417in}}%
\pgfpathlineto{\pgfqpoint{5.732584in}{0.695145in}}%
\pgfpathlineto{\pgfqpoint{5.735940in}{0.723245in}}%
\pgfpathlineto{\pgfqpoint{5.739296in}{0.732578in}}%
\pgfpathlineto{\pgfqpoint{5.742652in}{0.725645in}}%
\pgfpathlineto{\pgfqpoint{5.746007in}{0.704280in}}%
\pgfpathlineto{\pgfqpoint{5.749363in}{0.651390in}}%
\pgfpathlineto{\pgfqpoint{5.752719in}{0.667045in}}%
\pgfpathlineto{\pgfqpoint{5.756075in}{0.688365in}}%
\pgfpathlineto{\pgfqpoint{5.759430in}{0.685319in}}%
\pgfpathlineto{\pgfqpoint{5.762786in}{0.651045in}}%
\pgfpathlineto{\pgfqpoint{5.766142in}{0.663517in}}%
\pgfpathlineto{\pgfqpoint{5.769498in}{0.702699in}}%
\pgfpathlineto{\pgfqpoint{5.772854in}{0.716606in}}%
\pgfpathlineto{\pgfqpoint{5.776209in}{0.714734in}}%
\pgfpathlineto{\pgfqpoint{5.779565in}{0.694729in}}%
\pgfpathlineto{\pgfqpoint{5.782921in}{0.628869in}}%
\pgfpathlineto{\pgfqpoint{5.786277in}{0.675319in}}%
\pgfpathlineto{\pgfqpoint{5.789632in}{0.696992in}}%
\pgfpathlineto{\pgfqpoint{5.792988in}{0.700466in}}%
\pgfpathlineto{\pgfqpoint{5.796344in}{0.690133in}}%
\pgfpathlineto{\pgfqpoint{5.799700in}{0.646293in}}%
\pgfpathlineto{\pgfqpoint{5.803056in}{0.671392in}}%
\pgfpathlineto{\pgfqpoint{5.806411in}{0.705763in}}%
\pgfpathlineto{\pgfqpoint{5.809767in}{0.718109in}}%
\pgfpathlineto{\pgfqpoint{5.813123in}{0.715823in}}%
\pgfpathlineto{\pgfqpoint{5.816479in}{0.692585in}}%
\pgfpathlineto{\pgfqpoint{5.819834in}{0.624287in}}%
\pgfpathlineto{\pgfqpoint{5.823190in}{0.699131in}}%
\pgfpathlineto{\pgfqpoint{5.826546in}{0.716879in}}%
\pgfpathlineto{\pgfqpoint{5.829902in}{0.721274in}}%
\pgfpathlineto{\pgfqpoint{5.833258in}{0.715551in}}%
\pgfpathlineto{\pgfqpoint{5.836613in}{0.693926in}}%
\pgfpathlineto{\pgfqpoint{5.839969in}{0.620019in}}%
\pgfpathlineto{\pgfqpoint{5.843325in}{0.706862in}}%
\pgfpathlineto{\pgfqpoint{5.846681in}{0.728321in}}%
\pgfpathlineto{\pgfqpoint{5.850036in}{0.735165in}}%
\pgfpathlineto{\pgfqpoint{5.853392in}{0.727479in}}%
\pgfpathlineto{\pgfqpoint{5.860104in}{0.668295in}}%
\pgfpathlineto{\pgfqpoint{5.863460in}{0.716268in}}%
\pgfpathlineto{\pgfqpoint{5.866815in}{0.729597in}}%
\pgfpathlineto{\pgfqpoint{5.870171in}{0.730587in}}%
\pgfpathlineto{\pgfqpoint{5.873527in}{0.719796in}}%
\pgfpathlineto{\pgfqpoint{5.876883in}{0.680585in}}%
\pgfpathlineto{\pgfqpoint{5.880238in}{0.693186in}}%
\pgfpathlineto{\pgfqpoint{5.883594in}{0.730696in}}%
\pgfpathlineto{\pgfqpoint{5.886950in}{0.744435in}}%
\pgfpathlineto{\pgfqpoint{5.890306in}{0.744700in}}%
\pgfpathlineto{\pgfqpoint{5.893662in}{0.729149in}}%
\pgfpathlineto{\pgfqpoint{5.897017in}{0.691159in}}%
\pgfpathlineto{\pgfqpoint{5.900373in}{0.671763in}}%
\pgfpathlineto{\pgfqpoint{5.903729in}{0.709159in}}%
\pgfpathlineto{\pgfqpoint{5.907085in}{0.715306in}}%
\pgfpathlineto{\pgfqpoint{5.910441in}{0.704203in}}%
\pgfpathlineto{\pgfqpoint{5.913796in}{0.642443in}}%
\pgfpathlineto{\pgfqpoint{5.917152in}{0.699563in}}%
\pgfpathlineto{\pgfqpoint{5.920508in}{0.728517in}}%
\pgfpathlineto{\pgfqpoint{5.923864in}{0.738197in}}%
\pgfpathlineto{\pgfqpoint{5.927219in}{0.731830in}}%
\pgfpathlineto{\pgfqpoint{5.930575in}{0.711555in}}%
\pgfpathlineto{\pgfqpoint{5.933931in}{0.665349in}}%
\pgfpathlineto{\pgfqpoint{5.937287in}{0.660209in}}%
\pgfpathlineto{\pgfqpoint{5.940643in}{0.686469in}}%
\pgfpathlineto{\pgfqpoint{5.943998in}{0.682507in}}%
\pgfpathlineto{\pgfqpoint{5.947354in}{0.635048in}}%
\pgfpathlineto{\pgfqpoint{5.950710in}{0.675863in}}%
\pgfpathlineto{\pgfqpoint{5.954066in}{0.708380in}}%
\pgfpathlineto{\pgfqpoint{5.957421in}{0.719643in}}%
\pgfpathlineto{\pgfqpoint{5.960777in}{0.714195in}}%
\pgfpathlineto{\pgfqpoint{5.964133in}{0.692116in}}%
\pgfpathlineto{\pgfqpoint{5.967489in}{0.622704in}}%
\pgfpathlineto{\pgfqpoint{5.970845in}{0.668196in}}%
\pgfpathlineto{\pgfqpoint{5.974200in}{0.687628in}}%
\pgfpathlineto{\pgfqpoint{5.977556in}{0.686796in}}%
\pgfpathlineto{\pgfqpoint{5.980912in}{0.664882in}}%
\pgfpathlineto{\pgfqpoint{5.984268in}{0.626615in}}%
\pgfpathlineto{\pgfqpoint{5.987623in}{0.689512in}}%
\pgfpathlineto{\pgfqpoint{5.990979in}{0.708429in}}%
\pgfpathlineto{\pgfqpoint{5.994335in}{0.711684in}}%
\pgfpathlineto{\pgfqpoint{5.997691in}{0.696229in}}%
\pgfpathlineto{\pgfqpoint{6.001047in}{0.637106in}}%
\pgfpathlineto{\pgfqpoint{6.004402in}{0.676645in}}%
\pgfpathlineto{\pgfqpoint{6.007758in}{0.702624in}}%
\pgfpathlineto{\pgfqpoint{6.011114in}{0.710107in}}%
\pgfpathlineto{\pgfqpoint{6.014470in}{0.706674in}}%
\pgfpathlineto{\pgfqpoint{6.017825in}{0.688774in}}%
\pgfpathlineto{\pgfqpoint{6.021181in}{0.562724in}}%
\pgfpathlineto{\pgfqpoint{6.024537in}{0.692073in}}%
\pgfpathlineto{\pgfqpoint{6.027893in}{0.716549in}}%
\pgfpathlineto{\pgfqpoint{6.031249in}{0.724838in}}%
\pgfpathlineto{\pgfqpoint{6.034604in}{0.718464in}}%
\pgfpathlineto{\pgfqpoint{6.037960in}{0.686624in}}%
\pgfpathlineto{\pgfqpoint{6.041316in}{0.669231in}}%
\pgfpathlineto{\pgfqpoint{6.044672in}{0.713745in}}%
\pgfpathlineto{\pgfqpoint{6.048027in}{0.727438in}}%
\pgfpathlineto{\pgfqpoint{6.051383in}{0.729730in}}%
\pgfpathlineto{\pgfqpoint{6.054739in}{0.721925in}}%
\pgfpathlineto{\pgfqpoint{6.058095in}{0.695525in}}%
\pgfpathlineto{\pgfqpoint{6.061451in}{0.661855in}}%
\pgfpathlineto{\pgfqpoint{6.064806in}{0.719913in}}%
\pgfpathlineto{\pgfqpoint{6.068162in}{0.738417in}}%
\pgfpathlineto{\pgfqpoint{6.071518in}{0.743342in}}%
\pgfpathlineto{\pgfqpoint{6.074874in}{0.733350in}}%
\pgfpathlineto{\pgfqpoint{6.078230in}{0.702281in}}%
\pgfpathlineto{\pgfqpoint{6.081585in}{0.662367in}}%
\pgfpathlineto{\pgfqpoint{6.084941in}{0.714446in}}%
\pgfpathlineto{\pgfqpoint{6.088297in}{0.726320in}}%
\pgfpathlineto{\pgfqpoint{6.091653in}{0.724073in}}%
\pgfpathlineto{\pgfqpoint{6.095008in}{0.704910in}}%
\pgfpathlineto{\pgfqpoint{6.098364in}{0.618191in}}%
\pgfpathlineto{\pgfqpoint{6.101720in}{0.716938in}}%
\pgfpathlineto{\pgfqpoint{6.105076in}{0.738881in}}%
\pgfpathlineto{\pgfqpoint{6.108432in}{0.744975in}}%
\pgfpathlineto{\pgfqpoint{6.111787in}{0.735329in}}%
\pgfpathlineto{\pgfqpoint{6.115143in}{0.711757in}}%
\pgfpathlineto{\pgfqpoint{6.118499in}{0.646573in}}%
\pgfpathlineto{\pgfqpoint{6.121855in}{0.678794in}}%
\pgfpathlineto{\pgfqpoint{6.125210in}{0.695947in}}%
\pgfpathlineto{\pgfqpoint{6.128566in}{0.687433in}}%
\pgfpathlineto{\pgfqpoint{6.131922in}{0.610394in}}%
\pgfpathlineto{\pgfqpoint{6.135278in}{0.691392in}}%
\pgfpathlineto{\pgfqpoint{6.138634in}{0.718796in}}%
\pgfpathlineto{\pgfqpoint{6.141989in}{0.727455in}}%
\pgfpathlineto{\pgfqpoint{6.145345in}{0.719147in}}%
\pgfpathlineto{\pgfqpoint{6.148701in}{0.696436in}}%
\pgfpathlineto{\pgfqpoint{6.152057in}{0.635727in}}%
\pgfpathlineto{\pgfqpoint{6.155412in}{0.661112in}}%
\pgfpathlineto{\pgfqpoint{6.158768in}{0.680124in}}%
\pgfpathlineto{\pgfqpoint{6.162124in}{0.673588in}}%
\pgfpathlineto{\pgfqpoint{6.165480in}{0.617585in}}%
\pgfpathlineto{\pgfqpoint{6.168836in}{0.670333in}}%
\pgfpathlineto{\pgfqpoint{6.172191in}{0.700659in}}%
\pgfpathlineto{\pgfqpoint{6.175547in}{0.710831in}}%
\pgfpathlineto{\pgfqpoint{6.178903in}{0.704000in}}%
\pgfpathlineto{\pgfqpoint{6.182259in}{0.676398in}}%
\pgfpathlineto{\pgfqpoint{6.185614in}{0.615536in}}%
\pgfpathlineto{\pgfqpoint{6.188970in}{0.682346in}}%
\pgfpathlineto{\pgfqpoint{6.192326in}{0.696532in}}%
\pgfpathlineto{\pgfqpoint{6.195682in}{0.696395in}}%
\pgfpathlineto{\pgfqpoint{6.199038in}{0.681679in}}%
\pgfpathlineto{\pgfqpoint{6.202393in}{0.606707in}}%
\pgfpathlineto{\pgfqpoint{6.205749in}{0.678875in}}%
\pgfpathlineto{\pgfqpoint{6.209105in}{0.705912in}}%
\pgfpathlineto{\pgfqpoint{6.212461in}{0.715053in}}%
\pgfpathlineto{\pgfqpoint{6.215816in}{0.708724in}}%
\pgfpathlineto{\pgfqpoint{6.219172in}{0.675009in}}%
\pgfpathlineto{\pgfqpoint{6.222528in}{0.665925in}}%
\pgfpathlineto{\pgfqpoint{6.225884in}{0.707331in}}%
\pgfpathlineto{\pgfqpoint{6.229240in}{0.720686in}}%
\pgfpathlineto{\pgfqpoint{6.232595in}{0.723078in}}%
\pgfpathlineto{\pgfqpoint{6.235951in}{0.715768in}}%
\pgfpathlineto{\pgfqpoint{6.239307in}{0.691309in}}%
\pgfpathlineto{\pgfqpoint{6.242663in}{0.643594in}}%
\pgfpathlineto{\pgfqpoint{6.246019in}{0.710564in}}%
\pgfpathlineto{\pgfqpoint{6.249374in}{0.730320in}}%
\pgfpathlineto{\pgfqpoint{6.252730in}{0.736292in}}%
\pgfpathlineto{\pgfqpoint{6.256086in}{0.727736in}}%
\pgfpathlineto{\pgfqpoint{6.259442in}{0.694067in}}%
\pgfpathlineto{\pgfqpoint{6.262797in}{0.677693in}}%
\pgfpathlineto{\pgfqpoint{6.266153in}{0.720439in}}%
\pgfpathlineto{\pgfqpoint{6.269509in}{0.732836in}}%
\pgfpathlineto{\pgfqpoint{6.272865in}{0.733366in}}%
\pgfpathlineto{\pgfqpoint{6.276221in}{0.722145in}}%
\pgfpathlineto{\pgfqpoint{6.279576in}{0.681514in}}%
\pgfpathlineto{\pgfqpoint{6.282932in}{0.697161in}}%
\pgfpathlineto{\pgfqpoint{6.286288in}{0.733727in}}%
\pgfpathlineto{\pgfqpoint{6.289644in}{0.747157in}}%
\pgfpathlineto{\pgfqpoint{6.292999in}{0.747131in}}%
\pgfpathlineto{\pgfqpoint{6.296355in}{0.731570in}}%
\pgfpathlineto{\pgfqpoint{6.303067in}{0.663157in}}%
\pgfpathlineto{\pgfqpoint{6.306423in}{0.706028in}}%
\pgfpathlineto{\pgfqpoint{6.309778in}{0.711811in}}%
\pgfpathlineto{\pgfqpoint{6.313134in}{0.697899in}}%
\pgfpathlineto{\pgfqpoint{6.316490in}{0.539589in}}%
\pgfpathlineto{\pgfqpoint{6.319846in}{0.706245in}}%
\pgfpathlineto{\pgfqpoint{6.323201in}{0.730581in}}%
\pgfpathlineto{\pgfqpoint{6.326557in}{0.737528in}}%
\pgfpathlineto{\pgfqpoint{6.329913in}{0.727905in}}%
\pgfpathlineto{\pgfqpoint{6.333269in}{0.705504in}}%
\pgfpathlineto{\pgfqpoint{6.336625in}{0.653372in}}%
\pgfpathlineto{\pgfqpoint{6.339980in}{0.656898in}}%
\pgfpathlineto{\pgfqpoint{6.343336in}{0.677642in}}%
\pgfpathlineto{\pgfqpoint{6.346692in}{0.665321in}}%
\pgfpathlineto{\pgfqpoint{6.350048in}{0.605474in}}%
\pgfpathlineto{\pgfqpoint{6.353403in}{0.688051in}}%
\pgfpathlineto{\pgfqpoint{6.356759in}{0.710044in}}%
\pgfpathlineto{\pgfqpoint{6.360115in}{0.715199in}}%
\pgfpathlineto{\pgfqpoint{6.363471in}{0.704943in}}%
\pgfpathlineto{\pgfqpoint{6.373538in}{0.666743in}}%
\pgfpathlineto{\pgfqpoint{6.376894in}{0.658751in}}%
\pgfpathlineto{\pgfqpoint{6.380250in}{0.653978in}}%
\pgfpathlineto{\pgfqpoint{6.383605in}{0.647083in}}%
\pgfpathlineto{\pgfqpoint{6.383605in}{0.647083in}}%
\pgfusepath{stroke}%
\end{pgfscope}%
\begin{pgfscope}%
\pgfpathrectangle{\pgfqpoint{4.621821in}{0.320625in}}{\pgfqpoint{1.845679in}{1.164375in}}%
\pgfusepath{clip}%
\pgfsetrectcap%
\pgfsetroundjoin%
\pgfsetlinewidth{0.602250pt}%
\definecolor{currentstroke}{rgb}{0.819608,0.286275,0.356863}%
\pgfsetstrokecolor{currentstroke}%
\pgfsetstrokeopacity{0.120000}%
\pgfsetdash{}{0pt}%
\pgfpathmoveto{\pgfqpoint{4.705715in}{1.469196in}}%
\pgfpathlineto{\pgfqpoint{4.725850in}{1.433697in}}%
\pgfpathlineto{\pgfqpoint{4.745985in}{1.394185in}}%
\pgfpathlineto{\pgfqpoint{4.766120in}{1.355343in}}%
\pgfpathlineto{\pgfqpoint{4.782898in}{1.325411in}}%
\pgfpathlineto{\pgfqpoint{4.803033in}{1.292118in}}%
\pgfpathlineto{\pgfqpoint{4.826524in}{1.255806in}}%
\pgfpathlineto{\pgfqpoint{4.853370in}{1.216465in}}%
\pgfpathlineto{\pgfqpoint{4.886928in}{1.169333in}}%
\pgfpathlineto{\pgfqpoint{4.930553in}{1.110113in}}%
\pgfpathlineto{\pgfqpoint{4.987601in}{1.034652in}}%
\pgfpathlineto{\pgfqpoint{5.064784in}{0.934470in}}%
\pgfpathlineto{\pgfqpoint{5.155390in}{0.818221in}}%
\pgfpathlineto{\pgfqpoint{5.175525in}{0.793136in}}%
\pgfpathlineto{\pgfqpoint{5.199015in}{0.761843in}}%
\pgfpathlineto{\pgfqpoint{5.209082in}{0.751148in}}%
\pgfpathlineto{\pgfqpoint{5.219150in}{0.741227in}}%
\pgfpathlineto{\pgfqpoint{5.225861in}{0.732366in}}%
\pgfpathlineto{\pgfqpoint{5.232573in}{0.721873in}}%
\pgfpathlineto{\pgfqpoint{5.256063in}{0.682552in}}%
\pgfpathlineto{\pgfqpoint{5.259419in}{0.679432in}}%
\pgfpathlineto{\pgfqpoint{5.262775in}{0.677664in}}%
\pgfpathlineto{\pgfqpoint{5.266131in}{0.677206in}}%
\pgfpathlineto{\pgfqpoint{5.269487in}{0.677800in}}%
\pgfpathlineto{\pgfqpoint{5.279554in}{0.681303in}}%
\pgfpathlineto{\pgfqpoint{5.282910in}{0.681317in}}%
\pgfpathlineto{\pgfqpoint{5.286265in}{0.679752in}}%
\pgfpathlineto{\pgfqpoint{5.289621in}{0.675802in}}%
\pgfpathlineto{\pgfqpoint{5.296333in}{0.664514in}}%
\pgfpathlineto{\pgfqpoint{5.299689in}{0.662426in}}%
\pgfpathlineto{\pgfqpoint{5.303044in}{0.666369in}}%
\pgfpathlineto{\pgfqpoint{5.309756in}{0.680627in}}%
\pgfpathlineto{\pgfqpoint{5.313112in}{0.686310in}}%
\pgfpathlineto{\pgfqpoint{5.316467in}{0.690371in}}%
\pgfpathlineto{\pgfqpoint{5.319823in}{0.692842in}}%
\pgfpathlineto{\pgfqpoint{5.323179in}{0.693756in}}%
\pgfpathlineto{\pgfqpoint{5.326535in}{0.693074in}}%
\pgfpathlineto{\pgfqpoint{5.329891in}{0.690629in}}%
\pgfpathlineto{\pgfqpoint{5.333246in}{0.686047in}}%
\pgfpathlineto{\pgfqpoint{5.336602in}{0.678546in}}%
\pgfpathlineto{\pgfqpoint{5.339958in}{0.666329in}}%
\pgfpathlineto{\pgfqpoint{5.343314in}{0.644014in}}%
\pgfpathlineto{\pgfqpoint{5.346669in}{0.592770in}}%
\pgfpathlineto{\pgfqpoint{5.350025in}{0.643297in}}%
\pgfpathlineto{\pgfqpoint{5.353381in}{0.669565in}}%
\pgfpathlineto{\pgfqpoint{5.356737in}{0.684709in}}%
\pgfpathlineto{\pgfqpoint{5.360093in}{0.694767in}}%
\pgfpathlineto{\pgfqpoint{5.363448in}{0.701626in}}%
\pgfpathlineto{\pgfqpoint{5.366804in}{0.706019in}}%
\pgfpathlineto{\pgfqpoint{5.370160in}{0.708156in}}%
\pgfpathlineto{\pgfqpoint{5.373516in}{0.707870in}}%
\pgfpathlineto{\pgfqpoint{5.376871in}{0.704493in}}%
\pgfpathlineto{\pgfqpoint{5.380227in}{0.696303in}}%
\pgfpathlineto{\pgfqpoint{5.386939in}{0.665916in}}%
\pgfpathlineto{\pgfqpoint{5.390295in}{0.667989in}}%
\pgfpathlineto{\pgfqpoint{5.393650in}{0.674756in}}%
\pgfpathlineto{\pgfqpoint{5.397006in}{0.673864in}}%
\pgfpathlineto{\pgfqpoint{5.400362in}{0.661944in}}%
\pgfpathlineto{\pgfqpoint{5.403718in}{0.638159in}}%
\pgfpathlineto{\pgfqpoint{5.410429in}{0.682559in}}%
\pgfpathlineto{\pgfqpoint{5.413785in}{0.691157in}}%
\pgfpathlineto{\pgfqpoint{5.417141in}{0.685561in}}%
\pgfpathlineto{\pgfqpoint{5.420497in}{0.663448in}}%
\pgfpathlineto{\pgfqpoint{5.423852in}{0.635231in}}%
\pgfpathlineto{\pgfqpoint{5.427208in}{0.663513in}}%
\pgfpathlineto{\pgfqpoint{5.430564in}{0.676816in}}%
\pgfpathlineto{\pgfqpoint{5.433920in}{0.677289in}}%
\pgfpathlineto{\pgfqpoint{5.437276in}{0.664479in}}%
\pgfpathlineto{\pgfqpoint{5.440631in}{0.614462in}}%
\pgfpathlineto{\pgfqpoint{5.443987in}{0.654085in}}%
\pgfpathlineto{\pgfqpoint{5.447343in}{0.683708in}}%
\pgfpathlineto{\pgfqpoint{5.450699in}{0.694249in}}%
\pgfpathlineto{\pgfqpoint{5.454054in}{0.689666in}}%
\pgfpathlineto{\pgfqpoint{5.457410in}{0.630078in}}%
\pgfpathlineto{\pgfqpoint{5.460766in}{0.703415in}}%
\pgfpathlineto{\pgfqpoint{5.464122in}{0.730577in}}%
\pgfpathlineto{\pgfqpoint{5.467478in}{0.743633in}}%
\pgfpathlineto{\pgfqpoint{5.470833in}{0.749242in}}%
\pgfpathlineto{\pgfqpoint{5.474189in}{0.748960in}}%
\pgfpathlineto{\pgfqpoint{5.477545in}{0.742168in}}%
\pgfpathlineto{\pgfqpoint{5.480901in}{0.724548in}}%
\pgfpathlineto{\pgfqpoint{5.484256in}{0.660414in}}%
\pgfpathlineto{\pgfqpoint{5.487612in}{0.714048in}}%
\pgfpathlineto{\pgfqpoint{5.490968in}{0.743511in}}%
\pgfpathlineto{\pgfqpoint{5.494324in}{0.756444in}}%
\pgfpathlineto{\pgfqpoint{5.497680in}{0.760026in}}%
\pgfpathlineto{\pgfqpoint{5.501035in}{0.752580in}}%
\pgfpathlineto{\pgfqpoint{5.504391in}{0.734685in}}%
\pgfpathlineto{\pgfqpoint{5.507747in}{0.704145in}}%
\pgfpathlineto{\pgfqpoint{5.511103in}{0.653071in}}%
\pgfpathlineto{\pgfqpoint{5.514458in}{0.665459in}}%
\pgfpathlineto{\pgfqpoint{5.517814in}{0.641278in}}%
\pgfpathlineto{\pgfqpoint{5.524526in}{0.711833in}}%
\pgfpathlineto{\pgfqpoint{5.527882in}{0.724893in}}%
\pgfpathlineto{\pgfqpoint{5.531237in}{0.719881in}}%
\pgfpathlineto{\pgfqpoint{5.534593in}{0.701582in}}%
\pgfpathlineto{\pgfqpoint{5.537949in}{0.669425in}}%
\pgfpathlineto{\pgfqpoint{5.541305in}{0.604689in}}%
\pgfpathlineto{\pgfqpoint{5.544660in}{0.632689in}}%
\pgfpathlineto{\pgfqpoint{5.548016in}{0.604035in}}%
\pgfpathlineto{\pgfqpoint{5.554728in}{0.680088in}}%
\pgfpathlineto{\pgfqpoint{5.558084in}{0.693116in}}%
\pgfpathlineto{\pgfqpoint{5.561439in}{0.687476in}}%
\pgfpathlineto{\pgfqpoint{5.564795in}{0.651755in}}%
\pgfpathlineto{\pgfqpoint{5.568151in}{0.648575in}}%
\pgfpathlineto{\pgfqpoint{5.571507in}{0.686537in}}%
\pgfpathlineto{\pgfqpoint{5.574863in}{0.698193in}}%
\pgfpathlineto{\pgfqpoint{5.578218in}{0.698310in}}%
\pgfpathlineto{\pgfqpoint{5.581574in}{0.686480in}}%
\pgfpathlineto{\pgfqpoint{5.584930in}{0.642845in}}%
\pgfpathlineto{\pgfqpoint{5.588286in}{0.664805in}}%
\pgfpathlineto{\pgfqpoint{5.591641in}{0.699407in}}%
\pgfpathlineto{\pgfqpoint{5.594997in}{0.711762in}}%
\pgfpathlineto{\pgfqpoint{5.598353in}{0.709804in}}%
\pgfpathlineto{\pgfqpoint{5.601709in}{0.684784in}}%
\pgfpathlineto{\pgfqpoint{5.605065in}{0.640923in}}%
\pgfpathlineto{\pgfqpoint{5.608420in}{0.700480in}}%
\pgfpathlineto{\pgfqpoint{5.611776in}{0.716350in}}%
\pgfpathlineto{\pgfqpoint{5.615132in}{0.719448in}}%
\pgfpathlineto{\pgfqpoint{5.618488in}{0.711809in}}%
\pgfpathlineto{\pgfqpoint{5.625199in}{0.657790in}}%
\pgfpathlineto{\pgfqpoint{5.628555in}{0.711742in}}%
\pgfpathlineto{\pgfqpoint{5.631911in}{0.729103in}}%
\pgfpathlineto{\pgfqpoint{5.635267in}{0.732181in}}%
\pgfpathlineto{\pgfqpoint{5.638622in}{0.716767in}}%
\pgfpathlineto{\pgfqpoint{5.641978in}{0.631487in}}%
\pgfpathlineto{\pgfqpoint{5.645334in}{0.709813in}}%
\pgfpathlineto{\pgfqpoint{5.648690in}{0.732723in}}%
\pgfpathlineto{\pgfqpoint{5.652045in}{0.739443in}}%
\pgfpathlineto{\pgfqpoint{5.655401in}{0.735661in}}%
\pgfpathlineto{\pgfqpoint{5.658757in}{0.717457in}}%
\pgfpathlineto{\pgfqpoint{5.662113in}{0.582257in}}%
\pgfpathlineto{\pgfqpoint{5.665469in}{0.720912in}}%
\pgfpathlineto{\pgfqpoint{5.668824in}{0.744953in}}%
\pgfpathlineto{\pgfqpoint{5.672180in}{0.752604in}}%
\pgfpathlineto{\pgfqpoint{5.675536in}{0.744897in}}%
\pgfpathlineto{\pgfqpoint{5.678892in}{0.716497in}}%
\pgfpathlineto{\pgfqpoint{5.682247in}{0.657811in}}%
\pgfpathlineto{\pgfqpoint{5.685603in}{0.721470in}}%
\pgfpathlineto{\pgfqpoint{5.688959in}{0.733678in}}%
\pgfpathlineto{\pgfqpoint{5.692315in}{0.730015in}}%
\pgfpathlineto{\pgfqpoint{5.699026in}{0.678450in}}%
\pgfpathlineto{\pgfqpoint{5.702382in}{0.733263in}}%
\pgfpathlineto{\pgfqpoint{5.705738in}{0.750221in}}%
\pgfpathlineto{\pgfqpoint{5.709094in}{0.750959in}}%
\pgfpathlineto{\pgfqpoint{5.712449in}{0.734959in}}%
\pgfpathlineto{\pgfqpoint{5.715805in}{0.701373in}}%
\pgfpathlineto{\pgfqpoint{5.719161in}{0.636891in}}%
\pgfpathlineto{\pgfqpoint{5.722517in}{0.694266in}}%
\pgfpathlineto{\pgfqpoint{5.725873in}{0.694906in}}%
\pgfpathlineto{\pgfqpoint{5.729228in}{0.651779in}}%
\pgfpathlineto{\pgfqpoint{5.735940in}{0.721739in}}%
\pgfpathlineto{\pgfqpoint{5.739296in}{0.732470in}}%
\pgfpathlineto{\pgfqpoint{5.742652in}{0.724697in}}%
\pgfpathlineto{\pgfqpoint{5.746007in}{0.701127in}}%
\pgfpathlineto{\pgfqpoint{5.749363in}{0.633736in}}%
\pgfpathlineto{\pgfqpoint{5.752719in}{0.668513in}}%
\pgfpathlineto{\pgfqpoint{5.756075in}{0.684114in}}%
\pgfpathlineto{\pgfqpoint{5.759430in}{0.671945in}}%
\pgfpathlineto{\pgfqpoint{5.762786in}{0.582769in}}%
\pgfpathlineto{\pgfqpoint{5.766142in}{0.688257in}}%
\pgfpathlineto{\pgfqpoint{5.769498in}{0.710665in}}%
\pgfpathlineto{\pgfqpoint{5.772854in}{0.714811in}}%
\pgfpathlineto{\pgfqpoint{5.776209in}{0.698530in}}%
\pgfpathlineto{\pgfqpoint{5.779565in}{0.639154in}}%
\pgfpathlineto{\pgfqpoint{5.782921in}{0.675763in}}%
\pgfpathlineto{\pgfqpoint{5.786277in}{0.700340in}}%
\pgfpathlineto{\pgfqpoint{5.789632in}{0.705159in}}%
\pgfpathlineto{\pgfqpoint{5.792988in}{0.696317in}}%
\pgfpathlineto{\pgfqpoint{5.796344in}{0.658646in}}%
\pgfpathlineto{\pgfqpoint{5.799700in}{0.670197in}}%
\pgfpathlineto{\pgfqpoint{5.803056in}{0.707983in}}%
\pgfpathlineto{\pgfqpoint{5.806411in}{0.720820in}}%
\pgfpathlineto{\pgfqpoint{5.809767in}{0.717758in}}%
\pgfpathlineto{\pgfqpoint{5.813123in}{0.687628in}}%
\pgfpathlineto{\pgfqpoint{5.816479in}{0.671080in}}%
\pgfpathlineto{\pgfqpoint{5.819834in}{0.716224in}}%
\pgfpathlineto{\pgfqpoint{5.823190in}{0.730223in}}%
\pgfpathlineto{\pgfqpoint{5.826546in}{0.732653in}}%
\pgfpathlineto{\pgfqpoint{5.829902in}{0.724795in}}%
\pgfpathlineto{\pgfqpoint{5.836613in}{0.667974in}}%
\pgfpathlineto{\pgfqpoint{5.839969in}{0.723626in}}%
\pgfpathlineto{\pgfqpoint{5.843325in}{0.741403in}}%
\pgfpathlineto{\pgfqpoint{5.846681in}{0.745142in}}%
\pgfpathlineto{\pgfqpoint{5.850036in}{0.733249in}}%
\pgfpathlineto{\pgfqpoint{5.853392in}{0.704145in}}%
\pgfpathlineto{\pgfqpoint{5.856748in}{0.605958in}}%
\pgfpathlineto{\pgfqpoint{5.860104in}{0.694787in}}%
\pgfpathlineto{\pgfqpoint{5.863460in}{0.703286in}}%
\pgfpathlineto{\pgfqpoint{5.866815in}{0.687959in}}%
\pgfpathlineto{\pgfqpoint{5.870171in}{0.618490in}}%
\pgfpathlineto{\pgfqpoint{5.873527in}{0.705905in}}%
\pgfpathlineto{\pgfqpoint{5.876883in}{0.727003in}}%
\pgfpathlineto{\pgfqpoint{5.880238in}{0.730250in}}%
\pgfpathlineto{\pgfqpoint{5.883594in}{0.715100in}}%
\pgfpathlineto{\pgfqpoint{5.890306in}{0.640524in}}%
\pgfpathlineto{\pgfqpoint{5.893662in}{0.685160in}}%
\pgfpathlineto{\pgfqpoint{5.897017in}{0.688531in}}%
\pgfpathlineto{\pgfqpoint{5.903729in}{0.645401in}}%
\pgfpathlineto{\pgfqpoint{5.907085in}{0.698758in}}%
\pgfpathlineto{\pgfqpoint{5.910441in}{0.715472in}}%
\pgfpathlineto{\pgfqpoint{5.913796in}{0.714635in}}%
\pgfpathlineto{\pgfqpoint{5.917152in}{0.692170in}}%
\pgfpathlineto{\pgfqpoint{5.920508in}{0.555581in}}%
\pgfpathlineto{\pgfqpoint{5.923864in}{0.689210in}}%
\pgfpathlineto{\pgfqpoint{5.927219in}{0.706491in}}%
\pgfpathlineto{\pgfqpoint{5.930575in}{0.708034in}}%
\pgfpathlineto{\pgfqpoint{5.933931in}{0.695379in}}%
\pgfpathlineto{\pgfqpoint{5.937287in}{0.638237in}}%
\pgfpathlineto{\pgfqpoint{5.940643in}{0.685676in}}%
\pgfpathlineto{\pgfqpoint{5.943998in}{0.715220in}}%
\pgfpathlineto{\pgfqpoint{5.947354in}{0.724862in}}%
\pgfpathlineto{\pgfqpoint{5.950710in}{0.717855in}}%
\pgfpathlineto{\pgfqpoint{5.954066in}{0.679131in}}%
\pgfpathlineto{\pgfqpoint{5.957421in}{0.684743in}}%
\pgfpathlineto{\pgfqpoint{5.960777in}{0.720301in}}%
\pgfpathlineto{\pgfqpoint{5.964133in}{0.731587in}}%
\pgfpathlineto{\pgfqpoint{5.967489in}{0.731739in}}%
\pgfpathlineto{\pgfqpoint{5.970845in}{0.720353in}}%
\pgfpathlineto{\pgfqpoint{5.974200in}{0.679764in}}%
\pgfpathlineto{\pgfqpoint{5.977556in}{0.694829in}}%
\pgfpathlineto{\pgfqpoint{5.980912in}{0.731305in}}%
\pgfpathlineto{\pgfqpoint{5.984268in}{0.744398in}}%
\pgfpathlineto{\pgfqpoint{5.987623in}{0.743506in}}%
\pgfpathlineto{\pgfqpoint{5.990979in}{0.723091in}}%
\pgfpathlineto{\pgfqpoint{5.994335in}{0.595707in}}%
\pgfpathlineto{\pgfqpoint{5.997691in}{0.716852in}}%
\pgfpathlineto{\pgfqpoint{6.001047in}{0.735981in}}%
\pgfpathlineto{\pgfqpoint{6.004402in}{0.738980in}}%
\pgfpathlineto{\pgfqpoint{6.007758in}{0.728872in}}%
\pgfpathlineto{\pgfqpoint{6.011114in}{0.688012in}}%
\pgfpathlineto{\pgfqpoint{6.014470in}{0.705927in}}%
\pgfpathlineto{\pgfqpoint{6.017825in}{0.741717in}}%
\pgfpathlineto{\pgfqpoint{6.021181in}{0.754012in}}%
\pgfpathlineto{\pgfqpoint{6.024537in}{0.750757in}}%
\pgfpathlineto{\pgfqpoint{6.027893in}{0.731191in}}%
\pgfpathlineto{\pgfqpoint{6.031249in}{0.681795in}}%
\pgfpathlineto{\pgfqpoint{6.034604in}{0.686968in}}%
\pgfpathlineto{\pgfqpoint{6.037960in}{0.711832in}}%
\pgfpathlineto{\pgfqpoint{6.041316in}{0.708830in}}%
\pgfpathlineto{\pgfqpoint{6.044672in}{0.670241in}}%
\pgfpathlineto{\pgfqpoint{6.048027in}{0.693036in}}%
\pgfpathlineto{\pgfqpoint{6.051383in}{0.728894in}}%
\pgfpathlineto{\pgfqpoint{6.054739in}{0.740521in}}%
\pgfpathlineto{\pgfqpoint{6.058095in}{0.734158in}}%
\pgfpathlineto{\pgfqpoint{6.061451in}{0.712796in}}%
\pgfpathlineto{\pgfqpoint{6.064806in}{0.662648in}}%
\pgfpathlineto{\pgfqpoint{6.068162in}{0.662602in}}%
\pgfpathlineto{\pgfqpoint{6.071518in}{0.684485in}}%
\pgfpathlineto{\pgfqpoint{6.074874in}{0.672303in}}%
\pgfpathlineto{\pgfqpoint{6.078230in}{0.613606in}}%
\pgfpathlineto{\pgfqpoint{6.081585in}{0.694982in}}%
\pgfpathlineto{\pgfqpoint{6.084941in}{0.716146in}}%
\pgfpathlineto{\pgfqpoint{6.088297in}{0.719216in}}%
\pgfpathlineto{\pgfqpoint{6.091653in}{0.702574in}}%
\pgfpathlineto{\pgfqpoint{6.095008in}{0.654356in}}%
\pgfpathlineto{\pgfqpoint{6.098364in}{0.664898in}}%
\pgfpathlineto{\pgfqpoint{6.101720in}{0.692515in}}%
\pgfpathlineto{\pgfqpoint{6.105076in}{0.695808in}}%
\pgfpathlineto{\pgfqpoint{6.108432in}{0.681268in}}%
\pgfpathlineto{\pgfqpoint{6.111787in}{0.562765in}}%
\pgfpathlineto{\pgfqpoint{6.115143in}{0.686539in}}%
\pgfpathlineto{\pgfqpoint{6.118499in}{0.710791in}}%
\pgfpathlineto{\pgfqpoint{6.121855in}{0.716897in}}%
\pgfpathlineto{\pgfqpoint{6.125210in}{0.703405in}}%
\pgfpathlineto{\pgfqpoint{6.128566in}{0.636901in}}%
\pgfpathlineto{\pgfqpoint{6.131922in}{0.690390in}}%
\pgfpathlineto{\pgfqpoint{6.135278in}{0.715408in}}%
\pgfpathlineto{\pgfqpoint{6.138634in}{0.722866in}}%
\pgfpathlineto{\pgfqpoint{6.141989in}{0.719599in}}%
\pgfpathlineto{\pgfqpoint{6.145345in}{0.702156in}}%
\pgfpathlineto{\pgfqpoint{6.148701in}{0.596701in}}%
\pgfpathlineto{\pgfqpoint{6.152057in}{0.703500in}}%
\pgfpathlineto{\pgfqpoint{6.155412in}{0.728283in}}%
\pgfpathlineto{\pgfqpoint{6.158768in}{0.736273in}}%
\pgfpathlineto{\pgfqpoint{6.162124in}{0.728603in}}%
\pgfpathlineto{\pgfqpoint{6.165480in}{0.689842in}}%
\pgfpathlineto{\pgfqpoint{6.168836in}{0.694596in}}%
\pgfpathlineto{\pgfqpoint{6.172191in}{0.730152in}}%
\pgfpathlineto{\pgfqpoint{6.175547in}{0.741236in}}%
\pgfpathlineto{\pgfqpoint{6.178903in}{0.741037in}}%
\pgfpathlineto{\pgfqpoint{6.182259in}{0.728876in}}%
\pgfpathlineto{\pgfqpoint{6.185614in}{0.683896in}}%
\pgfpathlineto{\pgfqpoint{6.188970in}{0.708454in}}%
\pgfpathlineto{\pgfqpoint{6.192326in}{0.742403in}}%
\pgfpathlineto{\pgfqpoint{6.195682in}{0.754547in}}%
\pgfpathlineto{\pgfqpoint{6.199038in}{0.752415in}}%
\pgfpathlineto{\pgfqpoint{6.202393in}{0.732706in}}%
\pgfpathlineto{\pgfqpoint{6.205749in}{0.670013in}}%
\pgfpathlineto{\pgfqpoint{6.209105in}{0.706376in}}%
\pgfpathlineto{\pgfqpoint{6.212461in}{0.728037in}}%
\pgfpathlineto{\pgfqpoint{6.215816in}{0.728538in}}%
\pgfpathlineto{\pgfqpoint{6.219172in}{0.709241in}}%
\pgfpathlineto{\pgfqpoint{6.222528in}{0.649908in}}%
\pgfpathlineto{\pgfqpoint{6.225884in}{0.727278in}}%
\pgfpathlineto{\pgfqpoint{6.229240in}{0.747258in}}%
\pgfpathlineto{\pgfqpoint{6.232595in}{0.750242in}}%
\pgfpathlineto{\pgfqpoint{6.235951in}{0.736051in}}%
\pgfpathlineto{\pgfqpoint{6.239307in}{0.707092in}}%
\pgfpathlineto{\pgfqpoint{6.242663in}{0.579189in}}%
\pgfpathlineto{\pgfqpoint{6.246019in}{0.683020in}}%
\pgfpathlineto{\pgfqpoint{6.249374in}{0.687134in}}%
\pgfpathlineto{\pgfqpoint{6.252730in}{0.642928in}}%
\pgfpathlineto{\pgfqpoint{6.256086in}{0.684316in}}%
\pgfpathlineto{\pgfqpoint{6.259442in}{0.717441in}}%
\pgfpathlineto{\pgfqpoint{6.262797in}{0.728108in}}%
\pgfpathlineto{\pgfqpoint{6.266153in}{0.720098in}}%
\pgfpathlineto{\pgfqpoint{6.269509in}{0.695607in}}%
\pgfpathlineto{\pgfqpoint{6.272865in}{0.614664in}}%
\pgfpathlineto{\pgfqpoint{6.276221in}{0.670122in}}%
\pgfpathlineto{\pgfqpoint{6.279576in}{0.684205in}}%
\pgfpathlineto{\pgfqpoint{6.282932in}{0.673106in}}%
\pgfpathlineto{\pgfqpoint{6.286288in}{0.508979in}}%
\pgfpathlineto{\pgfqpoint{6.289644in}{0.683386in}}%
\pgfpathlineto{\pgfqpoint{6.292999in}{0.707506in}}%
\pgfpathlineto{\pgfqpoint{6.296355in}{0.712959in}}%
\pgfpathlineto{\pgfqpoint{6.299711in}{0.698035in}}%
\pgfpathlineto{\pgfqpoint{6.303067in}{0.637689in}}%
\pgfpathlineto{\pgfqpoint{6.306423in}{0.678341in}}%
\pgfpathlineto{\pgfqpoint{6.309778in}{0.703449in}}%
\pgfpathlineto{\pgfqpoint{6.313134in}{0.709549in}}%
\pgfpathlineto{\pgfqpoint{6.316490in}{0.703392in}}%
\pgfpathlineto{\pgfqpoint{6.323201in}{0.649770in}}%
\pgfpathlineto{\pgfqpoint{6.326557in}{0.704358in}}%
\pgfpathlineto{\pgfqpoint{6.329913in}{0.721527in}}%
\pgfpathlineto{\pgfqpoint{6.333269in}{0.723535in}}%
\pgfpathlineto{\pgfqpoint{6.336625in}{0.704164in}}%
\pgfpathlineto{\pgfqpoint{6.339980in}{0.625761in}}%
\pgfpathlineto{\pgfqpoint{6.343336in}{0.712965in}}%
\pgfpathlineto{\pgfqpoint{6.346692in}{0.732129in}}%
\pgfpathlineto{\pgfqpoint{6.350048in}{0.737793in}}%
\pgfpathlineto{\pgfqpoint{6.353403in}{0.733794in}}%
\pgfpathlineto{\pgfqpoint{6.356759in}{0.716302in}}%
\pgfpathlineto{\pgfqpoint{6.360115in}{0.623215in}}%
\pgfpathlineto{\pgfqpoint{6.363471in}{0.715569in}}%
\pgfpathlineto{\pgfqpoint{6.366827in}{0.741069in}}%
\pgfpathlineto{\pgfqpoint{6.370182in}{0.749884in}}%
\pgfpathlineto{\pgfqpoint{6.373538in}{0.744591in}}%
\pgfpathlineto{\pgfqpoint{6.383605in}{0.708917in}}%
\pgfpathlineto{\pgfqpoint{6.383605in}{0.708917in}}%
\pgfusepath{stroke}%
\end{pgfscope}%
\begin{pgfscope}%
\pgfpathrectangle{\pgfqpoint{4.621821in}{0.320625in}}{\pgfqpoint{1.845679in}{1.164375in}}%
\pgfusepath{clip}%
\pgfsetrectcap%
\pgfsetroundjoin%
\pgfsetlinewidth{1.154312pt}%
\definecolor{currentstroke}{rgb}{0.819608,0.286275,0.356863}%
\pgfsetstrokecolor{currentstroke}%
\pgfsetdash{}{0pt}%
\pgfpathmoveto{\pgfqpoint{4.705715in}{1.471665in}}%
\pgfpathlineto{\pgfqpoint{4.729206in}{1.429105in}}%
\pgfpathlineto{\pgfqpoint{4.772831in}{1.345188in}}%
\pgfpathlineto{\pgfqpoint{4.789610in}{1.316203in}}%
\pgfpathlineto{\pgfqpoint{4.809745in}{1.283688in}}%
\pgfpathlineto{\pgfqpoint{4.833235in}{1.248033in}}%
\pgfpathlineto{\pgfqpoint{4.860081in}{1.209475in}}%
\pgfpathlineto{\pgfqpoint{4.886928in}{1.172754in}}%
\pgfpathlineto{\pgfqpoint{4.917130in}{1.133333in}}%
\pgfpathlineto{\pgfqpoint{4.947332in}{1.095721in}}%
\pgfpathlineto{\pgfqpoint{4.980889in}{1.055802in}}%
\pgfpathlineto{\pgfqpoint{5.017803in}{1.013749in}}%
\pgfpathlineto{\pgfqpoint{5.061428in}{0.965865in}}%
\pgfpathlineto{\pgfqpoint{5.118476in}{0.905069in}}%
\pgfpathlineto{\pgfqpoint{5.202371in}{0.817502in}}%
\pgfpathlineto{\pgfqpoint{5.235929in}{0.782739in}}%
\pgfpathlineto{\pgfqpoint{5.256063in}{0.761975in}}%
\pgfpathlineto{\pgfqpoint{5.282910in}{0.735825in}}%
\pgfpathlineto{\pgfqpoint{5.309756in}{0.705947in}}%
\pgfpathlineto{\pgfqpoint{5.316467in}{0.700707in}}%
\pgfpathlineto{\pgfqpoint{5.323179in}{0.697124in}}%
\pgfpathlineto{\pgfqpoint{5.326535in}{0.696559in}}%
\pgfpathlineto{\pgfqpoint{5.329891in}{0.696988in}}%
\pgfpathlineto{\pgfqpoint{5.339958in}{0.700290in}}%
\pgfpathlineto{\pgfqpoint{5.343314in}{0.700487in}}%
\pgfpathlineto{\pgfqpoint{5.350025in}{0.699523in}}%
\pgfpathlineto{\pgfqpoint{5.353381in}{0.699973in}}%
\pgfpathlineto{\pgfqpoint{5.356737in}{0.701648in}}%
\pgfpathlineto{\pgfqpoint{5.366804in}{0.708951in}}%
\pgfpathlineto{\pgfqpoint{5.370160in}{0.710533in}}%
\pgfpathlineto{\pgfqpoint{5.376871in}{0.712021in}}%
\pgfpathlineto{\pgfqpoint{5.383583in}{0.713227in}}%
\pgfpathlineto{\pgfqpoint{5.390295in}{0.716195in}}%
\pgfpathlineto{\pgfqpoint{5.393650in}{0.716928in}}%
\pgfpathlineto{\pgfqpoint{5.397006in}{0.716191in}}%
\pgfpathlineto{\pgfqpoint{5.400362in}{0.714527in}}%
\pgfpathlineto{\pgfqpoint{5.403718in}{0.711627in}}%
\pgfpathlineto{\pgfqpoint{5.407074in}{0.705673in}}%
\pgfpathlineto{\pgfqpoint{5.410429in}{0.695702in}}%
\pgfpathlineto{\pgfqpoint{5.413785in}{0.687765in}}%
\pgfpathlineto{\pgfqpoint{5.417141in}{0.691640in}}%
\pgfpathlineto{\pgfqpoint{5.420497in}{0.700800in}}%
\pgfpathlineto{\pgfqpoint{5.423852in}{0.707994in}}%
\pgfpathlineto{\pgfqpoint{5.427208in}{0.712061in}}%
\pgfpathlineto{\pgfqpoint{5.433920in}{0.714831in}}%
\pgfpathlineto{\pgfqpoint{5.437276in}{0.715499in}}%
\pgfpathlineto{\pgfqpoint{5.440631in}{0.712521in}}%
\pgfpathlineto{\pgfqpoint{5.443987in}{0.703793in}}%
\pgfpathlineto{\pgfqpoint{5.450699in}{0.678712in}}%
\pgfpathlineto{\pgfqpoint{5.454054in}{0.674593in}}%
\pgfpathlineto{\pgfqpoint{5.457410in}{0.668999in}}%
\pgfpathlineto{\pgfqpoint{5.460766in}{0.683369in}}%
\pgfpathlineto{\pgfqpoint{5.464122in}{0.702787in}}%
\pgfpathlineto{\pgfqpoint{5.467478in}{0.715448in}}%
\pgfpathlineto{\pgfqpoint{5.470833in}{0.722120in}}%
\pgfpathlineto{\pgfqpoint{5.474189in}{0.722908in}}%
\pgfpathlineto{\pgfqpoint{5.477545in}{0.717432in}}%
\pgfpathlineto{\pgfqpoint{5.484256in}{0.696662in}}%
\pgfpathlineto{\pgfqpoint{5.487612in}{0.708670in}}%
\pgfpathlineto{\pgfqpoint{5.490968in}{0.723742in}}%
\pgfpathlineto{\pgfqpoint{5.494324in}{0.732500in}}%
\pgfpathlineto{\pgfqpoint{5.497680in}{0.734568in}}%
\pgfpathlineto{\pgfqpoint{5.501035in}{0.728227in}}%
\pgfpathlineto{\pgfqpoint{5.504391in}{0.718142in}}%
\pgfpathlineto{\pgfqpoint{5.507747in}{0.714806in}}%
\pgfpathlineto{\pgfqpoint{5.511103in}{0.717596in}}%
\pgfpathlineto{\pgfqpoint{5.514458in}{0.717492in}}%
\pgfpathlineto{\pgfqpoint{5.517814in}{0.713449in}}%
\pgfpathlineto{\pgfqpoint{5.521170in}{0.713975in}}%
\pgfpathlineto{\pgfqpoint{5.524526in}{0.715184in}}%
\pgfpathlineto{\pgfqpoint{5.527882in}{0.713064in}}%
\pgfpathlineto{\pgfqpoint{5.531237in}{0.703180in}}%
\pgfpathlineto{\pgfqpoint{5.534593in}{0.683858in}}%
\pgfpathlineto{\pgfqpoint{5.537949in}{0.684370in}}%
\pgfpathlineto{\pgfqpoint{5.541305in}{0.704920in}}%
\pgfpathlineto{\pgfqpoint{5.544660in}{0.715899in}}%
\pgfpathlineto{\pgfqpoint{5.548016in}{0.715047in}}%
\pgfpathlineto{\pgfqpoint{5.551372in}{0.704191in}}%
\pgfpathlineto{\pgfqpoint{5.554728in}{0.700278in}}%
\pgfpathlineto{\pgfqpoint{5.558084in}{0.706363in}}%
\pgfpathlineto{\pgfqpoint{5.561439in}{0.705409in}}%
\pgfpathlineto{\pgfqpoint{5.564795in}{0.698445in}}%
\pgfpathlineto{\pgfqpoint{5.568151in}{0.688386in}}%
\pgfpathlineto{\pgfqpoint{5.571507in}{0.685469in}}%
\pgfpathlineto{\pgfqpoint{5.578218in}{0.708280in}}%
\pgfpathlineto{\pgfqpoint{5.581574in}{0.711261in}}%
\pgfpathlineto{\pgfqpoint{5.584930in}{0.706389in}}%
\pgfpathlineto{\pgfqpoint{5.588286in}{0.703024in}}%
\pgfpathlineto{\pgfqpoint{5.594997in}{0.712912in}}%
\pgfpathlineto{\pgfqpoint{5.598353in}{0.711345in}}%
\pgfpathlineto{\pgfqpoint{5.601709in}{0.705081in}}%
\pgfpathlineto{\pgfqpoint{5.605065in}{0.700714in}}%
\pgfpathlineto{\pgfqpoint{5.608420in}{0.705845in}}%
\pgfpathlineto{\pgfqpoint{5.611776in}{0.714727in}}%
\pgfpathlineto{\pgfqpoint{5.615132in}{0.719989in}}%
\pgfpathlineto{\pgfqpoint{5.618488in}{0.718464in}}%
\pgfpathlineto{\pgfqpoint{5.621843in}{0.714280in}}%
\pgfpathlineto{\pgfqpoint{5.625199in}{0.715995in}}%
\pgfpathlineto{\pgfqpoint{5.628555in}{0.723118in}}%
\pgfpathlineto{\pgfqpoint{5.631911in}{0.727759in}}%
\pgfpathlineto{\pgfqpoint{5.635267in}{0.725758in}}%
\pgfpathlineto{\pgfqpoint{5.638622in}{0.715868in}}%
\pgfpathlineto{\pgfqpoint{5.641978in}{0.708713in}}%
\pgfpathlineto{\pgfqpoint{5.648690in}{0.732190in}}%
\pgfpathlineto{\pgfqpoint{5.652045in}{0.737098in}}%
\pgfpathlineto{\pgfqpoint{5.655401in}{0.732731in}}%
\pgfpathlineto{\pgfqpoint{5.658757in}{0.722579in}}%
\pgfpathlineto{\pgfqpoint{5.662113in}{0.721595in}}%
\pgfpathlineto{\pgfqpoint{5.665469in}{0.731673in}}%
\pgfpathlineto{\pgfqpoint{5.668824in}{0.739796in}}%
\pgfpathlineto{\pgfqpoint{5.672180in}{0.739614in}}%
\pgfpathlineto{\pgfqpoint{5.675536in}{0.727148in}}%
\pgfpathlineto{\pgfqpoint{5.678892in}{0.704403in}}%
\pgfpathlineto{\pgfqpoint{5.682247in}{0.717385in}}%
\pgfpathlineto{\pgfqpoint{5.685603in}{0.734180in}}%
\pgfpathlineto{\pgfqpoint{5.688959in}{0.737477in}}%
\pgfpathlineto{\pgfqpoint{5.692315in}{0.728463in}}%
\pgfpathlineto{\pgfqpoint{5.695671in}{0.714506in}}%
\pgfpathlineto{\pgfqpoint{5.699026in}{0.714053in}}%
\pgfpathlineto{\pgfqpoint{5.702382in}{0.729941in}}%
\pgfpathlineto{\pgfqpoint{5.705738in}{0.739423in}}%
\pgfpathlineto{\pgfqpoint{5.709094in}{0.737481in}}%
\pgfpathlineto{\pgfqpoint{5.715805in}{0.710338in}}%
\pgfpathlineto{\pgfqpoint{5.719161in}{0.715426in}}%
\pgfpathlineto{\pgfqpoint{5.722517in}{0.714934in}}%
\pgfpathlineto{\pgfqpoint{5.729228in}{0.696294in}}%
\pgfpathlineto{\pgfqpoint{5.732584in}{0.700615in}}%
\pgfpathlineto{\pgfqpoint{5.735940in}{0.713555in}}%
\pgfpathlineto{\pgfqpoint{5.739296in}{0.722667in}}%
\pgfpathlineto{\pgfqpoint{5.742652in}{0.720974in}}%
\pgfpathlineto{\pgfqpoint{5.746007in}{0.711222in}}%
\pgfpathlineto{\pgfqpoint{5.749363in}{0.704077in}}%
\pgfpathlineto{\pgfqpoint{5.752719in}{0.703530in}}%
\pgfpathlineto{\pgfqpoint{5.756075in}{0.694565in}}%
\pgfpathlineto{\pgfqpoint{5.759430in}{0.681124in}}%
\pgfpathlineto{\pgfqpoint{5.762786in}{0.684925in}}%
\pgfpathlineto{\pgfqpoint{5.766142in}{0.699803in}}%
\pgfpathlineto{\pgfqpoint{5.769498in}{0.708436in}}%
\pgfpathlineto{\pgfqpoint{5.772854in}{0.709831in}}%
\pgfpathlineto{\pgfqpoint{5.779565in}{0.700294in}}%
\pgfpathlineto{\pgfqpoint{5.782921in}{0.704242in}}%
\pgfpathlineto{\pgfqpoint{5.786277in}{0.711246in}}%
\pgfpathlineto{\pgfqpoint{5.789632in}{0.708808in}}%
\pgfpathlineto{\pgfqpoint{5.792988in}{0.696596in}}%
\pgfpathlineto{\pgfqpoint{5.796344in}{0.690288in}}%
\pgfpathlineto{\pgfqpoint{5.803056in}{0.716617in}}%
\pgfpathlineto{\pgfqpoint{5.806411in}{0.721627in}}%
\pgfpathlineto{\pgfqpoint{5.809767in}{0.717969in}}%
\pgfpathlineto{\pgfqpoint{5.813123in}{0.707364in}}%
\pgfpathlineto{\pgfqpoint{5.816479in}{0.706386in}}%
\pgfpathlineto{\pgfqpoint{5.819834in}{0.720241in}}%
\pgfpathlineto{\pgfqpoint{5.823190in}{0.729666in}}%
\pgfpathlineto{\pgfqpoint{5.826546in}{0.728821in}}%
\pgfpathlineto{\pgfqpoint{5.833258in}{0.711596in}}%
\pgfpathlineto{\pgfqpoint{5.836613in}{0.716616in}}%
\pgfpathlineto{\pgfqpoint{5.839969in}{0.726911in}}%
\pgfpathlineto{\pgfqpoint{5.843325in}{0.733483in}}%
\pgfpathlineto{\pgfqpoint{5.846681in}{0.734040in}}%
\pgfpathlineto{\pgfqpoint{5.850036in}{0.728750in}}%
\pgfpathlineto{\pgfqpoint{5.853392in}{0.727905in}}%
\pgfpathlineto{\pgfqpoint{5.856748in}{0.732651in}}%
\pgfpathlineto{\pgfqpoint{5.860104in}{0.734888in}}%
\pgfpathlineto{\pgfqpoint{5.863460in}{0.727711in}}%
\pgfpathlineto{\pgfqpoint{5.866815in}{0.718724in}}%
\pgfpathlineto{\pgfqpoint{5.870171in}{0.717668in}}%
\pgfpathlineto{\pgfqpoint{5.873527in}{0.721589in}}%
\pgfpathlineto{\pgfqpoint{5.876883in}{0.723862in}}%
\pgfpathlineto{\pgfqpoint{5.883594in}{0.722492in}}%
\pgfpathlineto{\pgfqpoint{5.890306in}{0.739107in}}%
\pgfpathlineto{\pgfqpoint{5.893662in}{0.738817in}}%
\pgfpathlineto{\pgfqpoint{5.897017in}{0.724570in}}%
\pgfpathlineto{\pgfqpoint{5.900373in}{0.697791in}}%
\pgfpathlineto{\pgfqpoint{5.903729in}{0.685889in}}%
\pgfpathlineto{\pgfqpoint{5.907085in}{0.703970in}}%
\pgfpathlineto{\pgfqpoint{5.910441in}{0.710452in}}%
\pgfpathlineto{\pgfqpoint{5.917152in}{0.704827in}}%
\pgfpathlineto{\pgfqpoint{5.923864in}{0.727300in}}%
\pgfpathlineto{\pgfqpoint{5.927219in}{0.728419in}}%
\pgfpathlineto{\pgfqpoint{5.930575in}{0.718730in}}%
\pgfpathlineto{\pgfqpoint{5.937287in}{0.678236in}}%
\pgfpathlineto{\pgfqpoint{5.940643in}{0.686353in}}%
\pgfpathlineto{\pgfqpoint{5.943998in}{0.697569in}}%
\pgfpathlineto{\pgfqpoint{5.947354in}{0.702344in}}%
\pgfpathlineto{\pgfqpoint{5.950710in}{0.702471in}}%
\pgfpathlineto{\pgfqpoint{5.954066in}{0.705710in}}%
\pgfpathlineto{\pgfqpoint{5.957421in}{0.713662in}}%
\pgfpathlineto{\pgfqpoint{5.960777in}{0.716401in}}%
\pgfpathlineto{\pgfqpoint{5.964133in}{0.713645in}}%
\pgfpathlineto{\pgfqpoint{5.967489in}{0.709371in}}%
\pgfpathlineto{\pgfqpoint{5.974200in}{0.703469in}}%
\pgfpathlineto{\pgfqpoint{5.977556in}{0.702270in}}%
\pgfpathlineto{\pgfqpoint{5.980912in}{0.708124in}}%
\pgfpathlineto{\pgfqpoint{5.984268in}{0.717026in}}%
\pgfpathlineto{\pgfqpoint{5.987623in}{0.721571in}}%
\pgfpathlineto{\pgfqpoint{5.990979in}{0.721160in}}%
\pgfpathlineto{\pgfqpoint{5.994335in}{0.719533in}}%
\pgfpathlineto{\pgfqpoint{5.997691in}{0.714330in}}%
\pgfpathlineto{\pgfqpoint{6.001047in}{0.712016in}}%
\pgfpathlineto{\pgfqpoint{6.007758in}{0.724267in}}%
\pgfpathlineto{\pgfqpoint{6.011114in}{0.725297in}}%
\pgfpathlineto{\pgfqpoint{6.014470in}{0.722760in}}%
\pgfpathlineto{\pgfqpoint{6.017825in}{0.721854in}}%
\pgfpathlineto{\pgfqpoint{6.027893in}{0.735868in}}%
\pgfpathlineto{\pgfqpoint{6.031249in}{0.738296in}}%
\pgfpathlineto{\pgfqpoint{6.034604in}{0.730443in}}%
\pgfpathlineto{\pgfqpoint{6.037960in}{0.713251in}}%
\pgfpathlineto{\pgfqpoint{6.041316in}{0.702792in}}%
\pgfpathlineto{\pgfqpoint{6.044672in}{0.717613in}}%
\pgfpathlineto{\pgfqpoint{6.048027in}{0.727899in}}%
\pgfpathlineto{\pgfqpoint{6.051383in}{0.730621in}}%
\pgfpathlineto{\pgfqpoint{6.054739in}{0.727415in}}%
\pgfpathlineto{\pgfqpoint{6.058095in}{0.721191in}}%
\pgfpathlineto{\pgfqpoint{6.061451in}{0.727802in}}%
\pgfpathlineto{\pgfqpoint{6.064806in}{0.738387in}}%
\pgfpathlineto{\pgfqpoint{6.068162in}{0.739785in}}%
\pgfpathlineto{\pgfqpoint{6.071518in}{0.735169in}}%
\pgfpathlineto{\pgfqpoint{6.074874in}{0.729349in}}%
\pgfpathlineto{\pgfqpoint{6.078230in}{0.719939in}}%
\pgfpathlineto{\pgfqpoint{6.081585in}{0.708019in}}%
\pgfpathlineto{\pgfqpoint{6.084941in}{0.707137in}}%
\pgfpathlineto{\pgfqpoint{6.091653in}{0.717185in}}%
\pgfpathlineto{\pgfqpoint{6.095008in}{0.721827in}}%
\pgfpathlineto{\pgfqpoint{6.098364in}{0.722734in}}%
\pgfpathlineto{\pgfqpoint{6.101720in}{0.719365in}}%
\pgfpathlineto{\pgfqpoint{6.105076in}{0.721795in}}%
\pgfpathlineto{\pgfqpoint{6.108432in}{0.728200in}}%
\pgfpathlineto{\pgfqpoint{6.111787in}{0.729449in}}%
\pgfpathlineto{\pgfqpoint{6.115143in}{0.722171in}}%
\pgfpathlineto{\pgfqpoint{6.118499in}{0.707097in}}%
\pgfpathlineto{\pgfqpoint{6.121855in}{0.695904in}}%
\pgfpathlineto{\pgfqpoint{6.125210in}{0.697218in}}%
\pgfpathlineto{\pgfqpoint{6.128566in}{0.702807in}}%
\pgfpathlineto{\pgfqpoint{6.131922in}{0.702530in}}%
\pgfpathlineto{\pgfqpoint{6.135278in}{0.703952in}}%
\pgfpathlineto{\pgfqpoint{6.141989in}{0.716535in}}%
\pgfpathlineto{\pgfqpoint{6.145345in}{0.716839in}}%
\pgfpathlineto{\pgfqpoint{6.148701in}{0.709571in}}%
\pgfpathlineto{\pgfqpoint{6.152057in}{0.698225in}}%
\pgfpathlineto{\pgfqpoint{6.155412in}{0.702725in}}%
\pgfpathlineto{\pgfqpoint{6.158768in}{0.712185in}}%
\pgfpathlineto{\pgfqpoint{6.162124in}{0.712965in}}%
\pgfpathlineto{\pgfqpoint{6.168836in}{0.695468in}}%
\pgfpathlineto{\pgfqpoint{6.175547in}{0.722351in}}%
\pgfpathlineto{\pgfqpoint{6.178903in}{0.726241in}}%
\pgfpathlineto{\pgfqpoint{6.182259in}{0.722360in}}%
\pgfpathlineto{\pgfqpoint{6.185614in}{0.709856in}}%
\pgfpathlineto{\pgfqpoint{6.188970in}{0.702068in}}%
\pgfpathlineto{\pgfqpoint{6.192326in}{0.718424in}}%
\pgfpathlineto{\pgfqpoint{6.195682in}{0.731861in}}%
\pgfpathlineto{\pgfqpoint{6.199038in}{0.735320in}}%
\pgfpathlineto{\pgfqpoint{6.202393in}{0.729268in}}%
\pgfpathlineto{\pgfqpoint{6.205749in}{0.714289in}}%
\pgfpathlineto{\pgfqpoint{6.209105in}{0.709728in}}%
\pgfpathlineto{\pgfqpoint{6.212461in}{0.718177in}}%
\pgfpathlineto{\pgfqpoint{6.215816in}{0.718864in}}%
\pgfpathlineto{\pgfqpoint{6.219172in}{0.708830in}}%
\pgfpathlineto{\pgfqpoint{6.222528in}{0.691615in}}%
\pgfpathlineto{\pgfqpoint{6.229240in}{0.729823in}}%
\pgfpathlineto{\pgfqpoint{6.232595in}{0.736172in}}%
\pgfpathlineto{\pgfqpoint{6.235951in}{0.730401in}}%
\pgfpathlineto{\pgfqpoint{6.242663in}{0.695127in}}%
\pgfpathlineto{\pgfqpoint{6.249374in}{0.723885in}}%
\pgfpathlineto{\pgfqpoint{6.252730in}{0.728910in}}%
\pgfpathlineto{\pgfqpoint{6.256086in}{0.725801in}}%
\pgfpathlineto{\pgfqpoint{6.259442in}{0.720504in}}%
\pgfpathlineto{\pgfqpoint{6.262797in}{0.716732in}}%
\pgfpathlineto{\pgfqpoint{6.266153in}{0.718014in}}%
\pgfpathlineto{\pgfqpoint{6.269509in}{0.720455in}}%
\pgfpathlineto{\pgfqpoint{6.272865in}{0.720914in}}%
\pgfpathlineto{\pgfqpoint{6.276221in}{0.718668in}}%
\pgfpathlineto{\pgfqpoint{6.279576in}{0.714735in}}%
\pgfpathlineto{\pgfqpoint{6.282932in}{0.715515in}}%
\pgfpathlineto{\pgfqpoint{6.289644in}{0.730598in}}%
\pgfpathlineto{\pgfqpoint{6.292999in}{0.734403in}}%
\pgfpathlineto{\pgfqpoint{6.296355in}{0.734465in}}%
\pgfpathlineto{\pgfqpoint{6.299711in}{0.728164in}}%
\pgfpathlineto{\pgfqpoint{6.303067in}{0.716555in}}%
\pgfpathlineto{\pgfqpoint{6.306423in}{0.700094in}}%
\pgfpathlineto{\pgfqpoint{6.309778in}{0.696253in}}%
\pgfpathlineto{\pgfqpoint{6.313134in}{0.708073in}}%
\pgfpathlineto{\pgfqpoint{6.316490in}{0.716273in}}%
\pgfpathlineto{\pgfqpoint{6.319846in}{0.719953in}}%
\pgfpathlineto{\pgfqpoint{6.326557in}{0.721796in}}%
\pgfpathlineto{\pgfqpoint{6.333269in}{0.726913in}}%
\pgfpathlineto{\pgfqpoint{6.336625in}{0.727329in}}%
\pgfpathlineto{\pgfqpoint{6.339980in}{0.725229in}}%
\pgfpathlineto{\pgfqpoint{6.343336in}{0.713289in}}%
\pgfpathlineto{\pgfqpoint{6.346692in}{0.704514in}}%
\pgfpathlineto{\pgfqpoint{6.353403in}{0.723190in}}%
\pgfpathlineto{\pgfqpoint{6.356759in}{0.724017in}}%
\pgfpathlineto{\pgfqpoint{6.363471in}{0.709760in}}%
\pgfpathlineto{\pgfqpoint{6.366827in}{0.720387in}}%
\pgfpathlineto{\pgfqpoint{6.370182in}{0.733638in}}%
\pgfpathlineto{\pgfqpoint{6.373538in}{0.738254in}}%
\pgfpathlineto{\pgfqpoint{6.376894in}{0.732327in}}%
\pgfpathlineto{\pgfqpoint{6.383605in}{0.708423in}}%
\pgfpathlineto{\pgfqpoint{6.383605in}{0.708423in}}%
\pgfusepath{stroke}%
\end{pgfscope}%
\begin{pgfscope}%
\pgfpathrectangle{\pgfqpoint{4.621821in}{0.320625in}}{\pgfqpoint{1.845679in}{1.164375in}}%
\pgfusepath{clip}%
\pgfsetrectcap%
\pgfsetroundjoin%
\pgfsetlinewidth{0.602250pt}%
\definecolor{currentstroke}{rgb}{0.000000,0.529412,0.423529}%
\pgfsetstrokecolor{currentstroke}%
\pgfsetstrokeopacity{0.120000}%
\pgfsetdash{}{0pt}%
\pgfpathmoveto{\pgfqpoint{4.705715in}{1.468694in}}%
\pgfpathlineto{\pgfqpoint{4.719139in}{1.443565in}}%
\pgfpathlineto{\pgfqpoint{4.729206in}{1.422261in}}%
\pgfpathlineto{\pgfqpoint{4.749341in}{1.376161in}}%
\pgfpathlineto{\pgfqpoint{4.776187in}{1.315436in}}%
\pgfpathlineto{\pgfqpoint{4.799677in}{1.265399in}}%
\pgfpathlineto{\pgfqpoint{4.826524in}{1.211211in}}%
\pgfpathlineto{\pgfqpoint{4.856726in}{1.152950in}}%
\pgfpathlineto{\pgfqpoint{4.890283in}{1.090600in}}%
\pgfpathlineto{\pgfqpoint{4.927197in}{1.024254in}}%
\pgfpathlineto{\pgfqpoint{4.967466in}{0.954099in}}%
\pgfpathlineto{\pgfqpoint{5.011091in}{0.880300in}}%
\pgfpathlineto{\pgfqpoint{5.048005in}{0.819148in}}%
\pgfpathlineto{\pgfqpoint{5.074851in}{0.775440in}}%
\pgfpathlineto{\pgfqpoint{5.084919in}{0.758470in}}%
\pgfpathlineto{\pgfqpoint{5.091630in}{0.749057in}}%
\pgfpathlineto{\pgfqpoint{5.101698in}{0.737119in}}%
\pgfpathlineto{\pgfqpoint{5.108409in}{0.726010in}}%
\pgfpathlineto{\pgfqpoint{5.115121in}{0.711964in}}%
\pgfpathlineto{\pgfqpoint{5.128544in}{0.681915in}}%
\pgfpathlineto{\pgfqpoint{5.131900in}{0.676335in}}%
\pgfpathlineto{\pgfqpoint{5.135255in}{0.672485in}}%
\pgfpathlineto{\pgfqpoint{5.138611in}{0.670882in}}%
\pgfpathlineto{\pgfqpoint{5.141967in}{0.671709in}}%
\pgfpathlineto{\pgfqpoint{5.152034in}{0.679429in}}%
\pgfpathlineto{\pgfqpoint{5.155390in}{0.679232in}}%
\pgfpathlineto{\pgfqpoint{5.158746in}{0.675024in}}%
\pgfpathlineto{\pgfqpoint{5.165457in}{0.656851in}}%
\pgfpathlineto{\pgfqpoint{5.168813in}{0.657362in}}%
\pgfpathlineto{\pgfqpoint{5.175525in}{0.679925in}}%
\pgfpathlineto{\pgfqpoint{5.178880in}{0.687911in}}%
\pgfpathlineto{\pgfqpoint{5.182236in}{0.692947in}}%
\pgfpathlineto{\pgfqpoint{5.185592in}{0.695407in}}%
\pgfpathlineto{\pgfqpoint{5.188948in}{0.695436in}}%
\pgfpathlineto{\pgfqpoint{5.192304in}{0.692888in}}%
\pgfpathlineto{\pgfqpoint{5.195659in}{0.687191in}}%
\pgfpathlineto{\pgfqpoint{5.199015in}{0.676867in}}%
\pgfpathlineto{\pgfqpoint{5.202371in}{0.657560in}}%
\pgfpathlineto{\pgfqpoint{5.205727in}{0.607973in}}%
\pgfpathlineto{\pgfqpoint{5.209082in}{0.644755in}}%
\pgfpathlineto{\pgfqpoint{5.212438in}{0.674746in}}%
\pgfpathlineto{\pgfqpoint{5.215794in}{0.690639in}}%
\pgfpathlineto{\pgfqpoint{5.219150in}{0.700375in}}%
\pgfpathlineto{\pgfqpoint{5.222506in}{0.706055in}}%
\pgfpathlineto{\pgfqpoint{5.225861in}{0.708190in}}%
\pgfpathlineto{\pgfqpoint{5.229217in}{0.706303in}}%
\pgfpathlineto{\pgfqpoint{5.232573in}{0.698515in}}%
\pgfpathlineto{\pgfqpoint{5.235929in}{0.684674in}}%
\pgfpathlineto{\pgfqpoint{5.239285in}{0.678482in}}%
\pgfpathlineto{\pgfqpoint{5.245996in}{0.698279in}}%
\pgfpathlineto{\pgfqpoint{5.249352in}{0.700874in}}%
\pgfpathlineto{\pgfqpoint{5.252708in}{0.696420in}}%
\pgfpathlineto{\pgfqpoint{5.256063in}{0.681747in}}%
\pgfpathlineto{\pgfqpoint{5.259419in}{0.642148in}}%
\pgfpathlineto{\pgfqpoint{5.266131in}{0.697265in}}%
\pgfpathlineto{\pgfqpoint{5.269487in}{0.710040in}}%
\pgfpathlineto{\pgfqpoint{5.272842in}{0.712372in}}%
\pgfpathlineto{\pgfqpoint{5.276198in}{0.700267in}}%
\pgfpathlineto{\pgfqpoint{5.279554in}{0.653958in}}%
\pgfpathlineto{\pgfqpoint{5.282910in}{0.691144in}}%
\pgfpathlineto{\pgfqpoint{5.286265in}{0.717286in}}%
\pgfpathlineto{\pgfqpoint{5.289621in}{0.727987in}}%
\pgfpathlineto{\pgfqpoint{5.292977in}{0.730300in}}%
\pgfpathlineto{\pgfqpoint{5.296333in}{0.724949in}}%
\pgfpathlineto{\pgfqpoint{5.299689in}{0.707997in}}%
\pgfpathlineto{\pgfqpoint{5.303044in}{0.641852in}}%
\pgfpathlineto{\pgfqpoint{5.306400in}{0.699429in}}%
\pgfpathlineto{\pgfqpoint{5.309756in}{0.728153in}}%
\pgfpathlineto{\pgfqpoint{5.313112in}{0.740031in}}%
\pgfpathlineto{\pgfqpoint{5.316467in}{0.741217in}}%
\pgfpathlineto{\pgfqpoint{5.319823in}{0.729040in}}%
\pgfpathlineto{\pgfqpoint{5.323179in}{0.700368in}}%
\pgfpathlineto{\pgfqpoint{5.326535in}{0.643102in}}%
\pgfpathlineto{\pgfqpoint{5.329891in}{0.698305in}}%
\pgfpathlineto{\pgfqpoint{5.333246in}{0.710104in}}%
\pgfpathlineto{\pgfqpoint{5.336602in}{0.705462in}}%
\pgfpathlineto{\pgfqpoint{5.339958in}{0.676769in}}%
\pgfpathlineto{\pgfqpoint{5.343314in}{0.667153in}}%
\pgfpathlineto{\pgfqpoint{5.346669in}{0.714025in}}%
\pgfpathlineto{\pgfqpoint{5.350025in}{0.730385in}}%
\pgfpathlineto{\pgfqpoint{5.353381in}{0.732344in}}%
\pgfpathlineto{\pgfqpoint{5.356737in}{0.718226in}}%
\pgfpathlineto{\pgfqpoint{5.363448in}{0.647209in}}%
\pgfpathlineto{\pgfqpoint{5.366804in}{0.694069in}}%
\pgfpathlineto{\pgfqpoint{5.370160in}{0.703161in}}%
\pgfpathlineto{\pgfqpoint{5.373516in}{0.696243in}}%
\pgfpathlineto{\pgfqpoint{5.376871in}{0.660512in}}%
\pgfpathlineto{\pgfqpoint{5.380227in}{0.670634in}}%
\pgfpathlineto{\pgfqpoint{5.383583in}{0.709873in}}%
\pgfpathlineto{\pgfqpoint{5.386939in}{0.724087in}}%
\pgfpathlineto{\pgfqpoint{5.390295in}{0.723886in}}%
\pgfpathlineto{\pgfqpoint{5.393650in}{0.705566in}}%
\pgfpathlineto{\pgfqpoint{5.397006in}{0.638239in}}%
\pgfpathlineto{\pgfqpoint{5.400362in}{0.688494in}}%
\pgfpathlineto{\pgfqpoint{5.403718in}{0.712375in}}%
\pgfpathlineto{\pgfqpoint{5.407074in}{0.718797in}}%
\pgfpathlineto{\pgfqpoint{5.410429in}{0.714251in}}%
\pgfpathlineto{\pgfqpoint{5.413785in}{0.693976in}}%
\pgfpathlineto{\pgfqpoint{5.417141in}{0.604116in}}%
\pgfpathlineto{\pgfqpoint{5.420497in}{0.704813in}}%
\pgfpathlineto{\pgfqpoint{5.423852in}{0.727150in}}%
\pgfpathlineto{\pgfqpoint{5.427208in}{0.734542in}}%
\pgfpathlineto{\pgfqpoint{5.430564in}{0.727548in}}%
\pgfpathlineto{\pgfqpoint{5.437276in}{0.666774in}}%
\pgfpathlineto{\pgfqpoint{5.440631in}{0.717498in}}%
\pgfpathlineto{\pgfqpoint{5.443987in}{0.732088in}}%
\pgfpathlineto{\pgfqpoint{5.447343in}{0.734861in}}%
\pgfpathlineto{\pgfqpoint{5.450699in}{0.727629in}}%
\pgfpathlineto{\pgfqpoint{5.454054in}{0.702834in}}%
\pgfpathlineto{\pgfqpoint{5.457410in}{0.659010in}}%
\pgfpathlineto{\pgfqpoint{5.460766in}{0.723577in}}%
\pgfpathlineto{\pgfqpoint{5.464122in}{0.743188in}}%
\pgfpathlineto{\pgfqpoint{5.467478in}{0.749229in}}%
\pgfpathlineto{\pgfqpoint{5.470833in}{0.741254in}}%
\pgfpathlineto{\pgfqpoint{5.474189in}{0.714642in}}%
\pgfpathlineto{\pgfqpoint{5.477545in}{0.638526in}}%
\pgfpathlineto{\pgfqpoint{5.480901in}{0.716832in}}%
\pgfpathlineto{\pgfqpoint{5.484256in}{0.732245in}}%
\pgfpathlineto{\pgfqpoint{5.487612in}{0.733290in}}%
\pgfpathlineto{\pgfqpoint{5.490968in}{0.721039in}}%
\pgfpathlineto{\pgfqpoint{5.494324in}{0.670355in}}%
\pgfpathlineto{\pgfqpoint{5.501035in}{0.739004in}}%
\pgfpathlineto{\pgfqpoint{5.504391in}{0.750825in}}%
\pgfpathlineto{\pgfqpoint{5.507747in}{0.748890in}}%
\pgfpathlineto{\pgfqpoint{5.511103in}{0.732250in}}%
\pgfpathlineto{\pgfqpoint{5.514458in}{0.695210in}}%
\pgfpathlineto{\pgfqpoint{5.517814in}{0.664759in}}%
\pgfpathlineto{\pgfqpoint{5.521170in}{0.706313in}}%
\pgfpathlineto{\pgfqpoint{5.524526in}{0.711794in}}%
\pgfpathlineto{\pgfqpoint{5.527882in}{0.697657in}}%
\pgfpathlineto{\pgfqpoint{5.531237in}{0.527072in}}%
\pgfpathlineto{\pgfqpoint{5.534593in}{0.706691in}}%
\pgfpathlineto{\pgfqpoint{5.537949in}{0.731227in}}%
\pgfpathlineto{\pgfqpoint{5.541305in}{0.738954in}}%
\pgfpathlineto{\pgfqpoint{5.544660in}{0.730825in}}%
\pgfpathlineto{\pgfqpoint{5.548016in}{0.708570in}}%
\pgfpathlineto{\pgfqpoint{5.551372in}{0.647785in}}%
\pgfpathlineto{\pgfqpoint{5.554728in}{0.675651in}}%
\pgfpathlineto{\pgfqpoint{5.558084in}{0.696352in}}%
\pgfpathlineto{\pgfqpoint{5.561439in}{0.694036in}}%
\pgfpathlineto{\pgfqpoint{5.564795in}{0.664702in}}%
\pgfpathlineto{\pgfqpoint{5.568151in}{0.662169in}}%
\pgfpathlineto{\pgfqpoint{5.571507in}{0.706950in}}%
\pgfpathlineto{\pgfqpoint{5.574863in}{0.722787in}}%
\pgfpathlineto{\pgfqpoint{5.578218in}{0.723760in}}%
\pgfpathlineto{\pgfqpoint{5.581574in}{0.707349in}}%
\pgfpathlineto{\pgfqpoint{5.584930in}{0.659870in}}%
\pgfpathlineto{\pgfqpoint{5.588286in}{0.672093in}}%
\pgfpathlineto{\pgfqpoint{5.591641in}{0.701787in}}%
\pgfpathlineto{\pgfqpoint{5.594997in}{0.708893in}}%
\pgfpathlineto{\pgfqpoint{5.598353in}{0.703291in}}%
\pgfpathlineto{\pgfqpoint{5.601709in}{0.677802in}}%
\pgfpathlineto{\pgfqpoint{5.605065in}{0.647169in}}%
\pgfpathlineto{\pgfqpoint{5.608420in}{0.704297in}}%
\pgfpathlineto{\pgfqpoint{5.611776in}{0.722633in}}%
\pgfpathlineto{\pgfqpoint{5.615132in}{0.726827in}}%
\pgfpathlineto{\pgfqpoint{5.618488in}{0.714446in}}%
\pgfpathlineto{\pgfqpoint{5.621843in}{0.668205in}}%
\pgfpathlineto{\pgfqpoint{5.625199in}{0.686092in}}%
\pgfpathlineto{\pgfqpoint{5.628555in}{0.717596in}}%
\pgfpathlineto{\pgfqpoint{5.631911in}{0.727943in}}%
\pgfpathlineto{\pgfqpoint{5.635267in}{0.728009in}}%
\pgfpathlineto{\pgfqpoint{5.638622in}{0.717492in}}%
\pgfpathlineto{\pgfqpoint{5.641978in}{0.682741in}}%
\pgfpathlineto{\pgfqpoint{5.645334in}{0.682007in}}%
\pgfpathlineto{\pgfqpoint{5.648690in}{0.724076in}}%
\pgfpathlineto{\pgfqpoint{5.652045in}{0.739603in}}%
\pgfpathlineto{\pgfqpoint{5.655401in}{0.742620in}}%
\pgfpathlineto{\pgfqpoint{5.658757in}{0.730261in}}%
\pgfpathlineto{\pgfqpoint{5.662113in}{0.688743in}}%
\pgfpathlineto{\pgfqpoint{5.665469in}{0.693344in}}%
\pgfpathlineto{\pgfqpoint{5.668824in}{0.727718in}}%
\pgfpathlineto{\pgfqpoint{5.672180in}{0.738189in}}%
\pgfpathlineto{\pgfqpoint{5.675536in}{0.737606in}}%
\pgfpathlineto{\pgfqpoint{5.678892in}{0.725201in}}%
\pgfpathlineto{\pgfqpoint{5.682247in}{0.679896in}}%
\pgfpathlineto{\pgfqpoint{5.685603in}{0.705135in}}%
\pgfpathlineto{\pgfqpoint{5.688959in}{0.739329in}}%
\pgfpathlineto{\pgfqpoint{5.692315in}{0.752311in}}%
\pgfpathlineto{\pgfqpoint{5.695671in}{0.752519in}}%
\pgfpathlineto{\pgfqpoint{5.699026in}{0.737614in}}%
\pgfpathlineto{\pgfqpoint{5.705738in}{0.671921in}}%
\pgfpathlineto{\pgfqpoint{5.709094in}{0.715634in}}%
\pgfpathlineto{\pgfqpoint{5.712449in}{0.723844in}}%
\pgfpathlineto{\pgfqpoint{5.715805in}{0.716054in}}%
\pgfpathlineto{\pgfqpoint{5.719161in}{0.676903in}}%
\pgfpathlineto{\pgfqpoint{5.722517in}{0.694786in}}%
\pgfpathlineto{\pgfqpoint{5.725873in}{0.731665in}}%
\pgfpathlineto{\pgfqpoint{5.729228in}{0.745076in}}%
\pgfpathlineto{\pgfqpoint{5.732584in}{0.743952in}}%
\pgfpathlineto{\pgfqpoint{5.735940in}{0.728103in}}%
\pgfpathlineto{\pgfqpoint{5.739296in}{0.695699in}}%
\pgfpathlineto{\pgfqpoint{5.742652in}{0.626950in}}%
\pgfpathlineto{\pgfqpoint{5.746007in}{0.690701in}}%
\pgfpathlineto{\pgfqpoint{5.749363in}{0.697038in}}%
\pgfpathlineto{\pgfqpoint{5.752719in}{0.679264in}}%
\pgfpathlineto{\pgfqpoint{5.756075in}{0.632252in}}%
\pgfpathlineto{\pgfqpoint{5.759430in}{0.703131in}}%
\pgfpathlineto{\pgfqpoint{5.762786in}{0.723817in}}%
\pgfpathlineto{\pgfqpoint{5.766142in}{0.728843in}}%
\pgfpathlineto{\pgfqpoint{5.769498in}{0.717272in}}%
\pgfpathlineto{\pgfqpoint{5.772854in}{0.688255in}}%
\pgfpathlineto{\pgfqpoint{5.776209in}{0.610943in}}%
\pgfpathlineto{\pgfqpoint{5.779565in}{0.686463in}}%
\pgfpathlineto{\pgfqpoint{5.782921in}{0.698747in}}%
\pgfpathlineto{\pgfqpoint{5.786277in}{0.694394in}}%
\pgfpathlineto{\pgfqpoint{5.789632in}{0.666430in}}%
\pgfpathlineto{\pgfqpoint{5.792988in}{0.653583in}}%
\pgfpathlineto{\pgfqpoint{5.796344in}{0.702110in}}%
\pgfpathlineto{\pgfqpoint{5.799700in}{0.718765in}}%
\pgfpathlineto{\pgfqpoint{5.803056in}{0.720918in}}%
\pgfpathlineto{\pgfqpoint{5.806411in}{0.705130in}}%
\pgfpathlineto{\pgfqpoint{5.809767in}{0.647320in}}%
\pgfpathlineto{\pgfqpoint{5.813123in}{0.684264in}}%
\pgfpathlineto{\pgfqpoint{5.816479in}{0.711032in}}%
\pgfpathlineto{\pgfqpoint{5.819834in}{0.719204in}}%
\pgfpathlineto{\pgfqpoint{5.823190in}{0.716937in}}%
\pgfpathlineto{\pgfqpoint{5.826546in}{0.702024in}}%
\pgfpathlineto{\pgfqpoint{5.829902in}{0.640308in}}%
\pgfpathlineto{\pgfqpoint{5.833258in}{0.692965in}}%
\pgfpathlineto{\pgfqpoint{5.836613in}{0.722210in}}%
\pgfpathlineto{\pgfqpoint{5.839969in}{0.733300in}}%
\pgfpathlineto{\pgfqpoint{5.843325in}{0.731504in}}%
\pgfpathlineto{\pgfqpoint{5.846681in}{0.710404in}}%
\pgfpathlineto{\pgfqpoint{5.850036in}{0.594505in}}%
\pgfpathlineto{\pgfqpoint{5.853392in}{0.712502in}}%
\pgfpathlineto{\pgfqpoint{5.856748in}{0.732451in}}%
\pgfpathlineto{\pgfqpoint{5.860104in}{0.738597in}}%
\pgfpathlineto{\pgfqpoint{5.863460in}{0.735436in}}%
\pgfpathlineto{\pgfqpoint{5.866815in}{0.720196in}}%
\pgfpathlineto{\pgfqpoint{5.870171in}{0.660297in}}%
\pgfpathlineto{\pgfqpoint{5.873527in}{0.709602in}}%
\pgfpathlineto{\pgfqpoint{5.876883in}{0.739295in}}%
\pgfpathlineto{\pgfqpoint{5.880238in}{0.750838in}}%
\pgfpathlineto{\pgfqpoint{5.883594in}{0.750174in}}%
\pgfpathlineto{\pgfqpoint{5.886950in}{0.733775in}}%
\pgfpathlineto{\pgfqpoint{5.890306in}{0.686242in}}%
\pgfpathlineto{\pgfqpoint{5.893662in}{0.698677in}}%
\pgfpathlineto{\pgfqpoint{5.897017in}{0.728380in}}%
\pgfpathlineto{\pgfqpoint{5.900373in}{0.735532in}}%
\pgfpathlineto{\pgfqpoint{5.903729in}{0.730029in}}%
\pgfpathlineto{\pgfqpoint{5.907085in}{0.704884in}}%
\pgfpathlineto{\pgfqpoint{5.910441in}{0.672240in}}%
\pgfpathlineto{\pgfqpoint{5.913796in}{0.730543in}}%
\pgfpathlineto{\pgfqpoint{5.917152in}{0.749068in}}%
\pgfpathlineto{\pgfqpoint{5.920508in}{0.753418in}}%
\pgfpathlineto{\pgfqpoint{5.923864in}{0.742729in}}%
\pgfpathlineto{\pgfqpoint{5.927219in}{0.718699in}}%
\pgfpathlineto{\pgfqpoint{5.930575in}{0.649817in}}%
\pgfpathlineto{\pgfqpoint{5.933931in}{0.688286in}}%
\pgfpathlineto{\pgfqpoint{5.937287in}{0.705523in}}%
\pgfpathlineto{\pgfqpoint{5.940643in}{0.698889in}}%
\pgfpathlineto{\pgfqpoint{5.943998in}{0.646639in}}%
\pgfpathlineto{\pgfqpoint{5.947354in}{0.693240in}}%
\pgfpathlineto{\pgfqpoint{5.950710in}{0.724736in}}%
\pgfpathlineto{\pgfqpoint{5.954066in}{0.736049in}}%
\pgfpathlineto{\pgfqpoint{5.957421in}{0.731756in}}%
\pgfpathlineto{\pgfqpoint{5.960777in}{0.713264in}}%
\pgfpathlineto{\pgfqpoint{5.964133in}{0.672217in}}%
\pgfpathlineto{\pgfqpoint{5.967489in}{0.653181in}}%
\pgfpathlineto{\pgfqpoint{5.970845in}{0.687602in}}%
\pgfpathlineto{\pgfqpoint{5.974200in}{0.689404in}}%
\pgfpathlineto{\pgfqpoint{5.980912in}{0.644391in}}%
\pgfpathlineto{\pgfqpoint{5.984268in}{0.698999in}}%
\pgfpathlineto{\pgfqpoint{5.987623in}{0.717131in}}%
\pgfpathlineto{\pgfqpoint{5.990979in}{0.720127in}}%
\pgfpathlineto{\pgfqpoint{5.994335in}{0.705639in}}%
\pgfpathlineto{\pgfqpoint{5.997691in}{0.662803in}}%
\pgfpathlineto{\pgfqpoint{6.001047in}{0.665813in}}%
\pgfpathlineto{\pgfqpoint{6.004402in}{0.698939in}}%
\pgfpathlineto{\pgfqpoint{6.007758in}{0.707653in}}%
\pgfpathlineto{\pgfqpoint{6.011114in}{0.704138in}}%
\pgfpathlineto{\pgfqpoint{6.014470in}{0.684453in}}%
\pgfpathlineto{\pgfqpoint{6.017825in}{0.593781in}}%
\pgfpathlineto{\pgfqpoint{6.021181in}{0.695622in}}%
\pgfpathlineto{\pgfqpoint{6.024537in}{0.717978in}}%
\pgfpathlineto{\pgfqpoint{6.027893in}{0.725102in}}%
\pgfpathlineto{\pgfqpoint{6.031249in}{0.717083in}}%
\pgfpathlineto{\pgfqpoint{6.034604in}{0.681375in}}%
\pgfpathlineto{\pgfqpoint{6.037960in}{0.676841in}}%
\pgfpathlineto{\pgfqpoint{6.041316in}{0.716453in}}%
\pgfpathlineto{\pgfqpoint{6.044672in}{0.729647in}}%
\pgfpathlineto{\pgfqpoint{6.048027in}{0.732316in}}%
\pgfpathlineto{\pgfqpoint{6.051383in}{0.725862in}}%
\pgfpathlineto{\pgfqpoint{6.054739in}{0.704389in}}%
\pgfpathlineto{\pgfqpoint{6.058095in}{0.617248in}}%
\pgfpathlineto{\pgfqpoint{6.061451in}{0.714879in}}%
\pgfpathlineto{\pgfqpoint{6.064806in}{0.737214in}}%
\pgfpathlineto{\pgfqpoint{6.068162in}{0.745300in}}%
\pgfpathlineto{\pgfqpoint{6.071518in}{0.740624in}}%
\pgfpathlineto{\pgfqpoint{6.074874in}{0.716633in}}%
\pgfpathlineto{\pgfqpoint{6.078230in}{0.634482in}}%
\pgfpathlineto{\pgfqpoint{6.081585in}{0.720923in}}%
\pgfpathlineto{\pgfqpoint{6.084941in}{0.738707in}}%
\pgfpathlineto{\pgfqpoint{6.088297in}{0.742931in}}%
\pgfpathlineto{\pgfqpoint{6.091653in}{0.736905in}}%
\pgfpathlineto{\pgfqpoint{6.095008in}{0.714446in}}%
\pgfpathlineto{\pgfqpoint{6.098364in}{0.651839in}}%
\pgfpathlineto{\pgfqpoint{6.101720in}{0.729917in}}%
\pgfpathlineto{\pgfqpoint{6.105076in}{0.750876in}}%
\pgfpathlineto{\pgfqpoint{6.108432in}{0.757640in}}%
\pgfpathlineto{\pgfqpoint{6.111787in}{0.750514in}}%
\pgfpathlineto{\pgfqpoint{6.115143in}{0.729099in}}%
\pgfpathlineto{\pgfqpoint{6.118499in}{0.669663in}}%
\pgfpathlineto{\pgfqpoint{6.121855in}{0.697059in}}%
\pgfpathlineto{\pgfqpoint{6.125210in}{0.719149in}}%
\pgfpathlineto{\pgfqpoint{6.128566in}{0.719084in}}%
\pgfpathlineto{\pgfqpoint{6.131922in}{0.697819in}}%
\pgfpathlineto{\pgfqpoint{6.135278in}{0.657446in}}%
\pgfpathlineto{\pgfqpoint{6.138634in}{0.722180in}}%
\pgfpathlineto{\pgfqpoint{6.141989in}{0.741753in}}%
\pgfpathlineto{\pgfqpoint{6.145345in}{0.746149in}}%
\pgfpathlineto{\pgfqpoint{6.148701in}{0.734872in}}%
\pgfpathlineto{\pgfqpoint{6.152057in}{0.711439in}}%
\pgfpathlineto{\pgfqpoint{6.155412in}{0.653827in}}%
\pgfpathlineto{\pgfqpoint{6.158768in}{0.669628in}}%
\pgfpathlineto{\pgfqpoint{6.162124in}{0.688310in}}%
\pgfpathlineto{\pgfqpoint{6.165480in}{0.677512in}}%
\pgfpathlineto{\pgfqpoint{6.168836in}{0.569859in}}%
\pgfpathlineto{\pgfqpoint{6.172191in}{0.693277in}}%
\pgfpathlineto{\pgfqpoint{6.175547in}{0.717555in}}%
\pgfpathlineto{\pgfqpoint{6.178903in}{0.725075in}}%
\pgfpathlineto{\pgfqpoint{6.182259in}{0.716067in}}%
\pgfpathlineto{\pgfqpoint{6.185614in}{0.690624in}}%
\pgfpathlineto{\pgfqpoint{6.188970in}{0.547750in}}%
\pgfpathlineto{\pgfqpoint{6.192326in}{0.678896in}}%
\pgfpathlineto{\pgfqpoint{6.195682in}{0.694602in}}%
\pgfpathlineto{\pgfqpoint{6.199038in}{0.692724in}}%
\pgfpathlineto{\pgfqpoint{6.202393in}{0.670790in}}%
\pgfpathlineto{\pgfqpoint{6.205749in}{0.628328in}}%
\pgfpathlineto{\pgfqpoint{6.209105in}{0.693819in}}%
\pgfpathlineto{\pgfqpoint{6.212461in}{0.713403in}}%
\pgfpathlineto{\pgfqpoint{6.215816in}{0.718002in}}%
\pgfpathlineto{\pgfqpoint{6.219172in}{0.705298in}}%
\pgfpathlineto{\pgfqpoint{6.222528in}{0.656526in}}%
\pgfpathlineto{\pgfqpoint{6.225884in}{0.679835in}}%
\pgfpathlineto{\pgfqpoint{6.229240in}{0.710200in}}%
\pgfpathlineto{\pgfqpoint{6.232595in}{0.720304in}}%
\pgfpathlineto{\pgfqpoint{6.235951in}{0.720308in}}%
\pgfpathlineto{\pgfqpoint{6.239307in}{0.709852in}}%
\pgfpathlineto{\pgfqpoint{6.242663in}{0.675626in}}%
\pgfpathlineto{\pgfqpoint{6.246019in}{0.673128in}}%
\pgfpathlineto{\pgfqpoint{6.249374in}{0.715862in}}%
\pgfpathlineto{\pgfqpoint{6.252730in}{0.731584in}}%
\pgfpathlineto{\pgfqpoint{6.256086in}{0.734831in}}%
\pgfpathlineto{\pgfqpoint{6.259442in}{0.722297in}}%
\pgfpathlineto{\pgfqpoint{6.262797in}{0.669515in}}%
\pgfpathlineto{\pgfqpoint{6.269509in}{0.731657in}}%
\pgfpathlineto{\pgfqpoint{6.272865in}{0.742031in}}%
\pgfpathlineto{\pgfqpoint{6.276221in}{0.742875in}}%
\pgfpathlineto{\pgfqpoint{6.279576in}{0.734332in}}%
\pgfpathlineto{\pgfqpoint{6.282932in}{0.707713in}}%
\pgfpathlineto{\pgfqpoint{6.286288in}{0.672444in}}%
\pgfpathlineto{\pgfqpoint{6.289644in}{0.731370in}}%
\pgfpathlineto{\pgfqpoint{6.292999in}{0.750365in}}%
\pgfpathlineto{\pgfqpoint{6.296355in}{0.756436in}}%
\pgfpathlineto{\pgfqpoint{6.299711in}{0.749241in}}%
\pgfpathlineto{\pgfqpoint{6.303067in}{0.725079in}}%
\pgfpathlineto{\pgfqpoint{6.306423in}{0.551771in}}%
\pgfpathlineto{\pgfqpoint{6.309778in}{0.717251in}}%
\pgfpathlineto{\pgfqpoint{6.313134in}{0.734656in}}%
\pgfpathlineto{\pgfqpoint{6.316490in}{0.736065in}}%
\pgfpathlineto{\pgfqpoint{6.319846in}{0.723272in}}%
\pgfpathlineto{\pgfqpoint{6.323201in}{0.665299in}}%
\pgfpathlineto{\pgfqpoint{6.326557in}{0.714417in}}%
\pgfpathlineto{\pgfqpoint{6.329913in}{0.744342in}}%
\pgfpathlineto{\pgfqpoint{6.333269in}{0.755216in}}%
\pgfpathlineto{\pgfqpoint{6.336625in}{0.751821in}}%
\pgfpathlineto{\pgfqpoint{6.339980in}{0.735174in}}%
\pgfpathlineto{\pgfqpoint{6.343336in}{0.702698in}}%
\pgfpathlineto{\pgfqpoint{6.346692in}{0.620118in}}%
\pgfpathlineto{\pgfqpoint{6.350048in}{0.691847in}}%
\pgfpathlineto{\pgfqpoint{6.353403in}{0.695587in}}%
\pgfpathlineto{\pgfqpoint{6.356759in}{0.665796in}}%
\pgfpathlineto{\pgfqpoint{6.360115in}{0.674449in}}%
\pgfpathlineto{\pgfqpoint{6.363471in}{0.716540in}}%
\pgfpathlineto{\pgfqpoint{6.366827in}{0.731966in}}%
\pgfpathlineto{\pgfqpoint{6.370182in}{0.732155in}}%
\pgfpathlineto{\pgfqpoint{6.383605in}{0.688136in}}%
\pgfpathlineto{\pgfqpoint{6.383605in}{0.688136in}}%
\pgfusepath{stroke}%
\end{pgfscope}%
\begin{pgfscope}%
\pgfpathrectangle{\pgfqpoint{4.621821in}{0.320625in}}{\pgfqpoint{1.845679in}{1.164375in}}%
\pgfusepath{clip}%
\pgfsetrectcap%
\pgfsetroundjoin%
\pgfsetlinewidth{0.602250pt}%
\definecolor{currentstroke}{rgb}{0.000000,0.529412,0.423529}%
\pgfsetstrokecolor{currentstroke}%
\pgfsetstrokeopacity{0.120000}%
\pgfsetdash{}{0pt}%
\pgfpathmoveto{\pgfqpoint{4.705715in}{1.469525in}}%
\pgfpathlineto{\pgfqpoint{4.719139in}{1.445279in}}%
\pgfpathlineto{\pgfqpoint{4.729206in}{1.424522in}}%
\pgfpathlineto{\pgfqpoint{4.745985in}{1.386302in}}%
\pgfpathlineto{\pgfqpoint{4.772831in}{1.324674in}}%
\pgfpathlineto{\pgfqpoint{4.792966in}{1.281277in}}%
\pgfpathlineto{\pgfqpoint{4.816456in}{1.233644in}}%
\pgfpathlineto{\pgfqpoint{4.843302in}{1.182095in}}%
\pgfpathlineto{\pgfqpoint{4.873504in}{1.126593in}}%
\pgfpathlineto{\pgfqpoint{4.910418in}{1.061125in}}%
\pgfpathlineto{\pgfqpoint{4.954043in}{0.985985in}}%
\pgfpathlineto{\pgfqpoint{5.004380in}{0.901341in}}%
\pgfpathlineto{\pgfqpoint{5.064784in}{0.802356in}}%
\pgfpathlineto{\pgfqpoint{5.078207in}{0.779425in}}%
\pgfpathlineto{\pgfqpoint{5.088274in}{0.762505in}}%
\pgfpathlineto{\pgfqpoint{5.094986in}{0.752895in}}%
\pgfpathlineto{\pgfqpoint{5.105053in}{0.739785in}}%
\pgfpathlineto{\pgfqpoint{5.111765in}{0.727982in}}%
\pgfpathlineto{\pgfqpoint{5.121832in}{0.706505in}}%
\pgfpathlineto{\pgfqpoint{5.128544in}{0.692092in}}%
\pgfpathlineto{\pgfqpoint{5.131900in}{0.686091in}}%
\pgfpathlineto{\pgfqpoint{5.135255in}{0.681577in}}%
\pgfpathlineto{\pgfqpoint{5.138611in}{0.678916in}}%
\pgfpathlineto{\pgfqpoint{5.141967in}{0.678125in}}%
\pgfpathlineto{\pgfqpoint{5.152034in}{0.680368in}}%
\pgfpathlineto{\pgfqpoint{5.155390in}{0.679049in}}%
\pgfpathlineto{\pgfqpoint{5.158746in}{0.674274in}}%
\pgfpathlineto{\pgfqpoint{5.162102in}{0.666389in}}%
\pgfpathlineto{\pgfqpoint{5.168813in}{0.644085in}}%
\pgfpathlineto{\pgfqpoint{5.172169in}{0.638524in}}%
\pgfpathlineto{\pgfqpoint{5.175525in}{0.643772in}}%
\pgfpathlineto{\pgfqpoint{5.178880in}{0.651626in}}%
\pgfpathlineto{\pgfqpoint{5.182236in}{0.657236in}}%
\pgfpathlineto{\pgfqpoint{5.185592in}{0.659875in}}%
\pgfpathlineto{\pgfqpoint{5.188948in}{0.659492in}}%
\pgfpathlineto{\pgfqpoint{5.192304in}{0.655764in}}%
\pgfpathlineto{\pgfqpoint{5.195659in}{0.647614in}}%
\pgfpathlineto{\pgfqpoint{5.199015in}{0.631841in}}%
\pgfpathlineto{\pgfqpoint{5.202371in}{0.594713in}}%
\pgfpathlineto{\pgfqpoint{5.205727in}{0.602358in}}%
\pgfpathlineto{\pgfqpoint{5.209082in}{0.639255in}}%
\pgfpathlineto{\pgfqpoint{5.212438in}{0.657410in}}%
\pgfpathlineto{\pgfqpoint{5.215794in}{0.668344in}}%
\pgfpathlineto{\pgfqpoint{5.219150in}{0.674796in}}%
\pgfpathlineto{\pgfqpoint{5.222506in}{0.677470in}}%
\pgfpathlineto{\pgfqpoint{5.225861in}{0.675956in}}%
\pgfpathlineto{\pgfqpoint{5.229217in}{0.668278in}}%
\pgfpathlineto{\pgfqpoint{5.232573in}{0.659187in}}%
\pgfpathlineto{\pgfqpoint{5.235929in}{0.672012in}}%
\pgfpathlineto{\pgfqpoint{5.239285in}{0.689221in}}%
\pgfpathlineto{\pgfqpoint{5.242640in}{0.700482in}}%
\pgfpathlineto{\pgfqpoint{5.245996in}{0.706671in}}%
\pgfpathlineto{\pgfqpoint{5.249352in}{0.708652in}}%
\pgfpathlineto{\pgfqpoint{5.252708in}{0.706503in}}%
\pgfpathlineto{\pgfqpoint{5.256063in}{0.699295in}}%
\pgfpathlineto{\pgfqpoint{5.259419in}{0.683671in}}%
\pgfpathlineto{\pgfqpoint{5.262775in}{0.642406in}}%
\pgfpathlineto{\pgfqpoint{5.266131in}{0.654908in}}%
\pgfpathlineto{\pgfqpoint{5.269487in}{0.691951in}}%
\pgfpathlineto{\pgfqpoint{5.272842in}{0.708648in}}%
\pgfpathlineto{\pgfqpoint{5.276198in}{0.717049in}}%
\pgfpathlineto{\pgfqpoint{5.279554in}{0.719037in}}%
\pgfpathlineto{\pgfqpoint{5.282910in}{0.712888in}}%
\pgfpathlineto{\pgfqpoint{5.286265in}{0.694043in}}%
\pgfpathlineto{\pgfqpoint{5.289621in}{0.654025in}}%
\pgfpathlineto{\pgfqpoint{5.296333in}{0.698040in}}%
\pgfpathlineto{\pgfqpoint{5.299689in}{0.703673in}}%
\pgfpathlineto{\pgfqpoint{5.303044in}{0.698801in}}%
\pgfpathlineto{\pgfqpoint{5.306400in}{0.678470in}}%
\pgfpathlineto{\pgfqpoint{5.309756in}{0.605831in}}%
\pgfpathlineto{\pgfqpoint{5.313112in}{0.689560in}}%
\pgfpathlineto{\pgfqpoint{5.316467in}{0.712111in}}%
\pgfpathlineto{\pgfqpoint{5.319823in}{0.720297in}}%
\pgfpathlineto{\pgfqpoint{5.323179in}{0.715502in}}%
\pgfpathlineto{\pgfqpoint{5.326535in}{0.693961in}}%
\pgfpathlineto{\pgfqpoint{5.329891in}{0.633732in}}%
\pgfpathlineto{\pgfqpoint{5.333246in}{0.686803in}}%
\pgfpathlineto{\pgfqpoint{5.336602in}{0.706399in}}%
\pgfpathlineto{\pgfqpoint{5.339958in}{0.711417in}}%
\pgfpathlineto{\pgfqpoint{5.343314in}{0.706059in}}%
\pgfpathlineto{\pgfqpoint{5.346669in}{0.684921in}}%
\pgfpathlineto{\pgfqpoint{5.350025in}{0.608315in}}%
\pgfpathlineto{\pgfqpoint{5.353381in}{0.697394in}}%
\pgfpathlineto{\pgfqpoint{5.356737in}{0.719494in}}%
\pgfpathlineto{\pgfqpoint{5.360093in}{0.727310in}}%
\pgfpathlineto{\pgfqpoint{5.363448in}{0.721896in}}%
\pgfpathlineto{\pgfqpoint{5.366804in}{0.697412in}}%
\pgfpathlineto{\pgfqpoint{5.370160in}{0.637001in}}%
\pgfpathlineto{\pgfqpoint{5.373516in}{0.703238in}}%
\pgfpathlineto{\pgfqpoint{5.376871in}{0.720840in}}%
\pgfpathlineto{\pgfqpoint{5.380227in}{0.725533in}}%
\pgfpathlineto{\pgfqpoint{5.383583in}{0.720709in}}%
\pgfpathlineto{\pgfqpoint{5.386939in}{0.702009in}}%
\pgfpathlineto{\pgfqpoint{5.390295in}{0.576009in}}%
\pgfpathlineto{\pgfqpoint{5.393650in}{0.705172in}}%
\pgfpathlineto{\pgfqpoint{5.397006in}{0.729995in}}%
\pgfpathlineto{\pgfqpoint{5.400362in}{0.739489in}}%
\pgfpathlineto{\pgfqpoint{5.403718in}{0.736688in}}%
\pgfpathlineto{\pgfqpoint{5.407074in}{0.718439in}}%
\pgfpathlineto{\pgfqpoint{5.410429in}{0.665202in}}%
\pgfpathlineto{\pgfqpoint{5.417141in}{0.716792in}}%
\pgfpathlineto{\pgfqpoint{5.420497in}{0.723194in}}%
\pgfpathlineto{\pgfqpoint{5.423852in}{0.717604in}}%
\pgfpathlineto{\pgfqpoint{5.427208in}{0.693364in}}%
\pgfpathlineto{\pgfqpoint{5.430564in}{0.653677in}}%
\pgfpathlineto{\pgfqpoint{5.433920in}{0.716285in}}%
\pgfpathlineto{\pgfqpoint{5.437276in}{0.735727in}}%
\pgfpathlineto{\pgfqpoint{5.440631in}{0.741337in}}%
\pgfpathlineto{\pgfqpoint{5.443987in}{0.732647in}}%
\pgfpathlineto{\pgfqpoint{5.447343in}{0.710521in}}%
\pgfpathlineto{\pgfqpoint{5.450699in}{0.651065in}}%
\pgfpathlineto{\pgfqpoint{5.454054in}{0.679042in}}%
\pgfpathlineto{\pgfqpoint{5.457410in}{0.700270in}}%
\pgfpathlineto{\pgfqpoint{5.460766in}{0.699786in}}%
\pgfpathlineto{\pgfqpoint{5.464122in}{0.677730in}}%
\pgfpathlineto{\pgfqpoint{5.467478in}{0.642764in}}%
\pgfpathlineto{\pgfqpoint{5.470833in}{0.704170in}}%
\pgfpathlineto{\pgfqpoint{5.474189in}{0.723626in}}%
\pgfpathlineto{\pgfqpoint{5.477545in}{0.728503in}}%
\pgfpathlineto{\pgfqpoint{5.480901in}{0.717994in}}%
\pgfpathlineto{\pgfqpoint{5.484256in}{0.693402in}}%
\pgfpathlineto{\pgfqpoint{5.487612in}{0.612636in}}%
\pgfpathlineto{\pgfqpoint{5.490968in}{0.672040in}}%
\pgfpathlineto{\pgfqpoint{5.494324in}{0.689005in}}%
\pgfpathlineto{\pgfqpoint{5.497680in}{0.686200in}}%
\pgfpathlineto{\pgfqpoint{5.501035in}{0.659074in}}%
\pgfpathlineto{\pgfqpoint{5.504391in}{0.646606in}}%
\pgfpathlineto{\pgfqpoint{5.507747in}{0.695598in}}%
\pgfpathlineto{\pgfqpoint{5.511103in}{0.712751in}}%
\pgfpathlineto{\pgfqpoint{5.514458in}{0.715741in}}%
\pgfpathlineto{\pgfqpoint{5.517814in}{0.702403in}}%
\pgfpathlineto{\pgfqpoint{5.521170in}{0.665615in}}%
\pgfpathlineto{\pgfqpoint{5.524526in}{0.649874in}}%
\pgfpathlineto{\pgfqpoint{5.527882in}{0.689949in}}%
\pgfpathlineto{\pgfqpoint{5.531237in}{0.700519in}}%
\pgfpathlineto{\pgfqpoint{5.534593in}{0.698673in}}%
\pgfpathlineto{\pgfqpoint{5.537949in}{0.682472in}}%
\pgfpathlineto{\pgfqpoint{5.541305in}{0.593120in}}%
\pgfpathlineto{\pgfqpoint{5.544660in}{0.682600in}}%
\pgfpathlineto{\pgfqpoint{5.548016in}{0.708846in}}%
\pgfpathlineto{\pgfqpoint{5.551372in}{0.718504in}}%
\pgfpathlineto{\pgfqpoint{5.554728in}{0.714632in}}%
\pgfpathlineto{\pgfqpoint{5.558084in}{0.691238in}}%
\pgfpathlineto{\pgfqpoint{5.561439in}{0.606784in}}%
\pgfpathlineto{\pgfqpoint{5.564795in}{0.695860in}}%
\pgfpathlineto{\pgfqpoint{5.568151in}{0.714561in}}%
\pgfpathlineto{\pgfqpoint{5.571507in}{0.720031in}}%
\pgfpathlineto{\pgfqpoint{5.574863in}{0.716236in}}%
\pgfpathlineto{\pgfqpoint{5.578218in}{0.699914in}}%
\pgfpathlineto{\pgfqpoint{5.581574in}{0.629672in}}%
\pgfpathlineto{\pgfqpoint{5.584930in}{0.693759in}}%
\pgfpathlineto{\pgfqpoint{5.588286in}{0.721831in}}%
\pgfpathlineto{\pgfqpoint{5.591641in}{0.732893in}}%
\pgfpathlineto{\pgfqpoint{5.594997in}{0.732065in}}%
\pgfpathlineto{\pgfqpoint{5.598353in}{0.715044in}}%
\pgfpathlineto{\pgfqpoint{5.601709in}{0.656393in}}%
\pgfpathlineto{\pgfqpoint{5.605065in}{0.693541in}}%
\pgfpathlineto{\pgfqpoint{5.608420in}{0.719986in}}%
\pgfpathlineto{\pgfqpoint{5.611776in}{0.727980in}}%
\pgfpathlineto{\pgfqpoint{5.615132in}{0.725591in}}%
\pgfpathlineto{\pgfqpoint{5.618488in}{0.710575in}}%
\pgfpathlineto{\pgfqpoint{5.621843in}{0.648319in}}%
\pgfpathlineto{\pgfqpoint{5.625199in}{0.701825in}}%
\pgfpathlineto{\pgfqpoint{5.628555in}{0.731107in}}%
\pgfpathlineto{\pgfqpoint{5.631911in}{0.742464in}}%
\pgfpathlineto{\pgfqpoint{5.635267in}{0.741431in}}%
\pgfpathlineto{\pgfqpoint{5.638622in}{0.725911in}}%
\pgfpathlineto{\pgfqpoint{5.641978in}{0.689789in}}%
\pgfpathlineto{\pgfqpoint{5.645334in}{0.661318in}}%
\pgfpathlineto{\pgfqpoint{5.648690in}{0.704343in}}%
\pgfpathlineto{\pgfqpoint{5.652045in}{0.712774in}}%
\pgfpathlineto{\pgfqpoint{5.655401in}{0.705800in}}%
\pgfpathlineto{\pgfqpoint{5.658757in}{0.671186in}}%
\pgfpathlineto{\pgfqpoint{5.662113in}{0.677817in}}%
\pgfpathlineto{\pgfqpoint{5.665469in}{0.718374in}}%
\pgfpathlineto{\pgfqpoint{5.668824in}{0.733296in}}%
\pgfpathlineto{\pgfqpoint{5.672180in}{0.734483in}}%
\pgfpathlineto{\pgfqpoint{5.675536in}{0.720869in}}%
\pgfpathlineto{\pgfqpoint{5.678892in}{0.693095in}}%
\pgfpathlineto{\pgfqpoint{5.682247in}{0.552465in}}%
\pgfpathlineto{\pgfqpoint{5.685603in}{0.677095in}}%
\pgfpathlineto{\pgfqpoint{5.688959in}{0.689314in}}%
\pgfpathlineto{\pgfqpoint{5.692315in}{0.680419in}}%
\pgfpathlineto{\pgfqpoint{5.695671in}{0.619483in}}%
\pgfpathlineto{\pgfqpoint{5.699026in}{0.678034in}}%
\pgfpathlineto{\pgfqpoint{5.702382in}{0.707875in}}%
\pgfpathlineto{\pgfqpoint{5.705738in}{0.718861in}}%
\pgfpathlineto{\pgfqpoint{5.709094in}{0.715006in}}%
\pgfpathlineto{\pgfqpoint{5.712449in}{0.695793in}}%
\pgfpathlineto{\pgfqpoint{5.715805in}{0.644031in}}%
\pgfpathlineto{\pgfqpoint{5.719161in}{0.660599in}}%
\pgfpathlineto{\pgfqpoint{5.722517in}{0.687301in}}%
\pgfpathlineto{\pgfqpoint{5.725873in}{0.691772in}}%
\pgfpathlineto{\pgfqpoint{5.729228in}{0.681459in}}%
\pgfpathlineto{\pgfqpoint{5.732584in}{0.634511in}}%
\pgfpathlineto{\pgfqpoint{5.739296in}{0.699987in}}%
\pgfpathlineto{\pgfqpoint{5.742652in}{0.712495in}}%
\pgfpathlineto{\pgfqpoint{5.746007in}{0.711264in}}%
\pgfpathlineto{\pgfqpoint{5.749363in}{0.691902in}}%
\pgfpathlineto{\pgfqpoint{5.752719in}{0.605462in}}%
\pgfpathlineto{\pgfqpoint{5.756075in}{0.681866in}}%
\pgfpathlineto{\pgfqpoint{5.759430in}{0.704517in}}%
\pgfpathlineto{\pgfqpoint{5.762786in}{0.711403in}}%
\pgfpathlineto{\pgfqpoint{5.766142in}{0.708478in}}%
\pgfpathlineto{\pgfqpoint{5.769498in}{0.693114in}}%
\pgfpathlineto{\pgfqpoint{5.772854in}{0.630025in}}%
\pgfpathlineto{\pgfqpoint{5.776209in}{0.684432in}}%
\pgfpathlineto{\pgfqpoint{5.779565in}{0.713584in}}%
\pgfpathlineto{\pgfqpoint{5.782921in}{0.724988in}}%
\pgfpathlineto{\pgfqpoint{5.786277in}{0.724256in}}%
\pgfpathlineto{\pgfqpoint{5.789632in}{0.706187in}}%
\pgfpathlineto{\pgfqpoint{5.792988in}{0.623491in}}%
\pgfpathlineto{\pgfqpoint{5.796344in}{0.696814in}}%
\pgfpathlineto{\pgfqpoint{5.799700in}{0.720419in}}%
\pgfpathlineto{\pgfqpoint{5.803056in}{0.728310in}}%
\pgfpathlineto{\pgfqpoint{5.806411in}{0.726906in}}%
\pgfpathlineto{\pgfqpoint{5.809767in}{0.714895in}}%
\pgfpathlineto{\pgfqpoint{5.813123in}{0.675452in}}%
\pgfpathlineto{\pgfqpoint{5.816479in}{0.686589in}}%
\pgfpathlineto{\pgfqpoint{5.819834in}{0.724630in}}%
\pgfpathlineto{\pgfqpoint{5.823190in}{0.739332in}}%
\pgfpathlineto{\pgfqpoint{5.826546in}{0.742222in}}%
\pgfpathlineto{\pgfqpoint{5.829902in}{0.731026in}}%
\pgfpathlineto{\pgfqpoint{5.833258in}{0.700567in}}%
\pgfpathlineto{\pgfqpoint{5.836613in}{0.652622in}}%
\pgfpathlineto{\pgfqpoint{5.839969in}{0.709903in}}%
\pgfpathlineto{\pgfqpoint{5.843325in}{0.723002in}}%
\pgfpathlineto{\pgfqpoint{5.846681in}{0.722476in}}%
\pgfpathlineto{\pgfqpoint{5.850036in}{0.707802in}}%
\pgfpathlineto{\pgfqpoint{5.853392in}{0.636657in}}%
\pgfpathlineto{\pgfqpoint{5.856748in}{0.703878in}}%
\pgfpathlineto{\pgfqpoint{5.860104in}{0.731782in}}%
\pgfpathlineto{\pgfqpoint{5.863460in}{0.742162in}}%
\pgfpathlineto{\pgfqpoint{5.866815in}{0.739006in}}%
\pgfpathlineto{\pgfqpoint{5.870171in}{0.722445in}}%
\pgfpathlineto{\pgfqpoint{5.873527in}{0.688082in}}%
\pgfpathlineto{\pgfqpoint{5.876883in}{0.637863in}}%
\pgfpathlineto{\pgfqpoint{5.880238in}{0.689773in}}%
\pgfpathlineto{\pgfqpoint{5.883594in}{0.695974in}}%
\pgfpathlineto{\pgfqpoint{5.886950in}{0.680827in}}%
\pgfpathlineto{\pgfqpoint{5.890306in}{0.589486in}}%
\pgfpathlineto{\pgfqpoint{5.893662in}{0.695609in}}%
\pgfpathlineto{\pgfqpoint{5.897017in}{0.718974in}}%
\pgfpathlineto{\pgfqpoint{5.900373in}{0.726443in}}%
\pgfpathlineto{\pgfqpoint{5.903729in}{0.718539in}}%
\pgfpathlineto{\pgfqpoint{5.907085in}{0.697282in}}%
\pgfpathlineto{\pgfqpoint{5.910441in}{0.643589in}}%
\pgfpathlineto{\pgfqpoint{5.913796in}{0.658365in}}%
\pgfpathlineto{\pgfqpoint{5.917152in}{0.682325in}}%
\pgfpathlineto{\pgfqpoint{5.920508in}{0.682121in}}%
\pgfpathlineto{\pgfqpoint{5.927219in}{0.632813in}}%
\pgfpathlineto{\pgfqpoint{5.930575in}{0.689282in}}%
\pgfpathlineto{\pgfqpoint{5.933931in}{0.707963in}}%
\pgfpathlineto{\pgfqpoint{5.937287in}{0.712140in}}%
\pgfpathlineto{\pgfqpoint{5.940643in}{0.699926in}}%
\pgfpathlineto{\pgfqpoint{5.943998in}{0.664214in}}%
\pgfpathlineto{\pgfqpoint{5.947354in}{0.646643in}}%
\pgfpathlineto{\pgfqpoint{5.950710in}{0.688336in}}%
\pgfpathlineto{\pgfqpoint{5.954066in}{0.699774in}}%
\pgfpathlineto{\pgfqpoint{5.957421in}{0.699140in}}%
\pgfpathlineto{\pgfqpoint{5.960777in}{0.685729in}}%
\pgfpathlineto{\pgfqpoint{5.964133in}{0.631610in}}%
\pgfpathlineto{\pgfqpoint{5.967489in}{0.673335in}}%
\pgfpathlineto{\pgfqpoint{5.970845in}{0.704421in}}%
\pgfpathlineto{\pgfqpoint{5.974200in}{0.716340in}}%
\pgfpathlineto{\pgfqpoint{5.977556in}{0.715462in}}%
\pgfpathlineto{\pgfqpoint{5.980912in}{0.696068in}}%
\pgfpathlineto{\pgfqpoint{5.984268in}{0.561952in}}%
\pgfpathlineto{\pgfqpoint{5.987623in}{0.693883in}}%
\pgfpathlineto{\pgfqpoint{5.990979in}{0.715761in}}%
\pgfpathlineto{\pgfqpoint{5.994335in}{0.723291in}}%
\pgfpathlineto{\pgfqpoint{5.997691in}{0.721973in}}%
\pgfpathlineto{\pgfqpoint{6.001047in}{0.710586in}}%
\pgfpathlineto{\pgfqpoint{6.004402in}{0.674802in}}%
\pgfpathlineto{\pgfqpoint{6.007758in}{0.675766in}}%
\pgfpathlineto{\pgfqpoint{6.011114in}{0.717196in}}%
\pgfpathlineto{\pgfqpoint{6.014470in}{0.732927in}}%
\pgfpathlineto{\pgfqpoint{6.017825in}{0.736901in}}%
\pgfpathlineto{\pgfqpoint{6.021181in}{0.727101in}}%
\pgfpathlineto{\pgfqpoint{6.024537in}{0.694721in}}%
\pgfpathlineto{\pgfqpoint{6.027893in}{0.669429in}}%
\pgfpathlineto{\pgfqpoint{6.031249in}{0.716313in}}%
\pgfpathlineto{\pgfqpoint{6.034604in}{0.729808in}}%
\pgfpathlineto{\pgfqpoint{6.037960in}{0.731617in}}%
\pgfpathlineto{\pgfqpoint{6.041316in}{0.722953in}}%
\pgfpathlineto{\pgfqpoint{6.044672in}{0.693683in}}%
\pgfpathlineto{\pgfqpoint{6.048027in}{0.674660in}}%
\pgfpathlineto{\pgfqpoint{6.051383in}{0.724970in}}%
\pgfpathlineto{\pgfqpoint{6.054739in}{0.742475in}}%
\pgfpathlineto{\pgfqpoint{6.058095in}{0.747199in}}%
\pgfpathlineto{\pgfqpoint{6.061451in}{0.738121in}}%
\pgfpathlineto{\pgfqpoint{6.064806in}{0.715253in}}%
\pgfpathlineto{\pgfqpoint{6.068162in}{0.646834in}}%
\pgfpathlineto{\pgfqpoint{6.071518in}{0.688531in}}%
\pgfpathlineto{\pgfqpoint{6.074874in}{0.708462in}}%
\pgfpathlineto{\pgfqpoint{6.078230in}{0.707829in}}%
\pgfpathlineto{\pgfqpoint{6.081585in}{0.686469in}}%
\pgfpathlineto{\pgfqpoint{6.084941in}{0.644775in}}%
\pgfpathlineto{\pgfqpoint{6.088297in}{0.710381in}}%
\pgfpathlineto{\pgfqpoint{6.091653in}{0.730344in}}%
\pgfpathlineto{\pgfqpoint{6.095008in}{0.735616in}}%
\pgfpathlineto{\pgfqpoint{6.098364in}{0.725773in}}%
\pgfpathlineto{\pgfqpoint{6.101720in}{0.704097in}}%
\pgfpathlineto{\pgfqpoint{6.105076in}{0.654209in}}%
\pgfpathlineto{\pgfqpoint{6.108432in}{0.655834in}}%
\pgfpathlineto{\pgfqpoint{6.111787in}{0.680369in}}%
\pgfpathlineto{\pgfqpoint{6.115143in}{0.676075in}}%
\pgfpathlineto{\pgfqpoint{6.118499in}{0.629276in}}%
\pgfpathlineto{\pgfqpoint{6.125210in}{0.702104in}}%
\pgfpathlineto{\pgfqpoint{6.128566in}{0.714715in}}%
\pgfpathlineto{\pgfqpoint{6.131922in}{0.712593in}}%
\pgfpathlineto{\pgfqpoint{6.135278in}{0.694739in}}%
\pgfpathlineto{\pgfqpoint{6.138634in}{0.648798in}}%
\pgfpathlineto{\pgfqpoint{6.141989in}{0.653004in}}%
\pgfpathlineto{\pgfqpoint{6.145345in}{0.683357in}}%
\pgfpathlineto{\pgfqpoint{6.148701in}{0.689267in}}%
\pgfpathlineto{\pgfqpoint{6.152057in}{0.680804in}}%
\pgfpathlineto{\pgfqpoint{6.155412in}{0.642974in}}%
\pgfpathlineto{\pgfqpoint{6.158768in}{0.656322in}}%
\pgfpathlineto{\pgfqpoint{6.162124in}{0.694474in}}%
\pgfpathlineto{\pgfqpoint{6.165480in}{0.708674in}}%
\pgfpathlineto{\pgfqpoint{6.168836in}{0.709315in}}%
\pgfpathlineto{\pgfqpoint{6.172191in}{0.691910in}}%
\pgfpathlineto{\pgfqpoint{6.175547in}{0.616012in}}%
\pgfpathlineto{\pgfqpoint{6.178903in}{0.680762in}}%
\pgfpathlineto{\pgfqpoint{6.182259in}{0.705096in}}%
\pgfpathlineto{\pgfqpoint{6.185614in}{0.713315in}}%
\pgfpathlineto{\pgfqpoint{6.188970in}{0.712219in}}%
\pgfpathlineto{\pgfqpoint{6.192326in}{0.700715in}}%
\pgfpathlineto{\pgfqpoint{6.195682in}{0.663535in}}%
\pgfpathlineto{\pgfqpoint{6.199038in}{0.668875in}}%
\pgfpathlineto{\pgfqpoint{6.202393in}{0.708795in}}%
\pgfpathlineto{\pgfqpoint{6.205749in}{0.724039in}}%
\pgfpathlineto{\pgfqpoint{6.209105in}{0.727410in}}%
\pgfpathlineto{\pgfqpoint{6.212461in}{0.716057in}}%
\pgfpathlineto{\pgfqpoint{6.215816in}{0.672801in}}%
\pgfpathlineto{\pgfqpoint{6.219172in}{0.686071in}}%
\pgfpathlineto{\pgfqpoint{6.222528in}{0.719702in}}%
\pgfpathlineto{\pgfqpoint{6.225884in}{0.731473in}}%
\pgfpathlineto{\pgfqpoint{6.229240in}{0.733433in}}%
\pgfpathlineto{\pgfqpoint{6.232595in}{0.726484in}}%
\pgfpathlineto{\pgfqpoint{6.235951in}{0.704390in}}%
\pgfpathlineto{\pgfqpoint{6.239307in}{0.625564in}}%
\pgfpathlineto{\pgfqpoint{6.242663in}{0.716145in}}%
\pgfpathlineto{\pgfqpoint{6.246019in}{0.738279in}}%
\pgfpathlineto{\pgfqpoint{6.249374in}{0.746590in}}%
\pgfpathlineto{\pgfqpoint{6.252730in}{0.742847in}}%
\pgfpathlineto{\pgfqpoint{6.256086in}{0.724193in}}%
\pgfpathlineto{\pgfqpoint{6.259442in}{0.670751in}}%
\pgfpathlineto{\pgfqpoint{6.262797in}{0.693061in}}%
\pgfpathlineto{\pgfqpoint{6.266153in}{0.719694in}}%
\pgfpathlineto{\pgfqpoint{6.269509in}{0.725377in}}%
\pgfpathlineto{\pgfqpoint{6.272865in}{0.718033in}}%
\pgfpathlineto{\pgfqpoint{6.276221in}{0.687056in}}%
\pgfpathlineto{\pgfqpoint{6.279576in}{0.679810in}}%
\pgfpathlineto{\pgfqpoint{6.282932in}{0.725284in}}%
\pgfpathlineto{\pgfqpoint{6.286288in}{0.741543in}}%
\pgfpathlineto{\pgfqpoint{6.289644in}{0.744395in}}%
\pgfpathlineto{\pgfqpoint{6.292999in}{0.732788in}}%
\pgfpathlineto{\pgfqpoint{6.296355in}{0.709513in}}%
\pgfpathlineto{\pgfqpoint{6.299711in}{0.652578in}}%
\pgfpathlineto{\pgfqpoint{6.303067in}{0.668042in}}%
\pgfpathlineto{\pgfqpoint{6.306423in}{0.687902in}}%
\pgfpathlineto{\pgfqpoint{6.309778in}{0.679967in}}%
\pgfpathlineto{\pgfqpoint{6.313134in}{0.600488in}}%
\pgfpathlineto{\pgfqpoint{6.316490in}{0.685509in}}%
\pgfpathlineto{\pgfqpoint{6.319846in}{0.713499in}}%
\pgfpathlineto{\pgfqpoint{6.323201in}{0.723657in}}%
\pgfpathlineto{\pgfqpoint{6.326557in}{0.718607in}}%
\pgfpathlineto{\pgfqpoint{6.329913in}{0.700141in}}%
\pgfpathlineto{\pgfqpoint{6.333269in}{0.659773in}}%
\pgfpathlineto{\pgfqpoint{6.336625in}{0.638498in}}%
\pgfpathlineto{\pgfqpoint{6.339980in}{0.674387in}}%
\pgfpathlineto{\pgfqpoint{6.343336in}{0.677345in}}%
\pgfpathlineto{\pgfqpoint{6.346692in}{0.657849in}}%
\pgfpathlineto{\pgfqpoint{6.350048in}{0.614375in}}%
\pgfpathlineto{\pgfqpoint{6.353403in}{0.682223in}}%
\pgfpathlineto{\pgfqpoint{6.356759in}{0.702690in}}%
\pgfpathlineto{\pgfqpoint{6.360115in}{0.708189in}}%
\pgfpathlineto{\pgfqpoint{6.363471in}{0.699442in}}%
\pgfpathlineto{\pgfqpoint{6.373538in}{0.664528in}}%
\pgfpathlineto{\pgfqpoint{6.376894in}{0.655621in}}%
\pgfpathlineto{\pgfqpoint{6.383605in}{0.640787in}}%
\pgfpathlineto{\pgfqpoint{6.383605in}{0.640787in}}%
\pgfusepath{stroke}%
\end{pgfscope}%
\begin{pgfscope}%
\pgfpathrectangle{\pgfqpoint{4.621821in}{0.320625in}}{\pgfqpoint{1.845679in}{1.164375in}}%
\pgfusepath{clip}%
\pgfsetrectcap%
\pgfsetroundjoin%
\pgfsetlinewidth{0.602250pt}%
\definecolor{currentstroke}{rgb}{0.000000,0.529412,0.423529}%
\pgfsetstrokecolor{currentstroke}%
\pgfsetstrokeopacity{0.120000}%
\pgfsetdash{}{0pt}%
\pgfpathmoveto{\pgfqpoint{4.705715in}{1.469311in}}%
\pgfpathlineto{\pgfqpoint{4.719139in}{1.444863in}}%
\pgfpathlineto{\pgfqpoint{4.729206in}{1.423893in}}%
\pgfpathlineto{\pgfqpoint{4.745985in}{1.385546in}}%
\pgfpathlineto{\pgfqpoint{4.772831in}{1.324120in}}%
\pgfpathlineto{\pgfqpoint{4.792966in}{1.280817in}}%
\pgfpathlineto{\pgfqpoint{4.816456in}{1.233087in}}%
\pgfpathlineto{\pgfqpoint{4.843302in}{1.181222in}}%
\pgfpathlineto{\pgfqpoint{4.873504in}{1.125290in}}%
\pgfpathlineto{\pgfqpoint{4.910418in}{1.059402in}}%
\pgfpathlineto{\pgfqpoint{4.954043in}{0.984005in}}%
\pgfpathlineto{\pgfqpoint{5.004380in}{0.899187in}}%
\pgfpathlineto{\pgfqpoint{5.068140in}{0.794029in}}%
\pgfpathlineto{\pgfqpoint{5.081563in}{0.773601in}}%
\pgfpathlineto{\pgfqpoint{5.094986in}{0.749142in}}%
\pgfpathlineto{\pgfqpoint{5.101698in}{0.737350in}}%
\pgfpathlineto{\pgfqpoint{5.108409in}{0.727903in}}%
\pgfpathlineto{\pgfqpoint{5.115121in}{0.721381in}}%
\pgfpathlineto{\pgfqpoint{5.121832in}{0.715264in}}%
\pgfpathlineto{\pgfqpoint{5.125188in}{0.710278in}}%
\pgfpathlineto{\pgfqpoint{5.131900in}{0.696204in}}%
\pgfpathlineto{\pgfqpoint{5.138611in}{0.677484in}}%
\pgfpathlineto{\pgfqpoint{5.148678in}{0.647482in}}%
\pgfpathlineto{\pgfqpoint{5.152034in}{0.640357in}}%
\pgfpathlineto{\pgfqpoint{5.155390in}{0.634822in}}%
\pgfpathlineto{\pgfqpoint{5.158746in}{0.631770in}}%
\pgfpathlineto{\pgfqpoint{5.162102in}{0.634449in}}%
\pgfpathlineto{\pgfqpoint{5.175525in}{0.666989in}}%
\pgfpathlineto{\pgfqpoint{5.178880in}{0.669671in}}%
\pgfpathlineto{\pgfqpoint{5.182236in}{0.668319in}}%
\pgfpathlineto{\pgfqpoint{5.185592in}{0.660923in}}%
\pgfpathlineto{\pgfqpoint{5.188948in}{0.657969in}}%
\pgfpathlineto{\pgfqpoint{5.195659in}{0.696887in}}%
\pgfpathlineto{\pgfqpoint{5.199015in}{0.710435in}}%
\pgfpathlineto{\pgfqpoint{5.202371in}{0.719893in}}%
\pgfpathlineto{\pgfqpoint{5.205727in}{0.726374in}}%
\pgfpathlineto{\pgfqpoint{5.209082in}{0.730495in}}%
\pgfpathlineto{\pgfqpoint{5.212438in}{0.732567in}}%
\pgfpathlineto{\pgfqpoint{5.215794in}{0.732689in}}%
\pgfpathlineto{\pgfqpoint{5.219150in}{0.730756in}}%
\pgfpathlineto{\pgfqpoint{5.222506in}{0.726401in}}%
\pgfpathlineto{\pgfqpoint{5.225861in}{0.718766in}}%
\pgfpathlineto{\pgfqpoint{5.229217in}{0.705752in}}%
\pgfpathlineto{\pgfqpoint{5.232573in}{0.680406in}}%
\pgfpathlineto{\pgfqpoint{5.235929in}{0.631094in}}%
\pgfpathlineto{\pgfqpoint{5.239285in}{0.690724in}}%
\pgfpathlineto{\pgfqpoint{5.242640in}{0.714427in}}%
\pgfpathlineto{\pgfqpoint{5.245996in}{0.728197in}}%
\pgfpathlineto{\pgfqpoint{5.249352in}{0.737030in}}%
\pgfpathlineto{\pgfqpoint{5.252708in}{0.742466in}}%
\pgfpathlineto{\pgfqpoint{5.256063in}{0.744946in}}%
\pgfpathlineto{\pgfqpoint{5.259419in}{0.744233in}}%
\pgfpathlineto{\pgfqpoint{5.262775in}{0.739170in}}%
\pgfpathlineto{\pgfqpoint{5.266131in}{0.728230in}}%
\pgfpathlineto{\pgfqpoint{5.272842in}{0.694978in}}%
\pgfpathlineto{\pgfqpoint{5.276198in}{0.685876in}}%
\pgfpathlineto{\pgfqpoint{5.279554in}{0.683782in}}%
\pgfpathlineto{\pgfqpoint{5.282910in}{0.675829in}}%
\pgfpathlineto{\pgfqpoint{5.286265in}{0.664190in}}%
\pgfpathlineto{\pgfqpoint{5.292977in}{0.700521in}}%
\pgfpathlineto{\pgfqpoint{5.296333in}{0.708680in}}%
\pgfpathlineto{\pgfqpoint{5.299689in}{0.703397in}}%
\pgfpathlineto{\pgfqpoint{5.303044in}{0.687574in}}%
\pgfpathlineto{\pgfqpoint{5.309756in}{0.634037in}}%
\pgfpathlineto{\pgfqpoint{5.313112in}{0.645660in}}%
\pgfpathlineto{\pgfqpoint{5.316467in}{0.647536in}}%
\pgfpathlineto{\pgfqpoint{5.319823in}{0.627053in}}%
\pgfpathlineto{\pgfqpoint{5.323179in}{0.624065in}}%
\pgfpathlineto{\pgfqpoint{5.326535in}{0.663316in}}%
\pgfpathlineto{\pgfqpoint{5.329891in}{0.680359in}}%
\pgfpathlineto{\pgfqpoint{5.333246in}{0.683927in}}%
\pgfpathlineto{\pgfqpoint{5.336602in}{0.667586in}}%
\pgfpathlineto{\pgfqpoint{5.339958in}{0.613987in}}%
\pgfpathlineto{\pgfqpoint{5.343314in}{0.683651in}}%
\pgfpathlineto{\pgfqpoint{5.346669in}{0.705825in}}%
\pgfpathlineto{\pgfqpoint{5.350025in}{0.716188in}}%
\pgfpathlineto{\pgfqpoint{5.353381in}{0.719733in}}%
\pgfpathlineto{\pgfqpoint{5.356737in}{0.717401in}}%
\pgfpathlineto{\pgfqpoint{5.360093in}{0.707813in}}%
\pgfpathlineto{\pgfqpoint{5.363448in}{0.683691in}}%
\pgfpathlineto{\pgfqpoint{5.366804in}{0.613368in}}%
\pgfpathlineto{\pgfqpoint{5.370160in}{0.696302in}}%
\pgfpathlineto{\pgfqpoint{5.373516in}{0.718602in}}%
\pgfpathlineto{\pgfqpoint{5.376871in}{0.728576in}}%
\pgfpathlineto{\pgfqpoint{5.380227in}{0.729717in}}%
\pgfpathlineto{\pgfqpoint{5.383583in}{0.719150in}}%
\pgfpathlineto{\pgfqpoint{5.386939in}{0.691557in}}%
\pgfpathlineto{\pgfqpoint{5.390295in}{0.633422in}}%
\pgfpathlineto{\pgfqpoint{5.393650in}{0.693992in}}%
\pgfpathlineto{\pgfqpoint{5.397006in}{0.708675in}}%
\pgfpathlineto{\pgfqpoint{5.400362in}{0.709407in}}%
\pgfpathlineto{\pgfqpoint{5.403718in}{0.696958in}}%
\pgfpathlineto{\pgfqpoint{5.407074in}{0.645853in}}%
\pgfpathlineto{\pgfqpoint{5.413785in}{0.715464in}}%
\pgfpathlineto{\pgfqpoint{5.417141in}{0.727775in}}%
\pgfpathlineto{\pgfqpoint{5.420497in}{0.727261in}}%
\pgfpathlineto{\pgfqpoint{5.423852in}{0.711263in}}%
\pgfpathlineto{\pgfqpoint{5.427208in}{0.668943in}}%
\pgfpathlineto{\pgfqpoint{5.430564in}{0.670315in}}%
\pgfpathlineto{\pgfqpoint{5.433920in}{0.702670in}}%
\pgfpathlineto{\pgfqpoint{5.437276in}{0.711121in}}%
\pgfpathlineto{\pgfqpoint{5.440631in}{0.707371in}}%
\pgfpathlineto{\pgfqpoint{5.443987in}{0.687263in}}%
\pgfpathlineto{\pgfqpoint{5.447343in}{0.606293in}}%
\pgfpathlineto{\pgfqpoint{5.450699in}{0.699871in}}%
\pgfpathlineto{\pgfqpoint{5.454054in}{0.722171in}}%
\pgfpathlineto{\pgfqpoint{5.457410in}{0.729764in}}%
\pgfpathlineto{\pgfqpoint{5.460766in}{0.723368in}}%
\pgfpathlineto{\pgfqpoint{5.464122in}{0.697415in}}%
\pgfpathlineto{\pgfqpoint{5.467478in}{0.637738in}}%
\pgfpathlineto{\pgfqpoint{5.470833in}{0.705158in}}%
\pgfpathlineto{\pgfqpoint{5.474189in}{0.722006in}}%
\pgfpathlineto{\pgfqpoint{5.477545in}{0.726072in}}%
\pgfpathlineto{\pgfqpoint{5.480901in}{0.720371in}}%
\pgfpathlineto{\pgfqpoint{5.484256in}{0.699478in}}%
\pgfpathlineto{\pgfqpoint{5.487612in}{0.608954in}}%
\pgfpathlineto{\pgfqpoint{5.490968in}{0.709892in}}%
\pgfpathlineto{\pgfqpoint{5.494324in}{0.732518in}}%
\pgfpathlineto{\pgfqpoint{5.497680in}{0.740808in}}%
\pgfpathlineto{\pgfqpoint{5.501035in}{0.736399in}}%
\pgfpathlineto{\pgfqpoint{5.504391in}{0.714783in}}%
\pgfpathlineto{\pgfqpoint{5.507747in}{0.620561in}}%
\pgfpathlineto{\pgfqpoint{5.511103in}{0.707609in}}%
\pgfpathlineto{\pgfqpoint{5.514458in}{0.728084in}}%
\pgfpathlineto{\pgfqpoint{5.517814in}{0.733487in}}%
\pgfpathlineto{\pgfqpoint{5.521170in}{0.728597in}}%
\pgfpathlineto{\pgfqpoint{5.524526in}{0.708633in}}%
\pgfpathlineto{\pgfqpoint{5.527882in}{0.600887in}}%
\pgfpathlineto{\pgfqpoint{5.531237in}{0.717425in}}%
\pgfpathlineto{\pgfqpoint{5.534593in}{0.740582in}}%
\pgfpathlineto{\pgfqpoint{5.537949in}{0.749001in}}%
\pgfpathlineto{\pgfqpoint{5.541305in}{0.744445in}}%
\pgfpathlineto{\pgfqpoint{5.544660in}{0.725347in}}%
\pgfpathlineto{\pgfqpoint{5.548016in}{0.674078in}}%
\pgfpathlineto{\pgfqpoint{5.551372in}{0.691112in}}%
\pgfpathlineto{\pgfqpoint{5.554728in}{0.717747in}}%
\pgfpathlineto{\pgfqpoint{5.558084in}{0.722664in}}%
\pgfpathlineto{\pgfqpoint{5.561439in}{0.713482in}}%
\pgfpathlineto{\pgfqpoint{5.564795in}{0.673584in}}%
\pgfpathlineto{\pgfqpoint{5.568151in}{0.691324in}}%
\pgfpathlineto{\pgfqpoint{5.571507in}{0.728157in}}%
\pgfpathlineto{\pgfqpoint{5.574863in}{0.742028in}}%
\pgfpathlineto{\pgfqpoint{5.578218in}{0.742555in}}%
\pgfpathlineto{\pgfqpoint{5.581574in}{0.728405in}}%
\pgfpathlineto{\pgfqpoint{5.584930in}{0.698375in}}%
\pgfpathlineto{\pgfqpoint{5.588286in}{0.623180in}}%
\pgfpathlineto{\pgfqpoint{5.591641in}{0.693710in}}%
\pgfpathlineto{\pgfqpoint{5.594997in}{0.704553in}}%
\pgfpathlineto{\pgfqpoint{5.598353in}{0.697322in}}%
\pgfpathlineto{\pgfqpoint{5.601709in}{0.655805in}}%
\pgfpathlineto{\pgfqpoint{5.605065in}{0.680678in}}%
\pgfpathlineto{\pgfqpoint{5.608420in}{0.716263in}}%
\pgfpathlineto{\pgfqpoint{5.611776in}{0.729674in}}%
\pgfpathlineto{\pgfqpoint{5.615132in}{0.729112in}}%
\pgfpathlineto{\pgfqpoint{5.618488in}{0.713007in}}%
\pgfpathlineto{\pgfqpoint{5.621843in}{0.674128in}}%
\pgfpathlineto{\pgfqpoint{5.625199in}{0.658825in}}%
\pgfpathlineto{\pgfqpoint{5.628555in}{0.696170in}}%
\pgfpathlineto{\pgfqpoint{5.631911in}{0.704132in}}%
\pgfpathlineto{\pgfqpoint{5.635267in}{0.697764in}}%
\pgfpathlineto{\pgfqpoint{5.638622in}{0.667304in}}%
\pgfpathlineto{\pgfqpoint{5.641978in}{0.660478in}}%
\pgfpathlineto{\pgfqpoint{5.645334in}{0.706096in}}%
\pgfpathlineto{\pgfqpoint{5.648690in}{0.722439in}}%
\pgfpathlineto{\pgfqpoint{5.652045in}{0.725231in}}%
\pgfpathlineto{\pgfqpoint{5.655401in}{0.711809in}}%
\pgfpathlineto{\pgfqpoint{5.658757in}{0.670589in}}%
\pgfpathlineto{\pgfqpoint{5.662113in}{0.672667in}}%
\pgfpathlineto{\pgfqpoint{5.665469in}{0.707414in}}%
\pgfpathlineto{\pgfqpoint{5.668824in}{0.717945in}}%
\pgfpathlineto{\pgfqpoint{5.672180in}{0.717475in}}%
\pgfpathlineto{\pgfqpoint{5.675536in}{0.705411in}}%
\pgfpathlineto{\pgfqpoint{5.678892in}{0.662301in}}%
\pgfpathlineto{\pgfqpoint{5.682247in}{0.683069in}}%
\pgfpathlineto{\pgfqpoint{5.685603in}{0.718532in}}%
\pgfpathlineto{\pgfqpoint{5.688959in}{0.732239in}}%
\pgfpathlineto{\pgfqpoint{5.692315in}{0.733693in}}%
\pgfpathlineto{\pgfqpoint{5.695671in}{0.719217in}}%
\pgfpathlineto{\pgfqpoint{5.699026in}{0.669924in}}%
\pgfpathlineto{\pgfqpoint{5.702382in}{0.692363in}}%
\pgfpathlineto{\pgfqpoint{5.705738in}{0.722384in}}%
\pgfpathlineto{\pgfqpoint{5.709094in}{0.732231in}}%
\pgfpathlineto{\pgfqpoint{5.712449in}{0.731998in}}%
\pgfpathlineto{\pgfqpoint{5.715805in}{0.721233in}}%
\pgfpathlineto{\pgfqpoint{5.719161in}{0.685920in}}%
\pgfpathlineto{\pgfqpoint{5.722517in}{0.686716in}}%
\pgfpathlineto{\pgfqpoint{5.725873in}{0.728329in}}%
\pgfpathlineto{\pgfqpoint{5.729228in}{0.743977in}}%
\pgfpathlineto{\pgfqpoint{5.732584in}{0.747521in}}%
\pgfpathlineto{\pgfqpoint{5.735940in}{0.736982in}}%
\pgfpathlineto{\pgfqpoint{5.739296in}{0.707423in}}%
\pgfpathlineto{\pgfqpoint{5.742652in}{0.655621in}}%
\pgfpathlineto{\pgfqpoint{5.746007in}{0.715686in}}%
\pgfpathlineto{\pgfqpoint{5.749363in}{0.729526in}}%
\pgfpathlineto{\pgfqpoint{5.752719in}{0.729878in}}%
\pgfpathlineto{\pgfqpoint{5.756075in}{0.717103in}}%
\pgfpathlineto{\pgfqpoint{5.759430in}{0.664550in}}%
\pgfpathlineto{\pgfqpoint{5.762786in}{0.704289in}}%
\pgfpathlineto{\pgfqpoint{5.766142in}{0.735820in}}%
\pgfpathlineto{\pgfqpoint{5.769498in}{0.747874in}}%
\pgfpathlineto{\pgfqpoint{5.772854in}{0.746960in}}%
\pgfpathlineto{\pgfqpoint{5.776209in}{0.731750in}}%
\pgfpathlineto{\pgfqpoint{5.779565in}{0.698991in}}%
\pgfpathlineto{\pgfqpoint{5.782921in}{0.646392in}}%
\pgfpathlineto{\pgfqpoint{5.786277in}{0.701798in}}%
\pgfpathlineto{\pgfqpoint{5.789632in}{0.710805in}}%
\pgfpathlineto{\pgfqpoint{5.792988in}{0.702211in}}%
\pgfpathlineto{\pgfqpoint{5.796344in}{0.655116in}}%
\pgfpathlineto{\pgfqpoint{5.803056in}{0.723268in}}%
\pgfpathlineto{\pgfqpoint{5.806411in}{0.735760in}}%
\pgfpathlineto{\pgfqpoint{5.809767in}{0.734127in}}%
\pgfpathlineto{\pgfqpoint{5.813123in}{0.717873in}}%
\pgfpathlineto{\pgfqpoint{5.819834in}{0.646993in}}%
\pgfpathlineto{\pgfqpoint{5.823190in}{0.691915in}}%
\pgfpathlineto{\pgfqpoint{5.826546in}{0.699256in}}%
\pgfpathlineto{\pgfqpoint{5.829902in}{0.689155in}}%
\pgfpathlineto{\pgfqpoint{5.833258in}{0.633626in}}%
\pgfpathlineto{\pgfqpoint{5.836613in}{0.682154in}}%
\pgfpathlineto{\pgfqpoint{5.839969in}{0.712947in}}%
\pgfpathlineto{\pgfqpoint{5.843325in}{0.724421in}}%
\pgfpathlineto{\pgfqpoint{5.846681in}{0.721589in}}%
\pgfpathlineto{\pgfqpoint{5.850036in}{0.702030in}}%
\pgfpathlineto{\pgfqpoint{5.853392in}{0.637250in}}%
\pgfpathlineto{\pgfqpoint{5.856748in}{0.680811in}}%
\pgfpathlineto{\pgfqpoint{5.860104in}{0.704713in}}%
\pgfpathlineto{\pgfqpoint{5.863460in}{0.710423in}}%
\pgfpathlineto{\pgfqpoint{5.866815in}{0.704452in}}%
\pgfpathlineto{\pgfqpoint{5.870171in}{0.679783in}}%
\pgfpathlineto{\pgfqpoint{5.873527in}{0.641644in}}%
\pgfpathlineto{\pgfqpoint{5.876883in}{0.703222in}}%
\pgfpathlineto{\pgfqpoint{5.880238in}{0.722529in}}%
\pgfpathlineto{\pgfqpoint{5.883594in}{0.728138in}}%
\pgfpathlineto{\pgfqpoint{5.886950in}{0.718937in}}%
\pgfpathlineto{\pgfqpoint{5.890306in}{0.685250in}}%
\pgfpathlineto{\pgfqpoint{5.893662in}{0.669094in}}%
\pgfpathlineto{\pgfqpoint{5.897017in}{0.712553in}}%
\pgfpathlineto{\pgfqpoint{5.900373in}{0.726125in}}%
\pgfpathlineto{\pgfqpoint{5.903729in}{0.728710in}}%
\pgfpathlineto{\pgfqpoint{5.907085in}{0.721861in}}%
\pgfpathlineto{\pgfqpoint{5.910441in}{0.699039in}}%
\pgfpathlineto{\pgfqpoint{5.913796in}{0.634783in}}%
\pgfpathlineto{\pgfqpoint{5.917152in}{0.713926in}}%
\pgfpathlineto{\pgfqpoint{5.920508in}{0.735241in}}%
\pgfpathlineto{\pgfqpoint{5.923864in}{0.742874in}}%
\pgfpathlineto{\pgfqpoint{5.927219in}{0.737827in}}%
\pgfpathlineto{\pgfqpoint{5.930575in}{0.714939in}}%
\pgfpathlineto{\pgfqpoint{5.933931in}{0.579531in}}%
\pgfpathlineto{\pgfqpoint{5.937287in}{0.712625in}}%
\pgfpathlineto{\pgfqpoint{5.940643in}{0.731848in}}%
\pgfpathlineto{\pgfqpoint{5.943998in}{0.736745in}}%
\pgfpathlineto{\pgfqpoint{5.947354in}{0.731452in}}%
\pgfpathlineto{\pgfqpoint{5.950710in}{0.710772in}}%
\pgfpathlineto{\pgfqpoint{5.954066in}{0.620641in}}%
\pgfpathlineto{\pgfqpoint{5.957421in}{0.721418in}}%
\pgfpathlineto{\pgfqpoint{5.960777in}{0.743974in}}%
\pgfpathlineto{\pgfqpoint{5.964133in}{0.752036in}}%
\pgfpathlineto{\pgfqpoint{5.967489in}{0.747022in}}%
\pgfpathlineto{\pgfqpoint{5.970845in}{0.727869in}}%
\pgfpathlineto{\pgfqpoint{5.974200in}{0.678357in}}%
\pgfpathlineto{\pgfqpoint{5.977556in}{0.689170in}}%
\pgfpathlineto{\pgfqpoint{5.980912in}{0.716808in}}%
\pgfpathlineto{\pgfqpoint{5.984268in}{0.721096in}}%
\pgfpathlineto{\pgfqpoint{5.987623in}{0.709841in}}%
\pgfpathlineto{\pgfqpoint{5.990979in}{0.654851in}}%
\pgfpathlineto{\pgfqpoint{5.994335in}{0.701003in}}%
\pgfpathlineto{\pgfqpoint{5.997691in}{0.731865in}}%
\pgfpathlineto{\pgfqpoint{6.001047in}{0.743405in}}%
\pgfpathlineto{\pgfqpoint{6.004402in}{0.740988in}}%
\pgfpathlineto{\pgfqpoint{6.007758in}{0.724768in}}%
\pgfpathlineto{\pgfqpoint{6.011114in}{0.691205in}}%
\pgfpathlineto{\pgfqpoint{6.014470in}{0.636541in}}%
\pgfpathlineto{\pgfqpoint{6.017825in}{0.691617in}}%
\pgfpathlineto{\pgfqpoint{6.021181in}{0.698669in}}%
\pgfpathlineto{\pgfqpoint{6.024537in}{0.685161in}}%
\pgfpathlineto{\pgfqpoint{6.027893in}{0.515409in}}%
\pgfpathlineto{\pgfqpoint{6.031249in}{0.694922in}}%
\pgfpathlineto{\pgfqpoint{6.034604in}{0.719908in}}%
\pgfpathlineto{\pgfqpoint{6.037960in}{0.728434in}}%
\pgfpathlineto{\pgfqpoint{6.041316in}{0.721770in}}%
\pgfpathlineto{\pgfqpoint{6.044672in}{0.699737in}}%
\pgfpathlineto{\pgfqpoint{6.048027in}{0.628198in}}%
\pgfpathlineto{\pgfqpoint{6.051383in}{0.677683in}}%
\pgfpathlineto{\pgfqpoint{6.054739in}{0.698642in}}%
\pgfpathlineto{\pgfqpoint{6.058095in}{0.700898in}}%
\pgfpathlineto{\pgfqpoint{6.061451in}{0.688043in}}%
\pgfpathlineto{\pgfqpoint{6.064806in}{0.624518in}}%
\pgfpathlineto{\pgfqpoint{6.068162in}{0.682638in}}%
\pgfpathlineto{\pgfqpoint{6.071518in}{0.711734in}}%
\pgfpathlineto{\pgfqpoint{6.074874in}{0.722537in}}%
\pgfpathlineto{\pgfqpoint{6.078230in}{0.719409in}}%
\pgfpathlineto{\pgfqpoint{6.081585in}{0.697641in}}%
\pgfpathlineto{\pgfqpoint{6.084941in}{0.535517in}}%
\pgfpathlineto{\pgfqpoint{6.088297in}{0.694821in}}%
\pgfpathlineto{\pgfqpoint{6.091653in}{0.715001in}}%
\pgfpathlineto{\pgfqpoint{6.095008in}{0.720803in}}%
\pgfpathlineto{\pgfqpoint{6.098364in}{0.716910in}}%
\pgfpathlineto{\pgfqpoint{6.101720in}{0.699830in}}%
\pgfpathlineto{\pgfqpoint{6.105076in}{0.615833in}}%
\pgfpathlineto{\pgfqpoint{6.108432in}{0.697716in}}%
\pgfpathlineto{\pgfqpoint{6.111787in}{0.724375in}}%
\pgfpathlineto{\pgfqpoint{6.115143in}{0.734708in}}%
\pgfpathlineto{\pgfqpoint{6.118499in}{0.732818in}}%
\pgfpathlineto{\pgfqpoint{6.121855in}{0.713115in}}%
\pgfpathlineto{\pgfqpoint{6.125210in}{0.603112in}}%
\pgfpathlineto{\pgfqpoint{6.128566in}{0.708080in}}%
\pgfpathlineto{\pgfqpoint{6.131922in}{0.730041in}}%
\pgfpathlineto{\pgfqpoint{6.135278in}{0.737207in}}%
\pgfpathlineto{\pgfqpoint{6.138634in}{0.735165in}}%
\pgfpathlineto{\pgfqpoint{6.141989in}{0.722138in}}%
\pgfpathlineto{\pgfqpoint{6.145345in}{0.677772in}}%
\pgfpathlineto{\pgfqpoint{6.148701in}{0.699921in}}%
\pgfpathlineto{\pgfqpoint{6.152057in}{0.734917in}}%
\pgfpathlineto{\pgfqpoint{6.155412in}{0.748690in}}%
\pgfpathlineto{\pgfqpoint{6.158768in}{0.750781in}}%
\pgfpathlineto{\pgfqpoint{6.162124in}{0.738629in}}%
\pgfpathlineto{\pgfqpoint{6.165480in}{0.706398in}}%
\pgfpathlineto{\pgfqpoint{6.168836in}{0.668696in}}%
\pgfpathlineto{\pgfqpoint{6.172191in}{0.719576in}}%
\pgfpathlineto{\pgfqpoint{6.175547in}{0.731746in}}%
\pgfpathlineto{\pgfqpoint{6.178903in}{0.730592in}}%
\pgfpathlineto{\pgfqpoint{6.182259in}{0.714910in}}%
\pgfpathlineto{\pgfqpoint{6.185614in}{0.629800in}}%
\pgfpathlineto{\pgfqpoint{6.188970in}{0.714536in}}%
\pgfpathlineto{\pgfqpoint{6.192326in}{0.741118in}}%
\pgfpathlineto{\pgfqpoint{6.195682in}{0.750906in}}%
\pgfpathlineto{\pgfqpoint{6.199038in}{0.747115in}}%
\pgfpathlineto{\pgfqpoint{6.202393in}{0.730122in}}%
\pgfpathlineto{\pgfqpoint{6.205749in}{0.694458in}}%
\pgfpathlineto{\pgfqpoint{6.209105in}{0.652664in}}%
\pgfpathlineto{\pgfqpoint{6.212461in}{0.699514in}}%
\pgfpathlineto{\pgfqpoint{6.215816in}{0.705106in}}%
\pgfpathlineto{\pgfqpoint{6.219172in}{0.689653in}}%
\pgfpathlineto{\pgfqpoint{6.222528in}{0.600503in}}%
\pgfpathlineto{\pgfqpoint{6.225884in}{0.704555in}}%
\pgfpathlineto{\pgfqpoint{6.229240in}{0.727839in}}%
\pgfpathlineto{\pgfqpoint{6.232595in}{0.735330in}}%
\pgfpathlineto{\pgfqpoint{6.235951in}{0.727559in}}%
\pgfpathlineto{\pgfqpoint{6.239307in}{0.706410in}}%
\pgfpathlineto{\pgfqpoint{6.242663in}{0.652850in}}%
\pgfpathlineto{\pgfqpoint{6.246019in}{0.667798in}}%
\pgfpathlineto{\pgfqpoint{6.249374in}{0.692002in}}%
\pgfpathlineto{\pgfqpoint{6.252730in}{0.692294in}}%
\pgfpathlineto{\pgfqpoint{6.256086in}{0.670739in}}%
\pgfpathlineto{\pgfqpoint{6.259442in}{0.634489in}}%
\pgfpathlineto{\pgfqpoint{6.262797in}{0.697008in}}%
\pgfpathlineto{\pgfqpoint{6.266153in}{0.716672in}}%
\pgfpathlineto{\pgfqpoint{6.269509in}{0.721715in}}%
\pgfpathlineto{\pgfqpoint{6.272865in}{0.710704in}}%
\pgfpathlineto{\pgfqpoint{6.279576in}{0.648859in}}%
\pgfpathlineto{\pgfqpoint{6.282932in}{0.696285in}}%
\pgfpathlineto{\pgfqpoint{6.286288in}{0.708924in}}%
\pgfpathlineto{\pgfqpoint{6.289644in}{0.709221in}}%
\pgfpathlineto{\pgfqpoint{6.292999in}{0.697339in}}%
\pgfpathlineto{\pgfqpoint{6.296355in}{0.652767in}}%
\pgfpathlineto{\pgfqpoint{6.299711in}{0.677305in}}%
\pgfpathlineto{\pgfqpoint{6.303067in}{0.711905in}}%
\pgfpathlineto{\pgfqpoint{6.306423in}{0.725188in}}%
\pgfpathlineto{\pgfqpoint{6.309778in}{0.725868in}}%
\pgfpathlineto{\pgfqpoint{6.313134in}{0.709210in}}%
\pgfpathlineto{\pgfqpoint{6.316490in}{0.635756in}}%
\pgfpathlineto{\pgfqpoint{6.319846in}{0.698145in}}%
\pgfpathlineto{\pgfqpoint{6.323201in}{0.723100in}}%
\pgfpathlineto{\pgfqpoint{6.326557in}{0.731902in}}%
\pgfpathlineto{\pgfqpoint{6.329913in}{0.731636in}}%
\pgfpathlineto{\pgfqpoint{6.333269in}{0.721780in}}%
\pgfpathlineto{\pgfqpoint{6.336625in}{0.691512in}}%
\pgfpathlineto{\pgfqpoint{6.339980in}{0.673721in}}%
\pgfpathlineto{\pgfqpoint{6.343336in}{0.723169in}}%
\pgfpathlineto{\pgfqpoint{6.346692in}{0.740734in}}%
\pgfpathlineto{\pgfqpoint{6.350048in}{0.746142in}}%
\pgfpathlineto{\pgfqpoint{6.353403in}{0.738991in}}%
\pgfpathlineto{\pgfqpoint{6.356759in}{0.723852in}}%
\pgfpathlineto{\pgfqpoint{6.360115in}{0.702036in}}%
\pgfpathlineto{\pgfqpoint{6.363471in}{0.670511in}}%
\pgfpathlineto{\pgfqpoint{6.366827in}{0.629531in}}%
\pgfpathlineto{\pgfqpoint{6.370182in}{0.628842in}}%
\pgfpathlineto{\pgfqpoint{6.373538in}{0.664388in}}%
\pgfpathlineto{\pgfqpoint{6.376894in}{0.688733in}}%
\pgfpathlineto{\pgfqpoint{6.380250in}{0.700953in}}%
\pgfpathlineto{\pgfqpoint{6.383605in}{0.699836in}}%
\pgfpathlineto{\pgfqpoint{6.383605in}{0.699836in}}%
\pgfusepath{stroke}%
\end{pgfscope}%
\begin{pgfscope}%
\pgfpathrectangle{\pgfqpoint{4.621821in}{0.320625in}}{\pgfqpoint{1.845679in}{1.164375in}}%
\pgfusepath{clip}%
\pgfsetrectcap%
\pgfsetroundjoin%
\pgfsetlinewidth{0.602250pt}%
\definecolor{currentstroke}{rgb}{0.000000,0.529412,0.423529}%
\pgfsetstrokecolor{currentstroke}%
\pgfsetstrokeopacity{0.120000}%
\pgfsetdash{}{0pt}%
\pgfpathmoveto{\pgfqpoint{4.705715in}{1.468955in}}%
\pgfpathlineto{\pgfqpoint{4.719139in}{1.444562in}}%
\pgfpathlineto{\pgfqpoint{4.729206in}{1.423574in}}%
\pgfpathlineto{\pgfqpoint{4.745985in}{1.384957in}}%
\pgfpathlineto{\pgfqpoint{4.776187in}{1.314994in}}%
\pgfpathlineto{\pgfqpoint{4.799677in}{1.263905in}}%
\pgfpathlineto{\pgfqpoint{4.823168in}{1.215450in}}%
\pgfpathlineto{\pgfqpoint{4.853370in}{1.155847in}}%
\pgfpathlineto{\pgfqpoint{4.890283in}{1.085674in}}%
\pgfpathlineto{\pgfqpoint{4.937264in}{0.998965in}}%
\pgfpathlineto{\pgfqpoint{4.997668in}{0.889908in}}%
\pgfpathlineto{\pgfqpoint{5.061428in}{0.776667in}}%
\pgfpathlineto{\pgfqpoint{5.074851in}{0.755168in}}%
\pgfpathlineto{\pgfqpoint{5.084919in}{0.734444in}}%
\pgfpathlineto{\pgfqpoint{5.091630in}{0.720301in}}%
\pgfpathlineto{\pgfqpoint{5.098342in}{0.709049in}}%
\pgfpathlineto{\pgfqpoint{5.101698in}{0.705351in}}%
\pgfpathlineto{\pgfqpoint{5.105053in}{0.702949in}}%
\pgfpathlineto{\pgfqpoint{5.111765in}{0.699686in}}%
\pgfpathlineto{\pgfqpoint{5.115121in}{0.696622in}}%
\pgfpathlineto{\pgfqpoint{5.118476in}{0.690672in}}%
\pgfpathlineto{\pgfqpoint{5.121832in}{0.682573in}}%
\pgfpathlineto{\pgfqpoint{5.125188in}{0.672270in}}%
\pgfpathlineto{\pgfqpoint{5.128544in}{0.659483in}}%
\pgfpathlineto{\pgfqpoint{5.135255in}{0.630478in}}%
\pgfpathlineto{\pgfqpoint{5.138611in}{0.624482in}}%
\pgfpathlineto{\pgfqpoint{5.141967in}{0.624537in}}%
\pgfpathlineto{\pgfqpoint{5.145323in}{0.622558in}}%
\pgfpathlineto{\pgfqpoint{5.148678in}{0.614294in}}%
\pgfpathlineto{\pgfqpoint{5.152034in}{0.598509in}}%
\pgfpathlineto{\pgfqpoint{5.155390in}{0.603074in}}%
\pgfpathlineto{\pgfqpoint{5.158746in}{0.627651in}}%
\pgfpathlineto{\pgfqpoint{5.162102in}{0.645059in}}%
\pgfpathlineto{\pgfqpoint{5.165457in}{0.656350in}}%
\pgfpathlineto{\pgfqpoint{5.168813in}{0.663034in}}%
\pgfpathlineto{\pgfqpoint{5.172169in}{0.665389in}}%
\pgfpathlineto{\pgfqpoint{5.175525in}{0.662363in}}%
\pgfpathlineto{\pgfqpoint{5.178880in}{0.651730in}}%
\pgfpathlineto{\pgfqpoint{5.182236in}{0.665825in}}%
\pgfpathlineto{\pgfqpoint{5.185592in}{0.691682in}}%
\pgfpathlineto{\pgfqpoint{5.188948in}{0.709412in}}%
\pgfpathlineto{\pgfqpoint{5.192304in}{0.721491in}}%
\pgfpathlineto{\pgfqpoint{5.195659in}{0.729826in}}%
\pgfpathlineto{\pgfqpoint{5.199015in}{0.735405in}}%
\pgfpathlineto{\pgfqpoint{5.202371in}{0.738739in}}%
\pgfpathlineto{\pgfqpoint{5.205727in}{0.740057in}}%
\pgfpathlineto{\pgfqpoint{5.209082in}{0.739383in}}%
\pgfpathlineto{\pgfqpoint{5.212438in}{0.736516in}}%
\pgfpathlineto{\pgfqpoint{5.215794in}{0.730923in}}%
\pgfpathlineto{\pgfqpoint{5.219150in}{0.721359in}}%
\pgfpathlineto{\pgfqpoint{5.222506in}{0.704475in}}%
\pgfpathlineto{\pgfqpoint{5.225861in}{0.665621in}}%
\pgfpathlineto{\pgfqpoint{5.229217in}{0.672965in}}%
\pgfpathlineto{\pgfqpoint{5.232573in}{0.710378in}}%
\pgfpathlineto{\pgfqpoint{5.235929in}{0.728629in}}%
\pgfpathlineto{\pgfqpoint{5.239285in}{0.739902in}}%
\pgfpathlineto{\pgfqpoint{5.242640in}{0.747087in}}%
\pgfpathlineto{\pgfqpoint{5.245996in}{0.751122in}}%
\pgfpathlineto{\pgfqpoint{5.249352in}{0.752116in}}%
\pgfpathlineto{\pgfqpoint{5.252708in}{0.749439in}}%
\pgfpathlineto{\pgfqpoint{5.256063in}{0.741322in}}%
\pgfpathlineto{\pgfqpoint{5.259419in}{0.728568in}}%
\pgfpathlineto{\pgfqpoint{5.266131in}{0.697051in}}%
\pgfpathlineto{\pgfqpoint{5.269487in}{0.687722in}}%
\pgfpathlineto{\pgfqpoint{5.276198in}{0.674263in}}%
\pgfpathlineto{\pgfqpoint{5.279554in}{0.684036in}}%
\pgfpathlineto{\pgfqpoint{5.282910in}{0.699717in}}%
\pgfpathlineto{\pgfqpoint{5.286265in}{0.707512in}}%
\pgfpathlineto{\pgfqpoint{5.289621in}{0.702541in}}%
\pgfpathlineto{\pgfqpoint{5.292977in}{0.688156in}}%
\pgfpathlineto{\pgfqpoint{5.296333in}{0.667035in}}%
\pgfpathlineto{\pgfqpoint{5.299689in}{0.641262in}}%
\pgfpathlineto{\pgfqpoint{5.303044in}{0.631801in}}%
\pgfpathlineto{\pgfqpoint{5.306400in}{0.628050in}}%
\pgfpathlineto{\pgfqpoint{5.309756in}{0.614648in}}%
\pgfpathlineto{\pgfqpoint{5.316467in}{0.662371in}}%
\pgfpathlineto{\pgfqpoint{5.319823in}{0.673013in}}%
\pgfpathlineto{\pgfqpoint{5.323179in}{0.669635in}}%
\pgfpathlineto{\pgfqpoint{5.326535in}{0.636168in}}%
\pgfpathlineto{\pgfqpoint{5.329891in}{0.645157in}}%
\pgfpathlineto{\pgfqpoint{5.333246in}{0.682899in}}%
\pgfpathlineto{\pgfqpoint{5.336602in}{0.698718in}}%
\pgfpathlineto{\pgfqpoint{5.339958in}{0.705675in}}%
\pgfpathlineto{\pgfqpoint{5.343314in}{0.706407in}}%
\pgfpathlineto{\pgfqpoint{5.346669in}{0.700793in}}%
\pgfpathlineto{\pgfqpoint{5.350025in}{0.685497in}}%
\pgfpathlineto{\pgfqpoint{5.353381in}{0.638717in}}%
\pgfpathlineto{\pgfqpoint{5.356737in}{0.662950in}}%
\pgfpathlineto{\pgfqpoint{5.360093in}{0.697515in}}%
\pgfpathlineto{\pgfqpoint{5.363448in}{0.712290in}}%
\pgfpathlineto{\pgfqpoint{5.366804in}{0.717511in}}%
\pgfpathlineto{\pgfqpoint{5.370160in}{0.712639in}}%
\pgfpathlineto{\pgfqpoint{5.373516in}{0.692357in}}%
\pgfpathlineto{\pgfqpoint{5.376871in}{0.620610in}}%
\pgfpathlineto{\pgfqpoint{5.380227in}{0.680414in}}%
\pgfpathlineto{\pgfqpoint{5.383583in}{0.702368in}}%
\pgfpathlineto{\pgfqpoint{5.386939in}{0.708517in}}%
\pgfpathlineto{\pgfqpoint{5.390295in}{0.704430in}}%
\pgfpathlineto{\pgfqpoint{5.393650in}{0.686239in}}%
\pgfpathlineto{\pgfqpoint{5.397006in}{0.567030in}}%
\pgfpathlineto{\pgfqpoint{5.400362in}{0.689365in}}%
\pgfpathlineto{\pgfqpoint{5.403718in}{0.714375in}}%
\pgfpathlineto{\pgfqpoint{5.407074in}{0.723882in}}%
\pgfpathlineto{\pgfqpoint{5.410429in}{0.720855in}}%
\pgfpathlineto{\pgfqpoint{5.413785in}{0.701183in}}%
\pgfpathlineto{\pgfqpoint{5.417141in}{0.628936in}}%
\pgfpathlineto{\pgfqpoint{5.420497in}{0.687446in}}%
\pgfpathlineto{\pgfqpoint{5.423852in}{0.710301in}}%
\pgfpathlineto{\pgfqpoint{5.427208in}{0.716856in}}%
\pgfpathlineto{\pgfqpoint{5.430564in}{0.713187in}}%
\pgfpathlineto{\pgfqpoint{5.433920in}{0.695869in}}%
\pgfpathlineto{\pgfqpoint{5.437276in}{0.602284in}}%
\pgfpathlineto{\pgfqpoint{5.440631in}{0.695904in}}%
\pgfpathlineto{\pgfqpoint{5.443987in}{0.721953in}}%
\pgfpathlineto{\pgfqpoint{5.447343in}{0.731966in}}%
\pgfpathlineto{\pgfqpoint{5.450699in}{0.729563in}}%
\pgfpathlineto{\pgfqpoint{5.454054in}{0.710497in}}%
\pgfpathlineto{\pgfqpoint{5.457410in}{0.640105in}}%
\pgfpathlineto{\pgfqpoint{5.460766in}{0.695489in}}%
\pgfpathlineto{\pgfqpoint{5.464122in}{0.719128in}}%
\pgfpathlineto{\pgfqpoint{5.467478in}{0.726119in}}%
\pgfpathlineto{\pgfqpoint{5.470833in}{0.722976in}}%
\pgfpathlineto{\pgfqpoint{5.474189in}{0.706819in}}%
\pgfpathlineto{\pgfqpoint{5.477545in}{0.633383in}}%
\pgfpathlineto{\pgfqpoint{5.480901in}{0.702250in}}%
\pgfpathlineto{\pgfqpoint{5.484256in}{0.729946in}}%
\pgfpathlineto{\pgfqpoint{5.487612in}{0.740750in}}%
\pgfpathlineto{\pgfqpoint{5.490968in}{0.739351in}}%
\pgfpathlineto{\pgfqpoint{5.494324in}{0.723074in}}%
\pgfpathlineto{\pgfqpoint{5.497680in}{0.680971in}}%
\pgfpathlineto{\pgfqpoint{5.501035in}{0.678611in}}%
\pgfpathlineto{\pgfqpoint{5.504391in}{0.712207in}}%
\pgfpathlineto{\pgfqpoint{5.507747in}{0.720384in}}%
\pgfpathlineto{\pgfqpoint{5.511103in}{0.715744in}}%
\pgfpathlineto{\pgfqpoint{5.514458in}{0.692656in}}%
\pgfpathlineto{\pgfqpoint{5.517814in}{0.647447in}}%
\pgfpathlineto{\pgfqpoint{5.521170in}{0.713996in}}%
\pgfpathlineto{\pgfqpoint{5.524526in}{0.734020in}}%
\pgfpathlineto{\pgfqpoint{5.527882in}{0.740034in}}%
\pgfpathlineto{\pgfqpoint{5.531237in}{0.731770in}}%
\pgfpathlineto{\pgfqpoint{5.534593in}{0.709910in}}%
\pgfpathlineto{\pgfqpoint{5.537949in}{0.649697in}}%
\pgfpathlineto{\pgfqpoint{5.541305in}{0.679187in}}%
\pgfpathlineto{\pgfqpoint{5.544660in}{0.701218in}}%
\pgfpathlineto{\pgfqpoint{5.548016in}{0.702054in}}%
\pgfpathlineto{\pgfqpoint{5.551372in}{0.684111in}}%
\pgfpathlineto{\pgfqpoint{5.554728in}{0.606601in}}%
\pgfpathlineto{\pgfqpoint{5.558084in}{0.699202in}}%
\pgfpathlineto{\pgfqpoint{5.561439in}{0.721693in}}%
\pgfpathlineto{\pgfqpoint{5.564795in}{0.728933in}}%
\pgfpathlineto{\pgfqpoint{5.568151in}{0.721321in}}%
\pgfpathlineto{\pgfqpoint{5.571507in}{0.699158in}}%
\pgfpathlineto{\pgfqpoint{5.574863in}{0.631667in}}%
\pgfpathlineto{\pgfqpoint{5.578218in}{0.674430in}}%
\pgfpathlineto{\pgfqpoint{5.581574in}{0.695509in}}%
\pgfpathlineto{\pgfqpoint{5.584930in}{0.697188in}}%
\pgfpathlineto{\pgfqpoint{5.588286in}{0.682602in}}%
\pgfpathlineto{\pgfqpoint{5.591641in}{0.590584in}}%
\pgfpathlineto{\pgfqpoint{5.594997in}{0.685048in}}%
\pgfpathlineto{\pgfqpoint{5.598353in}{0.711378in}}%
\pgfpathlineto{\pgfqpoint{5.601709in}{0.720897in}}%
\pgfpathlineto{\pgfqpoint{5.605065in}{0.716201in}}%
\pgfpathlineto{\pgfqpoint{5.608420in}{0.694313in}}%
\pgfpathlineto{\pgfqpoint{5.611776in}{0.596974in}}%
\pgfpathlineto{\pgfqpoint{5.615132in}{0.683999in}}%
\pgfpathlineto{\pgfqpoint{5.618488in}{0.704532in}}%
\pgfpathlineto{\pgfqpoint{5.621843in}{0.709326in}}%
\pgfpathlineto{\pgfqpoint{5.625199in}{0.703095in}}%
\pgfpathlineto{\pgfqpoint{5.628555in}{0.678914in}}%
\pgfpathlineto{\pgfqpoint{5.631911in}{0.635222in}}%
\pgfpathlineto{\pgfqpoint{5.635267in}{0.700215in}}%
\pgfpathlineto{\pgfqpoint{5.638622in}{0.720047in}}%
\pgfpathlineto{\pgfqpoint{5.641978in}{0.726269in}}%
\pgfpathlineto{\pgfqpoint{5.645334in}{0.718411in}}%
\pgfpathlineto{\pgfqpoint{5.648690in}{0.689435in}}%
\pgfpathlineto{\pgfqpoint{5.652045in}{0.649117in}}%
\pgfpathlineto{\pgfqpoint{5.655401in}{0.704920in}}%
\pgfpathlineto{\pgfqpoint{5.658757in}{0.720251in}}%
\pgfpathlineto{\pgfqpoint{5.662113in}{0.723605in}}%
\pgfpathlineto{\pgfqpoint{5.665469in}{0.717303in}}%
\pgfpathlineto{\pgfqpoint{5.668824in}{0.695304in}}%
\pgfpathlineto{\pgfqpoint{5.672180in}{0.622588in}}%
\pgfpathlineto{\pgfqpoint{5.675536in}{0.708528in}}%
\pgfpathlineto{\pgfqpoint{5.678892in}{0.730352in}}%
\pgfpathlineto{\pgfqpoint{5.682247in}{0.738242in}}%
\pgfpathlineto{\pgfqpoint{5.685603in}{0.733442in}}%
\pgfpathlineto{\pgfqpoint{5.688959in}{0.712200in}}%
\pgfpathlineto{\pgfqpoint{5.692315in}{0.625078in}}%
\pgfpathlineto{\pgfqpoint{5.695671in}{0.699623in}}%
\pgfpathlineto{\pgfqpoint{5.699026in}{0.720952in}}%
\pgfpathlineto{\pgfqpoint{5.702382in}{0.726093in}}%
\pgfpathlineto{\pgfqpoint{5.705738in}{0.720201in}}%
\pgfpathlineto{\pgfqpoint{5.709094in}{0.696824in}}%
\pgfpathlineto{\pgfqpoint{5.712449in}{0.647073in}}%
\pgfpathlineto{\pgfqpoint{5.715805in}{0.716191in}}%
\pgfpathlineto{\pgfqpoint{5.719161in}{0.736503in}}%
\pgfpathlineto{\pgfqpoint{5.722517in}{0.743037in}}%
\pgfpathlineto{\pgfqpoint{5.725873in}{0.735866in}}%
\pgfpathlineto{\pgfqpoint{5.729228in}{0.714734in}}%
\pgfpathlineto{\pgfqpoint{5.732584in}{0.655768in}}%
\pgfpathlineto{\pgfqpoint{5.735940in}{0.684047in}}%
\pgfpathlineto{\pgfqpoint{5.739296in}{0.707310in}}%
\pgfpathlineto{\pgfqpoint{5.742652in}{0.709600in}}%
\pgfpathlineto{\pgfqpoint{5.746007in}{0.695340in}}%
\pgfpathlineto{\pgfqpoint{5.749363in}{0.604613in}}%
\pgfpathlineto{\pgfqpoint{5.752719in}{0.697792in}}%
\pgfpathlineto{\pgfqpoint{5.756075in}{0.724242in}}%
\pgfpathlineto{\pgfqpoint{5.759430in}{0.733790in}}%
\pgfpathlineto{\pgfqpoint{5.762786in}{0.729094in}}%
\pgfpathlineto{\pgfqpoint{5.766142in}{0.710749in}}%
\pgfpathlineto{\pgfqpoint{5.769498in}{0.668552in}}%
\pgfpathlineto{\pgfqpoint{5.772854in}{0.657293in}}%
\pgfpathlineto{\pgfqpoint{5.776209in}{0.690719in}}%
\pgfpathlineto{\pgfqpoint{5.779565in}{0.695200in}}%
\pgfpathlineto{\pgfqpoint{5.782921in}{0.681706in}}%
\pgfpathlineto{\pgfqpoint{5.786277in}{0.585588in}}%
\pgfpathlineto{\pgfqpoint{5.789632in}{0.685945in}}%
\pgfpathlineto{\pgfqpoint{5.792988in}{0.712259in}}%
\pgfpathlineto{\pgfqpoint{5.796344in}{0.721672in}}%
\pgfpathlineto{\pgfqpoint{5.799700in}{0.716491in}}%
\pgfpathlineto{\pgfqpoint{5.803056in}{0.695367in}}%
\pgfpathlineto{\pgfqpoint{5.806411in}{0.624593in}}%
\pgfpathlineto{\pgfqpoint{5.809767in}{0.675022in}}%
\pgfpathlineto{\pgfqpoint{5.813123in}{0.697143in}}%
\pgfpathlineto{\pgfqpoint{5.816479in}{0.701260in}}%
\pgfpathlineto{\pgfqpoint{5.819834in}{0.692566in}}%
\pgfpathlineto{\pgfqpoint{5.823190in}{0.657765in}}%
\pgfpathlineto{\pgfqpoint{5.826546in}{0.661518in}}%
\pgfpathlineto{\pgfqpoint{5.829902in}{0.702633in}}%
\pgfpathlineto{\pgfqpoint{5.833258in}{0.717883in}}%
\pgfpathlineto{\pgfqpoint{5.836613in}{0.720097in}}%
\pgfpathlineto{\pgfqpoint{5.839969in}{0.706127in}}%
\pgfpathlineto{\pgfqpoint{5.843325in}{0.660430in}}%
\pgfpathlineto{\pgfqpoint{5.846681in}{0.674595in}}%
\pgfpathlineto{\pgfqpoint{5.850036in}{0.706514in}}%
\pgfpathlineto{\pgfqpoint{5.853392in}{0.716795in}}%
\pgfpathlineto{\pgfqpoint{5.856748in}{0.716752in}}%
\pgfpathlineto{\pgfqpoint{5.860104in}{0.706100in}}%
\pgfpathlineto{\pgfqpoint{5.863460in}{0.670909in}}%
\pgfpathlineto{\pgfqpoint{5.866815in}{0.671565in}}%
\pgfpathlineto{\pgfqpoint{5.870171in}{0.713274in}}%
\pgfpathlineto{\pgfqpoint{5.873527in}{0.728960in}}%
\pgfpathlineto{\pgfqpoint{5.876883in}{0.732535in}}%
\pgfpathlineto{\pgfqpoint{5.880238in}{0.721687in}}%
\pgfpathlineto{\pgfqpoint{5.883594in}{0.686678in}}%
\pgfpathlineto{\pgfqpoint{5.886950in}{0.671740in}}%
\pgfpathlineto{\pgfqpoint{5.890306in}{0.713836in}}%
\pgfpathlineto{\pgfqpoint{5.893662in}{0.726592in}}%
\pgfpathlineto{\pgfqpoint{5.897017in}{0.728168in}}%
\pgfpathlineto{\pgfqpoint{5.900373in}{0.719538in}}%
\pgfpathlineto{\pgfqpoint{5.903729in}{0.690962in}}%
\pgfpathlineto{\pgfqpoint{5.907085in}{0.668303in}}%
\pgfpathlineto{\pgfqpoint{5.910441in}{0.720459in}}%
\pgfpathlineto{\pgfqpoint{5.913796in}{0.738419in}}%
\pgfpathlineto{\pgfqpoint{5.917152in}{0.743708in}}%
\pgfpathlineto{\pgfqpoint{5.920508in}{0.735482in}}%
\pgfpathlineto{\pgfqpoint{5.923864in}{0.711704in}}%
\pgfpathlineto{\pgfqpoint{5.927219in}{0.608775in}}%
\pgfpathlineto{\pgfqpoint{5.930575in}{0.698585in}}%
\pgfpathlineto{\pgfqpoint{5.933931in}{0.717133in}}%
\pgfpathlineto{\pgfqpoint{5.937287in}{0.718989in}}%
\pgfpathlineto{\pgfqpoint{5.940643in}{0.706723in}}%
\pgfpathlineto{\pgfqpoint{5.943998in}{0.652212in}}%
\pgfpathlineto{\pgfqpoint{5.947354in}{0.696240in}}%
\pgfpathlineto{\pgfqpoint{5.950710in}{0.727248in}}%
\pgfpathlineto{\pgfqpoint{5.954066in}{0.739028in}}%
\pgfpathlineto{\pgfqpoint{5.957421in}{0.737479in}}%
\pgfpathlineto{\pgfqpoint{5.960777in}{0.721923in}}%
\pgfpathlineto{\pgfqpoint{5.964133in}{0.689372in}}%
\pgfpathlineto{\pgfqpoint{5.967489in}{0.630318in}}%
\pgfpathlineto{\pgfqpoint{5.970845in}{0.689500in}}%
\pgfpathlineto{\pgfqpoint{5.974200in}{0.698164in}}%
\pgfpathlineto{\pgfqpoint{5.977556in}{0.688127in}}%
\pgfpathlineto{\pgfqpoint{5.980912in}{0.627800in}}%
\pgfpathlineto{\pgfqpoint{5.984268in}{0.684241in}}%
\pgfpathlineto{\pgfqpoint{5.987623in}{0.714164in}}%
\pgfpathlineto{\pgfqpoint{5.990979in}{0.725337in}}%
\pgfpathlineto{\pgfqpoint{5.994335in}{0.722164in}}%
\pgfpathlineto{\pgfqpoint{5.997691in}{0.704336in}}%
\pgfpathlineto{\pgfqpoint{6.001047in}{0.661331in}}%
\pgfpathlineto{\pgfqpoint{6.004402in}{0.655626in}}%
\pgfpathlineto{\pgfqpoint{6.007758in}{0.688416in}}%
\pgfpathlineto{\pgfqpoint{6.011114in}{0.694226in}}%
\pgfpathlineto{\pgfqpoint{6.014470in}{0.684544in}}%
\pgfpathlineto{\pgfqpoint{6.017825in}{0.638327in}}%
\pgfpathlineto{\pgfqpoint{6.024537in}{0.703589in}}%
\pgfpathlineto{\pgfqpoint{6.027893in}{0.716334in}}%
\pgfpathlineto{\pgfqpoint{6.031249in}{0.715447in}}%
\pgfpathlineto{\pgfqpoint{6.034604in}{0.697511in}}%
\pgfpathlineto{\pgfqpoint{6.037960in}{0.636932in}}%
\pgfpathlineto{\pgfqpoint{6.041316in}{0.676119in}}%
\pgfpathlineto{\pgfqpoint{6.044672in}{0.701896in}}%
\pgfpathlineto{\pgfqpoint{6.048027in}{0.709398in}}%
\pgfpathlineto{\pgfqpoint{6.051383in}{0.706355in}}%
\pgfpathlineto{\pgfqpoint{6.054739in}{0.689886in}}%
\pgfpathlineto{\pgfqpoint{6.058095in}{0.609261in}}%
\pgfpathlineto{\pgfqpoint{6.061451in}{0.687451in}}%
\pgfpathlineto{\pgfqpoint{6.064806in}{0.714409in}}%
\pgfpathlineto{\pgfqpoint{6.068162in}{0.724773in}}%
\pgfpathlineto{\pgfqpoint{6.071518in}{0.722616in}}%
\pgfpathlineto{\pgfqpoint{6.074874in}{0.702539in}}%
\pgfpathlineto{\pgfqpoint{6.078230in}{0.587455in}}%
\pgfpathlineto{\pgfqpoint{6.081585in}{0.697729in}}%
\pgfpathlineto{\pgfqpoint{6.084941in}{0.719504in}}%
\pgfpathlineto{\pgfqpoint{6.088297in}{0.726580in}}%
\pgfpathlineto{\pgfqpoint{6.091653in}{0.724473in}}%
\pgfpathlineto{\pgfqpoint{6.095008in}{0.711371in}}%
\pgfpathlineto{\pgfqpoint{6.098364in}{0.666692in}}%
\pgfpathlineto{\pgfqpoint{6.101720in}{0.689484in}}%
\pgfpathlineto{\pgfqpoint{6.105076in}{0.724348in}}%
\pgfpathlineto{\pgfqpoint{6.108432in}{0.738126in}}%
\pgfpathlineto{\pgfqpoint{6.111787in}{0.740307in}}%
\pgfpathlineto{\pgfqpoint{6.115143in}{0.728169in}}%
\pgfpathlineto{\pgfqpoint{6.118499in}{0.693902in}}%
\pgfpathlineto{\pgfqpoint{6.121855in}{0.670204in}}%
\pgfpathlineto{\pgfqpoint{6.125210in}{0.714776in}}%
\pgfpathlineto{\pgfqpoint{6.128566in}{0.726839in}}%
\pgfpathlineto{\pgfqpoint{6.131922in}{0.726727in}}%
\pgfpathlineto{\pgfqpoint{6.135278in}{0.714244in}}%
\pgfpathlineto{\pgfqpoint{6.138634in}{0.666401in}}%
\pgfpathlineto{\pgfqpoint{6.145345in}{0.730449in}}%
\pgfpathlineto{\pgfqpoint{6.148701in}{0.743297in}}%
\pgfpathlineto{\pgfqpoint{6.152057in}{0.743639in}}%
\pgfpathlineto{\pgfqpoint{6.155412in}{0.729775in}}%
\pgfpathlineto{\pgfqpoint{6.158768in}{0.700068in}}%
\pgfpathlineto{\pgfqpoint{6.162124in}{0.614791in}}%
\pgfpathlineto{\pgfqpoint{6.165480in}{0.695117in}}%
\pgfpathlineto{\pgfqpoint{6.168836in}{0.706628in}}%
\pgfpathlineto{\pgfqpoint{6.172191in}{0.700430in}}%
\pgfpathlineto{\pgfqpoint{6.175547in}{0.664198in}}%
\pgfpathlineto{\pgfqpoint{6.178903in}{0.677144in}}%
\pgfpathlineto{\pgfqpoint{6.182259in}{0.716077in}}%
\pgfpathlineto{\pgfqpoint{6.185614in}{0.730655in}}%
\pgfpathlineto{\pgfqpoint{6.188970in}{0.731477in}}%
\pgfpathlineto{\pgfqpoint{6.192326in}{0.717500in}}%
\pgfpathlineto{\pgfqpoint{6.195682in}{0.688635in}}%
\pgfpathlineto{\pgfqpoint{6.199038in}{0.568860in}}%
\pgfpathlineto{\pgfqpoint{6.202393in}{0.678224in}}%
\pgfpathlineto{\pgfqpoint{6.205749in}{0.689893in}}%
\pgfpathlineto{\pgfqpoint{6.209105in}{0.682201in}}%
\pgfpathlineto{\pgfqpoint{6.212461in}{0.635139in}}%
\pgfpathlineto{\pgfqpoint{6.219172in}{0.704490in}}%
\pgfpathlineto{\pgfqpoint{6.222528in}{0.717106in}}%
\pgfpathlineto{\pgfqpoint{6.225884in}{0.715609in}}%
\pgfpathlineto{\pgfqpoint{6.229240in}{0.697918in}}%
\pgfpathlineto{\pgfqpoint{6.232595in}{0.648095in}}%
\pgfpathlineto{\pgfqpoint{6.235951in}{0.664317in}}%
\pgfpathlineto{\pgfqpoint{6.239307in}{0.692978in}}%
\pgfpathlineto{\pgfqpoint{6.242663in}{0.699976in}}%
\pgfpathlineto{\pgfqpoint{6.246019in}{0.694711in}}%
\pgfpathlineto{\pgfqpoint{6.249374in}{0.670888in}}%
\pgfpathlineto{\pgfqpoint{6.252730in}{0.629087in}}%
\pgfpathlineto{\pgfqpoint{6.256086in}{0.693224in}}%
\pgfpathlineto{\pgfqpoint{6.259442in}{0.712921in}}%
\pgfpathlineto{\pgfqpoint{6.262797in}{0.718769in}}%
\pgfpathlineto{\pgfqpoint{6.266153in}{0.709812in}}%
\pgfpathlineto{\pgfqpoint{6.269509in}{0.676458in}}%
\pgfpathlineto{\pgfqpoint{6.272865in}{0.659584in}}%
\pgfpathlineto{\pgfqpoint{6.276221in}{0.703627in}}%
\pgfpathlineto{\pgfqpoint{6.279576in}{0.717496in}}%
\pgfpathlineto{\pgfqpoint{6.282932in}{0.720467in}}%
\pgfpathlineto{\pgfqpoint{6.286288in}{0.714305in}}%
\pgfpathlineto{\pgfqpoint{6.289644in}{0.693485in}}%
\pgfpathlineto{\pgfqpoint{6.292999in}{0.590525in}}%
\pgfpathlineto{\pgfqpoint{6.296355in}{0.702241in}}%
\pgfpathlineto{\pgfqpoint{6.299711in}{0.725304in}}%
\pgfpathlineto{\pgfqpoint{6.303067in}{0.734074in}}%
\pgfpathlineto{\pgfqpoint{6.306423in}{0.730775in}}%
\pgfpathlineto{\pgfqpoint{6.309778in}{0.710011in}}%
\pgfpathlineto{\pgfqpoint{6.313134in}{0.591858in}}%
\pgfpathlineto{\pgfqpoint{6.316490in}{0.704576in}}%
\pgfpathlineto{\pgfqpoint{6.319846in}{0.725845in}}%
\pgfpathlineto{\pgfqpoint{6.323201in}{0.732322in}}%
\pgfpathlineto{\pgfqpoint{6.326557in}{0.729274in}}%
\pgfpathlineto{\pgfqpoint{6.329913in}{0.714052in}}%
\pgfpathlineto{\pgfqpoint{6.333269in}{0.653596in}}%
\pgfpathlineto{\pgfqpoint{6.336625in}{0.703973in}}%
\pgfpathlineto{\pgfqpoint{6.339980in}{0.733696in}}%
\pgfpathlineto{\pgfqpoint{6.343336in}{0.745483in}}%
\pgfpathlineto{\pgfqpoint{6.346692in}{0.745463in}}%
\pgfpathlineto{\pgfqpoint{6.350048in}{0.731899in}}%
\pgfpathlineto{\pgfqpoint{6.353403in}{0.704123in}}%
\pgfpathlineto{\pgfqpoint{6.356759in}{0.531721in}}%
\pgfpathlineto{\pgfqpoint{6.360115in}{0.689573in}}%
\pgfpathlineto{\pgfqpoint{6.363471in}{0.702282in}}%
\pgfpathlineto{\pgfqpoint{6.366827in}{0.694976in}}%
\pgfpathlineto{\pgfqpoint{6.370182in}{0.648392in}}%
\pgfpathlineto{\pgfqpoint{6.376894in}{0.717458in}}%
\pgfpathlineto{\pgfqpoint{6.380250in}{0.730149in}}%
\pgfpathlineto{\pgfqpoint{6.383605in}{0.728699in}}%
\pgfpathlineto{\pgfqpoint{6.383605in}{0.728699in}}%
\pgfusepath{stroke}%
\end{pgfscope}%
\begin{pgfscope}%
\pgfpathrectangle{\pgfqpoint{4.621821in}{0.320625in}}{\pgfqpoint{1.845679in}{1.164375in}}%
\pgfusepath{clip}%
\pgfsetrectcap%
\pgfsetroundjoin%
\pgfsetlinewidth{0.602250pt}%
\definecolor{currentstroke}{rgb}{0.000000,0.529412,0.423529}%
\pgfsetstrokecolor{currentstroke}%
\pgfsetstrokeopacity{0.120000}%
\pgfsetdash{}{0pt}%
\pgfpathmoveto{\pgfqpoint{4.705715in}{1.469553in}}%
\pgfpathlineto{\pgfqpoint{4.719139in}{1.443911in}}%
\pgfpathlineto{\pgfqpoint{4.729206in}{1.422090in}}%
\pgfpathlineto{\pgfqpoint{4.752696in}{1.366863in}}%
\pgfpathlineto{\pgfqpoint{4.772831in}{1.320967in}}%
\pgfpathlineto{\pgfqpoint{4.789610in}{1.285221in}}%
\pgfpathlineto{\pgfqpoint{4.809745in}{1.244987in}}%
\pgfpathlineto{\pgfqpoint{4.833235in}{1.200748in}}%
\pgfpathlineto{\pgfqpoint{4.860081in}{1.152570in}}%
\pgfpathlineto{\pgfqpoint{4.893639in}{1.094621in}}%
\pgfpathlineto{\pgfqpoint{4.937264in}{1.021620in}}%
\pgfpathlineto{\pgfqpoint{4.994313in}{0.928554in}}%
\pgfpathlineto{\pgfqpoint{5.071496in}{0.805003in}}%
\pgfpathlineto{\pgfqpoint{5.084919in}{0.784559in}}%
\pgfpathlineto{\pgfqpoint{5.108409in}{0.745762in}}%
\pgfpathlineto{\pgfqpoint{5.115121in}{0.737492in}}%
\pgfpathlineto{\pgfqpoint{5.121832in}{0.729990in}}%
\pgfpathlineto{\pgfqpoint{5.125188in}{0.725289in}}%
\pgfpathlineto{\pgfqpoint{5.131900in}{0.713053in}}%
\pgfpathlineto{\pgfqpoint{5.138611in}{0.698106in}}%
\pgfpathlineto{\pgfqpoint{5.148678in}{0.673826in}}%
\pgfpathlineto{\pgfqpoint{5.152034in}{0.667158in}}%
\pgfpathlineto{\pgfqpoint{5.155390in}{0.662401in}}%
\pgfpathlineto{\pgfqpoint{5.158746in}{0.660161in}}%
\pgfpathlineto{\pgfqpoint{5.162102in}{0.660566in}}%
\pgfpathlineto{\pgfqpoint{5.165457in}{0.663010in}}%
\pgfpathlineto{\pgfqpoint{5.172169in}{0.669385in}}%
\pgfpathlineto{\pgfqpoint{5.175525in}{0.671051in}}%
\pgfpathlineto{\pgfqpoint{5.178880in}{0.670246in}}%
\pgfpathlineto{\pgfqpoint{5.182236in}{0.665474in}}%
\pgfpathlineto{\pgfqpoint{5.188948in}{0.648329in}}%
\pgfpathlineto{\pgfqpoint{5.192304in}{0.647820in}}%
\pgfpathlineto{\pgfqpoint{5.202371in}{0.675285in}}%
\pgfpathlineto{\pgfqpoint{5.205727in}{0.680332in}}%
\pgfpathlineto{\pgfqpoint{5.209082in}{0.682970in}}%
\pgfpathlineto{\pgfqpoint{5.212438in}{0.683333in}}%
\pgfpathlineto{\pgfqpoint{5.215794in}{0.681322in}}%
\pgfpathlineto{\pgfqpoint{5.219150in}{0.676496in}}%
\pgfpathlineto{\pgfqpoint{5.222506in}{0.667744in}}%
\pgfpathlineto{\pgfqpoint{5.225861in}{0.652082in}}%
\pgfpathlineto{\pgfqpoint{5.229217in}{0.617508in}}%
\pgfpathlineto{\pgfqpoint{5.232573in}{0.614031in}}%
\pgfpathlineto{\pgfqpoint{5.235929in}{0.654275in}}%
\pgfpathlineto{\pgfqpoint{5.239285in}{0.673325in}}%
\pgfpathlineto{\pgfqpoint{5.242640in}{0.684892in}}%
\pgfpathlineto{\pgfqpoint{5.245996in}{0.692113in}}%
\pgfpathlineto{\pgfqpoint{5.249352in}{0.695954in}}%
\pgfpathlineto{\pgfqpoint{5.252708in}{0.696451in}}%
\pgfpathlineto{\pgfqpoint{5.256063in}{0.692718in}}%
\pgfpathlineto{\pgfqpoint{5.259419in}{0.682318in}}%
\pgfpathlineto{\pgfqpoint{5.262775in}{0.666902in}}%
\pgfpathlineto{\pgfqpoint{5.266131in}{0.668114in}}%
\pgfpathlineto{\pgfqpoint{5.269487in}{0.681538in}}%
\pgfpathlineto{\pgfqpoint{5.272842in}{0.689351in}}%
\pgfpathlineto{\pgfqpoint{5.276198in}{0.690322in}}%
\pgfpathlineto{\pgfqpoint{5.279554in}{0.683819in}}%
\pgfpathlineto{\pgfqpoint{5.282910in}{0.665271in}}%
\pgfpathlineto{\pgfqpoint{5.286265in}{0.621202in}}%
\pgfpathlineto{\pgfqpoint{5.289621in}{0.667336in}}%
\pgfpathlineto{\pgfqpoint{5.292977in}{0.691216in}}%
\pgfpathlineto{\pgfqpoint{5.296333in}{0.701435in}}%
\pgfpathlineto{\pgfqpoint{5.299689in}{0.701300in}}%
\pgfpathlineto{\pgfqpoint{5.303044in}{0.685215in}}%
\pgfpathlineto{\pgfqpoint{5.306400in}{0.632505in}}%
\pgfpathlineto{\pgfqpoint{5.309756in}{0.682247in}}%
\pgfpathlineto{\pgfqpoint{5.313112in}{0.705429in}}%
\pgfpathlineto{\pgfqpoint{5.316467in}{0.714702in}}%
\pgfpathlineto{\pgfqpoint{5.319823in}{0.715869in}}%
\pgfpathlineto{\pgfqpoint{5.323179in}{0.709126in}}%
\pgfpathlineto{\pgfqpoint{5.326535in}{0.689223in}}%
\pgfpathlineto{\pgfqpoint{5.329891in}{0.569761in}}%
\pgfpathlineto{\pgfqpoint{5.333246in}{0.691990in}}%
\pgfpathlineto{\pgfqpoint{5.336602in}{0.716844in}}%
\pgfpathlineto{\pgfqpoint{5.339958in}{0.727120in}}%
\pgfpathlineto{\pgfqpoint{5.343314in}{0.726822in}}%
\pgfpathlineto{\pgfqpoint{5.346669in}{0.712939in}}%
\pgfpathlineto{\pgfqpoint{5.350025in}{0.679605in}}%
\pgfpathlineto{\pgfqpoint{5.353381in}{0.652749in}}%
\pgfpathlineto{\pgfqpoint{5.356737in}{0.693598in}}%
\pgfpathlineto{\pgfqpoint{5.360093in}{0.703734in}}%
\pgfpathlineto{\pgfqpoint{5.363448in}{0.700143in}}%
\pgfpathlineto{\pgfqpoint{5.366804in}{0.678323in}}%
\pgfpathlineto{\pgfqpoint{5.370160in}{0.627218in}}%
\pgfpathlineto{\pgfqpoint{5.373516in}{0.697758in}}%
\pgfpathlineto{\pgfqpoint{5.376871in}{0.718259in}}%
\pgfpathlineto{\pgfqpoint{5.380227in}{0.724351in}}%
\pgfpathlineto{\pgfqpoint{5.383583in}{0.715647in}}%
\pgfpathlineto{\pgfqpoint{5.386939in}{0.690467in}}%
\pgfpathlineto{\pgfqpoint{5.390295in}{0.611700in}}%
\pgfpathlineto{\pgfqpoint{5.393650in}{0.683524in}}%
\pgfpathlineto{\pgfqpoint{5.397006in}{0.699871in}}%
\pgfpathlineto{\pgfqpoint{5.400362in}{0.700622in}}%
\pgfpathlineto{\pgfqpoint{5.403718in}{0.686806in}}%
\pgfpathlineto{\pgfqpoint{5.407074in}{0.619801in}}%
\pgfpathlineto{\pgfqpoint{5.410429in}{0.682253in}}%
\pgfpathlineto{\pgfqpoint{5.413785in}{0.710647in}}%
\pgfpathlineto{\pgfqpoint{5.417141in}{0.721148in}}%
\pgfpathlineto{\pgfqpoint{5.420497in}{0.717810in}}%
\pgfpathlineto{\pgfqpoint{5.423852in}{0.696967in}}%
\pgfpathlineto{\pgfqpoint{5.427208in}{0.614894in}}%
\pgfpathlineto{\pgfqpoint{5.430564in}{0.686204in}}%
\pgfpathlineto{\pgfqpoint{5.433920in}{0.707370in}}%
\pgfpathlineto{\pgfqpoint{5.437276in}{0.712703in}}%
\pgfpathlineto{\pgfqpoint{5.440631in}{0.707252in}}%
\pgfpathlineto{\pgfqpoint{5.443987in}{0.685268in}}%
\pgfpathlineto{\pgfqpoint{5.447343in}{0.620894in}}%
\pgfpathlineto{\pgfqpoint{5.450699in}{0.700600in}}%
\pgfpathlineto{\pgfqpoint{5.454054in}{0.721831in}}%
\pgfpathlineto{\pgfqpoint{5.457410in}{0.728856in}}%
\pgfpathlineto{\pgfqpoint{5.460766in}{0.721976in}}%
\pgfpathlineto{\pgfqpoint{5.464122in}{0.695975in}}%
\pgfpathlineto{\pgfqpoint{5.467478in}{0.628196in}}%
\pgfpathlineto{\pgfqpoint{5.470833in}{0.701143in}}%
\pgfpathlineto{\pgfqpoint{5.474189in}{0.717653in}}%
\pgfpathlineto{\pgfqpoint{5.477545in}{0.720786in}}%
\pgfpathlineto{\pgfqpoint{5.480901in}{0.713131in}}%
\pgfpathlineto{\pgfqpoint{5.487612in}{0.659269in}}%
\pgfpathlineto{\pgfqpoint{5.490968in}{0.713641in}}%
\pgfpathlineto{\pgfqpoint{5.494324in}{0.731817in}}%
\pgfpathlineto{\pgfqpoint{5.497680in}{0.736802in}}%
\pgfpathlineto{\pgfqpoint{5.501035in}{0.727493in}}%
\pgfpathlineto{\pgfqpoint{5.504391in}{0.700776in}}%
\pgfpathlineto{\pgfqpoint{5.507747in}{0.614985in}}%
\pgfpathlineto{\pgfqpoint{5.511103in}{0.699347in}}%
\pgfpathlineto{\pgfqpoint{5.514458in}{0.714716in}}%
\pgfpathlineto{\pgfqpoint{5.517814in}{0.715173in}}%
\pgfpathlineto{\pgfqpoint{5.521170in}{0.701361in}}%
\pgfpathlineto{\pgfqpoint{5.524526in}{0.636149in}}%
\pgfpathlineto{\pgfqpoint{5.527882in}{0.695689in}}%
\pgfpathlineto{\pgfqpoint{5.531237in}{0.724400in}}%
\pgfpathlineto{\pgfqpoint{5.534593in}{0.735038in}}%
\pgfpathlineto{\pgfqpoint{5.537949in}{0.731921in}}%
\pgfpathlineto{\pgfqpoint{5.541305in}{0.714249in}}%
\pgfpathlineto{\pgfqpoint{5.544660in}{0.672009in}}%
\pgfpathlineto{\pgfqpoint{5.548016in}{0.663811in}}%
\pgfpathlineto{\pgfqpoint{5.551372in}{0.696963in}}%
\pgfpathlineto{\pgfqpoint{5.554728in}{0.702338in}}%
\pgfpathlineto{\pgfqpoint{5.558084in}{0.691302in}}%
\pgfpathlineto{\pgfqpoint{5.561439in}{0.633893in}}%
\pgfpathlineto{\pgfqpoint{5.564795in}{0.684398in}}%
\pgfpathlineto{\pgfqpoint{5.568151in}{0.714651in}}%
\pgfpathlineto{\pgfqpoint{5.571507in}{0.725774in}}%
\pgfpathlineto{\pgfqpoint{5.574863in}{0.722471in}}%
\pgfpathlineto{\pgfqpoint{5.578218in}{0.703977in}}%
\pgfpathlineto{\pgfqpoint{5.581574in}{0.656368in}}%
\pgfpathlineto{\pgfqpoint{5.584930in}{0.663669in}}%
\pgfpathlineto{\pgfqpoint{5.588286in}{0.692488in}}%
\pgfpathlineto{\pgfqpoint{5.591641in}{0.697353in}}%
\pgfpathlineto{\pgfqpoint{5.594997in}{0.686970in}}%
\pgfpathlineto{\pgfqpoint{5.598353in}{0.638037in}}%
\pgfpathlineto{\pgfqpoint{5.605065in}{0.706639in}}%
\pgfpathlineto{\pgfqpoint{5.608420in}{0.718764in}}%
\pgfpathlineto{\pgfqpoint{5.611776in}{0.716862in}}%
\pgfpathlineto{\pgfqpoint{5.615132in}{0.697981in}}%
\pgfpathlineto{\pgfqpoint{5.618488in}{0.635399in}}%
\pgfpathlineto{\pgfqpoint{5.621843in}{0.676162in}}%
\pgfpathlineto{\pgfqpoint{5.625199in}{0.700711in}}%
\pgfpathlineto{\pgfqpoint{5.628555in}{0.706828in}}%
\pgfpathlineto{\pgfqpoint{5.631911in}{0.701400in}}%
\pgfpathlineto{\pgfqpoint{5.635267in}{0.678275in}}%
\pgfpathlineto{\pgfqpoint{5.638622in}{0.628589in}}%
\pgfpathlineto{\pgfqpoint{5.641978in}{0.697750in}}%
\pgfpathlineto{\pgfqpoint{5.645334in}{0.717933in}}%
\pgfpathlineto{\pgfqpoint{5.648690in}{0.724055in}}%
\pgfpathlineto{\pgfqpoint{5.652045in}{0.715535in}}%
\pgfpathlineto{\pgfqpoint{5.658757in}{0.650457in}}%
\pgfpathlineto{\pgfqpoint{5.662113in}{0.702463in}}%
\pgfpathlineto{\pgfqpoint{5.665469in}{0.716771in}}%
\pgfpathlineto{\pgfqpoint{5.668824in}{0.719008in}}%
\pgfpathlineto{\pgfqpoint{5.672180in}{0.710762in}}%
\pgfpathlineto{\pgfqpoint{5.675536in}{0.682550in}}%
\pgfpathlineto{\pgfqpoint{5.678892in}{0.659113in}}%
\pgfpathlineto{\pgfqpoint{5.682247in}{0.711698in}}%
\pgfpathlineto{\pgfqpoint{5.685603in}{0.729587in}}%
\pgfpathlineto{\pgfqpoint{5.688959in}{0.734501in}}%
\pgfpathlineto{\pgfqpoint{5.692315in}{0.725151in}}%
\pgfpathlineto{\pgfqpoint{5.695671in}{0.696157in}}%
\pgfpathlineto{\pgfqpoint{5.699026in}{0.644732in}}%
\pgfpathlineto{\pgfqpoint{5.702382in}{0.705485in}}%
\pgfpathlineto{\pgfqpoint{5.705738in}{0.719839in}}%
\pgfpathlineto{\pgfqpoint{5.709094in}{0.720960in}}%
\pgfpathlineto{\pgfqpoint{5.712449in}{0.709836in}}%
\pgfpathlineto{\pgfqpoint{5.715805in}{0.667998in}}%
\pgfpathlineto{\pgfqpoint{5.719161in}{0.687186in}}%
\pgfpathlineto{\pgfqpoint{5.722517in}{0.723146in}}%
\pgfpathlineto{\pgfqpoint{5.725873in}{0.736766in}}%
\pgfpathlineto{\pgfqpoint{5.729228in}{0.737516in}}%
\pgfpathlineto{\pgfqpoint{5.732584in}{0.723297in}}%
\pgfpathlineto{\pgfqpoint{5.735940in}{0.689939in}}%
\pgfpathlineto{\pgfqpoint{5.739296in}{0.648684in}}%
\pgfpathlineto{\pgfqpoint{5.742652in}{0.699042in}}%
\pgfpathlineto{\pgfqpoint{5.746007in}{0.708927in}}%
\pgfpathlineto{\pgfqpoint{5.749363in}{0.703355in}}%
\pgfpathlineto{\pgfqpoint{5.752719in}{0.673424in}}%
\pgfpathlineto{\pgfqpoint{5.756075in}{0.666255in}}%
\pgfpathlineto{\pgfqpoint{5.759430in}{0.712169in}}%
\pgfpathlineto{\pgfqpoint{5.762786in}{0.728459in}}%
\pgfpathlineto{\pgfqpoint{5.766142in}{0.730837in}}%
\pgfpathlineto{\pgfqpoint{5.769498in}{0.717829in}}%
\pgfpathlineto{\pgfqpoint{5.772854in}{0.688500in}}%
\pgfpathlineto{\pgfqpoint{5.776209in}{0.604559in}}%
\pgfpathlineto{\pgfqpoint{5.779565in}{0.684302in}}%
\pgfpathlineto{\pgfqpoint{5.782921in}{0.696171in}}%
\pgfpathlineto{\pgfqpoint{5.786277in}{0.690674in}}%
\pgfpathlineto{\pgfqpoint{5.789632in}{0.657735in}}%
\pgfpathlineto{\pgfqpoint{5.792988in}{0.661807in}}%
\pgfpathlineto{\pgfqpoint{5.796344in}{0.703557in}}%
\pgfpathlineto{\pgfqpoint{5.799700in}{0.718769in}}%
\pgfpathlineto{\pgfqpoint{5.803056in}{0.719915in}}%
\pgfpathlineto{\pgfqpoint{5.806411in}{0.704777in}}%
\pgfpathlineto{\pgfqpoint{5.809767in}{0.665255in}}%
\pgfpathlineto{\pgfqpoint{5.813123in}{0.655143in}}%
\pgfpathlineto{\pgfqpoint{5.816479in}{0.691798in}}%
\pgfpathlineto{\pgfqpoint{5.819834in}{0.700562in}}%
\pgfpathlineto{\pgfqpoint{5.823190in}{0.696117in}}%
\pgfpathlineto{\pgfqpoint{5.826546in}{0.672918in}}%
\pgfpathlineto{\pgfqpoint{5.829902in}{0.629561in}}%
\pgfpathlineto{\pgfqpoint{5.833258in}{0.694889in}}%
\pgfpathlineto{\pgfqpoint{5.836613in}{0.714603in}}%
\pgfpathlineto{\pgfqpoint{5.839969in}{0.720081in}}%
\pgfpathlineto{\pgfqpoint{5.843325in}{0.710173in}}%
\pgfpathlineto{\pgfqpoint{5.846681in}{0.676927in}}%
\pgfpathlineto{\pgfqpoint{5.850036in}{0.654819in}}%
\pgfpathlineto{\pgfqpoint{5.853392in}{0.699922in}}%
\pgfpathlineto{\pgfqpoint{5.856748in}{0.712933in}}%
\pgfpathlineto{\pgfqpoint{5.860104in}{0.714257in}}%
\pgfpathlineto{\pgfqpoint{5.863460in}{0.704748in}}%
\pgfpathlineto{\pgfqpoint{5.866815in}{0.672291in}}%
\pgfpathlineto{\pgfqpoint{5.870171in}{0.665420in}}%
\pgfpathlineto{\pgfqpoint{5.873527in}{0.710140in}}%
\pgfpathlineto{\pgfqpoint{5.876883in}{0.726374in}}%
\pgfpathlineto{\pgfqpoint{5.880238in}{0.729952in}}%
\pgfpathlineto{\pgfqpoint{5.883594in}{0.718535in}}%
\pgfpathlineto{\pgfqpoint{5.886950in}{0.682035in}}%
\pgfpathlineto{\pgfqpoint{5.890306in}{0.671219in}}%
\pgfpathlineto{\pgfqpoint{5.893662in}{0.711102in}}%
\pgfpathlineto{\pgfqpoint{5.897017in}{0.722957in}}%
\pgfpathlineto{\pgfqpoint{5.900373in}{0.723448in}}%
\pgfpathlineto{\pgfqpoint{5.903729in}{0.712750in}}%
\pgfpathlineto{\pgfqpoint{5.907085in}{0.675716in}}%
\pgfpathlineto{\pgfqpoint{5.910441in}{0.682292in}}%
\pgfpathlineto{\pgfqpoint{5.913796in}{0.721923in}}%
\pgfpathlineto{\pgfqpoint{5.917152in}{0.736822in}}%
\pgfpathlineto{\pgfqpoint{5.920508in}{0.739230in}}%
\pgfpathlineto{\pgfqpoint{5.923864in}{0.726889in}}%
\pgfpathlineto{\pgfqpoint{5.927219in}{0.695201in}}%
\pgfpathlineto{\pgfqpoint{5.930575in}{0.650867in}}%
\pgfpathlineto{\pgfqpoint{5.933931in}{0.704824in}}%
\pgfpathlineto{\pgfqpoint{5.937287in}{0.716655in}}%
\pgfpathlineto{\pgfqpoint{5.940643in}{0.714309in}}%
\pgfpathlineto{\pgfqpoint{5.943998in}{0.695024in}}%
\pgfpathlineto{\pgfqpoint{5.947354in}{0.610952in}}%
\pgfpathlineto{\pgfqpoint{5.950710in}{0.707754in}}%
\pgfpathlineto{\pgfqpoint{5.954066in}{0.730081in}}%
\pgfpathlineto{\pgfqpoint{5.957421in}{0.737239in}}%
\pgfpathlineto{\pgfqpoint{5.960777in}{0.729799in}}%
\pgfpathlineto{\pgfqpoint{5.964133in}{0.708542in}}%
\pgfpathlineto{\pgfqpoint{5.967489in}{0.652079in}}%
\pgfpathlineto{\pgfqpoint{5.974200in}{0.696972in}}%
\pgfpathlineto{\pgfqpoint{5.977556in}{0.697513in}}%
\pgfpathlineto{\pgfqpoint{5.980912in}{0.677596in}}%
\pgfpathlineto{\pgfqpoint{5.984268in}{0.627216in}}%
\pgfpathlineto{\pgfqpoint{5.987623in}{0.698961in}}%
\pgfpathlineto{\pgfqpoint{5.990979in}{0.719629in}}%
\pgfpathlineto{\pgfqpoint{5.994335in}{0.725277in}}%
\pgfpathlineto{\pgfqpoint{5.997691in}{0.715306in}}%
\pgfpathlineto{\pgfqpoint{6.001047in}{0.689517in}}%
\pgfpathlineto{\pgfqpoint{6.004402in}{0.537378in}}%
\pgfpathlineto{\pgfqpoint{6.007758in}{0.678375in}}%
\pgfpathlineto{\pgfqpoint{6.011114in}{0.694118in}}%
\pgfpathlineto{\pgfqpoint{6.014470in}{0.692661in}}%
\pgfpathlineto{\pgfqpoint{6.017825in}{0.672314in}}%
\pgfpathlineto{\pgfqpoint{6.021181in}{0.615337in}}%
\pgfpathlineto{\pgfqpoint{6.024537in}{0.691089in}}%
\pgfpathlineto{\pgfqpoint{6.027893in}{0.712025in}}%
\pgfpathlineto{\pgfqpoint{6.031249in}{0.718040in}}%
\pgfpathlineto{\pgfqpoint{6.034604in}{0.708412in}}%
\pgfpathlineto{\pgfqpoint{6.037960in}{0.677543in}}%
\pgfpathlineto{\pgfqpoint{6.041316in}{0.640408in}}%
\pgfpathlineto{\pgfqpoint{6.044672in}{0.692669in}}%
\pgfpathlineto{\pgfqpoint{6.048027in}{0.706229in}}%
\pgfpathlineto{\pgfqpoint{6.051383in}{0.707217in}}%
\pgfpathlineto{\pgfqpoint{6.054739in}{0.696386in}}%
\pgfpathlineto{\pgfqpoint{6.058095in}{0.656990in}}%
\pgfpathlineto{\pgfqpoint{6.061451in}{0.670255in}}%
\pgfpathlineto{\pgfqpoint{6.064806in}{0.707849in}}%
\pgfpathlineto{\pgfqpoint{6.068162in}{0.722043in}}%
\pgfpathlineto{\pgfqpoint{6.071518in}{0.723507in}}%
\pgfpathlineto{\pgfqpoint{6.074874in}{0.708490in}}%
\pgfpathlineto{\pgfqpoint{6.078230in}{0.655609in}}%
\pgfpathlineto{\pgfqpoint{6.084941in}{0.713158in}}%
\pgfpathlineto{\pgfqpoint{6.088297in}{0.722400in}}%
\pgfpathlineto{\pgfqpoint{6.091653in}{0.721526in}}%
\pgfpathlineto{\pgfqpoint{6.095008in}{0.709580in}}%
\pgfpathlineto{\pgfqpoint{6.098364in}{0.668756in}}%
\pgfpathlineto{\pgfqpoint{6.101720in}{0.683828in}}%
\pgfpathlineto{\pgfqpoint{6.105076in}{0.720730in}}%
\pgfpathlineto{\pgfqpoint{6.108432in}{0.734922in}}%
\pgfpathlineto{\pgfqpoint{6.111787in}{0.736973in}}%
\pgfpathlineto{\pgfqpoint{6.115143in}{0.724067in}}%
\pgfpathlineto{\pgfqpoint{6.118499in}{0.687228in}}%
\pgfpathlineto{\pgfqpoint{6.121855in}{0.673016in}}%
\pgfpathlineto{\pgfqpoint{6.125210in}{0.712961in}}%
\pgfpathlineto{\pgfqpoint{6.128566in}{0.723789in}}%
\pgfpathlineto{\pgfqpoint{6.131922in}{0.722444in}}%
\pgfpathlineto{\pgfqpoint{6.135278in}{0.707516in}}%
\pgfpathlineto{\pgfqpoint{6.138634in}{0.638879in}}%
\pgfpathlineto{\pgfqpoint{6.141989in}{0.702180in}}%
\pgfpathlineto{\pgfqpoint{6.145345in}{0.730369in}}%
\pgfpathlineto{\pgfqpoint{6.148701in}{0.740944in}}%
\pgfpathlineto{\pgfqpoint{6.152057in}{0.738249in}}%
\pgfpathlineto{\pgfqpoint{6.155412in}{0.721197in}}%
\pgfpathlineto{\pgfqpoint{6.158768in}{0.681641in}}%
\pgfpathlineto{\pgfqpoint{6.162124in}{0.663830in}}%
\pgfpathlineto{\pgfqpoint{6.165480in}{0.700758in}}%
\pgfpathlineto{\pgfqpoint{6.168836in}{0.706796in}}%
\pgfpathlineto{\pgfqpoint{6.172191in}{0.696000in}}%
\pgfpathlineto{\pgfqpoint{6.175547in}{0.638307in}}%
\pgfpathlineto{\pgfqpoint{6.178903in}{0.689530in}}%
\pgfpathlineto{\pgfqpoint{6.182259in}{0.719746in}}%
\pgfpathlineto{\pgfqpoint{6.185614in}{0.730829in}}%
\pgfpathlineto{\pgfqpoint{6.188970in}{0.727389in}}%
\pgfpathlineto{\pgfqpoint{6.192326in}{0.709697in}}%
\pgfpathlineto{\pgfqpoint{6.195682in}{0.669667in}}%
\pgfpathlineto{\pgfqpoint{6.199038in}{0.650857in}}%
\pgfpathlineto{\pgfqpoint{6.202393in}{0.687138in}}%
\pgfpathlineto{\pgfqpoint{6.205749in}{0.691870in}}%
\pgfpathlineto{\pgfqpoint{6.209105in}{0.677772in}}%
\pgfpathlineto{\pgfqpoint{6.212461in}{0.544612in}}%
\pgfpathlineto{\pgfqpoint{6.215816in}{0.685242in}}%
\pgfpathlineto{\pgfqpoint{6.219172in}{0.710394in}}%
\pgfpathlineto{\pgfqpoint{6.222528in}{0.718882in}}%
\pgfpathlineto{\pgfqpoint{6.225884in}{0.712062in}}%
\pgfpathlineto{\pgfqpoint{6.229240in}{0.688114in}}%
\pgfpathlineto{\pgfqpoint{6.232595in}{0.550528in}}%
\pgfpathlineto{\pgfqpoint{6.235951in}{0.679253in}}%
\pgfpathlineto{\pgfqpoint{6.239307in}{0.697305in}}%
\pgfpathlineto{\pgfqpoint{6.242663in}{0.699453in}}%
\pgfpathlineto{\pgfqpoint{6.246019in}{0.688276in}}%
\pgfpathlineto{\pgfqpoint{6.249374in}{0.642464in}}%
\pgfpathlineto{\pgfqpoint{6.252730in}{0.670866in}}%
\pgfpathlineto{\pgfqpoint{6.256086in}{0.704707in}}%
\pgfpathlineto{\pgfqpoint{6.259442in}{0.717456in}}%
\pgfpathlineto{\pgfqpoint{6.262797in}{0.716856in}}%
\pgfpathlineto{\pgfqpoint{6.266153in}{0.698355in}}%
\pgfpathlineto{\pgfqpoint{6.269509in}{0.619014in}}%
\pgfpathlineto{\pgfqpoint{6.272865in}{0.686935in}}%
\pgfpathlineto{\pgfqpoint{6.276221in}{0.710353in}}%
\pgfpathlineto{\pgfqpoint{6.279576in}{0.717647in}}%
\pgfpathlineto{\pgfqpoint{6.282932in}{0.715161in}}%
\pgfpathlineto{\pgfqpoint{6.286288in}{0.700651in}}%
\pgfpathlineto{\pgfqpoint{6.289644in}{0.644942in}}%
\pgfpathlineto{\pgfqpoint{6.292999in}{0.688040in}}%
\pgfpathlineto{\pgfqpoint{6.296355in}{0.718771in}}%
\pgfpathlineto{\pgfqpoint{6.299711in}{0.730795in}}%
\pgfpathlineto{\pgfqpoint{6.303067in}{0.730704in}}%
\pgfpathlineto{\pgfqpoint{6.306423in}{0.714235in}}%
\pgfpathlineto{\pgfqpoint{6.309778in}{0.656234in}}%
\pgfpathlineto{\pgfqpoint{6.313134in}{0.692935in}}%
\pgfpathlineto{\pgfqpoint{6.316490in}{0.719719in}}%
\pgfpathlineto{\pgfqpoint{6.319846in}{0.727999in}}%
\pgfpathlineto{\pgfqpoint{6.323201in}{0.726009in}}%
\pgfpathlineto{\pgfqpoint{6.326557in}{0.711866in}}%
\pgfpathlineto{\pgfqpoint{6.329913in}{0.657299in}}%
\pgfpathlineto{\pgfqpoint{6.333269in}{0.698828in}}%
\pgfpathlineto{\pgfqpoint{6.336625in}{0.729833in}}%
\pgfpathlineto{\pgfqpoint{6.339980in}{0.741881in}}%
\pgfpathlineto{\pgfqpoint{6.343336in}{0.741612in}}%
\pgfpathlineto{\pgfqpoint{6.346692in}{0.726494in}}%
\pgfpathlineto{\pgfqpoint{6.350048in}{0.689398in}}%
\pgfpathlineto{\pgfqpoint{6.353403in}{0.668031in}}%
\pgfpathlineto{\pgfqpoint{6.356759in}{0.708576in}}%
\pgfpathlineto{\pgfqpoint{6.360115in}{0.717342in}}%
\pgfpathlineto{\pgfqpoint{6.363471in}{0.711769in}}%
\pgfpathlineto{\pgfqpoint{6.366827in}{0.683713in}}%
\pgfpathlineto{\pgfqpoint{6.370182in}{0.668115in}}%
\pgfpathlineto{\pgfqpoint{6.373538in}{0.717691in}}%
\pgfpathlineto{\pgfqpoint{6.376894in}{0.734804in}}%
\pgfpathlineto{\pgfqpoint{6.380250in}{0.738055in}}%
\pgfpathlineto{\pgfqpoint{6.383605in}{0.728425in}}%
\pgfpathlineto{\pgfqpoint{6.383605in}{0.728425in}}%
\pgfusepath{stroke}%
\end{pgfscope}%
\begin{pgfscope}%
\pgfpathrectangle{\pgfqpoint{4.621821in}{0.320625in}}{\pgfqpoint{1.845679in}{1.164375in}}%
\pgfusepath{clip}%
\pgfsetrectcap%
\pgfsetroundjoin%
\pgfsetlinewidth{1.154312pt}%
\definecolor{currentstroke}{rgb}{0.000000,0.529412,0.423529}%
\pgfsetstrokecolor{currentstroke}%
\pgfsetdash{}{0pt}%
\pgfpathmoveto{\pgfqpoint{4.705715in}{1.469210in}}%
\pgfpathlineto{\pgfqpoint{4.719139in}{1.444446in}}%
\pgfpathlineto{\pgfqpoint{4.729206in}{1.423292in}}%
\pgfpathlineto{\pgfqpoint{4.745985in}{1.384736in}}%
\pgfpathlineto{\pgfqpoint{4.772831in}{1.323071in}}%
\pgfpathlineto{\pgfqpoint{4.792966in}{1.279649in}}%
\pgfpathlineto{\pgfqpoint{4.816456in}{1.231891in}}%
\pgfpathlineto{\pgfqpoint{4.843302in}{1.180129in}}%
\pgfpathlineto{\pgfqpoint{4.873504in}{1.124407in}}%
\pgfpathlineto{\pgfqpoint{4.910418in}{1.058822in}}%
\pgfpathlineto{\pgfqpoint{4.954043in}{0.983823in}}%
\pgfpathlineto{\pgfqpoint{5.004380in}{0.899644in}}%
\pgfpathlineto{\pgfqpoint{5.064784in}{0.800999in}}%
\pgfpathlineto{\pgfqpoint{5.111765in}{0.726902in}}%
\pgfpathlineto{\pgfqpoint{5.121832in}{0.712699in}}%
\pgfpathlineto{\pgfqpoint{5.128544in}{0.700793in}}%
\pgfpathlineto{\pgfqpoint{5.141967in}{0.674829in}}%
\pgfpathlineto{\pgfqpoint{5.145323in}{0.671023in}}%
\pgfpathlineto{\pgfqpoint{5.148678in}{0.669038in}}%
\pgfpathlineto{\pgfqpoint{5.155390in}{0.666309in}}%
\pgfpathlineto{\pgfqpoint{5.158746in}{0.662719in}}%
\pgfpathlineto{\pgfqpoint{5.162102in}{0.658265in}}%
\pgfpathlineto{\pgfqpoint{5.165457in}{0.656023in}}%
\pgfpathlineto{\pgfqpoint{5.168813in}{0.658268in}}%
\pgfpathlineto{\pgfqpoint{5.175525in}{0.668109in}}%
\pgfpathlineto{\pgfqpoint{5.182236in}{0.674654in}}%
\pgfpathlineto{\pgfqpoint{5.185592in}{0.680521in}}%
\pgfpathlineto{\pgfqpoint{5.195659in}{0.704713in}}%
\pgfpathlineto{\pgfqpoint{5.199015in}{0.710719in}}%
\pgfpathlineto{\pgfqpoint{5.202371in}{0.715070in}}%
\pgfpathlineto{\pgfqpoint{5.205727in}{0.717810in}}%
\pgfpathlineto{\pgfqpoint{5.209082in}{0.718984in}}%
\pgfpathlineto{\pgfqpoint{5.212438in}{0.718590in}}%
\pgfpathlineto{\pgfqpoint{5.215794in}{0.716553in}}%
\pgfpathlineto{\pgfqpoint{5.219150in}{0.712718in}}%
\pgfpathlineto{\pgfqpoint{5.222506in}{0.706894in}}%
\pgfpathlineto{\pgfqpoint{5.229217in}{0.691593in}}%
\pgfpathlineto{\pgfqpoint{5.232573in}{0.691279in}}%
\pgfpathlineto{\pgfqpoint{5.242640in}{0.722420in}}%
\pgfpathlineto{\pgfqpoint{5.245996in}{0.728442in}}%
\pgfpathlineto{\pgfqpoint{5.249352in}{0.731463in}}%
\pgfpathlineto{\pgfqpoint{5.252708in}{0.731398in}}%
\pgfpathlineto{\pgfqpoint{5.256063in}{0.727852in}}%
\pgfpathlineto{\pgfqpoint{5.259419in}{0.721853in}}%
\pgfpathlineto{\pgfqpoint{5.266131in}{0.704928in}}%
\pgfpathlineto{\pgfqpoint{5.269487in}{0.700761in}}%
\pgfpathlineto{\pgfqpoint{5.272842in}{0.700776in}}%
\pgfpathlineto{\pgfqpoint{5.276198in}{0.699311in}}%
\pgfpathlineto{\pgfqpoint{5.279554in}{0.696363in}}%
\pgfpathlineto{\pgfqpoint{5.282910in}{0.695960in}}%
\pgfpathlineto{\pgfqpoint{5.286265in}{0.699623in}}%
\pgfpathlineto{\pgfqpoint{5.289621in}{0.704129in}}%
\pgfpathlineto{\pgfqpoint{5.292977in}{0.707673in}}%
\pgfpathlineto{\pgfqpoint{5.296333in}{0.707671in}}%
\pgfpathlineto{\pgfqpoint{5.299689in}{0.700236in}}%
\pgfpathlineto{\pgfqpoint{5.303044in}{0.682634in}}%
\pgfpathlineto{\pgfqpoint{5.306400in}{0.677001in}}%
\pgfpathlineto{\pgfqpoint{5.309756in}{0.699736in}}%
\pgfpathlineto{\pgfqpoint{5.313112in}{0.713858in}}%
\pgfpathlineto{\pgfqpoint{5.316467in}{0.718362in}}%
\pgfpathlineto{\pgfqpoint{5.319823in}{0.713463in}}%
\pgfpathlineto{\pgfqpoint{5.323179in}{0.700556in}}%
\pgfpathlineto{\pgfqpoint{5.326535in}{0.677562in}}%
\pgfpathlineto{\pgfqpoint{5.329891in}{0.675441in}}%
\pgfpathlineto{\pgfqpoint{5.333246in}{0.694293in}}%
\pgfpathlineto{\pgfqpoint{5.336602in}{0.704373in}}%
\pgfpathlineto{\pgfqpoint{5.339958in}{0.708202in}}%
\pgfpathlineto{\pgfqpoint{5.343314in}{0.707763in}}%
\pgfpathlineto{\pgfqpoint{5.346669in}{0.706208in}}%
\pgfpathlineto{\pgfqpoint{5.353381in}{0.711981in}}%
\pgfpathlineto{\pgfqpoint{5.360093in}{0.709615in}}%
\pgfpathlineto{\pgfqpoint{5.363448in}{0.705292in}}%
\pgfpathlineto{\pgfqpoint{5.366804in}{0.697905in}}%
\pgfpathlineto{\pgfqpoint{5.370160in}{0.695930in}}%
\pgfpathlineto{\pgfqpoint{5.376871in}{0.713767in}}%
\pgfpathlineto{\pgfqpoint{5.380227in}{0.717934in}}%
\pgfpathlineto{\pgfqpoint{5.383583in}{0.714658in}}%
\pgfpathlineto{\pgfqpoint{5.390295in}{0.699567in}}%
\pgfpathlineto{\pgfqpoint{5.393650in}{0.697107in}}%
\pgfpathlineto{\pgfqpoint{5.397006in}{0.707795in}}%
\pgfpathlineto{\pgfqpoint{5.400362in}{0.716443in}}%
\pgfpathlineto{\pgfqpoint{5.403718in}{0.717201in}}%
\pgfpathlineto{\pgfqpoint{5.407074in}{0.711119in}}%
\pgfpathlineto{\pgfqpoint{5.410429in}{0.703857in}}%
\pgfpathlineto{\pgfqpoint{5.413785in}{0.704755in}}%
\pgfpathlineto{\pgfqpoint{5.417141in}{0.712942in}}%
\pgfpathlineto{\pgfqpoint{5.420497in}{0.716598in}}%
\pgfpathlineto{\pgfqpoint{5.423852in}{0.715154in}}%
\pgfpathlineto{\pgfqpoint{5.427208in}{0.712385in}}%
\pgfpathlineto{\pgfqpoint{5.430564in}{0.706681in}}%
\pgfpathlineto{\pgfqpoint{5.433920in}{0.705502in}}%
\pgfpathlineto{\pgfqpoint{5.437276in}{0.714237in}}%
\pgfpathlineto{\pgfqpoint{5.440631in}{0.721061in}}%
\pgfpathlineto{\pgfqpoint{5.443987in}{0.721133in}}%
\pgfpathlineto{\pgfqpoint{5.447343in}{0.718903in}}%
\pgfpathlineto{\pgfqpoint{5.454054in}{0.712443in}}%
\pgfpathlineto{\pgfqpoint{5.460766in}{0.716414in}}%
\pgfpathlineto{\pgfqpoint{5.467478in}{0.723981in}}%
\pgfpathlineto{\pgfqpoint{5.470833in}{0.721888in}}%
\pgfpathlineto{\pgfqpoint{5.474189in}{0.717837in}}%
\pgfpathlineto{\pgfqpoint{5.477545in}{0.715929in}}%
\pgfpathlineto{\pgfqpoint{5.480901in}{0.715073in}}%
\pgfpathlineto{\pgfqpoint{5.487612in}{0.720394in}}%
\pgfpathlineto{\pgfqpoint{5.494324in}{0.720987in}}%
\pgfpathlineto{\pgfqpoint{5.497680in}{0.724272in}}%
\pgfpathlineto{\pgfqpoint{5.501035in}{0.725756in}}%
\pgfpathlineto{\pgfqpoint{5.504391in}{0.726249in}}%
\pgfpathlineto{\pgfqpoint{5.507747in}{0.723344in}}%
\pgfpathlineto{\pgfqpoint{5.511103in}{0.716944in}}%
\pgfpathlineto{\pgfqpoint{5.514458in}{0.713790in}}%
\pgfpathlineto{\pgfqpoint{5.517814in}{0.712169in}}%
\pgfpathlineto{\pgfqpoint{5.521170in}{0.711725in}}%
\pgfpathlineto{\pgfqpoint{5.524526in}{0.712439in}}%
\pgfpathlineto{\pgfqpoint{5.527882in}{0.714328in}}%
\pgfpathlineto{\pgfqpoint{5.531237in}{0.717454in}}%
\pgfpathlineto{\pgfqpoint{5.534593in}{0.725279in}}%
\pgfpathlineto{\pgfqpoint{5.537949in}{0.730099in}}%
\pgfpathlineto{\pgfqpoint{5.541305in}{0.727061in}}%
\pgfpathlineto{\pgfqpoint{5.544660in}{0.714323in}}%
\pgfpathlineto{\pgfqpoint{5.548016in}{0.698822in}}%
\pgfpathlineto{\pgfqpoint{5.551372in}{0.698633in}}%
\pgfpathlineto{\pgfqpoint{5.554728in}{0.703956in}}%
\pgfpathlineto{\pgfqpoint{5.558084in}{0.704594in}}%
\pgfpathlineto{\pgfqpoint{5.561439in}{0.703421in}}%
\pgfpathlineto{\pgfqpoint{5.564795in}{0.704387in}}%
\pgfpathlineto{\pgfqpoint{5.568151in}{0.709370in}}%
\pgfpathlineto{\pgfqpoint{5.571507in}{0.718982in}}%
\pgfpathlineto{\pgfqpoint{5.574863in}{0.724536in}}%
\pgfpathlineto{\pgfqpoint{5.578218in}{0.721482in}}%
\pgfpathlineto{\pgfqpoint{5.584930in}{0.688838in}}%
\pgfpathlineto{\pgfqpoint{5.588286in}{0.698086in}}%
\pgfpathlineto{\pgfqpoint{5.591641in}{0.709887in}}%
\pgfpathlineto{\pgfqpoint{5.594997in}{0.711767in}}%
\pgfpathlineto{\pgfqpoint{5.601709in}{0.694861in}}%
\pgfpathlineto{\pgfqpoint{5.605065in}{0.699716in}}%
\pgfpathlineto{\pgfqpoint{5.608420in}{0.712999in}}%
\pgfpathlineto{\pgfqpoint{5.611776in}{0.720754in}}%
\pgfpathlineto{\pgfqpoint{5.615132in}{0.719504in}}%
\pgfpathlineto{\pgfqpoint{5.618488in}{0.706906in}}%
\pgfpathlineto{\pgfqpoint{5.621843in}{0.686433in}}%
\pgfpathlineto{\pgfqpoint{5.625199in}{0.695383in}}%
\pgfpathlineto{\pgfqpoint{5.628555in}{0.713720in}}%
\pgfpathlineto{\pgfqpoint{5.631911in}{0.722310in}}%
\pgfpathlineto{\pgfqpoint{5.635267in}{0.721560in}}%
\pgfpathlineto{\pgfqpoint{5.641978in}{0.703575in}}%
\pgfpathlineto{\pgfqpoint{5.645334in}{0.706756in}}%
\pgfpathlineto{\pgfqpoint{5.648690in}{0.716969in}}%
\pgfpathlineto{\pgfqpoint{5.652045in}{0.722325in}}%
\pgfpathlineto{\pgfqpoint{5.655401in}{0.720237in}}%
\pgfpathlineto{\pgfqpoint{5.658757in}{0.709880in}}%
\pgfpathlineto{\pgfqpoint{5.662113in}{0.702630in}}%
\pgfpathlineto{\pgfqpoint{5.668824in}{0.722308in}}%
\pgfpathlineto{\pgfqpoint{5.672180in}{0.724322in}}%
\pgfpathlineto{\pgfqpoint{5.678892in}{0.712669in}}%
\pgfpathlineto{\pgfqpoint{5.682247in}{0.713571in}}%
\pgfpathlineto{\pgfqpoint{5.688959in}{0.728212in}}%
\pgfpathlineto{\pgfqpoint{5.692315in}{0.731287in}}%
\pgfpathlineto{\pgfqpoint{5.695671in}{0.726976in}}%
\pgfpathlineto{\pgfqpoint{5.699026in}{0.714893in}}%
\pgfpathlineto{\pgfqpoint{5.702382in}{0.710305in}}%
\pgfpathlineto{\pgfqpoint{5.705738in}{0.716536in}}%
\pgfpathlineto{\pgfqpoint{5.709094in}{0.719321in}}%
\pgfpathlineto{\pgfqpoint{5.712449in}{0.715828in}}%
\pgfpathlineto{\pgfqpoint{5.715805in}{0.708963in}}%
\pgfpathlineto{\pgfqpoint{5.719161in}{0.709741in}}%
\pgfpathlineto{\pgfqpoint{5.722517in}{0.720764in}}%
\pgfpathlineto{\pgfqpoint{5.725873in}{0.729776in}}%
\pgfpathlineto{\pgfqpoint{5.729228in}{0.734497in}}%
\pgfpathlineto{\pgfqpoint{5.732584in}{0.731369in}}%
\pgfpathlineto{\pgfqpoint{5.735940in}{0.717704in}}%
\pgfpathlineto{\pgfqpoint{5.739296in}{0.699323in}}%
\pgfpathlineto{\pgfqpoint{5.742652in}{0.698770in}}%
\pgfpathlineto{\pgfqpoint{5.746007in}{0.706658in}}%
\pgfpathlineto{\pgfqpoint{5.749363in}{0.707701in}}%
\pgfpathlineto{\pgfqpoint{5.752719in}{0.704616in}}%
\pgfpathlineto{\pgfqpoint{5.756075in}{0.705516in}}%
\pgfpathlineto{\pgfqpoint{5.762786in}{0.721738in}}%
\pgfpathlineto{\pgfqpoint{5.766142in}{0.725937in}}%
\pgfpathlineto{\pgfqpoint{5.769498in}{0.724875in}}%
\pgfpathlineto{\pgfqpoint{5.772854in}{0.718566in}}%
\pgfpathlineto{\pgfqpoint{5.776209in}{0.704935in}}%
\pgfpathlineto{\pgfqpoint{5.779565in}{0.698860in}}%
\pgfpathlineto{\pgfqpoint{5.782921in}{0.703355in}}%
\pgfpathlineto{\pgfqpoint{5.786277in}{0.703735in}}%
\pgfpathlineto{\pgfqpoint{5.789632in}{0.695420in}}%
\pgfpathlineto{\pgfqpoint{5.792988in}{0.692106in}}%
\pgfpathlineto{\pgfqpoint{5.799700in}{0.715472in}}%
\pgfpathlineto{\pgfqpoint{5.803056in}{0.720235in}}%
\pgfpathlineto{\pgfqpoint{5.806411in}{0.718711in}}%
\pgfpathlineto{\pgfqpoint{5.809767in}{0.710758in}}%
\pgfpathlineto{\pgfqpoint{5.813123in}{0.696946in}}%
\pgfpathlineto{\pgfqpoint{5.816479in}{0.697571in}}%
\pgfpathlineto{\pgfqpoint{5.819834in}{0.709436in}}%
\pgfpathlineto{\pgfqpoint{5.823190in}{0.716712in}}%
\pgfpathlineto{\pgfqpoint{5.826546in}{0.716346in}}%
\pgfpathlineto{\pgfqpoint{5.829902in}{0.706348in}}%
\pgfpathlineto{\pgfqpoint{5.833258in}{0.700392in}}%
\pgfpathlineto{\pgfqpoint{5.836613in}{0.710681in}}%
\pgfpathlineto{\pgfqpoint{5.839969in}{0.719176in}}%
\pgfpathlineto{\pgfqpoint{5.843325in}{0.719704in}}%
\pgfpathlineto{\pgfqpoint{5.846681in}{0.710543in}}%
\pgfpathlineto{\pgfqpoint{5.850036in}{0.696467in}}%
\pgfpathlineto{\pgfqpoint{5.853392in}{0.701668in}}%
\pgfpathlineto{\pgfqpoint{5.856748in}{0.715899in}}%
\pgfpathlineto{\pgfqpoint{5.860104in}{0.724117in}}%
\pgfpathlineto{\pgfqpoint{5.863460in}{0.725370in}}%
\pgfpathlineto{\pgfqpoint{5.866815in}{0.717430in}}%
\pgfpathlineto{\pgfqpoint{5.870171in}{0.703302in}}%
\pgfpathlineto{\pgfqpoint{5.873527in}{0.710278in}}%
\pgfpathlineto{\pgfqpoint{5.876883in}{0.725145in}}%
\pgfpathlineto{\pgfqpoint{5.880238in}{0.731679in}}%
\pgfpathlineto{\pgfqpoint{5.883594in}{0.728622in}}%
\pgfpathlineto{\pgfqpoint{5.886950in}{0.712570in}}%
\pgfpathlineto{\pgfqpoint{5.890306in}{0.691908in}}%
\pgfpathlineto{\pgfqpoint{5.897017in}{0.723102in}}%
\pgfpathlineto{\pgfqpoint{5.900373in}{0.727003in}}%
\pgfpathlineto{\pgfqpoint{5.903729in}{0.720431in}}%
\pgfpathlineto{\pgfqpoint{5.907085in}{0.702888in}}%
\pgfpathlineto{\pgfqpoint{5.910441in}{0.698423in}}%
\pgfpathlineto{\pgfqpoint{5.913796in}{0.722073in}}%
\pgfpathlineto{\pgfqpoint{5.917152in}{0.735728in}}%
\pgfpathlineto{\pgfqpoint{5.920508in}{0.738494in}}%
\pgfpathlineto{\pgfqpoint{5.923864in}{0.730762in}}%
\pgfpathlineto{\pgfqpoint{5.927219in}{0.714868in}}%
\pgfpathlineto{\pgfqpoint{5.930575in}{0.695119in}}%
\pgfpathlineto{\pgfqpoint{5.933931in}{0.702961in}}%
\pgfpathlineto{\pgfqpoint{5.937287in}{0.713729in}}%
\pgfpathlineto{\pgfqpoint{5.940643in}{0.714756in}}%
\pgfpathlineto{\pgfqpoint{5.943998in}{0.709181in}}%
\pgfpathlineto{\pgfqpoint{5.947354in}{0.706150in}}%
\pgfpathlineto{\pgfqpoint{5.954066in}{0.726766in}}%
\pgfpathlineto{\pgfqpoint{5.957421in}{0.729568in}}%
\pgfpathlineto{\pgfqpoint{5.964133in}{0.724494in}}%
\pgfpathlineto{\pgfqpoint{5.967489in}{0.717514in}}%
\pgfpathlineto{\pgfqpoint{5.970845in}{0.706375in}}%
\pgfpathlineto{\pgfqpoint{5.974200in}{0.700123in}}%
\pgfpathlineto{\pgfqpoint{5.977556in}{0.697674in}}%
\pgfpathlineto{\pgfqpoint{5.980912in}{0.694048in}}%
\pgfpathlineto{\pgfqpoint{5.984268in}{0.697963in}}%
\pgfpathlineto{\pgfqpoint{5.987623in}{0.708767in}}%
\pgfpathlineto{\pgfqpoint{5.990979in}{0.716489in}}%
\pgfpathlineto{\pgfqpoint{5.994335in}{0.717890in}}%
\pgfpathlineto{\pgfqpoint{5.997691in}{0.716974in}}%
\pgfpathlineto{\pgfqpoint{6.001047in}{0.717561in}}%
\pgfpathlineto{\pgfqpoint{6.004402in}{0.713391in}}%
\pgfpathlineto{\pgfqpoint{6.007758in}{0.704581in}}%
\pgfpathlineto{\pgfqpoint{6.011114in}{0.702916in}}%
\pgfpathlineto{\pgfqpoint{6.014470in}{0.707304in}}%
\pgfpathlineto{\pgfqpoint{6.017825in}{0.708990in}}%
\pgfpathlineto{\pgfqpoint{6.021181in}{0.703984in}}%
\pgfpathlineto{\pgfqpoint{6.024537in}{0.702293in}}%
\pgfpathlineto{\pgfqpoint{6.027893in}{0.709510in}}%
\pgfpathlineto{\pgfqpoint{6.031249in}{0.713972in}}%
\pgfpathlineto{\pgfqpoint{6.037960in}{0.714072in}}%
\pgfpathlineto{\pgfqpoint{6.044672in}{0.709608in}}%
\pgfpathlineto{\pgfqpoint{6.048027in}{0.710886in}}%
\pgfpathlineto{\pgfqpoint{6.054739in}{0.718267in}}%
\pgfpathlineto{\pgfqpoint{6.058095in}{0.718740in}}%
\pgfpathlineto{\pgfqpoint{6.061451in}{0.714958in}}%
\pgfpathlineto{\pgfqpoint{6.064806in}{0.718250in}}%
\pgfpathlineto{\pgfqpoint{6.068162in}{0.724689in}}%
\pgfpathlineto{\pgfqpoint{6.071518in}{0.724008in}}%
\pgfpathlineto{\pgfqpoint{6.074874in}{0.713078in}}%
\pgfpathlineto{\pgfqpoint{6.078230in}{0.697876in}}%
\pgfpathlineto{\pgfqpoint{6.081585in}{0.702608in}}%
\pgfpathlineto{\pgfqpoint{6.084941in}{0.717610in}}%
\pgfpathlineto{\pgfqpoint{6.088297in}{0.725810in}}%
\pgfpathlineto{\pgfqpoint{6.091653in}{0.727148in}}%
\pgfpathlineto{\pgfqpoint{6.095008in}{0.721068in}}%
\pgfpathlineto{\pgfqpoint{6.098364in}{0.705923in}}%
\pgfpathlineto{\pgfqpoint{6.101720in}{0.709191in}}%
\pgfpathlineto{\pgfqpoint{6.105076in}{0.727679in}}%
\pgfpathlineto{\pgfqpoint{6.108432in}{0.737383in}}%
\pgfpathlineto{\pgfqpoint{6.111787in}{0.736265in}}%
\pgfpathlineto{\pgfqpoint{6.115143in}{0.725492in}}%
\pgfpathlineto{\pgfqpoint{6.118499in}{0.706872in}}%
\pgfpathlineto{\pgfqpoint{6.121855in}{0.693309in}}%
\pgfpathlineto{\pgfqpoint{6.125210in}{0.709043in}}%
\pgfpathlineto{\pgfqpoint{6.128566in}{0.719624in}}%
\pgfpathlineto{\pgfqpoint{6.131922in}{0.720950in}}%
\pgfpathlineto{\pgfqpoint{6.135278in}{0.716053in}}%
\pgfpathlineto{\pgfqpoint{6.138634in}{0.713122in}}%
\pgfpathlineto{\pgfqpoint{6.145345in}{0.728263in}}%
\pgfpathlineto{\pgfqpoint{6.148701in}{0.731577in}}%
\pgfpathlineto{\pgfqpoint{6.152057in}{0.731332in}}%
\pgfpathlineto{\pgfqpoint{6.155412in}{0.727402in}}%
\pgfpathlineto{\pgfqpoint{6.158768in}{0.722613in}}%
\pgfpathlineto{\pgfqpoint{6.165480in}{0.700404in}}%
\pgfpathlineto{\pgfqpoint{6.168836in}{0.698823in}}%
\pgfpathlineto{\pgfqpoint{6.172191in}{0.703485in}}%
\pgfpathlineto{\pgfqpoint{6.175547in}{0.709256in}}%
\pgfpathlineto{\pgfqpoint{6.178903in}{0.713281in}}%
\pgfpathlineto{\pgfqpoint{6.182259in}{0.714965in}}%
\pgfpathlineto{\pgfqpoint{6.185614in}{0.717882in}}%
\pgfpathlineto{\pgfqpoint{6.192326in}{0.720032in}}%
\pgfpathlineto{\pgfqpoint{6.195682in}{0.722683in}}%
\pgfpathlineto{\pgfqpoint{6.199038in}{0.718345in}}%
\pgfpathlineto{\pgfqpoint{6.202393in}{0.707891in}}%
\pgfpathlineto{\pgfqpoint{6.205749in}{0.702142in}}%
\pgfpathlineto{\pgfqpoint{6.209105in}{0.702818in}}%
\pgfpathlineto{\pgfqpoint{6.212461in}{0.701470in}}%
\pgfpathlineto{\pgfqpoint{6.215816in}{0.699359in}}%
\pgfpathlineto{\pgfqpoint{6.219172in}{0.701409in}}%
\pgfpathlineto{\pgfqpoint{6.222528in}{0.709264in}}%
\pgfpathlineto{\pgfqpoint{6.225884in}{0.715087in}}%
\pgfpathlineto{\pgfqpoint{6.229240in}{0.718644in}}%
\pgfpathlineto{\pgfqpoint{6.232595in}{0.718916in}}%
\pgfpathlineto{\pgfqpoint{6.235951in}{0.711125in}}%
\pgfpathlineto{\pgfqpoint{6.239307in}{0.698731in}}%
\pgfpathlineto{\pgfqpoint{6.242663in}{0.698804in}}%
\pgfpathlineto{\pgfqpoint{6.246019in}{0.711875in}}%
\pgfpathlineto{\pgfqpoint{6.249374in}{0.720916in}}%
\pgfpathlineto{\pgfqpoint{6.252730in}{0.721958in}}%
\pgfpathlineto{\pgfqpoint{6.256086in}{0.716787in}}%
\pgfpathlineto{\pgfqpoint{6.259442in}{0.708947in}}%
\pgfpathlineto{\pgfqpoint{6.262797in}{0.706094in}}%
\pgfpathlineto{\pgfqpoint{6.266153in}{0.711091in}}%
\pgfpathlineto{\pgfqpoint{6.269509in}{0.717581in}}%
\pgfpathlineto{\pgfqpoint{6.272865in}{0.719858in}}%
\pgfpathlineto{\pgfqpoint{6.276221in}{0.718777in}}%
\pgfpathlineto{\pgfqpoint{6.279576in}{0.716034in}}%
\pgfpathlineto{\pgfqpoint{6.282932in}{0.715463in}}%
\pgfpathlineto{\pgfqpoint{6.292999in}{0.727643in}}%
\pgfpathlineto{\pgfqpoint{6.296355in}{0.730840in}}%
\pgfpathlineto{\pgfqpoint{6.299711in}{0.728806in}}%
\pgfpathlineto{\pgfqpoint{6.303067in}{0.723144in}}%
\pgfpathlineto{\pgfqpoint{6.306423in}{0.715878in}}%
\pgfpathlineto{\pgfqpoint{6.309778in}{0.710240in}}%
\pgfpathlineto{\pgfqpoint{6.313134in}{0.710355in}}%
\pgfpathlineto{\pgfqpoint{6.319846in}{0.720400in}}%
\pgfpathlineto{\pgfqpoint{6.323201in}{0.722645in}}%
\pgfpathlineto{\pgfqpoint{6.326557in}{0.722931in}}%
\pgfpathlineto{\pgfqpoint{6.329913in}{0.725310in}}%
\pgfpathlineto{\pgfqpoint{6.333269in}{0.728912in}}%
\pgfpathlineto{\pgfqpoint{6.339980in}{0.728122in}}%
\pgfpathlineto{\pgfqpoint{6.343336in}{0.730488in}}%
\pgfpathlineto{\pgfqpoint{6.346692in}{0.729662in}}%
\pgfpathlineto{\pgfqpoint{6.350048in}{0.724468in}}%
\pgfpathlineto{\pgfqpoint{6.356759in}{0.704827in}}%
\pgfpathlineto{\pgfqpoint{6.360115in}{0.703465in}}%
\pgfpathlineto{\pgfqpoint{6.366827in}{0.706921in}}%
\pgfpathlineto{\pgfqpoint{6.370182in}{0.703746in}}%
\pgfpathlineto{\pgfqpoint{6.373538in}{0.704573in}}%
\pgfpathlineto{\pgfqpoint{6.376894in}{0.715270in}}%
\pgfpathlineto{\pgfqpoint{6.380250in}{0.720194in}}%
\pgfpathlineto{\pgfqpoint{6.383605in}{0.714345in}}%
\pgfpathlineto{\pgfqpoint{6.383605in}{0.714345in}}%
\pgfusepath{stroke}%
\end{pgfscope}%
\begin{pgfscope}%
\pgfpathrectangle{\pgfqpoint{4.621821in}{0.320625in}}{\pgfqpoint{1.845679in}{1.164375in}}%
\pgfusepath{clip}%
\pgfsetrectcap%
\pgfsetroundjoin%
\pgfsetlinewidth{0.602250pt}%
\definecolor{currentstroke}{rgb}{0.478431,0.243137,0.615686}%
\pgfsetstrokecolor{currentstroke}%
\pgfsetstrokeopacity{0.120000}%
\pgfsetdash{}{0pt}%
\pgfpathmoveto{\pgfqpoint{4.705715in}{1.469034in}}%
\pgfpathlineto{\pgfqpoint{4.715783in}{1.450518in}}%
\pgfpathlineto{\pgfqpoint{4.725850in}{1.429517in}}%
\pgfpathlineto{\pgfqpoint{4.739273in}{1.398106in}}%
\pgfpathlineto{\pgfqpoint{4.779543in}{1.301177in}}%
\pgfpathlineto{\pgfqpoint{4.799677in}{1.255951in}}%
\pgfpathlineto{\pgfqpoint{4.823168in}{1.205836in}}%
\pgfpathlineto{\pgfqpoint{4.853370in}{1.144320in}}%
\pgfpathlineto{\pgfqpoint{4.890283in}{1.071929in}}%
\pgfpathlineto{\pgfqpoint{4.937264in}{0.982479in}}%
\pgfpathlineto{\pgfqpoint{4.994313in}{0.876411in}}%
\pgfpathlineto{\pgfqpoint{5.041293in}{0.790300in}}%
\pgfpathlineto{\pgfqpoint{5.051361in}{0.771953in}}%
\pgfpathlineto{\pgfqpoint{5.068140in}{0.744699in}}%
\pgfpathlineto{\pgfqpoint{5.078207in}{0.722462in}}%
\pgfpathlineto{\pgfqpoint{5.084919in}{0.707580in}}%
\pgfpathlineto{\pgfqpoint{5.088274in}{0.701455in}}%
\pgfpathlineto{\pgfqpoint{5.091630in}{0.696886in}}%
\pgfpathlineto{\pgfqpoint{5.094986in}{0.694117in}}%
\pgfpathlineto{\pgfqpoint{5.098342in}{0.692917in}}%
\pgfpathlineto{\pgfqpoint{5.105053in}{0.691648in}}%
\pgfpathlineto{\pgfqpoint{5.108409in}{0.688670in}}%
\pgfpathlineto{\pgfqpoint{5.111765in}{0.682029in}}%
\pgfpathlineto{\pgfqpoint{5.115121in}{0.672642in}}%
\pgfpathlineto{\pgfqpoint{5.118476in}{0.660278in}}%
\pgfpathlineto{\pgfqpoint{5.121832in}{0.645420in}}%
\pgfpathlineto{\pgfqpoint{5.125188in}{0.633731in}}%
\pgfpathlineto{\pgfqpoint{5.128544in}{0.634356in}}%
\pgfpathlineto{\pgfqpoint{5.131900in}{0.639967in}}%
\pgfpathlineto{\pgfqpoint{5.135255in}{0.642681in}}%
\pgfpathlineto{\pgfqpoint{5.138611in}{0.640612in}}%
\pgfpathlineto{\pgfqpoint{5.141967in}{0.632254in}}%
\pgfpathlineto{\pgfqpoint{5.145323in}{0.613028in}}%
\pgfpathlineto{\pgfqpoint{5.148678in}{0.585166in}}%
\pgfpathlineto{\pgfqpoint{5.152034in}{0.621932in}}%
\pgfpathlineto{\pgfqpoint{5.155390in}{0.645401in}}%
\pgfpathlineto{\pgfqpoint{5.158746in}{0.659026in}}%
\pgfpathlineto{\pgfqpoint{5.162102in}{0.666890in}}%
\pgfpathlineto{\pgfqpoint{5.165457in}{0.670008in}}%
\pgfpathlineto{\pgfqpoint{5.168813in}{0.667684in}}%
\pgfpathlineto{\pgfqpoint{5.172169in}{0.657265in}}%
\pgfpathlineto{\pgfqpoint{5.175525in}{0.661053in}}%
\pgfpathlineto{\pgfqpoint{5.178880in}{0.685393in}}%
\pgfpathlineto{\pgfqpoint{5.182236in}{0.703130in}}%
\pgfpathlineto{\pgfqpoint{5.185592in}{0.714717in}}%
\pgfpathlineto{\pgfqpoint{5.188948in}{0.722096in}}%
\pgfpathlineto{\pgfqpoint{5.192304in}{0.726265in}}%
\pgfpathlineto{\pgfqpoint{5.195659in}{0.727650in}}%
\pgfpathlineto{\pgfqpoint{5.199015in}{0.726268in}}%
\pgfpathlineto{\pgfqpoint{5.202371in}{0.721691in}}%
\pgfpathlineto{\pgfqpoint{5.205727in}{0.712674in}}%
\pgfpathlineto{\pgfqpoint{5.209082in}{0.695566in}}%
\pgfpathlineto{\pgfqpoint{5.212438in}{0.651537in}}%
\pgfpathlineto{\pgfqpoint{5.215794in}{0.668542in}}%
\pgfpathlineto{\pgfqpoint{5.219150in}{0.704791in}}%
\pgfpathlineto{\pgfqpoint{5.222506in}{0.722112in}}%
\pgfpathlineto{\pgfqpoint{5.225861in}{0.732225in}}%
\pgfpathlineto{\pgfqpoint{5.229217in}{0.737673in}}%
\pgfpathlineto{\pgfqpoint{5.232573in}{0.738886in}}%
\pgfpathlineto{\pgfqpoint{5.235929in}{0.734811in}}%
\pgfpathlineto{\pgfqpoint{5.239285in}{0.723225in}}%
\pgfpathlineto{\pgfqpoint{5.242640in}{0.704086in}}%
\pgfpathlineto{\pgfqpoint{5.245996in}{0.679417in}}%
\pgfpathlineto{\pgfqpoint{5.249352in}{0.677751in}}%
\pgfpathlineto{\pgfqpoint{5.252708in}{0.684854in}}%
\pgfpathlineto{\pgfqpoint{5.256063in}{0.679443in}}%
\pgfpathlineto{\pgfqpoint{5.259419in}{0.653421in}}%
\pgfpathlineto{\pgfqpoint{5.262775in}{0.665191in}}%
\pgfpathlineto{\pgfqpoint{5.266131in}{0.697093in}}%
\pgfpathlineto{\pgfqpoint{5.269487in}{0.711121in}}%
\pgfpathlineto{\pgfqpoint{5.272842in}{0.712372in}}%
\pgfpathlineto{\pgfqpoint{5.276198in}{0.699571in}}%
\pgfpathlineto{\pgfqpoint{5.279554in}{0.675072in}}%
\pgfpathlineto{\pgfqpoint{5.282910in}{0.630724in}}%
\pgfpathlineto{\pgfqpoint{5.286265in}{0.654449in}}%
\pgfpathlineto{\pgfqpoint{5.289621in}{0.667143in}}%
\pgfpathlineto{\pgfqpoint{5.292977in}{0.661396in}}%
\pgfpathlineto{\pgfqpoint{5.296333in}{0.625249in}}%
\pgfpathlineto{\pgfqpoint{5.299689in}{0.643981in}}%
\pgfpathlineto{\pgfqpoint{5.303044in}{0.680212in}}%
\pgfpathlineto{\pgfqpoint{5.306400in}{0.694447in}}%
\pgfpathlineto{\pgfqpoint{5.309756in}{0.695300in}}%
\pgfpathlineto{\pgfqpoint{5.313112in}{0.676149in}}%
\pgfpathlineto{\pgfqpoint{5.316467in}{0.610354in}}%
\pgfpathlineto{\pgfqpoint{5.319823in}{0.685616in}}%
\pgfpathlineto{\pgfqpoint{5.323179in}{0.706947in}}%
\pgfpathlineto{\pgfqpoint{5.326535in}{0.715890in}}%
\pgfpathlineto{\pgfqpoint{5.329891in}{0.717419in}}%
\pgfpathlineto{\pgfqpoint{5.333246in}{0.711856in}}%
\pgfpathlineto{\pgfqpoint{5.336602in}{0.695475in}}%
\pgfpathlineto{\pgfqpoint{5.339958in}{0.637885in}}%
\pgfpathlineto{\pgfqpoint{5.343314in}{0.682020in}}%
\pgfpathlineto{\pgfqpoint{5.346669in}{0.712750in}}%
\pgfpathlineto{\pgfqpoint{5.350025in}{0.725857in}}%
\pgfpathlineto{\pgfqpoint{5.353381in}{0.729025in}}%
\pgfpathlineto{\pgfqpoint{5.356737in}{0.720247in}}%
\pgfpathlineto{\pgfqpoint{5.360093in}{0.695550in}}%
\pgfpathlineto{\pgfqpoint{5.363448in}{0.611065in}}%
\pgfpathlineto{\pgfqpoint{5.366804in}{0.689528in}}%
\pgfpathlineto{\pgfqpoint{5.370160in}{0.707007in}}%
\pgfpathlineto{\pgfqpoint{5.373516in}{0.709490in}}%
\pgfpathlineto{\pgfqpoint{5.376871in}{0.699631in}}%
\pgfpathlineto{\pgfqpoint{5.380227in}{0.662099in}}%
\pgfpathlineto{\pgfqpoint{5.383583in}{0.672153in}}%
\pgfpathlineto{\pgfqpoint{5.386939in}{0.711029in}}%
\pgfpathlineto{\pgfqpoint{5.390295in}{0.725788in}}%
\pgfpathlineto{\pgfqpoint{5.393650in}{0.727926in}}%
\pgfpathlineto{\pgfqpoint{5.397006in}{0.715121in}}%
\pgfpathlineto{\pgfqpoint{5.403718in}{0.650954in}}%
\pgfpathlineto{\pgfqpoint{5.407074in}{0.697686in}}%
\pgfpathlineto{\pgfqpoint{5.410429in}{0.709570in}}%
\pgfpathlineto{\pgfqpoint{5.413785in}{0.708621in}}%
\pgfpathlineto{\pgfqpoint{5.417141in}{0.693885in}}%
\pgfpathlineto{\pgfqpoint{5.420497in}{0.625156in}}%
\pgfpathlineto{\pgfqpoint{5.423852in}{0.689013in}}%
\pgfpathlineto{\pgfqpoint{5.427208in}{0.717335in}}%
\pgfpathlineto{\pgfqpoint{5.430564in}{0.728201in}}%
\pgfpathlineto{\pgfqpoint{5.433920in}{0.726209in}}%
\pgfpathlineto{\pgfqpoint{5.437276in}{0.707745in}}%
\pgfpathlineto{\pgfqpoint{5.440631in}{0.645115in}}%
\pgfpathlineto{\pgfqpoint{5.443987in}{0.687966in}}%
\pgfpathlineto{\pgfqpoint{5.447343in}{0.713053in}}%
\pgfpathlineto{\pgfqpoint{5.450699in}{0.720407in}}%
\pgfpathlineto{\pgfqpoint{5.454054in}{0.717379in}}%
\pgfpathlineto{\pgfqpoint{5.457410in}{0.701160in}}%
\pgfpathlineto{\pgfqpoint{5.460766in}{0.625418in}}%
\pgfpathlineto{\pgfqpoint{5.464122in}{0.697422in}}%
\pgfpathlineto{\pgfqpoint{5.467478in}{0.724898in}}%
\pgfpathlineto{\pgfqpoint{5.470833in}{0.735621in}}%
\pgfpathlineto{\pgfqpoint{5.474189in}{0.734130in}}%
\pgfpathlineto{\pgfqpoint{5.477545in}{0.716929in}}%
\pgfpathlineto{\pgfqpoint{5.480901in}{0.664507in}}%
\pgfpathlineto{\pgfqpoint{5.487612in}{0.716945in}}%
\pgfpathlineto{\pgfqpoint{5.490968in}{0.724987in}}%
\pgfpathlineto{\pgfqpoint{5.494324in}{0.722132in}}%
\pgfpathlineto{\pgfqpoint{5.497680in}{0.705634in}}%
\pgfpathlineto{\pgfqpoint{5.501035in}{0.622046in}}%
\pgfpathlineto{\pgfqpoint{5.504391in}{0.704088in}}%
\pgfpathlineto{\pgfqpoint{5.507747in}{0.730849in}}%
\pgfpathlineto{\pgfqpoint{5.511103in}{0.741150in}}%
\pgfpathlineto{\pgfqpoint{5.514458in}{0.738930in}}%
\pgfpathlineto{\pgfqpoint{5.517814in}{0.722067in}}%
\pgfpathlineto{\pgfqpoint{5.521170in}{0.679390in}}%
\pgfpathlineto{\pgfqpoint{5.524526in}{0.676785in}}%
\pgfpathlineto{\pgfqpoint{5.527882in}{0.710124in}}%
\pgfpathlineto{\pgfqpoint{5.531237in}{0.717811in}}%
\pgfpathlineto{\pgfqpoint{5.534593in}{0.712276in}}%
\pgfpathlineto{\pgfqpoint{5.541305in}{0.659847in}}%
\pgfpathlineto{\pgfqpoint{5.544660in}{0.714958in}}%
\pgfpathlineto{\pgfqpoint{5.548016in}{0.733379in}}%
\pgfpathlineto{\pgfqpoint{5.551372in}{0.738267in}}%
\pgfpathlineto{\pgfqpoint{5.554728in}{0.728688in}}%
\pgfpathlineto{\pgfqpoint{5.558084in}{0.704697in}}%
\pgfpathlineto{\pgfqpoint{5.561439in}{0.619640in}}%
\pgfpathlineto{\pgfqpoint{5.564795in}{0.685867in}}%
\pgfpathlineto{\pgfqpoint{5.568151in}{0.704293in}}%
\pgfpathlineto{\pgfqpoint{5.571507in}{0.704487in}}%
\pgfpathlineto{\pgfqpoint{5.574863in}{0.687375in}}%
\pgfpathlineto{\pgfqpoint{5.578218in}{0.569971in}}%
\pgfpathlineto{\pgfqpoint{5.581574in}{0.697686in}}%
\pgfpathlineto{\pgfqpoint{5.584930in}{0.721430in}}%
\pgfpathlineto{\pgfqpoint{5.588286in}{0.729631in}}%
\pgfpathlineto{\pgfqpoint{5.591641in}{0.723507in}}%
\pgfpathlineto{\pgfqpoint{5.594997in}{0.702233in}}%
\pgfpathlineto{\pgfqpoint{5.598353in}{0.635039in}}%
\pgfpathlineto{\pgfqpoint{5.601709in}{0.679205in}}%
\pgfpathlineto{\pgfqpoint{5.605065in}{0.701570in}}%
\pgfpathlineto{\pgfqpoint{5.608420in}{0.705300in}}%
\pgfpathlineto{\pgfqpoint{5.611776in}{0.695536in}}%
\pgfpathlineto{\pgfqpoint{5.615132in}{0.655066in}}%
\pgfpathlineto{\pgfqpoint{5.618488in}{0.673291in}}%
\pgfpathlineto{\pgfqpoint{5.621843in}{0.709955in}}%
\pgfpathlineto{\pgfqpoint{5.625199in}{0.723922in}}%
\pgfpathlineto{\pgfqpoint{5.628555in}{0.724930in}}%
\pgfpathlineto{\pgfqpoint{5.631911in}{0.710007in}}%
\pgfpathlineto{\pgfqpoint{5.635267in}{0.666699in}}%
\pgfpathlineto{\pgfqpoint{5.638622in}{0.671863in}}%
\pgfpathlineto{\pgfqpoint{5.641978in}{0.705160in}}%
\pgfpathlineto{\pgfqpoint{5.645334in}{0.714929in}}%
\pgfpathlineto{\pgfqpoint{5.648690in}{0.713568in}}%
\pgfpathlineto{\pgfqpoint{5.652045in}{0.699746in}}%
\pgfpathlineto{\pgfqpoint{5.655401in}{0.645150in}}%
\pgfpathlineto{\pgfqpoint{5.658757in}{0.687272in}}%
\pgfpathlineto{\pgfqpoint{5.662113in}{0.718355in}}%
\pgfpathlineto{\pgfqpoint{5.665469in}{0.730540in}}%
\pgfpathlineto{\pgfqpoint{5.668824in}{0.730512in}}%
\pgfpathlineto{\pgfqpoint{5.672180in}{0.714614in}}%
\pgfpathlineto{\pgfqpoint{5.675536in}{0.664627in}}%
\pgfpathlineto{\pgfqpoint{5.678892in}{0.685846in}}%
\pgfpathlineto{\pgfqpoint{5.682247in}{0.715400in}}%
\pgfpathlineto{\pgfqpoint{5.685603in}{0.724527in}}%
\pgfpathlineto{\pgfqpoint{5.688959in}{0.723197in}}%
\pgfpathlineto{\pgfqpoint{5.692315in}{0.710137in}}%
\pgfpathlineto{\pgfqpoint{5.695671in}{0.662530in}}%
\pgfpathlineto{\pgfqpoint{5.702382in}{0.725447in}}%
\pgfpathlineto{\pgfqpoint{5.705738in}{0.738608in}}%
\pgfpathlineto{\pgfqpoint{5.709094in}{0.739826in}}%
\pgfpathlineto{\pgfqpoint{5.712449in}{0.726399in}}%
\pgfpathlineto{\pgfqpoint{5.715805in}{0.691275in}}%
\pgfpathlineto{\pgfqpoint{5.719161in}{0.667466in}}%
\pgfpathlineto{\pgfqpoint{5.722517in}{0.711157in}}%
\pgfpathlineto{\pgfqpoint{5.725873in}{0.722397in}}%
\pgfpathlineto{\pgfqpoint{5.729228in}{0.721035in}}%
\pgfpathlineto{\pgfqpoint{5.732584in}{0.705694in}}%
\pgfpathlineto{\pgfqpoint{5.735940in}{0.630070in}}%
\pgfpathlineto{\pgfqpoint{5.739296in}{0.702857in}}%
\pgfpathlineto{\pgfqpoint{5.742652in}{0.730363in}}%
\pgfpathlineto{\pgfqpoint{5.746007in}{0.740839in}}%
\pgfpathlineto{\pgfqpoint{5.749363in}{0.738363in}}%
\pgfpathlineto{\pgfqpoint{5.752719in}{0.721732in}}%
\pgfpathlineto{\pgfqpoint{5.756075in}{0.683092in}}%
\pgfpathlineto{\pgfqpoint{5.759430in}{0.663797in}}%
\pgfpathlineto{\pgfqpoint{5.762786in}{0.702377in}}%
\pgfpathlineto{\pgfqpoint{5.766142in}{0.709955in}}%
\pgfpathlineto{\pgfqpoint{5.769498in}{0.702444in}}%
\pgfpathlineto{\pgfqpoint{5.772854in}{0.666559in}}%
\pgfpathlineto{\pgfqpoint{5.776209in}{0.676186in}}%
\pgfpathlineto{\pgfqpoint{5.779565in}{0.715813in}}%
\pgfpathlineto{\pgfqpoint{5.782921in}{0.730681in}}%
\pgfpathlineto{\pgfqpoint{5.786277in}{0.732244in}}%
\pgfpathlineto{\pgfqpoint{5.789632in}{0.718801in}}%
\pgfpathlineto{\pgfqpoint{5.792988in}{0.688217in}}%
\pgfpathlineto{\pgfqpoint{5.796344in}{0.625189in}}%
\pgfpathlineto{\pgfqpoint{5.799700in}{0.690063in}}%
\pgfpathlineto{\pgfqpoint{5.803056in}{0.701983in}}%
\pgfpathlineto{\pgfqpoint{5.806411in}{0.698332in}}%
\pgfpathlineto{\pgfqpoint{5.809767in}{0.674059in}}%
\pgfpathlineto{\pgfqpoint{5.813123in}{0.643043in}}%
\pgfpathlineto{\pgfqpoint{5.816479in}{0.701337in}}%
\pgfpathlineto{\pgfqpoint{5.819834in}{0.720328in}}%
\pgfpathlineto{\pgfqpoint{5.823190in}{0.725414in}}%
\pgfpathlineto{\pgfqpoint{5.826546in}{0.715416in}}%
\pgfpathlineto{\pgfqpoint{5.829902in}{0.686815in}}%
\pgfpathlineto{\pgfqpoint{5.833258in}{0.626325in}}%
\pgfpathlineto{\pgfqpoint{5.836613in}{0.692846in}}%
\pgfpathlineto{\pgfqpoint{5.839969in}{0.707572in}}%
\pgfpathlineto{\pgfqpoint{5.843325in}{0.708784in}}%
\pgfpathlineto{\pgfqpoint{5.846681in}{0.697676in}}%
\pgfpathlineto{\pgfqpoint{5.850036in}{0.655662in}}%
\pgfpathlineto{\pgfqpoint{5.853392in}{0.675442in}}%
\pgfpathlineto{\pgfqpoint{5.856748in}{0.711406in}}%
\pgfpathlineto{\pgfqpoint{5.860104in}{0.725245in}}%
\pgfpathlineto{\pgfqpoint{5.863460in}{0.726568in}}%
\pgfpathlineto{\pgfqpoint{5.866815in}{0.712051in}}%
\pgfpathlineto{\pgfqpoint{5.870171in}{0.666578in}}%
\pgfpathlineto{\pgfqpoint{5.873527in}{0.679143in}}%
\pgfpathlineto{\pgfqpoint{5.876883in}{0.711179in}}%
\pgfpathlineto{\pgfqpoint{5.880238in}{0.721293in}}%
\pgfpathlineto{\pgfqpoint{5.883594in}{0.720949in}}%
\pgfpathlineto{\pgfqpoint{5.886950in}{0.709663in}}%
\pgfpathlineto{\pgfqpoint{5.890306in}{0.671494in}}%
\pgfpathlineto{\pgfqpoint{5.893662in}{0.680425in}}%
\pgfpathlineto{\pgfqpoint{5.897017in}{0.719329in}}%
\pgfpathlineto{\pgfqpoint{5.900373in}{0.734278in}}%
\pgfpathlineto{\pgfqpoint{5.903729in}{0.737250in}}%
\pgfpathlineto{\pgfqpoint{5.907085in}{0.725851in}}%
\pgfpathlineto{\pgfqpoint{5.910441in}{0.692710in}}%
\pgfpathlineto{\pgfqpoint{5.913796in}{0.665729in}}%
\pgfpathlineto{\pgfqpoint{5.917152in}{0.712469in}}%
\pgfpathlineto{\pgfqpoint{5.920508in}{0.725356in}}%
\pgfpathlineto{\pgfqpoint{5.923864in}{0.726267in}}%
\pgfpathlineto{\pgfqpoint{5.927219in}{0.715860in}}%
\pgfpathlineto{\pgfqpoint{5.930575in}{0.679405in}}%
\pgfpathlineto{\pgfqpoint{5.933931in}{0.684884in}}%
\pgfpathlineto{\pgfqpoint{5.937287in}{0.725055in}}%
\pgfpathlineto{\pgfqpoint{5.940643in}{0.740267in}}%
\pgfpathlineto{\pgfqpoint{5.943998in}{0.743176in}}%
\pgfpathlineto{\pgfqpoint{5.947354in}{0.731906in}}%
\pgfpathlineto{\pgfqpoint{5.950710in}{0.704351in}}%
\pgfpathlineto{\pgfqpoint{5.954066in}{0.607357in}}%
\pgfpathlineto{\pgfqpoint{5.957421in}{0.700626in}}%
\pgfpathlineto{\pgfqpoint{5.960777in}{0.715322in}}%
\pgfpathlineto{\pgfqpoint{5.964133in}{0.714501in}}%
\pgfpathlineto{\pgfqpoint{5.967489in}{0.697445in}}%
\pgfpathlineto{\pgfqpoint{5.970845in}{0.486993in}}%
\pgfpathlineto{\pgfqpoint{5.974200in}{0.704898in}}%
\pgfpathlineto{\pgfqpoint{5.977556in}{0.729250in}}%
\pgfpathlineto{\pgfqpoint{5.980912in}{0.737896in}}%
\pgfpathlineto{\pgfqpoint{5.984268in}{0.732615in}}%
\pgfpathlineto{\pgfqpoint{5.987623in}{0.713764in}}%
\pgfpathlineto{\pgfqpoint{5.990979in}{0.668649in}}%
\pgfpathlineto{\pgfqpoint{5.994335in}{0.666821in}}%
\pgfpathlineto{\pgfqpoint{5.997691in}{0.697485in}}%
\pgfpathlineto{\pgfqpoint{6.001047in}{0.701935in}}%
\pgfpathlineto{\pgfqpoint{6.004402in}{0.689690in}}%
\pgfpathlineto{\pgfqpoint{6.007758in}{0.621311in}}%
\pgfpathlineto{\pgfqpoint{6.011114in}{0.687421in}}%
\pgfpathlineto{\pgfqpoint{6.014470in}{0.716017in}}%
\pgfpathlineto{\pgfqpoint{6.017825in}{0.726717in}}%
\pgfpathlineto{\pgfqpoint{6.021181in}{0.723459in}}%
\pgfpathlineto{\pgfqpoint{6.024537in}{0.705037in}}%
\pgfpathlineto{\pgfqpoint{6.027893in}{0.655532in}}%
\pgfpathlineto{\pgfqpoint{6.031249in}{0.668966in}}%
\pgfpathlineto{\pgfqpoint{6.034604in}{0.697438in}}%
\pgfpathlineto{\pgfqpoint{6.037960in}{0.703588in}}%
\pgfpathlineto{\pgfqpoint{6.041316in}{0.696566in}}%
\pgfpathlineto{\pgfqpoint{6.044672in}{0.666311in}}%
\pgfpathlineto{\pgfqpoint{6.048027in}{0.656985in}}%
\pgfpathlineto{\pgfqpoint{6.051383in}{0.703533in}}%
\pgfpathlineto{\pgfqpoint{6.054739in}{0.720229in}}%
\pgfpathlineto{\pgfqpoint{6.058095in}{0.723742in}}%
\pgfpathlineto{\pgfqpoint{6.061451in}{0.711814in}}%
\pgfpathlineto{\pgfqpoint{6.064806in}{0.675744in}}%
\pgfpathlineto{\pgfqpoint{6.068162in}{0.662048in}}%
\pgfpathlineto{\pgfqpoint{6.071518in}{0.703014in}}%
\pgfpathlineto{\pgfqpoint{6.074874in}{0.715158in}}%
\pgfpathlineto{\pgfqpoint{6.078230in}{0.715993in}}%
\pgfpathlineto{\pgfqpoint{6.081585in}{0.706011in}}%
\pgfpathlineto{\pgfqpoint{6.084941in}{0.672250in}}%
\pgfpathlineto{\pgfqpoint{6.088297in}{0.669380in}}%
\pgfpathlineto{\pgfqpoint{6.091653in}{0.712558in}}%
\pgfpathlineto{\pgfqpoint{6.095008in}{0.728614in}}%
\pgfpathlineto{\pgfqpoint{6.098364in}{0.732402in}}%
\pgfpathlineto{\pgfqpoint{6.101720in}{0.721778in}}%
\pgfpathlineto{\pgfqpoint{6.105076in}{0.687915in}}%
\pgfpathlineto{\pgfqpoint{6.108432in}{0.668219in}}%
\pgfpathlineto{\pgfqpoint{6.111787in}{0.712461in}}%
\pgfpathlineto{\pgfqpoint{6.115143in}{0.725671in}}%
\pgfpathlineto{\pgfqpoint{6.118499in}{0.727587in}}%
\pgfpathlineto{\pgfqpoint{6.121855in}{0.719458in}}%
\pgfpathlineto{\pgfqpoint{6.128566in}{0.661748in}}%
\pgfpathlineto{\pgfqpoint{6.131922in}{0.718380in}}%
\pgfpathlineto{\pgfqpoint{6.135278in}{0.737165in}}%
\pgfpathlineto{\pgfqpoint{6.138634in}{0.743151in}}%
\pgfpathlineto{\pgfqpoint{6.141989in}{0.735942in}}%
\pgfpathlineto{\pgfqpoint{6.145345in}{0.712869in}}%
\pgfpathlineto{\pgfqpoint{6.148701in}{0.611311in}}%
\pgfpathlineto{\pgfqpoint{6.152057in}{0.700871in}}%
\pgfpathlineto{\pgfqpoint{6.155412in}{0.720389in}}%
\pgfpathlineto{\pgfqpoint{6.158768in}{0.723734in}}%
\pgfpathlineto{\pgfqpoint{6.162124in}{0.714728in}}%
\pgfpathlineto{\pgfqpoint{6.165480in}{0.680040in}}%
\pgfpathlineto{\pgfqpoint{6.168836in}{0.682778in}}%
\pgfpathlineto{\pgfqpoint{6.172191in}{0.724205in}}%
\pgfpathlineto{\pgfqpoint{6.175547in}{0.739646in}}%
\pgfpathlineto{\pgfqpoint{6.178903in}{0.742303in}}%
\pgfpathlineto{\pgfqpoint{6.182259in}{0.730590in}}%
\pgfpathlineto{\pgfqpoint{6.185614in}{0.704399in}}%
\pgfpathlineto{\pgfqpoint{6.188970in}{0.578343in}}%
\pgfpathlineto{\pgfqpoint{6.192326in}{0.690222in}}%
\pgfpathlineto{\pgfqpoint{6.195682in}{0.705666in}}%
\pgfpathlineto{\pgfqpoint{6.199038in}{0.703086in}}%
\pgfpathlineto{\pgfqpoint{6.205749in}{0.651586in}}%
\pgfpathlineto{\pgfqpoint{6.209105in}{0.708134in}}%
\pgfpathlineto{\pgfqpoint{6.212461in}{0.726906in}}%
\pgfpathlineto{\pgfqpoint{6.215816in}{0.731687in}}%
\pgfpathlineto{\pgfqpoint{6.219172in}{0.721494in}}%
\pgfpathlineto{\pgfqpoint{6.222528in}{0.696727in}}%
\pgfpathlineto{\pgfqpoint{6.225884in}{0.599601in}}%
\pgfpathlineto{\pgfqpoint{6.229240in}{0.680034in}}%
\pgfpathlineto{\pgfqpoint{6.232595in}{0.697404in}}%
\pgfpathlineto{\pgfqpoint{6.235951in}{0.696810in}}%
\pgfpathlineto{\pgfqpoint{6.239307in}{0.678002in}}%
\pgfpathlineto{\pgfqpoint{6.242663in}{0.603715in}}%
\pgfpathlineto{\pgfqpoint{6.246019in}{0.693281in}}%
\pgfpathlineto{\pgfqpoint{6.249374in}{0.715542in}}%
\pgfpathlineto{\pgfqpoint{6.252730in}{0.722825in}}%
\pgfpathlineto{\pgfqpoint{6.256086in}{0.715430in}}%
\pgfpathlineto{\pgfqpoint{6.259442in}{0.690326in}}%
\pgfpathlineto{\pgfqpoint{6.262797in}{0.582123in}}%
\pgfpathlineto{\pgfqpoint{6.266153in}{0.688685in}}%
\pgfpathlineto{\pgfqpoint{6.269509in}{0.706161in}}%
\pgfpathlineto{\pgfqpoint{6.272865in}{0.709337in}}%
\pgfpathlineto{\pgfqpoint{6.276221in}{0.701222in}}%
\pgfpathlineto{\pgfqpoint{6.279576in}{0.671664in}}%
\pgfpathlineto{\pgfqpoint{6.282932in}{0.655984in}}%
\pgfpathlineto{\pgfqpoint{6.286288in}{0.704979in}}%
\pgfpathlineto{\pgfqpoint{6.289644in}{0.722274in}}%
\pgfpathlineto{\pgfqpoint{6.292999in}{0.726738in}}%
\pgfpathlineto{\pgfqpoint{6.296355in}{0.716538in}}%
\pgfpathlineto{\pgfqpoint{6.299711in}{0.682299in}}%
\pgfpathlineto{\pgfqpoint{6.303067in}{0.665966in}}%
\pgfpathlineto{\pgfqpoint{6.306423in}{0.709164in}}%
\pgfpathlineto{\pgfqpoint{6.309778in}{0.722519in}}%
\pgfpathlineto{\pgfqpoint{6.313134in}{0.724876in}}%
\pgfpathlineto{\pgfqpoint{6.316490in}{0.717689in}}%
\pgfpathlineto{\pgfqpoint{6.319846in}{0.693937in}}%
\pgfpathlineto{\pgfqpoint{6.323201in}{0.639665in}}%
\pgfpathlineto{\pgfqpoint{6.326557in}{0.711462in}}%
\pgfpathlineto{\pgfqpoint{6.329913in}{0.732196in}}%
\pgfpathlineto{\pgfqpoint{6.333269in}{0.739585in}}%
\pgfpathlineto{\pgfqpoint{6.336625in}{0.734398in}}%
\pgfpathlineto{\pgfqpoint{6.339980in}{0.712464in}}%
\pgfpathlineto{\pgfqpoint{6.343336in}{0.603973in}}%
\pgfpathlineto{\pgfqpoint{6.346692in}{0.703894in}}%
\pgfpathlineto{\pgfqpoint{6.350048in}{0.724458in}}%
\pgfpathlineto{\pgfqpoint{6.353403in}{0.729691in}}%
\pgfpathlineto{\pgfqpoint{6.356759in}{0.724439in}}%
\pgfpathlineto{\pgfqpoint{6.360115in}{0.703418in}}%
\pgfpathlineto{\pgfqpoint{6.363471in}{0.623062in}}%
\pgfpathlineto{\pgfqpoint{6.366827in}{0.715671in}}%
\pgfpathlineto{\pgfqpoint{6.370182in}{0.737913in}}%
\pgfpathlineto{\pgfqpoint{6.373538in}{0.745868in}}%
\pgfpathlineto{\pgfqpoint{6.376894in}{0.740988in}}%
\pgfpathlineto{\pgfqpoint{6.383605in}{0.719746in}}%
\pgfpathlineto{\pgfqpoint{6.383605in}{0.719746in}}%
\pgfusepath{stroke}%
\end{pgfscope}%
\begin{pgfscope}%
\pgfpathrectangle{\pgfqpoint{4.621821in}{0.320625in}}{\pgfqpoint{1.845679in}{1.164375in}}%
\pgfusepath{clip}%
\pgfsetrectcap%
\pgfsetroundjoin%
\pgfsetlinewidth{0.602250pt}%
\definecolor{currentstroke}{rgb}{0.478431,0.243137,0.615686}%
\pgfsetstrokecolor{currentstroke}%
\pgfsetstrokeopacity{0.120000}%
\pgfsetdash{}{0pt}%
\pgfpathmoveto{\pgfqpoint{4.705715in}{1.468362in}}%
\pgfpathlineto{\pgfqpoint{4.715783in}{1.449591in}}%
\pgfpathlineto{\pgfqpoint{4.725850in}{1.428347in}}%
\pgfpathlineto{\pgfqpoint{4.739273in}{1.396827in}}%
\pgfpathlineto{\pgfqpoint{4.779543in}{1.300125in}}%
\pgfpathlineto{\pgfqpoint{4.803033in}{1.247570in}}%
\pgfpathlineto{\pgfqpoint{4.826524in}{1.197745in}}%
\pgfpathlineto{\pgfqpoint{4.853370in}{1.143258in}}%
\pgfpathlineto{\pgfqpoint{4.886928in}{1.077698in}}%
\pgfpathlineto{\pgfqpoint{4.927197in}{1.001516in}}%
\pgfpathlineto{\pgfqpoint{4.977534in}{0.908753in}}%
\pgfpathlineto{\pgfqpoint{5.017803in}{0.836030in}}%
\pgfpathlineto{\pgfqpoint{5.031226in}{0.811625in}}%
\pgfpathlineto{\pgfqpoint{5.044649in}{0.787776in}}%
\pgfpathlineto{\pgfqpoint{5.054717in}{0.771338in}}%
\pgfpathlineto{\pgfqpoint{5.064784in}{0.751660in}}%
\pgfpathlineto{\pgfqpoint{5.074851in}{0.731747in}}%
\pgfpathlineto{\pgfqpoint{5.081563in}{0.721660in}}%
\pgfpathlineto{\pgfqpoint{5.088274in}{0.715075in}}%
\pgfpathlineto{\pgfqpoint{5.091630in}{0.712042in}}%
\pgfpathlineto{\pgfqpoint{5.094986in}{0.707753in}}%
\pgfpathlineto{\pgfqpoint{5.098342in}{0.701273in}}%
\pgfpathlineto{\pgfqpoint{5.101698in}{0.693307in}}%
\pgfpathlineto{\pgfqpoint{5.108409in}{0.672874in}}%
\pgfpathlineto{\pgfqpoint{5.118476in}{0.638745in}}%
\pgfpathlineto{\pgfqpoint{5.121832in}{0.631145in}}%
\pgfpathlineto{\pgfqpoint{5.125188in}{0.625503in}}%
\pgfpathlineto{\pgfqpoint{5.128544in}{0.624086in}}%
\pgfpathlineto{\pgfqpoint{5.131900in}{0.631152in}}%
\pgfpathlineto{\pgfqpoint{5.138611in}{0.653426in}}%
\pgfpathlineto{\pgfqpoint{5.141967in}{0.661036in}}%
\pgfpathlineto{\pgfqpoint{5.145323in}{0.665049in}}%
\pgfpathlineto{\pgfqpoint{5.148678in}{0.664557in}}%
\pgfpathlineto{\pgfqpoint{5.152034in}{0.657210in}}%
\pgfpathlineto{\pgfqpoint{5.155390in}{0.646908in}}%
\pgfpathlineto{\pgfqpoint{5.158746in}{0.660521in}}%
\pgfpathlineto{\pgfqpoint{5.162102in}{0.680764in}}%
\pgfpathlineto{\pgfqpoint{5.165457in}{0.695082in}}%
\pgfpathlineto{\pgfqpoint{5.168813in}{0.704742in}}%
\pgfpathlineto{\pgfqpoint{5.172169in}{0.711047in}}%
\pgfpathlineto{\pgfqpoint{5.175525in}{0.714699in}}%
\pgfpathlineto{\pgfqpoint{5.178880in}{0.716011in}}%
\pgfpathlineto{\pgfqpoint{5.182236in}{0.715008in}}%
\pgfpathlineto{\pgfqpoint{5.185592in}{0.711411in}}%
\pgfpathlineto{\pgfqpoint{5.188948in}{0.704441in}}%
\pgfpathlineto{\pgfqpoint{5.192304in}{0.692129in}}%
\pgfpathlineto{\pgfqpoint{5.195659in}{0.668013in}}%
\pgfpathlineto{\pgfqpoint{5.199015in}{0.611711in}}%
\pgfpathlineto{\pgfqpoint{5.202371in}{0.674670in}}%
\pgfpathlineto{\pgfqpoint{5.205727in}{0.699160in}}%
\pgfpathlineto{\pgfqpoint{5.209082in}{0.713002in}}%
\pgfpathlineto{\pgfqpoint{5.212438in}{0.721552in}}%
\pgfpathlineto{\pgfqpoint{5.215794in}{0.726338in}}%
\pgfpathlineto{\pgfqpoint{5.219150in}{0.727629in}}%
\pgfpathlineto{\pgfqpoint{5.222506in}{0.724707in}}%
\pgfpathlineto{\pgfqpoint{5.225861in}{0.715594in}}%
\pgfpathlineto{\pgfqpoint{5.232573in}{0.684363in}}%
\pgfpathlineto{\pgfqpoint{5.235929in}{0.682761in}}%
\pgfpathlineto{\pgfqpoint{5.239285in}{0.687919in}}%
\pgfpathlineto{\pgfqpoint{5.242640in}{0.686438in}}%
\pgfpathlineto{\pgfqpoint{5.245996in}{0.673609in}}%
\pgfpathlineto{\pgfqpoint{5.249352in}{0.646728in}}%
\pgfpathlineto{\pgfqpoint{5.256063in}{0.696944in}}%
\pgfpathlineto{\pgfqpoint{5.259419in}{0.706563in}}%
\pgfpathlineto{\pgfqpoint{5.262775in}{0.703395in}}%
\pgfpathlineto{\pgfqpoint{5.266131in}{0.686197in}}%
\pgfpathlineto{\pgfqpoint{5.269487in}{0.650247in}}%
\pgfpathlineto{\pgfqpoint{5.272842in}{0.654411in}}%
\pgfpathlineto{\pgfqpoint{5.276198in}{0.677356in}}%
\pgfpathlineto{\pgfqpoint{5.279554in}{0.682976in}}%
\pgfpathlineto{\pgfqpoint{5.282910in}{0.675980in}}%
\pgfpathlineto{\pgfqpoint{5.286265in}{0.646858in}}%
\pgfpathlineto{\pgfqpoint{5.289621in}{0.636689in}}%
\pgfpathlineto{\pgfqpoint{5.292977in}{0.682309in}}%
\pgfpathlineto{\pgfqpoint{5.296333in}{0.699013in}}%
\pgfpathlineto{\pgfqpoint{5.299689in}{0.702519in}}%
\pgfpathlineto{\pgfqpoint{5.303044in}{0.688289in}}%
\pgfpathlineto{\pgfqpoint{5.306400in}{0.612358in}}%
\pgfpathlineto{\pgfqpoint{5.309756in}{0.692125in}}%
\pgfpathlineto{\pgfqpoint{5.313112in}{0.716327in}}%
\pgfpathlineto{\pgfqpoint{5.316467in}{0.726865in}}%
\pgfpathlineto{\pgfqpoint{5.319823in}{0.729854in}}%
\pgfpathlineto{\pgfqpoint{5.323179in}{0.726189in}}%
\pgfpathlineto{\pgfqpoint{5.326535in}{0.713570in}}%
\pgfpathlineto{\pgfqpoint{5.329891in}{0.677864in}}%
\pgfpathlineto{\pgfqpoint{5.333246in}{0.675667in}}%
\pgfpathlineto{\pgfqpoint{5.336602in}{0.718170in}}%
\pgfpathlineto{\pgfqpoint{5.339958in}{0.734807in}}%
\pgfpathlineto{\pgfqpoint{5.343314in}{0.740828in}}%
\pgfpathlineto{\pgfqpoint{5.346669in}{0.736162in}}%
\pgfpathlineto{\pgfqpoint{5.350025in}{0.719239in}}%
\pgfpathlineto{\pgfqpoint{5.356737in}{0.652138in}}%
\pgfpathlineto{\pgfqpoint{5.360093in}{0.692591in}}%
\pgfpathlineto{\pgfqpoint{5.363448in}{0.698428in}}%
\pgfpathlineto{\pgfqpoint{5.366804in}{0.685347in}}%
\pgfpathlineto{\pgfqpoint{5.370160in}{0.590672in}}%
\pgfpathlineto{\pgfqpoint{5.373516in}{0.690842in}}%
\pgfpathlineto{\pgfqpoint{5.376871in}{0.716898in}}%
\pgfpathlineto{\pgfqpoint{5.380227in}{0.726027in}}%
\pgfpathlineto{\pgfqpoint{5.383583in}{0.720251in}}%
\pgfpathlineto{\pgfqpoint{5.386939in}{0.700934in}}%
\pgfpathlineto{\pgfqpoint{5.390295in}{0.655992in}}%
\pgfpathlineto{\pgfqpoint{5.393650in}{0.652261in}}%
\pgfpathlineto{\pgfqpoint{5.397006in}{0.681793in}}%
\pgfpathlineto{\pgfqpoint{5.400362in}{0.684269in}}%
\pgfpathlineto{\pgfqpoint{5.403718in}{0.666556in}}%
\pgfpathlineto{\pgfqpoint{5.407074in}{0.599330in}}%
\pgfpathlineto{\pgfqpoint{5.410429in}{0.684071in}}%
\pgfpathlineto{\pgfqpoint{5.413785in}{0.706107in}}%
\pgfpathlineto{\pgfqpoint{5.417141in}{0.712797in}}%
\pgfpathlineto{\pgfqpoint{5.420497in}{0.703642in}}%
\pgfpathlineto{\pgfqpoint{5.427208in}{0.644964in}}%
\pgfpathlineto{\pgfqpoint{5.430564in}{0.692546in}}%
\pgfpathlineto{\pgfqpoint{5.433920in}{0.706233in}}%
\pgfpathlineto{\pgfqpoint{5.437276in}{0.708215in}}%
\pgfpathlineto{\pgfqpoint{5.440631in}{0.699830in}}%
\pgfpathlineto{\pgfqpoint{5.443987in}{0.671522in}}%
\pgfpathlineto{\pgfqpoint{5.447343in}{0.648176in}}%
\pgfpathlineto{\pgfqpoint{5.450699in}{0.700705in}}%
\pgfpathlineto{\pgfqpoint{5.454054in}{0.718647in}}%
\pgfpathlineto{\pgfqpoint{5.457410in}{0.723734in}}%
\pgfpathlineto{\pgfqpoint{5.460766in}{0.714288in}}%
\pgfpathlineto{\pgfqpoint{5.464122in}{0.674377in}}%
\pgfpathlineto{\pgfqpoint{5.467478in}{0.682340in}}%
\pgfpathlineto{\pgfqpoint{5.470833in}{0.718271in}}%
\pgfpathlineto{\pgfqpoint{5.474189in}{0.731184in}}%
\pgfpathlineto{\pgfqpoint{5.477545in}{0.734380in}}%
\pgfpathlineto{\pgfqpoint{5.480901in}{0.729384in}}%
\pgfpathlineto{\pgfqpoint{5.484256in}{0.712319in}}%
\pgfpathlineto{\pgfqpoint{5.487612in}{0.640884in}}%
\pgfpathlineto{\pgfqpoint{5.490968in}{0.705890in}}%
\pgfpathlineto{\pgfqpoint{5.494324in}{0.733915in}}%
\pgfpathlineto{\pgfqpoint{5.497680in}{0.745334in}}%
\pgfpathlineto{\pgfqpoint{5.501035in}{0.745703in}}%
\pgfpathlineto{\pgfqpoint{5.504391in}{0.733004in}}%
\pgfpathlineto{\pgfqpoint{5.507747in}{0.707351in}}%
\pgfpathlineto{\pgfqpoint{5.511103in}{0.625044in}}%
\pgfpathlineto{\pgfqpoint{5.514458in}{0.682188in}}%
\pgfpathlineto{\pgfqpoint{5.517814in}{0.697174in}}%
\pgfpathlineto{\pgfqpoint{5.521170in}{0.689419in}}%
\pgfpathlineto{\pgfqpoint{5.524526in}{0.632074in}}%
\pgfpathlineto{\pgfqpoint{5.527882in}{0.686200in}}%
\pgfpathlineto{\pgfqpoint{5.531237in}{0.716780in}}%
\pgfpathlineto{\pgfqpoint{5.534593in}{0.728142in}}%
\pgfpathlineto{\pgfqpoint{5.537949in}{0.724709in}}%
\pgfpathlineto{\pgfqpoint{5.541305in}{0.707746in}}%
\pgfpathlineto{\pgfqpoint{5.544660in}{0.672950in}}%
\pgfpathlineto{\pgfqpoint{5.548016in}{0.623165in}}%
\pgfpathlineto{\pgfqpoint{5.551372in}{0.674056in}}%
\pgfpathlineto{\pgfqpoint{5.554728in}{0.679317in}}%
\pgfpathlineto{\pgfqpoint{5.558084in}{0.661528in}}%
\pgfpathlineto{\pgfqpoint{5.561439in}{0.610503in}}%
\pgfpathlineto{\pgfqpoint{5.564795in}{0.684135in}}%
\pgfpathlineto{\pgfqpoint{5.568151in}{0.705389in}}%
\pgfpathlineto{\pgfqpoint{5.571507in}{0.711453in}}%
\pgfpathlineto{\pgfqpoint{5.574863in}{0.701294in}}%
\pgfpathlineto{\pgfqpoint{5.581574in}{0.641242in}}%
\pgfpathlineto{\pgfqpoint{5.584930in}{0.688517in}}%
\pgfpathlineto{\pgfqpoint{5.588286in}{0.701579in}}%
\pgfpathlineto{\pgfqpoint{5.591641in}{0.702609in}}%
\pgfpathlineto{\pgfqpoint{5.594997in}{0.692321in}}%
\pgfpathlineto{\pgfqpoint{5.598353in}{0.656224in}}%
\pgfpathlineto{\pgfqpoint{5.601709in}{0.660716in}}%
\pgfpathlineto{\pgfqpoint{5.605065in}{0.701142in}}%
\pgfpathlineto{\pgfqpoint{5.608420in}{0.716277in}}%
\pgfpathlineto{\pgfqpoint{5.611776in}{0.718867in}}%
\pgfpathlineto{\pgfqpoint{5.615132in}{0.704858in}}%
\pgfpathlineto{\pgfqpoint{5.618488in}{0.639517in}}%
\pgfpathlineto{\pgfqpoint{5.621843in}{0.693281in}}%
\pgfpathlineto{\pgfqpoint{5.625199in}{0.720167in}}%
\pgfpathlineto{\pgfqpoint{5.628555in}{0.730439in}}%
\pgfpathlineto{\pgfqpoint{5.631911in}{0.732013in}}%
\pgfpathlineto{\pgfqpoint{5.635267in}{0.725384in}}%
\pgfpathlineto{\pgfqpoint{5.638622in}{0.705189in}}%
\pgfpathlineto{\pgfqpoint{5.641978in}{0.398674in}}%
\pgfpathlineto{\pgfqpoint{5.645334in}{0.709631in}}%
\pgfpathlineto{\pgfqpoint{5.648690in}{0.733923in}}%
\pgfpathlineto{\pgfqpoint{5.652045in}{0.743704in}}%
\pgfpathlineto{\pgfqpoint{5.655401in}{0.742442in}}%
\pgfpathlineto{\pgfqpoint{5.658757in}{0.728167in}}%
\pgfpathlineto{\pgfqpoint{5.662113in}{0.699015in}}%
\pgfpathlineto{\pgfqpoint{5.665469in}{0.580977in}}%
\pgfpathlineto{\pgfqpoint{5.668824in}{0.688048in}}%
\pgfpathlineto{\pgfqpoint{5.672180in}{0.698921in}}%
\pgfpathlineto{\pgfqpoint{5.675536in}{0.689408in}}%
\pgfpathlineto{\pgfqpoint{5.678892in}{0.626108in}}%
\pgfpathlineto{\pgfqpoint{5.682247in}{0.687591in}}%
\pgfpathlineto{\pgfqpoint{5.685603in}{0.717019in}}%
\pgfpathlineto{\pgfqpoint{5.688959in}{0.727823in}}%
\pgfpathlineto{\pgfqpoint{5.692315in}{0.723844in}}%
\pgfpathlineto{\pgfqpoint{5.695671in}{0.706249in}}%
\pgfpathlineto{\pgfqpoint{5.702382in}{0.637082in}}%
\pgfpathlineto{\pgfqpoint{5.705738in}{0.678234in}}%
\pgfpathlineto{\pgfqpoint{5.709094in}{0.682739in}}%
\pgfpathlineto{\pgfqpoint{5.712449in}{0.665976in}}%
\pgfpathlineto{\pgfqpoint{5.715805in}{0.596313in}}%
\pgfpathlineto{\pgfqpoint{5.719161in}{0.683764in}}%
\pgfpathlineto{\pgfqpoint{5.722517in}{0.706076in}}%
\pgfpathlineto{\pgfqpoint{5.725873in}{0.712892in}}%
\pgfpathlineto{\pgfqpoint{5.729228in}{0.703809in}}%
\pgfpathlineto{\pgfqpoint{5.732584in}{0.673335in}}%
\pgfpathlineto{\pgfqpoint{5.735940in}{0.635852in}}%
\pgfpathlineto{\pgfqpoint{5.739296in}{0.688801in}}%
\pgfpathlineto{\pgfqpoint{5.742652in}{0.702841in}}%
\pgfpathlineto{\pgfqpoint{5.746007in}{0.704541in}}%
\pgfpathlineto{\pgfqpoint{5.749363in}{0.695187in}}%
\pgfpathlineto{\pgfqpoint{5.752719in}{0.662636in}}%
\pgfpathlineto{\pgfqpoint{5.756075in}{0.656583in}}%
\pgfpathlineto{\pgfqpoint{5.759430in}{0.701047in}}%
\pgfpathlineto{\pgfqpoint{5.762786in}{0.717233in}}%
\pgfpathlineto{\pgfqpoint{5.766142in}{0.720740in}}%
\pgfpathlineto{\pgfqpoint{5.769498in}{0.708436in}}%
\pgfpathlineto{\pgfqpoint{5.772854in}{0.656614in}}%
\pgfpathlineto{\pgfqpoint{5.779565in}{0.718466in}}%
\pgfpathlineto{\pgfqpoint{5.782921in}{0.729597in}}%
\pgfpathlineto{\pgfqpoint{5.786277in}{0.731589in}}%
\pgfpathlineto{\pgfqpoint{5.789632in}{0.725228in}}%
\pgfpathlineto{\pgfqpoint{5.792988in}{0.705305in}}%
\pgfpathlineto{\pgfqpoint{5.796344in}{0.522861in}}%
\pgfpathlineto{\pgfqpoint{5.799700in}{0.709327in}}%
\pgfpathlineto{\pgfqpoint{5.803056in}{0.733751in}}%
\pgfpathlineto{\pgfqpoint{5.806411in}{0.743505in}}%
\pgfpathlineto{\pgfqpoint{5.809767in}{0.742023in}}%
\pgfpathlineto{\pgfqpoint{5.813123in}{0.727426in}}%
\pgfpathlineto{\pgfqpoint{5.816479in}{0.697389in}}%
\pgfpathlineto{\pgfqpoint{5.819834in}{0.604836in}}%
\pgfpathlineto{\pgfqpoint{5.823190in}{0.689579in}}%
\pgfpathlineto{\pgfqpoint{5.826546in}{0.699735in}}%
\pgfpathlineto{\pgfqpoint{5.829902in}{0.689967in}}%
\pgfpathlineto{\pgfqpoint{5.833258in}{0.626476in}}%
\pgfpathlineto{\pgfqpoint{5.836613in}{0.687976in}}%
\pgfpathlineto{\pgfqpoint{5.839969in}{0.717349in}}%
\pgfpathlineto{\pgfqpoint{5.843325in}{0.728124in}}%
\pgfpathlineto{\pgfqpoint{5.846681in}{0.724144in}}%
\pgfpathlineto{\pgfqpoint{5.850036in}{0.706477in}}%
\pgfpathlineto{\pgfqpoint{5.853392in}{0.668259in}}%
\pgfpathlineto{\pgfqpoint{5.856748in}{0.639692in}}%
\pgfpathlineto{\pgfqpoint{5.860104in}{0.679586in}}%
\pgfpathlineto{\pgfqpoint{5.863460in}{0.684086in}}%
\pgfpathlineto{\pgfqpoint{5.866815in}{0.667818in}}%
\pgfpathlineto{\pgfqpoint{5.870171in}{0.588160in}}%
\pgfpathlineto{\pgfqpoint{5.873527in}{0.683773in}}%
\pgfpathlineto{\pgfqpoint{5.876883in}{0.706559in}}%
\pgfpathlineto{\pgfqpoint{5.880238in}{0.713693in}}%
\pgfpathlineto{\pgfqpoint{5.883594in}{0.705065in}}%
\pgfpathlineto{\pgfqpoint{5.886950in}{0.675679in}}%
\pgfpathlineto{\pgfqpoint{5.890306in}{0.631174in}}%
\pgfpathlineto{\pgfqpoint{5.893662in}{0.688313in}}%
\pgfpathlineto{\pgfqpoint{5.897017in}{0.702858in}}%
\pgfpathlineto{\pgfqpoint{5.900373in}{0.704783in}}%
\pgfpathlineto{\pgfqpoint{5.903729in}{0.695609in}}%
\pgfpathlineto{\pgfqpoint{5.907085in}{0.663468in}}%
\pgfpathlineto{\pgfqpoint{5.910441in}{0.656227in}}%
\pgfpathlineto{\pgfqpoint{5.913796in}{0.701210in}}%
\pgfpathlineto{\pgfqpoint{5.917152in}{0.717509in}}%
\pgfpathlineto{\pgfqpoint{5.920508in}{0.721075in}}%
\pgfpathlineto{\pgfqpoint{5.923864in}{0.708874in}}%
\pgfpathlineto{\pgfqpoint{5.927219in}{0.658704in}}%
\pgfpathlineto{\pgfqpoint{5.933931in}{0.717691in}}%
\pgfpathlineto{\pgfqpoint{5.937287in}{0.728878in}}%
\pgfpathlineto{\pgfqpoint{5.940643in}{0.730801in}}%
\pgfpathlineto{\pgfqpoint{5.943998in}{0.724237in}}%
\pgfpathlineto{\pgfqpoint{5.947354in}{0.703678in}}%
\pgfpathlineto{\pgfqpoint{5.950710in}{0.570303in}}%
\pgfpathlineto{\pgfqpoint{5.954066in}{0.710011in}}%
\pgfpathlineto{\pgfqpoint{5.957421in}{0.733723in}}%
\pgfpathlineto{\pgfqpoint{5.960777in}{0.743066in}}%
\pgfpathlineto{\pgfqpoint{5.964133in}{0.740990in}}%
\pgfpathlineto{\pgfqpoint{5.967489in}{0.725766in}}%
\pgfpathlineto{\pgfqpoint{5.970845in}{0.693993in}}%
\pgfpathlineto{\pgfqpoint{5.974200in}{0.627066in}}%
\pgfpathlineto{\pgfqpoint{5.977556in}{0.691662in}}%
\pgfpathlineto{\pgfqpoint{5.980912in}{0.700571in}}%
\pgfpathlineto{\pgfqpoint{5.984268in}{0.690163in}}%
\pgfpathlineto{\pgfqpoint{5.987623in}{0.623774in}}%
\pgfpathlineto{\pgfqpoint{5.990979in}{0.688866in}}%
\pgfpathlineto{\pgfqpoint{5.994335in}{0.717772in}}%
\pgfpathlineto{\pgfqpoint{5.997691in}{0.728323in}}%
\pgfpathlineto{\pgfqpoint{6.001047in}{0.724125in}}%
\pgfpathlineto{\pgfqpoint{6.004402in}{0.706192in}}%
\pgfpathlineto{\pgfqpoint{6.007758in}{0.666655in}}%
\pgfpathlineto{\pgfqpoint{6.011114in}{0.644138in}}%
\pgfpathlineto{\pgfqpoint{6.014470in}{0.681406in}}%
\pgfpathlineto{\pgfqpoint{6.017825in}{0.685651in}}%
\pgfpathlineto{\pgfqpoint{6.021181in}{0.669706in}}%
\pgfpathlineto{\pgfqpoint{6.024537in}{0.578713in}}%
\pgfpathlineto{\pgfqpoint{6.027893in}{0.683939in}}%
\pgfpathlineto{\pgfqpoint{6.031249in}{0.707141in}}%
\pgfpathlineto{\pgfqpoint{6.034604in}{0.714553in}}%
\pgfpathlineto{\pgfqpoint{6.037960in}{0.706337in}}%
\pgfpathlineto{\pgfqpoint{6.041316in}{0.677841in}}%
\pgfpathlineto{\pgfqpoint{6.044672in}{0.626805in}}%
\pgfpathlineto{\pgfqpoint{6.048027in}{0.688149in}}%
\pgfpathlineto{\pgfqpoint{6.051383in}{0.703161in}}%
\pgfpathlineto{\pgfqpoint{6.054739in}{0.705312in}}%
\pgfpathlineto{\pgfqpoint{6.058095in}{0.696355in}}%
\pgfpathlineto{\pgfqpoint{6.061451in}{0.664812in}}%
\pgfpathlineto{\pgfqpoint{6.064806in}{0.655653in}}%
\pgfpathlineto{\pgfqpoint{6.068162in}{0.701475in}}%
\pgfpathlineto{\pgfqpoint{6.071518in}{0.717964in}}%
\pgfpathlineto{\pgfqpoint{6.074874in}{0.721665in}}%
\pgfpathlineto{\pgfqpoint{6.078230in}{0.709710in}}%
\pgfpathlineto{\pgfqpoint{6.081585in}{0.661835in}}%
\pgfpathlineto{\pgfqpoint{6.084941in}{0.685569in}}%
\pgfpathlineto{\pgfqpoint{6.088297in}{0.716895in}}%
\pgfpathlineto{\pgfqpoint{6.091653in}{0.728228in}}%
\pgfpathlineto{\pgfqpoint{6.095008in}{0.730131in}}%
\pgfpathlineto{\pgfqpoint{6.098364in}{0.723404in}}%
\pgfpathlineto{\pgfqpoint{6.101720in}{0.702263in}}%
\pgfpathlineto{\pgfqpoint{6.105076in}{0.598595in}}%
\pgfpathlineto{\pgfqpoint{6.108432in}{0.710679in}}%
\pgfpathlineto{\pgfqpoint{6.111787in}{0.733771in}}%
\pgfpathlineto{\pgfqpoint{6.115143in}{0.742731in}}%
\pgfpathlineto{\pgfqpoint{6.118499in}{0.740065in}}%
\pgfpathlineto{\pgfqpoint{6.121855in}{0.720843in}}%
\pgfpathlineto{\pgfqpoint{6.125210in}{0.643411in}}%
\pgfpathlineto{\pgfqpoint{6.128566in}{0.707566in}}%
\pgfpathlineto{\pgfqpoint{6.131922in}{0.730731in}}%
\pgfpathlineto{\pgfqpoint{6.135278in}{0.737474in}}%
\pgfpathlineto{\pgfqpoint{6.138634in}{0.734027in}}%
\pgfpathlineto{\pgfqpoint{6.141989in}{0.717178in}}%
\pgfpathlineto{\pgfqpoint{6.145345in}{0.632875in}}%
\pgfpathlineto{\pgfqpoint{6.148701in}{0.715370in}}%
\pgfpathlineto{\pgfqpoint{6.152057in}{0.741978in}}%
\pgfpathlineto{\pgfqpoint{6.155412in}{0.752151in}}%
\pgfpathlineto{\pgfqpoint{6.158768in}{0.749758in}}%
\pgfpathlineto{\pgfqpoint{6.162124in}{0.733761in}}%
\pgfpathlineto{\pgfqpoint{6.165480in}{0.699715in}}%
\pgfpathlineto{\pgfqpoint{6.168836in}{0.651940in}}%
\pgfpathlineto{\pgfqpoint{6.172191in}{0.703585in}}%
\pgfpathlineto{\pgfqpoint{6.175547in}{0.711107in}}%
\pgfpathlineto{\pgfqpoint{6.178903in}{0.699725in}}%
\pgfpathlineto{\pgfqpoint{6.182259in}{0.625755in}}%
\pgfpathlineto{\pgfqpoint{6.185614in}{0.700422in}}%
\pgfpathlineto{\pgfqpoint{6.188970in}{0.728375in}}%
\pgfpathlineto{\pgfqpoint{6.192326in}{0.738459in}}%
\pgfpathlineto{\pgfqpoint{6.195682in}{0.733781in}}%
\pgfpathlineto{\pgfqpoint{6.199038in}{0.716594in}}%
\pgfpathlineto{\pgfqpoint{6.202393in}{0.683074in}}%
\pgfpathlineto{\pgfqpoint{6.205749in}{0.613671in}}%
\pgfpathlineto{\pgfqpoint{6.209105in}{0.675804in}}%
\pgfpathlineto{\pgfqpoint{6.212461in}{0.679675in}}%
\pgfpathlineto{\pgfqpoint{6.215816in}{0.653273in}}%
\pgfpathlineto{\pgfqpoint{6.219172in}{0.650936in}}%
\pgfpathlineto{\pgfqpoint{6.222528in}{0.697298in}}%
\pgfpathlineto{\pgfqpoint{6.225884in}{0.714252in}}%
\pgfpathlineto{\pgfqpoint{6.229240in}{0.716767in}}%
\pgfpathlineto{\pgfqpoint{6.232595in}{0.702958in}}%
\pgfpathlineto{\pgfqpoint{6.239307in}{0.635137in}}%
\pgfpathlineto{\pgfqpoint{6.242663in}{0.682362in}}%
\pgfpathlineto{\pgfqpoint{6.246019in}{0.692754in}}%
\pgfpathlineto{\pgfqpoint{6.249374in}{0.688928in}}%
\pgfpathlineto{\pgfqpoint{6.252730in}{0.665942in}}%
\pgfpathlineto{\pgfqpoint{6.256086in}{0.623854in}}%
\pgfpathlineto{\pgfqpoint{6.259442in}{0.688633in}}%
\pgfpathlineto{\pgfqpoint{6.262797in}{0.708349in}}%
\pgfpathlineto{\pgfqpoint{6.266153in}{0.713805in}}%
\pgfpathlineto{\pgfqpoint{6.269509in}{0.703448in}}%
\pgfpathlineto{\pgfqpoint{6.272865in}{0.662984in}}%
\pgfpathlineto{\pgfqpoint{6.276221in}{0.669746in}}%
\pgfpathlineto{\pgfqpoint{6.279576in}{0.705312in}}%
\pgfpathlineto{\pgfqpoint{6.282932in}{0.717620in}}%
\pgfpathlineto{\pgfqpoint{6.286288in}{0.719944in}}%
\pgfpathlineto{\pgfqpoint{6.289644in}{0.713431in}}%
\pgfpathlineto{\pgfqpoint{6.292999in}{0.692372in}}%
\pgfpathlineto{\pgfqpoint{6.296355in}{0.590605in}}%
\pgfpathlineto{\pgfqpoint{6.299711in}{0.701075in}}%
\pgfpathlineto{\pgfqpoint{6.303067in}{0.724085in}}%
\pgfpathlineto{\pgfqpoint{6.306423in}{0.732893in}}%
\pgfpathlineto{\pgfqpoint{6.309778in}{0.729807in}}%
\pgfpathlineto{\pgfqpoint{6.313134in}{0.707509in}}%
\pgfpathlineto{\pgfqpoint{6.316490in}{0.624597in}}%
\pgfpathlineto{\pgfqpoint{6.319846in}{0.714063in}}%
\pgfpathlineto{\pgfqpoint{6.323201in}{0.733544in}}%
\pgfpathlineto{\pgfqpoint{6.326557in}{0.740092in}}%
\pgfpathlineto{\pgfqpoint{6.329913in}{0.738054in}}%
\pgfpathlineto{\pgfqpoint{6.333269in}{0.725734in}}%
\pgfpathlineto{\pgfqpoint{6.336625in}{0.686410in}}%
\pgfpathlineto{\pgfqpoint{6.339980in}{0.696593in}}%
\pgfpathlineto{\pgfqpoint{6.343336in}{0.734889in}}%
\pgfpathlineto{\pgfqpoint{6.346692in}{0.749805in}}%
\pgfpathlineto{\pgfqpoint{6.350048in}{0.753203in}}%
\pgfpathlineto{\pgfqpoint{6.353403in}{0.743965in}}%
\pgfpathlineto{\pgfqpoint{6.356759in}{0.727162in}}%
\pgfpathlineto{\pgfqpoint{6.360115in}{0.701756in}}%
\pgfpathlineto{\pgfqpoint{6.363471in}{0.659769in}}%
\pgfpathlineto{\pgfqpoint{6.366827in}{0.559813in}}%
\pgfpathlineto{\pgfqpoint{6.370182in}{0.610316in}}%
\pgfpathlineto{\pgfqpoint{6.373538in}{0.670952in}}%
\pgfpathlineto{\pgfqpoint{6.376894in}{0.698737in}}%
\pgfpathlineto{\pgfqpoint{6.380250in}{0.711653in}}%
\pgfpathlineto{\pgfqpoint{6.383605in}{0.710640in}}%
\pgfpathlineto{\pgfqpoint{6.383605in}{0.710640in}}%
\pgfusepath{stroke}%
\end{pgfscope}%
\begin{pgfscope}%
\pgfpathrectangle{\pgfqpoint{4.621821in}{0.320625in}}{\pgfqpoint{1.845679in}{1.164375in}}%
\pgfusepath{clip}%
\pgfsetrectcap%
\pgfsetroundjoin%
\pgfsetlinewidth{0.602250pt}%
\definecolor{currentstroke}{rgb}{0.478431,0.243137,0.615686}%
\pgfsetstrokecolor{currentstroke}%
\pgfsetstrokeopacity{0.120000}%
\pgfsetdash{}{0pt}%
\pgfpathmoveto{\pgfqpoint{4.705715in}{1.468904in}}%
\pgfpathlineto{\pgfqpoint{4.715783in}{1.450600in}}%
\pgfpathlineto{\pgfqpoint{4.725850in}{1.429890in}}%
\pgfpathlineto{\pgfqpoint{4.739273in}{1.399049in}}%
\pgfpathlineto{\pgfqpoint{4.779543in}{1.304329in}}%
\pgfpathlineto{\pgfqpoint{4.803033in}{1.252902in}}%
\pgfpathlineto{\pgfqpoint{4.829879in}{1.197114in}}%
\pgfpathlineto{\pgfqpoint{4.863437in}{1.130244in}}%
\pgfpathlineto{\pgfqpoint{4.903707in}{1.052533in}}%
\pgfpathlineto{\pgfqpoint{4.954043in}{0.957770in}}%
\pgfpathlineto{\pgfqpoint{5.014447in}{0.846328in}}%
\pgfpathlineto{\pgfqpoint{5.031226in}{0.815563in}}%
\pgfpathlineto{\pgfqpoint{5.044649in}{0.791484in}}%
\pgfpathlineto{\pgfqpoint{5.054717in}{0.774394in}}%
\pgfpathlineto{\pgfqpoint{5.068140in}{0.747436in}}%
\pgfpathlineto{\pgfqpoint{5.074851in}{0.734927in}}%
\pgfpathlineto{\pgfqpoint{5.081563in}{0.725536in}}%
\pgfpathlineto{\pgfqpoint{5.091630in}{0.714424in}}%
\pgfpathlineto{\pgfqpoint{5.094986in}{0.708222in}}%
\pgfpathlineto{\pgfqpoint{5.101698in}{0.691663in}}%
\pgfpathlineto{\pgfqpoint{5.108409in}{0.670113in}}%
\pgfpathlineto{\pgfqpoint{5.115121in}{0.647781in}}%
\pgfpathlineto{\pgfqpoint{5.118476in}{0.639558in}}%
\pgfpathlineto{\pgfqpoint{5.121832in}{0.635074in}}%
\pgfpathlineto{\pgfqpoint{5.125188in}{0.635966in}}%
\pgfpathlineto{\pgfqpoint{5.128544in}{0.642373in}}%
\pgfpathlineto{\pgfqpoint{5.135255in}{0.658896in}}%
\pgfpathlineto{\pgfqpoint{5.138611in}{0.664239in}}%
\pgfpathlineto{\pgfqpoint{5.141967in}{0.666070in}}%
\pgfpathlineto{\pgfqpoint{5.145323in}{0.662713in}}%
\pgfpathlineto{\pgfqpoint{5.148678in}{0.652627in}}%
\pgfpathlineto{\pgfqpoint{5.152034in}{0.638927in}}%
\pgfpathlineto{\pgfqpoint{5.155390in}{0.643217in}}%
\pgfpathlineto{\pgfqpoint{5.158746in}{0.660709in}}%
\pgfpathlineto{\pgfqpoint{5.162102in}{0.674040in}}%
\pgfpathlineto{\pgfqpoint{5.165457in}{0.682802in}}%
\pgfpathlineto{\pgfqpoint{5.168813in}{0.688123in}}%
\pgfpathlineto{\pgfqpoint{5.172169in}{0.690642in}}%
\pgfpathlineto{\pgfqpoint{5.175525in}{0.690580in}}%
\pgfpathlineto{\pgfqpoint{5.178880in}{0.687790in}}%
\pgfpathlineto{\pgfqpoint{5.182236in}{0.681628in}}%
\pgfpathlineto{\pgfqpoint{5.185592in}{0.670368in}}%
\pgfpathlineto{\pgfqpoint{5.188948in}{0.648535in}}%
\pgfpathlineto{\pgfqpoint{5.192304in}{0.586399in}}%
\pgfpathlineto{\pgfqpoint{5.195659in}{0.646886in}}%
\pgfpathlineto{\pgfqpoint{5.199015in}{0.673583in}}%
\pgfpathlineto{\pgfqpoint{5.202371in}{0.688146in}}%
\pgfpathlineto{\pgfqpoint{5.205727in}{0.696956in}}%
\pgfpathlineto{\pgfqpoint{5.209082in}{0.701721in}}%
\pgfpathlineto{\pgfqpoint{5.212438in}{0.702685in}}%
\pgfpathlineto{\pgfqpoint{5.215794in}{0.698908in}}%
\pgfpathlineto{\pgfqpoint{5.219150in}{0.688047in}}%
\pgfpathlineto{\pgfqpoint{5.222506in}{0.673757in}}%
\pgfpathlineto{\pgfqpoint{5.225861in}{0.678659in}}%
\pgfpathlineto{\pgfqpoint{5.229217in}{0.692743in}}%
\pgfpathlineto{\pgfqpoint{5.232573in}{0.701076in}}%
\pgfpathlineto{\pgfqpoint{5.235929in}{0.703266in}}%
\pgfpathlineto{\pgfqpoint{5.239285in}{0.699242in}}%
\pgfpathlineto{\pgfqpoint{5.242640in}{0.686566in}}%
\pgfpathlineto{\pgfqpoint{5.245996in}{0.652882in}}%
\pgfpathlineto{\pgfqpoint{5.249352in}{0.652370in}}%
\pgfpathlineto{\pgfqpoint{5.252708in}{0.691631in}}%
\pgfpathlineto{\pgfqpoint{5.256063in}{0.708238in}}%
\pgfpathlineto{\pgfqpoint{5.259419in}{0.714799in}}%
\pgfpathlineto{\pgfqpoint{5.262775in}{0.711698in}}%
\pgfpathlineto{\pgfqpoint{5.269487in}{0.664897in}}%
\pgfpathlineto{\pgfqpoint{5.272842in}{0.711908in}}%
\pgfpathlineto{\pgfqpoint{5.276198in}{0.730453in}}%
\pgfpathlineto{\pgfqpoint{5.279554in}{0.738679in}}%
\pgfpathlineto{\pgfqpoint{5.282910in}{0.740229in}}%
\pgfpathlineto{\pgfqpoint{5.286265in}{0.735335in}}%
\pgfpathlineto{\pgfqpoint{5.289621in}{0.721066in}}%
\pgfpathlineto{\pgfqpoint{5.292977in}{0.679051in}}%
\pgfpathlineto{\pgfqpoint{5.296333in}{0.693239in}}%
\pgfpathlineto{\pgfqpoint{5.299689in}{0.730370in}}%
\pgfpathlineto{\pgfqpoint{5.303044in}{0.745836in}}%
\pgfpathlineto{\pgfqpoint{5.306400in}{0.751425in}}%
\pgfpathlineto{\pgfqpoint{5.309756in}{0.746846in}}%
\pgfpathlineto{\pgfqpoint{5.313112in}{0.731791in}}%
\pgfpathlineto{\pgfqpoint{5.316467in}{0.705115in}}%
\pgfpathlineto{\pgfqpoint{5.319823in}{0.651441in}}%
\pgfpathlineto{\pgfqpoint{5.323179in}{0.677642in}}%
\pgfpathlineto{\pgfqpoint{5.326535in}{0.686558in}}%
\pgfpathlineto{\pgfqpoint{5.333246in}{0.647664in}}%
\pgfpathlineto{\pgfqpoint{5.336602in}{0.701461in}}%
\pgfpathlineto{\pgfqpoint{5.339958in}{0.721022in}}%
\pgfpathlineto{\pgfqpoint{5.343314in}{0.726364in}}%
\pgfpathlineto{\pgfqpoint{5.346669in}{0.716708in}}%
\pgfpathlineto{\pgfqpoint{5.350025in}{0.696545in}}%
\pgfpathlineto{\pgfqpoint{5.353381in}{0.656655in}}%
\pgfpathlineto{\pgfqpoint{5.356737in}{0.633577in}}%
\pgfpathlineto{\pgfqpoint{5.360093in}{0.662916in}}%
\pgfpathlineto{\pgfqpoint{5.363448in}{0.659341in}}%
\pgfpathlineto{\pgfqpoint{5.366804in}{0.608660in}}%
\pgfpathlineto{\pgfqpoint{5.370160in}{0.658160in}}%
\pgfpathlineto{\pgfqpoint{5.373516in}{0.689811in}}%
\pgfpathlineto{\pgfqpoint{5.376871in}{0.702197in}}%
\pgfpathlineto{\pgfqpoint{5.380227in}{0.700322in}}%
\pgfpathlineto{\pgfqpoint{5.383583in}{0.680161in}}%
\pgfpathlineto{\pgfqpoint{5.386939in}{0.596467in}}%
\pgfpathlineto{\pgfqpoint{5.390295in}{0.671969in}}%
\pgfpathlineto{\pgfqpoint{5.393650in}{0.693936in}}%
\pgfpathlineto{\pgfqpoint{5.397006in}{0.700740in}}%
\pgfpathlineto{\pgfqpoint{5.400362in}{0.698012in}}%
\pgfpathlineto{\pgfqpoint{5.403718in}{0.683390in}}%
\pgfpathlineto{\pgfqpoint{5.407074in}{0.627857in}}%
\pgfpathlineto{\pgfqpoint{5.410429in}{0.670607in}}%
\pgfpathlineto{\pgfqpoint{5.413785in}{0.701499in}}%
\pgfpathlineto{\pgfqpoint{5.417141in}{0.713823in}}%
\pgfpathlineto{\pgfqpoint{5.420497in}{0.714485in}}%
\pgfpathlineto{\pgfqpoint{5.423852in}{0.696859in}}%
\pgfpathlineto{\pgfqpoint{5.427208in}{0.601933in}}%
\pgfpathlineto{\pgfqpoint{5.430564in}{0.704777in}}%
\pgfpathlineto{\pgfqpoint{5.433920in}{0.727175in}}%
\pgfpathlineto{\pgfqpoint{5.437276in}{0.736701in}}%
\pgfpathlineto{\pgfqpoint{5.440631in}{0.738777in}}%
\pgfpathlineto{\pgfqpoint{5.443987in}{0.733950in}}%
\pgfpathlineto{\pgfqpoint{5.447343in}{0.719102in}}%
\pgfpathlineto{\pgfqpoint{5.450699in}{0.672291in}}%
\pgfpathlineto{\pgfqpoint{5.454054in}{0.697221in}}%
\pgfpathlineto{\pgfqpoint{5.457410in}{0.731571in}}%
\pgfpathlineto{\pgfqpoint{5.460766in}{0.745933in}}%
\pgfpathlineto{\pgfqpoint{5.464122in}{0.750188in}}%
\pgfpathlineto{\pgfqpoint{5.467478in}{0.743142in}}%
\pgfpathlineto{\pgfqpoint{5.470833in}{0.725680in}}%
\pgfpathlineto{\pgfqpoint{5.474189in}{0.692520in}}%
\pgfpathlineto{\pgfqpoint{5.477545in}{0.618714in}}%
\pgfpathlineto{\pgfqpoint{5.480901in}{0.682198in}}%
\pgfpathlineto{\pgfqpoint{5.484256in}{0.685525in}}%
\pgfpathlineto{\pgfqpoint{5.487612in}{0.655297in}}%
\pgfpathlineto{\pgfqpoint{5.490968in}{0.665308in}}%
\pgfpathlineto{\pgfqpoint{5.494324in}{0.707267in}}%
\pgfpathlineto{\pgfqpoint{5.497680in}{0.723327in}}%
\pgfpathlineto{\pgfqpoint{5.501035in}{0.725386in}}%
\pgfpathlineto{\pgfqpoint{5.504391in}{0.712728in}}%
\pgfpathlineto{\pgfqpoint{5.507747in}{0.689082in}}%
\pgfpathlineto{\pgfqpoint{5.511103in}{0.633300in}}%
\pgfpathlineto{\pgfqpoint{5.514458in}{0.645035in}}%
\pgfpathlineto{\pgfqpoint{5.517814in}{0.664400in}}%
\pgfpathlineto{\pgfqpoint{5.521170in}{0.654056in}}%
\pgfpathlineto{\pgfqpoint{5.524526in}{0.540370in}}%
\pgfpathlineto{\pgfqpoint{5.527882in}{0.669408in}}%
\pgfpathlineto{\pgfqpoint{5.531237in}{0.694519in}}%
\pgfpathlineto{\pgfqpoint{5.534593in}{0.703440in}}%
\pgfpathlineto{\pgfqpoint{5.537949in}{0.697078in}}%
\pgfpathlineto{\pgfqpoint{5.541305in}{0.670767in}}%
\pgfpathlineto{\pgfqpoint{5.544660in}{0.607457in}}%
\pgfpathlineto{\pgfqpoint{5.548016in}{0.678545in}}%
\pgfpathlineto{\pgfqpoint{5.551372in}{0.695372in}}%
\pgfpathlineto{\pgfqpoint{5.554728in}{0.699399in}}%
\pgfpathlineto{\pgfqpoint{5.558084in}{0.693655in}}%
\pgfpathlineto{\pgfqpoint{5.561439in}{0.672684in}}%
\pgfpathlineto{\pgfqpoint{5.564795in}{0.583609in}}%
\pgfpathlineto{\pgfqpoint{5.568151in}{0.683340in}}%
\pgfpathlineto{\pgfqpoint{5.571507in}{0.705944in}}%
\pgfpathlineto{\pgfqpoint{5.574863in}{0.714287in}}%
\pgfpathlineto{\pgfqpoint{5.578218in}{0.710003in}}%
\pgfpathlineto{\pgfqpoint{5.581574in}{0.680294in}}%
\pgfpathlineto{\pgfqpoint{5.584930in}{0.666691in}}%
\pgfpathlineto{\pgfqpoint{5.588286in}{0.712999in}}%
\pgfpathlineto{\pgfqpoint{5.591641in}{0.729593in}}%
\pgfpathlineto{\pgfqpoint{5.594997in}{0.736154in}}%
\pgfpathlineto{\pgfqpoint{5.598353in}{0.735722in}}%
\pgfpathlineto{\pgfqpoint{5.601709in}{0.727691in}}%
\pgfpathlineto{\pgfqpoint{5.605065in}{0.705892in}}%
\pgfpathlineto{\pgfqpoint{5.608420in}{0.596563in}}%
\pgfpathlineto{\pgfqpoint{5.611776in}{0.713297in}}%
\pgfpathlineto{\pgfqpoint{5.615132in}{0.736816in}}%
\pgfpathlineto{\pgfqpoint{5.618488in}{0.746802in}}%
\pgfpathlineto{\pgfqpoint{5.621843in}{0.746884in}}%
\pgfpathlineto{\pgfqpoint{5.625199in}{0.734947in}}%
\pgfpathlineto{\pgfqpoint{5.628555in}{0.711921in}}%
\pgfpathlineto{\pgfqpoint{5.631911in}{0.657042in}}%
\pgfpathlineto{\pgfqpoint{5.635267in}{0.668474in}}%
\pgfpathlineto{\pgfqpoint{5.638622in}{0.689533in}}%
\pgfpathlineto{\pgfqpoint{5.641978in}{0.682700in}}%
\pgfpathlineto{\pgfqpoint{5.645334in}{0.617141in}}%
\pgfpathlineto{\pgfqpoint{5.648690in}{0.684663in}}%
\pgfpathlineto{\pgfqpoint{5.652045in}{0.714247in}}%
\pgfpathlineto{\pgfqpoint{5.655401in}{0.725483in}}%
\pgfpathlineto{\pgfqpoint{5.658757in}{0.722181in}}%
\pgfpathlineto{\pgfqpoint{5.662113in}{0.706061in}}%
\pgfpathlineto{\pgfqpoint{5.665469in}{0.675339in}}%
\pgfpathlineto{\pgfqpoint{5.668824in}{0.561035in}}%
\pgfpathlineto{\pgfqpoint{5.672180in}{0.661271in}}%
\pgfpathlineto{\pgfqpoint{5.675536in}{0.668631in}}%
\pgfpathlineto{\pgfqpoint{5.678892in}{0.649489in}}%
\pgfpathlineto{\pgfqpoint{5.682247in}{0.617711in}}%
\pgfpathlineto{\pgfqpoint{5.685603in}{0.679462in}}%
\pgfpathlineto{\pgfqpoint{5.688959in}{0.699641in}}%
\pgfpathlineto{\pgfqpoint{5.692315in}{0.705226in}}%
\pgfpathlineto{\pgfqpoint{5.695671in}{0.694618in}}%
\pgfpathlineto{\pgfqpoint{5.699026in}{0.660219in}}%
\pgfpathlineto{\pgfqpoint{5.702382in}{0.642257in}}%
\pgfpathlineto{\pgfqpoint{5.705738in}{0.685579in}}%
\pgfpathlineto{\pgfqpoint{5.709094in}{0.698520in}}%
\pgfpathlineto{\pgfqpoint{5.712449in}{0.700173in}}%
\pgfpathlineto{\pgfqpoint{5.715805in}{0.691601in}}%
\pgfpathlineto{\pgfqpoint{5.719161in}{0.663145in}}%
\pgfpathlineto{\pgfqpoint{5.722517in}{0.640047in}}%
\pgfpathlineto{\pgfqpoint{5.725873in}{0.692469in}}%
\pgfpathlineto{\pgfqpoint{5.729228in}{0.710514in}}%
\pgfpathlineto{\pgfqpoint{5.732584in}{0.715919in}}%
\pgfpathlineto{\pgfqpoint{5.735940in}{0.707166in}}%
\pgfpathlineto{\pgfqpoint{5.739296in}{0.664040in}}%
\pgfpathlineto{\pgfqpoint{5.742652in}{0.683494in}}%
\pgfpathlineto{\pgfqpoint{5.746007in}{0.717663in}}%
\pgfpathlineto{\pgfqpoint{5.749363in}{0.731112in}}%
\pgfpathlineto{\pgfqpoint{5.752719in}{0.735621in}}%
\pgfpathlineto{\pgfqpoint{5.756075in}{0.733074in}}%
\pgfpathlineto{\pgfqpoint{5.759430in}{0.721736in}}%
\pgfpathlineto{\pgfqpoint{5.762786in}{0.689856in}}%
\pgfpathlineto{\pgfqpoint{5.766142in}{0.675285in}}%
\pgfpathlineto{\pgfqpoint{5.769498in}{0.723029in}}%
\pgfpathlineto{\pgfqpoint{5.772854in}{0.740696in}}%
\pgfpathlineto{\pgfqpoint{5.776209in}{0.747169in}}%
\pgfpathlineto{\pgfqpoint{5.779565in}{0.742689in}}%
\pgfpathlineto{\pgfqpoint{5.782921in}{0.726774in}}%
\pgfpathlineto{\pgfqpoint{5.786277in}{0.696069in}}%
\pgfpathlineto{\pgfqpoint{5.789632in}{0.589568in}}%
\pgfpathlineto{\pgfqpoint{5.792988in}{0.683679in}}%
\pgfpathlineto{\pgfqpoint{5.796344in}{0.691651in}}%
\pgfpathlineto{\pgfqpoint{5.799700in}{0.675084in}}%
\pgfpathlineto{\pgfqpoint{5.803056in}{0.622118in}}%
\pgfpathlineto{\pgfqpoint{5.806411in}{0.697945in}}%
\pgfpathlineto{\pgfqpoint{5.809767in}{0.719771in}}%
\pgfpathlineto{\pgfqpoint{5.813123in}{0.726538in}}%
\pgfpathlineto{\pgfqpoint{5.816479in}{0.718207in}}%
\pgfpathlineto{\pgfqpoint{5.819834in}{0.698284in}}%
\pgfpathlineto{\pgfqpoint{5.823190in}{0.655879in}}%
\pgfpathlineto{\pgfqpoint{5.826546in}{0.636914in}}%
\pgfpathlineto{\pgfqpoint{5.829902in}{0.669647in}}%
\pgfpathlineto{\pgfqpoint{5.833258in}{0.669131in}}%
\pgfpathlineto{\pgfqpoint{5.836613in}{0.637068in}}%
\pgfpathlineto{\pgfqpoint{5.839969in}{0.647049in}}%
\pgfpathlineto{\pgfqpoint{5.843325in}{0.688294in}}%
\pgfpathlineto{\pgfqpoint{5.846681in}{0.703916in}}%
\pgfpathlineto{\pgfqpoint{5.850036in}{0.705585in}}%
\pgfpathlineto{\pgfqpoint{5.853392in}{0.690089in}}%
\pgfpathlineto{\pgfqpoint{5.856748in}{0.640520in}}%
\pgfpathlineto{\pgfqpoint{5.860104in}{0.661411in}}%
\pgfpathlineto{\pgfqpoint{5.863460in}{0.691236in}}%
\pgfpathlineto{\pgfqpoint{5.866815in}{0.700585in}}%
\pgfpathlineto{\pgfqpoint{5.870171in}{0.699565in}}%
\pgfpathlineto{\pgfqpoint{5.873527in}{0.687155in}}%
\pgfpathlineto{\pgfqpoint{5.876883in}{0.643508in}}%
\pgfpathlineto{\pgfqpoint{5.880238in}{0.665030in}}%
\pgfpathlineto{\pgfqpoint{5.883594in}{0.700314in}}%
\pgfpathlineto{\pgfqpoint{5.886950in}{0.714114in}}%
\pgfpathlineto{\pgfqpoint{5.890306in}{0.715979in}}%
\pgfpathlineto{\pgfqpoint{5.893662in}{0.701035in}}%
\pgfpathlineto{\pgfqpoint{5.897017in}{0.617114in}}%
\pgfpathlineto{\pgfqpoint{5.900373in}{0.696629in}}%
\pgfpathlineto{\pgfqpoint{5.903729in}{0.721979in}}%
\pgfpathlineto{\pgfqpoint{5.907085in}{0.732272in}}%
\pgfpathlineto{\pgfqpoint{5.910441in}{0.734483in}}%
\pgfpathlineto{\pgfqpoint{5.913796in}{0.729327in}}%
\pgfpathlineto{\pgfqpoint{5.917152in}{0.713232in}}%
\pgfpathlineto{\pgfqpoint{5.920508in}{0.656216in}}%
\pgfpathlineto{\pgfqpoint{5.923864in}{0.699666in}}%
\pgfpathlineto{\pgfqpoint{5.927219in}{0.730453in}}%
\pgfpathlineto{\pgfqpoint{5.930575in}{0.743399in}}%
\pgfpathlineto{\pgfqpoint{5.933931in}{0.746079in}}%
\pgfpathlineto{\pgfqpoint{5.937287in}{0.736651in}}%
\pgfpathlineto{\pgfqpoint{5.940643in}{0.716022in}}%
\pgfpathlineto{\pgfqpoint{5.943998in}{0.670081in}}%
\pgfpathlineto{\pgfqpoint{5.947354in}{0.663000in}}%
\pgfpathlineto{\pgfqpoint{5.950710in}{0.691546in}}%
\pgfpathlineto{\pgfqpoint{5.954066in}{0.690474in}}%
\pgfpathlineto{\pgfqpoint{5.957421in}{0.659291in}}%
\pgfpathlineto{\pgfqpoint{5.960777in}{0.665701in}}%
\pgfpathlineto{\pgfqpoint{5.964133in}{0.707971in}}%
\pgfpathlineto{\pgfqpoint{5.967489in}{0.723836in}}%
\pgfpathlineto{\pgfqpoint{5.970845in}{0.725772in}}%
\pgfpathlineto{\pgfqpoint{5.974200in}{0.712890in}}%
\pgfpathlineto{\pgfqpoint{5.977556in}{0.687609in}}%
\pgfpathlineto{\pgfqpoint{5.980912in}{0.612988in}}%
\pgfpathlineto{\pgfqpoint{5.984268in}{0.658829in}}%
\pgfpathlineto{\pgfqpoint{5.987623in}{0.674684in}}%
\pgfpathlineto{\pgfqpoint{5.990979in}{0.666894in}}%
\pgfpathlineto{\pgfqpoint{5.994335in}{0.606250in}}%
\pgfpathlineto{\pgfqpoint{5.997691in}{0.665602in}}%
\pgfpathlineto{\pgfqpoint{6.001047in}{0.695762in}}%
\pgfpathlineto{\pgfqpoint{6.004402in}{0.707199in}}%
\pgfpathlineto{\pgfqpoint{6.007758in}{0.704209in}}%
\pgfpathlineto{\pgfqpoint{6.011114in}{0.683228in}}%
\pgfpathlineto{\pgfqpoint{6.014470in}{0.583379in}}%
\pgfpathlineto{\pgfqpoint{6.017825in}{0.674488in}}%
\pgfpathlineto{\pgfqpoint{6.021181in}{0.695910in}}%
\pgfpathlineto{\pgfqpoint{6.024537in}{0.701972in}}%
\pgfpathlineto{\pgfqpoint{6.027893in}{0.698024in}}%
\pgfpathlineto{\pgfqpoint{6.031249in}{0.680434in}}%
\pgfpathlineto{\pgfqpoint{6.034604in}{0.582630in}}%
\pgfpathlineto{\pgfqpoint{6.037960in}{0.680940in}}%
\pgfpathlineto{\pgfqpoint{6.041316in}{0.706807in}}%
\pgfpathlineto{\pgfqpoint{6.044672in}{0.716799in}}%
\pgfpathlineto{\pgfqpoint{6.048027in}{0.714516in}}%
\pgfpathlineto{\pgfqpoint{6.051383in}{0.691549in}}%
\pgfpathlineto{\pgfqpoint{6.054739in}{0.636228in}}%
\pgfpathlineto{\pgfqpoint{6.058095in}{0.706156in}}%
\pgfpathlineto{\pgfqpoint{6.061451in}{0.725398in}}%
\pgfpathlineto{\pgfqpoint{6.064806in}{0.732837in}}%
\pgfpathlineto{\pgfqpoint{6.068162in}{0.732653in}}%
\pgfpathlineto{\pgfqpoint{6.071518in}{0.724362in}}%
\pgfpathlineto{\pgfqpoint{6.074874in}{0.701029in}}%
\pgfpathlineto{\pgfqpoint{6.078230in}{0.630755in}}%
\pgfpathlineto{\pgfqpoint{6.081585in}{0.714053in}}%
\pgfpathlineto{\pgfqpoint{6.084941in}{0.735986in}}%
\pgfpathlineto{\pgfqpoint{6.088297in}{0.744865in}}%
\pgfpathlineto{\pgfqpoint{6.091653in}{0.743162in}}%
\pgfpathlineto{\pgfqpoint{6.095008in}{0.729108in}}%
\pgfpathlineto{\pgfqpoint{6.098364in}{0.701628in}}%
\pgfpathlineto{\pgfqpoint{6.101720in}{0.586916in}}%
\pgfpathlineto{\pgfqpoint{6.105076in}{0.682421in}}%
\pgfpathlineto{\pgfqpoint{6.108432in}{0.695017in}}%
\pgfpathlineto{\pgfqpoint{6.111787in}{0.685625in}}%
\pgfpathlineto{\pgfqpoint{6.115143in}{0.616050in}}%
\pgfpathlineto{\pgfqpoint{6.118499in}{0.686796in}}%
\pgfpathlineto{\pgfqpoint{6.121855in}{0.715606in}}%
\pgfpathlineto{\pgfqpoint{6.125210in}{0.726384in}}%
\pgfpathlineto{\pgfqpoint{6.128566in}{0.722691in}}%
\pgfpathlineto{\pgfqpoint{6.131922in}{0.705945in}}%
\pgfpathlineto{\pgfqpoint{6.135278in}{0.672203in}}%
\pgfpathlineto{\pgfqpoint{6.138634in}{0.614035in}}%
\pgfpathlineto{\pgfqpoint{6.141989in}{0.670599in}}%
\pgfpathlineto{\pgfqpoint{6.145345in}{0.676875in}}%
\pgfpathlineto{\pgfqpoint{6.148701in}{0.660826in}}%
\pgfpathlineto{\pgfqpoint{6.152057in}{0.591301in}}%
\pgfpathlineto{\pgfqpoint{6.155412in}{0.679275in}}%
\pgfpathlineto{\pgfqpoint{6.158768in}{0.701924in}}%
\pgfpathlineto{\pgfqpoint{6.162124in}{0.709242in}}%
\pgfpathlineto{\pgfqpoint{6.165480in}{0.701102in}}%
\pgfpathlineto{\pgfqpoint{6.168836in}{0.672608in}}%
\pgfpathlineto{\pgfqpoint{6.172191in}{0.623571in}}%
\pgfpathlineto{\pgfqpoint{6.175547in}{0.684127in}}%
\pgfpathlineto{\pgfqpoint{6.178903in}{0.699521in}}%
\pgfpathlineto{\pgfqpoint{6.182259in}{0.702429in}}%
\pgfpathlineto{\pgfqpoint{6.185614in}{0.695070in}}%
\pgfpathlineto{\pgfqpoint{6.188970in}{0.669546in}}%
\pgfpathlineto{\pgfqpoint{6.192326in}{0.631639in}}%
\pgfpathlineto{\pgfqpoint{6.195682in}{0.692664in}}%
\pgfpathlineto{\pgfqpoint{6.199038in}{0.712057in}}%
\pgfpathlineto{\pgfqpoint{6.202393in}{0.718308in}}%
\pgfpathlineto{\pgfqpoint{6.205749in}{0.710814in}}%
\pgfpathlineto{\pgfqpoint{6.209105in}{0.675604in}}%
\pgfpathlineto{\pgfqpoint{6.212461in}{0.673246in}}%
\pgfpathlineto{\pgfqpoint{6.215816in}{0.713470in}}%
\pgfpathlineto{\pgfqpoint{6.219172in}{0.727983in}}%
\pgfpathlineto{\pgfqpoint{6.222528in}{0.732699in}}%
\pgfpathlineto{\pgfqpoint{6.225884in}{0.729882in}}%
\pgfpathlineto{\pgfqpoint{6.229240in}{0.717527in}}%
\pgfpathlineto{\pgfqpoint{6.232595in}{0.680633in}}%
\pgfpathlineto{\pgfqpoint{6.235951in}{0.683211in}}%
\pgfpathlineto{\pgfqpoint{6.239307in}{0.724019in}}%
\pgfpathlineto{\pgfqpoint{6.242663in}{0.739963in}}%
\pgfpathlineto{\pgfqpoint{6.246019in}{0.744900in}}%
\pgfpathlineto{\pgfqpoint{6.249374in}{0.737954in}}%
\pgfpathlineto{\pgfqpoint{6.252730in}{0.719393in}}%
\pgfpathlineto{\pgfqpoint{6.256086in}{0.679723in}}%
\pgfpathlineto{\pgfqpoint{6.259442in}{0.654671in}}%
\pgfpathlineto{\pgfqpoint{6.262797in}{0.692054in}}%
\pgfpathlineto{\pgfqpoint{6.266153in}{0.695183in}}%
\pgfpathlineto{\pgfqpoint{6.269509in}{0.675605in}}%
\pgfpathlineto{\pgfqpoint{6.272865in}{0.633746in}}%
\pgfpathlineto{\pgfqpoint{6.276221in}{0.700723in}}%
\pgfpathlineto{\pgfqpoint{6.279576in}{0.721300in}}%
\pgfpathlineto{\pgfqpoint{6.282932in}{0.727207in}}%
\pgfpathlineto{\pgfqpoint{6.286288in}{0.718047in}}%
\pgfpathlineto{\pgfqpoint{6.289644in}{0.696676in}}%
\pgfpathlineto{\pgfqpoint{6.292999in}{0.645776in}}%
\pgfpathlineto{\pgfqpoint{6.296355in}{0.652555in}}%
\pgfpathlineto{\pgfqpoint{6.299711in}{0.677551in}}%
\pgfpathlineto{\pgfqpoint{6.303067in}{0.676448in}}%
\pgfpathlineto{\pgfqpoint{6.306423in}{0.648483in}}%
\pgfpathlineto{\pgfqpoint{6.309778in}{0.643431in}}%
\pgfpathlineto{\pgfqpoint{6.313134in}{0.689896in}}%
\pgfpathlineto{\pgfqpoint{6.316490in}{0.706788in}}%
\pgfpathlineto{\pgfqpoint{6.319846in}{0.709738in}}%
\pgfpathlineto{\pgfqpoint{6.323201in}{0.696107in}}%
\pgfpathlineto{\pgfqpoint{6.326557in}{0.654864in}}%
\pgfpathlineto{\pgfqpoint{6.329913in}{0.656370in}}%
\pgfpathlineto{\pgfqpoint{6.333269in}{0.691427in}}%
\pgfpathlineto{\pgfqpoint{6.336625in}{0.702093in}}%
\pgfpathlineto{\pgfqpoint{6.339980in}{0.701835in}}%
\pgfpathlineto{\pgfqpoint{6.343336in}{0.690256in}}%
\pgfpathlineto{\pgfqpoint{6.346692in}{0.649868in}}%
\pgfpathlineto{\pgfqpoint{6.350048in}{0.664544in}}%
\pgfpathlineto{\pgfqpoint{6.353403in}{0.701723in}}%
\pgfpathlineto{\pgfqpoint{6.356759in}{0.716089in}}%
\pgfpathlineto{\pgfqpoint{6.360115in}{0.718370in}}%
\pgfpathlineto{\pgfqpoint{6.363471in}{0.708673in}}%
\pgfpathlineto{\pgfqpoint{6.373538in}{0.674296in}}%
\pgfpathlineto{\pgfqpoint{6.376894in}{0.665632in}}%
\pgfpathlineto{\pgfqpoint{6.383605in}{0.651079in}}%
\pgfpathlineto{\pgfqpoint{6.383605in}{0.651079in}}%
\pgfusepath{stroke}%
\end{pgfscope}%
\begin{pgfscope}%
\pgfpathrectangle{\pgfqpoint{4.621821in}{0.320625in}}{\pgfqpoint{1.845679in}{1.164375in}}%
\pgfusepath{clip}%
\pgfsetrectcap%
\pgfsetroundjoin%
\pgfsetlinewidth{0.602250pt}%
\definecolor{currentstroke}{rgb}{0.478431,0.243137,0.615686}%
\pgfsetstrokecolor{currentstroke}%
\pgfsetstrokeopacity{0.120000}%
\pgfsetdash{}{0pt}%
\pgfpathmoveto{\pgfqpoint{4.705715in}{1.468747in}}%
\pgfpathlineto{\pgfqpoint{4.715783in}{1.450186in}}%
\pgfpathlineto{\pgfqpoint{4.725850in}{1.429202in}}%
\pgfpathlineto{\pgfqpoint{4.739273in}{1.398027in}}%
\pgfpathlineto{\pgfqpoint{4.776187in}{1.310285in}}%
\pgfpathlineto{\pgfqpoint{4.796322in}{1.265658in}}%
\pgfpathlineto{\pgfqpoint{4.819812in}{1.216530in}}%
\pgfpathlineto{\pgfqpoint{4.846658in}{1.163239in}}%
\pgfpathlineto{\pgfqpoint{4.876860in}{1.105745in}}%
\pgfpathlineto{\pgfqpoint{4.917130in}{1.031652in}}%
\pgfpathlineto{\pgfqpoint{4.967466in}{0.941491in}}%
\pgfpathlineto{\pgfqpoint{5.041293in}{0.812060in}}%
\pgfpathlineto{\pgfqpoint{5.054717in}{0.789122in}}%
\pgfpathlineto{\pgfqpoint{5.071496in}{0.758771in}}%
\pgfpathlineto{\pgfqpoint{5.078207in}{0.748966in}}%
\pgfpathlineto{\pgfqpoint{5.084919in}{0.740051in}}%
\pgfpathlineto{\pgfqpoint{5.091630in}{0.728202in}}%
\pgfpathlineto{\pgfqpoint{5.098342in}{0.713428in}}%
\pgfpathlineto{\pgfqpoint{5.108409in}{0.689407in}}%
\pgfpathlineto{\pgfqpoint{5.111765in}{0.682773in}}%
\pgfpathlineto{\pgfqpoint{5.115121in}{0.677957in}}%
\pgfpathlineto{\pgfqpoint{5.118476in}{0.675445in}}%
\pgfpathlineto{\pgfqpoint{5.121832in}{0.675182in}}%
\pgfpathlineto{\pgfqpoint{5.131900in}{0.678568in}}%
\pgfpathlineto{\pgfqpoint{5.135255in}{0.676550in}}%
\pgfpathlineto{\pgfqpoint{5.138611in}{0.670331in}}%
\pgfpathlineto{\pgfqpoint{5.141967in}{0.660729in}}%
\pgfpathlineto{\pgfqpoint{5.148678in}{0.635955in}}%
\pgfpathlineto{\pgfqpoint{5.152034in}{0.636366in}}%
\pgfpathlineto{\pgfqpoint{5.158746in}{0.653251in}}%
\pgfpathlineto{\pgfqpoint{5.162102in}{0.657458in}}%
\pgfpathlineto{\pgfqpoint{5.165457in}{0.658215in}}%
\pgfpathlineto{\pgfqpoint{5.168813in}{0.655373in}}%
\pgfpathlineto{\pgfqpoint{5.172169in}{0.648007in}}%
\pgfpathlineto{\pgfqpoint{5.175525in}{0.633154in}}%
\pgfpathlineto{\pgfqpoint{5.178880in}{0.597748in}}%
\pgfpathlineto{\pgfqpoint{5.182236in}{0.597228in}}%
\pgfpathlineto{\pgfqpoint{5.185592in}{0.637804in}}%
\pgfpathlineto{\pgfqpoint{5.188948in}{0.656651in}}%
\pgfpathlineto{\pgfqpoint{5.192304in}{0.667716in}}%
\pgfpathlineto{\pgfqpoint{5.195659in}{0.673965in}}%
\pgfpathlineto{\pgfqpoint{5.199015in}{0.676044in}}%
\pgfpathlineto{\pgfqpoint{5.202371in}{0.673211in}}%
\pgfpathlineto{\pgfqpoint{5.205727in}{0.662860in}}%
\pgfpathlineto{\pgfqpoint{5.209082in}{0.661104in}}%
\pgfpathlineto{\pgfqpoint{5.212438in}{0.682026in}}%
\pgfpathlineto{\pgfqpoint{5.215794in}{0.698616in}}%
\pgfpathlineto{\pgfqpoint{5.219150in}{0.708881in}}%
\pgfpathlineto{\pgfqpoint{5.222506in}{0.714488in}}%
\pgfpathlineto{\pgfqpoint{5.225861in}{0.716247in}}%
\pgfpathlineto{\pgfqpoint{5.229217in}{0.714221in}}%
\pgfpathlineto{\pgfqpoint{5.232573in}{0.707610in}}%
\pgfpathlineto{\pgfqpoint{5.235929in}{0.693684in}}%
\pgfpathlineto{\pgfqpoint{5.239285in}{0.660384in}}%
\pgfpathlineto{\pgfqpoint{5.242640in}{0.646472in}}%
\pgfpathlineto{\pgfqpoint{5.245996in}{0.693448in}}%
\pgfpathlineto{\pgfqpoint{5.249352in}{0.712587in}}%
\pgfpathlineto{\pgfqpoint{5.252708in}{0.722721in}}%
\pgfpathlineto{\pgfqpoint{5.256063in}{0.726835in}}%
\pgfpathlineto{\pgfqpoint{5.259419in}{0.724583in}}%
\pgfpathlineto{\pgfqpoint{5.262775in}{0.712685in}}%
\pgfpathlineto{\pgfqpoint{5.269487in}{0.660894in}}%
\pgfpathlineto{\pgfqpoint{5.272842in}{0.690063in}}%
\pgfpathlineto{\pgfqpoint{5.276198in}{0.701762in}}%
\pgfpathlineto{\pgfqpoint{5.279554in}{0.701057in}}%
\pgfpathlineto{\pgfqpoint{5.282910in}{0.686745in}}%
\pgfpathlineto{\pgfqpoint{5.286265in}{0.627928in}}%
\pgfpathlineto{\pgfqpoint{5.289621in}{0.681660in}}%
\pgfpathlineto{\pgfqpoint{5.292977in}{0.709913in}}%
\pgfpathlineto{\pgfqpoint{5.296333in}{0.720869in}}%
\pgfpathlineto{\pgfqpoint{5.299689in}{0.719022in}}%
\pgfpathlineto{\pgfqpoint{5.303044in}{0.701825in}}%
\pgfpathlineto{\pgfqpoint{5.306400in}{0.657394in}}%
\pgfpathlineto{\pgfqpoint{5.309756in}{0.670288in}}%
\pgfpathlineto{\pgfqpoint{5.313112in}{0.697381in}}%
\pgfpathlineto{\pgfqpoint{5.316467in}{0.704558in}}%
\pgfpathlineto{\pgfqpoint{5.319823in}{0.699967in}}%
\pgfpathlineto{\pgfqpoint{5.323179in}{0.678463in}}%
\pgfpathlineto{\pgfqpoint{5.326535in}{0.616955in}}%
\pgfpathlineto{\pgfqpoint{5.329891in}{0.694458in}}%
\pgfpathlineto{\pgfqpoint{5.333246in}{0.715747in}}%
\pgfpathlineto{\pgfqpoint{5.336602in}{0.722782in}}%
\pgfpathlineto{\pgfqpoint{5.339958in}{0.715665in}}%
\pgfpathlineto{\pgfqpoint{5.346669in}{0.659572in}}%
\pgfpathlineto{\pgfqpoint{5.350025in}{0.707751in}}%
\pgfpathlineto{\pgfqpoint{5.353381in}{0.722907in}}%
\pgfpathlineto{\pgfqpoint{5.356737in}{0.726931in}}%
\pgfpathlineto{\pgfqpoint{5.360093in}{0.722212in}}%
\pgfpathlineto{\pgfqpoint{5.363448in}{0.704858in}}%
\pgfpathlineto{\pgfqpoint{5.366804in}{0.625132in}}%
\pgfpathlineto{\pgfqpoint{5.370160in}{0.701262in}}%
\pgfpathlineto{\pgfqpoint{5.373516in}{0.728346in}}%
\pgfpathlineto{\pgfqpoint{5.376871in}{0.739193in}}%
\pgfpathlineto{\pgfqpoint{5.380227in}{0.738586in}}%
\pgfpathlineto{\pgfqpoint{5.383583in}{0.722018in}}%
\pgfpathlineto{\pgfqpoint{5.386939in}{0.664005in}}%
\pgfpathlineto{\pgfqpoint{5.390295in}{0.702087in}}%
\pgfpathlineto{\pgfqpoint{5.393650in}{0.728637in}}%
\pgfpathlineto{\pgfqpoint{5.397006in}{0.737191in}}%
\pgfpathlineto{\pgfqpoint{5.400362in}{0.735830in}}%
\pgfpathlineto{\pgfqpoint{5.403718in}{0.723248in}}%
\pgfpathlineto{\pgfqpoint{5.407074in}{0.679605in}}%
\pgfpathlineto{\pgfqpoint{5.410429in}{0.700796in}}%
\pgfpathlineto{\pgfqpoint{5.413785in}{0.736085in}}%
\pgfpathlineto{\pgfqpoint{5.417141in}{0.749867in}}%
\pgfpathlineto{\pgfqpoint{5.420497in}{0.751733in}}%
\pgfpathlineto{\pgfqpoint{5.423852in}{0.739589in}}%
\pgfpathlineto{\pgfqpoint{5.427208in}{0.711602in}}%
\pgfpathlineto{\pgfqpoint{5.430564in}{0.620446in}}%
\pgfpathlineto{\pgfqpoint{5.433920in}{0.705753in}}%
\pgfpathlineto{\pgfqpoint{5.437276in}{0.719280in}}%
\pgfpathlineto{\pgfqpoint{5.440631in}{0.716167in}}%
\pgfpathlineto{\pgfqpoint{5.443987in}{0.691985in}}%
\pgfpathlineto{\pgfqpoint{5.447343in}{0.662353in}}%
\pgfpathlineto{\pgfqpoint{5.450699in}{0.719928in}}%
\pgfpathlineto{\pgfqpoint{5.454054in}{0.738706in}}%
\pgfpathlineto{\pgfqpoint{5.457410in}{0.743334in}}%
\pgfpathlineto{\pgfqpoint{5.460766in}{0.733139in}}%
\pgfpathlineto{\pgfqpoint{5.464122in}{0.711234in}}%
\pgfpathlineto{\pgfqpoint{5.467478in}{0.660406in}}%
\pgfpathlineto{\pgfqpoint{5.470833in}{0.664893in}}%
\pgfpathlineto{\pgfqpoint{5.474189in}{0.688597in}}%
\pgfpathlineto{\pgfqpoint{5.477545in}{0.684357in}}%
\pgfpathlineto{\pgfqpoint{5.480901in}{0.639350in}}%
\pgfpathlineto{\pgfqpoint{5.487612in}{0.709518in}}%
\pgfpathlineto{\pgfqpoint{5.490968in}{0.722463in}}%
\pgfpathlineto{\pgfqpoint{5.494324in}{0.720930in}}%
\pgfpathlineto{\pgfqpoint{5.497680in}{0.704096in}}%
\pgfpathlineto{\pgfqpoint{5.501035in}{0.664243in}}%
\pgfpathlineto{\pgfqpoint{5.504391in}{0.649595in}}%
\pgfpathlineto{\pgfqpoint{5.507747in}{0.686025in}}%
\pgfpathlineto{\pgfqpoint{5.511103in}{0.692955in}}%
\pgfpathlineto{\pgfqpoint{5.514458in}{0.684576in}}%
\pgfpathlineto{\pgfqpoint{5.517814in}{0.644739in}}%
\pgfpathlineto{\pgfqpoint{5.521170in}{0.663638in}}%
\pgfpathlineto{\pgfqpoint{5.524526in}{0.700375in}}%
\pgfpathlineto{\pgfqpoint{5.527882in}{0.714166in}}%
\pgfpathlineto{\pgfqpoint{5.531237in}{0.714336in}}%
\pgfpathlineto{\pgfqpoint{5.534593in}{0.696651in}}%
\pgfpathlineto{\pgfqpoint{5.537949in}{0.625431in}}%
\pgfpathlineto{\pgfqpoint{5.541305in}{0.682987in}}%
\pgfpathlineto{\pgfqpoint{5.544660in}{0.707585in}}%
\pgfpathlineto{\pgfqpoint{5.548016in}{0.715614in}}%
\pgfpathlineto{\pgfqpoint{5.551372in}{0.714068in}}%
\pgfpathlineto{\pgfqpoint{5.554728in}{0.701519in}}%
\pgfpathlineto{\pgfqpoint{5.558084in}{0.658826in}}%
\pgfpathlineto{\pgfqpoint{5.561439in}{0.677707in}}%
\pgfpathlineto{\pgfqpoint{5.564795in}{0.713574in}}%
\pgfpathlineto{\pgfqpoint{5.568151in}{0.727588in}}%
\pgfpathlineto{\pgfqpoint{5.571507in}{0.729787in}}%
\pgfpathlineto{\pgfqpoint{5.574863in}{0.716405in}}%
\pgfpathlineto{\pgfqpoint{5.578218in}{0.663989in}}%
\pgfpathlineto{\pgfqpoint{5.584930in}{0.724882in}}%
\pgfpathlineto{\pgfqpoint{5.588286in}{0.735517in}}%
\pgfpathlineto{\pgfqpoint{5.591641in}{0.736832in}}%
\pgfpathlineto{\pgfqpoint{5.594997in}{0.729291in}}%
\pgfpathlineto{\pgfqpoint{5.598353in}{0.706043in}}%
\pgfpathlineto{\pgfqpoint{5.601709in}{0.641923in}}%
\pgfpathlineto{\pgfqpoint{5.605065in}{0.720720in}}%
\pgfpathlineto{\pgfqpoint{5.608420in}{0.742103in}}%
\pgfpathlineto{\pgfqpoint{5.611776in}{0.750100in}}%
\pgfpathlineto{\pgfqpoint{5.615132in}{0.746166in}}%
\pgfpathlineto{\pgfqpoint{5.618488in}{0.726904in}}%
\pgfpathlineto{\pgfqpoint{5.621843in}{0.665950in}}%
\pgfpathlineto{\pgfqpoint{5.625199in}{0.702917in}}%
\pgfpathlineto{\pgfqpoint{5.628555in}{0.727618in}}%
\pgfpathlineto{\pgfqpoint{5.631911in}{0.733423in}}%
\pgfpathlineto{\pgfqpoint{5.635267in}{0.727285in}}%
\pgfpathlineto{\pgfqpoint{5.638622in}{0.701747in}}%
\pgfpathlineto{\pgfqpoint{5.641978in}{0.669486in}}%
\pgfpathlineto{\pgfqpoint{5.645334in}{0.727553in}}%
\pgfpathlineto{\pgfqpoint{5.648690in}{0.746345in}}%
\pgfpathlineto{\pgfqpoint{5.652045in}{0.751542in}}%
\pgfpathlineto{\pgfqpoint{5.655401in}{0.742587in}}%
\pgfpathlineto{\pgfqpoint{5.658757in}{0.721873in}}%
\pgfpathlineto{\pgfqpoint{5.662113in}{0.675536in}}%
\pgfpathlineto{\pgfqpoint{5.665469in}{0.669236in}}%
\pgfpathlineto{\pgfqpoint{5.668824in}{0.697161in}}%
\pgfpathlineto{\pgfqpoint{5.672180in}{0.695328in}}%
\pgfpathlineto{\pgfqpoint{5.675536in}{0.660641in}}%
\pgfpathlineto{\pgfqpoint{5.678892in}{0.676092in}}%
\pgfpathlineto{\pgfqpoint{5.682247in}{0.715324in}}%
\pgfpathlineto{\pgfqpoint{5.685603in}{0.730131in}}%
\pgfpathlineto{\pgfqpoint{5.688959in}{0.730704in}}%
\pgfpathlineto{\pgfqpoint{5.692315in}{0.716438in}}%
\pgfpathlineto{\pgfqpoint{5.695671in}{0.688133in}}%
\pgfpathlineto{\pgfqpoint{5.699026in}{0.526715in}}%
\pgfpathlineto{\pgfqpoint{5.702382in}{0.672000in}}%
\pgfpathlineto{\pgfqpoint{5.705738in}{0.683360in}}%
\pgfpathlineto{\pgfqpoint{5.709094in}{0.672743in}}%
\pgfpathlineto{\pgfqpoint{5.712449in}{0.586753in}}%
\pgfpathlineto{\pgfqpoint{5.715805in}{0.677586in}}%
\pgfpathlineto{\pgfqpoint{5.719161in}{0.704820in}}%
\pgfpathlineto{\pgfqpoint{5.722517in}{0.714619in}}%
\pgfpathlineto{\pgfqpoint{5.725873in}{0.709358in}}%
\pgfpathlineto{\pgfqpoint{5.729228in}{0.686222in}}%
\pgfpathlineto{\pgfqpoint{5.732584in}{0.532239in}}%
\pgfpathlineto{\pgfqpoint{5.735940in}{0.680007in}}%
\pgfpathlineto{\pgfqpoint{5.739296in}{0.699100in}}%
\pgfpathlineto{\pgfqpoint{5.742652in}{0.703118in}}%
\pgfpathlineto{\pgfqpoint{5.746007in}{0.695936in}}%
\pgfpathlineto{\pgfqpoint{5.752719in}{0.641676in}}%
\pgfpathlineto{\pgfqpoint{5.756075in}{0.696690in}}%
\pgfpathlineto{\pgfqpoint{5.759430in}{0.715013in}}%
\pgfpathlineto{\pgfqpoint{5.762786in}{0.720073in}}%
\pgfpathlineto{\pgfqpoint{5.766142in}{0.710068in}}%
\pgfpathlineto{\pgfqpoint{5.769498in}{0.668521in}}%
\pgfpathlineto{\pgfqpoint{5.772854in}{0.679938in}}%
\pgfpathlineto{\pgfqpoint{5.776209in}{0.714791in}}%
\pgfpathlineto{\pgfqpoint{5.779565in}{0.727402in}}%
\pgfpathlineto{\pgfqpoint{5.782921in}{0.730411in}}%
\pgfpathlineto{\pgfqpoint{5.786277in}{0.725233in}}%
\pgfpathlineto{\pgfqpoint{5.789632in}{0.707842in}}%
\pgfpathlineto{\pgfqpoint{5.792988in}{0.632454in}}%
\pgfpathlineto{\pgfqpoint{5.796344in}{0.702662in}}%
\pgfpathlineto{\pgfqpoint{5.799700in}{0.730239in}}%
\pgfpathlineto{\pgfqpoint{5.803056in}{0.741499in}}%
\pgfpathlineto{\pgfqpoint{5.806411in}{0.741760in}}%
\pgfpathlineto{\pgfqpoint{5.809767in}{0.728741in}}%
\pgfpathlineto{\pgfqpoint{5.813123in}{0.700789in}}%
\pgfpathlineto{\pgfqpoint{5.816479in}{0.561599in}}%
\pgfpathlineto{\pgfqpoint{5.819834in}{0.690206in}}%
\pgfpathlineto{\pgfqpoint{5.823190in}{0.703175in}}%
\pgfpathlineto{\pgfqpoint{5.826546in}{0.697683in}}%
\pgfpathlineto{\pgfqpoint{5.829902in}{0.662872in}}%
\pgfpathlineto{\pgfqpoint{5.833258in}{0.672894in}}%
\pgfpathlineto{\pgfqpoint{5.836613in}{0.712831in}}%
\pgfpathlineto{\pgfqpoint{5.839969in}{0.727667in}}%
\pgfpathlineto{\pgfqpoint{5.843325in}{0.728595in}}%
\pgfpathlineto{\pgfqpoint{5.846681in}{0.714417in}}%
\pgfpathlineto{\pgfqpoint{5.850036in}{0.684127in}}%
\pgfpathlineto{\pgfqpoint{5.853392in}{0.605633in}}%
\pgfpathlineto{\pgfqpoint{5.856748in}{0.680151in}}%
\pgfpathlineto{\pgfqpoint{5.860104in}{0.690937in}}%
\pgfpathlineto{\pgfqpoint{5.863460in}{0.683693in}}%
\pgfpathlineto{\pgfqpoint{5.866815in}{0.642233in}}%
\pgfpathlineto{\pgfqpoint{5.870171in}{0.666855in}}%
\pgfpathlineto{\pgfqpoint{5.873527in}{0.702510in}}%
\pgfpathlineto{\pgfqpoint{5.876883in}{0.715927in}}%
\pgfpathlineto{\pgfqpoint{5.880238in}{0.715351in}}%
\pgfpathlineto{\pgfqpoint{5.883594in}{0.697528in}}%
\pgfpathlineto{\pgfqpoint{5.886950in}{0.638297in}}%
\pgfpathlineto{\pgfqpoint{5.890306in}{0.675025in}}%
\pgfpathlineto{\pgfqpoint{5.893662in}{0.701021in}}%
\pgfpathlineto{\pgfqpoint{5.897017in}{0.708424in}}%
\pgfpathlineto{\pgfqpoint{5.900373in}{0.705076in}}%
\pgfpathlineto{\pgfqpoint{5.903729in}{0.687721in}}%
\pgfpathlineto{\pgfqpoint{5.907085in}{0.586617in}}%
\pgfpathlineto{\pgfqpoint{5.910441in}{0.688965in}}%
\pgfpathlineto{\pgfqpoint{5.913796in}{0.714613in}}%
\pgfpathlineto{\pgfqpoint{5.917152in}{0.724244in}}%
\pgfpathlineto{\pgfqpoint{5.920508in}{0.720968in}}%
\pgfpathlineto{\pgfqpoint{5.923864in}{0.697899in}}%
\pgfpathlineto{\pgfqpoint{5.927219in}{0.621657in}}%
\pgfpathlineto{\pgfqpoint{5.930575in}{0.705321in}}%
\pgfpathlineto{\pgfqpoint{5.933931in}{0.724348in}}%
\pgfpathlineto{\pgfqpoint{5.937287in}{0.730594in}}%
\pgfpathlineto{\pgfqpoint{5.940643in}{0.728211in}}%
\pgfpathlineto{\pgfqpoint{5.943998in}{0.715256in}}%
\pgfpathlineto{\pgfqpoint{5.947354in}{0.672774in}}%
\pgfpathlineto{\pgfqpoint{5.950710in}{0.690399in}}%
\pgfpathlineto{\pgfqpoint{5.954066in}{0.726549in}}%
\pgfpathlineto{\pgfqpoint{5.957421in}{0.740828in}}%
\pgfpathlineto{\pgfqpoint{5.960777in}{0.743680in}}%
\pgfpathlineto{\pgfqpoint{5.964133in}{0.732500in}}%
\pgfpathlineto{\pgfqpoint{5.967489in}{0.698229in}}%
\pgfpathlineto{\pgfqpoint{5.970845in}{0.678289in}}%
\pgfpathlineto{\pgfqpoint{5.974200in}{0.722006in}}%
\pgfpathlineto{\pgfqpoint{5.977556in}{0.734575in}}%
\pgfpathlineto{\pgfqpoint{5.980912in}{0.735485in}}%
\pgfpathlineto{\pgfqpoint{5.984268in}{0.725308in}}%
\pgfpathlineto{\pgfqpoint{5.987623in}{0.690208in}}%
\pgfpathlineto{\pgfqpoint{5.990979in}{0.691628in}}%
\pgfpathlineto{\pgfqpoint{5.994335in}{0.733145in}}%
\pgfpathlineto{\pgfqpoint{5.997691in}{0.748642in}}%
\pgfpathlineto{\pgfqpoint{6.001047in}{0.751699in}}%
\pgfpathlineto{\pgfqpoint{6.004402in}{0.740738in}}%
\pgfpathlineto{\pgfqpoint{6.007758in}{0.716042in}}%
\pgfpathlineto{\pgfqpoint{6.011114in}{0.632432in}}%
\pgfpathlineto{\pgfqpoint{6.014470in}{0.694134in}}%
\pgfpathlineto{\pgfqpoint{6.017825in}{0.711163in}}%
\pgfpathlineto{\pgfqpoint{6.021181in}{0.708209in}}%
\pgfpathlineto{\pgfqpoint{6.024537in}{0.680381in}}%
\pgfpathlineto{\pgfqpoint{6.027893in}{0.670645in}}%
\pgfpathlineto{\pgfqpoint{6.031249in}{0.718467in}}%
\pgfpathlineto{\pgfqpoint{6.034604in}{0.735414in}}%
\pgfpathlineto{\pgfqpoint{6.037960in}{0.738312in}}%
\pgfpathlineto{\pgfqpoint{6.041316in}{0.726232in}}%
\pgfpathlineto{\pgfqpoint{6.044672in}{0.701891in}}%
\pgfpathlineto{\pgfqpoint{6.048027in}{0.635522in}}%
\pgfpathlineto{\pgfqpoint{6.051383in}{0.668798in}}%
\pgfpathlineto{\pgfqpoint{6.054739in}{0.686205in}}%
\pgfpathlineto{\pgfqpoint{6.058095in}{0.678844in}}%
\pgfpathlineto{\pgfqpoint{6.061451in}{0.618305in}}%
\pgfpathlineto{\pgfqpoint{6.064806in}{0.677842in}}%
\pgfpathlineto{\pgfqpoint{6.068162in}{0.707958in}}%
\pgfpathlineto{\pgfqpoint{6.071518in}{0.719210in}}%
\pgfpathlineto{\pgfqpoint{6.074874in}{0.715668in}}%
\pgfpathlineto{\pgfqpoint{6.078230in}{0.696283in}}%
\pgfpathlineto{\pgfqpoint{6.081585in}{0.640507in}}%
\pgfpathlineto{\pgfqpoint{6.088297in}{0.691933in}}%
\pgfpathlineto{\pgfqpoint{6.091653in}{0.696948in}}%
\pgfpathlineto{\pgfqpoint{6.095008in}{0.688575in}}%
\pgfpathlineto{\pgfqpoint{6.098364in}{0.653477in}}%
\pgfpathlineto{\pgfqpoint{6.101720in}{0.658856in}}%
\pgfpathlineto{\pgfqpoint{6.105076in}{0.699486in}}%
\pgfpathlineto{\pgfqpoint{6.108432in}{0.714524in}}%
\pgfpathlineto{\pgfqpoint{6.111787in}{0.716305in}}%
\pgfpathlineto{\pgfqpoint{6.115143in}{0.700835in}}%
\pgfpathlineto{\pgfqpoint{6.118499in}{0.639164in}}%
\pgfpathlineto{\pgfqpoint{6.121855in}{0.684511in}}%
\pgfpathlineto{\pgfqpoint{6.125210in}{0.711301in}}%
\pgfpathlineto{\pgfqpoint{6.128566in}{0.720692in}}%
\pgfpathlineto{\pgfqpoint{6.131922in}{0.720785in}}%
\pgfpathlineto{\pgfqpoint{6.135278in}{0.711306in}}%
\pgfpathlineto{\pgfqpoint{6.138634in}{0.682036in}}%
\pgfpathlineto{\pgfqpoint{6.141989in}{0.660222in}}%
\pgfpathlineto{\pgfqpoint{6.145345in}{0.711714in}}%
\pgfpathlineto{\pgfqpoint{6.148701in}{0.729661in}}%
\pgfpathlineto{\pgfqpoint{6.152057in}{0.735307in}}%
\pgfpathlineto{\pgfqpoint{6.155412in}{0.727790in}}%
\pgfpathlineto{\pgfqpoint{6.162124in}{0.671601in}}%
\pgfpathlineto{\pgfqpoint{6.165480in}{0.720471in}}%
\pgfpathlineto{\pgfqpoint{6.168836in}{0.735526in}}%
\pgfpathlineto{\pgfqpoint{6.172191in}{0.739450in}}%
\pgfpathlineto{\pgfqpoint{6.175547in}{0.734587in}}%
\pgfpathlineto{\pgfqpoint{6.178903in}{0.716896in}}%
\pgfpathlineto{\pgfqpoint{6.182259in}{0.631209in}}%
\pgfpathlineto{\pgfqpoint{6.185614in}{0.714637in}}%
\pgfpathlineto{\pgfqpoint{6.188970in}{0.741228in}}%
\pgfpathlineto{\pgfqpoint{6.192326in}{0.751816in}}%
\pgfpathlineto{\pgfqpoint{6.195682in}{0.750856in}}%
\pgfpathlineto{\pgfqpoint{6.199038in}{0.735593in}}%
\pgfpathlineto{\pgfqpoint{6.202393in}{0.698154in}}%
\pgfpathlineto{\pgfqpoint{6.205749in}{0.678519in}}%
\pgfpathlineto{\pgfqpoint{6.209105in}{0.718554in}}%
\pgfpathlineto{\pgfqpoint{6.212461in}{0.727547in}}%
\pgfpathlineto{\pgfqpoint{6.215816in}{0.722646in}}%
\pgfpathlineto{\pgfqpoint{6.222528in}{0.668942in}}%
\pgfpathlineto{\pgfqpoint{6.225884in}{0.725218in}}%
\pgfpathlineto{\pgfqpoint{6.229240in}{0.743713in}}%
\pgfpathlineto{\pgfqpoint{6.232595in}{0.748351in}}%
\pgfpathlineto{\pgfqpoint{6.235951in}{0.738470in}}%
\pgfpathlineto{\pgfqpoint{6.239307in}{0.717311in}}%
\pgfpathlineto{\pgfqpoint{6.242663in}{0.670993in}}%
\pgfpathlineto{\pgfqpoint{6.246019in}{0.662162in}}%
\pgfpathlineto{\pgfqpoint{6.249374in}{0.689543in}}%
\pgfpathlineto{\pgfqpoint{6.252730in}{0.685290in}}%
\pgfpathlineto{\pgfqpoint{6.256086in}{0.633870in}}%
\pgfpathlineto{\pgfqpoint{6.259442in}{0.681755in}}%
\pgfpathlineto{\pgfqpoint{6.262797in}{0.713887in}}%
\pgfpathlineto{\pgfqpoint{6.266153in}{0.726122in}}%
\pgfpathlineto{\pgfqpoint{6.269509in}{0.723667in}}%
\pgfpathlineto{\pgfqpoint{6.272865in}{0.706904in}}%
\pgfpathlineto{\pgfqpoint{6.279576in}{0.634426in}}%
\pgfpathlineto{\pgfqpoint{6.282932in}{0.679242in}}%
\pgfpathlineto{\pgfqpoint{6.286288in}{0.685616in}}%
\pgfpathlineto{\pgfqpoint{6.289644in}{0.673071in}}%
\pgfpathlineto{\pgfqpoint{6.292999in}{0.582798in}}%
\pgfpathlineto{\pgfqpoint{6.296355in}{0.677045in}}%
\pgfpathlineto{\pgfqpoint{6.299711in}{0.703741in}}%
\pgfpathlineto{\pgfqpoint{6.303067in}{0.713254in}}%
\pgfpathlineto{\pgfqpoint{6.306423in}{0.707852in}}%
\pgfpathlineto{\pgfqpoint{6.309778in}{0.682874in}}%
\pgfpathlineto{\pgfqpoint{6.313134in}{0.604522in}}%
\pgfpathlineto{\pgfqpoint{6.316490in}{0.687403in}}%
\pgfpathlineto{\pgfqpoint{6.319846in}{0.704977in}}%
\pgfpathlineto{\pgfqpoint{6.323201in}{0.709250in}}%
\pgfpathlineto{\pgfqpoint{6.326557in}{0.703561in}}%
\pgfpathlineto{\pgfqpoint{6.329913in}{0.682380in}}%
\pgfpathlineto{\pgfqpoint{6.333269in}{0.599618in}}%
\pgfpathlineto{\pgfqpoint{6.336625in}{0.693982in}}%
\pgfpathlineto{\pgfqpoint{6.339980in}{0.716250in}}%
\pgfpathlineto{\pgfqpoint{6.343336in}{0.724233in}}%
\pgfpathlineto{\pgfqpoint{6.346692in}{0.719280in}}%
\pgfpathlineto{\pgfqpoint{6.350048in}{0.700184in}}%
\pgfpathlineto{\pgfqpoint{6.353403in}{0.650548in}}%
\pgfpathlineto{\pgfqpoint{6.356759in}{0.661834in}}%
\pgfpathlineto{\pgfqpoint{6.360115in}{0.689542in}}%
\pgfpathlineto{\pgfqpoint{6.363471in}{0.694032in}}%
\pgfpathlineto{\pgfqpoint{6.366827in}{0.683282in}}%
\pgfpathlineto{\pgfqpoint{6.370182in}{0.632438in}}%
\pgfpathlineto{\pgfqpoint{6.373538in}{0.671735in}}%
\pgfpathlineto{\pgfqpoint{6.376894in}{0.703749in}}%
\pgfpathlineto{\pgfqpoint{6.380250in}{0.715798in}}%
\pgfpathlineto{\pgfqpoint{6.383605in}{0.714071in}}%
\pgfpathlineto{\pgfqpoint{6.383605in}{0.714071in}}%
\pgfusepath{stroke}%
\end{pgfscope}%
\begin{pgfscope}%
\pgfpathrectangle{\pgfqpoint{4.621821in}{0.320625in}}{\pgfqpoint{1.845679in}{1.164375in}}%
\pgfusepath{clip}%
\pgfsetrectcap%
\pgfsetroundjoin%
\pgfsetlinewidth{0.602250pt}%
\definecolor{currentstroke}{rgb}{0.478431,0.243137,0.615686}%
\pgfsetstrokecolor{currentstroke}%
\pgfsetstrokeopacity{0.120000}%
\pgfsetdash{}{0pt}%
\pgfpathmoveto{\pgfqpoint{4.705715in}{1.468685in}}%
\pgfpathlineto{\pgfqpoint{4.715783in}{1.450011in}}%
\pgfpathlineto{\pgfqpoint{4.725850in}{1.428817in}}%
\pgfpathlineto{\pgfqpoint{4.739273in}{1.397199in}}%
\pgfpathlineto{\pgfqpoint{4.776187in}{1.308109in}}%
\pgfpathlineto{\pgfqpoint{4.796322in}{1.262917in}}%
\pgfpathlineto{\pgfqpoint{4.819812in}{1.213174in}}%
\pgfpathlineto{\pgfqpoint{4.846658in}{1.159034in}}%
\pgfpathlineto{\pgfqpoint{4.880216in}{1.093847in}}%
\pgfpathlineto{\pgfqpoint{4.923841in}{1.011522in}}%
\pgfpathlineto{\pgfqpoint{4.984245in}{0.900090in}}%
\pgfpathlineto{\pgfqpoint{5.044649in}{0.790481in}}%
\pgfpathlineto{\pgfqpoint{5.058072in}{0.767949in}}%
\pgfpathlineto{\pgfqpoint{5.078207in}{0.728137in}}%
\pgfpathlineto{\pgfqpoint{5.081563in}{0.723244in}}%
\pgfpathlineto{\pgfqpoint{5.088274in}{0.716381in}}%
\pgfpathlineto{\pgfqpoint{5.091630in}{0.713581in}}%
\pgfpathlineto{\pgfqpoint{5.094986in}{0.709939in}}%
\pgfpathlineto{\pgfqpoint{5.098342in}{0.704065in}}%
\pgfpathlineto{\pgfqpoint{5.101698in}{0.696468in}}%
\pgfpathlineto{\pgfqpoint{5.105053in}{0.687267in}}%
\pgfpathlineto{\pgfqpoint{5.111765in}{0.663927in}}%
\pgfpathlineto{\pgfqpoint{5.118476in}{0.638258in}}%
\pgfpathlineto{\pgfqpoint{5.121832in}{0.628839in}}%
\pgfpathlineto{\pgfqpoint{5.125188in}{0.622344in}}%
\pgfpathlineto{\pgfqpoint{5.128544in}{0.619982in}}%
\pgfpathlineto{\pgfqpoint{5.131900in}{0.626530in}}%
\pgfpathlineto{\pgfqpoint{5.138611in}{0.650875in}}%
\pgfpathlineto{\pgfqpoint{5.141967in}{0.659697in}}%
\pgfpathlineto{\pgfqpoint{5.145323in}{0.664980in}}%
\pgfpathlineto{\pgfqpoint{5.148678in}{0.666102in}}%
\pgfpathlineto{\pgfqpoint{5.152034in}{0.661270in}}%
\pgfpathlineto{\pgfqpoint{5.155390in}{0.649770in}}%
\pgfpathlineto{\pgfqpoint{5.158746in}{0.651698in}}%
\pgfpathlineto{\pgfqpoint{5.162102in}{0.671951in}}%
\pgfpathlineto{\pgfqpoint{5.165457in}{0.688209in}}%
\pgfpathlineto{\pgfqpoint{5.168813in}{0.699284in}}%
\pgfpathlineto{\pgfqpoint{5.172169in}{0.706630in}}%
\pgfpathlineto{\pgfqpoint{5.175525in}{0.711125in}}%
\pgfpathlineto{\pgfqpoint{5.178880in}{0.713198in}}%
\pgfpathlineto{\pgfqpoint{5.182236in}{0.712970in}}%
\pgfpathlineto{\pgfqpoint{5.185592in}{0.710271in}}%
\pgfpathlineto{\pgfqpoint{5.188948in}{0.704526in}}%
\pgfpathlineto{\pgfqpoint{5.192304in}{0.694293in}}%
\pgfpathlineto{\pgfqpoint{5.195659in}{0.675384in}}%
\pgfpathlineto{\pgfqpoint{5.199015in}{0.626458in}}%
\pgfpathlineto{\pgfqpoint{5.205727in}{0.690366in}}%
\pgfpathlineto{\pgfqpoint{5.209082in}{0.706673in}}%
\pgfpathlineto{\pgfqpoint{5.212438in}{0.716664in}}%
\pgfpathlineto{\pgfqpoint{5.215794in}{0.722585in}}%
\pgfpathlineto{\pgfqpoint{5.219150in}{0.725024in}}%
\pgfpathlineto{\pgfqpoint{5.222506in}{0.723610in}}%
\pgfpathlineto{\pgfqpoint{5.225861in}{0.716641in}}%
\pgfpathlineto{\pgfqpoint{5.229217in}{0.703241in}}%
\pgfpathlineto{\pgfqpoint{5.232573in}{0.685852in}}%
\pgfpathlineto{\pgfqpoint{5.235929in}{0.679347in}}%
\pgfpathlineto{\pgfqpoint{5.239285in}{0.685808in}}%
\pgfpathlineto{\pgfqpoint{5.242640in}{0.687947in}}%
\pgfpathlineto{\pgfqpoint{5.245996in}{0.680692in}}%
\pgfpathlineto{\pgfqpoint{5.249352in}{0.658170in}}%
\pgfpathlineto{\pgfqpoint{5.252708in}{0.654138in}}%
\pgfpathlineto{\pgfqpoint{5.256063in}{0.686815in}}%
\pgfpathlineto{\pgfqpoint{5.259419in}{0.702526in}}%
\pgfpathlineto{\pgfqpoint{5.262775in}{0.706444in}}%
\pgfpathlineto{\pgfqpoint{5.266131in}{0.696570in}}%
\pgfpathlineto{\pgfqpoint{5.269487in}{0.672379in}}%
\pgfpathlineto{\pgfqpoint{5.272842in}{0.634090in}}%
\pgfpathlineto{\pgfqpoint{5.276198in}{0.667352in}}%
\pgfpathlineto{\pgfqpoint{5.279554in}{0.681869in}}%
\pgfpathlineto{\pgfqpoint{5.282910in}{0.682826in}}%
\pgfpathlineto{\pgfqpoint{5.286265in}{0.670758in}}%
\pgfpathlineto{\pgfqpoint{5.289621in}{0.623507in}}%
\pgfpathlineto{\pgfqpoint{5.296333in}{0.688965in}}%
\pgfpathlineto{\pgfqpoint{5.299689in}{0.701595in}}%
\pgfpathlineto{\pgfqpoint{5.303044in}{0.701716in}}%
\pgfpathlineto{\pgfqpoint{5.306400in}{0.681329in}}%
\pgfpathlineto{\pgfqpoint{5.309756in}{0.634720in}}%
\pgfpathlineto{\pgfqpoint{5.313112in}{0.698646in}}%
\pgfpathlineto{\pgfqpoint{5.316467in}{0.718997in}}%
\pgfpathlineto{\pgfqpoint{5.319823in}{0.727924in}}%
\pgfpathlineto{\pgfqpoint{5.323179in}{0.729864in}}%
\pgfpathlineto{\pgfqpoint{5.326535in}{0.725239in}}%
\pgfpathlineto{\pgfqpoint{5.329891in}{0.711223in}}%
\pgfpathlineto{\pgfqpoint{5.333246in}{0.669915in}}%
\pgfpathlineto{\pgfqpoint{5.336602in}{0.682536in}}%
\pgfpathlineto{\pgfqpoint{5.339958in}{0.720112in}}%
\pgfpathlineto{\pgfqpoint{5.343314in}{0.735600in}}%
\pgfpathlineto{\pgfqpoint{5.346669in}{0.741023in}}%
\pgfpathlineto{\pgfqpoint{5.350025in}{0.735921in}}%
\pgfpathlineto{\pgfqpoint{5.353381in}{0.718968in}}%
\pgfpathlineto{\pgfqpoint{5.360093in}{0.652285in}}%
\pgfpathlineto{\pgfqpoint{5.363448in}{0.692456in}}%
\pgfpathlineto{\pgfqpoint{5.366804in}{0.698806in}}%
\pgfpathlineto{\pgfqpoint{5.370160in}{0.687098in}}%
\pgfpathlineto{\pgfqpoint{5.373516in}{0.616452in}}%
\pgfpathlineto{\pgfqpoint{5.376871in}{0.687156in}}%
\pgfpathlineto{\pgfqpoint{5.380227in}{0.715315in}}%
\pgfpathlineto{\pgfqpoint{5.383583in}{0.725816in}}%
\pgfpathlineto{\pgfqpoint{5.386939in}{0.722159in}}%
\pgfpathlineto{\pgfqpoint{5.390295in}{0.704952in}}%
\pgfpathlineto{\pgfqpoint{5.393650in}{0.667763in}}%
\pgfpathlineto{\pgfqpoint{5.397006in}{0.639784in}}%
\pgfpathlineto{\pgfqpoint{5.400362in}{0.679633in}}%
\pgfpathlineto{\pgfqpoint{5.403718in}{0.685707in}}%
\pgfpathlineto{\pgfqpoint{5.407074in}{0.673752in}}%
\pgfpathlineto{\pgfqpoint{5.410429in}{0.600063in}}%
\pgfpathlineto{\pgfqpoint{5.413785in}{0.674279in}}%
\pgfpathlineto{\pgfqpoint{5.417141in}{0.702217in}}%
\pgfpathlineto{\pgfqpoint{5.420497in}{0.712581in}}%
\pgfpathlineto{\pgfqpoint{5.423852in}{0.708684in}}%
\pgfpathlineto{\pgfqpoint{5.427208in}{0.686364in}}%
\pgfpathlineto{\pgfqpoint{5.430564in}{0.574386in}}%
\pgfpathlineto{\pgfqpoint{5.433920in}{0.683595in}}%
\pgfpathlineto{\pgfqpoint{5.437276in}{0.703500in}}%
\pgfpathlineto{\pgfqpoint{5.440631in}{0.709162in}}%
\pgfpathlineto{\pgfqpoint{5.443987in}{0.705160in}}%
\pgfpathlineto{\pgfqpoint{5.447343in}{0.687933in}}%
\pgfpathlineto{\pgfqpoint{5.450699in}{0.601768in}}%
\pgfpathlineto{\pgfqpoint{5.454054in}{0.686319in}}%
\pgfpathlineto{\pgfqpoint{5.457410in}{0.712884in}}%
\pgfpathlineto{\pgfqpoint{5.460766in}{0.723301in}}%
\pgfpathlineto{\pgfqpoint{5.464122in}{0.721739in}}%
\pgfpathlineto{\pgfqpoint{5.467478in}{0.701513in}}%
\pgfpathlineto{\pgfqpoint{5.470833in}{0.592483in}}%
\pgfpathlineto{\pgfqpoint{5.474189in}{0.705149in}}%
\pgfpathlineto{\pgfqpoint{5.477545in}{0.726334in}}%
\pgfpathlineto{\pgfqpoint{5.480901in}{0.734227in}}%
\pgfpathlineto{\pgfqpoint{5.484256in}{0.733951in}}%
\pgfpathlineto{\pgfqpoint{5.487612in}{0.724986in}}%
\pgfpathlineto{\pgfqpoint{5.490968in}{0.698966in}}%
\pgfpathlineto{\pgfqpoint{5.494324in}{0.656762in}}%
\pgfpathlineto{\pgfqpoint{5.497680in}{0.719938in}}%
\pgfpathlineto{\pgfqpoint{5.501035in}{0.739887in}}%
\pgfpathlineto{\pgfqpoint{5.504391in}{0.747375in}}%
\pgfpathlineto{\pgfqpoint{5.507747in}{0.743429in}}%
\pgfpathlineto{\pgfqpoint{5.511103in}{0.727445in}}%
\pgfpathlineto{\pgfqpoint{5.514458in}{0.695569in}}%
\pgfpathlineto{\pgfqpoint{5.517814in}{0.622123in}}%
\pgfpathlineto{\pgfqpoint{5.521170in}{0.689784in}}%
\pgfpathlineto{\pgfqpoint{5.524526in}{0.697714in}}%
\pgfpathlineto{\pgfqpoint{5.527882in}{0.684091in}}%
\pgfpathlineto{\pgfqpoint{5.531237in}{0.554433in}}%
\pgfpathlineto{\pgfqpoint{5.534593in}{0.696120in}}%
\pgfpathlineto{\pgfqpoint{5.537949in}{0.720763in}}%
\pgfpathlineto{\pgfqpoint{5.541305in}{0.729324in}}%
\pgfpathlineto{\pgfqpoint{5.544660in}{0.723081in}}%
\pgfpathlineto{\pgfqpoint{5.548016in}{0.704438in}}%
\pgfpathlineto{\pgfqpoint{5.551372in}{0.664878in}}%
\pgfpathlineto{\pgfqpoint{5.554728in}{0.638922in}}%
\pgfpathlineto{\pgfqpoint{5.558084in}{0.676375in}}%
\pgfpathlineto{\pgfqpoint{5.561439in}{0.679194in}}%
\pgfpathlineto{\pgfqpoint{5.564795in}{0.658446in}}%
\pgfpathlineto{\pgfqpoint{5.568151in}{0.625060in}}%
\pgfpathlineto{\pgfqpoint{5.571507in}{0.686672in}}%
\pgfpathlineto{\pgfqpoint{5.574863in}{0.706548in}}%
\pgfpathlineto{\pgfqpoint{5.578218in}{0.711940in}}%
\pgfpathlineto{\pgfqpoint{5.581574in}{0.701405in}}%
\pgfpathlineto{\pgfqpoint{5.584930in}{0.668720in}}%
\pgfpathlineto{\pgfqpoint{5.588286in}{0.641492in}}%
\pgfpathlineto{\pgfqpoint{5.591641in}{0.688750in}}%
\pgfpathlineto{\pgfqpoint{5.594997in}{0.702012in}}%
\pgfpathlineto{\pgfqpoint{5.598353in}{0.703441in}}%
\pgfpathlineto{\pgfqpoint{5.601709in}{0.694056in}}%
\pgfpathlineto{\pgfqpoint{5.605065in}{0.662034in}}%
\pgfpathlineto{\pgfqpoint{5.608420in}{0.653818in}}%
\pgfpathlineto{\pgfqpoint{5.611776in}{0.699224in}}%
\pgfpathlineto{\pgfqpoint{5.615132in}{0.715784in}}%
\pgfpathlineto{\pgfqpoint{5.618488in}{0.719899in}}%
\pgfpathlineto{\pgfqpoint{5.621843in}{0.708984in}}%
\pgfpathlineto{\pgfqpoint{5.625199in}{0.662670in}}%
\pgfpathlineto{\pgfqpoint{5.628555in}{0.685276in}}%
\pgfpathlineto{\pgfqpoint{5.631911in}{0.717650in}}%
\pgfpathlineto{\pgfqpoint{5.635267in}{0.729929in}}%
\pgfpathlineto{\pgfqpoint{5.638622in}{0.733112in}}%
\pgfpathlineto{\pgfqpoint{5.641978in}{0.728585in}}%
\pgfpathlineto{\pgfqpoint{5.645334in}{0.713087in}}%
\pgfpathlineto{\pgfqpoint{5.648690in}{0.658668in}}%
\pgfpathlineto{\pgfqpoint{5.652045in}{0.698184in}}%
\pgfpathlineto{\pgfqpoint{5.655401in}{0.729591in}}%
\pgfpathlineto{\pgfqpoint{5.658757in}{0.742602in}}%
\pgfpathlineto{\pgfqpoint{5.662113in}{0.745017in}}%
\pgfpathlineto{\pgfqpoint{5.665469in}{0.734899in}}%
\pgfpathlineto{\pgfqpoint{5.668824in}{0.712630in}}%
\pgfpathlineto{\pgfqpoint{5.672180in}{0.655553in}}%
\pgfpathlineto{\pgfqpoint{5.678892in}{0.697977in}}%
\pgfpathlineto{\pgfqpoint{5.682247in}{0.696452in}}%
\pgfpathlineto{\pgfqpoint{5.685603in}{0.669378in}}%
\pgfpathlineto{\pgfqpoint{5.688959in}{0.659736in}}%
\pgfpathlineto{\pgfqpoint{5.692315in}{0.708003in}}%
\pgfpathlineto{\pgfqpoint{5.695671in}{0.725273in}}%
\pgfpathlineto{\pgfqpoint{5.699026in}{0.728606in}}%
\pgfpathlineto{\pgfqpoint{5.702382in}{0.716905in}}%
\pgfpathlineto{\pgfqpoint{5.705738in}{0.692226in}}%
\pgfpathlineto{\pgfqpoint{5.709094in}{0.616535in}}%
\pgfpathlineto{\pgfqpoint{5.712449in}{0.666125in}}%
\pgfpathlineto{\pgfqpoint{5.715805in}{0.683304in}}%
\pgfpathlineto{\pgfqpoint{5.719161in}{0.679038in}}%
\pgfpathlineto{\pgfqpoint{5.722517in}{0.644345in}}%
\pgfpathlineto{\pgfqpoint{5.725873in}{0.656200in}}%
\pgfpathlineto{\pgfqpoint{5.729228in}{0.695999in}}%
\pgfpathlineto{\pgfqpoint{5.732584in}{0.711002in}}%
\pgfpathlineto{\pgfqpoint{5.735940in}{0.712226in}}%
\pgfpathlineto{\pgfqpoint{5.739296in}{0.696597in}}%
\pgfpathlineto{\pgfqpoint{5.742652in}{0.648402in}}%
\pgfpathlineto{\pgfqpoint{5.746007in}{0.665614in}}%
\pgfpathlineto{\pgfqpoint{5.749363in}{0.695956in}}%
\pgfpathlineto{\pgfqpoint{5.752719in}{0.705178in}}%
\pgfpathlineto{\pgfqpoint{5.756075in}{0.703746in}}%
\pgfpathlineto{\pgfqpoint{5.759430in}{0.690286in}}%
\pgfpathlineto{\pgfqpoint{5.762786in}{0.639602in}}%
\pgfpathlineto{\pgfqpoint{5.769498in}{0.707089in}}%
\pgfpathlineto{\pgfqpoint{5.772854in}{0.719777in}}%
\pgfpathlineto{\pgfqpoint{5.776209in}{0.720412in}}%
\pgfpathlineto{\pgfqpoint{5.779565in}{0.703799in}}%
\pgfpathlineto{\pgfqpoint{5.782921in}{0.619361in}}%
\pgfpathlineto{\pgfqpoint{5.786277in}{0.697341in}}%
\pgfpathlineto{\pgfqpoint{5.789632in}{0.721893in}}%
\pgfpathlineto{\pgfqpoint{5.792988in}{0.731228in}}%
\pgfpathlineto{\pgfqpoint{5.796344in}{0.732108in}}%
\pgfpathlineto{\pgfqpoint{5.799700in}{0.724664in}}%
\pgfpathlineto{\pgfqpoint{5.803056in}{0.702641in}}%
\pgfpathlineto{\pgfqpoint{5.806411in}{0.614683in}}%
\pgfpathlineto{\pgfqpoint{5.809767in}{0.712762in}}%
\pgfpathlineto{\pgfqpoint{5.813123in}{0.735451in}}%
\pgfpathlineto{\pgfqpoint{5.816479in}{0.744511in}}%
\pgfpathlineto{\pgfqpoint{5.819834in}{0.742582in}}%
\pgfpathlineto{\pgfqpoint{5.823190in}{0.727936in}}%
\pgfpathlineto{\pgfqpoint{5.826546in}{0.698184in}}%
\pgfpathlineto{\pgfqpoint{5.829902in}{0.599040in}}%
\pgfpathlineto{\pgfqpoint{5.833258in}{0.689611in}}%
\pgfpathlineto{\pgfqpoint{5.836613in}{0.700382in}}%
\pgfpathlineto{\pgfqpoint{5.839969in}{0.691702in}}%
\pgfpathlineto{\pgfqpoint{5.843325in}{0.638341in}}%
\pgfpathlineto{\pgfqpoint{5.846681in}{0.684894in}}%
\pgfpathlineto{\pgfqpoint{5.850036in}{0.716394in}}%
\pgfpathlineto{\pgfqpoint{5.853392in}{0.728365in}}%
\pgfpathlineto{\pgfqpoint{5.856748in}{0.726259in}}%
\pgfpathlineto{\pgfqpoint{5.860104in}{0.710290in}}%
\pgfpathlineto{\pgfqpoint{5.863460in}{0.677339in}}%
\pgfpathlineto{\pgfqpoint{5.866815in}{0.618888in}}%
\pgfpathlineto{\pgfqpoint{5.870171in}{0.677154in}}%
\pgfpathlineto{\pgfqpoint{5.873527in}{0.685245in}}%
\pgfpathlineto{\pgfqpoint{5.876883in}{0.673964in}}%
\pgfpathlineto{\pgfqpoint{5.880238in}{0.597912in}}%
\pgfpathlineto{\pgfqpoint{5.883594in}{0.675572in}}%
\pgfpathlineto{\pgfqpoint{5.886950in}{0.703503in}}%
\pgfpathlineto{\pgfqpoint{5.890306in}{0.713830in}}%
\pgfpathlineto{\pgfqpoint{5.893662in}{0.709709in}}%
\pgfpathlineto{\pgfqpoint{5.897017in}{0.687840in}}%
\pgfpathlineto{\pgfqpoint{5.900373in}{0.574891in}}%
\pgfpathlineto{\pgfqpoint{5.903729in}{0.680085in}}%
\pgfpathlineto{\pgfqpoint{5.907085in}{0.700637in}}%
\pgfpathlineto{\pgfqpoint{5.910441in}{0.706005in}}%
\pgfpathlineto{\pgfqpoint{5.913796in}{0.701044in}}%
\pgfpathlineto{\pgfqpoint{5.917152in}{0.680893in}}%
\pgfpathlineto{\pgfqpoint{5.920508in}{0.579137in}}%
\pgfpathlineto{\pgfqpoint{5.923864in}{0.690332in}}%
\pgfpathlineto{\pgfqpoint{5.927219in}{0.713364in}}%
\pgfpathlineto{\pgfqpoint{5.930575in}{0.721791in}}%
\pgfpathlineto{\pgfqpoint{5.933931in}{0.717303in}}%
\pgfpathlineto{\pgfqpoint{5.937287in}{0.691378in}}%
\pgfpathlineto{\pgfqpoint{5.940643in}{0.647165in}}%
\pgfpathlineto{\pgfqpoint{5.943998in}{0.708117in}}%
\pgfpathlineto{\pgfqpoint{5.947354in}{0.725685in}}%
\pgfpathlineto{\pgfqpoint{5.950710in}{0.731701in}}%
\pgfpathlineto{\pgfqpoint{5.954066in}{0.729680in}}%
\pgfpathlineto{\pgfqpoint{5.957421in}{0.718041in}}%
\pgfpathlineto{\pgfqpoint{5.960777in}{0.682930in}}%
\pgfpathlineto{\pgfqpoint{5.964133in}{0.680948in}}%
\pgfpathlineto{\pgfqpoint{5.967489in}{0.723471in}}%
\pgfpathlineto{\pgfqpoint{5.970845in}{0.739726in}}%
\pgfpathlineto{\pgfqpoint{5.974200in}{0.744632in}}%
\pgfpathlineto{\pgfqpoint{5.977556in}{0.737268in}}%
\pgfpathlineto{\pgfqpoint{5.980912in}{0.717736in}}%
\pgfpathlineto{\pgfqpoint{5.984268in}{0.672749in}}%
\pgfpathlineto{\pgfqpoint{5.987623in}{0.667367in}}%
\pgfpathlineto{\pgfqpoint{5.990979in}{0.697652in}}%
\pgfpathlineto{\pgfqpoint{5.994335in}{0.700159in}}%
\pgfpathlineto{\pgfqpoint{5.997691in}{0.682283in}}%
\pgfpathlineto{\pgfqpoint{6.001047in}{0.617234in}}%
\pgfpathlineto{\pgfqpoint{6.004402in}{0.700568in}}%
\pgfpathlineto{\pgfqpoint{6.007758in}{0.722630in}}%
\pgfpathlineto{\pgfqpoint{6.011114in}{0.729609in}}%
\pgfpathlineto{\pgfqpoint{6.014470in}{0.721712in}}%
\pgfpathlineto{\pgfqpoint{6.017825in}{0.701006in}}%
\pgfpathlineto{\pgfqpoint{6.021181in}{0.650868in}}%
\pgfpathlineto{\pgfqpoint{6.024537in}{0.658065in}}%
\pgfpathlineto{\pgfqpoint{6.027893in}{0.684218in}}%
\pgfpathlineto{\pgfqpoint{6.031249in}{0.685202in}}%
\pgfpathlineto{\pgfqpoint{6.034604in}{0.664757in}}%
\pgfpathlineto{\pgfqpoint{6.037960in}{0.622012in}}%
\pgfpathlineto{\pgfqpoint{6.041316in}{0.688939in}}%
\pgfpathlineto{\pgfqpoint{6.044672in}{0.709361in}}%
\pgfpathlineto{\pgfqpoint{6.048027in}{0.715216in}}%
\pgfpathlineto{\pgfqpoint{6.051383in}{0.705481in}}%
\pgfpathlineto{\pgfqpoint{6.054739in}{0.674803in}}%
\pgfpathlineto{\pgfqpoint{6.058095in}{0.637013in}}%
\pgfpathlineto{\pgfqpoint{6.061451in}{0.690046in}}%
\pgfpathlineto{\pgfqpoint{6.064806in}{0.704143in}}%
\pgfpathlineto{\pgfqpoint{6.068162in}{0.705993in}}%
\pgfpathlineto{\pgfqpoint{6.071518in}{0.697014in}}%
\pgfpathlineto{\pgfqpoint{6.074874in}{0.666095in}}%
\pgfpathlineto{\pgfqpoint{6.078230in}{0.654216in}}%
\pgfpathlineto{\pgfqpoint{6.081585in}{0.701267in}}%
\pgfpathlineto{\pgfqpoint{6.084941in}{0.718184in}}%
\pgfpathlineto{\pgfqpoint{6.088297in}{0.722540in}}%
\pgfpathlineto{\pgfqpoint{6.091653in}{0.712115in}}%
\pgfpathlineto{\pgfqpoint{6.095008in}{0.671203in}}%
\pgfpathlineto{\pgfqpoint{6.098364in}{0.680115in}}%
\pgfpathlineto{\pgfqpoint{6.101720in}{0.715511in}}%
\pgfpathlineto{\pgfqpoint{6.105076in}{0.728200in}}%
\pgfpathlineto{\pgfqpoint{6.108432in}{0.731232in}}%
\pgfpathlineto{\pgfqpoint{6.111787in}{0.726075in}}%
\pgfpathlineto{\pgfqpoint{6.115143in}{0.708735in}}%
\pgfpathlineto{\pgfqpoint{6.118499in}{0.634065in}}%
\pgfpathlineto{\pgfqpoint{6.121855in}{0.703368in}}%
\pgfpathlineto{\pgfqpoint{6.125210in}{0.731064in}}%
\pgfpathlineto{\pgfqpoint{6.128566in}{0.742441in}}%
\pgfpathlineto{\pgfqpoint{6.131922in}{0.742933in}}%
\pgfpathlineto{\pgfqpoint{6.135278in}{0.728762in}}%
\pgfpathlineto{\pgfqpoint{6.138634in}{0.684712in}}%
\pgfpathlineto{\pgfqpoint{6.141989in}{0.694114in}}%
\pgfpathlineto{\pgfqpoint{6.145345in}{0.727052in}}%
\pgfpathlineto{\pgfqpoint{6.148701in}{0.737418in}}%
\pgfpathlineto{\pgfqpoint{6.152057in}{0.737242in}}%
\pgfpathlineto{\pgfqpoint{6.155412in}{0.726166in}}%
\pgfpathlineto{\pgfqpoint{6.158768in}{0.688813in}}%
\pgfpathlineto{\pgfqpoint{6.162124in}{0.695686in}}%
\pgfpathlineto{\pgfqpoint{6.165480in}{0.735260in}}%
\pgfpathlineto{\pgfqpoint{6.168836in}{0.750383in}}%
\pgfpathlineto{\pgfqpoint{6.172191in}{0.753472in}}%
\pgfpathlineto{\pgfqpoint{6.175547in}{0.742955in}}%
\pgfpathlineto{\pgfqpoint{6.178903in}{0.719467in}}%
\pgfpathlineto{\pgfqpoint{6.182259in}{0.651047in}}%
\pgfpathlineto{\pgfqpoint{6.185614in}{0.691328in}}%
\pgfpathlineto{\pgfqpoint{6.188970in}{0.710656in}}%
\pgfpathlineto{\pgfqpoint{6.192326in}{0.708808in}}%
\pgfpathlineto{\pgfqpoint{6.195682in}{0.683266in}}%
\pgfpathlineto{\pgfqpoint{6.199038in}{0.665561in}}%
\pgfpathlineto{\pgfqpoint{6.202393in}{0.717422in}}%
\pgfpathlineto{\pgfqpoint{6.205749in}{0.735393in}}%
\pgfpathlineto{\pgfqpoint{6.209105in}{0.739386in}}%
\pgfpathlineto{\pgfqpoint{6.212461in}{0.728608in}}%
\pgfpathlineto{\pgfqpoint{6.215816in}{0.706759in}}%
\pgfpathlineto{\pgfqpoint{6.219172in}{0.657898in}}%
\pgfpathlineto{\pgfqpoint{6.222528in}{0.655602in}}%
\pgfpathlineto{\pgfqpoint{6.225884in}{0.680680in}}%
\pgfpathlineto{\pgfqpoint{6.229240in}{0.675599in}}%
\pgfpathlineto{\pgfqpoint{6.232595in}{0.620544in}}%
\pgfpathlineto{\pgfqpoint{6.235951in}{0.673775in}}%
\pgfpathlineto{\pgfqpoint{6.239307in}{0.705170in}}%
\pgfpathlineto{\pgfqpoint{6.242663in}{0.717258in}}%
\pgfpathlineto{\pgfqpoint{6.246019in}{0.714959in}}%
\pgfpathlineto{\pgfqpoint{6.249374in}{0.697176in}}%
\pgfpathlineto{\pgfqpoint{6.252730in}{0.651064in}}%
\pgfpathlineto{\pgfqpoint{6.256086in}{0.656548in}}%
\pgfpathlineto{\pgfqpoint{6.259442in}{0.687091in}}%
\pgfpathlineto{\pgfqpoint{6.262797in}{0.693729in}}%
\pgfpathlineto{\pgfqpoint{6.266153in}{0.686940in}}%
\pgfpathlineto{\pgfqpoint{6.269509in}{0.656903in}}%
\pgfpathlineto{\pgfqpoint{6.272865in}{0.647245in}}%
\pgfpathlineto{\pgfqpoint{6.276221in}{0.693990in}}%
\pgfpathlineto{\pgfqpoint{6.279576in}{0.710712in}}%
\pgfpathlineto{\pgfqpoint{6.282932in}{0.714170in}}%
\pgfpathlineto{\pgfqpoint{6.286288in}{0.701364in}}%
\pgfpathlineto{\pgfqpoint{6.289644in}{0.652852in}}%
\pgfpathlineto{\pgfqpoint{6.292999in}{0.676670in}}%
\pgfpathlineto{\pgfqpoint{6.296355in}{0.707726in}}%
\pgfpathlineto{\pgfqpoint{6.299711in}{0.718887in}}%
\pgfpathlineto{\pgfqpoint{6.303067in}{0.720641in}}%
\pgfpathlineto{\pgfqpoint{6.306423in}{0.713735in}}%
\pgfpathlineto{\pgfqpoint{6.309778in}{0.692198in}}%
\pgfpathlineto{\pgfqpoint{6.313134in}{0.599358in}}%
\pgfpathlineto{\pgfqpoint{6.316490in}{0.701850in}}%
\pgfpathlineto{\pgfqpoint{6.319846in}{0.724668in}}%
\pgfpathlineto{\pgfqpoint{6.323201in}{0.733589in}}%
\pgfpathlineto{\pgfqpoint{6.326557in}{0.731064in}}%
\pgfpathlineto{\pgfqpoint{6.329913in}{0.710378in}}%
\pgfpathlineto{\pgfqpoint{6.333269in}{0.578264in}}%
\pgfpathlineto{\pgfqpoint{6.336625in}{0.712287in}}%
\pgfpathlineto{\pgfqpoint{6.339980in}{0.733279in}}%
\pgfpathlineto{\pgfqpoint{6.343336in}{0.740729in}}%
\pgfpathlineto{\pgfqpoint{6.346692in}{0.739723in}}%
\pgfpathlineto{\pgfqpoint{6.350048in}{0.729298in}}%
\pgfpathlineto{\pgfqpoint{6.353403in}{0.698012in}}%
\pgfpathlineto{\pgfqpoint{6.356759in}{0.683457in}}%
\pgfpathlineto{\pgfqpoint{6.360115in}{0.731340in}}%
\pgfpathlineto{\pgfqpoint{6.363471in}{0.748740in}}%
\pgfpathlineto{\pgfqpoint{6.366827in}{0.754346in}}%
\pgfpathlineto{\pgfqpoint{6.370182in}{0.748002in}}%
\pgfpathlineto{\pgfqpoint{6.383605in}{0.705650in}}%
\pgfpathlineto{\pgfqpoint{6.383605in}{0.705650in}}%
\pgfusepath{stroke}%
\end{pgfscope}%
\begin{pgfscope}%
\pgfpathrectangle{\pgfqpoint{4.621821in}{0.320625in}}{\pgfqpoint{1.845679in}{1.164375in}}%
\pgfusepath{clip}%
\pgfsetrectcap%
\pgfsetroundjoin%
\pgfsetlinewidth{1.154312pt}%
\definecolor{currentstroke}{rgb}{0.478431,0.243137,0.615686}%
\pgfsetstrokecolor{currentstroke}%
\pgfsetdash{}{0pt}%
\pgfpathmoveto{\pgfqpoint{4.705715in}{1.468748in}}%
\pgfpathlineto{\pgfqpoint{4.715783in}{1.450185in}}%
\pgfpathlineto{\pgfqpoint{4.725850in}{1.429162in}}%
\pgfpathlineto{\pgfqpoint{4.739273in}{1.397858in}}%
\pgfpathlineto{\pgfqpoint{4.779543in}{1.301803in}}%
\pgfpathlineto{\pgfqpoint{4.799677in}{1.257107in}}%
\pgfpathlineto{\pgfqpoint{4.823168in}{1.207628in}}%
\pgfpathlineto{\pgfqpoint{4.850014in}{1.153585in}}%
\pgfpathlineto{\pgfqpoint{4.883572in}{1.088504in}}%
\pgfpathlineto{\pgfqpoint{4.927197in}{1.006508in}}%
\pgfpathlineto{\pgfqpoint{4.977534in}{0.914247in}}%
\pgfpathlineto{\pgfqpoint{5.024515in}{0.829849in}}%
\pgfpathlineto{\pgfqpoint{5.048005in}{0.788271in}}%
\pgfpathlineto{\pgfqpoint{5.058072in}{0.770693in}}%
\pgfpathlineto{\pgfqpoint{5.078207in}{0.733825in}}%
\pgfpathlineto{\pgfqpoint{5.094986in}{0.710118in}}%
\pgfpathlineto{\pgfqpoint{5.101698in}{0.696590in}}%
\pgfpathlineto{\pgfqpoint{5.108409in}{0.681229in}}%
\pgfpathlineto{\pgfqpoint{5.115121in}{0.664184in}}%
\pgfpathlineto{\pgfqpoint{5.118476in}{0.657028in}}%
\pgfpathlineto{\pgfqpoint{5.121832in}{0.652646in}}%
\pgfpathlineto{\pgfqpoint{5.125188in}{0.651760in}}%
\pgfpathlineto{\pgfqpoint{5.128544in}{0.653456in}}%
\pgfpathlineto{\pgfqpoint{5.135255in}{0.657966in}}%
\pgfpathlineto{\pgfqpoint{5.138611in}{0.658695in}}%
\pgfpathlineto{\pgfqpoint{5.141967in}{0.658804in}}%
\pgfpathlineto{\pgfqpoint{5.145323in}{0.657507in}}%
\pgfpathlineto{\pgfqpoint{5.148678in}{0.654006in}}%
\pgfpathlineto{\pgfqpoint{5.152034in}{0.648253in}}%
\pgfpathlineto{\pgfqpoint{5.155390in}{0.646296in}}%
\pgfpathlineto{\pgfqpoint{5.158746in}{0.657415in}}%
\pgfpathlineto{\pgfqpoint{5.162102in}{0.671737in}}%
\pgfpathlineto{\pgfqpoint{5.165457in}{0.682933in}}%
\pgfpathlineto{\pgfqpoint{5.168813in}{0.690751in}}%
\pgfpathlineto{\pgfqpoint{5.172169in}{0.695819in}}%
\pgfpathlineto{\pgfqpoint{5.175525in}{0.699083in}}%
\pgfpathlineto{\pgfqpoint{5.178880in}{0.701295in}}%
\pgfpathlineto{\pgfqpoint{5.182236in}{0.702742in}}%
\pgfpathlineto{\pgfqpoint{5.185592in}{0.703431in}}%
\pgfpathlineto{\pgfqpoint{5.188948in}{0.703281in}}%
\pgfpathlineto{\pgfqpoint{5.192304in}{0.702255in}}%
\pgfpathlineto{\pgfqpoint{5.202371in}{0.697283in}}%
\pgfpathlineto{\pgfqpoint{5.205727in}{0.698064in}}%
\pgfpathlineto{\pgfqpoint{5.209082in}{0.701583in}}%
\pgfpathlineto{\pgfqpoint{5.215794in}{0.711980in}}%
\pgfpathlineto{\pgfqpoint{5.219150in}{0.715766in}}%
\pgfpathlineto{\pgfqpoint{5.222506in}{0.717806in}}%
\pgfpathlineto{\pgfqpoint{5.225861in}{0.717924in}}%
\pgfpathlineto{\pgfqpoint{5.229217in}{0.717174in}}%
\pgfpathlineto{\pgfqpoint{5.232573in}{0.715499in}}%
\pgfpathlineto{\pgfqpoint{5.235929in}{0.711112in}}%
\pgfpathlineto{\pgfqpoint{5.239285in}{0.701888in}}%
\pgfpathlineto{\pgfqpoint{5.242640in}{0.689137in}}%
\pgfpathlineto{\pgfqpoint{5.245996in}{0.680012in}}%
\pgfpathlineto{\pgfqpoint{5.249352in}{0.687071in}}%
\pgfpathlineto{\pgfqpoint{5.252708in}{0.699156in}}%
\pgfpathlineto{\pgfqpoint{5.256063in}{0.707395in}}%
\pgfpathlineto{\pgfqpoint{5.259419in}{0.710084in}}%
\pgfpathlineto{\pgfqpoint{5.262775in}{0.705204in}}%
\pgfpathlineto{\pgfqpoint{5.266131in}{0.691756in}}%
\pgfpathlineto{\pgfqpoint{5.269487in}{0.686154in}}%
\pgfpathlineto{\pgfqpoint{5.272842in}{0.698151in}}%
\pgfpathlineto{\pgfqpoint{5.276198in}{0.707923in}}%
\pgfpathlineto{\pgfqpoint{5.279554in}{0.712967in}}%
\pgfpathlineto{\pgfqpoint{5.282910in}{0.712531in}}%
\pgfpathlineto{\pgfqpoint{5.286265in}{0.706235in}}%
\pgfpathlineto{\pgfqpoint{5.289621in}{0.694254in}}%
\pgfpathlineto{\pgfqpoint{5.292977in}{0.688043in}}%
\pgfpathlineto{\pgfqpoint{5.299689in}{0.713178in}}%
\pgfpathlineto{\pgfqpoint{5.303044in}{0.720119in}}%
\pgfpathlineto{\pgfqpoint{5.306400in}{0.722700in}}%
\pgfpathlineto{\pgfqpoint{5.309756in}{0.719106in}}%
\pgfpathlineto{\pgfqpoint{5.313112in}{0.712851in}}%
\pgfpathlineto{\pgfqpoint{5.316467in}{0.712125in}}%
\pgfpathlineto{\pgfqpoint{5.319823in}{0.714599in}}%
\pgfpathlineto{\pgfqpoint{5.323179in}{0.714840in}}%
\pgfpathlineto{\pgfqpoint{5.326535in}{0.710476in}}%
\pgfpathlineto{\pgfqpoint{5.329891in}{0.702061in}}%
\pgfpathlineto{\pgfqpoint{5.333246in}{0.698984in}}%
\pgfpathlineto{\pgfqpoint{5.339958in}{0.720270in}}%
\pgfpathlineto{\pgfqpoint{5.343314in}{0.726404in}}%
\pgfpathlineto{\pgfqpoint{5.346669in}{0.726302in}}%
\pgfpathlineto{\pgfqpoint{5.350025in}{0.721693in}}%
\pgfpathlineto{\pgfqpoint{5.356737in}{0.708012in}}%
\pgfpathlineto{\pgfqpoint{5.366804in}{0.682776in}}%
\pgfpathlineto{\pgfqpoint{5.370160in}{0.691260in}}%
\pgfpathlineto{\pgfqpoint{5.373516in}{0.707362in}}%
\pgfpathlineto{\pgfqpoint{5.376871in}{0.718154in}}%
\pgfpathlineto{\pgfqpoint{5.380227in}{0.720917in}}%
\pgfpathlineto{\pgfqpoint{5.383583in}{0.714301in}}%
\pgfpathlineto{\pgfqpoint{5.386939in}{0.704233in}}%
\pgfpathlineto{\pgfqpoint{5.390295in}{0.705488in}}%
\pgfpathlineto{\pgfqpoint{5.393650in}{0.713049in}}%
\pgfpathlineto{\pgfqpoint{5.397006in}{0.714742in}}%
\pgfpathlineto{\pgfqpoint{5.400362in}{0.710636in}}%
\pgfpathlineto{\pgfqpoint{5.403718in}{0.698352in}}%
\pgfpathlineto{\pgfqpoint{5.407074in}{0.677356in}}%
\pgfpathlineto{\pgfqpoint{5.410429in}{0.692502in}}%
\pgfpathlineto{\pgfqpoint{5.413785in}{0.715060in}}%
\pgfpathlineto{\pgfqpoint{5.417141in}{0.726199in}}%
\pgfpathlineto{\pgfqpoint{5.420497in}{0.727048in}}%
\pgfpathlineto{\pgfqpoint{5.423852in}{0.715587in}}%
\pgfpathlineto{\pgfqpoint{5.427208in}{0.699707in}}%
\pgfpathlineto{\pgfqpoint{5.430564in}{0.704871in}}%
\pgfpathlineto{\pgfqpoint{5.433920in}{0.715585in}}%
\pgfpathlineto{\pgfqpoint{5.437276in}{0.719502in}}%
\pgfpathlineto{\pgfqpoint{5.440631in}{0.717859in}}%
\pgfpathlineto{\pgfqpoint{5.443987in}{0.710674in}}%
\pgfpathlineto{\pgfqpoint{5.447343in}{0.701794in}}%
\pgfpathlineto{\pgfqpoint{5.450699in}{0.706679in}}%
\pgfpathlineto{\pgfqpoint{5.454054in}{0.719912in}}%
\pgfpathlineto{\pgfqpoint{5.457410in}{0.727845in}}%
\pgfpathlineto{\pgfqpoint{5.460766in}{0.728677in}}%
\pgfpathlineto{\pgfqpoint{5.464122in}{0.726627in}}%
\pgfpathlineto{\pgfqpoint{5.467478in}{0.721047in}}%
\pgfpathlineto{\pgfqpoint{5.470833in}{0.718530in}}%
\pgfpathlineto{\pgfqpoint{5.474189in}{0.719288in}}%
\pgfpathlineto{\pgfqpoint{5.477545in}{0.718274in}}%
\pgfpathlineto{\pgfqpoint{5.480901in}{0.715755in}}%
\pgfpathlineto{\pgfqpoint{5.484256in}{0.711711in}}%
\pgfpathlineto{\pgfqpoint{5.487612in}{0.708953in}}%
\pgfpathlineto{\pgfqpoint{5.490968in}{0.712193in}}%
\pgfpathlineto{\pgfqpoint{5.497680in}{0.726170in}}%
\pgfpathlineto{\pgfqpoint{5.501035in}{0.729324in}}%
\pgfpathlineto{\pgfqpoint{5.504391in}{0.727455in}}%
\pgfpathlineto{\pgfqpoint{5.507747in}{0.723853in}}%
\pgfpathlineto{\pgfqpoint{5.511103in}{0.719366in}}%
\pgfpathlineto{\pgfqpoint{5.514458in}{0.712336in}}%
\pgfpathlineto{\pgfqpoint{5.521170in}{0.679907in}}%
\pgfpathlineto{\pgfqpoint{5.524526in}{0.684807in}}%
\pgfpathlineto{\pgfqpoint{5.527882in}{0.699808in}}%
\pgfpathlineto{\pgfqpoint{5.531237in}{0.708649in}}%
\pgfpathlineto{\pgfqpoint{5.534593in}{0.711562in}}%
\pgfpathlineto{\pgfqpoint{5.541305in}{0.706378in}}%
\pgfpathlineto{\pgfqpoint{5.544660in}{0.707328in}}%
\pgfpathlineto{\pgfqpoint{5.548016in}{0.712692in}}%
\pgfpathlineto{\pgfqpoint{5.551372in}{0.714800in}}%
\pgfpathlineto{\pgfqpoint{5.554728in}{0.706432in}}%
\pgfpathlineto{\pgfqpoint{5.561439in}{0.667634in}}%
\pgfpathlineto{\pgfqpoint{5.564795in}{0.690888in}}%
\pgfpathlineto{\pgfqpoint{5.568151in}{0.707298in}}%
\pgfpathlineto{\pgfqpoint{5.571507in}{0.712764in}}%
\pgfpathlineto{\pgfqpoint{5.574863in}{0.707706in}}%
\pgfpathlineto{\pgfqpoint{5.578218in}{0.695505in}}%
\pgfpathlineto{\pgfqpoint{5.581574in}{0.691260in}}%
\pgfpathlineto{\pgfqpoint{5.584930in}{0.708363in}}%
\pgfpathlineto{\pgfqpoint{5.588286in}{0.720147in}}%
\pgfpathlineto{\pgfqpoint{5.591641in}{0.723147in}}%
\pgfpathlineto{\pgfqpoint{5.594997in}{0.720036in}}%
\pgfpathlineto{\pgfqpoint{5.598353in}{0.711913in}}%
\pgfpathlineto{\pgfqpoint{5.601709in}{0.702162in}}%
\pgfpathlineto{\pgfqpoint{5.605065in}{0.705411in}}%
\pgfpathlineto{\pgfqpoint{5.608420in}{0.718280in}}%
\pgfpathlineto{\pgfqpoint{5.611776in}{0.726906in}}%
\pgfpathlineto{\pgfqpoint{5.615132in}{0.728264in}}%
\pgfpathlineto{\pgfqpoint{5.625199in}{0.719666in}}%
\pgfpathlineto{\pgfqpoint{5.628555in}{0.721286in}}%
\pgfpathlineto{\pgfqpoint{5.631911in}{0.721615in}}%
\pgfpathlineto{\pgfqpoint{5.635267in}{0.718559in}}%
\pgfpathlineto{\pgfqpoint{5.641978in}{0.704946in}}%
\pgfpathlineto{\pgfqpoint{5.645334in}{0.713479in}}%
\pgfpathlineto{\pgfqpoint{5.648690in}{0.726525in}}%
\pgfpathlineto{\pgfqpoint{5.652045in}{0.733572in}}%
\pgfpathlineto{\pgfqpoint{5.655401in}{0.732376in}}%
\pgfpathlineto{\pgfqpoint{5.665469in}{0.716421in}}%
\pgfpathlineto{\pgfqpoint{5.668824in}{0.709829in}}%
\pgfpathlineto{\pgfqpoint{5.672180in}{0.696397in}}%
\pgfpathlineto{\pgfqpoint{5.675536in}{0.674827in}}%
\pgfpathlineto{\pgfqpoint{5.678892in}{0.680155in}}%
\pgfpathlineto{\pgfqpoint{5.682247in}{0.703091in}}%
\pgfpathlineto{\pgfqpoint{5.685603in}{0.715941in}}%
\pgfpathlineto{\pgfqpoint{5.688959in}{0.719440in}}%
\pgfpathlineto{\pgfqpoint{5.692315in}{0.713985in}}%
\pgfpathlineto{\pgfqpoint{5.695671in}{0.705390in}}%
\pgfpathlineto{\pgfqpoint{5.699026in}{0.702005in}}%
\pgfpathlineto{\pgfqpoint{5.702382in}{0.705375in}}%
\pgfpathlineto{\pgfqpoint{5.705738in}{0.712420in}}%
\pgfpathlineto{\pgfqpoint{5.709094in}{0.712891in}}%
\pgfpathlineto{\pgfqpoint{5.712449in}{0.701552in}}%
\pgfpathlineto{\pgfqpoint{5.715805in}{0.682673in}}%
\pgfpathlineto{\pgfqpoint{5.719161in}{0.685908in}}%
\pgfpathlineto{\pgfqpoint{5.722517in}{0.701690in}}%
\pgfpathlineto{\pgfqpoint{5.725873in}{0.708112in}}%
\pgfpathlineto{\pgfqpoint{5.729228in}{0.707190in}}%
\pgfpathlineto{\pgfqpoint{5.732584in}{0.702569in}}%
\pgfpathlineto{\pgfqpoint{5.735940in}{0.694797in}}%
\pgfpathlineto{\pgfqpoint{5.739296in}{0.694097in}}%
\pgfpathlineto{\pgfqpoint{5.742652in}{0.708524in}}%
\pgfpathlineto{\pgfqpoint{5.746007in}{0.718820in}}%
\pgfpathlineto{\pgfqpoint{5.749363in}{0.720313in}}%
\pgfpathlineto{\pgfqpoint{5.756075in}{0.709740in}}%
\pgfpathlineto{\pgfqpoint{5.759430in}{0.706877in}}%
\pgfpathlineto{\pgfqpoint{5.766142in}{0.706215in}}%
\pgfpathlineto{\pgfqpoint{5.769498in}{0.708267in}}%
\pgfpathlineto{\pgfqpoint{5.776209in}{0.724411in}}%
\pgfpathlineto{\pgfqpoint{5.779565in}{0.726262in}}%
\pgfpathlineto{\pgfqpoint{5.782921in}{0.725262in}}%
\pgfpathlineto{\pgfqpoint{5.786277in}{0.722384in}}%
\pgfpathlineto{\pgfqpoint{5.792988in}{0.707852in}}%
\pgfpathlineto{\pgfqpoint{5.796344in}{0.707253in}}%
\pgfpathlineto{\pgfqpoint{5.803056in}{0.723473in}}%
\pgfpathlineto{\pgfqpoint{5.806411in}{0.727164in}}%
\pgfpathlineto{\pgfqpoint{5.809767in}{0.725169in}}%
\pgfpathlineto{\pgfqpoint{5.813123in}{0.721992in}}%
\pgfpathlineto{\pgfqpoint{5.819834in}{0.719531in}}%
\pgfpathlineto{\pgfqpoint{5.823190in}{0.713318in}}%
\pgfpathlineto{\pgfqpoint{5.826546in}{0.700331in}}%
\pgfpathlineto{\pgfqpoint{5.829902in}{0.676423in}}%
\pgfpathlineto{\pgfqpoint{5.833258in}{0.670612in}}%
\pgfpathlineto{\pgfqpoint{5.836613in}{0.696948in}}%
\pgfpathlineto{\pgfqpoint{5.839969in}{0.711227in}}%
\pgfpathlineto{\pgfqpoint{5.843325in}{0.715210in}}%
\pgfpathlineto{\pgfqpoint{5.846681in}{0.709607in}}%
\pgfpathlineto{\pgfqpoint{5.850036in}{0.702396in}}%
\pgfpathlineto{\pgfqpoint{5.853392in}{0.702109in}}%
\pgfpathlineto{\pgfqpoint{5.856748in}{0.704419in}}%
\pgfpathlineto{\pgfqpoint{5.860104in}{0.705374in}}%
\pgfpathlineto{\pgfqpoint{5.863460in}{0.703106in}}%
\pgfpathlineto{\pgfqpoint{5.870171in}{0.678813in}}%
\pgfpathlineto{\pgfqpoint{5.873527in}{0.689477in}}%
\pgfpathlineto{\pgfqpoint{5.876883in}{0.703018in}}%
\pgfpathlineto{\pgfqpoint{5.880238in}{0.707899in}}%
\pgfpathlineto{\pgfqpoint{5.883594in}{0.705016in}}%
\pgfpathlineto{\pgfqpoint{5.886950in}{0.701146in}}%
\pgfpathlineto{\pgfqpoint{5.893662in}{0.698760in}}%
\pgfpathlineto{\pgfqpoint{5.897017in}{0.703644in}}%
\pgfpathlineto{\pgfqpoint{5.900373in}{0.712361in}}%
\pgfpathlineto{\pgfqpoint{5.903729in}{0.716700in}}%
\pgfpathlineto{\pgfqpoint{5.907085in}{0.714348in}}%
\pgfpathlineto{\pgfqpoint{5.910441in}{0.711031in}}%
\pgfpathlineto{\pgfqpoint{5.913796in}{0.711528in}}%
\pgfpathlineto{\pgfqpoint{5.917152in}{0.713948in}}%
\pgfpathlineto{\pgfqpoint{5.920508in}{0.713196in}}%
\pgfpathlineto{\pgfqpoint{5.923864in}{0.709137in}}%
\pgfpathlineto{\pgfqpoint{5.927219in}{0.712183in}}%
\pgfpathlineto{\pgfqpoint{5.930575in}{0.721209in}}%
\pgfpathlineto{\pgfqpoint{5.933931in}{0.726802in}}%
\pgfpathlineto{\pgfqpoint{5.937287in}{0.727156in}}%
\pgfpathlineto{\pgfqpoint{5.940643in}{0.726595in}}%
\pgfpathlineto{\pgfqpoint{5.943998in}{0.723653in}}%
\pgfpathlineto{\pgfqpoint{5.947354in}{0.714808in}}%
\pgfpathlineto{\pgfqpoint{5.950710in}{0.708642in}}%
\pgfpathlineto{\pgfqpoint{5.954066in}{0.715304in}}%
\pgfpathlineto{\pgfqpoint{5.957421in}{0.724672in}}%
\pgfpathlineto{\pgfqpoint{5.960777in}{0.728661in}}%
\pgfpathlineto{\pgfqpoint{5.964133in}{0.724659in}}%
\pgfpathlineto{\pgfqpoint{5.967489in}{0.717623in}}%
\pgfpathlineto{\pgfqpoint{5.970845in}{0.718378in}}%
\pgfpathlineto{\pgfqpoint{5.974200in}{0.723340in}}%
\pgfpathlineto{\pgfqpoint{5.977556in}{0.725560in}}%
\pgfpathlineto{\pgfqpoint{5.980912in}{0.723816in}}%
\pgfpathlineto{\pgfqpoint{5.984268in}{0.713902in}}%
\pgfpathlineto{\pgfqpoint{5.987623in}{0.691222in}}%
\pgfpathlineto{\pgfqpoint{5.990979in}{0.686500in}}%
\pgfpathlineto{\pgfqpoint{5.994335in}{0.712290in}}%
\pgfpathlineto{\pgfqpoint{5.997691in}{0.725379in}}%
\pgfpathlineto{\pgfqpoint{6.001047in}{0.727187in}}%
\pgfpathlineto{\pgfqpoint{6.004402in}{0.718283in}}%
\pgfpathlineto{\pgfqpoint{6.007758in}{0.706668in}}%
\pgfpathlineto{\pgfqpoint{6.011114in}{0.702961in}}%
\pgfpathlineto{\pgfqpoint{6.017825in}{0.708220in}}%
\pgfpathlineto{\pgfqpoint{6.021181in}{0.703798in}}%
\pgfpathlineto{\pgfqpoint{6.024537in}{0.689659in}}%
\pgfpathlineto{\pgfqpoint{6.027893in}{0.683359in}}%
\pgfpathlineto{\pgfqpoint{6.031249in}{0.700521in}}%
\pgfpathlineto{\pgfqpoint{6.034604in}{0.712654in}}%
\pgfpathlineto{\pgfqpoint{6.037960in}{0.714292in}}%
\pgfpathlineto{\pgfqpoint{6.041316in}{0.706800in}}%
\pgfpathlineto{\pgfqpoint{6.044672in}{0.701415in}}%
\pgfpathlineto{\pgfqpoint{6.048027in}{0.700326in}}%
\pgfpathlineto{\pgfqpoint{6.051383in}{0.698263in}}%
\pgfpathlineto{\pgfqpoint{6.054739in}{0.700169in}}%
\pgfpathlineto{\pgfqpoint{6.058095in}{0.703793in}}%
\pgfpathlineto{\pgfqpoint{6.061451in}{0.704952in}}%
\pgfpathlineto{\pgfqpoint{6.064806in}{0.708074in}}%
\pgfpathlineto{\pgfqpoint{6.068162in}{0.712675in}}%
\pgfpathlineto{\pgfqpoint{6.071518in}{0.714933in}}%
\pgfpathlineto{\pgfqpoint{6.074874in}{0.710914in}}%
\pgfpathlineto{\pgfqpoint{6.078230in}{0.699849in}}%
\pgfpathlineto{\pgfqpoint{6.081585in}{0.698993in}}%
\pgfpathlineto{\pgfqpoint{6.084941in}{0.713584in}}%
\pgfpathlineto{\pgfqpoint{6.088297in}{0.723586in}}%
\pgfpathlineto{\pgfqpoint{6.091653in}{0.725246in}}%
\pgfpathlineto{\pgfqpoint{6.095008in}{0.720693in}}%
\pgfpathlineto{\pgfqpoint{6.098364in}{0.714161in}}%
\pgfpathlineto{\pgfqpoint{6.101720in}{0.705565in}}%
\pgfpathlineto{\pgfqpoint{6.105076in}{0.704624in}}%
\pgfpathlineto{\pgfqpoint{6.108432in}{0.713333in}}%
\pgfpathlineto{\pgfqpoint{6.111787in}{0.720438in}}%
\pgfpathlineto{\pgfqpoint{6.115143in}{0.722225in}}%
\pgfpathlineto{\pgfqpoint{6.118499in}{0.718488in}}%
\pgfpathlineto{\pgfqpoint{6.121855in}{0.712612in}}%
\pgfpathlineto{\pgfqpoint{6.125210in}{0.716224in}}%
\pgfpathlineto{\pgfqpoint{6.128566in}{0.723618in}}%
\pgfpathlineto{\pgfqpoint{6.131922in}{0.727862in}}%
\pgfpathlineto{\pgfqpoint{6.135278in}{0.727205in}}%
\pgfpathlineto{\pgfqpoint{6.138634in}{0.722922in}}%
\pgfpathlineto{\pgfqpoint{6.141989in}{0.713671in}}%
\pgfpathlineto{\pgfqpoint{6.145345in}{0.709750in}}%
\pgfpathlineto{\pgfqpoint{6.152057in}{0.729648in}}%
\pgfpathlineto{\pgfqpoint{6.155412in}{0.732229in}}%
\pgfpathlineto{\pgfqpoint{6.158768in}{0.726264in}}%
\pgfpathlineto{\pgfqpoint{6.162124in}{0.714617in}}%
\pgfpathlineto{\pgfqpoint{6.165480in}{0.716309in}}%
\pgfpathlineto{\pgfqpoint{6.168836in}{0.727839in}}%
\pgfpathlineto{\pgfqpoint{6.172191in}{0.733845in}}%
\pgfpathlineto{\pgfqpoint{6.175547in}{0.731439in}}%
\pgfpathlineto{\pgfqpoint{6.178903in}{0.722724in}}%
\pgfpathlineto{\pgfqpoint{6.182259in}{0.704555in}}%
\pgfpathlineto{\pgfqpoint{6.185614in}{0.703005in}}%
\pgfpathlineto{\pgfqpoint{6.188970in}{0.721251in}}%
\pgfpathlineto{\pgfqpoint{6.192326in}{0.730820in}}%
\pgfpathlineto{\pgfqpoint{6.195682in}{0.729080in}}%
\pgfpathlineto{\pgfqpoint{6.202393in}{0.705833in}}%
\pgfpathlineto{\pgfqpoint{6.205749in}{0.710592in}}%
\pgfpathlineto{\pgfqpoint{6.209105in}{0.717858in}}%
\pgfpathlineto{\pgfqpoint{6.212461in}{0.718937in}}%
\pgfpathlineto{\pgfqpoint{6.215816in}{0.716935in}}%
\pgfpathlineto{\pgfqpoint{6.219172in}{0.710261in}}%
\pgfpathlineto{\pgfqpoint{6.222528in}{0.708136in}}%
\pgfpathlineto{\pgfqpoint{6.229240in}{0.721705in}}%
\pgfpathlineto{\pgfqpoint{6.232595in}{0.721347in}}%
\pgfpathlineto{\pgfqpoint{6.235951in}{0.711996in}}%
\pgfpathlineto{\pgfqpoint{6.239307in}{0.708289in}}%
\pgfpathlineto{\pgfqpoint{6.242663in}{0.715710in}}%
\pgfpathlineto{\pgfqpoint{6.246019in}{0.720232in}}%
\pgfpathlineto{\pgfqpoint{6.249374in}{0.716187in}}%
\pgfpathlineto{\pgfqpoint{6.252730in}{0.706007in}}%
\pgfpathlineto{\pgfqpoint{6.256086in}{0.688889in}}%
\pgfpathlineto{\pgfqpoint{6.259442in}{0.683841in}}%
\pgfpathlineto{\pgfqpoint{6.262797in}{0.700121in}}%
\pgfpathlineto{\pgfqpoint{6.266153in}{0.708665in}}%
\pgfpathlineto{\pgfqpoint{6.269509in}{0.704920in}}%
\pgfpathlineto{\pgfqpoint{6.272865in}{0.692328in}}%
\pgfpathlineto{\pgfqpoint{6.276221in}{0.691882in}}%
\pgfpathlineto{\pgfqpoint{6.279576in}{0.704822in}}%
\pgfpathlineto{\pgfqpoint{6.282932in}{0.711809in}}%
\pgfpathlineto{\pgfqpoint{6.286288in}{0.709640in}}%
\pgfpathlineto{\pgfqpoint{6.292999in}{0.700660in}}%
\pgfpathlineto{\pgfqpoint{6.296355in}{0.697196in}}%
\pgfpathlineto{\pgfqpoint{6.299711in}{0.702492in}}%
\pgfpathlineto{\pgfqpoint{6.303067in}{0.711225in}}%
\pgfpathlineto{\pgfqpoint{6.306423in}{0.714917in}}%
\pgfpathlineto{\pgfqpoint{6.309778in}{0.711719in}}%
\pgfpathlineto{\pgfqpoint{6.313134in}{0.702977in}}%
\pgfpathlineto{\pgfqpoint{6.316490in}{0.702257in}}%
\pgfpathlineto{\pgfqpoint{6.319846in}{0.712141in}}%
\pgfpathlineto{\pgfqpoint{6.323201in}{0.719937in}}%
\pgfpathlineto{\pgfqpoint{6.326557in}{0.722810in}}%
\pgfpathlineto{\pgfqpoint{6.329913in}{0.721112in}}%
\pgfpathlineto{\pgfqpoint{6.336625in}{0.713949in}}%
\pgfpathlineto{\pgfqpoint{6.339980in}{0.716667in}}%
\pgfpathlineto{\pgfqpoint{6.343336in}{0.725478in}}%
\pgfpathlineto{\pgfqpoint{6.346692in}{0.731392in}}%
\pgfpathlineto{\pgfqpoint{6.350048in}{0.731416in}}%
\pgfpathlineto{\pgfqpoint{6.356759in}{0.714521in}}%
\pgfpathlineto{\pgfqpoint{6.360115in}{0.714540in}}%
\pgfpathlineto{\pgfqpoint{6.363471in}{0.721721in}}%
\pgfpathlineto{\pgfqpoint{6.366827in}{0.727557in}}%
\pgfpathlineto{\pgfqpoint{6.370182in}{0.726907in}}%
\pgfpathlineto{\pgfqpoint{6.373538in}{0.725565in}}%
\pgfpathlineto{\pgfqpoint{6.380250in}{0.715956in}}%
\pgfpathlineto{\pgfqpoint{6.383605in}{0.709220in}}%
\pgfpathlineto{\pgfqpoint{6.383605in}{0.709220in}}%
\pgfusepath{stroke}%
\end{pgfscope}%
\begin{pgfscope}%
\pgfsetrectcap%
\pgfsetmiterjoin%
\pgfsetlinewidth{0.803000pt}%
\definecolor{currentstroke}{rgb}{0.000000,0.000000,0.000000}%
\pgfsetstrokecolor{currentstroke}%
\pgfsetdash{}{0pt}%
\pgfpathmoveto{\pgfqpoint{4.621821in}{0.320625in}}%
\pgfpathlineto{\pgfqpoint{4.621821in}{1.485000in}}%
\pgfusepath{stroke}%
\end{pgfscope}%
\begin{pgfscope}%
\pgfsetrectcap%
\pgfsetmiterjoin%
\pgfsetlinewidth{0.803000pt}%
\definecolor{currentstroke}{rgb}{0.000000,0.000000,0.000000}%
\pgfsetstrokecolor{currentstroke}%
\pgfsetdash{}{0pt}%
\pgfpathmoveto{\pgfqpoint{6.467500in}{0.320625in}}%
\pgfpathlineto{\pgfqpoint{6.467500in}{1.485000in}}%
\pgfusepath{stroke}%
\end{pgfscope}%
\begin{pgfscope}%
\pgfsetrectcap%
\pgfsetmiterjoin%
\pgfsetlinewidth{0.803000pt}%
\definecolor{currentstroke}{rgb}{0.000000,0.000000,0.000000}%
\pgfsetstrokecolor{currentstroke}%
\pgfsetdash{}{0pt}%
\pgfpathmoveto{\pgfqpoint{4.621821in}{0.320625in}}%
\pgfpathlineto{\pgfqpoint{6.467500in}{0.320625in}}%
\pgfusepath{stroke}%
\end{pgfscope}%
\begin{pgfscope}%
\pgfsetrectcap%
\pgfsetmiterjoin%
\pgfsetlinewidth{0.803000pt}%
\definecolor{currentstroke}{rgb}{0.000000,0.000000,0.000000}%
\pgfsetstrokecolor{currentstroke}%
\pgfsetdash{}{0pt}%
\pgfpathmoveto{\pgfqpoint{4.621821in}{1.485000in}}%
\pgfpathlineto{\pgfqpoint{6.467500in}{1.485000in}}%
\pgfusepath{stroke}%
\end{pgfscope}%
\begin{pgfscope}%
\definecolor{textcolor}{rgb}{0.000000,0.000000,0.000000}%
\pgfsetstrokecolor{textcolor}%
\pgfsetfillcolor{textcolor}%
\pgftext[x=5.544660in,y=1.568333in,,base]{\color{textcolor}{\rmfamily\fontsize{8.000000}{9.600000}\selectfont\catcode`\^=\active\def^{\ifmmode\sp\else\^{}\fi}\catcode`\%=\active\def%{\%}(c) Finite width}}%
\end{pgfscope}%
\begin{pgfscope}%
\pgfsetbuttcap%
\pgfsetmiterjoin%
\definecolor{currentfill}{rgb}{1.000000,1.000000,1.000000}%
\pgfsetfillcolor{currentfill}%
\pgfsetfillopacity{0.800000}%
\pgfsetlinewidth{1.003750pt}%
\definecolor{currentstroke}{rgb}{0.800000,0.800000,0.800000}%
\pgfsetstrokecolor{currentstroke}%
\pgfsetstrokeopacity{0.800000}%
\pgfsetdash{}{0pt}%
\pgfpathmoveto{\pgfqpoint{5.911892in}{1.084181in}}%
\pgfpathlineto{\pgfqpoint{6.424722in}{1.084181in}}%
\pgfpathquadraticcurveto{\pgfqpoint{6.436944in}{1.084181in}}{\pgfqpoint{6.436944in}{1.096403in}}%
\pgfpathlineto{\pgfqpoint{6.436944in}{1.442222in}}%
\pgfpathquadraticcurveto{\pgfqpoint{6.436944in}{1.454444in}}{\pgfqpoint{6.424722in}{1.454444in}}%
\pgfpathlineto{\pgfqpoint{5.911892in}{1.454444in}}%
\pgfpathquadraticcurveto{\pgfqpoint{5.899669in}{1.454444in}}{\pgfqpoint{5.899669in}{1.442222in}}%
\pgfpathlineto{\pgfqpoint{5.899669in}{1.096403in}}%
\pgfpathquadraticcurveto{\pgfqpoint{5.899669in}{1.084181in}}{\pgfqpoint{5.911892in}{1.084181in}}%
\pgfpathlineto{\pgfqpoint{5.911892in}{1.084181in}}%
\pgfpathclose%
\pgfusepath{stroke,fill}%
\end{pgfscope}%
\begin{pgfscope}%
\pgfsetrectcap%
\pgfsetroundjoin%
\pgfsetlinewidth{1.304875pt}%
\definecolor{currentstroke}{rgb}{0.000000,0.000000,0.000000}%
\pgfsetstrokecolor{currentstroke}%
\pgfsetdash{}{0pt}%
\pgfpathmoveto{\pgfqpoint{5.914947in}{1.414125in}}%
\pgfpathlineto{\pgfqpoint{5.954669in}{1.414125in}}%
\pgfpathlineto{\pgfqpoint{5.994392in}{1.414125in}}%
\pgfusepath{stroke}%
\end{pgfscope}%
\begin{pgfscope}%
\definecolor{textcolor}{rgb}{0.000000,0.000000,0.000000}%
\pgfsetstrokecolor{textcolor}%
\pgfsetfillcolor{textcolor}%
\pgftext[x=6.043281in,y=1.392737in,left,base]{\color{textcolor}{\rmfamily\fontsize{4.400000}{5.280000}\selectfont\catcode`\^=\active\def^{\ifmmode\sp\else\^{}\fi}\catcode`\%=\active\def%{\%}mean field}}%
\end{pgfscope}%
\begin{pgfscope}%
\pgfsetrectcap%
\pgfsetroundjoin%
\pgfsetlinewidth{1.154312pt}%
\definecolor{currentstroke}{rgb}{0.141176,0.345098,0.650980}%
\pgfsetstrokecolor{currentstroke}%
\pgfsetdash{}{0pt}%
\pgfpathmoveto{\pgfqpoint{5.914947in}{1.343984in}}%
\pgfpathlineto{\pgfqpoint{5.954669in}{1.343984in}}%
\pgfpathlineto{\pgfqpoint{5.994392in}{1.343984in}}%
\pgfusepath{stroke}%
\end{pgfscope}%
\begin{pgfscope}%
\definecolor{textcolor}{rgb}{0.000000,0.000000,0.000000}%
\pgfsetstrokecolor{textcolor}%
\pgfsetfillcolor{textcolor}%
\pgftext[x=6.043281in,y=1.322595in,left,base]{\color{textcolor}{\rmfamily\fontsize{4.400000}{5.280000}\selectfont\catcode`\^=\active\def^{\ifmmode\sp\else\^{}\fi}\catcode`\%=\active\def%{\%}$n=32$}}%
\end{pgfscope}%
\begin{pgfscope}%
\pgfsetrectcap%
\pgfsetroundjoin%
\pgfsetlinewidth{1.154312pt}%
\definecolor{currentstroke}{rgb}{0.819608,0.286275,0.356863}%
\pgfsetstrokecolor{currentstroke}%
\pgfsetdash{}{0pt}%
\pgfpathmoveto{\pgfqpoint{5.914947in}{1.273842in}}%
\pgfpathlineto{\pgfqpoint{5.954669in}{1.273842in}}%
\pgfpathlineto{\pgfqpoint{5.994392in}{1.273842in}}%
\pgfusepath{stroke}%
\end{pgfscope}%
\begin{pgfscope}%
\definecolor{textcolor}{rgb}{0.000000,0.000000,0.000000}%
\pgfsetstrokecolor{textcolor}%
\pgfsetfillcolor{textcolor}%
\pgftext[x=6.043281in,y=1.252453in,left,base]{\color{textcolor}{\rmfamily\fontsize{4.400000}{5.280000}\selectfont\catcode`\^=\active\def^{\ifmmode\sp\else\^{}\fi}\catcode`\%=\active\def%{\%}$n=128$}}%
\end{pgfscope}%
\begin{pgfscope}%
\pgfsetrectcap%
\pgfsetroundjoin%
\pgfsetlinewidth{1.154312pt}%
\definecolor{currentstroke}{rgb}{0.000000,0.529412,0.423529}%
\pgfsetstrokecolor{currentstroke}%
\pgfsetdash{}{0pt}%
\pgfpathmoveto{\pgfqpoint{5.914947in}{1.203701in}}%
\pgfpathlineto{\pgfqpoint{5.954669in}{1.203701in}}%
\pgfpathlineto{\pgfqpoint{5.994392in}{1.203701in}}%
\pgfusepath{stroke}%
\end{pgfscope}%
\begin{pgfscope}%
\definecolor{textcolor}{rgb}{0.000000,0.000000,0.000000}%
\pgfsetstrokecolor{textcolor}%
\pgfsetfillcolor{textcolor}%
\pgftext[x=6.043281in,y=1.182312in,left,base]{\color{textcolor}{\rmfamily\fontsize{4.400000}{5.280000}\selectfont\catcode`\^=\active\def^{\ifmmode\sp\else\^{}\fi}\catcode`\%=\active\def%{\%}$n=512$}}%
\end{pgfscope}%
\begin{pgfscope}%
\pgfsetrectcap%
\pgfsetroundjoin%
\pgfsetlinewidth{1.154312pt}%
\definecolor{currentstroke}{rgb}{0.478431,0.243137,0.615686}%
\pgfsetstrokecolor{currentstroke}%
\pgfsetdash{}{0pt}%
\pgfpathmoveto{\pgfqpoint{5.914947in}{1.133559in}}%
\pgfpathlineto{\pgfqpoint{5.954669in}{1.133559in}}%
\pgfpathlineto{\pgfqpoint{5.994392in}{1.133559in}}%
\pgfusepath{stroke}%
\end{pgfscope}%
\begin{pgfscope}%
\definecolor{textcolor}{rgb}{0.000000,0.000000,0.000000}%
\pgfsetstrokecolor{textcolor}%
\pgfsetfillcolor{textcolor}%
\pgftext[x=6.043281in,y=1.112170in,left,base]{\color{textcolor}{\rmfamily\fontsize{4.400000}{5.280000}\selectfont\catcode`\^=\active\def^{\ifmmode\sp\else\^{}\fi}\catcode`\%=\active\def%{\%}$n=2048$}}%
\end{pgfscope}%
\end{pgfpicture}%
\makeatother%
\endgroup%
}
    \captionsetup{font=footnotesize,skip=2pt}
    \caption{
      Numerical convergence for quadratic loss with $d_u=16$, $d_v=9$, a random teacher, and whitened input.
      \textbf{(a)} Bures flow from $\Sigma_0=\alpha\Id_m$ at several white scales. Dashed lines show the spectral-margin rate $4\alpha\kappa$ from Theorem~\ref{thm:balanced-convergence}; dotted lines show the coarse rate $4\kappa c_{\rm coarse}$ from the determinant comparison in Appendix~\ref{app:gf} (Remark~\ref{rem:coarse-rate}), using the exact initial sublevel radius $M=2\norm{\Svu_\star}_{\F}$.
      \textbf{(b)} Rate deterioration when one matched diagonal pair of $\Sigma_0$ is scaled by $r$: $r=1$ is the white reference and $r=0$ is rank-deficient.
      \textbf{(c)} Bures flow from the reference covariance $\tSigma_0=\Id_m$ compared with finite-width flows initialized by i.i.d. samples from $\mathcal N(0,\tSigma_0)$, for $n=32,128,512,2048$ and five independent trials.
    }
    \label{fig:bures-convergence-triptych}
    \label{fig:bures-scale-sweep}
    \label{fig:perturb-scale-sweep}
    \label{fig:bures-finite-width-comparison}
\end{figure}

\paragraph{Convergence under finite width sampling}
Theorem~\ref{thm:gf} applies generally to neuron law covariances with positive support gap.
We now focus on finite-width networks, whose initial empirical covariance is $\Sigma_0:=W_0W_0^\top/n$.
For random initialization, the goal is to make this
empirical covariance close to a prescribed population covariance
$\tSigma_0$, thereby retaining its support gap and convergence rate.
Draw the columns $w_i:=(u_i,v_i)$ of $W_0:=\binom{U_0}{V_0}$ independently
from the same \emph{subgaussian} distribution:
$w_i:=\tSigma_0^{1/2}\xi_i$, where $\mathbb E[\xi_i\xi_i^\top]=\Id_m$ and $\sup_{\norm z=1}\norm{\langle z,\xi_i\rangle}_{\psi_2}\le K$.
Here $\norm X_{\psi_2}:=\inf\{c>0:\mathbb E e^{X^2/c^2}\le2\}$ denotes the standard subgaussian norm \citep{Vershynin2018};
standard Gaussian sampling has dimension-independent $K$.
Thus $\mathbb E[w_iw_i^\top]=\tSigma_0$, whereas $\Sigma_0=W_0W_0^\top/n$ is the empirical covariance of finitely many sampled neurons.
The loss gap along the sampled trajectory remains $\Delta_t=\ell(\Svu_t)-\ell_\star$ and we write $\widetilde\Delta_0:=\ell(\tSvu_0)-\ell_\star$ for its population reference.

\begin{corollary}[Finite width preserves the population rate]
\label{cor:finite-width}\label{thm:finite-width}
Under Assumption~\ref{ass:loss}, let $\tSigma_0\succeq0$ with
$\gap{\tSigma_0}>0$, set $r_0:=\rankop\tSigma_0$, and fix
$\varepsilon,\delta\in(0,1)$.
There is a universal $C>0$ such that, if
\begin{equation}\label{eq:width}
 n\ge CK^4\varepsilon^{-2}\bigl(r_0+\log(2/\delta)\bigr),
\end{equation}
then with probability at least $1-\delta$, it holds that 
\((1-\varepsilon)\tSigma_0\preceq\Sigma_0 \preceq(1+\varepsilon)\tSigma_0\).
On this event,
$\gap{\Sigma_0}\ge(1-\varepsilon)\gap{\tSigma_0}$,
the empirical factor flow \eqref{eq:mean-field-gradient-flow} is global, and
\(\Delta_t\le\Delta_0e^{-2\kappa(1-\varepsilon)\gap{\tSigma_0}t}\)
for all $t\ge0$.
The factors $(U_t, V_t)$ converge to a limit with globally optimal predictor $\frac{1}{n} V_\infty U_\infty^\top = \Svu_\infty$.
\end{corollary}

\begin{comment}
\emph{Proof sketch.}
Concentration is needed only on the $r_0$-dimensional range of $\tSigma_0$.
Expressed in an orthonormal basis of this range, the projections of the
whitened samples $\xi_i$ have second-moment matrix $\Id_{r_0}$ and
subgaussian bound at most $K$.
The resulting relative covariance
bound gives $\Sigma_0\succeq(1-\varepsilon)\tSigma_0$, hence
$a_\pm(\Sigma_0)\ge(1-\varepsilon)a_\pm(\tSigma_0)$ and
$\gap{\Sigma_0}\ge(1-\varepsilon)\gap{\tSigma_0}>0$.
This positive empirical support gap supplies the hypothesis needed to apply Theorem~\ref{thm:gf}.
Appendix~\ref{app:sampling} gives the full argument.\qed
\end{comment}

\begin{remark}[Width and initialization geometry]\label{rem:rank-adaptive-width}
The condition $\gap{\tSigma_0}>0$ forces $r_0\ge\min(d_u,d_v)$,
and this minimum is attainable. If $d_u\le d_v$, take
$\tSigma_0=\diag(2\sigma^2\Id_{d_u},0_{d_v})$, $\sigma>0$:
then $r_0=d_u$ and $\gap{\tSigma_0}=2\sigma^2$.
With Gaussian sampling, Corollary~\ref{cor:finite-width} gives width proportional to
$\varepsilon^{-2}(d_u+\log(2/\delta))$ and rate $4(1-\varepsilon)\kappa\sigma^2$, the same as for full-rank
population covariance $\sigma^2\Id_m$. Here every $v_i$ starts at zero,
so $\Svu_0=0$ and $\Delta_0=\widetilde\Delta_0$ exactly.
Exchanging the blocks gives dependence on $\min(d_u,d_v)$ in general.
This improvement changes the initialization: for $\tSigma_0=\sigma^2\Id_m$,
$r_0=m$. Its empirical support gap can vanish at small width even though its
population support gap is positive (Example~\ref{ex:white-small-width}).
\end{remark}

For the sampling scheme in Corollary~\ref{cor:finite-width}, the prefactor is
the empirical initial loss gap $\Delta_0$, which need not equal the population
loss gap $\widetilde\Delta_0$. Controlling it is a separate
accuracy requirement: for isotropic Gaussian sampling and the standard quadratic loss,
the expected excess initial error is $\sigma^4d_ud_v/(2n)$
(Proposition~\ref{prop:gaussian}). Theorem~\ref{thm:sampling-budgets}
gives a separate accuracy budget and a deterministic loss envelope.
For population variance $\sigma^2 \Id_m$ and $\varepsilon=1/2$,
the certified exponent is $2\kappa\sigma^2$.

%%% Local Variables:
%%% mode: LaTeX
%%% TeX-master: "../main"
%%% End:
